\documentclass[12pt]{report}
\usepackage{latexsym}
\usepackage{amssymb}
\usepackage{amsmath}
\usepackage{mathrsfs}
\usepackage{graphicx}
\usepackage{pstricks} 
%\usepackage{txfonts}
\usepackage{verbatim}
\usepackage[all]{xy}
\usepackage{hyperref}
\usepackage{longtable}
%\usepackage[pdftex]{graphicx}
%\usepackage[pdftex]{color}

\newtheorem{theo}{Theorem}[section]
\newenvironment{example}{ \smallskip \addtocounter{theo}{1} \noindent \textbf{Example  \arabic{section}.\arabic{theo}~}}{}%[section]
\newtheorem{lemma}[theo]{Lemma}
\newtheorem{sublemma}[theo]{\indent Lemma}
\newtheorem{corollary}[theo]{Corollary}
\newtheorem{prop}[theo]{Proposition}
\newtheorem{exercise}[theo]{Exercise}
\newtheorem{eduhypo}[theo]{Educational Hypothesis}
\newtheorem{hypo}[theo]{Hypothesis}
\newtheorem{stat}[theo]{Statement}
\newtheorem{conjecture}[theo]{Conjecture}
\newenvironment{proof}[1][]{ \textbf{Proof#1. }}{$\Box$\medskip}
\newtheorem{defi}[theo]{Definition}
\newtheorem{notation}[theo]{Notation}

\newcommand {\eqdef}{:=}% {\stackrel {def}{=}}
\newcommand{\Stab}{\mathrm{Stab}}
\newcommand{\Orb}{\mathrm{Orb}}

\newcommand{\WW}{\mathbf{W}}
\newcommand {\ZZ} {\mathbb {Z}}
\newcommand {\QQ} {\mathbb {Q}}
\newcommand {\BB} {\mathbf {B}}
\newcommand {\CC} {\mathbb {C}}
\newcommand {\CCC} {\mathfrak {C}}
\newcommand {\HC} {\mathfrak {HC}}
\newcommand {\RR} {\mathbb {R}}
\newcommand {\NN} {\mathbb {N}}
\newcommand {\KK} {\mathbb {K}}
\newcommand {\MM} {\mathcal {M}}
\newcommand {\Div} {\mathrm {Div}}

\newcommand{\sS}{\mathfrak{s}}
%\newcommand{\cc}{\mathfrak{c}}
\newcommand{\tT}{\mathfrak{t}}
\newcommand{\cC}{\mathfrak{c}}
\newcommand{\kk}{\mathfrak{k}}
\newcommand{\id}{\mathrm{id}}
\newcommand{\nn}{\mathfrak{n}}
\newcommand{\mm}{\mathfrak{m}}
\newcommand{\aaa}{\mathfrak{a}}
\newcommand{\hh}{\mathfrak{h}}
\newcommand{\SSS}{\mathbf{\mathfrak{S}}}
\newcommand{\qq}{\mathfrak{q}}
\newcommand{\pp}{\mathfrak{p}}
\newcommand{\gG}{\mathfrak{g}}
\newcommand{\bb}{\mathfrak{b}}
\newcommand{\dd}{\mathfrak{d}}
\newcommand{\LL}{\mathcal{L}}
\newcommand{\DD}{\mathcal{D}}
\newcommand{\lL}{\mathfrak{l}}
\newcommand{\binomial}[2]{\left(\begin{array}{c}#1 \\ #2\end{array} \right)}

\newcommand{\AAA}{\mathscr{A}}

\newcommand {\bphi}{\phi}
\renewcommand {\phi} {\varphi}
\newcommand{\ann}{\mathrm{Ann}}

\newcommand{\sign}{\mathrm{sign}}
\newcommand{\bigtab}{\hspace{2cm}}
\newcommand{\smalltab}{\hspace{1cm}}
\newcommand{\half}{\frac{1}{2}}
\newcommand{\doublebrace}[4]{\left\{\begin{array}{ll}#1 & #2 \\ #3 & #4\end{array}\right.}
\newcommand{\refth}[1]{Theorem \ref{#1}}
\newcommand{\refle}[1]{Lemma \ref{#1}}
\newcommand{\refsec}[1]{Section \ref{#1}}
\newcommand{\refcor}[1]{Corollary \ref{#1}}
\newcommand{\refprop}[1]{Proposition \ref{#1}}
\newcommand{\refcon}[1]{Conjecture \ref{#1}}
\newcommand{\refnot}[1]{Notation \ref{#1}}
\newcommand{\refeq}[1]{(\ref{#1})}
\newcommand{\refdef}[1]{Definition \ref{#1}}
\newcommand{\stackunder}[2]{\begin{array}{l}#1 \\ #2\end{array}}
\newcommand{\supp}{\mathrm{index\_\_supp~}}
\newcommand{\tr}{\mathrm{Tr}}
\renewcommand{\Im}{\mathrm{Im}}
\newcommand{\im}{\mathrm{Im}}
\newcommand{\sL}{\mathrm{sl}}
\newcommand{\oalg}{\mathrm{o}}
\newcommand{\mult}{\mathrm{mult}}
\newcommand{\Jac}{\mathrm{Jac}}
\newcommand{\Res}{\mathrm{Res}}
\newcommand{\sP}{\mathrm{sp}}
\newcommand{\linspan}{\mathrm{span}}
\newcommand{\gl}{\mathrm{gl}}
\newcommand{\so}{\mathrm{so}}
\newcommand{\ad}{\mathrm{ad~}}
\newcommand{\Cone}{\mathrm{Cone}}
\newcommand{\weights}{\mathrm{Weights}}
\newcommand{\hw}{\mathrm{Sing}}
\newcommand{\ch}{\mathrm{ch}}


%\newenvironment {proof} {$\mathbf{Proof}$ \newline} {$\Box$}
\newcommand {\ball} [2]{B_{#2}(#1)}
\newcommand {\sPh} [2]{S_{#2}(#1)}
\newcommand {\ints}[2]{\begin{array}{c} #2 \\ \int \\ #1 \end{array}}
\newcommand {\dint}{\int\!\!\!\int}
\newcommand {\rntwo} {\left( \frac{1}{r^{n-2}} \right)}
\newcommand {\dvol} {dx_1 \land \ldots \land dx_n}
\newcommand {\del} {\partial}
\newcommand{\eps}{\varepsilon}
\newcommand{\veps}{\epsilon}
%\addtolength{\hoffset}{-3.5cm}
%\addtolength{\textwidth}{6.8cm}
%\addtolength{\voffset}{-3.3cm}
%\addtolength{\textheight}{6.3cm}
\addtolength{\hoffset}{-0.5cm}\addtolength{\textwidth}{1cm}\addtolength{\voffset}{-0.6cm}\addtolength{\textheight}{1.2cm}

\newcommand{\ind}{\mathrm{ind}}
\newcommand{\extendMargins}{
\addtolength{\hoffset}{-4.5cm}
\addtolength{\textwidth}{7cm}}
\newcommand{\shrinkMargins}{
\addtolength{\hoffset}{4.5cm}
\addtolength{\textwidth}{-7cm}}

\newcommand{\rk}{\mathrm{rk}}
\def\cplus{\hbox{$\subset${\raise0.3ex\hbox{\kern -0.55em ${\scriptscriptstyle +}$}}\ }}

\def\clplus{\hbox{$\subset${\raise0.3ex\hbox{\kern -0.55em ${\scriptscriptstyle +}$}}\ }}
\def\crplus{\hbox{$\supset${\raise1.05pt\hbox{\kern -0.60em ${\scriptscriptstyle +}$}}\ }}
\def\nperp{{\hbox{$\perp${\raise1.05pt\hbox{\kern -0.55em $/$}}\ }}}
\def\sperp{{\hbox{$\perp${\hbox{\kern -.77em $-$}\raise-1pt\hbox{\kern -.77em $-$}}\ }}}
\def\nsimeq{\hbox{$\simeq${\raise1.05pt\hbox{\kern -0.55em $/$}}\ }}
\newcommand{\pma}{\pm}%{}^{+}_-}

\newcommand{\temp}{\tilde{a}}
\newcommand{\tempA}{|t\in [t_0,1] }

\newcommand{\DAthree}{$\mathrm{DA}_3$}
\newcommand{\DAtwo}{$\mathrm{DA}_2$}
\newcommand{\BCtwo}{$\mathrm{BC}_2$}
\newcommand{\BAone}{$\mathrm{BA}_1$}
\newcommand{\CAone}{$\mathrm{CA}_1$}
\newcommand{\Hom}{\mathrm{Hom}}
\newcommand{\pmf}[1]{ \pm}%(-1)^{\bar{#1}}}
\newcommand{\arxivVersion}[1]{#1}
\author{Todor Milev}
\title{On Root Fernando-Kac subalgebras of finite type \\~ \\ \Large Jacobs University Bremen \\\normalsize Ph. D. thesis}

\begin{document}
\begin{center}
\Huge{
~\\~\\~\\~\\
\textbf{On root Fernando-Kac subalgebras of finite type }
}

~\\~\\~\\

\Large  {Ph. D. thesis, Todor Milev}
~\\~\\

\Large { Jacobs University Bremen, \today} 

~\\~\\

\Large Advisor: {Professor Dr. Ivan Penkov}

~\\~\\~\\
\end{center}
\Large
Committee members:

\noindent Professor Dr. Peter Fiebig

\noindent Professor Dr. Ivan Penkov

\noindent Professor Dr. Vera Serganova

\normalsize
\thispagestyle{empty}

\newpage
~\\~\\~\\~\\~\\~\\~\\~\\~\\~\\~\\~\\~\\~\\~\\~\\~\\~\\~\\~\\~\\~\\~\\~\\~\\~\\~\\~\\~\\~\\~\\~\\~\\~\\

\thispagestyle{empty}


To my love, Que-Tien Tran
\newpage
\tableofcontents
\chapter{Introduction and Notation}
\section{Introduction}

Let $\gG$ be a finite-dimensional complex Lie algebra and $M$ a $\gG$-module. In the 1980's, S. Fernando and V. Kac independently proved that the set of elements $\gG[M]$ of $\gG$ having locally finite action on $M$ is in fact a Lie subalgebra, which we call \emph{Fernando-Kac subalgebra associated to $M$} (\cite{Kac:FernandoKacSubalgebra},  \cite{Fernando1}).

Let $\lL$ be an arbitrary Lie subalgebra of $\gG$. A \emph{$(\gG,\lL)$-module} is a module $M$ for which $\lL$ is a subset of $\gG[M]$. If the Jordan-H\"older multiplicities of any finite-dimensional irreducible $\lL$-module inside $M$ are finite, we say that $M$ is a \emph{$(\gG,\lL)$-module of finite type}. If $\lL=\gG[M]$ we say that $M$ is a \emph{strict $(\gG,\lL)$-module}. We say that a Lie subalgebra $\lL$ of $\gG$ is \emph{a Fernando-Kac subalgebra of $\gG$} if there exists an irreducible strict $(\gG,\lL)$-module $M$. Furthermore, we say that $\lL$ is a Fernando-Kac subalgebra of $\gG$ of \emph{finite type} if $M$ can be chosen of finite type. Otherwise $\lL$ is a Fernando-Kac subalgebra of \emph{infinite type}.

In 2004, I. Penkov, V. Serganova and G. Zuckerman gave a construction of irreducible $(\gG,\kk)$-modules of finite type, where $\kk$ is a reductive in $\gG$ subalgebra. For $\gG=\sL(n)$, this construction remarkably turned out to realize all possible root Fernando-Kac subalgebras $\lL=\kk\crplus\nn$ of finite type. Under a root subalgebra we understand a subalgebra which contains a Cartan subalgebra of $\gG$. In 2007, I. Penkov conjectured a necessary and sufficient condition for a root subalgebra $\lL$ to be Fernando-Kac of finite type. More precisely, his conjecture was that a combination of the already established in \cite{PSZ} combinatorial criterion for $\gG\simeq\sL(n)$ (cone condition) and a second criterion related to results of S. Fernando, \cite{Fernando1} (centralizer condition) is necessary and sufficient for $\lL$ to be Fernando-Kac subalgebra of finite type.

The present study proves this conjecture for $\so(2n)$, $\so(2n+1)$ and $\sP(2n)$, and for the exceptional simple Lie algebras except $ E_8$. We note that the result for simple Lie algebras carries over directly to reductive Lie algebras. %Therefore this study proves Penkov's conjecture for all reductive Lie algebras that have no simple component of type $E_8$.

It is important to note that the finite type of a module $M$ over $\gG[M]$ plays a central role in this study, as the problem of classifying root Fernando-Kac subalgebras of possibly infinite type has already been solved. In fact, in \cite{PS:GenHarishChandra}, V. Serganova and I. Penkov prove that all root subalgebras are Fernando-Kac. The purpose of the present study is to decide which root subalgebras are Fernando-Kac subalgebras of finite type, and which are Fernando-Kac subalgebras of infinite type.

The proof of Penkov's conjecture has three essential steps. The first step is to show that for a prescribed root subalgebra $\lL$, the failure of the cone condition implies that  $\lL$ is a Fernando-Kac subalgebra of infinite type. This is the main contribution of the present thesis, presented in chapters \ref{chapterConeConditionFails-NoFKFT}, \ref{secComputerProof} and the appendix. The second step is to verify that the failure of the centralizer condition implies that  $\lL$ is Fernando-Kac subalgebra of infinite type. This is a direct consequence of the deep result of S. Fernando \cite{Fernando1}. If $\hh$ is a Cartan subalgebra of $\gG$, S. Fernando's result states that a simple Lie algebra $\gG$ admits irreducible strict $(\gG,\hh)$-modules of finite type only if $\gG$ is isomorphic to $\sL(n)$ or $\sP(2n)$.

The final step in the proof is to show that the cone and centralizer conditions imply the existence of an irreducible strict $(\gG,\lL)$-module of finite type. This is an application of a construction of I. Penkov, V. Serganova and G. Zuckerman given in \cite{PSZ}. In chapter \ref{secTheConstruction} we prove some technical statements needed to relate the cone condition to the latter construction.

Chapter \ref{chapterPreliminaries} presents definitions, examples and statements from the general theory of $(\gG,\lL)$-modules, as well as some preliminary results needed in the proof of Penkov's conjecture. The appendix presents information on root subsystems that could be of interest beyond the study of root Fernando-Kac subalgebras of finite type.

\section{Overview of some results on $(\gG,\kk)$-modules}
Let $\gG$ be a finite-dimensional Lie algebra and $\kk$ be any subalgebra. Classically, $(\gG,\kk)$-modules have been studied for reductive Lie algebras $\gG$ and very special subalgebras $\kk$. The best known case is when $\kk$ coincides with the fixed points of an involution (Harish-Chandra modules). This is the celebrated classification of Harish-Chandra modules by  A. Beilinson - J.Bernstein \cite{BeilinsonBerstein}, A. Knapp - G. Zuckerman \cite{KnappZuckerman}, R. Langlands \cite{Langlands:IrredReps}, Vogan \cite{Vogan:SingularUnitary,Vogan:RepsRealRed}. Another type of well-known $(\gG,\kk)$-modules are weight modules, where $\kk=\hh$ is a Cartan subalgebra. Weight modules of finite type have been studied in the last 30 years, and play a key role in the present study. One of the important results in this area was given by the work of D. Britten and F. Lemire, \cite{BrittenLemire}. Further important results in the area followed in papers of G. Benkart, D. Britten, S. Fernando, V. Futorny, F. Lemire and others. In 2000, O. Mathieu finished the classification of $(\gG,\hh)$-modules of finite type, \cite{Mathieu}.

\nocite{BenkartBrittenLemire} \nocite{Futorny:thesis}

The present thesis uses two key results in the theory of $(\gG,\hh)$-modules of finite type: the existence of strict $(\sL(n),\hh)$- and  $(\sP(2n),\hh)$- modules of finite type (given for example in \cite{BrittenLemire}), and the fact that strict $(\gG,\hh)$-modules of finite type do not exist for a simple Lie algebra $\gG\nsimeq \sL(n), \sP(2n)$ (S. Fernando, \cite{Fernando1}). The latter result in turn builds upon Duflo's theorem (see \cite[Chapter 8]{Dixmier}), as well as results of  W. Borho, A. Joseph, O. Gabber, H. Kraft and others (\cite{BorhoKraft}, \cite{Gabber:IntegrabilityCharVar}, \cite{Joseph:MinDimNilpotentOrbit}, \cite{Joseph:MinimalRealizations}).

In the late 1990's, I. Penkov, V. Serganova and G. Zuckerman proposed a program for the study and possible classification of all irreducible $(\gG,\kk)$-modules of finite type for a reductive Lie algebra $\gG$ and a reductive in $\gG$ subalgebra $\kk$. In the recent past, they have made a substantial progress in both $(\gG,\kk)$-modules classification and Fernando-Kac subalgebra classification. In the case when $\kk$ is reductive in $\gG$, \cite{PZ} obtained classification results for $(\gG,\kk)$-modules of finite type under additional assumptions for the modules. More precisely, \cite{PZ} described all such modules with ``generic'' minimal (with respect to the Vogan norm) $\kk$-type. When the minimal $\kk$-type of the irreducible $(\gG,\kk)$-module of finite type is not generic, no general results are available beyond the case of Harish-Chandra modules and the case $\kk=\hh$. This is an active direction of current joint work of I. Penkov and G. Zuckerman.

Under the additional assumption that the irreducible $(\gG,\kk)$-module $M$ is bounded (i.e. the multiplicities of all $\kk$-types in $M$ are bounded by the same constant), in \cite{PS:Bounded}, I. Penkov and V. Serganova established combinatorial bounds on the possibilities for $\gG$ and $\kk$. They showed that if $\kk\simeq \sL(2)$, there are only 5 possible semisimple  Lie algebras $\gG$ admitting an infinite-dimensional bounded $(\gG,\sL(2))$-module: $\gG=\sL(2),\sL(3),\sL(2)\oplus\sL(2),\sP(4)$. I. Penkov and V. Serganova also restricted greatly the candidates for pairs $(\gG,\kk)$ for which infinite-dimensional bounded $\kk$-modules exist by proving an inequality relating $\gG$ and $\kk$. Recently, A. Petukhov described explicitly all bounded reductive subalgebras of $\gG=\sL(n)$, \cite{Petukhov:boundedRedSAssl(n)}.


In 2004, I. Penkov, V. Serganova and G. Zuckerman obtained strong results towards the classification of Fernando-Kac subalgebras of finite type. More precisely, they gave a general necessary and sufficient condition for a reductive in $\gG$ subalgebra to be the reductive part of a Fernando-Kac subalgebra of finite type - in particular they showed that such a reductive part always exists. They further proved that for a reductive subalgebra $\kk$ to be the reductive part of a Fernando-Kac subalgebra of finite type, it is necessary that the center of $\kk$ coincide with its centralizer. Under the assumption that $\kk$ is a root subalgebra, both of these statements are trivially satisfied. However, for a general $\kk$, these statements are a first important step towards a possible future classification of all Fernando-Kac subalgebras of finite type.

%As a first step in this direction, \cite{PSZ} classifies all root Fernando-Kac subalgebras of finite type of $\sL(n)$. The next logical step was then to the generalize this classification to the remaining simple Lie algebras, which is the topic of the present Ph.D. thesis. Such a classification was conjectured in 2007 by I. Penkov, and is proved in the present thesis for $\gG\nsimeq E_8$. 

\section{Notation}
The base field is $\CC$ unless stated otherwise and all algebras and Lie algebras are defined over $\CC$. Unless specified otherwise, the symbol $\gG$ denotes a finite-dimensional reductive Lie algebra (``reductive'' means that if an ideal of $\gG$ is solvable then this ideal must consist of central elements). $[\bullet,\bullet]$ denotes the Lie bracket, and $\ad x$ stands for the adjoint action of $x$: $\ad x (y):=[x,y]$. We denote by  $U(\gG)$ the universal enveloping algebra of $\gG$. For a subalgebra $\sS$ of $\gG$, $Z(\sS)$ denotes the center of $\sS$, $N(\sS)$ - the normalizer of $\sS$ in $\gG$ (i.e. $N(\sS)=\{g\in\gG~|~[g,\sS]\subset\sS\}$), and $C(\sS)$ - the centralizer of $\sS$ in $\gG$ (i.e. $C(\sS)=\{g\in\gG~|~[g,\sS]=0\}$). 

We denote the action of a $\CC$-algebra on a module by $\cdot$~~. The notation $\langle a\rangle\cdot m$ stands for $\linspan_\CC\{m,a\cdot m, a^2\cdot m,\dots\}$.

Let $\sS$ be a Lie algebra and $\kk,\nn\subset\sS$ be Lie subalgebras such that $\kk\cap\nn=\{0\}$, $\nn$ is an ideal in $\sS$, and $\sS$ is the direct sum as a vector space of $\kk$ and $\nn$. Then $\sS$ is the \emph{semi-direct} sum of $\kk$ and $\nn$ and we write $\sS=\kk\crplus\nn$. The sign $\crplus$ is rounded towards the ideal (if we write $\kk\oplus\nn$ instead then both $\kk$ and $\nn$ are ideals).

A subalgebra $\kk$ is \textit{reductive in} $\gG$ if $\gG$ is a semi-simple $\kk$-module under the adjoint action of $\kk$. When $\gG$ is reductive, the \textit{semisimple part} $\gG_{ss}$ of $\gG$ is canonically defined and equals $[\gG,\gG]$. We will use the notation $\gG_{ss}$ only for reductive Lie algebras $\gG$.


%\textbf{To be adjusted:} The standard filtration of $U(\gG)$, given by total degree, is denoted by $\CC\subset U(\gG)_1\subset U(\gG)_2\subset\dots=U(\gG)$. By $\bar\gG$ we denote the image of the embedding of $\gG$ in $U(\gG)$. In general we identify $\gG$ with $\bar\gG$ without explicitly saying so. However, to avoid confusion, when we consider an element $a\in\gG$ as an element of $\bar\gG$, we write $\bar a$. $C_{U(\gG)}(\bar\kk)$ denotes the centralizer of $\bar\kk\subset\bar\gG$ in $U(\gG)$ (which is of course larger than the centralizer of $\bar \kk$ in $\bar\gG$).

Let $\lL\subset\gG$ be a Lie subalgebra.  Let $\mm$ be the unique maximal ideal of $\lL$ consisting of elements that have nilpotent adjoint action on $\gG$, and set $\nn:=\mm\cap[\gG,\gG]$. If $\nn$ admits a complement subalgebra $\lL_{red}$ which is reductive in $\gG$, we call $\lL_{red}$ a \textit{reductive in $\gG$ part} of $\lL$. We have $\lL=\lL_{red}\crplus\nn$.

In what follows we fix a Cartan subalgebra $\hh$ of $\gG$. Under a root subalgebra we understand a subalgebra of $\gG$ containing $\hh$. By $\lL$ we denote a variable root subalgebra of $\gG$ with nilradical $\nn$. The unique reductive part of $\lL$ which contains $\hh$ is denoted by $\kk$. The set of $\hh$-roots of $\lL$ is denoted by $\Delta(\lL)$ (each element of $\Delta(\lL)$ is automatically a root of $\gG$) . We also put $\Delta(\nn):=\Delta(\lL)\backslash \Delta(\kk)$. There are vector space decompositions 
\[
\begin{array}{cc}
\displaystyle\gG=\hh\oplus\bigoplus_{\alpha\in\Delta(\gG)}\gG^\alpha,& \displaystyle\lL=\hh\oplus\bigoplus_{\alpha \in \Delta(\lL)} \gG^{\alpha},\\
\displaystyle \kk=\hh\oplus\bigoplus_{\substack{ \alpha\in\Delta(\lL) :\\-\alpha \in\Delta(\lL)}}\gG^\alpha,&\displaystyle\nn=\bigoplus_{ \substack{\alpha\in \Delta(\lL) :\\ -\alpha \notin \Delta(\lL)}}\gG^{\alpha}\\
\end{array}.
\]
We fix a Borel subalgebra $\bb\supset\hh$ whose roots are by definition the \emph{positive} roots; we denote them by $\Delta^+(\gG)$. We define the element $\rho\in\hh^*$ to be the half-sum of the positive roots. Given a set of roots $I$, we denote by $\Cone_{\ZZ}(I)$ (respectively, $\Cone_{\QQ}(I)$) the $\ZZ_{\geq 0}$-span (respectively, $\QQ_{\geq 0}$-span) of $I$.

The form on $\hh^*$ induced by the Killing form is denoted by $\langle\bullet,\bullet\rangle$. The sign $\sperp$ stands for \emph{strongly orthogonal}; two roots $\alpha,\beta$ are defined to be strongly orthogonal if neither $\alpha+\beta$ nor $\alpha-\beta$ is a root or zero (which implies $\langle\alpha,\beta\rangle=0$).  We say that a root $\alpha$ is \emph{linked} to an arbitrary set of roots $I$ if there is an element of $I$ that is not orthogonal to $\alpha$. The Weyl group of $\gG$ is denoted by $\WW$. By $s_\alpha\in\WW$ we denote the reflection with respect to root $\alpha$. For two roots $\alpha,\beta\in\Delta(\gG)$ we say that $\alpha\preceq\beta$ if $\beta-\alpha$ is a non-negative linear combination of positive roots. 



We fix the conventional expressions for the positive roots of the classical root systems:
\begin{eqnarray*}
A_n, n\geq 2&:& \Delta^+(\gG)=\{\eps_i-\eps_j|i< j\in\{1,\dots,n+1\}\};\\
B_n, n\geq 2&:& \Delta^+(\gG)=\{\eps_i\pma\eps_j|i< j\in\{1,\dots,n\}\}\cup\{\eps_i|i\in\{1,\dots n\}\};\\
C_n, n\geq 2&:& \Delta^+(\gG)=\{\eps_i\pma\eps_j|i\leq j\in\{1,\dots,n\}\}\backslash\{0\};\\
D_n, n\geq 4&:& \Delta^+(\gG)=\{\eps_i\pma\eps_j|i< j\in\{1,\dots,n\}\}.
\end{eqnarray*}

It is possible to choose a representative $g^\alpha$ from each root space $\gG^\alpha$ such that there exist integers $n_{\alpha\beta}$ with $[g^\alpha,g^\beta]=n_{\alpha\beta}g^{\alpha+\beta}$ whenever $\alpha+\beta$ is a root and with $[g^\alpha,g^{-\alpha}] =\frac{2h_{\alpha}} {\langle \alpha,\alpha\rangle}$, where $h_\alpha$ is the element for which $\langle\alpha,\beta\rangle=\beta(h_\alpha)$ for all $\beta$. We assume one such choice to be fixed; this choice is in general not unique and is known as Chevalley basis. A detailed discussion of the subject can be found in \cite[\S 2.8, 2.9]{Samelson}. The numbers $n_{\alpha\beta}$ can be chosen to be equal to $\pm$ the number $(\max \{t|\beta-t\alpha\in\Delta(\gG)\})+1$.

\chapter{Preliminaries on $(\gG,\kk)$-modules}\label{chapterPreliminaries}
\section{Fernando-Kac subalgebras} \label{secOrigin}
\begin{defi}
Let $\kk$ be a finite-dimensional Lie algebra. Then a $\kk$-module $M$ is a locally finite (integrable) $\kk$-module if $\kk$ acts locally finitely on $M$, i.e. the $\kk$ submodule generated by an arbitrary vector $m\in M$ is finite dimensional. In other words, $M$ is locally finite if
\[
\dim_\CC U(\kk)\cdot m<\infty \quad \mathrm{for~all~}m\in M,
\]
where  $U(\kk)$ is the universal enveloping algebra of $\kk$.
\end{defi}
\begin{defi}\label{defg[M]}
Let $\gG$ be a Lie algebra, and $M$ a $\gG$-module. We define $\gG[M]\subset\gG$ to be the set of all elements of $\gG$ which act locally finitely on $M$, i.e.
\[
\gG[M]:=\{g\in\gG~|~\dim_\CC (\langle g\rangle \cdot m)<\infty ~\forall m\in M\}.
\]
\end{defi}
The set $\gG[M]$ turns out to be a Lie algebra itself. This fact was noted by Bertram Kostant in the 1960's and informally communicated to a few of his colleagues. The first published proof over $\CC$ is due to V. Kac (\cite{Kac:FernandoKacSubalgebra}); the statement was proved independently a few years later by S. Fernando over an arbitrary algebraically closed field of zero characteristic. Kac's proof, presented below, uses an elementary observation but is valid only over a field that is a complete metric space - for example $\CC$ or $\RR$.
\begin{theo}\label{thFernando-Kac}
Let $\gG$ be a finite-dimensional Lie algebra and $M$ a $\gG$-module over the field $\KK$= $\RR$ or $\CC$. Then $\gG[M]$ is a Lie subalgebra of $\gG$. 
\end{theo}
\begin{proof} (Kac). 
Let $z_1,\dots,z_n\in\gG[M]$ be a basis of $\linspan_{\KK}\gG[M]$, the linear span of $\gG[M]$.

We claim that for any $m\in M$ and any $t\in\KK$ the equality $e^{t\ad x}y\cdot m=e^{x}ye^{-x}\cdot m$ holds. Indeed, in $U(\gG)$ we have that $\ad x=L_x-R_x$, where $L_x$ stands for left multiplication by $x$ and $R_x$ - for right multiplication by $x$. The linear operators $L_x$ and $R_x$ commute. Therefore $e^{\ad x}y\cdot m=\sum_{n=0}^\infty \frac{(L_x-R_x)^n}{n!}y\cdot m= \sum_{n,k} \frac{(-1)^{n-k} L_x^k R_x^{n-k}} {k!(n-k)!} y\cdot m=\sum_{k}\frac{L_x^k}{k!}\sum_{l} \frac{(-1)^lR_x^l}{l!}y\cdot m=e^xye^{-x}\cdot m$. Since $tx,ty\in\gG[M]$ for any $t\in\KK$, we have more generally $e^{t\ad x}y\cdot m=e^{tx}ye^{-tx}\cdot m$. Now $(e^{tx}ye^{-tx})^n=e^{tx}y^ne^{-tx}$ implies that $e^{t\ad x}y\in \gG[M]$. 

Consider the element
\[
F(t):=\frac{e^{t\ad x} y-y}{t} \in\gG\quad ,
\]
and set $A:=\linspan \left( \{y\}\cup\bigcup_{t\in\mathbb{K}} \{F(t)\}\right)$. The space $A$ is finite dimensional (being a subspace of $\gG$) hence it is spanned by $y$ and finitely many elements of the form $e^{t\ad x}y$. Therefore $A$ has a basis of elements lying in $\gG[M]$. In particular, $[x,y]=\lim_{t\to 0}F(t)\in A$ is a linear combination of $z_1, \dots, z_n$.

The preceding discussion shows that $\linspan_{\KK}\gG[M]$ is a Lie subalgebra. It remains to show that $\gG[M]=\linspan_{\KK}\gG[M]$. This is equivalent to showing that $a_1z_1+\dots a_nz_n$ acts locally finitely for all $a_i\in\KK$. By the Poincare-Birkhoff-Witt theorem, for any $N$, one can express an arbitrary element $(a_1z_1 +\dots+ a_nz_n)^N$ as a linear combination of monomials in the form $z_1^{l_1}\dots z_n^{l_n} $, where $l_1+\dots+l_n\leq N$. This makes it clear that  $(a_1z_1+\dots+ a_nz_n)$ acts locally finitely on $M$. Hence $\linspan_{\KK}\gG[M]=\gG[M] $.
\end{proof}
\begin{defi}
Let $\kk$ be a subalgebra of $\gG$. We say that $\gG[M]$ is the Fernando-Kac subalgebra associated to the $\gG$-module $M$. Further, we define a subalgebra $\lL\subset\gG$ to be a Fernando-Kac subalgebra $\gG$ if there exists an irreducible $\gG$-module $M$ with $\gG[M]=\lL$.
\end{defi}
\begin{defi}
Let $\kk$ be a subalgebra of $\gG$. We say that the $\gG$-module $M$ is a $(\gG,\kk)$-module if $\kk\subset\gG[M]$. If $\kk=\gG[M]$ we say that $M$ is a strict $(\gG,\kk)$-module.
\end{defi}

Example \ref{exFK} in section \ref{secGammaS} (taken from \cite{PS:GenHarishChandra}) shows that not all subalgebras of a reductive Lie algebra $\gG$ are Fernando-Kac subalgebras. In fact, the classification of Fernando-Kac subalgebras is not known.

\section{The functor $\Gamma_S$}\label{secGammaS}
In the following section we review some well known facts about the functor of finite vectors.
\begin{defi} Let $\gG$ be a finite-dimensional Lie algebra and $M$ a $\gG$-module. Let $S$ be any subset of $\gG$. Define $\Gamma_S(M)$ as
\begin{equation}\label{eqZuck}
\Gamma_S(M)\eqdef\{m\in M ~|~\dim (\langle s\rangle\cdot m)<\infty\quad \forall s\in S\}.
\end{equation}
\end{defi}
\begin{lemma}\label{propGamma}
$\Gamma_S(M)$ is a $\gG$-submodule.
\end{lemma}
\begin{proof}
Let $s\in S$, $g\in\gG$, $m\in \Gamma_S(M)$. We have to show that $\dim(\langle s\rangle\cdot g\cdot m) < \infty$. The Poincare-Birkhoff-Witt theorem implies $(\langle s\rangle\cdot g\cdot m)\subset \gG\cdot\langle s\rangle \cdot m$. Since $\gG$ is finite dimensional and $m$ is in $\Gamma_S(M)$, we have $\dim(\langle s\rangle\cdot g\cdot m) < \infty$.
\end{proof}
\begin{prop}\label{propExactZuckerman}
Let $\gG$ be a Lie algebra and let $S\subset\gG$.  $\Gamma_S$ is a left exact functor. In other words, if $0\to L\overset{\phi}{\to} M\overset{\psi}{\to} N$ is an exact sequence of $\gG$-modules, then the sequence
\begin{equation}\label{eqExactZuck}
0\to\Gamma_S(L)\overset{\phi_\Gamma}{\to}\Gamma_S(M)\overset{\psi_\Gamma}{\to}\Gamma_S(N),
\end{equation}
where $\phi_\Gamma:=\phi_{|\Gamma_S(L)}$ and $\psi_\Gamma:=\psi_{|\Gamma_S(M)}$, is well defined and is exact.
\end{prop}
\begin{proof}
To prove that $\Gamma_S$ is indeed a functor we must prove that if $L\overset{\phi}{\to}M$ is a module homomorphism, then $\phi$ maps $\Gamma_S(L)$ to $\Gamma_S(M)$. That is clear: for $l\in L$, $\dim (\langle s\rangle \cdot l)<\infty$ implies that $\dim (\langle s\rangle \cdot \phi(l))=\dim \phi(\langle s\rangle \cdot(l))<\infty$.

The proof of the exactness of \refeq{eqExactZuck} can be done by the following straightforward diagram chasing  argument.
\begin{itemize}
\item $\ker \phi_\Gamma=0$: clear as $\phi_\Gamma=\phi_{|\Gamma_S(L)}$ and $\ker \phi=0$.
\item $\im\phi_\Gamma\subset\ker \psi_\Gamma$: clear as $\phi_\Gamma$ and $\psi_\Gamma$ are restrictions of $\phi$ and $\psi$.
\item $\ker \psi_\Gamma\subset\im\phi_\Gamma$: $m\in\ker\psi_\Gamma$. Suppose on the contrary that there exists $l\notin\Gamma_S(L)$ such that $\phi(l)=m$ and $\exists s: \dim(\langle s\rangle \cdot l)=\infty$. Then $\dim\phi(\langle s\rangle \cdot l)<\infty$ $\Rightarrow$ $\phi$ is not injective, contradiction.
\end{itemize}
\end{proof}

A quick way to show that $\Gamma_S$ is not right exact is to show that a finite-dimensional module can be a quotient of an infinite-dimensional module $M$ with $\Gamma_S(M)=\{0\}$. Here is one more example.
 
\begin{example}
Let $\gG:=\CC\del_x$, $M:=\frac{1}{x}\CC[\frac{1}{x}]$, $N:=M/(\frac{1}{x^2}\CC[\frac{1}{x}])\simeq \CC$, $M\overset{\psi}{\to}N\to 0$, where $\psi$ is the projection map. Then $\Gamma_{\{\del_x\}}(M)=\{0\}$, $\Gamma_{\{\del_x\}}(N)=N$ (since $N$ is finite dimensional), and hence $\psi_\Gamma$ cannot be surjective.
\end{example}

Although we will not use the right derived functors of $\Gamma_S$ in this thesis, we would like to mention that the higher right derived functors of $\Gamma_S$ are called Zuckerman functors. These functors play a key role in the recent work on $(\gG,\kk)$-modules of I. Penkov and G. Zuckerman, \cite{PZ}, \cite{PZ:ArbitraryMinKtype}.

The following observation will be used throughout this thesis.
\begin{lemma}\label{leActFree}
Let $\gG$ be a finite-dimensional Lie algebra. Let $M$ be an irreducible module and $\lL:=\gG[M]$ be its Fernando-Kac subalgebra. If $x\in\gG\backslash\lL$, then $x$ acts freely on any non-zero vector in $M$. 
\end{lemma}
\begin{proof}
%Let $0\neq v\in M$. Suppose $x\in\gG$ such that  $\langle x\rangle\cdot v$ is finite dimensional. We will prove that, independent of the choice of $ v\neq 0$, we have that $x\in \lL$. Let $w\in M$ be an arbitrary vector. Since $M$ is irreducible and $U(\gG)\cdot w$ is an invariant $\gG$-module, we have that $w=\lambda v+\sum_ig_i\cdot v$ for some $g_i\in U(\gG)$ and some constant $\lambda$. Pick $k$ large enough so that $g_i\in U(\gG)_k$. Applying \refle{lecommute} (b) with $a=x,d=g_i$, we obtain that $\langle x\rangle\cdot (g_i\cdot v) \subset U(\gG)_k\cdot(\langle x\rangle \cdot v)$. Since $\langle x\rangle\cdot v$ and $U(\gG)_k$ are finite dimensional, we obtain that $\langle x\rangle\cdot (g_i\cdot v)$ is finite dimensional, and therefore so is $\langle x\rangle \cdot w$. Hence $x\in \lL$, which proves the Lemma.
According to \refprop{propGamma}, $\Gamma_{\{x\}}(M)$ is a $\gG$-submodule, and by the irreducibility of $M$ either $\Gamma_{\{x\}}(M)=M$ or $\Gamma_{\{x\}}(M)=\{0\}$. The first case is not possible since $x\notin \lL$, hence $\Gamma_{\{x\}}(M)=\{0\}$. This means that $x$ acts freely on $M$.
\end{proof}

\begin{example}\label{exFK} Let  $\gG=\sL(n)$ and $\lL:=\bigoplus_{j>i} \gG^{\eps_i-\eps_j}\oplus\{h\in\hh |(\eps_1- \frac{1}{n-1} (\eps_2+ \dots+ \eps_n))(h)=0 \}$ (in matrix form $\lL$ is the subalgebra of all upper triangular matrices with zero entry in the top left corner). Then $\lL$ is not a Fernando-Kac subalgebra of $\gG$, i.e. for every irreducible $(\gG,\lL)$-module $M$ we have $\gG[M]\supsetneq\lL$.  
\end{example}

\begin{proof}
Assume on the contrary there is an irreducible module $M$ such that $\gG[M]=\lL$. Let $\nn$ denote the nilradical of $\lL$. Let  $\{h_1,\dots,h_{n-1}\}$ be an orthonormal basis of $\hh$ with respect to the trace form, such that $h_2,\dots, h_{n-1}\in\hh\cap\lL$. This implies $h_1\notin \lL$. Set $g^{\eps_i-\eps_j}=E_{ij}$, where $E_{ij}$ is the elementary matrix $(a_{ij})$ with $a_{i'j'}=\doublebrace{1}{i'=i, j'=j}{0}{\mathrm{otherwise}}$. Then the Casimir element $c\in U(\gG)$ has the form
\[
c:=h_1^2+\dots +h_{n-1}^2+\sum_{i<j}g^{\eps_i-\eps_j}g^{\eps_j-\eps_i}+\sum_{i>j}g^{\eps_j-\eps_i}g^{\eps_i-\eps_j}.
\]
We then compute
\begin{eqnarray}\label{eq2}
c&=&h_1^2+h_2^2+\dots+h_{n-1}^2+\sum_{i<j}g^{\eps_i-\eps_j}g^{\eps_j-\eps_i}+\sum_{i>j}g^{\eps_j-\eps_i}g^{\eps_i-\eps_j}\notag\\
&=&h_1^2+h_2^2+\dots+h_{n-1}^2+2\sum_{i<j}g^{\eps_j-\eps_i}g^{\eps_i-\eps_j}+\sum_{i>j}[g^{\eps_j-\eps_i},g^{\eps_i-\eps_j}]\notag\\
&=&h_1^2+h_2^2+\dots+h_{n-1}^2+2\sum_{i<j}g^{\eps_j-\eps_i}g^{\eps_i-\eps_j}+l(h_1,\dots,h_n),
\end{eqnarray}
where $l$ is some linear function without constant term. Since $\lL$ is a solvable Lie algebra and $\gG[M]=\lL$, we can apply Lie's Theorem to a finite-dimensional submodule of $M$. Therefore there is a $\lL$-eigenvector $m\in M$ with $\nn\cdot m =0$ and $h_i\cdot m=\lambda_im$ for some numbers $\lambda_i$ and $2\leq i\leq n-1$. Recall the infinite-dimensional version of Schur's Lemma (\cite[2.6.8]{Dixmier}), which states that $c$ acts on the irreducible module $M$ via a constant $\mu$. Thus \refeq{eq2} implies that for some number $\alpha$ we have that $m$ is an eigenvector of $h_1^2+\alpha h_1$. Therefore $h_1$ has a locally finite action on $m\in M$, which is impossible by Lemma \ref{leActFree}. Contradiction.
\end{proof}

\section{$\kk$-multiplicities of $(\gG,\kk)$-modules}
\begin{defi} \textbf{(Multiplicity in finite-dimensional modules)}
Let $\kk$ be a Lie algebra, $M$ be a finite-dimensional $\kk$-module, and $N$ be an irreducible $\kk$-module. Let $M_1\subset M_2\subset\dots\subset M$ be  Jordan-H\"older series of $M$ as a $\kk$-module. The $\kk$-multiplicity of $N$ in $M$ is defined as the number of irreducible finite-dimensional $\kk$-modules $M_i/M_{i+1}$ isomorphic to $N$.
\end{defi}
\begin{defi}\textbf{(Multiplicity in locally finite modules)}
Let $M$ be a $\kk$-module such that $\kk=\gG[M]$. Let $N$ be a finite-dimensional irreducible $\kk$-module. The $\kk$-multiplicity of $N$ in $M$ is defined as the supremum of the multiplicities of $N$ in all finite-dimensional $\kk$-submodules $M'\subset M$.
\end{defi}

An important subcategory of the category of all $(\gG,\kk)$-modules is the category of $\kk$-semisimple $(\gG,\kk)$-modules, i.e. $(\gG,\kk)$-modules $M$ that admit a decomposition
\begin{equation}\label{eqgkm}
M=\bigoplus_{\lambda\in I} M_\lambda,
\end{equation}
where  $M_\lambda$ are simple finite-dimensional $\kk$-modules of highest weight $\lambda$ (for some Borel subalgebra of $\kk$) and $I$ is some indexing set (multiplicities greater than one are allowed). %If $ x\in C_{U(\gG)}(\kk)$, $\ker x$ is $\kk$-invariant and therefore $ x$ either acts by zero on each $M_\lambda$ or sends it to an isomorphic (due to $[x,\kk]=0$) copy. 

We can rewrite \refeq{eqgkm} as 
\begin{equation}\label{eqgkm2}
M\simeq \bigoplus_{\lambda} M_\lambda\otimes M^\lambda,
\end{equation}
where $M^\lambda:=\Hom_\kk(M_\lambda,M)$. The simple $\kk$-modules $M_\lambda$ with $M^\lambda\neq 0$ are called the \textit{$\kk$-types of $M$}.

In the case when $\kk\subset\gG$ is a semisimple Lie subalgebra and $M$ is a $(\gG,\kk)$-module, Weyl's semisimplicity theorem (\cite[Ch.~1,\S 6,~Th.~2]{Bourbaki}) implies that $M$ always has a decomposition \refeq{eqgkm}. If $\kk$ is reductive in $\gG$ and $M$ is simple, then $M$ is also necessarily $\kk$-semisimple. This follows from the fact that $\kk$ acts semisimply on $U(\gG)$, and hence on any quotient of $U(\gG)$.

\begin{defi}\label{defgkmft}
A $(\gG,\kk)$-module $M$ is said to be of finite type if every irreducible finite-dimensional $\kk$-module has finite $\kk$-multiplicity in $M$. A $(\gG,\kk)$-module is said to be of infinite type if the multiplicity of every finite-dimensional irreducible $\kk$-module is either infinite or zero.
\end{defi}
In the case when $\gG$ is semisimple, $\kk$ is reductive in $\gG$ subalgebra, and $M$ is an irreducible $(\gG,\kk)$-module, Lemma 6 from \cite{PS:GenHarishChandra} (see also the Erratum to \cite{PS:GenHarishChandra}) proves that $M$ is either of finite or of infinite type, i.e. that there are no irreducible $(\gG,\kk)$-modules of ``mixed type''.
\begin{defi}\label{defFKft}
A Lie subalgebra $\lL\subset\gG$ is Fernando-Kac subalgebra of finite type if there exists an irreducible $(\gG,\kk)$-module $M$ of finite type for which $\lL=\gG[M]$. If $\lL$ is a Fernando-Kac subalgebra and $\lL$ is not of finite type, we say that $\lL$ is of infinite type.
\end{defi}
\begin{defi}\label{defB}
A $\kk$-module $M$ is bounded  if $\kk=\kk[M]$ and there exists an integer $C$ such that every irreducible $\kk$-module has multiplicity at most $C$. A $(\gG,\kk)$-module is bounded if it is bounded as a $\kk$-module. 
\end{defi}
\begin{defi}\label{defFKB}
A Lie subalgebra $\lL\subset\gG$ is a bounded Fernando-Kac subalgebra if there exists a bounded irreducible $(\gG,\kk)$-module $M$ for which $\lL=\gG[M]$. 
\end{defi}

The present thesis focuses only on $(\gG, \kk)$-modules of finite type. Our main interest lies in the study of the irreducible ones and the corresponding Fernando-Kac subalgebras of finite type. 

Before we concentrate on irreducible $(\gG,\kk)$-modules, let us note that there is an obvious way to construct reducible strict $(\gG,\kk)$-modules. Indeed, if $M_0$ is a finite-dimensional $\kk$-module, then the induced module $U(\gG)\otimes_{U(\kk)} M_0$ is a strict $(\gG,\kk)$-module. %, called the induced module. Here $U(\gG)$ acts by left multiplication and the tensor product is taken with respect to the right ${U(\kk)}$-module structure of $U(\gG)$. Note that any quotient of the induced module is also a $(\gG,\kk)$-module. 
Any quotient of the induced module is also a $(\gG,\kk)$-module. Moreover, all cyclic (and hence all irreducible) $(\gG,\kk)$-modules $M$ arise as quotients of $U(\gG)\otimes_{U(\kk)} M_0$ for some finite-dimensional $M_0$. Indeed, let $M$ be a cyclic $(\gG,\kk)$-module generated by the vector $m\in M$. Since $M$ is $\kk$-locally finite, $\kk$ generates a finite-dimensional $\kk$-submodule $M_0\subset M$. Furthermore, there is a natural surjective homomorphism from $U(\gG)\otimes_{U(\kk)} M_0$ to $M$, and hence $M$ is the quotient of $U(\gG)\otimes_{U(\kk)} M_0$ by the kernel of this homomorphism. However, for a fixed irreducible $M_0$, it is not clear how to find all possible such kernels. In addition, it is in general not clear how to compute the Fernando-Kac subalgebra associated to a given quotient of  $U(\gG)\otimes_{U(\kk)} M_0$.

\section{Some known results on Fernando-Kac subalgebras of finite type and an existence theorem}
We start by recalling a general characterization of Fernando-Kac subalgebras of finite type given by I. Penkov, V. Serganova and G. Zuckerman.
\begin{theo} \label{thNec}

\textbf{\cite[Theorem 3.1(1),(3)]{PSZ}}
Let $\lL\subset\gG$ be a Fernando-Kac subalgebra of finite type.
\begin{itemize}
\item[(a)] $N(\lL)=\lL$. 
\item[(b)] $\lL$ has a well-defined reductive in $\gG$ part $\kk=\lL_{red}$.
\item[(c)] Any irreducible $(\gG,\kk)$-module $M$ of finite type over $\lL$ has finite type over $\kk$ and $\kk$ acts semi-simply on $M$.
%\item[(c)] $C(\lL_{red})=Z(\lL_{red})$ and $Z(\lL_{red})$ is a Cartan subalgebra of finite type over $C(\lL_{ss})$.
\end{itemize}
\end{theo}
\begin{proof}

(a) Let $M$ be an irreducible strict $(\gG,\lL)$-module and let $M_0\subset M$ be an irreducible finite-dimensional $\lL$-submodule (to make sure such exists take an irreducible submodule of any finite-dimensional $\lL$-submodule of $M$). %Assume on the contrary that $N(\lL)\neq\lL$. Let $l\in\lL$. Since $\ad(\l_{ss})=(\ad(l))_{ss}$, $\ad(l_{ss})$ can be expressed as a polynomial with zero constant term of $\ad l$ we have $\ad(\lL_{ss})(\lL)\subset\lL$. By the definition of semisimple algebra the $\ad(\lL_{ss})$-invariant space $\lL$ has an invariant complement space $M$, $N(\lL)=M\oplus\lL$. However, $(\ad\lL_{ss}) (M)\subset \lL$ because $M$ is in the normalizer of $\lL$ and again because $\ad\lL_{ss}$ is a polynomial with zero constant term in $\ad\lL$. Therefore $\ad\lL_{ss}(M)=0$. Pick a non-zero element in $x\in M$.
Assume on the contrary that $N(\lL)\neq\lL$, and pick $x\in N(\lL)\backslash\lL$. Set
\[
M_n:=M_0+x\cdot M_0+\dots+x^n\cdot M_0.
\]
We are proving that $M_n$ is $\lL$ invariant, using the fact that $M_0$ is $\lL$-invariant and that commuting $x$ with $\lL$ produces elements in $\lL$. Indeed, we proceed by induction on $n$. Assume that $M_{n-1}$ is $\lL$-invariant. Pick any $l\in\lL$. Then $l\cdot M_n= l\cdot ((x^n\cdot M_0)+M_{n-1})\subset x\cdot l\cdot x^{n-1}M_0+[l,x]\cdot x^{n-1}M_0+M_{n-1}$. But $[l,x]\in\lL$ ($x$ is in the normalizer of $\lL$) and so by the induction hypothesis $[l,x]\cdot x^{n-1}M_0\subset M_{n-1}$ and $l\cdot x^{n-1}M_0\subset M_{n-1}$. Therefore $l\cdot M_n\subset M_n$.

We claim next that $M_{n}/M_{n-1}\simeq M_0$. More precisely we claim that the epimorphism $\phi:M_0\to M_n/M_{n-1},\quad m\mapsto x^n\cdot m/M_{n-1}$ is injective. Indeed, since $M_0$ is an irreducible $\lL$-module, $\phi$ is injective unless $\phi=0$. However the latter would imply $M_{n}=M_{n-1}$ which in turn would imply that the action of $x$ on $M_0$ is locally finite. On the other hand, by \refle{leActFree} and by the assumption that $x\notin\lL$, the action of $x$ on any vector in $M_0$ is free. Contradiction. Hence $\phi$ is injective and we have established that $M_{n}/M_{n-1}\simeq M_0$.

Now the assumption $N(\lL)\neq\lL$ implies that the multiplicity of $M_0\simeq M_n/M_{n-1}$ for all $n$ is infinite in $M$. This is impossible as $M$ has finite type of $\lL$. Contradiction.

(b) It is well known that a self-normalizing subalgebra over an algebraically closed field has a well-defined reductive part. This follows, for example, from Corollary 1 and Proposition 7 in \cite[Chapter 7,\S 5]{Bourbaki}.

(c) Since $M$ is generated by any of its non-zero vectors ($M$ is irreducible), we have that $M$ is a quotient of the induced module $U(\gG)\otimes_{U(\lL)}M_0$. By the Poincare-Birkhoff-Witt Theorem, $U(\gG)_k$ is a finite-dimensional  $\kk$-module for all $k$, where $U(\gG)_k$ denotes the subspace of $U(\gG)$ generated by monomials of total degree less than or equal to $k$. 

It remains to prove that $M$ has finite type over $\kk$. By \cite[Ch.1 par. 6.8]{Bourbaki} we know that $\nn_\lL\subset \nn+ Z(\lL)$, where $\nn$ is the nilradical of $\lL$. This, together with Schur's Lemma (applied to the irreducible module $M_0$) implies that there exists $\lambda\in\nn_\lL$ such that 
\[x\cdot m=\lambda(x)m\quad \forall x\in\nn_\lL, \forall m\in M_0.\]
The fact that $\nn_\lL$ has nilpotent action on $\gG$ together with the equality $(\ad x) ( g_1\dots g_k)=[ x,g_1]g_2\dots g_k+\dots+ g_1\dots  g_{k-1}[ x,g_k]$ imply that for any $x\in\nn_\lL$, $x$  has nilpotent adjoint action on $U(\gG)_k$. Therefore $x-\lambda(x)$ acts nilpotently on any $\lL$-subquotient of $U(\gG)\otimes_{U(\kk)}M_0$ and hence acts by zero on any irreducible subquotient. Therefore two irreducible $\lL$-subquotients of $M$ are isomorphic if and only if they are isomorphic as $\kk$-modules. Since $M$ has finite type over $\lL$, it necessarily has also finite type over $\kk$.
\end{proof}


We recall next a key theorem from \cite{PSZ} which ensures the existence of irreducible $(\gG,\kk)$-modules of finite type with a certain prescribed Fernando-Kac subalgebra. Let $\gG$  be a reductive Lie algebra with Cartan subalgebra $\hh$ and let $\kk\supset\hh$ be a reductive root subalgebra of $\gG$. Let $h\in\hh$ be an element such that $\gamma(h)\neq 0$ for all $\gamma\in\Delta(\kk)$ and let $\pp:= \bigoplus_{\gamma(h)\geq 0} g^{\gamma}$ be a parabolic subalgebra of $\gG$. Then $\pp_{red}=\bigoplus_{\gamma(h)= 0} g^{\gamma}$ and let the nilradical of $\pp$ is $\nn_{\pp}:=\bigoplus_{\gamma(h)> 0} \gG^{\gamma}$. Let $L$ be an irreducible $(\pp, \hh\cap\kk)$-module of finite type over $\hh\cap\kk$ with trivial action of $\nn_{\pp}+ (Z(\pp_{red}) \cap \kk_{ss})$ and with $\pp_{red}$ -central character $\theta_{\pp_{red}}^\nu$ for some $\pp_{ss}\cap\bb$-dominant weight $\nu\in\hh^*$.
\begin{theo}\label{thPSZmain}
\textbf{(\cite[Theorem 4.3]{PSZ})} Then there exists a module $M$ for such that
\begin{itemize}
\item[(a)] $M$ is an infinite-dimensional irreducible $\gG$-module;
\item[(b)] $\gG[M] = \kk_{ss}\oplus \mm_L$, where $\mm_L$ is the maximal $\kk_{ss}$-stable subspace in $\pp[L]$; moreover $\gG [M ]$ is the unique maximal subalgebra in $\pp [L] + \kk$ which contains $\kk$;
\item[(c)] $M$ is a $(\gG, \kk)$-module of finite type over $\kk$.
\end{itemize}
\end{theo}
In \cite{PSZ} there is an explicit formula for the module $M$ in terms of two additional parameters. We do not present it as that would require extra notation which would not be otherwise used in this thesis. The proof of the theorem uses the Beilinson-Bernstein localization theorem and Kashiwara's theorem (\cite{Kashiwara:holonomic}, \cite{BeilinsonBerstein}).


We present some corollaries of Theorems \ref{thNec} and \ref{thPSZmain}.
\begin{theo}[\cite{PSZ}]\label{thFK}
Let $\kk$ be a reductive in $\gG$ subalgebra for a semisimple algebra $\gG$ such that the centralizer $C(\kk)$ of $\kk$ equals the center $Z(\kk)$ of $\kk$. Then there exists a Fernando-Kac subalgebra of finite type $\lL\subset\gG$ such that $\lL_{red}=\kk$. Conversely, if $\lL$ is a Fernando-Kac subalgebra of finite type, then $\lL_{red}$ is well defined and $C(\lL_{red})=Z(\lL_{red})$.
\end{theo}
The necessity of the condition $C(\lL_{red})=Z(\lL_{red})$ is the much easier part of the theorem, as it uses only Theorem \ref{thNec}. \refth{thFK} reduces the classification of Fernando-Kac subalgebras of finite type to the classification of all possible subalgebras $\nn$ for a fixed reductive in $\gG$ subalgebra $\kk=\lL_{red}\subset\gG$ with $C(\lL_{red})=Z(\lL_{red})$. In the specific case of $\gG=\gl(n)$, the following two theorems characterize the reductive in $\gG$ Fernando-Kac subalgebras of finite type (here $\lL=\kk$, $\nn=\{0\}$).
\begin{theo}[\cite{PSZ}]\label{thgl(n)}
A reductive in $\gl(n)$ subalgebra $\lL$ is a Fernando-Kac subalgebra of finite type if and only if $C(\lL)=Z(\lL)$.
\end{theo}
\begin{corollary}[\cite{PSZ}]\label{corgl(n)red}
A reductive in $\gl(n)$ subalgebra $\kk$ is a Fernando-Kac subalgebra of finite type if and only if the defining $\gG$-module is multiplicity free.
\end{corollary}
The problem of describing all nilpotent Lie algebras $\nn$ so that $\kk\crplus\nn$ is Fernando-Kac of finite type turns out to be much more complicated, and is still open. In the present work we solve this problem under the assumption that $\kk$ is a root subalgebra of a simple Lie algebra $\gG\nsimeq E_8$.

The following theorem is the main motivation for Penkov's conjecture. We note that its formulation of the cone condition is slightly different from the one we give in the next chapter.
\begin{theo}\label{throots} (\cite{PSZ})
A root subalgebra $\lL=(\lL_{red}\crplus\nn)\subset\gl(n)$ is a Fernando-Kac subalgebra of finite type if and only if $\Cone_{\QQ}(\hw_{\bb\cap\kk}(\gG/\lL))\cap \Cone_{\QQ}(\Delta(\nn))=\{0\}$, where $\hw_{\bb\cap\kk}(\gG/\lL):= \{\alpha \in \Delta (\gG)\backslash\Delta(\lL) | \alpha + \delta \notin \Delta(\gG), \forall \delta \in\Delta(\bb)\cap\Delta(\kk)\}$ (cf. Definition \ref{defCone} below).
\end{theo}

Theorem \ref{thPSZmain} uses as input data an irreducible strict $(\pp_{red},\hh)$-module $L$ of finite type over $\hh$, where $\pp$ is a parabolic subalgebra of $\gG$. Furthermore, determining which values are allowed for $\pp_{red}[L]$ in an irreducible $(\pp_{red},\hh)$-module $L$ of finite type is a partial case of the problem studied here. In particular, it is important to know whether the equality $\pp_{red}[L]=\hh$ is possible for a given $\pp$. A definitive answer is given by S. Fernando, \cite[Theorem 5.2]{Fernando1}: an irreducible strict $(\pp_{red},\hh)$-module $L$ of finite type exists if and onlyif $\pp_{red}$ has simple components of types $A$ and $C$ only. 

In \cite{Mathieu}, O. Mathieu classified the irreducible weight modules (i.e. the $(\gG,\hh)$-modules of finite type). However, for our purposes, it suffices to present any construction of strict irreducible $(\sL(n),\hh)$- and $(\sP(2n),\hh)$-modules of finite type. One rather elementary way to obtain such a construction comes from the natural embeddings of $\sL(n)$ and $\sP(2n)$ in the Weyl algebra in $n$ variables $W_n$. We recall that $W_n$ is the algebra over $\CC$ generated by letters $x_1,\dots, x_n$ and differential operators $ \partial_{x_1},\dots,\partial_{x_n}$ with the usual commutation relations. The construction we describe below can be found, for example, in \cite{BrittenLemire}.

Let the natural module of $\gl(n)$ be $V'$, let $\{d_1,\dots d_n\}$ be an arbitrary basis of $V'$, and let $\{d_1^*,\dots d_n^*\}$ be the corresponding dual basis of $(V')^*$. Write $\gl(n)$ in tensor notation as $\linspan\{d_i\otimes d_j^* \} $. Let the natural module of $\sP(2n)$ be $V''$ and let $B$ be the defining invariant symplectic form on $V''$. Let $\{e_1,\dots e_n, f_1, \dots f_n\}$ be a basis of $V''$ such that $B(e_i, e_j)=0$, $B(e_i,f_j)= -B(f_j,e_i) =\doublebrace{1 }{i=j}{0}{\mathrm{otherwise}}$, $B(f_i,f_j)=0$. Let $\{e_1^*,\dots, e_n^*, f_1^*,\dots f_n^* \}$ be its corresponding dual basis in ${V''}^*$.  Write $\sP(2n)$ in tensor notation as $\linspan\bigcup_{i,j} \{e_i \otimes e_j^*-f_j\otimes f_i^*\}\cup \bigcup_{i\neq j}\{e_i\otimes f_j^*- e_j\otimes f_i^*\}\cup \bigcup_{i\neq j}\{f_i\otimes e_j^*- f_j\otimes e_i^*\}_{i\neq j} $. We can embed $\gl(n)$ in $\sP(2n)$ via the map $\theta: \gl(n) \hookrightarrow \sP(2n)$ given by $\theta( d_i\otimes d_j^*):=e_i\otimes e_j^* - f_i\otimes f_j^* $. 

We can embed $\gl(n)\stackrel{\phi}{\hookrightarrow} W_n$ via the map $\phi (e_i\otimes e_j^*):= x_i\partial_{x_j}$. Set $\psi:=\phi_{|\sL(n)}$. Following \cite{Mathieu}, we can write an embedding $\bar{\psi}$ of $\sP(2n)$ in $W_n$ such that $\bar\psi=\theta_{|\sL(n)}\circ\bar{\psi}$. Indeed, we can set $\bar \psi(e_i\otimes e_j^*-f_i\otimes f_j^*):=x_i\partial_{x_j}$, where  $i\neq j$,  $\bar\psi(e_i\otimes f_j^*+e_j\otimes f_i^*):= x_ix_j$, where $i=j$ is allowed, $\bar \psi(f_i\otimes e_j^*+f_j\otimes e_i^*):= -\partial_{x_i}\partial_{x_j}$, where $i=j$ is allowed, and $\bar\psi(e_i\otimes e_i^*-f_i\otimes f_i^*):=x_1\partial_{x_1}+\frac{1}{2} $. Note that $\bar \psi$ and $\phi$ have the same action on $\theta(\sL(n))$, but not on $\theta(\gl(n))$. The embedding $\psi$ is called the Shale-Weil representation of $\sP(2n)$ in $\CC[x_1,\dots, x_n]$, \cite{Mathieu}.

The embeddings of $\sL(n)$ and $\sP(2n)$ in $W_n$ give a natural way to construct strict irreducible $\sL(n)$- and $\sP(2n)$-modules of finite type. Consider any formal monomial $x_1^{\alpha_1}\dots x_n^{\alpha_n}$, with not necessarily integral exponents $\alpha_i\in\CC$. The realization of $\gG:=\sL(n)$ in $W_n$ induces a natural action on the vector space $M:=\linspan_{\CC}$ $\{x_1^{\beta_1}\dots x_n^{\beta_n} $ $\in$ $x_1^{\alpha_1} \dots x_n^{\alpha_n}  \CC[x_1, \dots, x_n, x_1^{-1},\dots,x_n^{-1}]$ $| \sum_i\alpha_i=\sum\beta_i\}$. The explicit form of the embeddings implies that the Cartan subalgebra of $\sL(n)$ acts locally finitely and semisimply on $M$. If $\alpha_i-\alpha_j$ is not an integer for all $i,j$, then none of the root spaces of $\sL(n)$ acts locally finitely. Furthermore, these conditions imply that the module is irreducible. By Theorem \ref{thFernando-Kac}, we have that $\gG[M]$ equals the Cartan subalgebra $\hh$ of $\sL(n)$. We call the $\sL(n)$-modules constructed above Britten-Lemire modules. %The Britten-Lemire modules constructed here parametrize all irreducible multiplicity free $(\sL(n),\hh)$-modules (\cite{BrittenLemire}).

Similarly, the realization of $\gG=\sP(2n)$  in $W_n$ induces a natural action on the vector space $M':=x_1^{\alpha_1}\dots x_n^{\alpha_n}\CC[x_1,\dots, x_n, x_1^{-1},\dots,x_n^{-1}]$. The Cartan subalgebra of $\gG$ is a subalgebra of $\gG[M']$. If none of the $\alpha_i$ and none of the possible sums $ \alpha_i\pm\alpha_j\neq 0$ is integral for all $i,j$, then none of the root spaces of $\sP(n)$ acts locally finitely. Furthermore, these conditions imply that $M'$ is irreducible. By theorem \ref{thFernando-Kac} we have that the Cartan subalgebra equals $\gG[M']$.

In both cases -- $\sL(n)$ acting on $M$ and $\sP(2n)$ acting on $M'$ -- we see that for any two monomials $x_1^{\beta_1}\dots x_n^{\beta_n}\in M$ (respectively, $x_1^{\beta_1}\dots x_n^{\beta_n}\in M'$) and $x_1^{\beta_1'}\dots x_n^{\beta_n'}\in M$ (respectively, $x_1^{\beta_1'}\dots x_n^{\beta_n'}\in M'$) there is an element of the Cartan subalgebra having different eigenvalues on the two monomials. This shows that, in addition,  both modules $M$ and $M'$ are multiplicity-free as weight modules.

\begin{comment}
In this section, we give explicit embeddings of $\so(2n)$, $\sP(2n)$ in $\sL(2n)$ and an embedding of $\so(2n+1)$ in $\sL(2n+1)$ following the book of E. Vinberg and A. Onishchik \cite{VinbergOnishchik}. The embeddings of $\so(2n)$, $\sP(2n)$ in $\sL(2n)$ and $\so(2n+1)$ in $\sL(2n+1)$ are given by specifying a vector space $V$ (of dimension $2n$ for $\so(2n)$ and $\sP(2n)$, and of dimension $2n+1$ for $\so(2n+1)$) and a non-degenerate form $B$, symmetric for $\so(2n), \so(2n+1)$, and antisymmetric for $\sP(2n)$. The corresponding subalgebras are then explicitly realized as $\{g\in\gl(V)| B(gx,y)+B(x,gy)=0, \mathrm{~for~all~}x,y\in V\}$; the latter is automatically a subset of $\sL(V)\subset\gl(V)$. 

We now give one possible matrix of the embeddings of $\so(2n)$, $\sP(2n)$ in $\sL(2n)$ by specifying concrete values for $B$. We denote by  $B_{mat}$ the matrix of $B$ in the bases of $V$ we fix below. For these bases, the matrices of each  $\so(2n)$, $\sP(2n)$, $\so(2n+1)$ are then given by $\{g|~ g^TB_{mat}+B_{mat}g=0\}$.
\begin{itemize}
\item $\so(2n)$. Set $V:=\linspan_\CC\{e_1,\dots e_n, f_1, \dots f_n\}$ and $B(e_i, e_j):=0$, $B(e_i,f_j):=B(f_j,e_i):=\delta_{ij}$, $B(f_i,f_j):=0$. Then the matrix realization of $\so(2n)$ is the following:
\[
\so(2n)=\left\{\left(\begin{array}{c|c}A& C \\\hline D& -A^T\end{array}\right)|~ C=-C^T, D=-D^T\right\}, \quad B_{mat}=\left(\begin{array}{c|c}0& I \\\hline I& 0\end{array}\right).
\]
If $\{e_1^*,\dots, e_n^*, f_1,\dots f_n^* \}$ is the basis dual to $\{e_1,\dots e_n, f_1, \dots f_n\}$, then the above embedding of $\so(2n)$ can be written in tensor notation as $\linspan\bigcup_{i,j}\{e_i\otimes e_j^*-f_j\otimes f_i^*\}\cup \bigcup_{i\neq j}\{e_i\otimes f_j^*- e_j\otimes f_i^*\}\cup \bigcup_{i\neq j}\{f_i\otimes e_j^*- f_j\otimes e_i^*\}_{i\neq j} $.
\item $\sP(2n)$. Set $V:=\linspan_\CC\{e_1,\dots e_n, f_1, \dots f_n\}$ and $B(e_i, e_j):=0$, $B(e_i,f_j):=-B(f_j,e_i):=-\delta_{ij}$, $B(f_i,f_j):=0$. Then the matrix realization of $\sP(2n)$ is the following:
\[
\sP(2n)=\left\{\left(\begin{array}{c|c}A& C \\\hline D& -A^T\end{array}\right)| C=C^T, D=D^T\right\}, \quad B_{mat}=\left(\begin{array}{c|c}0& -I \\\hline I& 0\end{array}\right).
\]
If $e_1^*,\dots, e_n^*, f_1,\dots f_n^* $ are the basis dual to $\{e_1,\dots e_n, f_1, \dots f_n\}$, then the above embedding of $\sP(2n)$ can be written in tensor notation as $\linspan\bigcup_{i,j}\{e_i\otimes e_j^*-f_j\otimes f_i^*\}\cup \bigcup_{i,j}\{e_i\otimes f_j^*+ e_j\otimes f_i^*\}\cup \bigcup_{i,j}\{f_i\otimes e_j^*+ f_j\otimes e_i^*\} $.
\item $\so(2n+1)$. Set $V:=\linspan_\CC\{e_1,\dots e_n, f_1, \dots f_n, d\}$, $B(e_i, e_j):=0$, $B(e_i,f_j):=B(f_j,e_i):=\delta_{ij}$, $B(f_i,f_j):=0$, $B(e_i,d):=B(d,e_i):=B(f_i,d):=B(d,f_i):=0$, $B(d,d)=1$. Then the matrix realization of $\so(2n+1)$ is the following:
\[
\so(2n+1)=\left\{\left(\begin{array}{c|c|c}A&\begin{array}{c}v_1\\ \vdots \\ v_n\end{array} &C=-C^T \\\hline \begin{array}{ccc}w_1 &\dots&  w_n\end{array} &0& \begin{array}{ccc}-v_n &\dots&  -v_1\end{array} \\\hline D=-D^T&\begin{array}{c}-w_n\\ \vdots \\ -w_1\end{array} & -A^T\end{array}\right)\right\}, \]
\[B_{mat}=\left(\begin{array}{c|c|c}0& \begin{array}{c}0\\ \vdots \\ 0\end{array}&I \\\hline \begin{array}{ccc}0 &\dots&  0\end{array}&1&\begin{array}{ccc}0 &\dots&  0\end{array}\\ \hline I& \begin{array}{c}0\\ \vdots \\ 0\end{array}&0\end{array}\right).
\]
\end{itemize}

From the above it is clear how one can embed $\gl(n)\hookrightarrow\sP(2n)$ and $\gl(n)\hookrightarrow\so(2n)$ (the latter in turn naturally embeds in $\so(2n+1)$). Indeed, we note that in both cases the matrices of the embedding  are $A\mapsto\left(\begin{array}{c|c}A&0\\\hline 0&-A^T\end{array}\right)$ (for the already fixed choices of basis and bilinear form).

We can embed $\gl(n)\stackrel{\phi}{\hookrightarrow} W_n$, where $W_n$ is the Weyl algebra $\CC[x_1,\dots, x_n, \partial_{x_1}, \dots, \partial_{x_n}]$ via the map $\phi(e_i\otimes e_j^*):=x_i\partial_{x_j}$. This map induces an embedding $\sL(n) \stackrel{\phi|{\sL(n)}} {\hookrightarrow} W_n$. Set $\psi:=\phi_{|\sL(n)}$. Following \cite{Mathieu}, we can extend $\psi$ to an embedding $\sL(n)\hookrightarrow \sP(2n) \stackrel{\bar \psi}{\hookrightarrow} W_n$ (coming from a Shale-Weil representation in $\CC[x_1,\dots, x_n]$)  by setting $\bar \psi(e_i\otimes e_j^*-f_i\otimes f_j^*)=x_i\partial_{x_j}$, where  $i\neq j$,  $\bar\psi(e_i\otimes f_j^*+e_j\otimes f_i^*):= x_ix_j$, where $i=j$ is allowed, $\bar \psi(f_i\otimes e_j^*+f_j\otimes e_i^*):= -\partial_{x_i}\partial_{x_j}$, where $i=j$ is allowed, and $\bar\psi(e_i\otimes e_i^*-f_i\otimes f_i^*):=x_1\partial_{x_1}+\frac{1}{2} $. Note that $\bar \psi$ and $\phi$ have the same action on $\sL(n)\hookrightarrow\sP(2n)$, but not on $\gl(n)\hookrightarrow\sP(2n)$.

The coefficients chosen in the definition of $\bar{\psi}$ we justify by the following equalities.
\begin{eqnarray*}&&[{x_1}{x_2} , -\partial_{x_1}\partial_{x_2} ]= {x_1}\partial_{x_1}+{x_2}\partial_{x_2}+1
\\&&[{x_1}\partial_{x_1}+{x_2}\partial_{x_2}+1 , {x_1}{x_2} ]= 2{x_1}{x_2}
\\&&[{x_1}\partial_{x_1}+{x_2}\partial_{x_2}+1 , -\partial_{x_1}\partial_{x_2} ]= 2\partial_{x_1}\partial_{x_2}
\\&&[{x_1}^{2} , -\partial_{x_1}^{2} ]= 4{x_1}\partial_{x_1}+2
\\&&[4{x_1}\partial_{x_1}+2 , -\partial_{x_1}^{2} ]= 8\partial_{x_1}^{2}
\\&&[4{x_1}\partial_{x_1}+2 , {x_1}^{2} ]= 8{x_1}^{2}
\\&&[{x_1}\partial_{x_2} , {x_2}\partial_{x_1} ]= {x_1}\partial_{x_1}-{x_2}\partial_{x_2}\end{eqnarray*}
\end{comment}
\chapter{Fernando-Kac subalgebras arising through a construction of Penkov, Serganova and Zuckerman}\label{secTheConstruction}
\section{Main result of the thesis}
In this section we formulate Penkov's conjecture for all simple Lie algebras. The rest of the thesis is dedicated to proving this conjecture in all cases except $E_8$.

Recall that $\lL$ is a root subalgebra of $\gG$.
\begin{defi}\label{defCone}
(I. Penkov) 
\begin{itemize}
\item[(a)] \emph{Cone condition}. We say that $\lL$ satisfies the cone condition if $\Cone_{\QQ}(\Delta(\nn))$ $\cap$ $ \Cone_{\QQ}( \hw_{ \bb \cap \kk} (\gG/\lL))$ $=\{0\}$, where $\hw_{\bb\cap\kk}(\gG/\lL):= \{\alpha \in \Delta (\gG)\backslash\Delta(\lL) | \alpha + \delta \notin \Delta(\gG), \forall \delta \in\Delta(\bb)\cap\Delta(\kk)\}$ are the weights of the $\bb\cap\kk$-singular vectors of the $\kk$-module $\gG/\lL$.
\item[(b)] \emph{Centralizer condition}. We say that $\lL$ satisfies the centralizer condition if a (equivalently any) Levi subalgebra of the Lie algebra $C(\kk_{ss})\cap N(\nn)$ has simple constituents of type $A$ and $C$ only.
\end{itemize}
\end{defi}
\textbf{Remark.} The cone condition (a) holds as stated if and only if it holds with $\QQ$ replaced by $\ZZ$. %We will use this without further explicit mention.

\begin{conjecture}[Penkov's  conjecture]\label{conjPenkov}
Let $\lL=\kk\crplus\nn$ be a root subalgebra  of a simple Lie algebra. Then $\lL$ is a Fernando-Kac subalgebra of finite type if and only if the cone condition and the centralizer condition are satisfied.
\end{conjecture}
The main result of this thesis is the following.
\begin{theo}\label{thMainResult}
If $\gG\nsimeq E_8$, Penkov's conjecture holds. If $\gG\simeq E_8$, then $\lL$ is a Fernando-Kac subalgebra of finite type whenever the centralizer and cone conditions hold.
\end{theo}
In the current chapter, we construct via the theorem of I. Penkov, V. Serganova and G. Zuckerman (Theorem \ref{thPSZmain}) an irreducible $(\gG,\lL)$-module of finite type for any subalgebra $\lL$ that satisfies the cone and centralizer conditions. This completes one of the directions of the proof of Theorem \ref{thMainResult}.

The conjecture was formulated by I. Penkov for the classical Lie algebras. Due to its root-system formulation, his conjecture directly applied to the exceptional Lie algebras as well. At the time of writing of this thesis, Penkov's conjecture remains open for $E_8$ due to the computational challenge posed by the root system of $E_8$. 

Note that the criterion of Theorem \ref{thMainResult} is entirely combinatorial. This clearly applies to the cone condition. Checking the centralizer condition under the assumption that the cone condition holds is also an entirely combinatorial procedure. Indeed, in this latter case Proposition \ref{lePara} below gives that $ C(\kk_{ss})\cap N(\nn)= C(\kk_{ss}) \cap N(C(\kk_{ss})\cap\nn)$, i.e. $ C(\kk_{ss})\cap N(\nn)$ is a parabolic subalgebra of $C(\kk_{ss})$. Therefore checking the centralizer condition reduces to checking the type of the root subsystem $Q\cap -Q\subset\Delta(\gG)$, where $Q:=\{\alpha\in\Delta(\gG)~|~\alpha\sperp \Delta(\kk)$, such that for all $\beta\in\Delta(\nn)$ with $\beta\sperp\Delta(\kk)$, either $\alpha+\beta\in\Delta(\nn)$, or $\alpha+\beta$ is not a root $\}$. 


In the case when $\kk=\hh$ (i.e. $\lL$ is solvable), the cone condition is equivalent to the requirement that $\nn$ be the nilradical of a parabolic subalgebra containing $\hh$ (see \cite[Prop. 4]{PS:GenHarishChandra} and also Lemma \ref{leCentralizerAndNilradical} below). Furthermore, using Corollary 5.5 and Theorem 5.8 from \cite{PSZ}, it is not difficult to show that when $\gG$ is of type $A$, the cone condition holds if and only if $\nn$ is the nilradical of a parabolic subalgebra of $\gG$ which contains $\kk$. This is not the case in type $B,C$, and $D$. Here is an example for type $C$.

\begin{example}\label{exsp6} $\gG=\sP(6)$, $\Delta(\kk)=\{\pma2\eps_1,\pma2\eps_3\}$, $\Delta(\nn)=2\eps_2, \eps_2+ \eps_1, \eps_2- \eps_1 $.  A computation shows that $\hw_{ \bb\cap\kk } (\gG/\lL))$ = $\{\eps_3+ \eps_2,\eps_1+ \eps_3, \eps_3- \eps_2, -\eps_2+ \eps_1,-2\eps_2\}$. Thus the cone condition is satisfied, but there exists no parabolic subalgebra $\pp$ with $\nn=\nn_\pp$. Indeed, assume the contrary. Then there exists a vector $t\in h$ with $\alpha(t)>0$ for all $\alpha \in \Delta (\nn)$, and $\alpha(t)\leq 0$ for all other roots $\alpha$ of $\gG$. Therefore $\eps_1(t)=\eps_3(t)=0$,  $\eps_2(t)>0$, and $t(\eps_3+\eps_2)>0$. Contradiction.
\end{example}

To illustrate the cone condition in the non-solvable case, we present all non-solvable root subalgebras that fail the cone condition in types $B_3$ and $C_3$ (Table \ref{tableC3B3} below). These subalgebras are, up to conjugation, all non-solvable root subalgebras of infinite type in types $B_3$ and $C_3$. Indeed, the centralizer condition is trivially satisfied in type $C_3$. In type $B_3$, the centralizer condition holds for $\lL\nsimeq \hh$, as the root system $B_2$ is isomorphic to $C_2$. Up to conjugation, in $\so(7)$ (respectively, $\sP(6)$), there are 11 (respectively, 16) non-solvable root subalgebras that fail the cone condition, and 32 (respectively, 38) non-solvable root subalgebras that satisfy it. %\arxivVersion{One matrix form realization of these subalgebras is included in section \ref{secSubalgebrasB3C3} in the appendix.}


\begin{longtable}{r|l}
\multicolumn{2}{c}{ $ \mathfrak{g}\simeq \mathrm{so(7)}$} \\\hline\hline\begin{tabular}{r}$\Delta(\mathfrak{k})$ is of type $A_1$+$A_1$; \\ $\Delta^+(\mathfrak{k})=$ $\varepsilon_{1}$-$\varepsilon_{2}$, $\varepsilon_{1}$+$\varepsilon_{2}$\end{tabular} &
\begin{tabular}{l}$\Delta(\mathfrak{n})=$$\varepsilon_{1}$+$\varepsilon_{3}$, $\varepsilon_{2}$+$\varepsilon_{3}$, -$\varepsilon_{2}$+$\varepsilon_{3}$, -$\varepsilon_{1}$+$\varepsilon_{3}$, 
\end{tabular}

\\\hline\begin{tabular}{r}$\Delta(\mathfrak{k})$ is of type $A_1$; \\ $\Delta^+(\mathfrak{k})=$ $\varepsilon_{3}$\end{tabular} &
\begin{tabular}{l}$\Delta(\mathfrak{n})=$$\varepsilon_{1}$-$\varepsilon_{2}$, $\varepsilon_{1}$+$\varepsilon_{3}$, $\varepsilon_{1}$-$\varepsilon_{3}$, $\varepsilon_{1}$
\\$\Delta(\mathfrak{n})=$$\varepsilon_{1}$+$\varepsilon_{3}$, $\varepsilon_{1}$-$\varepsilon_{3}$, $\varepsilon_{1}$


\\\end{tabular}

\\\hline\begin{tabular}{r}$\Delta(\mathfrak{k})$ is of type $A_1$; \\ $\Delta^+(\mathfrak{k})=$ $\varepsilon_{1}$-$\varepsilon_{2}$\end{tabular} &
\begin{tabular}{l}$\Delta(\mathfrak{n})=$$\varepsilon_{3}$, -$\varepsilon_{2}$+$\varepsilon_{3}$, -$\varepsilon_{1}$+$\varepsilon_{3}$
\\$\Delta(\mathfrak{n})=$$\varepsilon_{1}$+$\varepsilon_{2}$, $\varepsilon_{1}$-$\varepsilon_{3}$, $\varepsilon_{2}$-$\varepsilon_{3}$, $\varepsilon_{1}$+$\varepsilon_{3}$, $\varepsilon_{2}$+$\varepsilon_{3}$


\\$\Delta(\mathfrak{n})=$$\varepsilon_{1}$+$\varepsilon_{2}$, $\varepsilon_{1}$-$\varepsilon_{3}$, $\varepsilon_{2}$-$\varepsilon_{3}$, -$\varepsilon_{2}$-$\varepsilon_{3}$, -$\varepsilon_{1}$-$\varepsilon_{3}$
\\$\Delta(\mathfrak{n})=$$\varepsilon_{1}$+$\varepsilon_{2}$, $\varepsilon_{1}$-$\varepsilon_{3}$, $\varepsilon_{2}$-$\varepsilon_{3}$, $\varepsilon_{1}$, $\varepsilon_{2}$


\\$\Delta(\mathfrak{n})=$$\varepsilon_{1}$+$\varepsilon_{2}$, $\varepsilon_{1}$-$\varepsilon_{3}$, $\varepsilon_{2}$-$\varepsilon_{3}$
\\$\Delta(\mathfrak{n})=$$\varepsilon_{1}$+$\varepsilon_{2}$, $\varepsilon_{1}$, $\varepsilon_{2}$


\\$\Delta(\mathfrak{n})=$$\varepsilon_{1}$-$\varepsilon_{3}$, $\varepsilon_{2}$-$\varepsilon_{3}$, -$\varepsilon_{2}$-$\varepsilon_{3}$, -$\varepsilon_{1}$-$\varepsilon_{3}$
\\$\Delta(\mathfrak{n})=$$\varepsilon_{1}$-$\varepsilon_{3}$, $\varepsilon_{2}$-$\varepsilon_{3}$


\\\end{tabular}

\\\hline\hline

\multicolumn{2}{c}{ $\mathfrak{g}\simeq \mathrm{sp(6)}$} \\\hline\hline\begin{tabular}{r}$\Delta(\mathfrak{k})$ is of type $A_1$; \\ $\Delta^+(\mathfrak{k})=$ -$\varepsilon_{2}$+$\varepsilon_{3}$\end{tabular} &
\begin{tabular}{l}$\Delta(\mathfrak{n})=$2$\varepsilon_{1}$, 2$\varepsilon_{3}$, 2$\varepsilon_{2}$, $\varepsilon_{2}$+$\varepsilon_{3}$
\\$\Delta(\mathfrak{n})=$2$\varepsilon_{1}$, -2$\varepsilon_{2}$, -2$\varepsilon_{3}$, -$\varepsilon_{2}$-$\varepsilon_{3}$


\\$\Delta(\mathfrak{n})=$$\varepsilon_{1}$+$\varepsilon_{3}$, $\varepsilon_{1}$+$\varepsilon_{2}$, 2$\varepsilon_{3}$, 2$\varepsilon_{2}$, $\varepsilon_{2}$+$\varepsilon_{3}$
\\$\Delta(\mathfrak{n})=$$\varepsilon_{1}$+$\varepsilon_{3}$, $\varepsilon_{1}$+$\varepsilon_{2}$


\\$\Delta(\mathfrak{n})=$2$\varepsilon_{3}$, 2$\varepsilon_{2}$, $\varepsilon_{2}$+$\varepsilon_{3}$
\\$\Delta(\mathfrak{n})=$$\varepsilon_{1}$-$\varepsilon_{2}$, $\varepsilon_{1}$-$\varepsilon_{3}$, -2$\varepsilon_{2}$, -2$\varepsilon_{3}$, -$\varepsilon_{2}$-$\varepsilon_{3}$


\\$\Delta(\mathfrak{n})=$$\varepsilon_{1}$-$\varepsilon_{2}$, $\varepsilon_{1}$-$\varepsilon_{3}$
\\$\Delta(\mathfrak{n})=$-2$\varepsilon_{2}$, -2$\varepsilon_{3}$, -$\varepsilon_{2}$-$\varepsilon_{3}$


\\
\end{tabular}

\\\hline\begin{tabular}{r}$\Delta(\mathfrak{k})$ is of type $A_1$; \\ $\Delta^+(\mathfrak{k})=$ 2$\varepsilon_{1}$\end{tabular} &
\begin{tabular}{l}$\Delta(\mathfrak{n})=$-$\varepsilon_{2}$+$\varepsilon_{3}$, 2$\varepsilon_{3}$, -2$\varepsilon_{2}$, $\varepsilon_{1}$+$\varepsilon_{3}$, -$\varepsilon_{1}$+$\varepsilon_{3}$
\\$\Delta(\mathfrak{n})=$-$\varepsilon_{2}$+$\varepsilon_{3}$, 2$\varepsilon_{3}$, $\varepsilon_{1}$+$\varepsilon_{3}$, -$\varepsilon_{1}$+$\varepsilon_{3}$


\\$\Delta(\mathfrak{n})=$-$\varepsilon_{2}$+$\varepsilon_{3}$, 2$\varepsilon_{3}$
\\$\Delta(\mathfrak{n})=$-$\varepsilon_{2}$+$\varepsilon_{3}$


\\$\Delta(\mathfrak{n})=$2$\varepsilon_{2}$, 2$\varepsilon_{3}$, $\varepsilon_{1}$+$\varepsilon_{2}$, -$\varepsilon_{1}$+$\varepsilon_{2}$
\\$\Delta(\mathfrak{n})=$2$\varepsilon_{2}$, 2$\varepsilon_{3}$


\\$\Delta(\mathfrak{n})=$2$\varepsilon_{2}$, $\varepsilon_{1}$+$\varepsilon_{2}$, -$\varepsilon_{1}$+$\varepsilon_{2}$
\\$\Delta(\mathfrak{n})=$2$\varepsilon_{2}$\quad .


\\\end{tabular}

\\
\multicolumn{2}{c}{Table  \refstepcounter{theo} \label{tableC3B3} \ref{tableC3B3} }
\\
\multicolumn{2}{c}{Non-solvable root subalgebras that fail the cone condition in types $B_3$ and $C_3$.}
\end{longtable}


\begin{comment}
\begin{example} $\gG=\so(8)$, $\Delta(\kk)=\{\pma(\eps_1-\eps_2), \pma\eps_3\pma\eps_4\}$, $\Delta(\nn)= \eps_1+\eps_2$. Then the cone condition is satisfied, but there exists no parabolic subalgebra $\pp$ with $\nn=\nn_\pp$. 
\end{example}
\begin{proof}
\end{proof}
\end{comment}

\section{Root subalgebra preliminaries}\label{secRootSA}
\begin{lemma}\label{leC(kss)}
Let $\kk\subset\gG$ be a reductive root subalgebra. Then $C(\kk_{ss})= \hh_1\oplus$ $\bigoplus_{ \alpha\sperp\Delta(\kk)}$ $\gG^{\alpha}$, where $\hh_1=\{h\in\hh~|~\gamma(h)=0, \forall \gamma \in \Delta( \kk)\}$. %In particular, if two roots are strongly orthogonal to $\Delta(\kk)$, their sum either remains strongly orthogonal to $\Delta(\kk)$ or is not a root.
\end{lemma}
\begin{proof}
Let $x:=h+\sum_{\alpha\in\Delta(\gG)}a_\alpha g^{\alpha}\in C(\kk_{ss})$, where $g^\alpha\in\gG^\alpha$. For any $\gamma\in\Delta(\kk)$ we have $0=[x,g^\gamma] =\gamma(h)g^\gamma+ \sum_{\alpha\in\Delta(\gG)} a_\alpha c_{\alpha \gamma} g^{\alpha +\gamma}$, where $g^{\alpha +\gamma} =0$ if $\alpha+\gamma$ is not a root and $c_{\alpha\gamma}\neq 0$ whenever $\alpha+\gamma$ is a root. Therefore $a_\alpha=0$ for all $\alpha$ which are not strongly orthogonal to $\Delta(\kk)$ and $\gamma(h)=0$. On the other hand, it is clear that when $\gamma(h)=0$ for all $\gamma\in \Delta(\kk)$, and $a_\alpha$ are arbitrary, then $h+\sum_{\alpha\sperp\Delta(\kk)}a_\alpha g^\alpha$ is an element of $C(\kk_{ss})$.
\end{proof}

\begin{prop}\label{lePara}
Suppose that $\lL$ satisfies the cone condition. %$\Cone_{\ZZ}(\hw_{\bb\cap\kk}(\gG/\lL))$ $\cap \Cone_{\ZZ} (\weights(\nn))$ $=\{0\}$. 
Then $C(\kk_{ss})\cap N(C(\kk_{ss})\cap\nn)$ is a parabolic subalgebra of $C(\kk_{ss})$. Equivalently, in view of \refle{leC(kss)}, there exists $h\in\hh$ such that 
\begin{equation}\label{eqPara} 
C(\kk_{ss})\cap N(C(\kk_{ss})\cap\nn)=\mathfrak{q}_h:= \hh_1\oplus\bigoplus_{\substack{\alpha(h)\geq 0\\ \alpha\sperp\Delta(\kk_{ss})}} \gG^{\alpha},
\end{equation} 
where $\hh_1=\{h\in\hh~|~\gamma(h)=0$ for all $ \gamma\in\Delta(\kk)\}$. In addition, $C(\kk_{ss})\cap N(\nn) = C(\kk_{ss})\cap N(C(\kk_{ss}) \cap \nn )$.
\end{prop}
\begin{proof}
Throughout the entire proof we use \refle{leC(kss)}. 

The equality $\Cone_{\ZZ}(\hw_{\bb\cap\kk}(\gG/\lL))\cap\Cone_{\ZZ}(\Delta(\nn))= \{0\}$ implies that there exists $h\in\hh$ for which $\beta(h) >0, \forall \beta\in\Delta(\nn)$ and $\alpha(h)\leq 0, \forall \alpha\in\hw_{\bb\cap\kk}(\gG/\lL)$. Let $\mathfrak{q}_h$ be defined as in \refeq{eqPara}. 

%In order to prove the equality $\mathfrak{q}_h=C(\kk_{ss})\cap N(\nn) = C(\kk_{ss})\cap N(C(\kk_{ss})\cap\nn)$, we will first prove the inclusion $\mathfrak{q}_h\supset \left(C(\kk_{ss})\cap N(C(\kk_{ss}) \cap \nn)\right)$ (and clearly $\left(C(\kk_{ss})\cap N(C(\kk_{ss}) \cap \nn) \right) \supset \left( C(\kk_{ss} ) \cap N(\nn)\right) $). Then we will show the reverse inclusion $\mathfrak{q}_h\subset \left( C(\kk_{ss} ) \cap N(\nn) \right)$; as an intermediate step we will show the weaker inclusion $\mathfrak{q}_h\subset \left(C(\kk_{ss})\cap N(C(\kk_{ss}) \cap \nn)\right)$.


We claim first that $\mathfrak{q}_h\supset \left(C(\kk_{ss})\cap N(C(\kk_{ss})\cap\nn)\right)\supset \left( C(\kk_{ss} ) \cap N(\nn)\right)$. Suppose on the contrary that there exists $x:=\underbrace{g}_{\in\hh_1}$ $+ \sum_{ \alpha \sperp\Delta (\kk_{ss})} a_{\alpha} g^\alpha\in C(\kk_{ss})\cap N(C(\kk_{ss})\cap\nn)$ for which there is a root $\gamma\in \Delta(C(\kk_{ss}))$ such that $\gamma(h)<0$ and $a_\gamma\neq 0$. Then $\hh\subset N(C(\kk_{ss})\cap\nn)$ implies that whenever $a_\alpha\neq 0$ we have $g^\alpha\in N(C(\kk_{ss})\cap\nn)$. In particular $g^\gamma\in N(C(\kk_{ss})\cap\nn)$. As $C(\kk_{ss})$ is reductive, ${-\gamma} \in\Delta (C(\kk_{ss}))$, and $-\gamma(h)>0$ implies $-\gamma\in\Delta(\nn)$. Therefore $\gG^{-\gamma}\in C(\kk_{ss})\cap\nn$ which contradicts the inclusion $\gG^{\gamma}\subset N(C(\kk_{ss})\cap\nn)$.

We claim next that $\mathfrak{q}_h\subset C(\kk_{ss})\cap N(C(\kk_{ss})\cap\nn)$. Fix $\alpha\in \Delta( C( \kk_{ss} ))$ for which $\alpha(h) \geq 0$. If $\beta\in\Delta(C(\kk_{ss})\cap\nn)$ and $(\alpha+\beta)$ is a root, then $(\alpha+\beta)(h)>0$. Therefore $\alpha + \beta \in \Delta(C(\kk_{ss}) \cap\nn)$ as all roots in $\Delta(C(\kk_{ss}))$ are $\bb\cap\kk$-singular. Therefore $g^{\alpha}\in N(C(\kk_{ss})\cap\nn)$. 

So far we have established that $\mathfrak{q}_h= C(\kk_{ss})\cap N(C(\kk_{ss})\cap\nn)$; we are left to prove that $\qq_h \subset C(\kk_{ss}) \cap N(\nn)$. Suppose, on the contrary, that there is $-\alpha\in\Delta(C(\kk_{ss}))$ such that  $-\alpha(h)= 0$ and $\gamma:=-\alpha+\beta\in\Delta(\gG)\backslash\Delta(\nn)$ for some $\beta\in\Delta(\nn)$. Since $-\alpha\sperp\Delta(\kk)$, $-\alpha,\alpha\in\hw_{\bb\cap\kk}(\gG/\lL)$ and we have the relation 
\begin{equation}\label{eqBadRel}
\alpha+\gamma=\beta.
\end{equation}
Clearly $\gamma\notin\Delta(\kk)$. For the already fixed choice of $\alpha$, assume that $\gamma\in\Delta(\gG) \backslash \Delta(\lL)$ is a root maximal with respect to the partial order defined by $\bb\cap \kk$, such that there exists a relation \refeq{eqBadRel} as above. If $\gamma\in \hw_{\bb\cap\kk}(\gG/\lL)$, this would contradict the cone condition; therefore there exists $\delta\in\Delta^+(\kk)$ such that $\delta+\gamma$ is a root. The requirement that $\gamma\in\Delta( C(\kk_{ss}))$ forces $\delta$ to be strongly orthogonal to $\alpha$. Therefore $\delta$ has the same scalar product with $\gamma$ as it does with $\beta$, but at the same time $\delta+\gamma$ is a root and $\delta+\beta$ isn't (due to the maximality of $\delta$). We will prove that these requirements are contradictory. Let the simple component of $\gG$ containing $\alpha$, $\gamma$ and $\beta$ be $\sS$.

 
\begin{itemize}
\item[Case 1] $\sS$ is of type $A$, $D$, $E$ or $G_2$. The inequality $\langle\delta, \beta\rangle= \langle \delta, \gamma \rangle< 0$ contradicts the maximality of $\gamma$ because if it held, we could add $\delta$ on both sides of \refeq{eqBadRel}. The inequality $\langle\delta, \gamma \rangle \geq 0$ implies $\langle\delta, \gamma \rangle = 0$ (the sum of two roots with positive scalar product is never a root). In turn, this contradicts the condition that $\delta+\gamma$ is a root since in root systems of type $A$, $D$, $E$ and $G_2$, strong orthogonality is equivalent to orthogonality.
\item[Case 2] $\sS$ is of type $C$. Without loss of generality we can assume that \refeq{eqBadRel} is $\underbrace{\eps_{j_1}+ \eps_{j_2}}_{\alpha}$ +$\underbrace{ (-\eps_{j_2} + \eps_{j_3})}_{\gamma}$ $= \underbrace{\eps_{j_1}+  \eps_{j_3}}_{\beta}$, where the indices $j_1, j_2, j_3$ are not assumed to be pairwise different. Then $\delta=-\eps_{j_3}+\eps_l$ contradicts the maximality of $\gamma$ for all possible choices of the indices $j_1, j_2, j_3, l$. Furthermore, $\delta=\eps_{j_2}+ \eps_l$ contradicts $\alpha \in \Delta(C(\kk_{ss}))$ for all possible choices of the indices $j_1, j_2, j_3, l$. Contradiction.
\item[Case 3] $\sS$ is of type $B$. 
\begin{itemize}
\item[Case 3.1] $\alpha$ and $\gamma$ are both short. Without loss of generality \refeq{eqBadRel} becomes $\underbrace{\eps_{1}}_{\alpha}$ $+ \underbrace{\eps_{2}}_{\gamma}$ $= \underbrace{\eps_{1} + \eps_{2}}_{\beta}$. The maximality of $\gamma$ implies $\delta=\eps_{1}-\eps_{2}$ which contradicts $\alpha\in\Delta(C(\kk_{ss}))$.
\item[Case 3.2] $\alpha$ is short and $\gamma$ is long. Without loss of generality \refeq{eqBadRel} becomes $\underbrace{\eps_{1}}_{\alpha} + \underbrace{(-\eps_{1}+\eps_{2})}_{\gamma}$ $= \underbrace{\eps_{2}}_{\beta}$. The maximality of $\gamma$ implies $\delta =\eps_{1} \pm\eps_l$ for some index $l$, which contradicts $\alpha\in\Delta(C(\kk_{ss}))$.
\item[Case 3.3] $\alpha$ is long and $\gamma$ is short. Without loss of generality \refeq{eqBadRel} becomes $\underbrace{\eps_{1}+\eps_{2}}_{\alpha}$ +$\underbrace{(-\eps_{2})}_{\gamma} $ $= \underbrace{ \eps_{1} }_{ \beta} $. Thus $\delta=$ $\eps_{2}$ $+\eps_l$ for some index $l$ and $\alpha \in \Delta( C( \kk_{ ss}))$ implies $\delta=\eps_{2}-(\eps_{1})$. Then $\beta + \beta$ $+\delta$ $=\alpha$ yields a contradiction.
\item[Case 3.4] Both $\alpha$ and $\gamma$ are long. Without loss of generality \refeq{eqBadRel} becomes
$\underbrace{ \eps_{1}- \eps_{2}}_{\alpha}$ $+\underbrace{ \eps_{2} +\eps_{3}}_{\gamma}$ $= \underbrace{\eps_{1} +\eps_{3}}_{\beta}$. The assumption that $\delta$ is short contradicts either $\alpha\in \Delta(C(\kk_{ss}))$ or the maximality of the choice of $\gamma$. The fact that all roots participating in \refeq{eqBadRel} together with the root $\delta$ are long is contradictory. Indeed, otherwise we could use the exact same data to obtain a relation (\ref{eqBadRel}) in type $D$.
\end{itemize}
\item[Case 4] $\sS$ is of type $F_4$. Suppose on the contrary that there exist roots $\alpha,\beta,\gamma,\delta$ for which \refeq{eqBadRel} holds and $\delta\sperp\alpha$, $\delta\sperp\beta$, $\delta\perp\gamma$. The same conditions would continue to hold in the root subsystem $\Delta '\supset \beta$ generated by $\alpha, \gamma, \delta $. Since $\Delta'$ is of rank 3, setting $\Delta(\nn'):=\{\beta\}$, $\Delta(\kk_{ss}') =\{\pma\delta\}$ we get data whose existence we proved impossible in the preceding cases. Contradiction.
\end{itemize}
\end{proof}


\section{Sufficiency of cone and centralizer conditions for finite type}
\begin{lemma}
Let $\lL$ be a root subalgebra of the semisimple Lie algebra $\gG$ such that the cone conditions holds. Then $\lL$ is a Fernando-Kac subalgebra of finite type if and only if the centralizer condition holds.
\end{lemma}
\begin{proof}
Suppose $\Cone_{\ZZ} (\hw_{\bb\cap\kk}(\gG/\lL))\cap \Cone_{\ZZ}(\Delta(\nn))= \{0\}$ but $[C(\kk_{ss})\cap N(\nn)]$ has a Levi subalgebra that has a simple component of type $B$, $D$, or $E$. Let $h\in\hh$ be such that $\gamma(h)>0$ for all $\gamma \in \Cone_{\ZZ}(\Delta(\nn))$ and $\gamma(h)\leq 0$ for all $\gamma\in\hw_{\bb\cap\kk}(\gG/\lL)$. According to Proposition \ref{lePara}, $[C(\kk_{ss}) \cap N(\nn)] =[C(\kk_{ss})\cap N(C(\kk_{ss})\cap\nn)]=\qq_h$, where $\qq_h$ is defined as in \refle{lePara}. Assume on the contrary that there exists an irreducible $(\gG,\lL)$-module with $\gG[M]=\lL$. Pick an arbitrary $\bb\cap\lL$-singular vector $v$ and consider the $\qq_h\cap C(\kk_{ss})_{ss}$-module $N$ generated by $v$. We have that $N$ is a strict $(\qq_h\cap C(\kk_{ss})_{ss}, \hh\cap C(\kk_{ss})_{ss})$-module (``torsion-free'' according to the terminology of \cite{Fernando1}). Then, according to \cite[Theorem 5.2]{Fernando1}, it cannot have finite-dimensional $\hh\cap C(\kk_{ss})_{ss}$-weight spaces. In particular there are infinitely many $u_1,\dots\in U(C(\kk_{ss}))$ such that $u_1\cdot v,\dots$ are linearly independent and of same $\hh$-weight. Then $u_i\in U(C(\kk_{ss}))$ implies $u_1\cdot v,\dots$ are all $\bb\cap\kk$-singular, which contradicts the fact that $M$ is of finite type over $\kk$ (\cite[Theorem 3.1]{PSZ}). 

%\textbf{Remark.} The above paragraph can be reformulated using that $\pp_h$ is precisely the opposite parabolic subalgebra to the parabolic subalgebra of $\cc:=C(\kk_{ss})$ from \cite[Prop. 4]{PS1}, whose reductive part according to the latter must have simple components of type $A$ and $C$ only.

Suppose next $\Cone_{\ZZ}(\hw_{\bb\cap\kk}(\gG/\lL))\cap \Cone_{\ZZ}(\Delta(\nn))=\{0\}$ and $C(\kk_{ss})\cap N(\nn)$ has simple Levi components of type $A$ and $C$ only.  We will prove that $\lL$ is Fernando-Kac subalgebra of finite type by the construction \cite[Theorem 4.3]{PSZ}. Since the cones do not intersect, there exists a hyperplane in $\hh^*$ given by an element $h\in\hh$ such that $\Delta(\nn))$ lies in the $h$-strictly positive half-space, and $\Cone_{\ZZ}(\hw_{\bb\cap\kk}(\gG/\lL))$ lies in the $h$-non-positive half-space. Clearly we can assume $h$ to have rational action on $\hh^*$.

We introduce now a ``small perturbation'' procedure for $h$ to produce an element $h'$  such that $\gamma(h')\neq 0$ for all $\gamma\in\Delta(\kk)$. Suppose $\gamma\in \Delta(\bb\cap\kk)$ is a root with $\gamma(h)=0$. Define $g\in\hh$ by the properties $\gamma(g)=1$, $\gamma'(g)=0$ for all $\gamma'\perp\gamma$. Now choose $t$ to be a sufficiently small positive rational number ($t\leq\half \min_{\beta \in \Delta(\gG), \beta(h)\neq 0}{|\beta(h)|}$ serves our purpose). Set $h_{1}:=h- tg$. Then all $h$-positive (respectively $h$-negative) vectors remain $h_1$-positive (respectively $h_1$-negative) vectors. The only roots $\alpha$ whose positivity would be affected by the change are those with $\alpha(h)=0$, $\langle\alpha,\gamma\rangle\neq 0$. By the preceding remarks, $\Delta(\nn)$ lies in the $h_1$-positive half-space. We show next that $\Cone_{\ZZ} (\hw_{\bb\cap\kk} (\gG/\lL))$ remains in the $h_1$-non-positive half-space. Suppose on the contrary we had a vector $\alpha\in \hw_{\bb \cap \kk}( \gG/ \lL)$ that now lies in the $h_1$-positive half space. By the preceding remarks $\alpha(h)=0$. Therefore $\alpha(g)= -\frac{1}{t}\alpha(h_1)<0$ which implies $\langle\alpha,\gamma\rangle <0$ and thus $\alpha+\gamma$ is a root. Contradiction. 

If there is a root of $\kk$ that still has zero action on the newly obtained $h_1$, we apply the above procedure again, and so on. The number of roots $\alpha\in\Delta(\kk)$ for which $\alpha(h_1)=0$ is smaller than the corresponding number for $h$. Therefore after finitely many iterations we will obtain an element, say $h'$, for which $\gamma(h')\neq 0$ for all $\gamma\in\Delta(\kk)$ and
\begin{equation}\label{eqTheCondition2}
\begin{array}{c}
\alpha(h')=0 \mathrm{~for~all~} \alpha \mathrm{~for~which~}\gG^\alpha\in C(\kk_{ss}).
\end{array}
\end{equation}

%It does not have to be true however that the entire $\Cone_{\ZZ}(\hw_{\bb\cap\kk}(\gG/\lL))$ remains in the $h_1$-non-positive half-space. Suppose we had a vector $\alpha\in\hw_{\bb\cap\kk}(\gG/\lL)$ that now lies in the $h_1$-positive half space. By the preceding remarks $\alpha(h)=0$. Therefore $\alpha(g)= -\frac{1}{t}\alpha(h_1)<0$ which implies $\langle\alpha,\gamma\rangle>0$ and thus $\alpha-\gamma$ is a root. Now $\alpha-\gamma$ is a $h_1$-negative root, for which $\gG^{\alpha-\gamma}/\lL$ lies in the $\kk$-submodule of $\gG/\lL$ generated by $\gG^{\alpha}/\lL$. In other words, if our perturbation $h_1$ of $h$ would make a root $\alpha\in\hw_{\bb\cap\kk}(\gG/\lL)$ lie in the $h_1$-positive half-space, it would also make another root from the same $\kk$-submodule as $\gG^\alpha/\lL$ be $h_1$-negative.

%We can repeat inductively the procedure described in the preceding paragraph to obtain in the end a functional $h'$ for which $\gamma(h')\neq 0$ for all $\gamma\in\Delta(\kk)$, with the additional key properties that
%\begin{equation}\label{eqTheCondition}
%\begin{array}{c}
%\mathrm{if~}\alpha(h')>0 \mathrm{~for~some~} \alpha\in\Delta(\gG/\lL),\mathrm{~then~there~exists~a~root~} \alpha'\in\Delta(\gG/\lL) \\ \mathrm{~such~that~} \gG^{\alpha'}/\lL \mathrm{~lies~in~the~}\kk-\mathrm{submodule} \mathrm{~generated~by~} \gG^{\alpha}/\lL \mathrm{~and~such~that~} \alpha'(h')<0
%\end{array}
%\end{equation}
%and
%\begin{equation}\label{eqTheCondition2}
%\begin{array}{c}
%\alpha(h')=0 \mathrm{~for~all~} \alpha \mathrm{~for~which~}\gG^\alpha\in C(\kk).
%\end{array}
%\end{equation}
Now define 
\[
\pp:=\bigoplus_{\alpha(h')\geq 0}\gG^{\alpha}, \pp_h:=\bigoplus_{\alpha(h)\geq 0}\gG^{\alpha}.
\]
Then $\pp_{red}\subset{(\pp_h)}_{red}=\hh+\qq_h$, where $\qq_h$ is the subalgebra defined in \refle{lePara}. By \refle{lePara}, we get ${\qq_h}=C(\kk_{ss})\cap N(\nn)$ and the latter is direct sum of simple components of type $A$ and $C$ by the centralizer condition. Thus $\pp_{red}$ is a sum of root systems of type $A$ and $C$ (since types $A$ and $C$ contain root subsystems of type $A$ and $C$ only). We can now pick a $(\pp_{ red} , \hh)$-module $L$  for which $\pp_{red}[L]=\hh$ (see \cite{BrittenLemire}, \cite[Sections 8,9]{Mathieu}), and we can extend $L$ to a $\pp$-module by choosing trivial action of the nilradical of $\pp$. The choice of $h'$ allows us to apply \cite[Theorem 4.3]{PSZ} to get a $\gG$-module $M$ for which $\gG[M]$ is the sum of $\kk$ and the maximal $\kk$-stable subspace of $\pp[L]=\hh\crplus\nn_\pp$. The fact that at least one weight of each irreducible direct summand of $\gG/\lL$ (namely, its $\bb\cap\kk$-singular weight) is outside of $\pp[L]$ implies that the maximal $\kk$-stable subspace of $\pp[L]$ is $\nn$. This completes the proof.
\end{proof}


The preceding proof completes the existence part in the proof of Penkov's conjecture. The remainder of the thesis is dedicated to proving that the failure of the cone condition implies that $\lL$ is Fernando-Kac subalgebra of infinite type.


\chapter{Cone condition fails $\Rightarrow$ $\lL$ is Fernando-Kac subalgebra of infinite type}\label{chapterConeConditionFails-NoFKFT}
Section \ref{secConeIntersections} gives a sufficient condition for a root subalgebra to be Fernando-Kac subalgebra of infinite type. More precisely, we show that the existence of a certain $\lL$-infinite weight implies that $\lL$ is Fernando-Kac subalgebra of infinite type. Section \ref{secBigTedious} proves that the failure of the cone condition is equivalent to the existence of an $\lL$-infinite weight. 
\section{A sufficient condition for infinite type}
\label{secConeIntersections}
Let $\lL=\kk\crplus\nn$ be a root subalgebra, let $M$ be a $(\gG,\lL)$-module. For every root $\alpha\in\Delta(\gG)$ choose a non-zero vector $g^\alpha\in\gG^\alpha$ such that $[g^{\alpha},g^{-\alpha}]=h^\alpha$, where $h^{\alpha}$ is the element of $\hh$ for which $[h^{\alpha},g^\beta]=\langle\alpha,\beta \rangle g^{\beta}$ for all $\beta\in\Delta(\gG)$.

By Lie's theorem, there exists an $\bb\cap\lL$-singular vector $v$ in $M$. Suppose that there exist roots $\alpha_i\in\hw_{\bb\cap\kk}(\gG/\lL)$ and $\beta_i\in\Delta(\nn)$, as well as numbers $a_i,b_j\in\ZZ_{>0}$, such that the vectors of the form $\left((g^{-\beta_1})^{b_1}\dots (g^{-\beta_k})^{b_k}\right)^t$ $\left((g^{\alpha_1})^{a_1}\dots (g^{\alpha_l})^{a_l}\right)^t$  $\cdot v $ for $t\in\ZZ_{>0}$ have the following three properties. First, these vectors have the same $\hh$-weight; second, they are linearly independent; third, each of them projects naturally to a $\bb\cap\lL$-singular vector in an appropriate $\lL$-subquotient of $M$. If all three properties hold, then $M$ is a $(\gG,\lL)$-module of infinite type as the irreducible $\lL$-module with highest weight equal to the weight of $v$ has infinite multiplicity in $M$.

The above summarizes our approach for proving that the failure of the cone condition implies $\lL$ is a Fernando-Kac subalgebra of infinite type. 
The present section establishes that the three properties in question hold under an additional assumption.

\begin{defi}\label{defStronglyPerp}\label{defKmult}~
\begin{itemize}
\item Let $I$ be a set of roots and $\omega$ be a weight. We say that $\omega$ \emph{has a strongly orthogonal decomposition with respect to $I$} if there exist roots  $\beta_i\in I$ and positive integers $b_i$ such that $\omega=b_1\beta_1+\dots +b_k\beta_k$ and $\beta_i\sperp\beta_j$ for all $i,j$. 
\item Fix $\lL=\kk\crplus\nn\subset\gG$ with $\nn\subset\bb$. Let $\omega$ be a weight. We say that $\omega$ is \emph{two-sided with respect to $\lL$}, or simply \emph{two-sided}, if the following two conditions hold:
\begin{itemize}
\item $\omega\in \Cone_{\ZZ}(\hw_{\bb\cap\kk}(\gG/\lL))\cap\Cone_{\ZZ}(\Delta(\nn))\backslash\{0\}$, i.e. there exist $a_i\in\ZZ_{>0}$, $\alpha_i \in \hw_{\bb\cap\kk} (\gG/\lL)$, $b_i\in\ZZ_{>0}$ and $\beta_i\in\Delta(\nn)$ with 
\begin{equation}\label{eqRelation}
\omega=\sum_{i=1}^la_i\alpha_i=\sum_{i=1}^k b_i\beta_i;
\end{equation}
\item among all expressions for $\omega$ of type \refeq{eqRelation}, there exists one for which $[\gG^{\alpha_1},\nn]\subset\nn,\dots, [\gG^{\alpha_l},\nn]\subset\nn$. 
\end{itemize}
%An expression of type \refeq{eqRelation} that satisfies the second requirement will be called a \emph{generating relation}.
\item Let $\omega$ be a weight. If $\omega$ is both two-sided and has a strongly orthogonal decomposition with respect to $\Delta(\nn)$, we say that $\omega$ is \emph{$\lL$-strictly infinite}.
\item If for a given weight $\omega$ there exists a root subalgebra $\tT$ containing $\kk$, such that $\omega$ is $\lL'$-strictly infinite in $\tT$, where $\lL':=\lL\cap\tT=\kk\crplus(\tT\cap\nn)$, we say that $\omega$ is \emph{$\lL$-infinite}.
\end{itemize}
\end{defi}


\label{secMultiplicities}

%The following lemma is a well known fact and we leave its proof to the reader.
%\begin{lemma}
%Let $\kk\crplus\nn$ be a root subalgebra of $\gG$. Then there exists a linear functional $h\in\hh$ on $\Delta(\gG)$ such that $\gamma(h)=0$ for all $\gamma\in\Delta(\kk)$ and $\beta(\hh)>0$ for all $\beta\in\Delta(\nn)$. 
%\end{lemma}

\begin{lemma}\label{leNilradicalIsSeparate}
Given $\lL=\kk\crplus\nn\subset\gG$, there exists $h\in\hh$ such that $\gamma(h)=0$ for all $\gamma\in\Delta(\kk)$ and $\beta(h)>0$ for all $\beta\in\Delta(\nn)$.
\end{lemma}
\begin{proof}
Since $\hh\crplus\nn$ is a solvable Lie algebra, it lies in a maximal solvable (i.e. Borel) subalgebra; assume without loss of generality that this Borel subalgebra is $\bb$. Fix $h'\in\hh$ such that $\gamma(h')>0$ for all $\gamma\in \Delta^+(\gG)$. Let $h''\in\hh$ be defined by $\gamma(h''):=\gamma(h')$ for all $\gamma\in\Delta(\kk)$ and $\alpha(h'')=0$ for all weights $\alpha \in \Delta(\kk)^\perp$. Set $h:=h'-h''$. 

We claim that $h$ has the properties stated in the lemma. Indeed, let $\nn'\subset\nn$ be a $\kk$-submodule of $\nn$. Since $\gamma(h)=0$ for all $\alpha \in \Delta(\kk)$, the value $r:=\beta(h)$ is the same for all roots $\beta\in\Delta(\nn')$. Our statement is now equivalent to showing that $r>0$. Assume on the contrary that $r\leq 0$. Let the sum of the weights of $\Delta(\nn')$ be  $\lambda$, i.e. $\lambda:=\sum_{ \beta_i \in \Delta(\nn')}\beta_i$. Then $\lambda(h)= \#(\Delta(\nn'))r\leq 0$. On the other hand, the sum of the weights of a finite-dimensional $\kk_{ss}$-module always equals zero, i.e. $\lambda\in \Delta(\kk)^\perp$. Therefore $\lambda(h)=\lambda(h')+\lambda(h'')=\lambda(h')>0$, contradiction.
\end{proof}

For an arbitrary weight $\mu\in\hh^*$, denote by $L_{\mu}(\kk)$ (respectively, $L_{\mu}(\gG)$) the irreducible highest weight $\kk$-module (respectively, $\gG$-module) with $\bb\cap\kk$-highest (respectively, $\bb$-highest) weight $\mu$.
\begin{lemma}\label{leMultMain}
%Let $\lL=\kk\crplus\nn$ be a root subalgebra of a reductive Lie algebra $\gG$ with a Borel subalgebra $\bb$ such that $\nn\subset\bb$. 
Let $\omega$ be a weight that has a strongly orthogonal decomposition with respect to $\Delta(\nn)$. Let $\lambda\in\hh^*$ be an arbitrary $\bb\cap\kk$-dominant and $\kk$-integral weight. Then there exists a number $t_0$ such that, for any $t>t_0$ and any $\gG$-module $M$ that has a $(\bb\cap\kk)\crplus\nn$-singular vector $v$ of weight $\lambda+t\omega$, it follows that $M$ has a $\kk$-subquotient in which there is a non-zero $\bb\cap\kk$-singular vector $\tilde w$ of weight $\lambda$.
\end{lemma}
Before we proceed with the proof we state the following.
\begin{corollary}
\label{corMultMain}
Let $\gG$, $\bb$, $\lL=\kk\crplus\nn$, $\lambda$ and $\omega$ be as above. Then there exists $t_0$ such that for any $t>t_0$ and any $(\gG,\kk)$-module that has a $(\bb\cap\kk)\crplus\nn$-singular vector $v\in M$ of weight $\lambda+t\omega$, it follows that $M$ has non-zero multiplicity of $L_\lambda(\kk)$. In particular, the existence of a $(\gG,\kk)$-module with the required singular vector implies that $\omega$ is $\bb\cap\kk$-dominant.
\end{corollary}
\textbf{Proof of \refle{leMultMain}.}Let $\nn^-$ be the subalgebra generated by the root spaces opposite to the root spaces of $\nn$. Let $b_1\beta_1+\dots +b_k\beta_k=\omega$ be a strongly orthogonal decomposition of $\omega$ with respect to $\Delta(\nn)$ (\refdef{defKmult}). Let $u:=(g^{\beta_1 })^{b_1}\dots (g^{ \beta_k})^{b_k} \in U(\nn)$ and $\bar u:= (g^{-\beta_1 })^{b_1}\dots (g^{ -\beta_k})^{b_k}$. 

Let $A$ be the linear subspace of $U(\nn^-)$ generated by all possible monomials $g^{-\gamma_1}\dots g^{-\gamma_k}$ that have strictly higher weight than $-t\omega$, where $-\gamma_i\in \Delta(\nn^-)$, in other words, $A:=\linspan \{g^{-\gamma_1}\dots g^{-\gamma_k}|\gamma_i\in\Delta(\nn), \sum\gamma_i\prec\omega\}$. Denote by $N$ the $\kk$-module generated by the vectors $\{A\cdot v\}$. To prove the lemma, we will show that the $\kk$-module $M/N$ has $\tilde w$  as a $\bb\cap\kk$-singular weight vector, where $\tilde w$ is the image in $M/N$ of $w:=\bar u^{t}\cdot v$. 

First, we will prove that $\tilde w$ is $\bb\cap\kk$-singular: indeed, $\nn^-$ is an ideal in the Lie subalgebra $\kk\crplus\nn^-$ and so $g^{\alpha} \bar u^t \in\left( \bar u^t g^{\alpha}+A\right)$ for all $\alpha \in \Delta^+(\kk)$; this, together with the fact that $v$ is $\bb\cap\kk$-singular, proves our claim. Second, we will prove that if $w$ is non-zero, then $w\notin N$ and therefore $\tilde w$ is non-zero. Indeed, the weight spaces of $N$ are a subset of the set 
\[
X:=\bigcup_{\substack{ \gamma\in\Cone_{\ZZ}(\Delta(\nn^-)) \\ \gamma\succ-t\omega }}(\lambda+t\omega+\gamma+ \linspan_\ZZ{ \Delta( \kk)} ).
\]
We claim that $X$ does not contain $\lambda$: indeed, choose $l\in h$ such that $\gamma(l)=0$ for all $\gamma\in\Delta(\kk)$ and $\beta'(l)>0$ for all $\beta'\in\Delta(\nn)$ (Lemma \ref{leNilradicalIsSeparate}). Therefore $-t\omega(l)\notin \{\mu(l)| \mu\in X\}$ and our claim is established.

To finish the proof of the lemma we are left to show that $\omega=\bar u^{t}\cdot v$ is non-zero, and this is the first and only place we will use the strongly orthogonal decomposition of $\omega$. To do that we will prove by direct computation that the vector $u^{t}\bar u^{t}\cdot v$ is a strictly positive multiple of $v$. For any $n\in\ZZ_{>0}$ we compute
\begin{eqnarray*}
(g^{\beta_i}) (g^{-\beta_i})^{n} \cdot v&=& \sum_{j=0}^{n-1} \langle\beta_i,-j\beta_i+\lambda+t\beta \rangle (g^{-\beta_i})^{n-1}\cdot v\\
&=&\left(nt^2\langle\beta_i,\beta\rangle+nt\langle\beta_i,\lambda\rangle - \frac{n(n-1)}{2}\langle\beta_i,\beta_i\rangle \right)(g^{-\beta_i})^{n-1}\cdot v\\
&=& \left(\langle\beta_i, \beta_i\rangle\left(b_int^2-\frac{n(n-1)}{2}\right)+ \langle\beta_i,\lambda\rangle\right) (g^{ - \beta_i })^{n-1}\cdot v .
\end{eqnarray*}
Therefore
\begin{eqnarray*}
(g^{\beta_i})^{t} (g^{-\beta_i})^{t}\cdot v&=&\prod_{k=0}^{t-1}\Bigg(  \langle\beta_i, \beta_i \rangle \left(b_i(t-k)t^2-\frac{(t-k)(t-k-1)}{2}\right)\\ &&~~~~~~~~~~+ \langle\beta_i,\lambda\rangle ~ \Bigg)\cdot v.
\end{eqnarray*}
Define $c_i(t,\lambda)$ to be the above computed coefficient of $ v$, in other words, set $ c_i(t,\lambda) v:= (g^{\beta_i})^{t-s} (g^{-\beta_i})^{t-s}\cdot v$. Since $b_i>0$, using the explicit form of $c_i(t,\lambda)$, we see that for a fixed $\lambda$, $c_i(t,\lambda)>0$ for all large enough $t$. Using that $g^{\pma\beta_i}$ and $g^{\pma\beta_j}$ commute whenever $i\neq j$, we get immediately that $u^{t}\bar u^{t}\cdot v=\prod_i c_i(t,\lambda) v$, which proves our claim that $u^{t}\bar u^{t}\cdot v$ is  a positive multiple of $v$. Therefore $\bar u^{t}\cdot v$ cannot be zero, which completes the proof of the lemma.
$\Box$

\begin{example} 
Let us illustrate Lemma \ref{leMultMain} in the case when $\gG\simeq \sL(3)$ and $M$ is an irreducible $(\gG,\lL)$-module of finite type. Consider first the case $\kk=\hh$. If $\nn=\{0\}$, the statement of the Lemma is a tautology. If $\nn\neq \{0\}$, the lemma asserts that a certain weight space of $M$ is non-zero. As the $\hh$-characters of all simple $\sL(3)$-modules of finite type are known (see for instance, \cite[Section 7]{Mathieu}), the claim of the lemma is a direct corollary of this result. 

The only other possibility for $\kk\neq\gG$ is $\kk\simeq \sL(2)+\hh$. Then there are 2 options for $\lL$: $\lL=\kk$ or $\lL=\kk\crplus \nn$, where $\dim \nn=2$. For $\lL=\kk$ the lemma is a tautology as $\nn=\{0\}$. Consider the case when $\lL=\kk\crplus\nn$ with $\dim \nn=2$, i.e. the case when $\lL$ is a parabolic subalgebra with Levi component isomorphic to $\sL(2)$. Here, there are two options for $M$: $\dim M<\infty$, and $\dim M=\infty$. In both cases $\kk$ acts semisimply on $M$ and the lemma asserts the existence of certain $\bb\cap\kk$ singular vector in $M$. More precisely, let $\gamma_1:=\eps_1-\eps_2, \gamma_2 :=\eps_2 - \eps_3$ be the positive simple basis of $\Delta(\gG)$ with respect to $\bb$, and let $\Delta(\kk)=\{\pm\gamma_1\}$. Then Lemma \ref{leMultMain} claims that if $\lambda$ is $\bb\cap\kk$-dominant and integral, then $L_{\lambda+t\omega}(\gG)$ has a $\bb\cap\kk$-singular vector of weight $\lambda$ for all large enough $t$. Up to multiplication by a positive integer, there are two different options for picking the weight $\omega$ - either  $\omega=\gamma_1+\gamma_3$ or $\omega=\gamma_2+\gamma_3$.

\begin{itemize}
\item Suppose $\omega=\gamma_1+\gamma_3$.  Let $x(t)$ and $y(t)$ be functions of $t$ and $\lambda$, defined by $\lambda+t\omega = \frac{x(t)}{3} ( 2\gamma_1 + \gamma_2) +\frac{y(t)}{3} (\gamma_1 + 2 \gamma_2) $. The requirement that $\lambda$ be $\bb\cap\kk$-dominant forces $ t\leq x(t) $. Then the lemma states that there exists a constant $t_0$, such that for all $t_0\leq t\leq x(t)$ we have that $L_{\lambda+t\omega}(\gG)$ has a $\bb\cap\kk$-singular vector of weight $\lambda$. The reader can verify that for both infinite and finite-dimensional $M$, that the constant $t_0$ can be chosen to be zero.

\item Suppose $\omega=\gamma_2$. Then the lemma states that there exists a constant $t_0$, such that for all $t\geq t_0$ we have that $L_{\lambda+t\omega}(\gG)$ has a $\bb\cap\kk$-singular vector of weight $\lambda$. As the reader can verify, when $\dim M=\infty$, the statement of the Lemma holds for $t_0=0$; in the case that $M$ is finite-dimensional, one must pick $t_0\geq -\langle\lambda,\gamma_2\rangle$.
\end{itemize}
 
\end{example}

\begin{lemma}\label{leNotFK}
%Let $\gG$ be a reductive Lie algebra and $\lL=\kk\crplus\nn$ be a root subalgebra. 
Suppose there exists an $\lL$-strictly infinite weight $\omega$. 
\begin{itemize}
\item[(a)] Any $(\gG,\lL)$-module $M$ for which any element in $\gG\backslash\lL$ acts freely is of infinite type over $\lL$.
\item[(b)] $\lL=\kk\crplus\nn$ is a Fernando-Kac subalgebra of infinite type. 
\end{itemize}
\end{lemma}
\begin{proof} 
As any irreducible strict $(\gG,\lL)$-module satisfies the conditions of (a), (a) implies (b); we will now show (a). Let $v_\lambda$ be a $(\bb\cap \lL)$-singular vector. 

Let $\gG$ be the Lie subalgebra generated by $\kk$ and $\gG^{\pma\beta_i}$, where $\omega:=\sum_{i=1}^l a_i\alpha_i= \sum_{i=1}^k b_i\beta_i$ is one decomposition \refeq{eqRelation}. Let $u^{ \bar \alpha}:= (g^{\alpha_1})^{a_1} \dots (g^{\alpha_l})^{a_l}\in U(\gG)$. The vector $v_{\lambda+ t\omega}:= (u^{ \bar\alpha} )^t\cdot v_\lambda$ is non-zero by the conditions of (a). We claim that $v_{\lambda+t\omega}$ is $\bb\cap \lL$-singular. Indeed, first note that since all $\alpha_i$ are $\kk \cap \bb$-singular, $v_{\lambda+t\omega}$ is $\bb\cap \kk$-singular. Second, let $g^{\beta} \in \gG^{\beta} \subset \nn$. By the second requirement for being two-sided we can commute $g^\beta$ with $u^{\bar \alpha}$ to obtain that $g^{\beta}(u^{\bar\alpha})^n \in U(\qq)  a$, where $a\in U(\nn)$ is an element with no constant term and $\qq$ is the Lie subalgebra generated by $\gG^{\alpha_1},\dots, \gG^{ \alpha_k }$. Since $a\cdot v_{\lambda}=0$, we get $g^{\beta}\cdot v_{\lambda+t\omega}=0$, which proves our claim.

All $v_{\lambda+t\omega}$ are linearly independent since they have pairwise non-coinciding weights. Let $M_t$ be the $\gG$-submodule of $M$ generated by $v_{\lambda+t\omega}$ and let $M'$ be the sum of the $M_t$'s as $t$ runs over the non-negative integers. \refcor{corMultMain} shows that the $\kk$-module $L_{\lambda }( \kk )$ has non-zero multiplicity in $M_t$ for all large enough $t$. Consider the vectors $\bar u^t\cdot v_{\lambda+t\omega}$ generating the $\kk$-subquotients isomorphic to $L_{\lambda }( \kk )$, where $\bar u$ is defined as in the proof of \refle{leMultMain}. Let $A_t$ be the linear subspace of $U(\nn^-)$ generated by all possible monomials $g^{-\gamma_1}\dots g^{-\gamma_k}$ that have strictly higher weight than $-t\omega$, where $\gamma_i\in\Delta(\nn)$. Let $N$ be the $\kk$-submodule generated by the vectors $\displaystyle \bigcup_{t> t_0}A_t \cdot v_{\lambda+t\omega}$, where $t_0$ is the number given by Lemma \ref{leMultMain}. Just as in the proof of \refle{leMultMain} we see that each vector $\bar u^t\cdot v_{\lambda+t\omega}$ is not in $N$ and is the image of a $\bb\cap\kk$-singular vector in the quotient $ M'/N $. 

We will now prove that $\bar u^t\cdot v_{\lambda+t\omega}$ are linearly independent. Indeed, let $u$ be defined as in the proof of \refle{leMultMain}. Now take a linear dependence $0=\sum_{i=1}^N c_i \bar u^{t_i}\cdot v_{\lambda+t_i\omega}$ such that $t_N\geq t_1, \dots, t_N\geq  t_{N-1}$ and apply $u^{t_N}$ to both sides. As the computations in the proof of \refle{leMultMain} show, $u^{t_N}$ kills all but the last summand; therefore the last summand has coefficient $c_N=0$. Arguing in a similar fashion for the remaining summands, we conclude that the starting linear dependence is trivial. This shows that the $\kk$-module $L_{\lambda }( \kk )$ has infinite multiplicity in the $\kk$-module $M'$. We conclude that $M$ has infinite type over $\kk$, hence, by \cite[Theorem 3.1]{PSZ}, $M$ has also infinite type over $\lL$. 
\end{proof}
%Let us consider all possible irreducible weight modules of $\sL(3)$ with finite-dimensional weight spaces. According to \cite[Chapter 8]{Mathieu}, the possible weight supports are all possible translates of a fixed weight $\omega$

\section{Existence of $\lL$-infinite weights}\label{secBigTedious}
\subsection{Existence of two-sided weights}\label{secTwoSided}

\begin{lemma}\label{lestreakOfRoots} Let $\alpha_1,\dots,\alpha_k,\alpha_{k+1},\gamma$ be vectors of a root system such that $\alpha_1+ \dots+ \alpha_k + \alpha_{k+1}+\gamma$ is a root different from $\gamma$ or is equal to zero, and $\alpha_i+\gamma$ is neither a root nor zero for $i=1,\dots,k$. Then $\alpha_{k+1}+\gamma$ is a root or zero.
\end{lemma}
\begin{proof}
We will establish the lemma only for an irreducible root system; the case of a reducible root system is an immediate corollary which we leave to the reader. For $G_2$ the statement is a straightforward check, so assume in addition that the root system is not of type $G_2$. 

Assume the contrary to the statement of the lemma. Let 
\begin{equation}\label{eqLeRootTricks}
\alpha_1+\dots+\alpha_k+\alpha_{k+1}+\gamma=\delta.
\end{equation} 
Then $\langle\beta,\gamma\rangle\geq 0$, $\langle\alpha_i,\gamma\rangle\geq 0$. Apply $\langle \bullet,\gamma\rangle$ to both sides of (\ref{eqLeRootTricks}). We get $\langle\alpha_1, \gamma\rangle+ \dots+\langle \alpha_k , \gamma \rangle+\langle\alpha_{k+1}, \gamma \rangle + \langle \gamma, \gamma \rangle = \langle \delta,\gamma\rangle$. If $\delta=0$, we immediately get that $\langle\gamma,\alpha_i\rangle<0$ for some $i$ and the statement of the lemma holds as the sum of two roots with negative scalar product is always a root. Therefore we can suppose until the end of the proof that $\delta\neq 0$. 

Since $\langle\alpha_i,\gamma\rangle\geq 0$ and $\delta\neq \gamma$, we must have $\langle \gamma, \alpha_1 \rangle = \dots= \langle \gamma, \alpha_k \rangle= \langle \gamma, \alpha_{k+1} \rangle= 0$ and $\langle \gamma, \gamma \rangle= \langle \delta, \gamma \rangle $. Since $\delta \neq \gamma$ by the conditions of the lemma, the only way for this to happen is to have that $\gamma$ is a short and $\delta$ is long, which gives the desired contradiction in types $A$, $D$ and $E$. Suppose now the given root system is of type $C$. Then without loss of generality we can assume that $\delta = 2\eps_1$ and $\gamma= \eps_1 + \eps_2$. But then there must be a summand on the left-hand side of (\ref{eqLeRootTricks}) which cancels the $+\eps_2$ term of $\gamma$. None of the $\alpha_i$'s have a $-\eps_2$ term (since $\alpha_i+\gamma$ is not a root) and therefore $\alpha_{k+1}+\gamma$ is a root, contradiction. Suppose next that the given root system is of type $B$. Then without loss of generality $\gamma$ can be assumed to be $\eps_1$ and $\delta$ to be $\eps_1+\eps_2$. Clearly  $\alpha_1+ \dots+\alpha_k +\alpha_{k+1} +\eps_1 = \eps_1+\eps_2$ wouldn't be possible if all $\alpha_i$'s and $\alpha_{k+1}$ were long. Therefore one of them is short, which implies that this root plus $\gamma$ is a root, contradiction.

Suppose finally that the given root system is of type $F_4$. Pick a minimal relation (\ref{eqLeRootTricks}) that contradicts the statement of the lemma, i.e. one with minimal number of $\alpha_i$'s. This number must be at least 3, since otherwise this relation would generate a root subsystem of rank 3 or less and this is impossible by the preceding cases. We claim that for all $i,j$, $\alpha_{ij}:=\alpha_i+\alpha_j$ is not a root or zero. Indeed, assume the contrary. If $\alpha_{ij}+\gamma= \alpha_i+ \alpha_j +\gamma$ is not a root or zero we could replace $\alpha_i+\alpha_j$ by $\alpha_{ij}$ in contradiction with the minimality of the initial relation. Therefore $\alpha_{ij}+\gamma= \alpha_i+ \alpha_j +\gamma$ is a root or zero, and since the three roots $\alpha_i,\alpha_j,\gamma$ generate a root subsystem of rank at most 3, the preceding cases imply that at least one of $\alpha_i+\gamma$ and $\alpha_j+\gamma$ is a root, contradiction.

So far, for all $i,j$, we established that $\alpha_i+\alpha_j$ is not a root or zero; therefore $\langle\alpha_i,\alpha_j\rangle\geq 0$ for all $i,j$.
% and also, by the minimality of our relation, $\delta\neq \alpha_i$ for all $i$. 
Taking $\langle\alpha_1,\bullet \rangle$ on each side of  $\alpha_1+ \dots+ \alpha_k +\alpha_{k+1}+\gamma=\delta$ we see that $2\leq\langle\alpha_1,\alpha_1\rangle\leq \langle\alpha_1,\delta\rangle$. Therefore $\delta-\alpha_1$ is a root or zero, and transferring $\alpha_1$ to the right-hand side we get a shorter relation than the initial one. Contradiction.
\end{proof}
%\begin{corollary}\label{corstreakOfRoots}
%Let $\alpha_1,\dots,\alpha_k$ be roots of a root system such that $\alpha_1+\dots+\alpha_k$ is a root or is equal to zero. Then they can be rearranged to a sum $\alpha_{\sigma(1)}+\dots+\alpha_{\sigma(k)}$ (where $\sigma$ denotes a permutation of $k$ elements) so that for all $l\leq k$ the sum of the first $l$ terms is a root or is equal to zero.
%\end{corollary}

%\begin{lemma}
%Let $\Delta(\kk)$ be the root system of a root subalgebra of the semisimple Lie algebra $\gG$. Let $\pi$ denote the orthogonal projection in $\hh^*$ onto $\Delta(\kk)^\perp$, i.e. the projection on the second summand in the decomposition $\hh^*=\linspan(\Delta(\kk))\oplus \Delta(\kk)^\perp$. Denote the set of weights of the $\kk$-submodules of $\gG$ by $\bar\Delta$. Then $\pi$ gives a bijection between $\pi(\bar\Delta)\backslash \{0\}$ and the $\kk$-submodules of $\gG$ whose weights lie outside of $\linspan(\Delta(\kk))$. The image $\pi(\bar\Delta)$ is is either a root system or the union of a root system with the zero vector with respect to the bilinear form $\langle\alpha,\beta\rangle_\pi:= \langle \pi( \alpha), \pi( \beta) \rangle$.
%\end{lemma}
%\textbf{Remark.} In other words, the lemma states that $\kk$-submodules of the $\kk$-module $\gG$ whose weights lie outside of $\linspan(\Delta(\kk))$ can naturally be considered as a root subsystem of $\Delta(\gG)$.

%\begin{proof}
%Let $\alpha,\beta$ be roots. Suppose that there exist integers $a'$ and $b'$ such that $\pi (a'\alpha)=\pi(b'\beta)\neq 0$. We then claim that $\alpha$ and $\beta$ are weights of the same $\kk$-submodule of the $\kk$-module $\gG$. $a'\alpha-b'\beta\in\linspan\Delta(\kk)$ implies that there exist roots $\gamma_1,\dots,\gamma_k\in\Delta(\kk)$ and integers $a,b,c_1,\dots,c_k$ with $a\alpha=b\beta+c_1\gamma_1+\dots+c_k\gamma_k$. Among all such possible choices for $a,b,c_1,\dots,c_k,\gamma_1,\dots,\gamma_k$ for which $a\neq 0, b\neq 0$, pick one with minimal value for  $|a|+ |b|$, and, fixing those $a,b$, minimal sum $\sum |c_i|$.    
%\end{proof}

%\begin{lemma}\label{le2alphaplusbeta}
%Let $2\alpha+\beta=\gamma $, where $\alpha, \beta,\gamma$ are roots.  Then $\alpha+\beta$ is also a root. 
%\end{lemma}
%\begin{proof}
%Taking $\langle\gamma,\bullet\rangle$ on both sides of our equality we get $\langle\gamma,\gamma\rangle=\langle\beta,\beta\rangle + 4\langle\alpha,\alpha\rangle+4\langle\alpha,\beta\rangle$. For any three roots $\alpha,\beta,\gamma$, we have that $\langle\gamma, \gamma\rangle < \langle \beta , \beta \rangle + 4\langle\alpha,\alpha\rangle$, and therefore $\langle\alpha,\beta\rangle<0$. 

%\textbf{Alternative proof.} Assume the contrary and apply the previous lemma for $\alpha+\alpha+\beta=\gamma$ to get a contradiction.
%\end{proof}

\begin{defi}\label{defMinRelation}
For a relation \refeq{eqRelation} we define \emph{the length of the relation} to be $\sum_i a_i$. We define a relation \refeq{eqRelation} to be \emph{minimal} if its length is minimal, there are no repeating summands on either side, and no two $\beta_i$'s sum up to a root. 
\end{defi}

\textbf{Remark.} Any relation of minimal length can be transformed to a minimal relation by combining the repeating summands on both sides and by replacing the $\beta_i$'s in \refeq{eqRelation} that sum up to roots by their sums. If in addition the initial relation of minimal length corresponds to a two-sided weight, Lemma \ref{lestreakOfRoots} implies that the resulting minimal relation again corresponds to a two-sided weight.


\begin{prop}\label{propabundanceHolds}
Let the cone condition fail. Then there exists a minimal relation (\ref{eqRelation}) correspoding to a two-sided weight $\omega$.
\end{prop}
\begin{proof}
The failure of the cone condition is equivalent to the existence of a relation \refeq{eqRelation}. Pick a minimal such relation. Assume that the weight arising in this way is not two-sided. Together with the minimality of the relation this implies that for one of the $\alpha_i$'s, say $\alpha_1$, there exist roots $\beta'\in\Delta(\nn)$ and $\delta\in\Delta(\gG)\backslash\Delta(\nn)$ such that $\delta= \alpha_{1}+\beta'$.

We claim that $\delta\notin\Delta(\lL)$. Indeed, assume on the contrary that $\delta\in\Delta(\kk)$. Then $\beta'-\delta$ is a root, and therefore lies in $\Delta(\nn)$. We get the relation $(a_1-1)\alpha_1+a_2\alpha_2+ \dots+ a_l\alpha_l = b_1\beta_1 +\dots+b_k\beta_k+(\beta'-\delta)$ which is shorter than the initial relation, contradiction. 

Now suppose that $\delta$ is not $\bb\cap\kk$-singular. Therefore there exists $\gamma_1\in\Delta^+(\kk)$ such that $\delta_1 :=\delta + \gamma_1$ is a root. If $\delta_1$ is not singular, continue picking in a similar fashion roots $\gamma_2,\dots , \gamma_s \in \Delta^+(\kk)$, such that $\delta_t= \delta+ \gamma_1+ \dots +\gamma_t$ is a root for any $t\leq s$. Since this process must be finite, $\delta_s$ is $\bb\cap\kk$-singular for some $s$. As $\alpha_1$ is $\bb\cap\kk$-singular, $\alpha_1 +\gamma_1$ is not a root.  Apply now \refle{lestreakOfRoots} to $\delta_1=\alpha_1+\beta'+\gamma_1$ to get that $\beta'':=\beta'+\gamma_1$ is a root. Therefore $\beta''\in\Delta(\nn)$. Arguing in a similar fashion, we obtain that $\beta''':=\beta''+ \gamma_2$ is a root of $\nn$, and so on. Finally, we obtain $\beta^{(s+1)}:=\beta'+\gamma_1+\dots +\gamma_{s}\in\Delta(\nn)$ and so we get a new relation \refeq{eqRelation}:
\begin{equation}\label{eqminMaxRel}
(a_1-1)\alpha_1+ \delta_s+a_2\alpha_2+\dots+a_l\alpha_l=\beta_1+\dots +\beta_k+\beta^{(s+1)}.
\end{equation}
We can reduce \refeq{eqminMaxRel} so that no two $\beta$'s add to a root (replace any such pairs by their sum) and so that if $\delta_s=\alpha_i$ for some $i$ then $\delta_s+a_i\alpha_i$ is replaced by $(a_i+1)\alpha_i$. %The reduction of (\ref{eqminMaxRel}) is a minimal relation (\ref{eqRelation}).

This reduction of (\ref{eqminMaxRel}) is a minimal relation. If this relation does not yield a two-sided weight, one applies the procedure again and obtains a new minimal relation, and so on. As this process adds vectors from $\Delta(\nn)$ to the right-hand side of the relation, while the length of the left-hand side remains constant, the process must be finite (cf. Lemma \ref{leNilradicalIsSeparate}). Therefore there exists a minimal relation corresponding to a two-sided weight.
% with strictly $t$-larger right-hand side. Furthermore, applying this procedure does not change the value of $\sum a_i$. Therefore, there are only finitely many possibilities for the left-hand side and thus only finitely many possible values for its evaluation by $t$. On the other hand the right-hand side strictly increases with respect to $t$, and so this process cannot continue infinitely. Since applying the described process depended on the failure of our relation to be two-sided, we see that at least one minimal relation \refeq{eqRelation} must be two-sided. 
\end{proof}
\subsection{From two-sided to $\lL$-infinite weights}
In the remainder of this section we prove that the failure of the cone condition implies the existence of an $\lL$-infinite weight: our proof is mathematical for the classical Lie algebras and $G_2$ and uses a computer program for the exceptional Lie algebras $F_4$, $E_6$ and $E_7$.

For the classical Lie algebras, our scheme of proof can be summarized as follows. First, we classify all minimal relations (\ref{eqRelation}). It turns out by direct observation that whenever the cones intersect, the minimal relations (\ref{eqRelation}) are always of length 2, in particular this minimal length does not depend on the rank of the root system. In type $A$ this was discovered in \cite{PSZ}. In types $A$, $B$ and $D$, a direct inspection of all minimal relations shows that each of them possesses a strongly orthogonal decomposition with respect to $\Delta(\nn)$. Since at least one minimal relation must be two-sided by Proposition \ref{propabundanceHolds}, we obtain the existence of an $\lL$-strictly infinite weight. 

In type $C$ we do not have that all minimal relations possess a strongly orthogonal decomposition. However, the ``discrepancy'' is small - there is only one minimal relation (\ref{eqRelation}) without such a decomposition. In this particular case, we exhibit a root subalgebra $\tT$ containing $\kk$ such that $\tT$ has an $\lL\cap\tT$-strictly infinite weight, i.e. there is an $\lL$-infinite weight.

The proof for the exceptional Lie algebras uses a mixture of combinatorics and computer brute force. If $C(\kk_{ss})\cap\nn$ is not the nilradical of a parabolic subalgebra of $C(\kk_{ss})$, we prove in \refle{leCentralizerAndNilradical} that an $\lL$-infinite weight always exists, involving only roots of the root system of $C(\kk_{ss})$. Then, using our computer program, we enumerate up to $\gG$-automorphisms all remaining cases - i.e. the root subalgebras for which $C(\kk_{ss})\cap\nn$ is the nilradical of a parabolic subalgebra of $C(\kk_{ss})$ containing $C(\kk_{ss})\cap \hh$. This direct computation shows the existence of $\lL$-strictly infinite weights in types $E_6$ and $E_7$. In type $F_4$, our program fails to exhibit an $\lL$-strictly infinite weight for only one (unique up to $\gG$-automorphism) choice of $\lL$; in this case we give an argument similar to that in the special case in type $C$.

In order to enumerate all possible subalgebras $\kk$ we use the classification of reductive root subalgebras given in the fundamental paper \cite{Dynkin:Semisimple}. The list of possible subalgebras $\kk$ is very short (it contains respectively 19, 45, 75 and 22 entries for $F_4$, $E_6$, $E_7$ and $E_8$\arxivVersion{, see section \ref{secRootSATables} in the appendix}). For a fixed $\kk$, we use all automorphisms of $\Delta(\gG)$ which preserve $\Delta(\bb\cap\kk)$\arxivVersion{ (see section \ref{secCardinalitiesIsos} in the appendix)} in order to generate only pairwise non-conjugate subalgebras $\lL$ and thus further decrease the size of the computation.

We include tables of the cardinalities for each of those groups in the appendix. An interesting side question arises: what is the structure of those groups? Since they contain a subgroup isomorphic to the Weyl group of $C(\kk_{ss})$, this question can often be answered by mere inspection of the cardinalities of the groups listed in Table \ref{tableGroups}, but not in all cases - for example, not for the group of order 168 of root system automorphisms of $E_7$ that permutes the positive roots of $7A_1\subset E_7$.


The following is an observation that is helpful in the proof of Lemma \ref{leBigTedious} (\textit{cf.} \cite[Lemma 5.4]{PSZ}).
\begin{lemma}\label{leminrelationgeneral}
Let the cone condition fail and let us have a minimal relation \refeq{eqRelation}. Then
\begin{itemize}
\item[(a)] The relation has the form 
\[
\alpha_1+\alpha_2=\beta_1,
\] 
or 
\item[(b)] $\alpha_i+\alpha_j$ is not a root for all $i,j$.
\end{itemize}
\end{lemma}
\begin{proof}
Pick a minimal relation \refeq{eqRelation}. Suppose there exist indices $i,j$ such that $\gamma:=\alpha_i+\alpha_j$ is a root. We claim that $\gamma\in \Delta(\nn)$. Indeed, assume the contrary. First, suppose $\gamma\in\Delta^-(\kk)$. Then $\alpha_i$ and $\alpha_j$ would both fail to be $\bb\cap\kk$-singular. 

Second, suppose $\gamma\in\Delta^+(\kk)$. We prove that $\alpha_3+\dots+\alpha_l=\beta_1+\dots +\beta_k-\gamma$ is a shorter relation than \refeq{eqRelation}. Indeed, $\beta_1+\dots +\beta_k-\gamma$ is clearly non-zero (positive linear combination of elements of $\Delta(\nn)$ cannot be in the span of the roots of the semisimple part). By the $\bb\cap\kk$-singularity of the $\alpha_i$'s,   $\langle\gamma, \beta_1+\dots+ \beta_k\rangle =\langle\gamma,\alpha_1 +\dots +\alpha_k\rangle =\langle \gamma, \gamma \rangle +\langle \gamma, \alpha_3 +\dots + \alpha_k \rangle>0$ and therefore, for some index $i$,  $\langle\gamma,\beta_i\rangle>0$. This shows that $\beta_i-\gamma$ is a root, which therefore belongs to $\Delta(\nn)$. Contradiction.

Third, suppose $\gamma\notin\Delta(\gG)\backslash\Delta(\lL)$. Then $\gamma$ is $\bb\cap\kk$-singular - if $\gamma+ \delta = \alpha_1 + \alpha_2+\delta$ were a root for some $\delta\in\Delta^+(\kk)$, then \refle{lestreakOfRoots} would imply that $\alpha_1+\delta$ is also a root. Therefore we can shorten the relation \refeq{eqRelation} by replacing $\alpha_1+\alpha_2$ by $\gamma$, and the obtained relation is non-trivial since the right-hand side is not zero. Contradiction.

Therefore $\gamma\in \Delta(\nn)$, and our lemma is proved.
\end{proof}


\subsection{Minimal relations \refeq{eqRelation} in the classical Lie algebras}
The following lemma describes all minimal relations (\ref{eqRelation}) up to automorphisms of $\Delta(\gG)$.
\begin{lemma}\label{leBigTedious}
Let $\gG\simeq\so(2n),\so(2n+1),\sP(2n)$. Suppose $\lL=\kk\crplus\nn$ does not satisfy the cone condition. 
\begin{itemize}
\item A minimal relation \refeq{eqRelation} has length 2 (Definition \ref{defMinRelation}).
\item All possibilities for minimal relations \refeq{eqRelation}, up to an automorphism of $\Delta(\gG)$, are given in the following table.
\end{itemize}
\end{lemma}
\noindent \begin{longtable}{rcl | c |c c}
\multicolumn{3}{c|}{$\omega$} & \begin{tabular}{p{4.5cm}}Scalar products. All non-listed scalar products are zero. All roots, unless stated otherwise, are assumed long in types $B,D$ and short in type $C$.\end{tabular}&  \multicolumn{2}{p{2cm}}{The roots from the relation generate }
\endhead
\hline
\multicolumn{5}{c}{$\gG\simeq\so(2n)$}
\\\hline
$\alpha_1+\alpha_2$ &$=$&$\beta_1$&
\begin{tabular}{l}
$\langle\alpha_1, \beta_1 \rangle$ $=$ $\langle\alpha_2,\beta_1\rangle$ $=1$\\
$\langle\alpha_1,\alpha_2\rangle$ $=-1$
\end{tabular}&
 $A_2$&
 \refstepcounter{equation} \label{eqDcontainsA2'} (\ref{eqDcontainsA2'})\\\hline
$\alpha_1+\alpha_2$ &$=$&$\beta_1+\beta_2$&
\begin{tabular}{l}
$\langle\alpha_1,\beta_1\rangle$ $=$ $\langle\alpha_1,\beta_2\rangle$ $=$\\ 
$\langle\alpha_2,\beta_1\rangle$ $=$ $\langle\alpha_2,\beta_2\rangle$ $=1$ 
\end{tabular}&
\begin{tabular}{l}
 $A_3\subset A_4$,\\ $n\geq 5$ 
\end{tabular}& \refstepcounter{equation} \label{eqDcontainsA3'} (\ref{eqDcontainsA3'})
\\\hline
$\alpha_1+\alpha_2$ &$=$&$\beta_1+\beta_2+\beta_3$&
\begin{tabular}{l}
$\langle\alpha_1,\alpha_2\rangle$ $=$  $\langle\alpha_1,\beta_1\rangle$ $=$\\
$\langle\alpha_1,\beta_2\rangle$ $=$ $\langle\alpha_1,\beta_3\rangle$ $=$ \\
$\langle\alpha_2,\beta_1\rangle$ $=$ $\langle\alpha_2,\beta_2\rangle$ $=$ \\
$\langle\alpha_2,\beta_3\rangle$ $=1$
\end{tabular}& 
$D_4$&
\refstepcounter{equation} \label{eqDcontainsD4'} (\ref{eqDcontainsD4'})
\\\hline
$2\alpha_1$ &$=$&$\beta_1+\beta_2+\beta_3+\beta_4$&
\begin{tabular}{l}
$\langle\alpha_1,\beta_1\rangle$ $=$ $\langle\alpha_1,\beta_2\rangle$ $=$ \\
$\langle\alpha_1,\beta_3\rangle$ $=$ $\langle\alpha_1,\beta_4\rangle$ $=1$ 
\end{tabular}&
$D_4$&
\refstepcounter{equation} \label{eqDcontainsD4''} (\ref{eqDcontainsD4''})\\\hline
$\alpha_1+\alpha_2$ &$=$&$\beta_1+\beta_2$&
\begin{tabular}{l}
$\langle\alpha_1,\beta_1\rangle$ $=$ $\langle\alpha_1,\beta_2\rangle$ $=$ \\
$\langle\alpha_2,\beta_1\rangle$ $=$ $\langle\alpha_2,\beta_2\rangle$ $=1$ 
\end{tabular}&
 $A_3$\footnote{\cite[Table 9]{Dynkin:Semisimple} uses the notation ``$D_3$'' for such subalgebras. $D_3$ is defined as a root subsystem of type $A_3$ of root system of type $B$ or $D$, which cannot be extended to a root subsystem of type $A_4$.}
&
\refstepcounter{equation} \label{eqDcontainsD3'} (\ref{eqDcontainsD3'}) \\\hline
\multicolumn{5}{c}{$\gG\simeq\so(2n+1)$}\\\hline
\multicolumn{3}{c|}{\begin{tabular}{c} all relations \\ listed for $\so(2n)$\end{tabular}}& - & - \\\hline
$\alpha_1+\alpha_2$ &$=$&$\beta_1$&
\begin{tabular}{l}
$\langle\alpha_1, \beta_1 \rangle$ $=$ $\langle\alpha_2,\beta_1\rangle$ $=1$,\\
$\|\alpha_1\|=\|\alpha_2\|=1 $ 
\end{tabular}&
 $B_2$ &
 \refstepcounter{equation} \label{eqBcontainsB2'} (\ref{eqBcontainsB2'})\\\hline
$\alpha_1+\alpha_2$ &$=$&$\beta_1$&  
\begin{tabular}{l}
$\langle\alpha_1, \alpha_2 \rangle$ $=-1$, $\langle\alpha_2,\beta_1\rangle$ $=1$, \\
$\|\alpha_1\|=\|\beta_1\|=1 $ 
\end{tabular}&
 $B_2$&
\refstepcounter{equation} \label{eqBcontainsB2''} (\ref{eqBcontainsB2''}) \\\hline
$2\alpha_1$ &$=$&$\beta_1+\beta_2$&  
\begin{tabular}{l}
$\|\alpha_1\|=1, \langle\alpha_1,\beta_1\rangle=\langle\alpha_1,\beta_2\rangle=1$ 
\end{tabular}&
 $B_2$&
\refstepcounter{equation} \label{eqBcontainsB2'''} (\ref{eqBcontainsB2'''}) \\\hline
$\alpha_1+\alpha_2$ &$=$&$2\beta_1$&  
\begin{tabular}{l}
$\|\beta_1\|=1, \langle\alpha_1,\beta_1\rangle=\langle\alpha_2,\beta_1\rangle=1$ 
\end{tabular}&
 $B_2$&
 \refstepcounter{equation} \label{eqBcontainsB2''''} (\ref{eqBcontainsB2''''}) \\\hline
$\alpha_1+\alpha_2$ &$=$&$\beta_1+\beta_2$&  
\begin{tabular}{l}
$\langle\alpha_1, \beta_1 \rangle=$  $\langle\alpha_2,\beta_1\rangle=$ $\langle\alpha_2,\beta_2\rangle$ $=1$, \\
$\|\alpha_1\|=\|\beta_2\|=1 $ 
\end{tabular}&
 $B_3$&
\refstepcounter{equation} \label{eqBcontainsB3'} (\ref{eqBcontainsB3'}) \\\hline
$\alpha_1+\alpha_2$ &$=$&$2\beta_1+\beta_2$&  
\begin{tabular}{l}
$\langle\alpha_1, \alpha_2 \rangle=$ $\langle\alpha_1,\beta_1\rangle=$ $\langle\alpha_2,\beta_1\rangle$ $=1$, \\
$\langle\alpha_1,\beta_2\rangle=$ $\langle\alpha_2,\beta_2\rangle=$ $\|\beta_1\|$ $=1$, 
\end{tabular}&
 $B_3$&
 \refstepcounter{equation} \label{eqBcontainsB3''} (\ref{eqBcontainsB3''}) \\\hline
$2\alpha_1$ &$=$&$2\beta_1+\beta_2+\beta_3$&  
\begin{tabular}{l}
$\langle\alpha_1, \beta_1 \rangle=$ $\langle\alpha_1,\beta_2\rangle=$ $\langle\alpha_1,\beta_3\rangle$ $=1$, \\
$\|\beta_1\|=1 $ 
\end{tabular}&
$B_3$&
\refstepcounter{equation} \label{eqBcontainsB3'''} (\ref{eqBcontainsB3'''}) \\\hline
\multicolumn{5}{c}{$\gG\simeq\sP(2n)$}\\\hline
$\alpha_1+\alpha_2$ &$=$&$\beta_1$&
\begin{tabular}{l}
$\langle\alpha_1, \beta_1 \rangle$ $=$ $\langle\alpha_2,\beta_1\rangle$ $=1$\\
$\langle\alpha_1,\alpha_2\rangle$ $=-1$
\end{tabular}&
$A_2$&
\refstepcounter{equation} \label{eqCcontainsA2'} (\ref{eqCcontainsA2'})\\\hline
$\alpha_1+\alpha_2$ &$=$&$\beta_1$&
\begin{tabular}{l}
$\langle\alpha_1, \beta_1 \rangle$ $=$ $\langle\alpha_2,\beta_1\rangle$ $=2$,\\
$\|\beta_1\|=2$
\end{tabular}&
$C_2$&
\refstepcounter{equation} \label{eqCcontainsC2'} (\ref{eqCcontainsC2'}) \\\hline
$\alpha_1+\alpha_2$ &$=$&$2\beta_1$&
\begin{tabular}{l}
$\langle\alpha_1,\beta_1\rangle=\langle\alpha_2,\beta_1\rangle$ $=1$ \\
$\|\alpha_1\|=\|\alpha_2\|$ $=2$
\end{tabular}& 
$C_2$&
\refstepcounter{equation} \label{eqCcontainsC2''} (\ref{eqCcontainsC2''}) \\\hline
$\alpha_1+\alpha_2$ &$=$&$\beta_1$&
\begin{tabular}{l}
$\langle\alpha_1,\beta_1\rangle=2$, $\langle\alpha_2,\alpha_1\rangle=-1$ \\
$\|\alpha_1\|$ $=2$
\end{tabular}& 
$C_2$&
\refstepcounter{equation} \label{eqCcontainsC2'''} (\ref{eqCcontainsC2'''}) \\\hline
$\alpha_1+\alpha_2$ &$=$&$\beta_1+\beta_2$&
\begin{tabular}{l}
$\langle\alpha_1,\beta_1\rangle$ $=$ $\langle\alpha_1,\beta_2\rangle$ $=$ \\
$\langle\alpha_2,\beta_1\rangle$ $=$ $\langle\alpha_2,\beta_2\rangle$ $=1$ 
\end{tabular}& 
 $A_3$&
 \refstepcounter{equation} \label{eqCcontainsA3'} (\ref{eqCcontainsA3'})\\\hline
$\alpha_1+\alpha_2$ &$=$&$\beta_1+\beta_2$&
\begin{tabular}{l}
$\langle\alpha_2,\beta_1\rangle$ $=$ $\langle\alpha_2,\beta_2\rangle$ $=2$,\\
$\langle\alpha_1,\beta_1\rangle=\langle\alpha_1,\beta_1\rangle$ $=$ \\
$\langle\alpha_1,\beta_2\rangle=1$,  $\langle\beta_1,\beta_2\rangle=1$, \\
$\|\alpha_2\|=2$ 
\end{tabular}& 
 $C_3$&
 \refstepcounter{equation} \label{eqCcontainsC3'} (\ref{eqCcontainsC3'})\\\hline
$\alpha_1+\alpha_2$ &$=$&$\beta_1+\beta_2$&
\begin{tabular}{l}
$\langle\alpha_1,\beta_2\rangle$ $=$ $\langle\alpha_1,\beta_2\rangle$ $=2$, \\
$\langle\alpha_1,\alpha_2\rangle=\langle\alpha_1,\beta_1\rangle$ $=$ \\
$\langle\alpha_2,\beta_1\rangle=1$, $\|\beta_2\|=2$ 
\end{tabular}& 
 $C_3$&
 \refstepcounter{equation} \label{eqCcontainsC3''} (\ref{eqCcontainsC3''})\\\hline
\end{longtable}

One possible way to realize the above data in $\eps$-notation follows. Equalities (\ref{eqDcontainsA2'})-(\ref{eqCcontainsC3''}) correspond to equalities (\ref{eqRelso(2n)Type2})-(\ref{eqRelsp(2n)Type3'}) in the same order.

\smallskip \noindent (a) $\gG\simeq\so(2n)$
\begin{eqnarray}
\underbrace{\eps_{1}+\eps_{2}}_{\alpha_1}+\underbrace{(- \eps_{2}+\eps_{3})}_{\alpha_2} &=& \underbrace{\eps_{1}+\eps_{3}}_{\beta_1},\label{eqRelso(2n)Type2}\\
\label{eqRelso(2n)Type1}
\underbrace{\eps_{1}+\eps_{2}}_{\alpha_1}+\underbrace{\eps_{3}+\eps_{4}}_{\alpha_2}  &=& \underbrace{ \eps_{1}+ \eps_{3}}_{\beta_1}+ \underbrace{\eps_{2}+\eps_{4}}_{\beta_2},\\
\label{eqRelso(2n)Type1'}
\underbrace{\eps_{1}+\eps_{2}}_{\alpha_1}+\underbrace{ \eps_{1}+\eps_{4}}_{\alpha_2}&=& \underbrace{  \eps_{1}+ \eps_{3}}_{\beta_1} + \underbrace{\eps_{1}- \eps_{3}  }_{\beta_2} +  \underbrace{\eps_{2 }+ \eps_{4}  }_{\beta_3},\\
\label{eqRelso(2n)Type1''}
2\underbrace{(\eps_{1}+\eps_{2})}_{\alpha_1}&=&
\underbrace{\eps_{1}+\eps_{3}}_{\beta_1}+\underbrace{\eps_{1}  -\eps_{3}}_{\beta_2} + \underbrace{\eps_{2} + \eps_{4} }_{\beta_3} +\underbrace{\eps_{2}- \eps_{4} }_{\beta_4},\\
\underbrace{\eps_{1}+\eps_{2}}_{\alpha_1}+ \underbrace{(-\eps_{2}+\eps_{1})}_{\alpha_2}&=&\label{eqRelso(2n)Type3} \underbrace{\eps_{1} +\eps_{3}}_{\beta_1} + \underbrace{\eps_{1}-\eps_{3}}_{\beta_2},
\end{eqnarray}

\smallskip \noindent (b) $\gG\simeq\so(2n+1)$
\begin{eqnarray}
\label{eqRelso(2n+1)Type1}
\underbrace{\eps_{1}}_{\alpha_1}+\underbrace{\eps_{2}}_{\alpha_2}&=& \underbrace{ \eps_{1} +\eps_{2} }_{\beta_1},\\
\label{eqRelso(2n+1)Type1'}
\underbrace{\eps_{1}}_{\alpha_1}+\underbrace{(-\eps_{1}+ \eps_{2})}_{\alpha_2}&=& \underbrace{ \eps_{2} }_{\beta_1},\\
\label{eqRelso(2n+1)Type3}
2\underbrace{\eps_{1}}_{\alpha_1}&=& \underbrace{ \eps_{ 1}+\eps_{2} }_{\beta_1}+  \underbrace{ \eps_{ 1}-\eps_{2} }_{\beta_2},\\
\label{eqRelso(2n+1)Type3.5}
\underbrace{\eps_{ 1}+\eps_{2} }_{\alpha_1}+  \underbrace{ \eps_{ 1}-\eps_{2} }_{\alpha_2}&=&2\underbrace{\eps_{1}}_{\beta_1},\\
\label{eqRelso(2n+1)Type2}
\underbrace{\eps_{1}}_{\alpha_1}+\underbrace{\eps_{2}+ \eps_{3}}_{\alpha_2}&=& \underbrace{\eps_{ 1}+\eps_{2} }_{\beta_1}+ \underbrace{ \eps_{3 }}_{\beta_2},\\
\label{eqRelso(2n+1)Type4}
\underbrace{\eps_{1}+\eps_{j_2}}_{\alpha_1}+\underbrace{\eps_{1}+ \eps_{3}}_{\alpha_1}&=& 2\underbrace{ \eps_{ 1} }_{\beta_1}+ \underbrace{ \eps_{ 2}+\eps_{3} }_{\beta_2},\\
\label{eqRelso(2n+1)Type5}
2\underbrace{(\eps_{1}+ \eps_{2})}_{\alpha_1}&=& 2\underbrace{ \eps_{ 1} }_{\beta_1}+\underbrace{\eps_{ 2}+\eps_{3} }_{\beta_2}+ \underbrace{ \eps_{ 2}-\eps_{3} }_{\beta_3} ,
\end{eqnarray}

\smallskip \noindent (c) $\gG\simeq\sP(2n)$
\begin{eqnarray}
\underbrace{\eps_{1}+\eps_{2}}_{\alpha_1}+\underbrace{(- \eps_{2}+\eps_{3})}_{\alpha_2} &=& \underbrace{ \eps_{1} + \eps_{3}}_{\beta_1}, \label{eqRelsp(2n)TypeA2}\\
\label{eqRelsp(2n)Type1}
\underbrace{\eps_{1}-\eps_{3}}_{\alpha_1} +\underbrace{\eps_{3}+\eps_{2}}_{\alpha_2} &=&\underbrace{\eps_{1}+\eps_{2}}_{\beta_1},\\
\label{eqRelsp(2n)Type4}
\underbrace{2\eps_{1}}_{\alpha_1}+\underbrace{ 2\eps_{2}}_{\alpha_2} &=& 2\underbrace{(\eps_{1}+\eps_{2})}_{\beta_{1}},\\
\label{eqRelsp(2n)Type1'}
\underbrace{\eps_{1}-\eps_{2}}_{\alpha_1} +\underbrace{\eps_{2}+\eps_{1}}_{\alpha_2} &=&\underbrace{2\eps_{1}}_{\beta_1},\\
\label{eqRelsp(2n)Type2}
\underbrace{\eps_{1}+\eps_{2}}_{\alpha_1}+\underbrace{\eps_{3}+\eps_{4}}_{\alpha_2}&=& \underbrace{\eps_{1}+ \eps_{3}}_{\beta_{1}}+\underbrace{ \eps_{2}+ \eps_{4}}_{\beta_{2}},\\
\label{eqRelsp(2n)Type3} 
\underbrace{\eps_{1}+\eps_{2}}_{\alpha_1}+\underbrace{ 2\eps_{3}}_{\alpha_2}&=& \underbrace{\eps_{1}+ \eps_{3}}_{\beta_{1}}+\underbrace{\eps_{2}+\eps_{3}}_{\beta_{2}},\\
\label{eqRelsp(2n)Type3'} 
\underbrace{\eps_{1}+\eps_{2}}_{\alpha_1}+\underbrace{ \eps_{2}+ \eps_{3}}_{\alpha_2}&=& \underbrace{\eps_{1}+  \eps_{3}}_{\beta_{1}}+ \underbrace{ 2\eps_{2}}_{\beta_{2}},
\end{eqnarray}

\noindent where $\alpha_{i}\in \hw_{\kk\cap\bb}(\gG/\lL)$, $\beta_i\in \Delta(\nn)$.

\subsection{Relations \refeq{eqRelation} with minimal support}\label{secAlgoMinRel}
Before we proceed with the proof of Lemma \ref{leBigTedious}, we present an algorithm for determining relations (\ref{eqRelation}) with minimal support. The algorithm is based on the main idea used in the proof of Lemma \ref{leBigTedious}. Nethertheless, the formal exposition of the proof of Lemma \ref{leBigTedious} is independent of the current section.

We define the \emph{$\alpha$-support} (respectively, the \emph{$\beta$-support}) of a relation \refeq{eqRelation} to be the set of roots $\{\alpha_i\}$ (respectively, $\{\beta_j\}$) participating in the relation. An $\alpha$-support (respectively $\beta$-support) is defined to be \emph{minimal} if it is the support of a minimal relation (\ref{eqRelation}). Denote $S:=\hw_{\bb\cap\kk}(\gG/\lL)$. Let us have a not necessarily minimal relation \refeq{eqRelation}
\begin{equation}\label{eqMinRelTemp}
a_1\alpha_{1}+\dots+a_L\alpha_L+\dots = b_1\beta_1+\dots+b_K\beta_K+\dots,
\end{equation}
where $\alpha_i\in S$, $\beta_j\in\Delta(\nn)$, $a_i,b_j>0$ and such that $a_1\alpha_{1}+\dots+a_L\alpha_L\neq b_1\beta_1+\dots+b_K\beta_K$. Suppose $\gamma:=c_1\alpha_{i_1}+\dots+ c_l\alpha_{i_l} -d_1\beta_{j_1} -\dots - d_k\beta_{j_k}$, $c_i,d_i>0$ is a root.
Then the following observations hold.
\begin{itemize}
\item If $l\geq 2$, $c_r\leq a_{i_r}$, $d_r\leq b_{j_r}$ and $\gamma\in S$ we can shorten \refeq{eqMinRelTemp} by transferring the $\beta_{j_i}$'s to the left hand side; $k=0$ is allowed.
\item If $-\gamma\in S$, $c_r\leq a_{i_r}$, $d_r\leq b_{j_r}$ and $k\geq 1$ then $\gamma+c_1\alpha_{i_1}+ \dots+c_l\alpha_{i_l}= d_1\beta_{j_1}+ \dots +d_k \beta_{j_k}$ is shorter than \refeq{eqMinRelTemp} unless $l=L$ and $k=K$ and all coefficient inequalities are equalities. The preceding equality has smaller support unless $l=L$ and $k=K$ independent of the inequalities $c_r\leq a_{i_r}$, $d_r\leq b_{j_r}$.
\item If $\gamma\in \Delta(\nn)$, $c_r\leq a_{i_r}$, $d_r\leq b_{j_r}$ then $\alpha_{i_1}+\dots+ \alpha_{i_l}= \beta_{j_1} +\dots + \beta_{j_k}+\gamma$ is shorter than \refeq{eqMinRelTemp} unless $l=L$ and $k=K$; note $l=0$ and $k=0$ are both allowed. The preceding equality has smaller support unless $l=L$ and $k=K$ independent of the inequalities $c_r\leq a_{i_r}$, $d_r\leq b_{j_r}$.
\item If $l\geq 1$, $c_r\leq a_{i_r}$, $d_r\leq b_{j_r}$ and $-\gamma\in \Delta(\nn)$ then we can shorten \refeq{eqMinRelTemp} by transferring the $\alpha_{i_j}$'s to the right hand side; note $k=0$ is allowed.
\end{itemize}
%\begin{itemize}
%\item If $l\geq 2$, $c_r\leq a_{i_r}$, $d_r\leq b_{j_r}$ and $\gamma\in S$ we can shorten \refeq{eqMinRelTemp} by transferring the $\beta_{j_i}$'s to the left hand side; $k=0$ is allowed.
%\item If $-\gamma\in S$, $c_r\leq a_{i_r}$, $d_r\leq b_{j_r}$ and $k\geq 1$ then $\gamma+c_1\alpha_{i_1}+ \dots+c_l\alpha_{i_l}= d_1\beta_{j_1}+ \dots +d_k \beta_{j_k}$ is shorter than \refeq{eqMinRelTemp} unless $l=L$ and $k=K$ and all coefficient inequalities are equalities. The preceding equality has smaller support unless $l=L$ and $k=K$ independent of the inequalities $c_r\leq a_{i_r}$, $d_r\leq b_{j_r}$.
%\item If $\gamma\in \Delta(\nn)$, $c_r\leq a_{i_r}$, $d_r\leq b_{j_r}$ then $\alpha_{i_1}+\dots+ \alpha_{i_l}= \beta_{j_1} +\dots + \beta_{j_k}+\gamma$ is shorter than \refeq{eqMinRelTemp} unless $l=L$ and $k=K$; note $l=0$ and $k=0$ are both allowed. The preceding equality has smaller support unless $l=L$ and $k=K$ independent of the inequalities $c_r\leq a_{i_r}$, $d_r\leq b_{j_r}$.
%\item If $l\geq 1$, $c_r\leq a_{i_r}$, $d_r\leq b_{j_r}$ and $-\gamma\in \Delta(\nn)$ then we can shorten \refeq{eqMinRelTemp} by transferring the $\alpha_{i_j}$'s to the right hand side; note $k=0$ is allowed.
%\end{itemize}
Using the preceding observation we will now formulate an algorithm for generating a list of possible minimal $\alpha$- and $\beta$-supports of the relations \refeq{eqRelation}. We keep note of four subsets $A,B,\bar A, \bar B \subset \Delta (\gG)$. The sets $A$ and $B$ parametrize respectively the $\alpha$-support and the $\beta$-support of a relation (\ref{eqRelation}). The sets $\bar A$ and $\bar B$ parametrize respectively the set of roots that do not belong to $S$ and the set of roots that do not belong to $\Delta(\nn)$. At the start of our algorithm, we set $A:=\{\alpha_1\}$ and $B:=\{\beta_1\}$ for two different roots with $\langle\alpha_1, \beta_1 \rangle >0$; we set $\bar A$ and $\bar B$ to be the empty set.
\begin{itemize}
\item[Step 1.] If $\Cone_{\QQ}(A)$ and $\Cone_{\QQ}(B)$ have a non-zero intersection, record the sets $A$ and $B$ respectively as possible $\alpha$- and $\beta$-supports. Terminate the current branch of the computation. If $\Cone_{\QQ}(A)$ and $\Cone_{\QQ}(B)$ do not intersect, go to Step 2.
\item[Step 2.] Since the cones do not intersect, there exists a weight $\mu$ such that $\langle\mu,\alpha\rangle\geq 0$ for all $\alpha\in A$, $\langle\mu,\alpha\rangle\leq 0$ for all $\beta\in B$ and such that $\mu$ has non-zero scalar product with at least one of the $\alpha$. Pick one such weight $\mu$.
\item[Step 3.] For each root $\alpha'$ with $\langle\mu,\alpha'\rangle<0$, add $\alpha'$ to $A$, and go to Step 4. For each root $\beta'$ with $\langle\mu,\alpha'\rangle>0$ add $\beta$ to $B$ and go to Step 4. Here the computation branches: in a computer program, Step 3 would correspond to two loops - one for each possibility for $\alpha'$ and one for each possibility for $\beta'$.
\item[Step 4.] For all possible subsets $\{\alpha_{i_1},\dots,\alpha_{i_l}\}$ of $A$ and $\{\beta_{j_1},\dots,\beta_{j_k}\}$ of $B$ check whether $\gamma:=\alpha_{i_1}+\dots +\alpha_{i_l}-\beta_{j_1}-\dots-\beta_{j_k}$ is a root. If that is the case, do the following four operations, corresponding to the four cases of the preceding observation. If $l\geq 2$ add $\gamma$ to $\bar A $; if $(l,k)\neq(\# A,\# B)$ add $-\gamma$ to  $\bar A$ and $\gamma$ to $\bar B$; finally, if $l\geq 1$ add $-\gamma$ to  $\bar B $.
\item[Step 5.] If the sets $A$ and $\bar A $ intersect, terminate the current branch of the computation. If the set of roots additively generated by $B$ intersects $\bar B $, terminate the current branch of the computation. In any other case go to Step 6.
\item[Step 6.] Check whether the current sets $A$ and $B$ have been considered, up to an automorphism of $\Delta(\gG)$, in any preceding branch of the computation. If that is the case terminate the current branch of the computation. %\textbf{Remark.} This step, as careful analysis of the proof of \refle{leBigTedious} below shows, is the main reason why the number of cases we get will not depend on $n$ for $\Delta(\gG)=A_n, B_n, C_n$ or $D_n$ for all large enough $n$.
\item[Step 7.] Go to Step 1.
\end{itemize}

\subsection{Proof of \refle{leBigTedious}} \label{secBigTediousProof}

\begin{proof}Pick a minimal relation \refeq{eqRelation} of the form $\omega:=a_1\alpha_1+\dots+a_l\alpha_l=b_1\beta_1+\dots+b_k\beta_k$ (see Definition \ref{defMinRelation}).

Throughout this proof, we will use the informal expression ``$\pm\eps_{i}$ appears with a positive (resp. non-positive) coefficient in the weight $\omega$'' to describe the $\pm\eps_i$-coordinate of $\omega$ in the basis $\{\eps_1,\dots, \eps_{i-1},\pmf{i}\eps_i, \eps_{i+1}, \dots, \eps_n\}$. 

\subsubsection{$\gG\simeq\so(2n)$}
\smallskip\noindent Case 1. There exists an index $i$, such that $\alpha_i =\pmf{j_1} \eps_{j_1} +(\pmf{j_2} \eps_{j_2})$, $j_1\neq j_2$ and both $\pmf{j_1} \eps_{j_1}$ and $\pmf{j_2}\eps_{j_2}$ appear with a positive coefficient in $\omega$. Without loss of generality we may assume $i=1$ and $\alpha_1= \eps_{1}+\eps_{2}$. Therefore there exist $\beta_{1}$ and $\beta_{2}$ on the right-hand side of the relation with $\beta_1= \eps_{1} +(\pmf{j_3}\eps_{j_3})$ and $\beta_2=\eps_{2}+(\pmf{j_4} \eps_{j_4})$. The minimality of the relation implies $\{1, 2\}\cap\{j_3,j_4\}=\emptyset $. The latter allows us to assume without loss of generality that $\beta_1= \eps_{1} +\eps_{3}$.
\[
\underbrace{\eps_{1}+\eps_{2}}_{\alpha_1}+\dots= \underbrace{\eps_{1}+\eps_{3}}_{\beta_1}+\underbrace{ \eps_{2} +\eps_{j_3}}_{\beta_2}+\dots .
\]
We will now prove $j_4\neq 3$. 

Assume on the contrary that $3=j_4$. As the relation is minimal, the choice of $\pmf{j_4}$ sign must be such that $\eps_{3}= \pmf{j_4} \eps_{j_4}$. The minimality of the relation implies that there can be no cancellation of the weight $\eps_{3}$ on the right-hand side. Therefore on the left-hand side there exists a root, say $\alpha_2$, such that $\alpha_2=\eps_{3}+(\pmf{j_5}\eps_{j_5})$, $j_5\neq 3$. The minimality of the relation implies that in addition $j_5\neq 1,2$. Thus we can assume without loss of generality that $j_5=5$ and $\alpha_2=\eps_3+\eps_5$. So far, the assumption that $j_4=3$ implies that the relation has the form
\begin{equation}\label{eqConesso(2n)}
\underbrace{\eps_{1}+\eps_{2}}_{\alpha_1}+\underbrace{ \eps_{3}+\eps_{5}}_{\alpha_{2}}+\underbrace{\dots}_{\gamma}= \underbrace{\eps_{1}+\eps_{3}}_{\beta_1}+\underbrace{ \eps_{2} +\eps_{3}}_{\beta_2}+\underbrace{\dots}_{\delta},
\end{equation}
where $\gamma$ and $\delta$ denote the omitted summands. Suppose at least one of the roots $\eps_{1}+\eps_{5}$ and $\eps_{2} +\eps_{5}$ belongs to $\Delta(\nn)$. Without loss of generality we may assume $\eps_{1} +\eps_{5} \in\Delta(\nn)$. Then the relation $\alpha_1+\alpha_2= \beta_{2}+ \eps_{ 1} +\eps_{5}$ is shorter than \refeq{eqConesso(2n)}. Contradiction. Suppose at least one of the roots $\eps_{1}+\eps_{5}$ and $\eps_{2}+\eps_{5}$ belongs to $\Delta(\kk)$. Without loss of generality we may assume $\eps_{1}+\eps_{5}\in\Delta(\kk)$. Then $\eps_{3} - \eps_{5}= \beta_1 -( \eps_{1} + \eps_{5}) \in \Delta(\nn)$ and the relation $\gamma= \eps_{3} - \eps_{5}+\delta$ is shorter than \refeq{eqConesso(2n)}. The latter relation is non-trivial since the right-hand side is a positive linear combination of roots of $\Delta(\nn)$. Contradiction.

So far we proved that $\eps_{1}+\eps_{5},\eps_{2} +\eps_{5}$ do not belong to $\Delta(\lL)$. If $\eps_{1} + \eps_{5}$ were a $\bb\cap\kk$-singular weight, we could replace $\alpha_1+\alpha_2$ by $\eps_{1} + \eps_{5}$ and remove $ \eps_{ 2} + \eps_{3}$ on the right-hand side of \refeq{eqConesso(2n)}, shortening the initial relation. Similarly, we reason that $\eps_{2}+\eps_{5}$ is not a $\bb\cap\kk$-singular weight. In order for $\eps_{ 1}+\eps_{5}$ not to be $\bb\cap\kk$-singular, there must exist an index $k$ and a choice of sign for which one of $\pmf{k}\eps_{k} -\eps_{ 1}$ and $\pmf{k}\eps_{k}-\eps_{5}$ is a positive root of $\kk$. Similarly, there exists an index $l$ and a choice of sign for which one of $\pmf{l}\eps_{l} -\eps_{ 2}$ and $\pmf{l}\eps_{l}-\eps_{5}$ is a positive root of $\kk$. As $\alpha_1$ and $ \alpha_2 $ are $\bb\cap\kk$-singular, a short consideration shows that the only possibility is $\pmf{k}\eps_k=-\eps_2$ and $\pmf{l}\eps_l=\eps_{3}$. Therefore $\beta_3 := \beta_2 -(\eps_{3}-\eps_{5})\in \Delta(\nn)$. Finally, we obtain the relation $\alpha_1+\alpha_2=\beta_1+\beta_3$ which is shorter than \refeq{eqConesso(2n)}. Contradiction.

So far, we have proved that $3\neq j_4$. Therefore we can assume without loss of generality that $\beta_2=\eps_2+\eps_4$. We have now established that the relation has the form
\[
\underbrace{\eps_{1}+\eps_{2}}_{\alpha_1}+\dots= \underbrace{\eps_{1}+\eps_{3}}_{\beta_1}+\underbrace{ \eps_{2} +\eps_{4}}_{\beta_2}+\underbrace{\dots}_{\mathrm{zero~allowed}}.
\]
\noindent Case 1.1. $\eps_{3}$ and $\eps_{4}$ both appear with positive coefficients in $\omega$. We claim that $\alpha_2:= \eps_{3}+ \eps_{4} \in \Delta(\gG)\backslash\Delta(\lL)$. Indeed, first, $\alpha_1=(\beta_1-\alpha_2)+\beta_2$ implies that $\alpha_2 \notin\Delta(\kk)$. Second, if $\alpha_2\in\Delta(\nn)$, we could remove $\alpha_1$ on the left-hand side of the relation and substitute $\beta_1+\beta_2$ by $\alpha_2$ to get a relation shorter than the initial one.

We will now prove that $\alpha_2$ is $\bb\cap\kk$-singular.

Assume on the contrary that there exists $\delta\in\Delta^+(\kk)$ such that $\alpha_2+\delta$ is a root. Then $\delta$ is either of the form $\pmf{k}\eps_{k}-\eps_{3}$ or $\pmf{k}\eps_{k}-\eps_{4}$; without loss of generality we may assume that $\delta=\pmf{k}\eps_{k}-\eps_{3}$. The requirement that $\eps_{3}$ and $\eps_{4}$ appear with positive coefficients in $\omega$ implies that there exist $\alpha_3$, $\alpha_4 \in\hw_{\bb\cap\kk}(\gG/\lL)$ such that $\alpha_3=\pmf{j_5}\eps_{j_5} +\eps_{3}$, $\alpha_4= \pmf{j_6} \eps_{j_6}+\eps_{4}$, $\{j_5,j_6\} \cap\{3,4\} = \emptyset$, $1\neq j_5$, and $2\neq j_6$. Furthermore, the preceding assumptions imply that there are at least three distinct roots on the left-hand side of the relation. Since $\alpha_3$ is $\bb\cap\kk$-singular, we have $j_5=k$ and $\delta= \pmf{k} \eps_{k} -\eps_{3} =  \pmf{j_5} \eps_{j_5}-\eps_{ 3}$. Then $\eps_{1}+ (\pmf{j_5} \eps_{j_5}) = \beta_1 + \delta \in \Delta(\nn)$ and therefore $k=j_5\neq 2$. We can now assume without loss of generality that ${j_5}=5$ and the choice of $\pm$ signs is such that $\alpha_3=\eps_5+\eps_3$ and $\delta=\eps_5-\eps_3\in\Delta^+(\kk)$. So far, the assumption that $\alpha_2$ is not $\bb\cap\kk$-singular implies that the relation has the form
\[
\underbrace{\eps_{1}+\eps_{2}}_{\alpha_1}+\underbrace{\eps_{3}+\eps_{5}}_{\alpha_3} +\underbrace{ \eps_{4}+ (\pmf{j_6} \eps_{j_6})}_{\alpha_4} +\underbrace{\dots}_{\mathrm{zero~allowed}} = \underbrace{\eps_{1} + \eps_{3} }_{\beta_1}+\underbrace{ \eps_{2} +\eps_{4}}_{\beta_2}+\underbrace{\dots}_{\mathrm{zero~allowed}}.
\]
We have that $\eps_{5} +\eps_{4}=\alpha_2+\delta\in\Delta(\gG)\backslash\Delta(\lL)$. We claim that $\eps_{5} +\eps_{4}$ is not $\bb\cap\kk$-singular: indeed, otherwise the relation $\alpha_1+\eps_{5} +\eps_{4}= \underbrace{\eps_{1} +\eps_{5}}_{\in\Delta(\nn)} +\beta_2$ would be shorter than the initial one. Therefore there is a root $\delta'\in \Delta^+(\kk)$ such that $\delta'+\eps_{5} +\eps_{4}$ is a root. The $\bb\cap\kk$-singularity of $\alpha_4$ together with $\delta=\eps_5-\eps_3\in\Delta^+(\kk)$ imply that $\delta' = \pmf{j_6}\eps_{j_6}-\eps_{4}$. Therefore $\eps_{2}+(\pmf{j_6} \eps_{j_6}) \in \Delta(\nn)$ and if the weight $\eps_{5}+ (\pmf{j_6}\eps_{j_6})$ is a root, it belongs to $\Delta(\gG) \backslash\Delta(\lL)$. We can write 
\begin{equation}\label{eqTempso(2n)} 
\alpha_1+\eps_{5}+(\pmf{j_6}\eps_{j_6})=\eps_{1}+\eps_{5} + \eps_{2}+ (\pmf{j_6} \eps_{j_6}). 
\end{equation} 
We will arrive at a contradiction for all possible choices of $j_6$. Indeed, if $5\neq j_6$, then $\eps_{5}+(\pmf{j_6}\eps_{j_6})$ is a root. The fact that  $\alpha_2,\alpha_3\in\hw_{\bb\cap\kk}(\gG/\lL)$ together with $\delta, \delta'\in \Delta^+(\kk)$ imply that $\eps_{5}+(\pmf{j_6}\eps_{j_6})$ is $\bb\cap\kk$-singular. Thus \refeq{eqTempso(2n)} is a relation of type (\ref{eqRelation}) which is shorter than the initial one. Contradiction. If $j_6=5$ and the choice of the sign $\pmf{j_6}$ is such that $\alpha_4=\eps_4-\eps_5 $, we get a contradiction as $-\eps_4+\eps_5= \delta' \in \Delta (\kk)$. Finally, if $\eps_{5}=\pmf{j_6}\eps_{j_6}$, then $\delta'':= -\delta+\delta'= \eps_{3}-\eps_{4}\in\Delta(\kk)$. Then depending on whether $\delta''$ is positive or negative we get a contradiction with the $\bb\cap\kk$-singularity of either $\alpha_4$ or $\alpha_3$.

We have now $\alpha_2\in\hw_{\bb\cap\kk}(\gG/\lL)$. Therefore the initial relation is $\alpha_1+\alpha_2=\beta_1+\beta_2$, of type \refeq{eqDcontainsA3'}.

\noindent Case 1.2. One of $\eps_{3}$, $\eps_{4}$ appears with positive coefficient in $\omega$ and the other with non-positive. Without loss of generality we may assume that $\eps_{4}$ appears with positive coefficient in $\omega$ and $\eps_3$ with non-positive. Then there exists a root on the right-hand side, say $\beta_3$, of the form $\pmf{k}\eps_{k}- \eps_{ 3}$. The minimality of the relation implies $\beta_3=\eps_{1} - \eps_{ 3}$. So far the relation is
\[
\underbrace{\eps_{1}+\eps_{2}}_{\alpha_1}+\dots=\underbrace{\eps_{1}+ \eps_{3}}_{\beta_1}+ \underbrace{\eps_{1}- \eps_{3}}_{\beta_3} + \underbrace{\eps_{2}+\eps_{4}}_{\beta_2}+\dots
\]
Now consider $\alpha_2:=\eps_{1}+\eps_{4}$. We claim, as in Case 1.1, that $\alpha_2\in\Delta(\gG)\backslash\Delta(\lL)$. Indeed, first, if we had that $\alpha_2\in\Delta(\nn)$, we could substitute $\beta_1+\beta_2+\beta_3$ by $\alpha_2$ on the right-hand side and remove $\alpha_1$ on the left-hand side to obtain a shorter relation than the initial one. Second, $\alpha_1=((\beta_1-\alpha_2)+\beta_2)+\beta_3$ implies $\alpha_2 \notin\Delta(\kk)$. 

Now, as in Case 1.1, we will show that $\alpha_2$ is $\bb\cap\kk$-singular. Indeed, assume the contrary. The fact that $\eps_{4}$ appears with positive coefficient in $\omega$ implies that on the left-hand side there is a $\bb\cap\kk$-singular weight, say $\alpha_3$, of the form $\alpha_3= \pmf{j_5}\eps_{j_5}+\eps_{ 4} $, where $j_5\neq 2$. 

We claim next that  $j_5\neq 1$. 

Indeed, first, if $\pmf{j_5}\eps_{j_5} =-\eps_{1}$, the relation $\alpha_1+\alpha_3=\beta_2$ is shorter than the initial one. Contradiction. Second, $\pmf{j_5}\eps_{j_5} =\eps_{1}$ contradicts the $\bb\cap\kk$-singularity of $\alpha_2$. Therefore $j_5\neq 1,2$ and we can assume without loss of generality that $j_5=5$ and $\alpha_3 =\eps_5 +\eps_4$. The assumption that $\alpha_2=\eps_{ 1}+\eps_{4}$ is not $\bb\cap\kk$-singular implies that there exists some index $l$ for which at least one of $\gamma:=\pmf{l}\eps_{l} -\eps_{4 }$ and $\delta:= \pmf{l} \eps_{l} - \eps_{1}$ belongs to $\Delta^+(\kk)$. The choice $\delta\in\Delta^+(\kk)$ contradicts the $\bb \cap \kk$-singularity of $\alpha_1$ unless $\delta=\eps_{2}-\eps_{1}$. The latter yields a contradiction as well, as it implies $\alpha_1 \in \Delta( \nn)$. The choice $\gamma \in \Delta^+(\kk)$ together with the $\bb \cap \kk$-singularity of $\alpha_3$ implies $\gamma= \eps_{5} -\eps_{ 4}$. Then $\beta_4 = \beta_2 +\gamma\in\Delta(\nn)$ and the relation $\alpha_1+ \alpha_3 = \beta_1+\beta_3+\beta_4$ is of type (\ref{eqRelation}) and is shorter than the initial one. Contradiction.

So far we have proved that $\alpha_2\in\hw_{\bb\cap\kk}(\gG/\lL)$. Therefore $\alpha_1+\alpha_2=\beta_1+\beta_3+\beta_2$ is the desired relation \refeq{eqDcontainsD4'}.

\noindent Case 1.3.  $\eps_{3}$ and $\eps_{4}$ both appear with non-positive coefficients in $\omega$. As $\eps_{3}$ and $\eps_{4}$ are canceled on the right-hand side without contradicting the minimality of the relation, we need to have $\beta_3:=\eps_{1}-\eps_{3}\in\Delta(\nn)$, $\beta_4: =\eps_{2}-\eps_{4}\in\Delta(\nn)$. Thus we have the desired relation \refeq{eqDcontainsD4''}.

\noindent Case 2. There is no index $i$ such that $\alpha_i=\pmf{j_1}\eps_{j_1}+(\pmf{j_2}\eps_{j_2})$ and both $\pmf{j_1} \eps_{j_1}$ and $\pmf{j_2} \eps_{j_2}$ appear with positive coefficients in $\omega$. As $\omega$ is non-trivial, it has at least one non-zero coordinate. Without loss of generality we may assume this to coordinate to be positive, corresponding to $\eps_1$. In addition, without loss of generality, assume that $\alpha_1 =\eps_{1} +\eps_{2}$. By our current assumption, $\eps_{2}$ appears in $\omega$ with non-positive coefficient. Then some $\alpha_i$, say $\alpha_2$, is of the form $\alpha_2 =-\eps_{2} +(\pmf{j_3} \eps_{j_3}) $. 

\noindent Case 2.1. $j_3\neq 1$. Without loss of generality we can assume that $j_3=3$ and $\alpha_2=-\eps_2+\eps_3$. Then $\beta_1:=\alpha_1+\alpha_2$ is a root and by \refle{leminrelationgeneral} we have the desired relation \refeq{eqDcontainsA2'}.

\noindent Case 2.2. $j_3=j_1$ and $\alpha_2=-\eps_{2}+\eps_{1} $. On the right-hand side, there is a root, say $\beta_1$, of the form $\beta_1=\eps_{1}+(\pmf{j_4}\eps_{j_4})$. A short consideration shows that $j_4\neq 1,2$, and so we assume without loss of generality that $\beta_1:=\eps_{1}+\eps_{4}$. The relation so far has the form
\[
\underbrace{\eps_{1}+\eps_{2}}_{\alpha_1} +\underbrace{(-\eps_{2}+\eps_{1})}_{\alpha_2}+\underbrace{\dots}_{\mathrm{allowed~to~be~zero}} =\underbrace{\eps_{1}+\eps_{4}}_{\beta_1}+\dots.
\]

We will now prove that $\eps_{4}$ appears with positive coefficient in $\omega$. Indeed, assume the contrary. Therefore there exists a root on the left-hand side, say $\alpha_3$, of the form $\alpha_3= \eps_{4}+ (\pmf{j_5}\eps_{j_5})$. By \refle{leminrelationgeneral} we get that $j_5\neq 1,2$, and therefore we can assume without loss of generality that $j_5=5$ and $\alpha_3= \eps_{4} +\eps_{5}$. By the requirement of Case 2, $\eps_5$ appears with a non-positive coefficient in $\omega$, and therefore there exists $\alpha_4=-\eps_5+(\pmf{j_6}\eps_{j_6})$. By \refle{leminrelationgeneral}, $\alpha_4+\alpha_3$ is not a root and therefore $\alpha_{4}= \eps_{4} -\eps_{5}$. Therefore we cannot have a shorter relation than
\begin{equation}\label{eqFakeRelationso(2n)}
\underbrace{\eps_{j_1}+\eps_{2}}_{\alpha_1}+ \underbrace{(-\eps_{2}+\eps_{1})}_{\alpha_2}+\underbrace{ \eps_{3}+ \eps_{4}}_{\alpha_3}+ \underbrace{(-\eps_{4}+\eps_{3})}_{\alpha_4}= 2(\underbrace{\eps_{1}+\eps_{3}}_{\beta_1}).
\end{equation}
We claim that the above expression cannot correspond to a minimal relation. Consider $\delta:= \eps_{1}+\eps_{4}$. First, the possibility $\delta \in \Delta(\kk)$ implies $\beta_1 \in\Delta(\gG)\backslash\Delta(\lL)$. Contradiction. Second, the possibility $\delta\in \Delta(\gG) \backslash \Delta(\lL)$ together with the $\bb \cap \kk$-singularity of $\alpha_1,\alpha_2,\alpha_3$ and $\alpha_4$ imply $\delta \in \hw_{\bb\cap\kk}(\gG/\lL)$. In turn this is contradictory since $\delta + \alpha_4 = \beta_1$ is shorter than (\ref{eqFakeRelationso(2n)}). We conclude $\delta\in\Delta(\nn)$. Since in (\ref{eqFakeRelationso(2n)}), the indices $(1,4)$ are symmetric to $(2,3)$, we conclude that $\delta':= \eps_{2} +\eps_{3} \in\Delta(\nn)$. Finally, $\alpha_1+ \alpha_3 = \delta + \delta'$ is a shorter relation than (\ref{eqFakeRelationso(2n)}). Contradiction.

So far, we have proved that  $\eps_{4}$ appears with a non-positive coefficient in $\omega$. Therefore, on the right-hand side there is a root, say $\beta_2$, of the form $\beta_2= -\eps_{4}+( \pmf{j_5}\eps_{j_5})$. The minimality of the relation implies $\beta_2=\eps_1-\eps_4$. Therefore we have the desired relation $\alpha_1+\alpha_2=\beta_1+\beta_2$ of type \refeq{eqDcontainsD3'}.

\subsubsection{$\gG\simeq\so(2n+1)$}

\noindent Case 1. The relation has a short root on the left-hand side, say $\alpha_1$. Without loss of generality we may assume $\alpha_1=\eps_{1}$. The $\bb\cap\kk$-singularity of $\alpha_1$ implies that $\kk$ has no short roots.

\noindent Case 1.1. $\eps_{1}$ appears with a positive coefficient in $\omega$ and therefore there is a root on the right-hand side, say $\beta_1$, of the form $\beta_1= \eps_{1}+ (\pmf{j_2}\eps_{j_2})$. Without loss of generality we may assume $\beta_1=\eps_1+\eps_2$.

\noindent Case 1.1.1. $\eps_{2}$ appears with a non-positive coefficient in $\omega$. As $\eps_2$ must be canceled out without contradicting the minimality of the relation, one of the roots on the right-hand side, say $\beta_2$, is of the form $\beta_2 = \eps_{1} -\eps_{2}$. It is now clear that we cannot have a relation shorter than (\ref{eqBcontainsB2'''}).
%\begin{equation*}
%2\underbrace{\eps_{1}}_{\alpha_1}=\underbrace{\eps_{1}+\eps_{2}}_{\beta_1}+ \underbrace{\eps_{1}-\eps_{2}}_{\beta_2}.
%\end{equation*}

\noindent Case 1.1.2. $\eps_{2}$ appears with a positive coefficient in $\omega$. The weight $\eps_{2}$ is not a root of $\kk$. Therefore on the left-hand side of the relation there exists a root, say $\alpha_2$, in which $\eps_{2}$ appears with a positive coefficient.

\noindent Case 1.1.2.1. $\alpha_2$ is short, i.e. $\alpha_2=\eps_2$. The relation is \refeq{eqBcontainsB2'}.

\noindent Case 1.1.2.2. $\alpha_2=\eps_2+(\pmf{j_3}\eps_{j_3})$ is long. We claim that $j_3\neq 1$. Indeed otherwise we would have $\alpha_2=\eps_2-\eps_1$, then $\alpha_1+\alpha_2$ would be a root, and by Lemma \ref{leminrelationgeneral} the relation would be $\alpha_1+\alpha_2=\beta_1$. This is impossible. Therefore $j_3\neq 1$, and without loss of generality we can assume $\alpha_2=\eps_2+\eps_3$. Consider $\beta_2:=\eps_3$;  we claim that $\beta_2\in\Delta(\nn)$. Indeed, we immediately see that $\beta_2\notin\Delta(\kk)$, as otherwise $\alpha_1$ would not be $\bb\cap\kk$-singular. Second, assume $\beta_2 \in\Delta(\gG) \backslash \Delta(\nn)$. If $\beta_2$ were $\bb\cap\kk$-singular, we could shorten the relation by removing $\beta_1$ and replacing $\alpha_1+\alpha_2$ by $\beta_2$. Therefore there exists a root $\gamma\in\Delta^+(\kk)$ such that $\beta_2+\gamma$ is a root. The $\bb\cap\kk$-singularity of $\alpha_1$ and $\alpha_2$ implies that $\gamma=\eps_2-\eps_3$. Therefore $\eps_2\in\Delta(\gG)\backslash\Delta(\lL)$. 

Now consider the relation $\eps_1+\eps_2=\beta_1$. If $\eps_2$ were not $\bb\cap\kk$-singular, there would be a positive root $\gamma\in\kk$ such that $\eps_2+\gamma$ is a root but $\eps_2+\eps_3+\gamma$ is not a root, which is impossible. Thus we have a minimal relation of length two of the form $\eps_1+\eps_2=\beta_1$. Hence the initial relation $\eps_1+(\eps_2+\eps_3)+\dots=\beta_1+\dots$ is also of length two. Therefore the unknowns on the right-hand side sum up to $\eps_3$, which together with Lemma \ref{lestreakOfRoots} implies that $\eps_3\in\Delta(\nn)$. Contradiction. Therefore the relation is $\eps_1+(\eps_2 +\eps_3) = (\eps_1+ \eps_2)+\eps_3$ of type \refeq{eqBcontainsB3'}.

\noindent Case 1.2. $\eps_{1}$ appears with a non-positive coefficient in $\omega$. Therefore there is a root, say $\alpha_2$, of the form $\alpha_2=-\eps_{1}+\eps_{2}$. Now \refle{leminrelationgeneral} implies $\alpha_1+\alpha_2\in\Delta(\nn)$ and we get the desired relation \refeq{eqBcontainsB2''}.

\noindent Case 2. Among all minimal relations there is no relation with short roots on the left-hand side. 

\noindent Case 2.1. On the right-hand side there is a short root, say $\beta_1$. Without loss of generality we may assume $\beta_1= \eps_{1}$. As the relation is minimal, $\eps_{1}$  appears with a positive coefficient in $\omega$. Therefore we can assume without loss of generality that $\alpha_1$ is of the form $\alpha_1= \eps_{1}+ \eps_{2}$.

\noindent Case 2.1.1. $\eps_{2}$ appears with a non-positive coefficient in $\omega$. Then there is a root on the left-hand side, say $\alpha_2$, of the form $\alpha_2= -\eps_{2}+(\pmf{j_3}\eps_{j_3})$. If $\alpha_2\neq \eps_2+\eps_1$, we can apply Lemma \ref{leminrelationgeneral} to get a shorter relation than the initial one. Therefore $\alpha_2=-\eps_2+\eps_1$ and the relation is $\alpha_1+\alpha_2=2\beta_1$, of type \refeq{eqBcontainsB2''''}.

\noindent Case 2.1.2. $\eps_{2}$ appears with a positive coefficient in $\omega$. Since $ \eps_{2}$ cannot be a root of $\nn$ (that would imply $\alpha_1\in\Delta(\nn)$), we have a root, say $\beta_2\in\Delta(\nn)$, of the form $\beta_2=\eps_{2}+(\pmf{j_3}\eps_{j_3})$. Since $j_3\neq 1$ we can assume without loss of generality that $\beta_2=\eps_2+\eps_3$. 

\noindent Case 2.1.2.1. $\eps_{3}$ appears with a positive coefficient in $\omega$. Therefore there is a root, say $\alpha_2$, of the from $\alpha_2=\eps_3+(\pmf{j_4}\eps_{j_4})$. We claim that $j_4=1$. Assume the contrary. Since $j_4\neq 2$, we can assume further without loss of generality that $\alpha_2=\eps_{3}+\eps_4$. A short consideration of all possibilities shows that $\alpha_1+ (\eps_3+\eps_4)+ \dots= \eps_1 + (\eps_2 +\eps_3)+\dots$ must be of length at least 3. Consider the root $\eps_3$. If it were in $\Delta(\nn)$ we could shorten the relation by removing $\alpha_1$ and replacing $\beta_1+ \beta_2 $ by $\eps_{ 3}$. If  $\eps_3$ were in $\Delta(\kk)$, we would get $\alpha_1\in\Delta(\nn)$, which is impossible. Therefore $\eps_{3}\in\Delta(\gG)\backslash\Delta(\lL)$. In a similar fashion, we conclude that $\eps_1+\eps_3 \in \Delta (\gG) \backslash \Delta(\lL)$. If at one of the two roots $\eps_{3}$ or $\eps_1+\eps_3$ were $\bb\cap\kk$-singular, we would get a minimal relation of length 2 - either $\alpha_1+\eps_{3}=\beta_1+\beta_2$ or $\alpha_1 +\eps_1 +\eps_{3} =2\beta_1+\beta_2$. Contradiction. Therefore both $\eps_{3}$ and $\eps_1+\eps_3$ are not $\bb\cap\kk$-singular. This shows that there exists $\gamma\in\Delta^+(\kk)$ such that $\gamma+\eps_1+\eps_3$ is a root. Since $\eps_2-\eps_1$ is not a root of $\Delta(\kk)$ and $\alpha_2$ is $\bb\cap\kk$-singular, we obtain that $\gamma= \eps_4 -\eps_3$. Consider $\alpha:=\eps_1+\eps_4 = \gamma+ \eps_1 +\eps_3$. One checks that the $\bb\cap\kk$-singularity of $\alpha_1$ and $\alpha_2$ implies that $\alpha$ is also $\bb\cap\kk$-singular. Therefore we can shorten the relation by replacing $\alpha_1+\alpha_2$ by $\alpha$ and removing $\beta_2$ on the right-hand side. Contradiction.

So far we have established that $j_4=1$. We have immediately a relation of length two, either $\alpha_1+\eps_3+\eps_{1}=2\eps_1+(\eps_2+\eps_3)$ or $\alpha_1+\eps_3-\eps_1=\eps_2+\eps_3$, and so the initial relation is also of length two.  As there can be no two roots on either side that sum up to a root (see Lemma \ref{leminrelationgeneral}), one quickly checks that the only possibility for the minimal relation (up to $\Delta (\gG )$-automorphism) is $(\eps_1+\eps_2)+(\eps_1+\eps_3)=2\eps_1+(\eps_2+\eps_3)$, i.e. type \refeq{eqBcontainsB3''}.

\noindent Case 2.1.2.2. $\eps_{3}$ appears with a non-positive coefficient in $\omega$. Therefore on the right-hand side there is a root, say $\beta_3$, of the form $\beta_3=\eps_2-\eps_3$. We have a relation of length two: $2(\eps_1+\eps_2)=2\eps_1+(\eps_2+\eps_3)+(\eps_2-\eps_3)$ of type \refeq{eqBcontainsB3'''}. In view of the already fixed data, one quickly checks that the only possibility for the initial relation to be of length two is to coincide with this relation.

\noindent Case 2.2. There is no short root on either side of the minimal relation. Therefore we can repeat verbatim the proof for the case $\gG\simeq \so(2n)$ to obtain that we have one of the relations described for this case. 
\subsubsection{$\gG\simeq\sP(2n)$}
\smallskip\noindent Case 1. $\alpha_i+\alpha_j\notin \Delta(\gG)$ for all $i,j$.

\smallskip\noindent Case 1.1 One of the roots $\alpha_i$, say $\alpha_1$, is short. Without loss of generality we may assume $\alpha_1= \eps_{1}+ \eps_{2}$. Since $\alpha_1+\alpha_j\notin\Delta(\gG)$ for all $j$, both $\eps_{1}$ and $\eps_2$ appear with a positive coefficient in $\omega$. Therefore on the right side of the relation there are roots, say $\beta_1$ and $\beta_2$, of the from $\beta_{1} =\eps_{1}+(\pmf{j_3}\eps_{j_3})$ and $\beta_2=\eps_{2} +(\pmf{j_4} \eps_{j_4})$. Consider the vector $\gamma:= \pmf{j_3} \eps_{j_3}+ (\pmf{j_4} \eps_{j_4})$. The minimality of the relation implies that $\gamma$ is non-zero, and therefore that $\gamma$ is a root. If $\gamma\in\Delta(\nn)$ we could shorten the relation by removing $\alpha_1$ on the left-hand side and replacing $\beta_1 + \beta_2$ by $\gamma$. If $\gamma\in\Delta(\kk)$ then $\alpha_1 = (\beta_1- \gamma ) +\beta_2\in\Delta(\nn)$, which is impossible. Therefore $\gamma\in\Delta(\gG)\backslash\Delta(\lL)$. 

As the relation is minimal, $\pmf{j_3}\eps_{j_3}$ and $\pmf{j_4}\eps_{j_4}$  appear with a positive coefficient in $\omega$. Therefore $\pmf{j_3}\eps_{j_3}$ (respectively, $\pmf{j_4}\eps_{j_4}$) appears also in some root, say $\alpha_3$ (respectively, $\alpha_4$) on the left-hand side. If there existed a root $\delta\in\Delta^+(\kk)$ for which $\gamma+\delta$ is a root, $\delta$ would have a negative coefficient in front of one of $\pmf{j_3}\eps_{j_3}$ or $\pmf{j_4}\eps_{j_4}$. This would contradict the $\bb\cap\kk$-singularity of either $\alpha_{3}$ or $\alpha_{4}$. Therefore we have a minimal relation $\alpha_1+\gamma=\beta_1+\beta_2$. Depending on whether $j_3=j_4$ and whether $j_4=2$ our relation is of type \refeq{eqCcontainsA3'}, \refeq{eqCcontainsC3'} or \refeq{eqCcontainsC3''}.

\smallskip\noindent Case 1.2 All roots $\alpha_i$ are long. Without loss of generality we may assume $\alpha_1=2\eps_{1}$. Since $\alpha_1+\alpha_j\notin\Delta(\gG)$ for all $j$, the weight  $\eps_{1}$ appears with a positive coefficient in $\omega$. Therefore there is a root on the right-hand side, say $\beta_1$, of the form $\beta_1 =\eps_{1}+(\pmf{j_2}\eps_{j_2})$. Without loss of generality we may assume $\beta_1=\eps_1+\eps_2$. If $\eps_{ 2}$ appeared with a non-positive coefficient in $\omega$, there would be a cancellation in the right-hand side of the relation. This is impossible. Thus $\eps_{2}$ appears on the left-hand side and we have the desired relation \refeq{eqCcontainsC2''}.

\smallskip\noindent Case 2. For some $\alpha_i$, $\alpha_j$, we have that $\alpha_i+\alpha_j=\gamma$ is a root. By \refle{leminrelationgeneral} $\gamma\in\Delta(\nn)$ and we get one of the relations \refeq{eqCcontainsA2'}, \refeq{eqCcontainsC2'},  (\ref{eqCcontainsC2''}), or (\ref{eqCcontainsC2'''}).
\end{proof}
\begin{corollary}\label{corConeImpliesNotFKFT}
Let $\gG$ be classical simple and suppose that $\lL$ does not satisfy the cone condition. Then the following statements hold.
\begin{itemize}
\item If $\gG\simeq \sL(n), \so(n),$ or $\so(2n)$, there exists an $\lL$-strictly infinite weight $\omega$.
\item If $\gG\simeq \sP(2n)$ there exists an $\lL$-infinite weight $\omega$.
\end{itemize}
%By Corollary \ref{leNotFK} $\lL$ is not a Fernando-Kac subalgebra of finite type for all classical Lie algebras $\gG$.
\end{corollary}
\begin{proof} 
The statement for $\sL(n)$ follows from \cite[Lemma 5.4]{PSZ}, so let $\gG\simeq$ $\so(2n)$, $\so(2n+1)$ or $\sP(2n)$. By \refprop{propabundanceHolds}, we can always pick a minimal relation corresponding to a two-sided weight. By direct observation of all possibilities for minimal relations given in Lemma \ref{leBigTedious} we see that all such relations have a strongly orthogonal decomposition with respect to $\Delta(\nn)$ except when $\gG\simeq\sP(2n)$ and the two-sided weight is given by \refeq{eqCcontainsC3'}. 

Suppose now $\gG\simeq\sP(2n)$ and relation \refeq{eqCcontainsC3'} holds. According to the proof of Lemma \ref{leBigTedious} we can assume the relation has the form
\[
\underbrace{\eps_{1}+\eps_{2}}_{\alpha_1}+ \underbrace{2\eps_{3}}_{\alpha_2} = \underbrace{ \eps_{1}+ \eps_{3}}_{\beta_{1}}+\underbrace{\eps_{2}+ \eps_{3} }_{\beta_{2}}.
\]
Consider the root $2\eps_{1}$. If $2 \eps_{1}$ belonged to $\Delta(\kk)$, we would have the contradictory $\alpha_2=\beta_1- (2 \eps_{1})+\beta_1\in\Delta(\nn)$. Similarly, we get $2\eps_{2} \notin \Delta (\kk)$. If both $2\eps_{2}, 2\eps_{1} \in\Delta(\nn)$, we get the new relation $2\alpha_1=2\eps_{1}+2\eps_{2}$ which corresponds to a two-sided weight (since the relation corresponds to a two-sided weight) and this new relation gives an $\lL$-strictly infinite weight. If one of $2\eps_{2}, 2\eps_{1}$, say $2\eps_{1}$, belongs to $\Delta(\gG/\lL)$, it is also $\bb\cap\kk$-singular (otherwise $\alpha_1$ would fail to be $\bb\cap\kk$-singular as well). Therefore we have a new relation 
\begin{equation}
\label{eqSoProblem}\omega':= 2\eps_{1}+ 2{\eps_{3}}=2\beta_1. 
\end{equation}  

We claim that $\omega'$ is $\lL$-infinite. Indeed, let $\tT$ be the subalgebra generated by $\kk$, $\gG^{\pma \beta_1}$, $\gG^{\pma 2\eps_{1}}$ and $\gG^{\pma 2\eps_{3}}$. Let $\nn ':=\nn\cap\tT$. Since $\tT$ contains the Cartan subalgebra $\hh$, $\nn '$ is a direct sum of root spaces and is therefore generated as a $\kk$-module by $\gG^{ \beta_1}$. Let $\sS_1$ be the simple component of $\kk$ whose roots are linked to $2\eps_{1}$; in case there is no such simple component, set $\sS_1:=\{0\}$. Define similarly $\sS_3$ using $2\eps_3$. Then $\sS_1\cap\sS_3=\{0\}$ as otherwise $\Delta(\nn)$ would contain $-\beta_2$. In addition, each $\sS_i$ must be of type $A$, (otherwise it would have a root $2\eps_{i}$). It follows that $\omega'$ is two-sided with respect to $\tT$, and therefore $\omega'$ is $\lL$-infinite.
\end{proof}

\begin{lemma}\label{leRank2Solvable}
Let $\gG$ be a simple Lie algebra of rank 2 and $\lL$ be a solvable root subalgebra (i.e. $\kk=\hh$) which does not satisfy the cone condition. Then there exists an $\lL$-strictly infinite weight.
\end{lemma}
\begin{proof}
We leave the proof of cases $A_2$, $B_2$ and $C_2$  to the reader. We note that in case of type $B_2$ all relations (\ref{eqBcontainsB2'})-(\ref{eqBcontainsB2''''}) appear; similarly, in case of type $C_2$, all relations (\ref{eqCcontainsA2'})-(\ref{eqCcontainsC2'''}) appear.
\begin{comment}
\noindent\begin{tabular}{ll}
\multicolumn{2}{c}{ $\gG\simeq \so(5)$}\\
$\Delta(\nn)$ & $\lL$-strictly infinite weight \\\hline
one long root & (\ref{eqBcontainsB2'})\\
one short and one long root with scalar product$>0$ & (\ref{eqBcontainsB2''})\\
two perpendicular long roots& (\ref{eqBcontainsB2'''})\\
one short root & (\ref{eqBcontainsB2''''})\\\hline
\multicolumn{2}{c}{ $\gG\simeq \sP(4)$}\\
$\Delta(\nn)$ & type of $\omega$ \\\hline
two perpendicular long roots & (\ref{eqCcontainsC2'})\\
one short root& (\ref{eqCcontainsC2''})\\
one short and one long root with scalar product$>0$ & (\ref{eqCcontainsC2'''})\\
one long root & (\ref{eqCcontainsC2''''})\\\hline
\end{tabular}
\end{comment}

Let now $\gG\simeq G_2$, and fix the scalar product in $\Delta(\gG)$ so that the length of the long root is $\sqrt{6}$. The following table exhibits one $\lL$-strictly infinite weight in each possible case for $\Delta(\nn)$.

\noindent \begin{tabular}{r|c|c}
\multicolumn{3}{c}{$\gG\simeq G_2$}\\
\multicolumn{1}{c|}{$\omega$} & \begin{tabular}{c}Scalar products. All non-listed \\ scalar products are zero. 
\end{tabular} 
&
\begin{tabular}{p{2cm}} The roots from the relation generate\end{tabular}
%\begin{tabular}{l}$\Delta(\nn)$\end{tabular}
\\\hline
$\alpha_1+\alpha_2=3\beta_1$& 
\begin{tabular}{p{4.5cm}}
$\langle\alpha_1,\alpha_2\rangle=\langle\alpha_1,\beta_1\rangle=\langle\alpha_2,\beta_1\rangle=3$, $\langle\alpha_2,\alpha_2\rangle=\langle\alpha_1,\alpha_1\rangle =6$,  $\langle\beta_1,\beta_1\rangle=2$
\end{tabular}
&
$G_2$
%\begin{tabular}{l}one short root \end{tabular}
\\\hline
$\alpha_1+\alpha_2=\beta_1$ & 
\begin{tabular}{p{4.5cm}}
$\langle\alpha_1,\alpha_2\rangle=1$,  $\langle\alpha_1,\beta_1\rangle=\langle\alpha_2,\beta_1\rangle=3$, $\langle\alpha_2,\alpha_2\rangle = \langle \alpha_1, \alpha_1 \rangle =2$, $\langle\beta_1,\beta_1\rangle=6$
\end{tabular}
&
$G_2$
%\begin{tabular}{l} one long root or\\ two long roots with \\ or three long roots\end{tabular}
\\\hline
$2\alpha_1=3\beta_1+\beta_2$ & 
\begin{tabular}{p{4.5cm}}
$\langle\alpha_1,\beta_1\rangle=\langle\alpha_1,\beta_2\rangle=3$, $\langle\beta_1,\beta_1\rangle =2$, $\langle\alpha_1,\alpha_1\rangle=\langle\beta_2,\beta_2\rangle=6$\\
\end{tabular}
&
$G_2$
%\begin{tabular}{l} one short and \\ one perpendicular\\ long root\end{tabular}
\\\hline
$\alpha_1+\alpha_2=\beta_1$ & 
\begin{tabular}{p{4.5cm}}
$\langle\alpha_1,\beta_1\rangle=\langle\alpha_2,\beta_1\rangle=3$, $\langle\alpha_1,\alpha_2\rangle =-3$,  $\langle\alpha_1,\alpha_1\rangle=\langle\alpha_2,\alpha_2\rangle=\langle\beta_2,\beta_2\rangle=6$\\
\end{tabular}
&
$A_2$
%\begin{tabular}{l} two short roots which\\ sum up to  a long root \end{tabular} 
\\\hline
$\alpha_1+\alpha_2=\beta_1$ &
\begin{tabular}{p{4.5cm}}
$\langle\alpha_1,\beta_1\rangle=\langle\alpha_2,\beta_1\rangle=1$, $\langle\alpha_1,\alpha_2\rangle =-3$, $\langle\alpha_2,\alpha_2\rangle=\langle\beta_1,\beta_1\rangle=2$, $\langle\alpha_1,\alpha_1\rangle=6$\\
\end{tabular}
&
$G_2$
%\begin{tabular}{l} one short root and\\ two long roots at\\ angles $\pi/3,\pi/6$ or \\at angles $\pi/2,\pi/3,\pi/6$ \end{tabular}
\\\hline
$\alpha_1+\alpha_2=\beta_1$ &
\begin{tabular}{p{4.5cm}}
$\langle\alpha_1,\beta_1\rangle=3$ $\langle\alpha_2,\beta_1\rangle=-1$, $\langle\alpha_1,\alpha_2\rangle =-3$, $\langle\alpha_1,\alpha_1\rangle=\langle\alpha_2,\alpha_2\rangle=\langle\beta_1,\beta_1\rangle=2$\\
\end{tabular}
&
$G_2$
%\begin{tabular}{l} two short and\\two long roots\\\end{tabular} 
\\\hline
\end{tabular}

\smallskip
\end{proof}


The statement of the following lemma is general, but we will make use of it only for the exceptional Lie algebras.
\begin{lemma}\label{leCentralizerAndNilradical}
Suppose $\nn \cap C(\kk_{ss})$ is not the nilradical of a parabolic subalgebra in $C(\kk_{ss})$ containing $\hh\cap C(\kk_{ss})$. Then the following hold. 
\begin{itemize}
\item[(a)] The cone condition fails. 
\item[(b)] There exists a relation \refeq{eqRelation} of the form given by \refle{leminrelationgeneral}(a) for which $\alpha_1,\alpha_2$ and $\beta_1$ all lie in $\Delta(C(\kk_{ss}))$.  
\item[(c)] There is a relation \refeq{eqRelation} that is $\lL$-infinite.
\end{itemize}
\end{lemma}
\begin{proof}
(a) Suppose on the contrary the cone condition holds. Then there exists $h\in\hh$ such that $h(\beta)>0$ for all $\beta\in\Delta(\nn)$ and $h(\alpha)\leq 0$ for all $\alpha\in\hw_{\bb\cap\kk}(\gG/\lL) \supset \Delta(C(\kk_{ss})$. The element $h$ defines a parabolic subalgebra $(\hh\cap C(\kk_{ss}))+\bigoplus_{\substack{ \gamma\in\Delta(C(\kk_{ss})) \\ \gamma(h)\geq 0}}\gG^\gamma $ of  $C(\kk_{ss})$ whose nilradical is $\nn\cap C(\kk_{ss})$, contradiction. 

(b) Using similar arguments to (a), we see that the cone condition fails when restricted to $\Delta(C(\kk_{ss}))$, i.e. the cones $\Cone_{\ZZ} (\Delta(C(\kk_{ss}))\cap$ $\hw_{\bb\cap\kk}(\gG/\lL))$ and  $\Cone_{\ZZ}$ $(\Delta(\nn) \cap$ $\Delta (C(\kk_{ss})))$ have non-zero intersection. 

Take now a relation \refeq{eqRelation}. Note that $\Delta(C(\kk_{ss}))$ $\cap \hw_{\bb\cap\kk}$ $(\gG/\lL)$ $= \Delta( C(\kk_{ss}))$ $\backslash\Delta(\nn)$. Therefore when we add $-\beta_i$ to both sides of \refeq{eqRelation} we still get a relation of the type \refeq{eqRelation} or zero; thus we can obtain a relation \refeq{eqRelation} with only one term $\beta_1$ on the right-hand side. If we have more than two terms on the left-hand side, by \refle{lestreakOfRoots} we get that the sum of two $\alpha_i$'s must be a root. If that root is in $\Delta(\nn)$, we have found a relation of type given by \refle{leminrelationgeneral}(a); else we can substitute the two roots with their sum and thus reduce the number of terms on the left-hand side. In this fashion, we can reduce the number of summands on the left-hand side to two, which gives the desired relation.

(c) Let $\alpha_1,\alpha_2, \beta_1$ be the roots obtained in (b) and let $\tT$ be the subalgebra generated by $\kk$, $\gG^{\pma\alpha_1}$, $\gG^{\pma\alpha_2}$ and $\gG^{\pma\beta_1}$. Lemma \ref{leRank2Solvable} implies that there exists an $\lL\cap\tT$-strictly infinite weight in $\tT$, which is the desired $\lL$-infinite weight.
\end{proof}

\subsection{Exceptional Lie algebras $G_2$, $F_4$, $E_6$ and $E_7$}
\subsubsection{Exceptional Lie algebra $G_2$}
If $\kk_{ss}=\{0\}$ the existence of an $\lL$-infinite weight is guaranteed by \refle{leCentralizerAndNilradical}. If $\kk_{ss} \neq \{0\}$ it is a straightforward check that, up to a $\gG$-automorphism, the only root subalgebra $\lL=\kk\crplus\nn$ for which the cone condition fails is given by $\Delta(\kk)=\{\pm\gamma_1\}$, $\Delta(\nn)=\{\gamma_1+3\gamma_2, 2\gamma_1+3\gamma_2\}$, where $\gamma_1, \gamma_2$ are positive simple roots of $G_2$ such that $\gamma_1$ is long. For this subalgebra, $(\gamma_1+2\gamma_2)+(\gamma_1+\gamma_2)=2\gamma_1+3\gamma_2$ is the desired $\lL$-(strictly) infinite weight.

Note that $G_2$ is the only simple Lie algebra of rank 2 which admits a non-solvable and non-reductive root Fernando-Kac subalgebra of infinite type (cf. \cite[Example 2]{PS:GenHarishChandra}).

\begin{comment}
\begin{pspicture}(-1,-1)(1,1)
\psline[linewidth=0.6pt, linecolor=blue](0,0)(0.5,0.87)
\psline[linewidth=0.2pt, linecolor=cyan](0,0)(-0.5,0.87)
\psline[linewidth=0.2pt, linecolor=red](0,0)(0.5,-0.87)
\psline[linewidth=0.2pt, linecolor=pink](0,0)(-0.5,-0.87)
\psline[linewidth=0.3pt, linecolor=black](0,0)(1,0)
\psline[linewidth=0.3pt, linecolor=gray](0,0)(-1,0)
\psline[linewidth=0.6pt, linecolor=red](0,0)(0.57,0.32) %0.58, 0.33
\psline[linewidth=0.2pt, linecolor=pink](0,0)(-0.57,0.32)
\psline[linewidth=0.2pt, linecolor=red](0,0)(0.57,-0.32)
\psline[linewidth=0.2pt, linecolor=pink](0,0)(-0.57,-0.32)
\psline[linewidth=0.6pt, linecolor=red](0,0)(0,0.67)
\psline[linewidth=0.2pt, linecolor=red](0,0)(0,-0.67)
\put(1,0.1){\tiny{$\alpha$}}
\put(-0.58,0.43){\tiny{$\beta$}}
\end{pspicture}
\end{comment}


\subsubsection{Exceptional Lie algebras $F_4$, $E_6$, $E_7$}\label{secExceptional}
For a fixed exceptional Lie algebra $\gG$, \refle{leCentralizerAndNilradical} allows us to assume that $\nn\cap C(\kk_{ss})$ is the nilradical of a parabolic subalgebra of $C(\kk_{ss})$ containing $C(\kk_{ss})\cap \hh$. 
The following two lemmas can be proved using a computer; the algorithm we used is described in the next section.
%This cuts down significantly the amount of cases we need to consider (via computer). %That is why, instead of classifying minimal relations, we consider it easier to directly check the existence of $\lL$-strictly infinite and $\lL$-infinite relations under the aforementioned assumption on the nilradical.

\begin{lemma}\label{leE6}~
Let $\gG\simeq E_6$ or $E_7$ with a root subalgebra $\lL=\kk\crplus\nn$ for which the cone condition fails. Suppose in addition that $\nn \cap C(\kk_{ss})$ is the nilradical of some parabolic subalgebra in $C(\kk_{ss})$ containing $\hh\cap C(\kk_{ss})$. Then there exists an $\lL$-strictly infinite relation \refeq{eqRelation} of one of the types listed for $\so(2n)$ in \refle{leBigTedious} or of the type

\noindent\begin{tabular}{rcl| c |c}
\multicolumn{3}{c|}{$\omega$} & \begin{tabular}{p{4.5cm}} Scalar products.\\ All non-listed scalar products are zero.\end{tabular}&  \begin{tabular}{p{1,3cm}}The roots from the relation generate\end{tabular}
\\\hline
\multicolumn{5}{c}{$\gG\simeq E_6, E_7$}\\\hline
$\alpha_1+\alpha_2+\alpha_3$&$=$&$\beta_1+\beta_1+\beta_3$ &
\begin{tabular}{c}
$\langle\alpha_1,\beta_2\rangle=\langle\alpha_1,\beta_3\rangle=$ \\
$\langle\alpha_2,\beta_1\rangle=\langle\alpha_2,\beta_3\rangle=$ \\
$\langle\alpha_3,\beta_1\rangle=\langle\alpha_3,\beta_2\rangle=1$ \\
\end{tabular} & $A_5$.\\\hline
\end{tabular}
\medskip

\noindent When $\gG\simeq E_6$, the above relation occurs only when $\Delta(\kk)\simeq A_1+A_1+A_1$.
\end{lemma}

For the next lemma, we need to define a special root subalgebra of $\gG\simeq F_4$. Fix the scalar product of the root system of $F_4$ so that the long roots have length $2$. Let $\kk$ be defined by the requirement that $\Delta(\kk_{ss})$ be of type $A_1+A_1$ where both $A_1$ roots are long (all such $\kk$ are conjugate, \cite{Dynkin:Semisimple}). Then $C(\kk_{ss})_{ss}$ is of type $C_2\simeq B_2$. Let $\gamma_1$ and $\gamma_2$ be the positive long roots of $\kk$ and $\beta_1$ and $\beta_2$ be the positive long roots of $C(\kk_{ss})$. Let $\beta_0$ be the unique short root of $\Delta(C(\kk_{ss}))$ which has positive scalar products with both $\beta_1$ and $\beta_2$. The roots $\beta_1,\beta_2,\gamma_1$ and $\gamma_2$ are linearly independent. Let $\beta_3$ be given by the requirement $\langle\beta_1,\beta_3 \rangle=0$, $\langle \beta_2, \beta_3 \rangle=2$, $\langle \gamma_1, \beta_3\rangle=0$, $\langle\gamma_2,\beta_3 \rangle=2$ and let $\beta_4$ be given by the requirement $\langle\beta_1,\beta_4 \rangle=0$, $\langle\beta_2,\beta_4\rangle=2$, $\langle \gamma_1, \beta_4 \rangle=2$, $\langle\gamma_2,\beta_4 \rangle=0$. Then $\gG^{\beta_3}$ and $\gG^{\beta_4}$ generate two $\kk$-submodules of $\gG$, say $\nn'$ and $\nn''$, each of dimension 2. Define $\nn$ as the linear span of $\nn',\nn''$, $\gG^{\beta_0}$, $\gG^{\beta_1}$ and $\gG^{\beta_2}$. Then $\nn$ is a nilpotent subalgebra of $\gG$, and is a $\kk$-module. Further, $\dim\nn = 2+2+(1+1+1)=7$ and $C(\kk_{ss})\cap\nn$ is the nilradical of a parabolic subalgebra of $C(\kk_{ss})$. Set $\lL_1:=\kk\crplus\nn$.  

\begin{lemma}\label{leF4} Let $\gG\simeq F_4$. Suppose in addition that $\nn \cap C(\kk_{ss})$ is the nilradical of some parabolic subalgebra of $C(\kk_{ss})$ containing $\hh\cap C(\kk_{ss})$.
\begin{itemize}
\item[(a)] If $\lL$ is not conjugate to $\lL_1$, there exists an $\lL$-strictly infinite relation \refeq{eqRelation} from the list of Lemma \ref{leBigTedious}. Moreover, all relations from Lemma \ref{leBigTedious} except (\ref{eqCcontainsC3'}) do appear.
\item[(b)] If $\lL$ is conjugate to $\lL_1$, there exists an $\lL$-(non-strictly) infinite relation \refeq{eqRelation}. This relation comes from an $\lL':=\lL\cap\tT$-strictly infinite relation in $\tT$, where $\tT$ is one of the two semisimple subalgebras of type $C_3+A_1$ generated by  $\kk$, $C(\kk_{ss})$ and the conjugate of either $\nn'\cup{\nn'}^-$ or $\nn''\cup{\nn''}^-$. The $\lL'$-strictly infinite relation in $\tT$ can be chosen to be isomorphic to relation (\ref{eqCcontainsA2'}).
 
\end{itemize}
\end{lemma}

%\begin{corollary}
%Let $\gG$ be a simple Lie algebra of type $F_4$ or $E_6$ and let the cone condition fail for a root subalgebra $\lL=\kk\crplus\nn$. A minimal relation of type \refeq{eqRelation} is of length 2 for type $F_4$ and of length 2 or 3 for $\gG$ of type $E_6$. A minimal relation of length 3 in $E_6$ occurs only when $\Delta(\kk)\simeq A_1+A_1+A_1$.
%\end{corollary}

Combining \refle{leCentralizerAndNilradical} with Lemmas \ref{leF4} and \ref{leE6} we get the following.

\begin{corollary}\label{corExceptional}
The failure of the cone condition for a root subalgebra $\lL$ of the exceptional Lie algebras of type $F_4, E_6, E_7$ implies the existence of an $\lL$-infinite weight.
\end{corollary} 







\chapter{Combinatorics and algorithms used for the exceptional Lie algebras}\label{secComputerProof}\label{chapterComputer}
This chapter describes the combinatorics and algorithms used in the proofs of the lemmas in section \ref{secExceptional}. The combinatorial facts are presented in a mathematical fashion. The algorithms are described without direct reference to the underlying implementation and programming details, with the only exception of algorithm modifications motivated by practical computer restrictions.

The main reference for this chapter is the fundamental paper of E. Dynkin \cite{Dynkin:Semisimple}. A more elementary discussion of root systems and their combinatorics can be found in J. Humphrey's book \cite{Humphreys}.

In this chapter, we use curly brackets $\{\}$ for non-ordered sets and usual parentheses $()$ for ordered sets.
\section{Outline of the computation for exceptional Lie algebras}\label{secComputationOutline}
In this section, we outline the computation used in the proofs of Lemmas \ref{leE6} and \ref{leF4}. In the subsequent sections, we elaborate on each step of the computation.

We note that for the exceptional Lie algebras, checking the existence of an $\lL$-infinite weight is a finite problem; the present chapter describes how to practically solve it with a computer.

The algorithm has as input the Cartan matrix of a semisimple Lie algebra $\gG$. For a given value of $\lL$, let $S$ be the set of weights of $\hw_{\bb\cap\kk}(\gG/\lL)$ for which $[\gG^{\alpha}, \nn]\subset\nn$ (see Definition \ref{defStronglyPerp}). The output is the following.
\begin{enumerate}
\item[(i)] A list of all possible (up to an automorphism of $\Delta(\gG)$) sets of roots of subalgebras $\lL=\kk\crplus\nn$, for which $C(\kk_{ss})\cap\nn$ is the niradical of a parabolic subalgebra of $C(\kk_{ss})$ containing $\hh\cap C(\kk_{ss})$.
\item[(ii)] A sublist of the list in (i) for which the corresponding subalgebras do not satisfy the cone condition but there exists no $\lL$-strictly infinite weight of length less than or equal to $\max\{\# S, \rk\gG\}$.

\textbf{Remark.} This sublist turns out to be empty for $\gG\simeq E_6, E_7$ and contains one entry for $\gG\simeq F_4$. This entry corresponds to subalgebras conjugate to $\lL_1$, where $\lL_1$ is the subalgebra defined in section \ref{secExceptional}.
\item[(iii)] A list complementary within (i) to the sublist (ii).

\arxivVersion{\textbf{Remark.} The actual list of $\lL$-strictly infinite weights we produced is more detailed; it includes information about the simple direct summands of $\kk$ whose roots are linked to the roots participating in the relation. %The tables in Appendix \ref{secAppendixExceptional} have no two identical entries.
}
\end{enumerate}
The algorithm follows. \arxivVersion{The actual tables printed out for $\gG\simeq F_4$, $E_6$ and $E_7$ are included in the appendix.}
\begin{itemize}
\item Enumerate (up to a $\gG$-automorphism) all reductive root subalgebras $\kk$ containing $\hh$, according to the classification in \cite{Dynkin:Semisimple}.
\item Fix $\kk$. Compute the $\kk$-module decomposition of $\Delta(\gG)$. Then $\nn$ is given by a set of $\kk$-submodules of $\gG$. %A vector subspace  $P$ is a $\kk$-submodule if either $\gamma+\eta\in\Delta(P)$ or $\gamma+\eta$ is not a root for all $\gamma\in\Delta(\kk)$, $\eta\in\Delta(P)$.
\item Compute $\Delta(C(\kk_{ss}))$ (Lemma \ref{leC(kss)}). Compute the group $W'$ of all root system automorphisms of $\Delta(\gG)$ which preserve $\Delta(\bb\cap\kk)$. Note that $W'=W'''\rtimes W''$ is the semidirect product of the Weyl group $W'''$ of $C(\kk_{ss})$ with the group $W''$ of graph automorphisms of $(\Delta(\kk_{ss})\oplus \Delta(C(\kk_{ss}))\cap\Delta(\bb)$ which preserve $\Delta(\kk_{ss})$ and $\Delta(C(\kk_{ss}))$ and extend to automorphisms of $\Delta(\gG)$. \arxivVersion{The tables in Appendix \ref{tableGroups} list the cardinalities of the groups $W''$.}
\item Introduce a total order $\prec$ on the set of all sets of $\kk$-submodules of $\gG$ in an arbitrary fashion.
\item Enumerate all relevant possibilities for $\nn$:
\begin{itemize} 
\item Discard all sets of submodules $\mathcal P$ for which there exists $w\in W'$ with $w\Delta(\mathcal P)\prec \Delta(\mathcal P)$ (act element-wise). 
\item Discard all sets of submodules $\mathcal P$ whose union, intersected with $C(\kk_{ss})$, does not correspond to a nilradical of a parabolic subalgebra of $C(\kk_{ss})$.
\end{itemize}
\item Fix $\nn$.
\item Intersect the two cones $\Cone_\QQ(\Delta(\nn))$ and $\Cone_\QQ(\hw_{\bb\cap\kk}(\gG/\lL))$ (by using the simplex algorithm over $\QQ$ to solve the corresponding linear system of inequalities). If the cones intersect, proceed with the remaining steps.
\item Generate the set of weights $S$.
\item Generate all possible couples $\alpha_1,\alpha_2\in S$ ($\alpha_1=\alpha_2$ is allowed) and compute whether $\alpha_1+\alpha_2$ has a strongly orthogonal decomposition with respect to $\Delta(\nn)$. If no such strongly orthogonal decomposition exists, proceed with all triples, quadruples, \dots, up to  $\max\{\# S, \rk\gG\}$-tuples, until reaching a weight with a strongly orthogonal decomposition with respect to $\Delta(\nn)$. If such a strongly orthogonal decomposition is found, add the found $\lL$-infinite weight and $\Delta(\lL)$ to the list (iii), else add it to the list (ii).
\begin{itemize}
\item We note that if $\omega=b\beta+b_2\beta_2+\dots+b_k\beta_k$ is a strongly orthogonal decomposition with respect to a set $I$ then $\omega-b\beta=b_2\beta+\dots+b_k\beta_k$ is a strongly orthogonal decomposition for $\omega-b\beta$ with respect to set $\beta^\perp\cap I$. The so described relation yields a recursive procedure for computing strongly orthogonal decompositions.
\end{itemize}

\end{itemize}


\section{Computing Weyl groups and Weyl group orbits in $\hh^*$}
To generate a root system from a symmetric Cartan matrix, we have the following short algorithm. Let the simple roots of $\Delta(\gG)$ be $\{\alpha_1,\dots \alpha_n\}$. The Weyl group $\WW$ of $\gG$ is generated by simple reflections (i.e. the reflections $s_{\alpha_i}$ with respect to the simple roots). All roots of the same length lie in the same orbit under the action of the $\WW$ (see for example \cite[\S 10.3, 10.4]{Humphreys}). In other words, $\Delta(\gG)=\{s_{\alpha_1}\dots s_{\alpha_k}(\alpha)\}$ if the root system is of type $A$, $D$ or $E$, and $\Delta(\gG)=\{s_{\alpha_{i_1}}\dots s_{\alpha_{i_k}}(\alpha)\}\cup \{s_{\alpha_{i_1}}\dots s_{\alpha_{i_k}}(\beta)\}$ for the other types, where $\alpha$ is a long root, $\beta$ is a short root, and $s_{\alpha_{i_j}}$ run over all possible simple reflections. This immediately gives the following algorithm.
\begin{itemize}
\item Step 0. Let $S$ be an ordered set of roots and set $S:=(\alpha)$ for root systems with one root length and  $S:=(\alpha,\beta)$ for root systems with two different root lengths, where $\alpha$ is an arbitrary long and $\beta$ an arbitrary short root. Mark the elements of $S$ as ``non-explored''.
\item Step 1. If all elements of $S$ are marked as ``explored'', terminate the program. Else pick the first element $\gamma\in S$ that is marked as ``non-explored''.
\item Step 2. Generate the $n$ different reflection $s_{\alpha_i}(\gamma),\dots,s_{\alpha_n}(\gamma)$ images. Add to $S$ each of these elements that does not already belong to $S$. Mark the newly added roots as ``non-explored''.
\item Step 3. Mark the weight $\gamma$ as ``explored'' and go to Step 1.
\end{itemize}
Step 2 of the above algorithm is executed only once for each root. If in Step 0 in the above algorithm, we set instead $S:=\{\mu\}$ for an arbitrary $\mu\in\hh^*$, the algorithm produces the orbit of $\mu$ under the action of the Weyl group $\WW$. Since $\WW$ acts transitively and irreducibly on the simple bases, the stabilizer of $\rho$ (=the half sum of the positive roots) is the identity and the orbit of $\rho$ enumerates all elements of the Weyl group. If for each weight in the orbit of $\rho$ we record the sequence of reflections sending $\rho$ onto it, we have written down expressions for all elements of $\WW$. Furthermore, if in Step 2 we append each newly found element of the orbit to the end of $S$, we get each element of $\WW$ in reduced form (i.e. written with minimal number of simple reflections). Indeed, this follows as the order of $S$ respects the partial order generated by the length function on the Weyl group.
\section{Computing reductive root subalgebras}
\subsection{Computing the $\kk$-module decomposition}\label{secComputingKmods}
In this section we explain how to compute the $\kk$-module decomposition of $\gG$ for a given a reductive root subalgebra $\kk=\kk_{red}\supset\hh$. Such subalgebras are in one to one correspondence with the root subsystems of $\Delta(\gG)$. In turn, the root subsystems of $\Delta(\gG)$ are parametrized by their respective positive simple bases. Under a positive simple basis we understand a simple basis consisting of positive roots of $\gG$. The following algorithm determines the positive simple basis of a root system from a given generating set of roots $S$. %The set $S$ is ordered.
\begin{itemize}
\item Step 0. If the set $S$ contains negative roots, substitute them with their opposite roots.
\item Step 1. If there exists no couple $\alpha,\beta\in S$ so that $\alpha-\beta$ is a root or zero, terminate the program. Else pick any such couple $\alpha$, $\beta$.
\item Step 2. If $\alpha=\beta$ remove one of the two roots from $S$. Else, if $\gamma=\alpha-\beta$ is a positive root, substitute $\alpha$ by $\gamma$. If $\gamma$ is a negative root, substitute $\beta$ by $-\gamma$.
\item Step 3. Go to Step 1.
\end{itemize}
\textbf{Remark.} We note that the notion of ``positive root'' is relative to a choice of a regular element $h$ of the Cartan subalgebra - i.e. a root $\alpha $ is positive if $\alpha(h)>0$. In a computer realization however, all roots are already given in simple basis coordinates with respect to a fixed simple basis of $\Delta(\gG)$, so a root $\alpha$ is positive if and only if it has positive coordinates.

We now explain how to compute the $\kk$-module decomposition of $\gG$ under the adjoint action of $\kk$. Since $\kk$ contains a Cartan subalgebra, a $\kk$-submodule $M$ of $\gG$ is parametrized by a set of root spaces. Whenever $\alpha+\gamma$ is a root $[g^{\alpha},g^{\gamma}]$ is a non-zero multiple of $g^{\alpha+\gamma}$. Therefore if $g^{\alpha}$ is an element of $M$ and $\gamma+\alpha$ is a root for some $\gamma\in\Delta(\kk)$, then $g ^{\alpha+\gamma}$ is also an element of $M$. We note that whenever $\alpha+\gamma$ is a root for some $\gamma\in\Delta^+(\kk)$, then $\gamma$ can always be chosen to be an element of the positive simple basis of $\kk$. This follows for example from Lemma \ref{lestreakOfRoots}. 

The preceding discussion yields the following algorithm, which, given a starting one-element set of roots $S:=\{\alpha\}$, generates the set of $\hh$-weights of the $\kk$-module generated by $g^\alpha$.
\begin{itemize}
\item Step 0. Mark the initial element of $S$ as ``non-explored''.
\item Step 1. If all elements of $S$ are marked as ``explored'' terminate the program. Else pick a root $\alpha\in S$ marked as ``non-explored''.
\item Step 2. For all simple roots $\gamma_i$ of $\kk$, check whether $\alpha\pm\gamma_i$ is a root. If so, check whether $\alpha\pm \gamma_i$ is already in $S$. If not, add $\alpha\pm \gamma_i$ to $S$ and mark it as ``non-explored''.
\item Step 3. Go to Step 1.
\end{itemize}
To fully decompose $\gG$ as a $\kk$-module, one repeats the preceding algorithm until all roots outside of $\Delta(\kk)$ have been recorded as weights of some module $M$. To that decomposition one adds $\hh\cap C(\kk_{ss})$ and the simple components of $\kk$. The simple components of $\kk$ correspond to the irreducible parts of the root system $\Delta(\kk)$. 
\subsection{Computing Dynkin diagrams of root subsystems}
Given a root subalgebra via its simple basis $(\alpha_1,\dots,\alpha_n)$ (ordered arbitrarily), we can split the simple basis into pairwise strongly orthogonal subsets (corresponding to irreducible root subsystems). Indeed, suppose we have already split the first $k$ roots into irreducible subsets. If $\alpha_{k+1}$ is orthogonal to all preceding roots, we put it into a new subset; if it is non-orthogonal to exactly one of the already computed subsets, we add it there. If $\alpha_{k+1}$ has non-zero scalar product with elements of two or three (more than three is impossible) of the already computed subsets we merge those subsets and add $\alpha_{k+1}$ there. We note that in the preceding discussion, we only need to consider whether $\alpha_{k+1}$ is orthogonal to a subset - strong orthogonality is not needed. The latter is due to the requirement that the starting set $\{\alpha_1,\dots,\alpha_n\}$ be a simple basis of $\kk$.

Once we have a simple basis $\alpha_1,\dots,\alpha_n$ such that the corresponding Dynkin diagram is connected, finding its Dynkin diagram is straightforward. If the diagram has no triple node (i.e. a root connected to three other roots), then it must be of types $A$, $B$, $C$, $F$ or $G$, and considering the number of roots of different lengths (and the lengths themselves), we can determine the Dynkin diagram. If the diagram has a triple node, then one counts the sizes of the three components of the Dynkin diagram obtained by removing this triple node. The so obtained sizes distinguish between types $D$ and $E$.

\subsection{Computing root subsystem isomorphisms}\label{secRootSubsystemIsos}
Proposition \ref{propRootSubsystemsUpToIso} below gives a necessary and sufficient criterion for the existence of automorphism of $\Delta(\gG)$ which maps a root subsystem $\Delta_1$ to another root subsystem $\Delta_2$. In this section we describe how to explicitly compute all such automorphisms. Our algorithm is natural, however somewhat ad-hoc; there could exist a computationally more effective way of doing this. 

We note that the computing the structure of the groups of automorphisms of $\Delta(\gG)$ which preserve a fixed root subsystem appears to be an interesting computational exercise, outside of the scope of the present thesis.

Let $\Delta_1$ and $\Delta_2$ be two root subsystems of $\Delta(\gG)$. Let $\phi$ be a an automorphism of $\Delta(\gG)$ mapping $\Delta_1$ to $\Delta_2$. Then $\phi$ must also map $\Delta_1^\sperp$ onto $\Delta_2^\sperp$, hence the Dynkin diagrams of $\Delta_1^\sperp$ and $\Delta_2^\sperp$ (as well as the diagrams of $\Delta_1$ and $\Delta_2$) must be isomorphic. This is in fact sufficient for such a $\phi$ to exist (Proposition \ref{propRootSubsystemsUpToIso} below). Suppose now $\Delta_1^\sperp$ and $\Delta_2^\sperp$, as well as $\Delta_1$ and $\Delta_2$, have the same Dynkin diagrams. 

The map $\phi$ can be decomposed uniquely as the composition of a graph isomorphism between $\Delta_1$ and $\Delta_2$ and an inner automorphism of $\Delta_2$. In turn, the graph isomorphism mapping $\Delta_1$ to $\Delta_2$ can be decomposed uniquely as the composition of a permutation of isotypic components and a sequence of graph automorphisms for each connected component of $\Delta_2$. We note that a non-trivial graph automorphism of irreducible root systems exist for types $A_n$ for $n\geq 2$, $D_n$ for $n\geq 4$ and $E_6$. Similar considerations hold for the action of $\phi$ on $\Delta_1^\sperp$. 

In the present thesis, we do not present a sufficiency criterion for which graph isomorphisms between two root subsystems $\Delta_1$ and $\Delta_2$ extend to an automorphism of $\Delta(\gG)$; as the tables in the appendix show, not all such graph isomorphisms do. Instead, for fixed data $\Delta_1, \Delta_2\subset\Delta(\gG)$, we summon our computational power and try extending all possible graph isomorphisms. As the current chapter is aimed at exceptional Lie algebras, the upper bound for the graph isomorphisms between two root subsystems of $E_8$ is only $8!=40320$ (corresponding to two root subsystems of type $8A_1\subset E_8$). Since this computation presents no challenge for a modern computer, we have not made any attempts to optimize our procedure.

We will now describe how to check whether a given graph automorphism between two root subsytems extends to an automorphism of $\Delta(\gG)$. Let $D':=\Delta_1\cup \Delta_1^\sperp $ and $R':=\Delta_2\cup \Delta_2^\sperp$. Suppose in addition we have a linear map $\phi:\linspan R'\to \linspan D'$ that maps $D'$ to $R'$ and corresponds to a graph isomorphism between $\Delta_1$ and $\Delta_2$ and  $\Delta_1^\sperp$ and $\Delta_2^\sperp$. If $\rk \Delta_1+\rk \Delta_1^\sperp=\dim \hh^*$, the map $\phi$ is determined as a linear map on the entire $\hh^*$ and is injective. In this case, $\phi$ is a root system isomorphism of $\Delta(\gG)$ if and only if for all $\alpha\in\Delta(\gG)$ we have that $\phi(\alpha)$ is a root. 

The preceding discussion does not indicate what should be done in case $\rk \Delta_1+\rk \Delta_1^\sperp<\rk\Delta(\gG)$. Let $D$ and $R$ be linearly independent ordered sets of roots of equal cardinality and let the map $\phi_{DR}:\linspan D\to \linspan R$ be defined by the requirement that $\phi$ maps the $i^{th}$ element of $D$ to the  $i^{th}$ element of $R$. The algorithm below returns ``TRUE'' whenever $\phi_{DR}$ can be extended to an automorphism of $\Delta(\gG)$ and ``FALSE'' otherwise. 

\begin{itemize}
\item Step 0. Input: an ordered set of roots $D=(\alpha_1,\dots,\alpha_k)$ and an ordered set of roots $R:=(\beta_1,\dots, \beta_k)$. %The sets $D$ and $R$ determine a linear map $\phi_{DR}$ by $\phi_{DR}(\alpha_i)=\beta_i$ between $\linspan_{\QQ} ~D$ and $\linspan_{\QQ} ~R$.
\item Step 1. If $\rk~D\neq \rk R$ return ``FALSE''.
\item Step 2. If $\#D\neq\rk(\Delta(\gG))$, proceed to Step 4. Otherwise, if $\#D\neq\rk(\Delta(\gG))$, then  $\phi_{DR}$ is completely determined by linearity. If $\phi_{DR}(\alpha)$ is a root for all $\alpha\in\Delta(\gG)$, return ``TRUE'', else return ``FALSE''.
\item Step 3. If the Dynkin diagram generated by $D$ is isomorphic to the Dynkin diagram generated by $R$ proceed to Step 4; else return ``FALSE''.
\item Step 4. Let the root subalgebras generated by $D$ and $R$ be respectively  $\kk_D$ and $\kk_R$. Compute the $\kk_D$-module (respectively, $\kk_R$-module) decomposition of $\gG$. If the two decompositions are different return ``FALSE''. Else fix a $\kk_D$-submodule $M$ of $\gG$ and a root that is a weight of $M$ with the requirement that $\alpha$ is linearly independent from $D$. For each $\kk_R$-module $M'$ for which $\dim M=\dim M'$ and for each root $\beta$ that is a weight of $M'$ for which $\beta$ is linearly independent from $R$, check (by invoking recursion) whether $\phi_{D\cup\{\alpha\}, R\cup\{\beta\} }$ can be extended to an automorphism of $\Delta(\gG)$. Upon discovery of a couple $\alpha, \beta$ for which such an extension exists, stop the current branch of the computation and return ``TRUE''. If no such couple exists, return ``FALSE''. We note that in this step the algorithm branches and calls itself recursively.
\end{itemize}

\subsection{Root subsystems up to isomorphism}
The fundamental paper \cite{Dynkin:Semisimple} classifies all regular semisimple subalgebras of a simple Lie algebra. Regular semisimple subalgebras are by definition subalgebras whose Cartan subalgebra can be extended to a Cartan subalgebra $\hh$ of the ambient Lie algebra, such that the vector space generated by $\hh$ and the starting subalgebra is closed under the Lie bracket. The regular subalgebras are conjugate to the semisimple part $\kk_{ss}$  of some reductive root Lie algebra $\kk$. The classification of regular semisimple Lie algebras therefore corresponds to the classification of reductive root subalgebras and hence to the classification of root subsystems. 

The classification of root subsystems of the irreducible root systems carries over directly to a classification for an arbitrary root system. That is why, until the end of this section, we fix $\Delta(\gG)$ to be an irreducible root system. 

In order to classify the root subsystems of $\Delta(\gG)$, Dynkin first classifies the root subsystems of maximal rank. It turns out that two maximal rank root subsystems are conjugate by an automorphism of $\Delta(\gG)$ if and only if their Dynkin diagrams are the same. For the classical Lie algebras the proofs are straightforward. Indeed, one sees that a root system of type $D_l$ that lies inside a root system of type $B_n$ or $D_n$ is of the form $\{\pm\eps_a\pm\eps_b|a,b\in I\subset \{1,\dots, n\} \}$ for some set $I$ with $\# I=l$. Similarly, a root system $B_l$ that lies inside $B_n$ is of the form $\{\pm\eps_a\pm\eps_b|a,b\in I\subset \{1,\dots, n\}\cup \}\{\pm\eps_a|a\in I\} $. Finally, a root system $C_l$ that lies inside $C_n$ is of the form $\{\pm\eps_a\pm\eps_b|a,b\in I\subset \{1,\dots, n\}\cup \{\pm2\eps_a|a\in I\} \}$. Therefore, up to an automorphism of $\Delta(\gG)$, the maximal root subsystems in types $B_n$, $C_n$ and $D_n$ are in a one to one correspondence with the partitions of $n$. Indeed, the reader can easily verify that the root subsystems corresponding to the same partition are indeed isomorphic via an automorphism of $\Delta(\gG)$ (the automorphism is not necessarily inner in type $D$, \cite{Dynkin:Semisimple}).

We note that Table 9 in \cite{Dynkin:Semisimple} allows root subsystems of type $B_1$, $C_1$, $D_2$ and $D_3$. These have no meaning as abstract root systems, but only as root subsystems. They can be defined in an invariant fashion as follows. $B_1$ is a rank 1 root subsystem whose positive root is short, and $C_1$ is a rank 1 root subsystem whose positive root is long. For $\Delta(\gG)$ of type $D_n$ or $B_n$ with $n\geq 5$, we can define $D_2$ (respectively, $D_3$) as the root subsystems of type $A_1+A_1$ (respectively, $A_3$) that cannot be extended to a root subsystem of type $A_4$. We note that this definition applies to the root subsystems of $D_4$, $B_4$, as well - as $A_4\nsubseteq D_4$, all root subsystems of type $A_2$ in $D_4$ or $B_4$ are also of type $D_2$, and all root subsystems of type $A_3$ in $D_4$ and $B_4$ are also of type $D_3$.

We can now summarize the algorithm in \cite{Dynkin:Semisimple} for enumerating all root subsystems of a given root system. At the start, one initializes the list of root subsystems with the Dynkin diagrams of the maximal rank root subsytems. In addition, one labels the vertices of the diagrams with explicit root realizations. For each Dynkin diagram one erases one point at a time, as well as the incident edges. In case the newly obtained diagram is not already contained (up to relabeling of the vertices) in the list, one adds it to the bottom of the list. If the diagram is already contained up to relabeling of the vertices in the list, one checks whether the labels correspond to root subsystems conjugate by an automorphism of $\Delta(\gG)$. If that is not the case, one adds the newly obtained diagram to the bottom of the list. Repeating the above procedure, one works ``top-down'' to exhaust all possible root subsystems of the starting root system.

The algorithm from \cite{Dynkin:Semisimple} is practical for a human, however we have chosen a more computational approach based on section \ref{secRootSubsystemIsos}. Instead of working ``top-down'', we describe an intuitive ``bottom-up'' algorithm for generating all root subsystems. In the algorithm that follows, $\Delta$ parametrizes the elements of a root subsystem and $R$ stands for the output list of root subalgebras. 
\begin{itemize}
\item Step 0. Initialize $\Delta$ to be the empty set, and $R$ to be the empty list.
\item Step 1. For each root $\alpha$ in $\Delta(\gG)\backslash\Delta$ set $\Delta=\Delta\cup \{\alpha\}$. Here the computation branches (corresponding to a ``for loop'').
\item Step 2. Transform $\Delta$ to a simple basis (see section \ref{secComputingKmods}).
\item Step 3. Check whether $\Delta$ is isomorphic via an automorphism of $\Delta(\gG)$ to one of the root subsystems already recorded in $R$ (see section \ref{secRootSubsystemIsos}). If so, terminate the current branch of the computation. Otherwise, add $\Delta$ to $R$ and go to Step 1.
\end{itemize}
We are now in a position to state the following necessary and sufficient condition for two root subsystems to be conjugate by an automorphism of $\Delta(\gG)$.
\begin{prop}\label{propRootSubsystemsUpToIso}
For two root subsystems $\Delta_1$ and $\Delta_2$ to be conjugate via an automorphism of $\Delta(\gG)$ it is necessary and sufficient that their Dynkin diagrams and the Dynkin diagrams of $\Delta_1^\sperp$ and $\Delta_2^\sperp$ be isomorphic.
\end{prop}
\textbf{Remark.} This fact was not noted in \cite{Dynkin:Semisimple}. 

For the classical Lie algebras, the proof of this lemma is a straightforward check using the explicit form of their root subsystems. For the exceptional Lie algebras, using the preceding algorithm one can enumerate all root subsystems up to an automorphism of $\Delta(\gG)$ and check the statement. As this straightforward check was already performed by our ``vector partition'' program, we see no need to write down a detailed proof.

The above proposition allows us to substitute Step 3 in the preceding algorithm by the simpler procedure of comparing the Dynkin diagram types of $\Delta$ and $\Delta^\sperp$. The so modified algorithm has no speed disadvantage compared to the original algorithm from \cite{Dynkin:Semisimple}, but is perhaps simpler to program.



\section{Enumerating root subalgebras up to isomorphism}\label{secEnumeratingRootSAs}
In this section we explain how to generate all possible nilradicals $\nn$ such that $\kk\crplus \nn$ is a Lie subalgebra with reductive part the root subalgebra $\kk$. By section \ref{secComputingKmods}, we can compute the $\kk$-module decomposition of $\gG/\kk$ into irreducible submodules. Let this decomposition be $M=M_1\oplus \dots\oplus M_k$. Thus the nilradical $\nn$ is parametrized by a subset of the set $\{M_1,\dots M_k\}$. Denote by $\weights(M_i)$ the $\hh$-weights of the module $M_i$; the weights of $M_i$ are also roots of $\Delta(\gG)$.

We introduce the notions of opposite irreducible $\kk$-submodules of $\gG/\kk$ and a pairing table of irreducible $\kk$-submodules of $\gG/\kk$.  We call $M_i$ and $M_j$ \emph{opposite} if $\weights(M_i)=-\weights(M_j)$. A $\kk$-submodule can be opposite to itself - for example  that is always the case if $\gG/\kk$ is irreducible. %One can quickly figure out of all possibilities for $\kk$ such that $\gG/\kk$ is an irreducible $\kk$-module. In the exceptional Lie algebras there are further interesting examples for which the $\kk$-module $\gG/\kk$ is irreducible, such as $\kk$ of type $A_5+A_1$ in $E_6$, $B_4$ in $F_4$, $A_7$ in $E_7$, $D_8$ in $E_8$, etc. We note that the requirement that $\kk \supset \hh$ is maximal in $\gG$ implies neither that $\gG/\kk$ has a unique irreducible $\kk$-submodule nor that there is an $M_i$ opposite to itself - for example $\sL(3) \times \sL(3) \times \sL(3)$ (type $3A_2$) is maximal in $E_6$ but $E_6/(\sL(3)\times\sL(3)\times \sL(3)) $ decomposes as the sum of two $\sL(3)\times\sL(3)\times \sL(3)$-modules opposite to each other.
We say that \emph{$M_i$ and $M_j$ pair to $M_k$} if there exist roots $\alpha\in \weights (M_i)$ and $\beta\in\weights (M_j)$ such that $\alpha+\beta$ is a root lying in $\weights (M_k)$. We make the following elementary observation.

\begin{lemma}
For two fixed indices $i,j$ ($i=j$ allowed) there exists at most one index $k$ such that $M_i$ and $M_j$ pair to $M_k$. 
\end{lemma}
\begin{proof}
Suppose $M_i$ and $M_j$ pair to $M_k$, i.e. there are roots $\alpha_1\in \weights( M_i)$, $\alpha_2\in\weights(M_j)$ and $\alpha_3\in\weights(M_k)$ such that $\alpha_1+\alpha_2=\alpha_3$. Let $\alpha_1'\in \weights(M_i)$ and $\alpha_2'\in \weights (M_j)$ be two roots such that $\alpha_1'+\alpha_2'=:\alpha_3'$ is a root. Clearly there exist roots $\gamma_i\in\Delta(\kk)$ and $\delta_j\in\Delta(\kk)$ such that $\alpha_1'=\alpha_1+\gamma_1+\dots +\gamma_l$ and $\alpha_2'=\alpha_2+\delta_1+\dots +\delta_m$. Then we have that
\[
\alpha_3+\gamma_1+\dots +\gamma_l+\delta_1+\dots +\delta_m=\alpha_3'.
\]
Applying Lemma \ref{lestreakOfRoots} consecutively, we see that the summands in the above expression can be reordered so that for any $s\geq 1$, the sum of the first $s$ summands yields a root (zero is not a possibility, as it would imply that $\alpha_3'\in \Delta(\kk)$). Furthermore, this reordering can be made so that $\alpha_3$ remains in the first position. This implies that ${\alpha_3'}\in \weights( M_k)$, as desired. 
\end{proof}

We define a \emph{pairing table} to be a $k\times k$ table of indices, such that on the $i^{th}$ row and $j^{th}$ column we write the index $k$  if $M_i$ and $M_j$ pair to $M_k$ and 0 otherwise. A subalgebra $\kk\crplus\nn$ is given by a subset $A\subset\{1,\dots,k\}$ of indices % $\{M_1,\dots, M_k\}$ %, parametrized by an ordered subset $A$ of the indices $\{1,\dots ,k\}$, 
with the following two properties. 
\begin{itemize}
\item $A$ contains no two indices $a,b$ such that $M_{a}$ and $M_{b}$ are opposite. 
\item If $M_{a}$ and $M_{b}$ pair to $M_c$, then $c\in A$ (i.e. $A$ is closed with respect to the pairing table).
\end{itemize}
\subsection{An algorithm enumerating all possible nilradicals}\label{secEnumerateAllNilradicals}
Before we present an algorithm for enumerating all possible nilradicals of $\lL$ for a fixed value of $\kk$, we note that the actual implementation is essential as the total number of subsets of $\{1,\dots,k\} $ is $2^k$. Therefore the most naive approach for generating nilradicals - enumeration of all subsets of $\{1,\dots,k\} $ and checking the two conditions for being a valid nilradical parametrization - would not work sufficiently fast enough for our purposes.

In the algorithm below, we parametrize a subset of $A$ of $\{1,\dots,k\}$ by a $k$-tuple of 0's and 1's. In our parametrization, the value 1 on $l^{th}$ position implies that $l\in A$, and the value 0 - that $l\notin A$. %In this fashion, we have also parametrized the irreducible submodules of $\gG/\kk$ by sequences of $k$ zeroes and ones.
We obtain a total order $\succ$ on the subsets of $\{1,\dots,k\}$ by setting $A\succ B$ if the number, recorded in base two by the parametrization of $A$ is larger than the corresponding number for the set $B$. This total order respects the partial order of subset inclusion. The order $\succ$ on the subsets of $\{1,\dots, k\}$ induces an order on the nilradicals $\nn$; we denote it using the same symbol $\succ$. 

The following algorithm takes as input a number $l\leq k$ and an subset $A$ of $\{1,\dots,k\}$. The algorithm then enumerates all subsets $A'\subset \{1,\dots,k\}$ such that $A\subset A'$, such that $A'$ is a valid nilradical parametrization, and such that $A'\cap \{1,\dots,l\}=A\cap\{1,\dots,l\}$. We note that setting as input $l=0$ makes the last requirement trivial. We give larger values to $l$ in section \ref{secComputingParabolics} below to make a computational optimization related to Lemma \ref{leCentralizerAndNilradical}. 
\begin{itemize}
\item Step 1. Input: A $k$-tuple $B=(f_1,\dots, f_k)$ of zeroes and ones and an index $0\leq l\leq k$. The $k$-tuple $B$ is required to be a valid nilradical parametrization. More precisely, $B$ must satisfy two conditions. First, for all pairs $(a,b)$, where $a=b$ is allowed, if $f_a=f_b=1$, then $M_a$ is not opposite to $M_b$. Second, if $f_a=f_b=1$ and the modules $M_a$ and $M_b$ pair to $M_c$, then $f_c=1$.
\item Step 2. 
\begin{itemize}
\item If $l=k$ then $B$ is a valid nilradical parametrization. We can perform a computation\footnote{ In practice, for simple Lie algebras of higher rank, we do not have enough RAM memory to store all values of $B$ produced in this step; we have to discard the so generated values of $B$.} using $B$. Terminate the current branch of the computation.
\item  If $l<k$, copy the value of $B$ to a new $k$-tuple $A=(e_1,\dots, e_k)$ and go to the next step.
\end{itemize}
\item Step 3. Branch the computation into two cases. First branch: mark the index $l+1$ as ``newly added'', set $e_{l+1}=1$ and go to Step 4. Second branch: set $e_{l+1}=0$ and go to Step 5.
\item Step 4. Ensure that $A$ is a valid nilradical parametrization and respects the total order $\succ$. More precisely, we transform $A$ with the following subroutine. Within the subroutine, we write ``terminate the subroutine'' to exit from the subroutine and go to Step 5.
\begin{itemize}
\item Step 4.1. If there exists no index marked as ``newly added'', terminate the subroutine; otherwise pick the first index $s$ marked as ``newly added''.
\item Step 4.2. For each index $i$ for which $e_i=1$, check whether $M_i$ and $M_s$ have a non-zero pairing. If $M_i$ and $M_s$ pair to an index $t$, and $e_t=0$, set $e_t=1$ and mark the index $t$ as ``newly added''. Then apply the following two rules.
\begin{itemize}
\item If $t<l+1$, we have that $A$ is larger with respect to $\succ$ than the starting value of $A$. Terminate the current branch of the computation. 
\item If $t'$ is the index of the module opposite to $M_t$ and $e_{t'}=1$, $A$ does not correspond to a nilradical. Terminate the current branch of the computation.
\end{itemize}
\item Step 4.3. Go to Step 4.1.
\end{itemize}
\item Step 5. Call Step 1 recursively with input $B:=A$ and $l:=l+1$.
\end{itemize}
It can be proved that the algorithm requires at most $O(k^2(\mathrm{~the~size~of~the~output}))$ operations, which is practically applicable for computations in exceptional Lie algebra $E_7$ (and possibly also in $E_8$ with the use of a powerful computer).

\subsection{Modification to get 1 element in each $W'$-orbit }\label{secComputation1elementInOrbit}
Let $W'$ be the group of automorphisms of $\Delta(\gG)$ which preserve $\Delta(\bb\cap\kk)$. The current section describes a modification of the algrorithm in section \ref{secEnumerateAllNilradicals} so that it produces exactly one representative in each $W'$-orbit of the nilradicals - namely, the maximal element with respect to the order $\succ$.

As follows from section \ref{secRootSubsystemIsos}, $W'$ is generated by graph automorphisms of the Dynkin diagram of $\Delta(\bb\cap\kk)$ and by reflections with respect to a simple basis of $\Delta(C(\kk_{ss}))$. Thus we can explicitly compute the elements of $W'$ as a sequence of graph automorphisms and simple reflections. In this way, we can compute the action of $W'$ on the set $M_1,\dots, M_k$ and thus on the set of all nilradicals for a fixed $\kk$. We note that this action is given by a subgroup of the symmetric group of $k$ elements.

In order to modify the algorithm in the preceding section, we need to do two things. First, we must first require that the input $B$ parametrize a set that is maximal with respect to $\succ$ in its $W'$-orbit. Second, we must check when going from Step 4 to Step 5 that the value of $A$ produced in Step 4 is maximal with respect to $\succ$. If not, we should terminate the branch of the computation. The so modified algorithm can be proved to require at most $O(\max\{k^2, \# W'\}(\mathrm{~the~size~of~the~output}))$ operations.
\subsection{Generating parabolic subalgebras of $C(\kk_{ss})$ }\label{secComputingParabolics}
Lemma \ref{leCentralizerAndNilradical} guarantees the existence of an $\lL$-prohibiting weight if $\nn\cap C(\kk_{ss})$ is not the nilradical of a parabolic subalgebra of $C(\kk_{ss})$. That is why, for the purpose of searching for $\lL$-infinite weights, we can use the algorithm in section \ref{secEnumerateAllNilradicals} to enumerate only the nilradicals $\nn$ for which $\nn\cap C(\kk_{ss})$ is the nilradical of a parabolic subalgebra of $C(\kk_{ss})$. In addition, we can request that the latter parabolic subalgebra contain $\bb\cap C(\kk_{ss})$. 

In order to do enumerate the nilradicals as above, we do not need to make any changes to the  algorithm in section \ref{secEnumerateAllNilradicals}. Rather, we need to do two things. First, we need to reorder the modules $M_i$ so that the one dimensional ones (corresponding to the roots of $C(\kk_{ss})$) come first, i.e. have indices $1,\dots, l$, where $l:=\#\Delta(C(\kk_{ss}))$. Second, we need to choose special values for the input data - the $k$-tuple $B$. Note that a parabolic subalgebra of $C(\kk_{ss})$ containing $\bb\cap C(\kk_{ss})$ is uniquely determined by the positive simple root spaces which lie in its nilradical. For each simple root of $\bb\cap C(\kk_{ss})$ we have two options and so we have a total of $2^{\rk \Delta(C(\kk_{ss}))}$ parabolic subalgebras of $C(\kk_{ss})$ containing $\bb\cap C(\kk_{ss})$. Thus each possible parabolic subalgebra of $C(\kk_{ss})$ containing $\bb\cap C(\kk_{ss})$ determines initial values for the first $l$ elements of the $k$-tuple $B$; we set the remaining elements of $B$ to be zero. Thus we obtain the desired initial values for $l$ and $B$. 

There can exist isomorphisms between the parabolic subalgebras of $C(\kk_{ss})$ induced by the group $W'$ defined in section \ref{secComputation1elementInOrbit}. That is why, in order to make the current optimization compatible with the one in section \ref{secComputation1elementInOrbit}, we need to additionally transform $B$ by the action of $W'$ to ensure $B$ is maximal with respect to the order $\succ$.

\section{Notes on the software implementation}
The preceding sections covered computational tools specific to Lie theory. In this section, we make a non-comprehensive list of some of the additional tools we needed in order to perform the computation outlined in section \ref{secComputationOutline}. The tools range from general computer science techniques to computer algebra specific algorithms. All computational modules were written from scratch by the author in the programming language C++. The resulting computer program was named ``vector partition'' program.

\begin{itemize}
\item A library for making computations with large rational numbers. Whenever the program multiplies two rational numbers so that the sizes of the numerators/denominators exceed the allocated memory for the purpose, the program dynamically allocates extra memory to fit the computation. 
\item Simplex algorithm. For computations up to $E_8$, we need to solve simplex problems given by matrices of sizes no larger than 240x9. That is why we did not need any special implementation; we have used standard Internet sources \\ (\url{http://en.wikipedia.org/wiki/Simplex_algorithm}) for description of the algorithm and have otherwise implemented it from scratch.
\item Algorithms for efficiently manipulating and enumerating permutations, combinations, variations, subsets with and without multiplicities.
\item Containers for matrices of elements of arbitrary type. 
\item Containers for data structures and RAM memory management. Those include hashed lists and arrays of elements of arbitrary type whose memory use expands on demand.
\item Routines for creating tables in \LaTeX~ and .html format. All tables in the appendices are automatically generated.
\end{itemize}
\subsubsection{Information on the actual software tools}
The ``vector partition'' program was compiled on two operational systems - Linux (Ubuntu 10.04) and Windows (XP), with two different compilers - gcc and the Microsoft C++ compiler. The source code was modified in two different C++ programming environments - Microsoft Visual Studio Express 2008 and Code::Blocks 8.02. The web-server of the ``vector partition'' program\\ (\url{http://vector-partition.jacobs-university.de/cgi-bin/vector_partition_linux_cgi?rootSAs})  currently runs on Apache 2.2.13 web server, running on Linux OpenSUSE operational system, on machine provided by Jacobs University. The installation and support of the Apache 2.2.13/ Linux OpenSUSE is courtesy of the IT support of Jacobs University, and in particular of A. Gelessus and S. Schmidt. The ``vector partition'' program is activated by the Apache server via CGI (Common Gateway Interface). 

The ``vector partition'' program also has a graphical user interface written with the cross-platform tool wxWidgets, \url{http://www.wxwidgets.org/}. The tool runs on Linux, Windows and Macintosh.

The program is licensed under the Library General Public License (LGPL) 3.0. This means that the program and its code can be freely downloaded, used and modified by anyone. The code is hosted on the Open Source hosting site  \url{http://sourceforge.net/}. The project page of the vector partition function is located at \url{http://vectorpartition.sourceforge.net/}. The large volume of source code (approximately 35 000 lines of code at the time of writing of this thesis) is handled via the Sub-VersioN (SVN) protocol and related Open Source software tools. The SVN source code management allows reverting of source code modifications, comparisons of different versions, and many other techniques used in general software development. The SVN tool provides an excellent environment for collaboration and would easily allow more developers to join the project and coordinate efforts.



%\begin{lemma}\label{leExternalAuts}
%Let $\kk\subset\gG$ be a reductive root Lie algebra and let $\phi$ be a root system automorphism of $\Delta(\kk)$. Then $\phi$ can be extended to a root system automorphism $\bar\phi$ of $\gG$. $\bar\phi$ induces a Lie algebra automorphism of $\gG$ that preserves $\kk$ and $C(\kk_{ss})$.  
%\end{lemma}
%\begin{proof}
%Extend the map $\phi$ to a linear map $\bar\phi$ by $\bar\phi(\gamma):=\gamma$ for all $\gamma\in \Delta(\kk)^\perp$. 

%Suppose we have already proved that $\bar\phi$ is a root system automorphism of $\Delta(\gG)$. By \cite[Corollary C, \S 2.9]{Samelson} $\bar\phi$ induces a Lie algebra automorphism of $\gG$ that preserves $\hh$. This automorphism maps the root spaces of $\kk$ onto themselves and the root spaces of $C(\kk_{ss})$ onto themselves, and the statement of the lemma follows.

%We are left to prove that $\bar\phi$ is indeed a root system automorphism. Since $\bar\phi$ is a linear map by definition, we only need to show that it maps the set $\Delta(\gG)$ onto itself.
%\end{proof}
 
%\begin{example}\label{exsp6} $\gG=\sP(6)$, $\Delta(\kk)=\{\pma2\eps_1,\pma2\eps_3\}$, $\Delta(\nn)=2\eps_2, \eps_2+ \eps_1, \eps_2- \eps_1 $. Then the cone condition is satisfied ($\hw_{ \bb\cap\kk }(\gG/\lL))$ = $\{\eps_3+ \eps_2,\eps_1+ \eps_3, \eps_3- \eps_2, -\eps_2+ \eps_1,-2\eps_2\}$), but there exists no parabolic $\pp$ with $\nn=\nn_\pp$. Indeed, assume the contrary. Then there exists a functional $t:\hh^*\to\RR$ with $t(\alpha)>0$ for all $\alpha \in \Delta (\nn)$, and $t(\alpha)\leq 0$ for all other roots $\alpha$. Therefore $t(\eps_1)=t(\eps_3)=0$, therefore $t(\eps_2)>0$, therefore $t(\eps_3+\eps_2)>0$, contradiction.
%\end{example}
%\begin{example}\label{exCase3.2.2.2}
%``Worst case scenario'' for $\gG=\sP(12)$. 
%\begin{eqnarray*}
%\Delta(\kk)&=&\pma\{\eps_1+ \eps_4,\eps_2-\eps_5,\eps_3-\eps_6\} \\
%\Delta(\nn)&=&\{\eps_1+\eps_3,\eps_2+\eps_4,\eps_2+\eps_3,2\eps_2,2\eps_3,\eps_2+\eps_6, \eps_1+\eps_6, \eps_5+\eps_6, 2\eps_5, 2\eps_6,\\
%&&~ \eps_3+\eps_5,\eps_4+\eps_5, \eps_6-\eps_4,\eps_5-\eps_1,\eps_3-\eps_4,\eps_2-\eps_1,\eps_3+\eps_6,\eps_2+\eps_5\}\\
%\hw_{\bb\cap\kk}(\gG/\lL)&=&\{\eps_1+\eps_2,\eps_3+\eps_4,2\eps_1,2\eps_4,\eps_4-\eps_5,\eps_4-\eps_6,\eps_1-\eps_5,\eps_1-\eps_6,\eps_2-\eps_6,\\
%&&~ \eps_3-\eps_5, -\eps_5-\eps_6,-2\eps_5,-2\eps_6\}
%\end{eqnarray*}
%\end{example}
%\begin{example}\label{extypeD}
%``Worst case scenario'' for $\gG=\so(14)$. 
%\begin{eqnarray*}
%\Delta(\kk)&=&\pma\{\eps_1- \eps_2,\eps_1-\eps_6,\eps_2-\eps_6,\eps_3-\eps_5,\eps_4-\eps_7\} \\
%\Delta(\nn)&=&\{\eps_1+\eps_3,\eps_1+\eps_5,\eps_2+\eps_3,\eps_2+\eps_5,\eps_6+\eps_3,\eps_5+ \eps_6,\\
%&&\eps_1+\eps_4,\eps_1+\eps_7,\eps_2+\eps_4,\eps_2+\eps_7,\eps_6+\eps_4,\eps_6+\eps_7,\eps_3+\eps_5,\eps_4+\eps_7\}\\
%\hw_{\bb\cap\kk}(\gG/\lL)&=&\{\eps_1+\eps_2,-\eps_6-\eps_2,\\
%&&\eps_1-\eps_5,\eps_3-\eps_6,-\eps_6-\eps_5, \\
%&&\eps_1-\eps_7,\eps_4-\eps_6,-\eps_6-\eps_7,\\
%&&-\eps_3-\eps_5,\\
%&&\eps_3+\eps_4,-\eps_5-\eps_7,\eps_3-\eps_7,\eps_4-\eps_5, \\
%&&-\eps_4-\eps_7 \}
%\end{eqnarray*}
%\end{example}

\nocite{Jantzen}
\nocite{VinbergOnishchik}
\nocite{Serre}
\nocite{Joseph:MinDimNilpotentOrbit}
\nocite{BorhoKraft}
\cleardoublepage
\addcontentsline{toc}{chapter}{Bibliography}
\bibliographystyle{hamsalpha}
\bibliography{./TodorMilevsBibliography}

\begin{appendix}

\arxivVersion{
\chapter{Tables for the exceptional Lie algebras}
\section{Note on table generation}
All tables in the appendix are generated by computer. The tables are also available in .html format from the author or from the web server of the ``vector partition'' program.
%\newpage
%\addtolength{\hoffset}{-4.5cm}
%\addtolength{\textwidth}{8cm}
%\addtolength{\voffset}{-3.3cm}
%\addtolength{\textheight}{6.3cm}
%\section{Non-solvable non-reductive root subalgebras of $\sP(6)$ and $\so(7)$ failing the cone condition}\label{secSubalgebrasB3C3}

%\subsection{$\so(7)$}\label{secSubalgebrasB3}


\section{Reductive root subalgebras of the exceptional Lie algebras}\label{secRootSATables}
The reductive root subalgebras of the exceptional Lie algebras are described and tabulated in \cite{Dynkin:Semisimple}. Our tables list in addition the type of $C(\kk_{ss})_{ss}$ and the inclusions between the root subsystems parametrizing the reductive root subalgebras.
\subsection{$F_4$}
All diagrams that consist of short roots are labeled by $'$. For example, $A_2'$ has 6 short roots; in the notation $C_3+A_1$, the root system $A_1$ has long roots.

\begin{longtable}{cccc}\multicolumn{4}{c}{$F_4$}\endhead
 \begin{tabular}{p{2cm}}
  $\mathfrak{k}_{ss}$: $B_4$\\
 $C(\mathfrak{k_{ss}})_{ss}$: -\\
 
 $\mathfrak{k}_{ss}$ lies in: \\
 $F_4$, \end{tabular} & \begin{tabular}{p{2cm}}
  $\mathfrak{k}_{ss}$: $D_4$\\
 $C(\mathfrak{k_{ss}})_{ss}$: -\\
 
 $\mathfrak{k}_{ss}$ lies in: \\
 $B_4$, \end{tabular} & \begin{tabular}{p{2cm}}
  $\mathfrak{k}_{ss}$: $C_3$+$A_1$\\
 $C(\mathfrak{k_{ss}})_{ss}$: -\\
 
 $\mathfrak{k}_{ss}$ lies in: \\
 $F_4$, \end{tabular} & \begin{tabular}{p{2cm}}
  $\mathfrak{k}_{ss}$: $A_3$+$A_1'$\\
 $C(\mathfrak{k_{ss}})_{ss}$: -\\
 
 $\mathfrak{k}_{ss}$ lies in: \\
 $F_4$, $B_4$, \end{tabular}\\
 \hline\begin{tabular}{p{2cm}}
  $\mathfrak{k}_{ss}$: $B_2\simeq C_2$+2$A_1$\\
 $C(\mathfrak{k_{ss}})_{ss}$: -\\
 
 $\mathfrak{k}_{ss}$ lies in: \\
 $C_3$+$A_1$, $B_4$, \end{tabular} & \begin{tabular}{p{2cm}}
  $\mathfrak{k}_{ss}$: $A_2'$+$A_2$\\
 $C(\mathfrak{k_{ss}})_{ss}$: -\\
 
 $\mathfrak{k}_{ss}$ lies in: \\
 $F_4$, \end{tabular} & \begin{tabular}{p{2cm}}
  $\mathfrak{k}_{ss}$: 4$A_1$\\
 $C(\mathfrak{k_{ss}})_{ss}$: -\\
 
 $\mathfrak{k}_{ss}$ lies in: \\
 $B_2\simeq C_2$+2$A_1$, \\
 $D_4$, \end{tabular} & \begin{tabular}{p{2cm}}
  $\mathfrak{k}_{ss}$: $C_3$\\
 $C(\mathfrak{k_{ss}})_{ss}$: $A_1$\\
 
 $\mathfrak{k}_{ss}$ lies in: \\
 $C_3$+$A_1$, $F_4$, \end{tabular}\\
 \hline\begin{tabular}{p{2cm}}
  $\mathfrak{k}_{ss}$: $B_3$\\
 $C(\mathfrak{k_{ss}})_{ss}$: -\\
 
 $\mathfrak{k}_{ss}$ lies in: \\
 $F_4$, $B_4$, \end{tabular} & \begin{tabular}{p{2cm}}
  $\mathfrak{k}_{ss}$: $A_3$\\
 $C(\mathfrak{k_{ss}})_{ss}$: $A_1'$\\
 
 $\mathfrak{k}_{ss}$ lies in: \\
 $A_3$+$A_1'$, $B_4$, $D_4$, \\
 $B_3$, \end{tabular} & \begin{tabular}{p{2cm}}
  $\mathfrak{k}_{ss}$: $B_2\simeq C_2$+$A_1$\\
 $C(\mathfrak{k_{ss}})_{ss}$: $A_1$\\
 
 $\mathfrak{k}_{ss}$ lies in: \\
 $B_2\simeq C_2$+2$A_1$, \\
 $C_3$+$A_1$, $C_3$, $B_4$, \\
 \end{tabular} & \begin{tabular}{p{2cm}}
  $\mathfrak{k}_{ss}$: $A_2'$+$A_1$\\
 $C(\mathfrak{k_{ss}})_{ss}$: -\\
 
 $\mathfrak{k}_{ss}$ lies in: \\
 $F_4$, $C_3$+$A_1$, $A_2'$+$A_2$, \\
 \end{tabular}\\
 \hline\begin{tabular}{p{2cm}}
  $\mathfrak{k}_{ss}$: $A_2$+$A_1'$\\
 $C(\mathfrak{k_{ss}})_{ss}$: -\\
 
 $\mathfrak{k}_{ss}$ lies in: \\
 $F_4$, $A_3$+$A_1'$, $B_4$, \\
 $A_2'$+$A_2$, \end{tabular} & \begin{tabular}{p{2cm}}
  $\mathfrak{k}_{ss}$: $A_1'$+2$A_1$\\
 $C(\mathfrak{k_{ss}})_{ss}$: -\\
 
 $\mathfrak{k}_{ss}$ lies in: \\
 $C_3$+$A_1$, $A_3$+$A_1'$, \\
 $B_2\simeq C_2$+2$A_1$, \\
 $B_3$, \end{tabular} & \begin{tabular}{p{2cm}}
  $\mathfrak{k}_{ss}$: 3$A_1$\\
 $C(\mathfrak{k_{ss}})_{ss}$: $A_1$\\
 
 $\mathfrak{k}_{ss}$ lies in: \\
 4$A_1$, $D_4$, $B_2\simeq C_2$+$A_1$, \\
 $B_2\simeq C_2$+2$A_1$, \\
 \end{tabular} & \begin{tabular}{p{2cm}}
  $\mathfrak{k}_{ss}$: $B_2\simeq C_2$\\
 $C(\mathfrak{k_{ss}})_{ss}$: 2$A_1$\\
 
 $\mathfrak{k}_{ss}$ lies in: \\
 $B_2\simeq C_2$+$A_1$, \\
 $B_3$, $C_3$, \end{tabular}\\
 \hline\begin{tabular}{p{2cm}}
  $\mathfrak{k}_{ss}$: $A_2'$\\
 $C(\mathfrak{k_{ss}})_{ss}$: $A_2$\\
 
 $\mathfrak{k}_{ss}$ lies in: \\
 $A_2'$+$A_1$, $C_3$, \end{tabular} & \begin{tabular}{p{2cm}}
  $\mathfrak{k}_{ss}$: $A_2$\\
 $C(\mathfrak{k_{ss}})_{ss}$: $A_2'$\\
 
 $\mathfrak{k}_{ss}$ lies in: \\
 $A_2$+$A_1'$, $A_3$, $B_3$, \\
 \end{tabular} & \begin{tabular}{p{2cm}}
  $\mathfrak{k}_{ss}$: $A_1'$+$A_1$\\
 $C(\mathfrak{k_{ss}})_{ss}$: $A_1$\\
 
 $\mathfrak{k}_{ss}$ lies in: \\
 $A_1'$+2$A_1$, $B_2\simeq C_2$+$A_1$, \\
 $A_2'$+$A_1$, $B_3$, $A_2$+$A_1'$, \\
 $C_3$, \end{tabular} & \begin{tabular}{p{2cm}}
  $\mathfrak{k}_{ss}$: 2$A_1$\\
 $C(\mathfrak{k_{ss}})_{ss}$: $B_2\simeq C_2$\\
 
 $\mathfrak{k}_{ss}$ lies in: \\
 $A_1'$+2$A_1$, 3$A_1$, $B_2\simeq C_2$+$A_1$, \\
 $A_3$, $B_2\simeq C_2$, \end{tabular}\\
 \hline\begin{tabular}{p{2cm}}
  $\mathfrak{k}_{ss}$: $A_1'$\\
 $C(\mathfrak{k_{ss}})_{ss}$: $A_3$\\
 
 $\mathfrak{k}_{ss}$ lies in: \\
 $A_1'$+$A_1$, $A_2'$, $B_2\simeq C_2$, \\
 \end{tabular} & \begin{tabular}{p{2cm}}
  $\mathfrak{k}_{ss}$: $A_1$\\
 $C(\mathfrak{k_{ss}})_{ss}$: $C_3$\\
 
 $\mathfrak{k}_{ss}$ lies in: \\
 $A_1'$+$A_1$, 2$A_1$, $A_2$, \\
 $B_2\simeq C_2$, \end{tabular} & \begin{tabular}{p{2cm}}
  $\mathfrak{k}_{ss}$: -\\
 $C(\mathfrak{k_{ss}})_{ss}$: $F_4$\\
 
 $\mathfrak{k}_{ss}$ lies in: \\
 $A_1'$, $A_1$, \end{tabular} & \end{longtable}



\subsection{$E_6$}



\begin{longtable}{cccc}\multicolumn{4}{c}{$E_6$}\endhead
\begin{tabular}{p{2cm}}
 $\mathfrak{k}_{ss}$: $A_5$+$A_1$\\
$C(\mathfrak{k_{ss}})_{ss}$: -\\

$\mathfrak{k}_{ss}$ lies in: \\
$E_6$, \end{tabular} & \begin{tabular}{p{2cm}}
 $\mathfrak{k}_{ss}$: 3$A_2$\\
$C(\mathfrak{k_{ss}})_{ss}$: -\\

$\mathfrak{k}_{ss}$ lies in: \\
$E_6$, \end{tabular} & \begin{tabular}{p{2cm}}
 $\mathfrak{k}_{ss}$: $D_5$\\
$C(\mathfrak{k_{ss}})_{ss}$: -\\

$\mathfrak{k}_{ss}$ lies in: \\
$E_6$, \end{tabular} & \begin{tabular}{p{2cm}}
 $\mathfrak{k}_{ss}$: $A_5$\\
$C(\mathfrak{k_{ss}})_{ss}$: $A_1$\\

$\mathfrak{k}_{ss}$ lies in: \\
$A_5$+$A_1$, $E_6$, \end{tabular}\\
\hline\begin{tabular}{p{2cm}}
 $\mathfrak{k}_{ss}$: $A_4$+$A_1$\\
$C(\mathfrak{k_{ss}})_{ss}$: -\\

$\mathfrak{k}_{ss}$ lies in: \\
$E_6$, $A_5$+$A_1$, \end{tabular} & \begin{tabular}{p{2cm}}
 $\mathfrak{k}_{ss}$: $A_3$+2$A_1$\\
$C(\mathfrak{k_{ss}})_{ss}$: -\\

$\mathfrak{k}_{ss}$ lies in: \\
$D_5$, $A_5$+$A_1$, \end{tabular} & \begin{tabular}{p{2cm}}
 $\mathfrak{k}_{ss}$: 2$A_2$+$A_1$\\
$C(\mathfrak{k_{ss}})_{ss}$: -\\

$\mathfrak{k}_{ss}$ lies in: \\
$E_6$, $A_5$+$A_1$, 3$A_2$, \\
\end{tabular} & \begin{tabular}{p{2cm}}
 $\mathfrak{k}_{ss}$: $D_4$\\
$C(\mathfrak{k_{ss}})_{ss}$: -\\

$\mathfrak{k}_{ss}$ lies in: \\
$D_5$, \end{tabular}\\
\hline\begin{tabular}{p{2cm}}
 $\mathfrak{k}_{ss}$: $A_4$\\
$C(\mathfrak{k_{ss}})_{ss}$: $A_1$\\

$\mathfrak{k}_{ss}$ lies in: \\
$A_4$+$A_1$, $D_5$, $A_5$, \\
\end{tabular} & \begin{tabular}{p{2cm}}
 $\mathfrak{k}_{ss}$: $A_3$+$A_1$\\
$C(\mathfrak{k_{ss}})_{ss}$: $A_1$\\

$\mathfrak{k}_{ss}$ lies in: \\
$A_3$+2$A_1$, $D_5$, $A_4$+$A_1$, \\
$A_5$, \end{tabular} & \begin{tabular}{p{2cm}}
 $\mathfrak{k}_{ss}$: 2$A_2$\\
$C(\mathfrak{k_{ss}})_{ss}$: $A_2$\\

$\mathfrak{k}_{ss}$ lies in: \\
2$A_2$+$A_1$, $A_5$, \end{tabular} & \begin{tabular}{p{2cm}}
 $\mathfrak{k}_{ss}$: $A_2$+2$A_1$\\
$C(\mathfrak{k_{ss}})_{ss}$: -\\

$\mathfrak{k}_{ss}$ lies in: \\
2$A_2$+$A_1$, $A_4$+$A_1$, \\
$D_5$, $A_3$+2$A_1$, \end{tabular}\\
\hline\begin{tabular}{p{2cm}}
 $\mathfrak{k}_{ss}$: 4$A_1$\\
$C(\mathfrak{k_{ss}})_{ss}$: -\\

$\mathfrak{k}_{ss}$ lies in: \\
$A_3$+2$A_1$, $D_4$, \end{tabular} & \begin{tabular}{p{2cm}}
 $\mathfrak{k}_{ss}$: $A_3$\\
$C(\mathfrak{k_{ss}})_{ss}$: 2$A_1$\\

$\mathfrak{k}_{ss}$ lies in: \\
$A_3$+$A_1$, $A_4$, $D_4$, \\
\end{tabular} & \begin{tabular}{p{2cm}}
 $\mathfrak{k}_{ss}$: $A_2$+$A_1$\\
$C(\mathfrak{k_{ss}})_{ss}$: $A_2$\\

$\mathfrak{k}_{ss}$ lies in: \\
$A_2$+2$A_1$, $A_3$+$A_1$, \\
2$A_2$, $A_4$, \end{tabular} & \begin{tabular}{p{2cm}}
 $\mathfrak{k}_{ss}$: 3$A_1$\\
$C(\mathfrak{k_{ss}})_{ss}$: $A_1$\\

$\mathfrak{k}_{ss}$ lies in: \\
4$A_1$, $A_2$+2$A_1$, $D_4$, \\
$A_3$+$A_1$, \end{tabular}\\
\hline\begin{tabular}{p{2cm}}
 $\mathfrak{k}_{ss}$: $A_2$\\
$C(\mathfrak{k_{ss}})_{ss}$: 2$A_2$\\

$\mathfrak{k}_{ss}$ lies in: \\
$A_2$+$A_1$, $A_3$, \end{tabular} & \begin{tabular}{p{2cm}}
 $\mathfrak{k}_{ss}$: 2$A_1$\\
$C(\mathfrak{k_{ss}})_{ss}$: $A_3$\\

$\mathfrak{k}_{ss}$ lies in: \\
3$A_1$, $A_2$+$A_1$, $A_3$, \\
\end{tabular} & \begin{tabular}{p{2cm}}
 $\mathfrak{k}_{ss}$: $A_1$\\
$C(\mathfrak{k_{ss}})_{ss}$: $A_5$\\

$\mathfrak{k}_{ss}$ lies in: \\
2$A_1$, $A_2$, \end{tabular} & \begin{tabular}{p{2cm}}
 $\mathfrak{k}_{ss}$: -\\
$C(\mathfrak{k_{ss}})_{ss}$: $E_6$\\

$\mathfrak{k}_{ss}$ lies in: \\
$A_1$, \end{tabular}\\
\hline\end{longtable}







\subsection{$E_7$}
 
 
\begin{longtable}{cccc}\multicolumn{4}{c}{$E_7$}\endhead
  \begin{tabular}{p{2cm}}
   $\mathfrak{k}_{ss}$: $A_7$\\
  $C(\mathfrak{k_{ss}})_{ss}$: -\\
  
  $\mathfrak{k}_{ss}$ lies in: \\
  $E_7$, \end{tabular} & \begin{tabular}{p{2cm}}
   $\mathfrak{k}_{ss}$: $D_6$+$A_1$\\
  $C(\mathfrak{k_{ss}})_{ss}$: -\\
  
  $\mathfrak{k}_{ss}$ lies in: \\
  $E_7$, \end{tabular} & \begin{tabular}{p{2cm}}
   $\mathfrak{k}_{ss}$: $A_5$+$A_2$\\
  $C(\mathfrak{k_{ss}})_{ss}$: -\\
  
  $\mathfrak{k}_{ss}$ lies in: \\
  $E_7$, \end{tabular} & \begin{tabular}{p{2cm}}
   $\mathfrak{k}_{ss}$: $D_4$+3$A_1$\\
  $C(\mathfrak{k_{ss}})_{ss}$: -\\
  
  $\mathfrak{k}_{ss}$ lies in: \\
  $D_6$+$A_1$, \end{tabular}\\
  \hline\begin{tabular}{p{2cm}}
   $\mathfrak{k}_{ss}$: 2$A_3$+$A_1$\\
  $C(\mathfrak{k_{ss}})_{ss}$: -\\
  
  $\mathfrak{k}_{ss}$ lies in: \\
  $D_6$+$A_1$, $E_7$, \end{tabular} & \begin{tabular}{p{2cm}}
   $\mathfrak{k}_{ss}$: 7$A_1$\\
  $C(\mathfrak{k_{ss}})_{ss}$: -\\
  
  $\mathfrak{k}_{ss}$ lies in: \\
  $D_4$+3$A_1$, \end{tabular} & \begin{tabular}{p{2cm}}
   $\mathfrak{k}_{ss}$: $E_6$\\
  $C(\mathfrak{k_{ss}})_{ss}$: -\\
  
  $\mathfrak{k}_{ss}$ lies in: \\
  $E_7$, \end{tabular} & \begin{tabular}{p{2cm}}
   $\mathfrak{k}_{ss}$: $D_6$\\
  $C(\mathfrak{k_{ss}})_{ss}$: $A_1$\\
  
  $\mathfrak{k}_{ss}$ lies in: \\
  $D_6$+$A_1$, $E_7$, \end{tabular}\\
  \hline\begin{tabular}{p{2cm}}
   $\mathfrak{k}_{ss}$: $A_6$\\
  $C(\mathfrak{k_{ss}})_{ss}$: -\\
  
  $\mathfrak{k}_{ss}$ lies in: \\
  $E_7$, $A_7$, \end{tabular} & \begin{tabular}{p{2cm}}
   $\mathfrak{k}_{ss}$: $D_5$+$A_1$\\
  $C(\mathfrak{k_{ss}})_{ss}$: -\\
  
  $\mathfrak{k}_{ss}$ lies in: \\
  $E_7$, $D_6$+$A_1$, \end{tabular} & \begin{tabular}{p{2cm}}
   $\mathfrak{k}_{ss}$: $A_5$+$A_1$\\
  $C(\mathfrak{k_{ss}})_{ss}$: -\\
  
  $\mathfrak{k}_{ss}$ lies in: \\
  $E_7$, $D_6$+$A_1$, $A_5$+$A_2$, \\
  \end{tabular} & \begin{tabular}{p{2cm}}
   $\mathfrak{k}_{ss}$: $D_4$+2$A_1$\\
  $C(\mathfrak{k_{ss}})_{ss}$: $A_1$\\
  
  $\mathfrak{k}_{ss}$ lies in: \\
  $D_4$+3$A_1$, $D_6$, $D_6$+$A_1$, \\
  \end{tabular}\\
  \hline\begin{tabular}{p{2cm}}
   $\mathfrak{k}_{ss}$: $A_4$+$A_2$\\
  $C(\mathfrak{k_{ss}})_{ss}$: -\\
  
  $\mathfrak{k}_{ss}$ lies in: \\
  $E_7$, $A_7$, $A_5$+$A_2$, \\
  \end{tabular} & \begin{tabular}{p{2cm}}
   $\mathfrak{k}_{ss}$: 2$A_3$\\
  $C(\mathfrak{k_{ss}})_{ss}$: $A_1$\\
  
  $\mathfrak{k}_{ss}$ lies in: \\
  2$A_3$+$A_1$, $D_6$, $A_7$, \\
  \end{tabular} & \begin{tabular}{p{2cm}}
   $\mathfrak{k}_{ss}$: $A_3$+$A_2$+$A_1$\\
  $C(\mathfrak{k_{ss}})_{ss}$: -\\
  
  $\mathfrak{k}_{ss}$ lies in: \\
  $D_6$+$A_1$, $E_7$, $A_5$+$A_2$, \\
  2$A_3$+$A_1$, \end{tabular} & \begin{tabular}{p{2cm}}
   $\mathfrak{k}_{ss}$: $A_3$+3$A_1$\\
  $C(\mathfrak{k_{ss}})_{ss}$: -\\
  
  $\mathfrak{k}_{ss}$ lies in: \\
  $D_5$+$A_1$, $D_6$+$A_1$, \\
  $D_4$+3$A_1$, 2$A_3$+$A_1$, \\
  \end{tabular}\\
  \hline\begin{tabular}{p{2cm}}
   $\mathfrak{k}_{ss}$: 3$A_2$\\
  $C(\mathfrak{k_{ss}})_{ss}$: -\\
  
  $\mathfrak{k}_{ss}$ lies in: \\
  $A_5$+$A_2$, $E_6$, \end{tabular} & \begin{tabular}{p{2cm}}
   $\mathfrak{k}_{ss}$: 6$A_1$\\
  $C(\mathfrak{k_{ss}})_{ss}$: $A_1$\\
  
  $\mathfrak{k}_{ss}$ lies in: \\
  7$A_1$, $D_4$+2$A_1$, $D_4$+3$A_1$, \\
  \end{tabular} & \begin{tabular}{p{2cm}}
   $\mathfrak{k}_{ss}$: $D_5$\\
  $C(\mathfrak{k_{ss}})_{ss}$: $A_1$\\
  
  $\mathfrak{k}_{ss}$ lies in: \\
  $D_5$+$A_1$, $E_6$, $D_6$, \\
  \end{tabular} & \begin{tabular}{p{2cm}}
   $\mathfrak{k}_{ss}$: $A_5$\\
  $C(\mathfrak{k_{ss}})_{ss}$: $A_1$\\
  
  $\mathfrak{k}_{ss}$ lies in: \\
  $A_5$+$A_1$, $A_6$, $E_6$, \\
  $D_6$, \end{tabular}\\
  \hline\begin{tabular}{p{2cm}}
   $\mathfrak{k}_{ss}$: $A_5$\\
  $C(\mathfrak{k_{ss}})_{ss}$: $A_2$\\
  
  $\mathfrak{k}_{ss}$ lies in: \\
  $A_5$+$A_1$, $D_6$, \end{tabular} & \begin{tabular}{p{2cm}}
   $\mathfrak{k}_{ss}$: $D_4$+$A_1$\\
  $C(\mathfrak{k_{ss}})_{ss}$: 2$A_1$\\
  
  $\mathfrak{k}_{ss}$ lies in: \\
  $D_4$+2$A_1$, $D_5$+$A_1$, \\
  $D_6$, \end{tabular} & \begin{tabular}{p{2cm}}
   $\mathfrak{k}_{ss}$: $A_4$+$A_1$\\
  $C(\mathfrak{k_{ss}})_{ss}$: -\\
  
  $\mathfrak{k}_{ss}$ lies in: \\
  $E_6$, $A_5$+$A_1$, $A_6$, \\
  $D_5$+$A_1$, $A_4$+$A_2$, \\
  \end{tabular} & \begin{tabular}{p{2cm}}
   $\mathfrak{k}_{ss}$: $A_3$+$A_2$\\
  $C(\mathfrak{k_{ss}})_{ss}$: $A_1$\\
  
  $\mathfrak{k}_{ss}$ lies in: \\
  $A_3$+$A_2$+$A_1$, $A_4$+$A_2$, \\
  $D_6$, $A_6$, 2$A_3$, \end{tabular}\\
  \hline\begin{tabular}{p{2cm}}
   $\mathfrak{k}_{ss}$: $A_3$+2$A_1$\\
  $C(\mathfrak{k_{ss}})_{ss}$: $A_1$\\
  
  $\mathfrak{k}_{ss}$ lies in: \\
  $A_3$+3$A_1$, $D_4$+2$A_1$, \\
  $D_6$, $A_5$+$A_1$, $A_3$+$A_2$+$A_1$, \\
  $D_5$+$A_1$, \end{tabular} & \begin{tabular}{p{2cm}}
   $\mathfrak{k}_{ss}$: 2$A_2$+$A_1$\\
  $C(\mathfrak{k_{ss}})_{ss}$: -\\
  
  $\mathfrak{k}_{ss}$ lies in: \\
  $A_4$+$A_2$, $A_5$+$A_1$, \\
  $E_6$, $A_3$+$A_2$+$A_1$, \\
  3$A_2$, \end{tabular} & \begin{tabular}{p{2cm}}
   $\mathfrak{k}_{ss}$: $A_2$+3$A_1$\\
  $C(\mathfrak{k_{ss}})_{ss}$: -\\
  
  $\mathfrak{k}_{ss}$ lies in: \\
  $A_3$+3$A_1$, $D_5$+$A_1$, \\
  $A_3$+$A_2$+$A_1$, \end{tabular} & \begin{tabular}{p{2cm}}
   $\mathfrak{k}_{ss}$: 5$A_1$\\
  $C(\mathfrak{k_{ss}})_{ss}$: 2$A_1$\\
  
  $\mathfrak{k}_{ss}$ lies in: \\
  6$A_1$, $D_4$+2$A_1$, $A_3$+3$A_1$, \\
  $D_4$+$A_1$, \end{tabular}\\
  \hline\begin{tabular}{p{2cm}}
   $\mathfrak{k}_{ss}$: $D_4$\\
  $C(\mathfrak{k_{ss}})_{ss}$: 3$A_1$\\
  
  $\mathfrak{k}_{ss}$ lies in: \\
  $D_4$+$A_1$, $D_5$, \end{tabular} & \begin{tabular}{p{2cm}}
   $\mathfrak{k}_{ss}$: $A_4$\\
  $C(\mathfrak{k_{ss}})_{ss}$: $A_2$\\
  
  $\mathfrak{k}_{ss}$ lies in: \\
  $A_4$+$A_1$, $A_5$, $D_5$, \\
  $A_5$, \end{tabular} & \begin{tabular}{p{2cm}}
   $\mathfrak{k}_{ss}$: $A_3$+$A_1$\\
  $C(\mathfrak{k_{ss}})_{ss}$: 2$A_1$\\
  
  $\mathfrak{k}_{ss}$ lies in: \\
  $A_3$+2$A_1$, $D_5$, $A_4$+$A_1$, \\
  $D_4$+$A_1$, $A_5$, $A_3$+$A_2$, \\
  \end{tabular} & \begin{tabular}{p{2cm}}
   $\mathfrak{k}_{ss}$: $A_3$+$A_1$\\
  $C(\mathfrak{k_{ss}})_{ss}$: $A_3$\\
  
  $\mathfrak{k}_{ss}$ lies in: \\
  $A_3$+2$A_1$, $A_5$, $D_4$+$A_1$, \\
  \end{tabular}\\
  \hline\begin{tabular}{p{2cm}}
   $\mathfrak{k}_{ss}$: 2$A_2$\\
  $C(\mathfrak{k_{ss}})_{ss}$: $A_2$\\
  
  $\mathfrak{k}_{ss}$ lies in: \\
  2$A_2$+$A_1$, $A_3$+$A_2$, \\
  $A_5$, $A_5$, \end{tabular} & \begin{tabular}{p{2cm}}
   $\mathfrak{k}_{ss}$: $A_2$+2$A_1$\\
  $C(\mathfrak{k_{ss}})_{ss}$: $A_1$\\
  
  $\mathfrak{k}_{ss}$ lies in: \\
  $A_2$+3$A_1$, 2$A_2$+$A_1$, \\
  $A_3$+2$A_1$, $D_5$, $A_4$+$A_1$, \\
  $A_3$+$A_2$, \end{tabular} & \begin{tabular}{p{2cm}}
   $\mathfrak{k}_{ss}$: 4$A_1$\\
  $C(\mathfrak{k_{ss}})_{ss}$: 3$A_1$\\
  
  $\mathfrak{k}_{ss}$ lies in: \\
  5$A_1$, $A_2$+3$A_1$, $D_4$+$A_1$, \\
  $A_3$+2$A_1$, \end{tabular} & \begin{tabular}{p{2cm}}
   $\mathfrak{k}_{ss}$: $A_3$\\
  $C(\mathfrak{k_{ss}})_{ss}$: $A_3$+$A_1$\\
  
  $\mathfrak{k}_{ss}$ lies in: \\
  $A_3$+$A_1$, $A_3$+$A_1$, \\
  $A_4$, $D_4$, \end{tabular}\\
  \hline\begin{tabular}{p{2cm}}
   $\mathfrak{k}_{ss}$: $A_2$+$A_1$\\
  $C(\mathfrak{k_{ss}})_{ss}$: $A_3$\\
  
  $\mathfrak{k}_{ss}$ lies in: \\
  $A_2$+2$A_1$, $A_4$, $A_3$+$A_1$, \\
  2$A_2$, $A_3$+$A_1$, \end{tabular} & \begin{tabular}{p{2cm}}
   $\mathfrak{k}_{ss}$: 3$A_1$\\
  $C(\mathfrak{k_{ss}})_{ss}$: 4$A_1$\\
  
  $\mathfrak{k}_{ss}$ lies in: \\
  4$A_1$, $A_2$+2$A_1$, $A_3$+$A_1$, \\
  $D_4$, \end{tabular} & \begin{tabular}{p{2cm}}
   $\mathfrak{k}_{ss}$: 3$A_1$\\
  $C(\mathfrak{k_{ss}})_{ss}$: $D_4$\\
  
  $\mathfrak{k}_{ss}$ lies in: \\
  4$A_1$, $A_3$+$A_1$, \end{tabular} & \begin{tabular}{p{2cm}}
   $\mathfrak{k}_{ss}$: $A_2$\\
  $C(\mathfrak{k_{ss}})_{ss}$: $A_5$\\
  
  $\mathfrak{k}_{ss}$ lies in: \\
  $A_2$+$A_1$, $A_3$, \end{tabular}\\
  \hline\begin{tabular}{p{2cm}}
   $\mathfrak{k}_{ss}$: 2$A_1$\\
  $C(\mathfrak{k_{ss}})_{ss}$: $D_4$+$A_1$\\
  
  $\mathfrak{k}_{ss}$ lies in: \\
  3$A_1$, 3$A_1$, $A_2$+$A_1$, \\
  $A_3$, \end{tabular} & \begin{tabular}{p{2cm}}
   $\mathfrak{k}_{ss}$: $A_1$\\
  $C(\mathfrak{k_{ss}})_{ss}$: $D_6$\\
  
  $\mathfrak{k}_{ss}$ lies in: \\
  2$A_1$, $A_2$, \end{tabular} & \begin{tabular}{p{2cm}}
   $\mathfrak{k}_{ss}$: -\\
  $C(\mathfrak{k_{ss}})_{ss}$: $E_7$\\
  
  $\mathfrak{k}_{ss}$ lies in: \\
  $A_1$, \end{tabular} & \end{longtable}

\subsection{$E_8$}

\begin{longtable}{cccc}\multicolumn{4}{c}{$E_8$}\endhead
\begin{tabular}{p{2cm}}
 $\mathfrak{k}_{ss}$: $D_8$\\
$C(\mathfrak{k_{ss}})_{ss}$: -\\

$\mathfrak{k}_{ss}$ lies in: \\
$E_8$, \end{tabular} & \begin{tabular}{p{2cm}}
 $\mathfrak{k}_{ss}$: $A_8$\\
$C(\mathfrak{k_{ss}})_{ss}$: -\\

$\mathfrak{k}_{ss}$ lies in: \\
$E_8$, \end{tabular} & \begin{tabular}{p{2cm}}
 $\mathfrak{k}_{ss}$: $E_7$+$A_1$\\
$C(\mathfrak{k_{ss}})_{ss}$: -\\

$\mathfrak{k}_{ss}$ lies in: \\
$E_8$, \end{tabular} & \begin{tabular}{p{2cm}}
 $\mathfrak{k}_{ss}$: $A_7$+$A_1$\\
$C(\mathfrak{k_{ss}})_{ss}$: -\\

$\mathfrak{k}_{ss}$ lies in: \\
$E_7$+$A_1$, $E_8$, \end{tabular}\\
\hline\begin{tabular}{p{2cm}}
 $\mathfrak{k}_{ss}$: $E_6$+$A_2$\\
$C(\mathfrak{k_{ss}})_{ss}$: -\\

$\mathfrak{k}_{ss}$ lies in: \\
$E_8$, \end{tabular} & \begin{tabular}{p{2cm}}
 $\mathfrak{k}_{ss}$: $D_6$+2$A_1$\\
$C(\mathfrak{k_{ss}})_{ss}$: -\\

$\mathfrak{k}_{ss}$ lies in: \\
$E_7$+$A_1$, $D_8$, \end{tabular} & \begin{tabular}{p{2cm}}
 $\mathfrak{k}_{ss}$: $D_5$+$A_3$\\
$C(\mathfrak{k_{ss}})_{ss}$: -\\

$\mathfrak{k}_{ss}$ lies in: \\
$E_8$, $D_8$, \end{tabular} & \begin{tabular}{p{2cm}}
 $\mathfrak{k}_{ss}$: $A_5$+$A_2$+$A_1$\\
$C(\mathfrak{k_{ss}})_{ss}$: -\\

$\mathfrak{k}_{ss}$ lies in: \\
$E_6$+$A_2$, $E_7$+$A_1$, \\
$E_8$, \end{tabular}\\
\hline\begin{tabular}{p{2cm}}
 $\mathfrak{k}_{ss}$: 2$D_4$\\
$C(\mathfrak{k_{ss}})_{ss}$: -\\

$\mathfrak{k}_{ss}$ lies in: \\
$D_8$, \end{tabular} & \begin{tabular}{p{2cm}}
 $\mathfrak{k}_{ss}$: $D_4$+4$A_1$\\
$C(\mathfrak{k_{ss}})_{ss}$: -\\

$\mathfrak{k}_{ss}$ lies in: \\
$D_6$+2$A_1$, 2$D_4$, \end{tabular} & \begin{tabular}{p{2cm}}
 $\mathfrak{k}_{ss}$: 2$A_4$\\
$C(\mathfrak{k_{ss}})_{ss}$: -\\

$\mathfrak{k}_{ss}$ lies in: \\
$E_8$, \end{tabular} & \begin{tabular}{p{2cm}}
 $\mathfrak{k}_{ss}$: 2$A_3$+2$A_1$\\
$C(\mathfrak{k_{ss}})_{ss}$: -\\

$\mathfrak{k}_{ss}$ lies in: \\
$E_7$+$A_1$, $D_6$+2$A_1$, \\
$D_5$+$A_3$, \end{tabular}\\
\hline\begin{tabular}{p{2cm}}
 $\mathfrak{k}_{ss}$: 4$A_2$\\
$C(\mathfrak{k_{ss}})_{ss}$: -\\

$\mathfrak{k}_{ss}$ lies in: \\
$E_6$+$A_2$, \end{tabular} & \begin{tabular}{p{2cm}}
 $\mathfrak{k}_{ss}$: 8$A_1$\\
$C(\mathfrak{k_{ss}})_{ss}$: -\\

$\mathfrak{k}_{ss}$ lies in: \\
$D_4$+4$A_1$, \end{tabular} & \begin{tabular}{p{2cm}}
 $\mathfrak{k}_{ss}$: $E_7$\\
$C(\mathfrak{k_{ss}})_{ss}$: $A_1$\\

$\mathfrak{k}_{ss}$ lies in: \\
$E_7$+$A_1$, $E_8$, \end{tabular} & \begin{tabular}{p{2cm}}
 $\mathfrak{k}_{ss}$: $D_7$\\
$C(\mathfrak{k_{ss}})_{ss}$: -\\

$\mathfrak{k}_{ss}$ lies in: \\
$E_8$, $D_8$, \end{tabular}\\
\hline\begin{tabular}{p{2cm}}
 $\mathfrak{k}_{ss}$: $A_7$\\
$C(\mathfrak{k_{ss}})_{ss}$: -\\

$\mathfrak{k}_{ss}$ lies in: \\
$E_8$, $D_8$, $A_8$, \end{tabular} & \begin{tabular}{p{2cm}}
 $\mathfrak{k}_{ss}$: $A_7$\\
$C(\mathfrak{k_{ss}})_{ss}$: $A_1$\\

$\mathfrak{k}_{ss}$ lies in: \\
$A_7$+$A_1$, $D_8$, $E_7$, \\
\end{tabular} & \begin{tabular}{p{2cm}}
 $\mathfrak{k}_{ss}$: $E_6$+$A_1$\\
$C(\mathfrak{k_{ss}})_{ss}$: -\\

$\mathfrak{k}_{ss}$ lies in: \\
$E_7$+$A_1$, $E_8$, $E_6$+$A_2$, \\
\end{tabular} & \begin{tabular}{p{2cm}}
 $\mathfrak{k}_{ss}$: $D_6$+$A_1$\\
$C(\mathfrak{k_{ss}})_{ss}$: $A_1$\\

$\mathfrak{k}_{ss}$ lies in: \\
$D_6$+2$A_1$, $E_7$+$A_1$, \\
$E_7$, $D_8$, \end{tabular}\\
\hline\begin{tabular}{p{2cm}}
 $\mathfrak{k}_{ss}$: $A_6$+$A_1$\\
$C(\mathfrak{k_{ss}})_{ss}$: -\\

$\mathfrak{k}_{ss}$ lies in: \\
$E_7$+$A_1$, $E_8$, $A_8$, \\
$A_7$+$A_1$, \end{tabular} & \begin{tabular}{p{2cm}}
 $\mathfrak{k}_{ss}$: $D_5$+$A_2$\\
$C(\mathfrak{k_{ss}})_{ss}$: -\\

$\mathfrak{k}_{ss}$ lies in: \\
$E_8$, $D_8$, $E_6$+$A_2$, \\
$D_5$+$A_3$, \end{tabular} & \begin{tabular}{p{2cm}}
 $\mathfrak{k}_{ss}$: $D_5$+2$A_1$\\
$C(\mathfrak{k_{ss}})_{ss}$: -\\

$\mathfrak{k}_{ss}$ lies in: \\
$E_7$+$A_1$, $D_7$, $D_6$+2$A_1$, \\
$D_5$+$A_3$, \end{tabular} & \begin{tabular}{p{2cm}}
 $\mathfrak{k}_{ss}$: $A_5$+$A_2$\\
$C(\mathfrak{k_{ss}})_{ss}$: $A_1$\\

$\mathfrak{k}_{ss}$ lies in: \\
$A_5$+$A_2$+$A_1$, $E_6$+$A_2$, \\
$A_8$, $E_7$, \end{tabular}\\
\hline\begin{tabular}{p{2cm}}
 $\mathfrak{k}_{ss}$: $A_5$+2$A_1$\\
$C(\mathfrak{k_{ss}})_{ss}$: -\\

$\mathfrak{k}_{ss}$ lies in: \\
$E_6$+$A_1$, $E_7$+$A_1$, \\
$D_8$, $D_6$+2$A_1$, $A_7$+$A_1$, \\
$A_5$+$A_2$+$A_1$, \end{tabular} & \begin{tabular}{p{2cm}}
 $\mathfrak{k}_{ss}$: $D_4$+$A_3$\\
$C(\mathfrak{k_{ss}})_{ss}$: -\\

$\mathfrak{k}_{ss}$ lies in: \\
$D_8$, $D_7$, 2$D_4$, $D_5$+$A_3$, \\
\end{tabular} & \begin{tabular}{p{2cm}}
 $\mathfrak{k}_{ss}$: $D_4$+3$A_1$\\
$C(\mathfrak{k_{ss}})_{ss}$: $A_1$\\

$\mathfrak{k}_{ss}$ lies in: \\
$D_4$+4$A_1$, $D_6$+2$A_1$, \\
$D_6$+$A_1$, 2$D_4$, \end{tabular} & \begin{tabular}{p{2cm}}
 $\mathfrak{k}_{ss}$: $A_4$+$A_3$\\
$C(\mathfrak{k_{ss}})_{ss}$: -\\

$\mathfrak{k}_{ss}$ lies in: \\
$D_8$, $E_8$, $A_8$, $D_5$+$A_3$, \\
2$A_4$, \end{tabular}\\
\hline
\begin{tabular}{p{2cm}}
 $\mathfrak{k}_{ss}$: $A_4$+$A_2$+$A_1$\\
$C(\mathfrak{k_{ss}})_{ss}$: -\\

$\mathfrak{k}_{ss}$ lies in: \\
$E_7$+$A_1$, $E_8$, $E_6$+$A_2$, \\
$A_7$+$A_1$, 2$A_4$, $A_5$+$A_2$+$A_1$, \\
\end{tabular} & \begin{tabular}{p{2cm}}
 $\mathfrak{k}_{ss}$: 2$A_3$+$A_1$\\
$C(\mathfrak{k_{ss}})_{ss}$: $A_1$\\

$\mathfrak{k}_{ss}$ lies in: \\
2$A_3$+2$A_1$, $D_5$+$A_3$, \\
$D_6$+$A_1$, $E_7$, $A_7$+$A_1$, \\
\end{tabular} & \begin{tabular}{p{2cm}}
 $\mathfrak{k}_{ss}$: $A_3$+$A_2$+2$A_1$\\
$C(\mathfrak{k_{ss}})_{ss}$: -\\

$\mathfrak{k}_{ss}$ lies in: \\
$E_7$+$A_1$, $D_6$+2$A_1$, \\
$D_5$+$A_2$, $D_5$+$A_3$, \\
$A_5$+$A_2$+$A_1$, 2$A_3$+2$A_1$, \\
\end{tabular} & \begin{tabular}{p{2cm}}
 $\mathfrak{k}_{ss}$: $A_3$+4$A_1$\\
$C(\mathfrak{k_{ss}})_{ss}$: -\\

$\mathfrak{k}_{ss}$ lies in: \\
$D_6$+2$A_1$, $D_5$+2$A_1$, \\
$D_4$+$A_3$, $D_4$+4$A_1$, \\
2$A_3$+2$A_1$, \end{tabular}\\
\hline\begin{tabular}{p{2cm}}
 $\mathfrak{k}_{ss}$: 3$A_2$+$A_1$\\
$C(\mathfrak{k_{ss}})_{ss}$: -\\

$\mathfrak{k}_{ss}$ lies in: \\
$A_5$+$A_2$+$A_1$, $E_6$+$A_2$, \\
$E_6$+$A_1$, 4$A_2$, \end{tabular} & \begin{tabular}{p{2cm}}
 $\mathfrak{k}_{ss}$: 7$A_1$\\
$C(\mathfrak{k_{ss}})_{ss}$: $A_1$\\

$\mathfrak{k}_{ss}$ lies in: \\
8$A_1$, $D_4$+3$A_1$, $D_4$+4$A_1$, \\
\end{tabular} & \begin{tabular}{p{2cm}}
 $\mathfrak{k}_{ss}$: $E_6$\\
$C(\mathfrak{k_{ss}})_{ss}$: $A_2$\\

$\mathfrak{k}_{ss}$ lies in: \\
$E_6$+$A_1$, $E_7$, \end{tabular} & \begin{tabular}{p{2cm}}
 $\mathfrak{k}_{ss}$: $D_6$\\
$C(\mathfrak{k_{ss}})_{ss}$: 2$A_1$\\

$\mathfrak{k}_{ss}$ lies in: \\
$D_6$+$A_1$, $E_7$, $D_7$, \\
\end{tabular}\\
\hline\begin{tabular}{p{2cm}}
 $\mathfrak{k}_{ss}$: $A_6$\\
$C(\mathfrak{k_{ss}})_{ss}$: $A_1$\\

$\mathfrak{k}_{ss}$ lies in: \\
$A_6$+$A_1$, $D_7$, $E_7$, \\
$A_7$, $A_7$, \end{tabular} & \begin{tabular}{p{2cm}}
 $\mathfrak{k}_{ss}$: $D_5$+$A_1$\\
$C(\mathfrak{k_{ss}})_{ss}$: $A_1$\\

$\mathfrak{k}_{ss}$ lies in: \\
$D_5$+2$A_1$, $D_7$, $E_7$, \\
$E_6$+$A_1$, $D_6$+$A_1$, \\
$D_5$+$A_2$, \end{tabular} & \begin{tabular}{p{2cm}}
 $\mathfrak{k}_{ss}$: $A_5$+$A_1$\\
$C(\mathfrak{k_{ss}})_{ss}$: $A_1$\\

$\mathfrak{k}_{ss}$ lies in: \\
$A_5$+2$A_1$, $A_7$, $E_6$+$A_1$, \\
$E_7$, $A_6$+$A_1$, $D_6$+$A_1$, \\
$A_5$+$A_2$, \end{tabular} & \begin{tabular}{p{2cm}}
 $\mathfrak{k}_{ss}$: $A_5$+$A_1$\\
$C(\mathfrak{k_{ss}})_{ss}$: $A_2$\\

$\mathfrak{k}_{ss}$ lies in: \\
$A_5$+2$A_1$, $D_6$+$A_1$, \\
$A_7$, $E_6$, \end{tabular}\\
\hline\begin{tabular}{p{2cm}}
 $\mathfrak{k}_{ss}$: $D_4$+$A_2$\\
$C(\mathfrak{k_{ss}})_{ss}$: -\\

$\mathfrak{k}_{ss}$ lies in: \\
$D_7$, $D_4$+$A_3$, $D_5$+$A_2$, \\
\end{tabular} & \begin{tabular}{p{2cm}}
 $\mathfrak{k}_{ss}$: $D_4$+2$A_1$\\
$C(\mathfrak{k_{ss}})_{ss}$: 2$A_1$\\

$\mathfrak{k}_{ss}$ lies in: \\
$D_4$+3$A_1$, $D_6$+$A_1$, \\
$D_5$+2$A_1$, $D_4$+$A_3$, \\
$D_6$, \end{tabular} & \begin{tabular}{p{2cm}}
 $\mathfrak{k}_{ss}$: $A_4$+$A_2$\\
$C(\mathfrak{k_{ss}})_{ss}$: $A_1$\\

$\mathfrak{k}_{ss}$ lies in: \\
$A_4$+$A_2$+$A_1$, $A_7$, \\
$E_7$, $D_5$+$A_2$, $A_4$+$A_3$, \\
$A_5$+$A_2$, $A_7$, \end{tabular} & \begin{tabular}{p{2cm}}
 $\mathfrak{k}_{ss}$: $A_4$+2$A_1$\\
$C(\mathfrak{k_{ss}})_{ss}$: -\\

$\mathfrak{k}_{ss}$ lies in: \\
$A_5$+2$A_1$, $E_6$+$A_1$, \\
$A_4$+$A_3$, $D_7$, $A_6$+$A_1$, \\
$D_5$+2$A_1$, $A_4$+$A_2$+$A_1$, \\
\end{tabular}\\
\hline\begin{tabular}{p{2cm}}
 $\mathfrak{k}_{ss}$: 2$A_3$\\
$C(\mathfrak{k_{ss}})_{ss}$: 2$A_1$\\

$\mathfrak{k}_{ss}$ lies in: \\
2$A_3$+$A_1$, $A_7$, $D_4$+$A_3$, \\
$D_6$, \end{tabular} & \begin{tabular}{p{2cm}}
 $\mathfrak{k}_{ss}$: 2$A_3$\\
$C(\mathfrak{k_{ss}})_{ss}$: -\\

$\mathfrak{k}_{ss}$ lies in: \\
$D_7$, $A_7$, $D_4$+$A_3$, \\
$A_4$+$A_3$, \end{tabular} & \begin{tabular}{p{2cm}}
 $\mathfrak{k}_{ss}$: $A_3$+$A_2$+$A_1$\\
$C(\mathfrak{k_{ss}})_{ss}$: $A_1$\\

$\mathfrak{k}_{ss}$ lies in: \\
$A_3$+$A_2$+2$A_1$, $E_7$, \\
$A_4$+$A_2$+$A_1$, $A_4$+$A_3$, \\
$A_6$+$A_1$, $A_5$+$A_2$, \\
$D_5$+$A_2$, 2$A_3$+$A_1$, \\
$D_6$+$A_1$, \end{tabular} & \begin{tabular}{p{2cm}}
 $\mathfrak{k}_{ss}$: $A_3$+3$A_1$\\
$C(\mathfrak{k_{ss}})_{ss}$: $A_1$\\

$\mathfrak{k}_{ss}$ lies in: \\
$A_3$+4$A_1$, $D_5$+2$A_1$, \\
$D_5$+$A_1$, $D_6$+$A_1$, \\
$D_4$+$A_3$, $A_5$+2$A_1$, \\
$D_4$+3$A_1$, $A_3$+$A_2$+2$A_1$, \\
2$A_3$+$A_1$, \end{tabular}\\
\hline\begin{tabular}{p{2cm}}
 $\mathfrak{k}_{ss}$: 3$A_2$\\
$C(\mathfrak{k_{ss}})_{ss}$: $A_2$\\

$\mathfrak{k}_{ss}$ lies in: \\
3$A_2$+$A_1$, $A_5$+$A_2$, \\
$E_6$, \end{tabular} & \begin{tabular}{p{2cm}}
 $\mathfrak{k}_{ss}$: 2$A_2$+2$A_1$\\
$C(\mathfrak{k_{ss}})_{ss}$: -\\

$\mathfrak{k}_{ss}$ lies in: \\
$A_3$+$A_2$+2$A_1$, $D_5$+$A_2$, \\
$E_6$+$A_1$, $A_4$+$A_2$+$A_1$, \\
$A_5$+2$A_1$, 3$A_2$+$A_1$, \\
\end{tabular} & \begin{tabular}{p{2cm}}
 $\mathfrak{k}_{ss}$: $A_2$+4$A_1$\\
$C(\mathfrak{k_{ss}})_{ss}$: -\\

$\mathfrak{k}_{ss}$ lies in: \\
$D_5$+2$A_1$, $A_3$+4$A_1$, \\
$D_4$+$A_2$, $A_3$+$A_2$+2$A_1$, \\
\end{tabular} & \begin{tabular}{p{2cm}}
 $\mathfrak{k}_{ss}$: 6$A_1$\\
$C(\mathfrak{k_{ss}})_{ss}$: 2$A_1$\\

$\mathfrak{k}_{ss}$ lies in: \\
7$A_1$, $D_4$+2$A_1$, $D_4$+3$A_1$, \\
$A_3$+4$A_1$, \end{tabular}\\
\hline\begin{tabular}{p{2cm}}
 $\mathfrak{k}_{ss}$: $D_5$\\
$C(\mathfrak{k_{ss}})_{ss}$: $A_3$\\

$\mathfrak{k}_{ss}$ lies in: \\
$D_5$+$A_1$, $E_6$, $D_6$, \\
\end{tabular} & \begin{tabular}{p{2cm}}
 $\mathfrak{k}_{ss}$: $A_5$\\
$C(\mathfrak{k_{ss}})_{ss}$: $A_2$+$A_1$\\

$\mathfrak{k}_{ss}$ lies in: \\
$A_5$+$A_1$, $A_5$+$A_1$, \\
$D_6$, $A_6$, $E_6$, \end{tabular} & \begin{tabular}{p{2cm}}
 $\mathfrak{k}_{ss}$: $D_4$+$A_1$\\
$C(\mathfrak{k_{ss}})_{ss}$: 3$A_1$\\

$\mathfrak{k}_{ss}$ lies in: \\
$D_4$+2$A_1$, $D_5$+$A_1$, \\
$D_6$, $D_4$+$A_2$, \end{tabular} & \begin{tabular}{p{2cm}}
 $\mathfrak{k}_{ss}$: $A_4$+$A_1$\\
$C(\mathfrak{k_{ss}})_{ss}$: $A_2$\\

$\mathfrak{k}_{ss}$ lies in: \\
$A_4$+2$A_1$, $A_5$+$A_1$, \\
$A_5$+$A_1$, $A_6$, $D_5$+$A_1$, \\
$A_4$+$A_2$, $E_6$, \end{tabular}\\
\hline\begin{tabular}{p{2cm}}
 $\mathfrak{k}_{ss}$: $A_3$+$A_2$\\
$C(\mathfrak{k_{ss}})_{ss}$: 2$A_1$\\

$\mathfrak{k}_{ss}$ lies in: \\
$A_3$+$A_2$+$A_1$, 2$A_3$, \\
$D_6$, $A_4$+$A_2$, $D_4$+$A_2$, \\
$A_6$, 2$A_3$, \end{tabular} & \begin{tabular}{p{2cm}}
 $\mathfrak{k}_{ss}$: $A_3$+2$A_1$\\
$C(\mathfrak{k_{ss}})_{ss}$: $A_3$\\

$\mathfrak{k}_{ss}$ lies in: \\
$A_3$+3$A_1$, $D_4$+2$A_1$, \\
2$A_3$, $A_5$+$A_1$, $D_5$, \\
\end{tabular} & \begin{tabular}{p{2cm}}
 $\mathfrak{k}_{ss}$: $A_3$+2$A_1$\\
$C(\mathfrak{k_{ss}})_{ss}$: 2$A_1$\\

$\mathfrak{k}_{ss}$ lies in: \\
$A_3$+3$A_1$, $D_5$+$A_1$, \\
$A_5$+$A_1$, $A_4$+2$A_1$, \\
$D_6$, 2$A_3$, $D_4$+2$A_1$, \\
$A_3$+$A_2$+$A_1$, \end{tabular} & \begin{tabular}{p{2cm}}
 $\mathfrak{k}_{ss}$: 2$A_2$+$A_1$\\
$C(\mathfrak{k_{ss}})_{ss}$: $A_2$\\

$\mathfrak{k}_{ss}$ lies in: \\
2$A_2$+2$A_1$, $A_3$+$A_2$+$A_1$, \\
$A_5$+$A_1$, $E_6$, $A_4$+$A_2$, \\
$A_5$+$A_1$, 3$A_2$, \end{tabular}\\
\hline\begin{tabular}{p{2cm}}
 $\mathfrak{k}_{ss}$: $A_2$+3$A_1$\\
$C(\mathfrak{k_{ss}})_{ss}$: $A_1$\\

$\mathfrak{k}_{ss}$ lies in: \\
$A_2$+4$A_1$, $D_5$+$A_1$, \\
2$A_2$+2$A_1$, $A_3$+$A_2$+$A_1$, \\
$A_4$+2$A_1$, $A_3$+3$A_1$, \\
$D_4$+$A_2$, \end{tabular} & \begin{tabular}{p{2cm}}
 $\mathfrak{k}_{ss}$: 5$A_1$\\
$C(\mathfrak{k_{ss}})_{ss}$: 3$A_1$\\

$\mathfrak{k}_{ss}$ lies in: \\
6$A_1$, $A_3$+3$A_1$, $D_4$+2$A_1$, \\
$A_2$+4$A_1$, $D_4$+$A_1$, \\
\end{tabular} & \begin{tabular}{p{2cm}}
 $\mathfrak{k}_{ss}$: $D_4$\\
$C(\mathfrak{k_{ss}})_{ss}$: $D_4$\\

$\mathfrak{k}_{ss}$ lies in: \\
$D_4$+$A_1$, $D_5$, \end{tabular} & \begin{tabular}{p{2cm}}
 $\mathfrak{k}_{ss}$: $A_4$\\
$C(\mathfrak{k_{ss}})_{ss}$: $A_4$\\

$\mathfrak{k}_{ss}$ lies in: \\
$A_4$+$A_1$, $A_5$, $D_5$, \\
\end{tabular}\\
\hline\begin{tabular}{p{2cm}}
 $\mathfrak{k}_{ss}$: $A_3$+$A_1$\\
$C(\mathfrak{k_{ss}})_{ss}$: $A_3$+$A_1$\\

$\mathfrak{k}_{ss}$ lies in: \\
$A_3$+2$A_1$, $A_3$+2$A_1$, \\
$A_4$+$A_1$, $A_3$+$A_2$, \\
$D_4$+$A_1$, $A_5$, $D_5$, \\
\end{tabular} & \begin{tabular}{p{2cm}}
 $\mathfrak{k}_{ss}$: 2$A_2$\\
$C(\mathfrak{k_{ss}})_{ss}$: 2$A_2$\\

$\mathfrak{k}_{ss}$ lies in: \\
2$A_2$+$A_1$, $A_3$+$A_2$, \\
$A_5$, \end{tabular} & \begin{tabular}{p{2cm}}
 $\mathfrak{k}_{ss}$: $A_2$+2$A_1$\\
$C(\mathfrak{k_{ss}})_{ss}$: $A_3$\\

$\mathfrak{k}_{ss}$ lies in: \\
$A_2$+3$A_1$, $A_4$+$A_1$, \\
$A_3$+2$A_1$, $A_3$+$A_2$, \\
$A_3$+2$A_1$, $D_5$, 2$A_2$+$A_1$, \\
\end{tabular} & \begin{tabular}{p{2cm}}
 $\mathfrak{k}_{ss}$: 4$A_1$\\
$C(\mathfrak{k_{ss}})_{ss}$: 4$A_1$\\

$\mathfrak{k}_{ss}$ lies in: \\
5$A_1$, $A_2$+3$A_1$, $D_4$+$A_1$, \\
$A_3$+2$A_1$, \end{tabular}\\
\hline\begin{tabular}{p{2cm}}
 $\mathfrak{k}_{ss}$: 4$A_1$\\
$C(\mathfrak{k_{ss}})_{ss}$: $D_4$\\

$\mathfrak{k}_{ss}$ lies in: \\
5$A_1$, $A_3$+2$A_1$, $D_4$, \\
\end{tabular} & \begin{tabular}{p{2cm}}
 $\mathfrak{k}_{ss}$: $A_3$\\
$C(\mathfrak{k_{ss}})_{ss}$: $D_5$\\

$\mathfrak{k}_{ss}$ lies in: \\
$A_3$+$A_1$, $A_4$, $D_4$, \\
\end{tabular} & \begin{tabular}{p{2cm}}
 $\mathfrak{k}_{ss}$: $A_2$+$A_1$\\
$C(\mathfrak{k_{ss}})_{ss}$: $A_5$\\

$\mathfrak{k}_{ss}$ lies in: \\
$A_2$+2$A_1$, $A_3$+$A_1$, \\
2$A_2$, $A_4$, \end{tabular} & \begin{tabular}{p{2cm}}
 $\mathfrak{k}_{ss}$: 3$A_1$\\
$C(\mathfrak{k_{ss}})_{ss}$: $D_4$+$A_1$\\

$\mathfrak{k}_{ss}$ lies in: \\
4$A_1$, 4$A_1$, $A_2$+2$A_1$, \\
$A_3$+$A_1$, $D_4$, \end{tabular}\\
\hline\begin{tabular}{p{2cm}}
 $\mathfrak{k}_{ss}$: $A_2$\\
$C(\mathfrak{k_{ss}})_{ss}$: $E_6$\\

$\mathfrak{k}_{ss}$ lies in: \\
$A_2$+$A_1$, $A_3$, \end{tabular} & \begin{tabular}{p{2cm}}
 $\mathfrak{k}_{ss}$: 2$A_1$\\
$C(\mathfrak{k_{ss}})_{ss}$: $D_6$\\

$\mathfrak{k}_{ss}$ lies in: \\
3$A_1$, $A_2$+$A_1$, $A_3$, \\
\end{tabular} & \begin{tabular}{p{2cm}}
 $\mathfrak{k}_{ss}$: $A_1$\\
$C(\mathfrak{k_{ss}})_{ss}$: $E_7$\\

$\mathfrak{k}_{ss}$ lies in: \\
2$A_1$, $A_2$, \end{tabular} & \begin{tabular}{p{2cm}}
 $\mathfrak{k}_{ss}$: -\\
$C(\mathfrak{k_{ss}})_{ss}$: $E_8$\\

$\mathfrak{k}_{ss}$ lies in: \\
$A_1$, \end{tabular}\\
\hline\end{longtable}



\section{Cardinalities of groups preserving $\Delta(\bb\cap\kk)$}\label{secCardinalitiesIsos}
Define $W'$ to be the group of root system automorphisms of $\Delta(\gG)$ that preserve $\Delta(\bb\cap\kk)$. $\Delta(\bb\cap\kk)$ $W'=W'''\rtimes W''$ is the semidirect sum of the Weyl group $W'''$ of $C(\kk_{ss})$ with the group $W''$ of graph automorphisms of $(\Delta(\kk_{ss})\oplus \Delta(C(\kk_{ss}))\cap\Delta(\bb)$ that preserve $\Delta(\kk_{ss})$ and $\Delta(C(\kk_{ss})$ and extend to root system automorphisms of $\Delta(\gG)$.\label{tableGroups}


\subsection{$F_4$}


\begin{tabular}{ccccc}$\mathfrak{k}_{ss}$& $C(k_{ss})_{ss}$ & $\#W''$ &$\#W'''$&$\#(W'''\rtimes W'')$\\\hline$F_4$ &  & 1 & 1 & 1 \\
$B_4$ &  & 1 & 1 & 1 \\
$D_4$ &  & 6 & 1 & 6 \\
$C_3$+$A_1$ &  & 1 & 1 & 1 \\
$A_3$+$A_1'$ &  & 2 & 1 & 2 \\
$B_2=C_2$+2$A_1$ &  & 2 & 1 & 2 \\
$A_2'$+$A_2$ &  & 2 & 1 & 2 \\
4$A_1$ &  & 24 & 1 & 24 \\
$C_3$ & $A_1$ & 1 & 2 & 2 \\
$B_3$ &  & 1 & 1 & 1 \\
$A_3$ & $A_1'$ & 2 & 2 & 4 \\
$B_2=C_2$+$A_1$ & $A_1$ & 1 & 2 & 2 \\
$A_2'$+$A_1$ &  & 2 & 1 & 2 \\
$A_2$+$A_1'$ &  & 2 & 1 & 2 \\
$A_1'$+2$A_1$ &  & 2 & 1 & 2 \\
3$A_1$ & $A_1$ & 6 & 2 & 12 \\
$B_2=C_2$ & 2$A_1$ & 2 & 4 & 8 \\
$A_2'$ & $A_2$ & 2 & 6 & 12 \\
$A_2$ & $A_2'$ & 2 & 6 & 12 \\
$A_1'$+$A_1$ & $A_1$ & 1 & 2 & 2 \\
2$A_1$ & $B_2=C_2$ & 2 & 8 & 16 \\
$A_1'$ & $A_3$ & 2 & 24 & 48 \\
$A_1$ & $C_3$ & 1 & 48 & 48 \\
- & $F_4$ & 1 & 1152 & 1152 \\
\end{tabular}


\subsection{$E_6$}


\begin{tabular}{ccccc}$\mathfrak{k}_{ss}$& $C(k_{ss})_{ss}$ & $\#W''$ &$\#W'''$&$\#(W'''\rtimes W'')$\\\hline$E_6$ &  & 2 & 1 & 2 \\
$A_5$+$A_1$ &  & 2 & 1 & 2 \\
3$A_2$ &  & 12 & 1 & 12 \\
$D_5$ &  & 2 & 1 & 2 \\
$A_5$ & $A_1$ & 2 & 2 & 4 \\
$A_4$+$A_1$ &  & 2 & 1 & 2 \\
$A_3$+2$A_1$ &  & 4 & 1 & 4 \\
2$A_2$+$A_1$ &  & 4 & 1 & 4 \\
$D_4$ &  & 12 & 1 & 12 \\
$A_4$ & $A_1$ & 2 & 2 & 4 \\
$A_3$+$A_1$ & $A_1$ & 2 & 2 & 4 \\
2$A_2$ & $A_2$ & 4 & 6 & 24 \\
$A_2$+2$A_1$ &  & 4 & 1 & 4 \\
4$A_1$ &  & 48 & 1 & 48 \\
$A_3$ & 2$A_1$ & 4 & 4 & 16 \\
$A_2$+$A_1$ & $A_2$ & 2 & 6 & 12 \\
3$A_1$ & $A_1$ & 12 & 2 & 24 \\
$A_2$ & 2$A_2$ & 4 & 36 & 144 \\
2$A_1$ & $A_3$ & 4 & 24 & 96 \\
$A_1$ & $A_5$ & 2 & 720 & 1440 \\
- & $E_6$ & 2 & 51840 & 103680 \\
\end{tabular}


\subsection{$E_7$}


\begin{tabular}{ccccc}
$\mathfrak{k}_{ss}$& $C(k_{ss})_{ss}$ & $\#W''$ &$\#W'''$&$\#(W'''\rtimes W'')$\\\hline$E_7$ &  & 1 & 1 & 1 \\
$A_7$ &  & 2 & 1 & 2 \\
$D_6$+$A_1$ &  & 1 & 1 & 1 \\
$A_5$+$A_2$ &  & 2 & 1 & 2 \\
$D_4$+3$A_1$ &  & 6 & 1 & 6 \\
2$A_3$+$A_1$ &  & 4 & 1 & 4 \\
7$A_1$ &  & 168 & 1 & 168 \\
$E_6$ &  & 2 & 1 & 2 \\
$D_6$ & $A_1$ & 1 & 2 & 2 \\
$A_6$ &  & 2 & 1 & 2 \\
$D_5$+$A_1$ &  & 2 & 1 & 2 \\
$A_5$+$A_1$ &  & 2 & 1 & 2 \\
$D_4$+2$A_1$ & $A_1$ & 2 & 2 & 4 \\
$A_4$+$A_2$ &  & 2 & 1 & 2 \\
2$A_3$ & $A_1$ & 4 & 2 & 8 \\
$A_3$+$A_2$+$A_1$ &  & 2 & 1 & 2 \\
$A_3$+3$A_1$ &  & 4 & 1 & 4 \\
3$A_2$ &  & 12 & 1 & 12 \\
6$A_1$ & $A_1$ & 24 & 2 & 48 \\
$D_5$ & $A_1$ & 2 & 2 & 4 \\
$A_5$ & $A_1$ & 2 & 2 & 4 \\
$A_5$ & $A_2$ & 2 & 6 & 12 \\
$D_4$+$A_1$ & 2$A_1$ & 2 & 4 & 8 \\
$A_4$+$A_1$ &  & 2 & 1 & 2 \\
$A_3$+$A_2$ & $A_1$ & 2 & 2 & 4 \\
$A_3$+2$A_1$ & $A_1$ & 2 & 2 & 4 \\
2$A_2$+$A_1$ &  & 4 & 1 & 4 \\
$A_2$+3$A_1$ &  & 12 & 1 & 12 \\
5$A_1$ & 2$A_1$ & 8 & 4 & 32 \\
$D_4$ & 3$A_1$ & 6 & 8 & 48 \\
$A_4$ & $A_2$ & 2 & 6 & 12 \\
$A_3$+$A_1$ & 2$A_1$ & 2 & 4 & 8 \\
$A_3$+$A_1$ & $A_3$ & 2 & 24 & 48 \\
2$A_2$ & $A_2$ & 4 & 6 & 24 \\
$A_2$+2$A_1$ & $A_1$ & 4 & 2 & 8 \\
4$A_1$ & 3$A_1$ & 6 & 8 & 48 \\
$A_3$ & $A_3$+$A_1$ & 2 & 48 & 96 \\
$A_2$+$A_1$ & $A_3$ & 2 & 24 & 48 \\
3$A_1$ & 4$A_1$ & 6 & 16 & 96 \\
3$A_1$ & $D_4$ & 6 & 192 & 1152 \\
$A_2$ & $A_5$ & 2 & 720 & 1440 \\
2$A_1$ & $D_4$+$A_1$ & 2 & 384 & 768 \\
$A_1$ & $D_6$ & 1 & 23040 & 23040 \\
- & $E_7$ & 1 & 2903040 & 2903040 \\
\end{tabular}

\subsection{$E_8$}
\begin{longtable}{ccccc}
$\mathfrak{k}_{ss}$& $C(k_{ss})_{ss}$ & $\#W''$ &$\#W'''$&$\#(W'''\rtimes W'')$ \endhead 
$E_8$ &  & 1 & 1 & 1\\
$D_8$ &  & 1 & 1 & 1 \\
$A_8$ &  & 2 & 1 & 2 \\
$E_7$+$A_1$ &  & 1 & 1 & 1 \\
$A_7$+$A_1$ &  & 2 & 1 & 2 \\
$E_6$+$A_2$ &  & 2 & 1 & 2 \\
$D_6$+2$A_1$ &  & 2 & 1 & 2 \\
$D_5$+$A_3$ &  & 2 & 1 & 2 \\
$A_5$+$A_2$+$A_1$ &  & 2 & 1 & 2 \\
2$D_4$ &  & 12 & 1 & 12 \\
$D_4$+4$A_1$ &  & 24 & 1 & 24 \\
2$A_4$ &  & 4 & 1 & 4 \\
2$A_3$+2$A_1$ &  & 8 & 1 & 8 \\
4$A_2$ &  & 48 & 1 & 48 \\
8$A_1$ &  & 1344 & 1 & 1344 \\
$E_7$ & $A_1$ & 1 & 2 & 2 \\
$D_7$ &  & 2 & 1 & 2 \\
$A_7$ &  & 2 & 1 & 2 \\
$A_7$ & $A_1$ & 2 & 2 & 4 \\
$E_6$+$A_1$ &  & 2 & 1 & 2 \\
$D_6$+$A_1$ & $A_1$ & 1 & 2 & 2 \\
$A_6$+$A_1$ &  & 2 & 1 & 2 \\
$D_5$+$A_2$ &  & 2 & 1 & 2 \\
$D_5$+2$A_1$ &  & 4 & 1 & 4 \\
$A_5$+$A_2$ & $A_1$ & 2 & 2 & 4 \\
$A_5$+2$A_1$ &  & 2 & 1 & 2 \\
$D_4$+$A_3$ &  & 4 & 1 & 4 \\
$D_4$+3$A_1$ & $A_1$ & 6 & 2 & 12 \\
$A_4$+$A_3$ &  & 2 & 1 & 2 \\
$A_4$+$A_2$+$A_1$ &  & 2 & 1 & 2 \\
2$A_3$+$A_1$ & $A_1$ & 4 & 2 & 8 \\
$A_3$+$A_2$+2$A_1$ &  & 4 & 1 & 4 \\
$A_3$+4$A_1$ &  & 16 & 1 & 16 \\
3$A_2$+$A_1$ &  & 12 & 1 & 12 \\
7$A_1$ & $A_1$ & 168 & 2 & 336 \\
$E_6$ & $A_2$ & 2 & 6 & 12 \\
$D_6$ & 2$A_1$ & 2 & 4 & 8 \\
$A_6$ & $A_1$ & 2 & 2 & 4 \\
$D_5$+$A_1$ & $A_1$ & 2 & 2 & 4 \\
$A_5$+$A_1$ & $A_1$ & 2 & 2 & 4 \\
$A_5$+$A_1$ & $A_2$ & 2 & 6 & 12 \\
$D_4$+$A_2$ &  & 12 & 1 & 12 \\
$D_4$+2$A_1$ & 2$A_1$ & 4 & 4 & 16 \\
$A_4$+$A_2$ & $A_1$ & 2 & 2 & 4 \\
$A_4$+2$A_1$ &  & 4 & 1 & 4 \\
2$A_3$ & 2$A_1$ & 8 & 4 & 32 \\
2$A_3$ &  & 8 & 1 & 8 \\
$A_3$+$A_2$+$A_1$ & $A_1$ & 2 & 2 & 4 \\
$A_3$+3$A_1$ & $A_1$ & 4 & 2 & 8 \\
3$A_2$ & $A_2$ & 12 & 6 & 72 \\
2$A_2$+2$A_1$ &  & 8 & 1 & 8 \\
$A_2$+4$A_1$ &  & 48 & 1 & 48 \\
6$A_1$ & 2$A_1$ & 48 & 4 & 192 \\
$D_5$ & $A_3$ & 2 & 24 & 48 \\
$A_5$ & $A_2$+$A_1$ & 2 & 12 & 24 \\
$D_4$+$A_1$ & 3$A_1$ & 6 & 8 & 48 \\
$A_4$+$A_1$ & $A_2$ & 2 & 6 & 12 \\
$A_3$+$A_2$ & 2$A_1$ & 4 & 4 & 16 \\
$A_3$+2$A_1$ & $A_3$ & 4 & 24 & 96 \\
$A_3$+2$A_1$ & 2$A_1$ & 4 & 4 & 16 \\
2$A_2$+$A_1$ & $A_2$ & 4 & 6 & 24 \\
$A_2$+3$A_1$ & $A_1$ & 12 & 2 & 24 \\
5$A_1$ & 3$A_1$ & 24 & 8 & 192 \\
$D_4$ & $D_4$ & 6 & 192 & 1152 \\
$A_4$ & $A_4$ & 2 & 120 & 240 \\
$A_3$+$A_1$ & $A_3$+$A_1$ & 2 & 48 & 96 \\
2$A_2$ & 2$A_2$ & 8 & 36 & 288 \\
$A_2$+2$A_1$ & $A_3$ & 4 & 24 & 96 \\
4$A_1$ & 4$A_1$ & 24 & 16 & 384 \\
4$A_1$ & $D_4$ & 24 & 192 & 4608 \\
$A_3$ & $D_5$ & 2 & 1920 & 3840 \\
$A_2$+$A_1$ & $A_5$ & 2 & 720 & 1440 \\
3$A_1$ & $D_4$+$A_1$ & 6 & 384 & 2304 \\
$A_2$ & $E_6$ & 2 & 51840 & 103680 \\
2$A_1$ & $D_6$ & 2 & 23040 & 46080 \\
$A_1$ & $E_7$ & 1 & 2903040 & 2903040 \\
- & $E_8$ & 1 & 696729600 & 696729600 \\
\end{longtable}
\section{$\lL$-infinite weights for the exceptional Lie algebras}\label{secAppendixExceptional}
This section lists all minimal $\lL$-infinite relations and the corresponding two-sided weights under the assumption that $C(\kk_{ss})\cap\nn$ is the nilradical of a parabolic subalgebra of  $C(\kk_{ss})\cap\nn$ containing $C(\kk_{ss})\cap \hh$. %We have also indicated the $\eps$-coordinates with respect to the reductive root subalgebra generated by the roots participating in the relation and the roots of $\kk$.

In the tables to follow, under each root $\alpha_i$ (respectively, $\beta_j$) we write the type of the semisimple component of $\kk$ whose roots are linked to $\alpha_i$ (respectively, $\beta_j$). The ~~~'~~~ sign is used to distinguish different components of $\kk$ that have the same Dynkin type. In type $F_4$, the sign ~~~$'$~~~ stands for a component of $\kk$ whose roots are short. For example, $A_1'+A_1$ represents the direct sum of two $\sL(2)$, one with long and one with short roots, and $A_1$+$A_1$' stands for two long-root $\sL(2)$'s. For example, if a root $\alpha_i$ is linked to $A_1$', and a root $\beta_j$ is linked to $A_1$', then $\alpha_i$ and $\alpha_j$ are linked to the same component of $\Delta(\kk)$; similarly if a root $\alpha_i$ is linked to $A_1$, and a root $\beta_j$ is linked to $A_1$', then the two roots are linked to two different components of $\Delta(\kk)$.
\newpage
\addtolength{\hoffset}{-2.5cm}
\addtolength{\textwidth}{8cm}
\addtolength{\voffset}{-3.3cm}
\addtolength{\textheight}{5.3cm}


\subsection{$F_4$}
\noindent Number of different non-solvable subalgebras up to $\gG$-automorphism such that $\nn\cap C(\kk_{ss})$ is a nilradical of a parabolic subalgebra of $\kk$ containing $C(\kk_{ss})\cap \hh$: 503


\noindent Among them 234 satisfy the cone condition and 269 do not.

\tiny

\begin{longtable}{rcl p{1cm}p{1cm}p{3cm} } \multicolumn{3}{c}{ Relation / linked $\mathfrak{k}$-components} &$\alpha_i$'s, $\beta_i$'s generate & adding $\Delta(\mathfrak{k})$ generates&Non-zero scalar products\endhead
\hline\multicolumn{5}{c}{$\mathfrak{k}$-semisimple type: $A_3$}\\
\hline\hline\begin{tabular}{c}2$\alpha_1$\\\tiny{  }\end{tabular} &\begin{tabular}{c} = \\~ \end{tabular} & \begin{tabular}{cc}$\beta_1$+ & $\beta_2$\\\tiny{ $A_3$ } & \tiny{ $A_3$ }\end{tabular} & $B_2=C_2$ & $B_4$ &  $\langle\alpha_1, \beta_1\rangle=2$, $\langle\alpha_1, \beta_2\rangle=2$, 
\\
\hline\multicolumn{5}{c}{$\mathfrak{k}$-semisimple type: $B_2=C_2$+$A_1$}\\
\hline\hline\begin{tabular}{c}2$\alpha_1$\\\tiny{ $B_2=C_2$ }\end{tabular} &\begin{tabular}{c} = \\~ \end{tabular} & \begin{tabular}{cc}$\beta_1$+ & $\beta_2$\\\tiny{ $B_2=C_2$+$A_1$ } & \tiny{ $B_2=C_2$+$A_1$ }\end{tabular} & $B_2=C_2$ & $F_4$ &  $\langle\alpha_1, \beta_1\rangle=2$, $\langle\alpha_1, \beta_2\rangle=2$, 
\\
\hline\multicolumn{5}{c}{$\mathfrak{k}$-semisimple type: $A_2$+$A_1'$}\\
\hline\hline\begin{tabular}{c}2$\alpha_1$\\\tiny{ $A_1'$ }\end{tabular} &\begin{tabular}{c} = \\~ \end{tabular} & \begin{tabular}{cc}$\beta_1$+ & $\beta_2$\\\tiny{ $A_2$+$A_1'$ } & \tiny{ $A_2$ }\end{tabular} & $B_2=C_2$ & $F_4$ &  $\langle\alpha_1, \beta_1\rangle=2$, $\langle\alpha_1, \beta_2\rangle=2$, 
\\
\hline\multicolumn{5}{c}{$\mathfrak{k}$-semisimple type: $A_1$+$A_1$+$A_1$}\\
\hline\hline\begin{tabular}{c}2$\alpha_1$\\\tiny{ $A_1$ }\end{tabular} &\begin{tabular}{c} = \\~ \end{tabular} & \begin{tabular}{cc}$\beta_1$+ & $\beta_2$\\\tiny{ $A_1$+$A_1$'+$A_1$'' } & \tiny{ $A_1$+$A_1$'+$A_1$'' }\end{tabular} & $B_2=C_2$ & $B_4$ &  $\langle\alpha_1, \beta_1\rangle=2$, $\langle\alpha_1, \beta_2\rangle=2$, 
\\
\hline\multicolumn{5}{c}{$\mathfrak{k}$-semisimple type: $B_2=C_2$}\\
\hline\hline\begin{tabular}{cc}$\alpha_1$+ & $\alpha_2$\\\tiny{  } & \tiny{  }\end{tabular} &\begin{tabular}{c} = \\~ \end{tabular} & \begin{tabular}{cc}$\beta_1$+ & $\beta_2$\\\tiny{ $B_2=C_2$ } & \tiny{ $B_2=C_2$ }\end{tabular} & $A_3$ & $B_4$ &  $\langle\alpha_1, \beta_1\rangle=2$, $\langle\alpha_1, \beta_2\rangle=2$, $\langle\alpha_2, \beta_1\rangle=2$, $\langle\alpha_2, \beta_2\rangle=2$, 
\\
\begin{tabular}{cc}$\alpha_1$+ & $\alpha_2$\\\tiny{ $B_2=C_2$ } & \tiny{  }\end{tabular} &\begin{tabular}{c} = \\~ \end{tabular} & \begin{tabular}{c}$\beta_1$\\\tiny{ $B_2=C_2$ }\end{tabular} & $A_2$ & $B_4$ & $\langle\alpha_1, \alpha_2\rangle=-2$,  $\langle\alpha_1, \beta_1\rangle=2$, $\langle\alpha_2, \beta_1\rangle=2$, 
\\
\begin{tabular}{c}2$\alpha_1$\\\tiny{ $B_2=C_2$ }\end{tabular} &\begin{tabular}{c} = \\~ \end{tabular} & \begin{tabular}{cc}$\beta_1$+ & $\beta_2$\\\tiny{ $B_2=C_2$ } & \tiny{ $B_2=C_2$ }\end{tabular} & $B_2=C_2$ & $F_4$ &  $\langle\alpha_1, \beta_1\rangle=2$, $\langle\alpha_1, \beta_2\rangle=2$, 
\\
\hline\multicolumn{5}{c}{$\mathfrak{k}$-semisimple type: $A_2'$}\\
\hline\hline\begin{tabular}{cc}$\alpha_1$+ & $\alpha_2$\\\tiny{  } & \tiny{  }\end{tabular} &\begin{tabular}{c} = \\~ \end{tabular} & \begin{tabular}{ccc}$\beta_1$+ & $\beta_2$+ & $\beta_3$\\\tiny{ $A_2'$ } & \tiny{ $A_2'$ } & \tiny{ $A_2'$ }\end{tabular} & $D_4$ & $F_4$ & $\langle\alpha_1, \alpha_2\rangle=2$,  $\langle\alpha_1, \beta_1\rangle=2$, $\langle\alpha_1, \beta_2\rangle=2$, $\langle\alpha_1, \beta_3\rangle=2$, $\langle\alpha_2, \beta_1\rangle=2$, $\langle\alpha_2, \beta_2\rangle=2$, $\langle\alpha_2, \beta_3\rangle=2$, 
\\
\begin{tabular}{cc}$\alpha_1$+ & $\alpha_2$\\\tiny{ $A_2'$ } & \tiny{  }\end{tabular} &\begin{tabular}{c} = \\~ \end{tabular} & \begin{tabular}{c}$\beta_1$\\\tiny{ $A_2'$ }\end{tabular} & $A_2$ & $F_4$ & $\langle\alpha_1, \alpha_2\rangle=-2$,  $\langle\alpha_1, \beta_1\rangle=2$, $\langle\alpha_2, \beta_1\rangle=2$, 
\\
\begin{tabular}{cc}$\alpha_1$+ & $\alpha_2$\\\tiny{ $A_2'$ } & \tiny{  }\end{tabular} &\begin{tabular}{c} = \\~ \end{tabular} & \begin{tabular}{cc}$\beta_1$+ & $\beta_2$\\\tiny{ $A_2'$ } & \tiny{ $A_2'$ }\end{tabular} & $A_3$ & $F_4$ &  $\langle\alpha_1, \beta_1\rangle=2$, $\langle\alpha_1, \beta_2\rangle=2$, $\langle\alpha_2, \beta_1\rangle=2$, $\langle\alpha_2, \beta_2\rangle=2$, 
\\
\hline\multicolumn{5}{c}{$\mathfrak{k}$-semisimple type: $A_2$}\\
\hline\hline\begin{tabular}{c}2$\alpha_1$\\\tiny{  }\end{tabular} &\begin{tabular}{c} = \\~ \end{tabular} & \begin{tabular}{cc}$\beta_1$+ & $\beta_2$\\\tiny{ $A_2$ } & \tiny{ $A_2$ }\end{tabular} & $B_2=C_2$ & $B_4$ &  $\langle\alpha_1, \beta_1\rangle=2$, $\langle\alpha_1, \beta_2\rangle=2$, 
\\
\begin{tabular}{c}2$\alpha_1$\\\tiny{ $A_2$ }\end{tabular} &\begin{tabular}{c} = \\~ \end{tabular} & \begin{tabular}{cc}$\beta_1$+ & $\beta_2$\\\tiny{ $A_2$ } & \tiny{ $A_2$ }\end{tabular} & $B_2=C_2$ & $B_4$ &  $\langle\alpha_1, \beta_1\rangle=2$, $\langle\alpha_1, \beta_2\rangle=2$, 
\\
\begin{tabular}{cc}$\alpha_1$+ & $\alpha_2$\\\tiny{ $A_2$ } & \tiny{  }\end{tabular} &\begin{tabular}{c} = \\~ \end{tabular} & \begin{tabular}{c}$\beta_1$\\\tiny{ $A_2$ }\end{tabular} & $B_2=C_2$ & $B_4$ & $\langle\alpha_1, \alpha_2\rangle=-2$,  $\langle\alpha_1, \beta_1\rangle=2$, 
\\
\begin{tabular}{cc}$\alpha_1$+ & $\alpha_2$\\\tiny{ $A_2$ } & \tiny{ $A_2$ }\end{tabular} &\begin{tabular}{c} = \\~ \end{tabular} & \begin{tabular}{c}2$\beta_1$\\\tiny{ $A_2$ }\end{tabular} & $B_2=C_2$ & $B_4$ &  $\langle\alpha_1, \beta_1\rangle=2$, $\langle\alpha_2, \beta_1\rangle=2$, 
\\
\begin{tabular}{cc}$\alpha_1$+ & $\alpha_2$\\\tiny{ $A_2$ } & \tiny{  }\end{tabular} &\begin{tabular}{c} = \\~ \end{tabular} & \begin{tabular}{c}$\beta_1$\\\tiny{ $A_2$ }\end{tabular} & $B_2=C_2$ & $B_4$ &  $\langle\alpha_1, \beta_1\rangle=2$, $\langle\alpha_2, \beta_1\rangle=2$, 
\\
\begin{tabular}{cc}$\alpha_1$+ & $\alpha_2$\\\tiny{ $A_2$ } & \tiny{  }\end{tabular} &\begin{tabular}{c} = \\~ \end{tabular} & \begin{tabular}{c}$\beta_1$\\\tiny{ $A_2$ }\end{tabular} & $A_2'$ & $F_4$ & $\langle\alpha_1, \alpha_2\rangle=-1$,  $\langle\alpha_1, \beta_1\rangle=1$, $\langle\alpha_2, \beta_1\rangle=1$, 
\\
\hline\multicolumn{5}{c}{$\mathfrak{k}$-semisimple type: $A_1'$+$A_1$}\\
\hline\hline\begin{tabular}{cc}$\alpha_1$+ & $\alpha_2$\\\tiny{ $A_1'$ } & \tiny{  }\end{tabular} &\begin{tabular}{c} = \\~ \end{tabular} & \begin{tabular}{c}$\beta_1$\\\tiny{ $A_1'$ }\end{tabular} & $B_2=C_2$ & $C_3$+$A_1$ & $\langle\alpha_1, \alpha_2\rangle=-2$,  $\langle\alpha_2, \beta_1\rangle=2$, 
\\
\begin{tabular}{cc}$\alpha_1$+ & $\alpha_2$\\\tiny{ $A_1'$+$A_1$ } & \tiny{ $A_1$ }\end{tabular} &\begin{tabular}{c} = \\~ \end{tabular} & \begin{tabular}{c}2$\beta_1$\\\tiny{ $A_1'$+$A_1$ }\end{tabular} & $B_2=C_2$ & $F_4$ &  $\langle\alpha_1, \beta_1\rangle=2$, $\langle\alpha_2, \beta_1\rangle=2$, 
\\
\begin{tabular}{cc}$\alpha_1$+ & $\alpha_2$\\\tiny{ $A_1$ } & \tiny{  }\end{tabular} &\begin{tabular}{c} = \\~ \end{tabular} & \begin{tabular}{c}$\beta_1$\\\tiny{ $A_1$ }\end{tabular} & $A_2$ & $A_3$+$A_1'$ & $\langle\alpha_1, \alpha_2\rangle=-2$,  $\langle\alpha_1, \beta_1\rangle=2$, $\langle\alpha_2, \beta_1\rangle=2$, 
\\
\begin{tabular}{cc}$\alpha_1$+ & $\alpha_2$\\\tiny{ $A_1'$ } & \tiny{ $A_1'$ }\end{tabular} &\begin{tabular}{c} = \\~ \end{tabular} & \begin{tabular}{c}$\beta_1$\\\tiny{ $A_1'$ }\end{tabular} & $B_2=C_2$ & $C_3$+$A_1$ &  $\langle\alpha_1, \beta_1\rangle=2$, $\langle\alpha_2, \beta_1\rangle=2$, 
\\
\begin{tabular}{cc}$\alpha_1$+ & $\alpha_2$\\\tiny{ $A_1$ } & \tiny{ $A_1'$ }\end{tabular} &\begin{tabular}{c} = \\~ \end{tabular} & \begin{tabular}{cc}$\beta_1$+ & $\beta_2$\\\tiny{ $A_1'$+$A_1$ } & \tiny{ $A_1'$ }\end{tabular} & $B_3$ & $F_4$ &  $\langle\alpha_1, \beta_1\rangle=2$, $\langle\alpha_1, \beta_2\rangle=2$, $\langle\alpha_2, \beta_2\rangle=2$, 
\\
\begin{tabular}{cc}$\alpha_1$+ & $\alpha_2$\\\tiny{ $A_1$ } & \tiny{ $A_1'$ }\end{tabular} &\begin{tabular}{c} = \\~ \end{tabular} & \begin{tabular}{cc}$\beta_1$+ & $\beta_2$\\\tiny{ $A_1$ } & \tiny{ $A_1'$ }\end{tabular} & $B_3$ & $F_4$ &  $\langle\alpha_1, \beta_1\rangle=2$, $\langle\alpha_1, \beta_2\rangle=2$, $\langle\alpha_2, \beta_1\rangle=2$, 
\\
\begin{tabular}{cc}$\alpha_1$+ & $\alpha_2$\\\tiny{ $A_1$ } & \tiny{ $A_1$ }\end{tabular} &\begin{tabular}{c} = \\~ \end{tabular} & \begin{tabular}{cc}2$\beta_1$+ & $\beta_2$\\\tiny{ $A_1'$+$A_1$ } & \tiny{ $A_1'$ }\end{tabular} & $B_3$ & $F_4$ & $\langle\alpha_1, \alpha_2\rangle=2$,  $\langle\alpha_1, \beta_1\rangle=2$, $\langle\alpha_1, \beta_2\rangle=2$, $\langle\alpha_2, \beta_1\rangle=2$, $\langle\alpha_2, \beta_2\rangle=2$, 
\\
\begin{tabular}{cc}$\alpha_1$+ & $\alpha_2$\\\tiny{ $A_1'$+$A_1$ } & \tiny{  }\end{tabular} &\begin{tabular}{c} = \\~ \end{tabular} & \begin{tabular}{cc}$\beta_1$+ & 2$\beta_2$\\\tiny{ $A_1$ } & \tiny{ $A_1'$ }\end{tabular} & $B_3$ & $F_4$ & $\langle\alpha_1, \alpha_2\rangle=2$,  $\langle\alpha_1, \beta_1\rangle=2$, $\langle\alpha_1, \beta_2\rangle=2$, $\langle\alpha_2, \beta_1\rangle=2$, $\langle\alpha_2, \beta_2\rangle=2$, 
\\
\begin{tabular}{cc}$\alpha_1$+ & $\alpha_2$\\\tiny{ $A_1'$ } & \tiny{  }\end{tabular} &\begin{tabular}{c} = \\~ \end{tabular} & \begin{tabular}{c}2$\beta_1$\\\tiny{ $A_1'$ }\end{tabular} & $B_2=C_2$ & $C_3$+$A_1$ &  $\langle\alpha_1, \beta_1\rangle=2$, $\langle\alpha_2, \beta_1\rangle=2$, 
\\
\begin{tabular}{cc}$\alpha_1$+ & $\alpha_2$\\\tiny{ $A_1$ } & \tiny{ $A_1'$ }\end{tabular} &\begin{tabular}{c} = \\~ \end{tabular} & \begin{tabular}{cc}$\beta_1$+ & 2$\beta_2$\\\tiny{ $A_1$ } & \tiny{ $A_1'$ }\end{tabular} & $B_3$ & $F_4$ & $\langle\alpha_1, \alpha_2\rangle=2$,  $\langle\alpha_1, \beta_1\rangle=2$, $\langle\alpha_1, \beta_2\rangle=2$, $\langle\alpha_2, \beta_1\rangle=2$, $\langle\alpha_2, \beta_2\rangle=2$, 
\\
\begin{tabular}{cc}$\alpha_1$+ & $\alpha_2$\\\tiny{ $A_1$ } & \tiny{ $A_1'$ }\end{tabular} &\begin{tabular}{c} = \\~ \end{tabular} & \begin{tabular}{c}$\beta_1$\\\tiny{ $A_1'$+$A_1$ }\end{tabular} & $B_2=C_2$ & $F_4$ & $\langle\alpha_1, \alpha_2\rangle=-2$,  $\langle\alpha_1, \beta_1\rangle=2$, 
\\
\begin{tabular}{c}2$\alpha_1$\\\tiny{ $A_1'$+$A_1$ }\end{tabular} &\begin{tabular}{c} = \\~ \end{tabular} & \begin{tabular}{cc}$\beta_1$+ & $\beta_2$\\\tiny{ $A_1'$+$A_1$ } & \tiny{ $A_1$ }\end{tabular} & $B_2=C_2$ & $F_4$ &  $\langle\alpha_1, \beta_1\rangle=2$, $\langle\alpha_1, \beta_2\rangle=2$, 
\\
\begin{tabular}{c}2$\alpha_1$\\\tiny{ $A_1'$ }\end{tabular} &\begin{tabular}{c} = \\~ \end{tabular} & \begin{tabular}{cc}$\beta_1$+ & $\beta_2$\\\tiny{ $A_1'$+$A_1$ } & \tiny{ $A_1$ }\end{tabular} & $B_2=C_2$ & $F_4$ &  $\langle\alpha_1, \beta_1\rangle=2$, $\langle\alpha_1, \beta_2\rangle=2$, 
\\
\begin{tabular}{cc}$\alpha_1$+ & $\alpha_2$\\\tiny{ $A_1'$+$A_1$ } & \tiny{ $A_1'$ }\end{tabular} &\begin{tabular}{c} = \\~ \end{tabular} & \begin{tabular}{c}$\beta_1$\\\tiny{ $A_1'$+$A_1$ }\end{tabular} & $B_2=C_2$ & $F_4$ &  $\langle\alpha_1, \beta_1\rangle=2$, $\langle\alpha_2, \beta_1\rangle=2$, 
\\
\begin{tabular}{c}2$\alpha_1$\\\tiny{ $A_1'$ }\end{tabular} &\begin{tabular}{c} = \\~ \end{tabular} & \begin{tabular}{cc}$\beta_1$+ & $\beta_2$\\\tiny{ $A_1'$ } & \tiny{  }\end{tabular} & $B_2=C_2$ & $C_3$+$A_1$ &  $\langle\alpha_1, \beta_1\rangle=2$, $\langle\alpha_1, \beta_2\rangle=2$, 
\\
\begin{tabular}{c}2$\alpha_1$\\\tiny{ $A_1$ }\end{tabular} &\begin{tabular}{c} = \\~ \end{tabular} & \begin{tabular}{ccc}2$\beta_1$+ & $\beta_2$+ & $\beta_3$\\\tiny{ $A_1'$+$A_1$ } & \tiny{ $A_1'$ } & \tiny{  }\end{tabular} & $B_3$ & $F_4$ &  $\langle\alpha_1, \beta_1\rangle=2$, $\langle\alpha_1, \beta_2\rangle=2$, $\langle\alpha_1, \beta_3\rangle=2$, 
\\
\begin{tabular}{cc}$\alpha_1$+ & $\alpha_2$\\\tiny{ $A_1$ } & \tiny{ $A_1'$ }\end{tabular} &\begin{tabular}{c} = \\~ \end{tabular} & \begin{tabular}{c}$\beta_1$\\\tiny{ $A_1'$+$A_1$ }\end{tabular} & $A_2$ & $B_4$ & $\langle\alpha_1, \alpha_2\rangle=-2$,  $\langle\alpha_1, \beta_1\rangle=2$, $\langle\alpha_2, \beta_1\rangle=2$, 
\\
\begin{tabular}{cc}$\alpha_1$+ & $\alpha_2$\\\tiny{ $A_1$ } & \tiny{ $A_1'$ }\end{tabular} &\begin{tabular}{c} = \\~ \end{tabular} & \begin{tabular}{cc}$\beta_1$+ & $\beta_2$\\\tiny{ $A_1'$+$A_1$ } & \tiny{ $A_1'$ }\end{tabular} & $B_3$ & $F_4$ &  $\langle\alpha_1, \beta_1\rangle=2$, $\langle\alpha_1, \beta_2\rangle=2$, $\langle\alpha_2, \beta_1\rangle=2$, 
\\
\begin{tabular}{cc}$\alpha_1$+ & $\alpha_2$\\\tiny{ $A_1$ } & \tiny{ $A_1$ }\end{tabular} &\begin{tabular}{c} = \\~ \end{tabular} & \begin{tabular}{cc}$\beta_1$+ & $\beta_2$\\\tiny{ $A_1'$+$A_1$ } & \tiny{ $A_1'$+$A_1$ }\end{tabular} & $A_3$ & $B_4$ &  $\langle\alpha_1, \beta_1\rangle=2$, $\langle\alpha_1, \beta_2\rangle=2$, $\langle\alpha_2, \beta_1\rangle=2$, $\langle\alpha_2, \beta_2\rangle=2$, 
\\
\begin{tabular}{cc}$\alpha_1$+ & 2$\alpha_2$\\\tiny{ $A_1$ } & \tiny{ $A_1'$ }\end{tabular} &\begin{tabular}{c} = \\~ \end{tabular} & \begin{tabular}{cc}$\beta_1$+ & 2$\beta_2$\\\tiny{ $A_1$ } & \tiny{ $A_1'$ }\end{tabular} & $C_3$ & $F_4$ &  $\langle\alpha_1, \beta_2\rangle=2$, $\langle\alpha_2, \beta_1\rangle=2$, $\langle\alpha_2, \beta_2\rangle=1$, 
\\
\hline\multicolumn{5}{c}{$\mathfrak{k}$-semisimple type: $A_1$+$A_1$}\\
\hline\hline\begin{tabular}{cc}$\alpha_1$+ & $\alpha_2$\\\tiny{ $A_1$ } & \tiny{ $A_1$' }\end{tabular} &\begin{tabular}{c} = \\~ \end{tabular} & \begin{tabular}{c}$\beta_1$\\\tiny{ $A_1$+$A_1$' }\end{tabular} & $B_2=C_2$ & $B_4$ &  $\langle\alpha_1, \beta_1\rangle=2$, $\langle\alpha_2, \beta_1\rangle=2$, 
\\
\begin{tabular}{c}2$\alpha_1$\\\tiny{  }\end{tabular} &\begin{tabular}{c} = \\~ \end{tabular} & \begin{tabular}{cc}$\beta_1$+ & $\beta_2$\\\tiny{ $A_1$+$A_1$' } & \tiny{ $A_1$+$A_1$' }\end{tabular} & $B_2=C_2$ & $B_3$ &  $\langle\alpha_1, \beta_1\rangle=2$, $\langle\alpha_1, \beta_2\rangle=2$, 
\\
\begin{tabular}{cc}$\alpha_1$+ & $\alpha_2$\\\tiny{ $A_1$ } & \tiny{  }\end{tabular} &\begin{tabular}{c} = \\~ \end{tabular} & \begin{tabular}{c}$\beta_1$\\\tiny{ $A_1$ }\end{tabular} & $A_2'$ & $C_3$+$A_1$ & $\langle\alpha_1, \alpha_2\rangle=-1$,  $\langle\alpha_1, \beta_1\rangle=1$, $\langle\alpha_2, \beta_1\rangle=1$, 
\\
\begin{tabular}{c}2$\alpha_1$\\\tiny{ $A_1$ }\end{tabular} &\begin{tabular}{c} = \\~ \end{tabular} & \begin{tabular}{cc}$\beta_1$+ & $\beta_2$\\\tiny{ $A_1$+$A_1$' } & \tiny{ $A_1$+$A_1$' }\end{tabular} & $B_2=C_2$ & $B_4$ &  $\langle\alpha_1, \beta_1\rangle=2$, $\langle\alpha_1, \beta_2\rangle=2$, 
\\
\begin{tabular}{cc}$\alpha_1$+ & $\alpha_2$\\\tiny{ $A_1$+$A_1$' } & \tiny{  }\end{tabular} &\begin{tabular}{c} = \\~ \end{tabular} & \begin{tabular}{c}$\beta_1$\\\tiny{ $A_1$+$A_1$' }\end{tabular} & $A_2$ & $D_4$ & $\langle\alpha_1, \alpha_2\rangle=-2$,  $\langle\alpha_1, \beta_1\rangle=2$, $\langle\alpha_2, \beta_1\rangle=2$, 
\\
\begin{tabular}{cc}$\alpha_1$+ & $\alpha_2$\\\tiny{ $A_1$ } & \tiny{ $A_1$' }\end{tabular} &\begin{tabular}{c} = \\~ \end{tabular} & \begin{tabular}{cc}$\beta_1$+ & $\beta_2$\\\tiny{ $A_1$+$A_1$' } & \tiny{  }\end{tabular} & $C_3$ & $F_4$ & $\langle\alpha_1, \alpha_2\rangle=1$,  $\langle\alpha_1, \beta_1\rangle=2$, $\langle\alpha_1, \beta_2\rangle=1$, $\langle\alpha_2, \beta_1\rangle=2$, $\langle\alpha_2, \beta_2\rangle=1$, 
\\
\hline\multicolumn{5}{c}{$\mathfrak{k}$-semisimple type: $A_1'$}\\
\hline\hline\begin{tabular}{cc}$\alpha_1$+ & $\alpha_2$\\\tiny{  } & \tiny{  }\end{tabular} &\begin{tabular}{c} = \\~ \end{tabular} & \begin{tabular}{cc}$\beta_1$+ & $\beta_2$\\\tiny{ $A_1'$ } & \tiny{ $A_1'$ }\end{tabular} & $A_3$ & $B_3$ &  $\langle\alpha_1, \beta_1\rangle=2$, $\langle\alpha_1, \beta_2\rangle=2$, $\langle\alpha_2, \beta_1\rangle=2$, $\langle\alpha_2, \beta_2\rangle=2$, 
\\
\begin{tabular}{cc}$\alpha_1$+ & $\alpha_2$\\\tiny{ $A_1'$ } & \tiny{  }\end{tabular} &\begin{tabular}{c} = \\~ \end{tabular} & \begin{tabular}{c}$\beta_1$\\\tiny{ $A_1'$ }\end{tabular} & $B_2=C_2$ & $C_3$ & $\langle\alpha_1, \alpha_2\rangle=-2$,  $\langle\alpha_2, \beta_1\rangle=2$, 
\\
\begin{tabular}{cc}$\alpha_1$+ & $\alpha_2$\\\tiny{ $A_1'$ } & \tiny{  }\end{tabular} &\begin{tabular}{c} = \\~ \end{tabular} & \begin{tabular}{c}2$\beta_1$\\\tiny{ $A_1'$ }\end{tabular} & $B_2=C_2$ & $C_3$ &  $\langle\alpha_1, \beta_1\rangle=2$, $\langle\alpha_2, \beta_1\rangle=2$, 
\\
\begin{tabular}{cc}$\alpha_1$+ & $\alpha_2$\\\tiny{ $A_1'$ } & \tiny{ $A_1'$ }\end{tabular} &\begin{tabular}{c} = \\~ \end{tabular} & \begin{tabular}{c}$\beta_1$\\\tiny{ $A_1'$ }\end{tabular} & $B_2=C_2$ & $C_3$ &  $\langle\alpha_1, \beta_1\rangle=2$, $\langle\alpha_2, \beta_1\rangle=2$, 
\\
\begin{tabular}{cc}$\alpha_1$+ & $\alpha_2$\\\tiny{ $A_1'$ } & \tiny{  }\end{tabular} &\begin{tabular}{c} = \\~ \end{tabular} & \begin{tabular}{c}$\beta_1$\\\tiny{ $A_1'$ }\end{tabular} & $A_2$ & $B_3$ & $\langle\alpha_1, \alpha_2\rangle=-2$,  $\langle\alpha_1, \beta_1\rangle=2$, $\langle\alpha_2, \beta_1\rangle=2$, 
\\
\begin{tabular}{cc}$\alpha_1$+ & $\alpha_2$\\\tiny{ $A_1'$ } & \tiny{ $A_1'$ }\end{tabular} &\begin{tabular}{c} = \\~ \end{tabular} & \begin{tabular}{cc}$\beta_1$+ & $\beta_2$\\\tiny{ $A_1'$ } & \tiny{ $A_1'$ }\end{tabular} & $B_3$ & $F_4$ &  $\langle\alpha_1, \beta_1\rangle=2$, $\langle\alpha_2, \beta_1\rangle=2$, $\langle\alpha_2, \beta_2\rangle=2$, 
\\
\begin{tabular}{cc}$\alpha_1$+ & $\alpha_2$\\\tiny{ $A_1'$ } & \tiny{ $A_1'$ }\end{tabular} &\begin{tabular}{c} = \\~ \end{tabular} & \begin{tabular}{cc}$\beta_1$+ & 2$\beta_2$\\\tiny{ $A_1'$ } & \tiny{ $A_1'$ }\end{tabular} & $B_3$ & $F_4$ & $\langle\alpha_1, \alpha_2\rangle=2$,  $\langle\alpha_1, \beta_1\rangle=2$, $\langle\alpha_1, \beta_2\rangle=2$, $\langle\alpha_2, \beta_1\rangle=2$, $\langle\alpha_2, \beta_2\rangle=2$, 
\\
\begin{tabular}{cc}$\alpha_1$+ & $\alpha_2$\\\tiny{  } & \tiny{  }\end{tabular} &\begin{tabular}{c} = \\~ \end{tabular} & \begin{tabular}{cc}$\beta_1$+ & 2$\beta_2$\\\tiny{ $A_1'$ } & \tiny{ $A_1'$ }\end{tabular} & $B_3$ & $F_4$ & $\langle\alpha_1, \alpha_2\rangle=2$,  $\langle\alpha_1, \beta_1\rangle=2$, $\langle\alpha_1, \beta_2\rangle=2$, $\langle\alpha_2, \beta_1\rangle=2$, $\langle\alpha_2, \beta_2\rangle=2$, 
\\
\begin{tabular}{c}2$\alpha_1$\\\tiny{ $A_1'$ }\end{tabular} &\begin{tabular}{c} = \\~ \end{tabular} & \begin{tabular}{cc}$\beta_1$+ & $\beta_2$\\\tiny{ $A_1'$ } & \tiny{  }\end{tabular} & $B_2=C_2$ & $C_3$ &  $\langle\alpha_1, \beta_1\rangle=2$, $\langle\alpha_1, \beta_2\rangle=2$, 
\\
\begin{tabular}{cc}$\alpha_1$+ & $\alpha_2$\\\tiny{ $A_1'$ } & \tiny{ $A_1'$ }\end{tabular} &\begin{tabular}{c} = \\~ \end{tabular} & \begin{tabular}{cc}$\beta_1$+ & $\beta_2$\\\tiny{ $A_1'$ } & \tiny{  }\end{tabular} & $C_3$ & $F_4$ & $\langle\alpha_1, \alpha_2\rangle=1$,  $\langle\alpha_1, \beta_1\rangle=2$, $\langle\alpha_1, \beta_2\rangle=1$, $\langle\alpha_2, \beta_1\rangle=2$, $\langle\alpha_2, \beta_2\rangle=1$, 
\\
\hline\multicolumn{5}{c}{$\mathfrak{k}$-semisimple type: $A_1$}\\
\hline\hline\begin{tabular}{cc}$\alpha_1$+ & $\alpha_2$\\\tiny{ $A_1$ } & \tiny{  }\end{tabular} &\begin{tabular}{c} = \\~ \end{tabular} & \begin{tabular}{c}$\beta_1$\\\tiny{ $A_1$ }\end{tabular} & $B_2=C_2$ & $B_3$ &  $\langle\alpha_1, \beta_1\rangle=2$, $\langle\alpha_2, \beta_1\rangle=2$, 
\\
\begin{tabular}{cc}$\alpha_1$+ & $\alpha_2$\\\tiny{ $A_1$ } & \tiny{  }\end{tabular} &\begin{tabular}{c} = \\~ \end{tabular} & \begin{tabular}{c}$\beta_1$\\\tiny{ $A_1$ }\end{tabular} & $A_2$ & $A_3$ & $\langle\alpha_1, \alpha_2\rangle=-2$,  $\langle\alpha_1, \beta_1\rangle=2$, $\langle\alpha_2, \beta_1\rangle=2$, 
\\
\begin{tabular}{c}2$\alpha_1$\\\tiny{ $A_1$ }\end{tabular} &\begin{tabular}{c} = \\~ \end{tabular} & \begin{tabular}{cc}$\beta_1$+ & $\beta_2$\\\tiny{ $A_1$ } & \tiny{ $A_1$ }\end{tabular} & $B_2=C_2$ & $B_3$ &  $\langle\alpha_1, \beta_1\rangle=2$, $\langle\alpha_1, \beta_2\rangle=2$, 
\\
\begin{tabular}{cc}$\alpha_1$+ & $\alpha_2$\\\tiny{ $A_1$ } & \tiny{ $A_1$ }\end{tabular} &\begin{tabular}{c} = \\~ \end{tabular} & \begin{tabular}{cc}$\beta_1$+ & $\beta_2$\\\tiny{ $A_1$ } & \tiny{ $A_1$ }\end{tabular} & $A_3$ & $D_4$ &  $\langle\alpha_1, \beta_1\rangle=2$, $\langle\alpha_1, \beta_2\rangle=2$, $\langle\alpha_2, \beta_1\rangle=2$, $\langle\alpha_2, \beta_2\rangle=2$, 
\\
\begin{tabular}{cc}$\alpha_1$+ & $\alpha_2$\\\tiny{ $A_1$ } & \tiny{  }\end{tabular} &\begin{tabular}{c} = \\~ \end{tabular} & \begin{tabular}{c}$\beta_1$\\\tiny{ $A_1$ }\end{tabular} & $B_2=C_2$ & $B_3$ & $\langle\alpha_1, \alpha_2\rangle=-2$,  $\langle\alpha_1, \beta_1\rangle=2$, 
\\
\begin{tabular}{cc}$\alpha_1$+ & $\alpha_2$\\\tiny{ $A_1$ } & \tiny{ $A_1$ }\end{tabular} &\begin{tabular}{c} = \\~ \end{tabular} & \begin{tabular}{c}2$\beta_1$\\\tiny{ $A_1$ }\end{tabular} & $B_2=C_2$ & $B_3$ &  $\langle\alpha_1, \beta_1\rangle=2$, $\langle\alpha_2, \beta_1\rangle=2$, 
\\
\begin{tabular}{c}2$\alpha_1$\\\tiny{  }\end{tabular} &\begin{tabular}{c} = \\~ \end{tabular} & \begin{tabular}{cc}$\beta_1$+ & $\beta_2$\\\tiny{ $A_1$ } & \tiny{ $A_1$ }\end{tabular} & $B_2=C_2$ & $B_3$ &  $\langle\alpha_1, \beta_1\rangle=2$, $\langle\alpha_1, \beta_2\rangle=2$, 
\\
\begin{tabular}{cc}$\alpha_1$+ & $\alpha_2$\\\tiny{ $A_1$ } & \tiny{ $A_1$ }\end{tabular} &\begin{tabular}{c} = \\~ \end{tabular} & \begin{tabular}{cc}$\beta_1$+ & $\beta_2$\\\tiny{ $A_1$ } & \tiny{ $A_1$ }\end{tabular} & $C_3$ & $F_4$ & $\langle\alpha_1, \alpha_2\rangle=1$,  $\langle\alpha_1, \beta_1\rangle=2$, $\langle\alpha_1, \beta_2\rangle=1$, $\langle\alpha_2, \beta_1\rangle=2$, $\langle\alpha_2, \beta_2\rangle=1$, 
\\
\begin{tabular}{cc}$\alpha_1$+ & $\alpha_2$\\\tiny{ $A_1$ } & \tiny{  }\end{tabular} &\begin{tabular}{c} = \\~ \end{tabular} & \begin{tabular}{c}$\beta_1$\\\tiny{ $A_1$ }\end{tabular} & $A_2'$ & $C_3$ & $\langle\alpha_1, \alpha_2\rangle=-1$,  $\langle\alpha_1, \beta_1\rangle=1$, $\langle\alpha_2, \beta_1\rangle=1$, 
\\
\end{longtable}

\newpage
\subsubsection[$F_4$ two-sided weights without $\sperp$ decomposition]{$F_4$ two-sided weights without a strongly orthogonal decomposition}
This section lists the only case up to $F_4$-automorphism in $F_4$ for which no two-sided weight with strongly orthogonal decomposition exists. The second table gives one $\lL$-(non-strictly) infinite weight in this case.

\noindent\begin{tabular}{rcl p{1cm}p{1cm}p{3cm} } \multicolumn{3}{c}{ Relation / linked $\mathfrak{k}$-components} &$\alpha_i$'s, $\beta_i$'s generate & adding $\mathfrak{k}$ generates&Non-zero scalar products
\\\hline\multicolumn{5}{c}{$\mathfrak{k}$-semisimple type: $A_1$+$A_1$}\\
\hline\hline\begin{tabular}{cc}$\alpha_1$+ & $\alpha_2$\\\tiny{ $A_1$+$A_1$' } & \tiny{  }\end{tabular} &\begin{tabular}{c} = \\~ \end{tabular} & \begin{tabular}{cc}$\beta_1$+ & $\beta_2$\\\tiny{ $A_1$ } & \tiny{ $A_1$' }\end{tabular} & $C_3$ & $F_4$ &  $\langle\beta_1,\beta_2\rangle=1$, $\langle\alpha_1,\beta_1\rangle=2$, $\langle\alpha_1,\beta_2\rangle=2$, $\langle\alpha_2,\beta_1\rangle=1$, $\langle\alpha_2,\beta_2\rangle=1$, 
\\
\multicolumn{5}{c}{$\varepsilon$-form~relative~to~the~subalgebra~generated~by~$\mathfrak{k}$~and~the~relation}\\
\multicolumn{5}{c}{(2$\varepsilon_{1}$) + (+$\varepsilon_{2}$+$\varepsilon_{3}$)=($\varepsilon_{1}$+$\varepsilon_{3}$) + ($\varepsilon_{1}$+$\varepsilon_{2}$)}\\\hline
\end{tabular}

\medskip ~

\medskip ~

\medskip ~

\medskip ~

\noindent Corresponding $\lL$-(non-strongly) infinite weight.

\medskip ~

\medskip ~

\medskip ~

\noindent\begin{tabular}{rcl p{1cm}p{1cm}p{3cm} } \multicolumn{3}{c}{ Relation / linked $\mathfrak{k}$-components} &$\alpha_i$'s, $\beta_i$'s generate & adding $\mathfrak{k}$ generates&Non-zero scalar products
\\\hline\multicolumn{5}{c}{$\mathfrak{k}$-semisimple type: $A_1$+$A_1$}\\
\hline\hline\begin{tabular}{cc}$\alpha_1$+ & $\alpha_2$\\\tiny{ $A_1$ } & \tiny{  }\end{tabular} &\begin{tabular}{c} = \\~ \end{tabular} & \begin{tabular}{c}$\beta_1$\\\tiny{ $A_1$ }\end{tabular} & $C_3$+$A_1$ & $C_3$+$A_1$ & $\langle\alpha_1,\alpha_2\rangle=-1$, $\langle\alpha_1,\beta_1\rangle=1$, $\langle\alpha_2,\beta_1\rangle=1$, 
\\
\multicolumn{5}{c}{$\varepsilon$-form~relative~to~the~subalgebra~generated~by~$\mathfrak{k}$~and~the~relation}\\
\multicolumn{5}{c}{($\varepsilon_{1}$-$\varepsilon_{2}$) + (+$\varepsilon_{2}$+$\varepsilon_{3}$)=($\varepsilon_{1}$+$\varepsilon_{3}$)}\\\hline
\end{tabular}



\normalsize
\subsection{$E_6$: $\lL$-strictly infinite weights and corresponding relations }
\noindent Number of different non-solvable subalgebras up to $\gG$-automorphism such that $\nn\cap C(\kk_{ss})$ is a nilradical of a parabolic subalgebra of $\kk$ containing $C(\kk_{ss})\cap \hh$: 2044


\noindent Among them 706 satisfy the cone condition and 1338 do not.

\tiny

\begin{longtable}{rcl p{1cm}p{1cm}p{3cm} } \multicolumn{3}{c}{ Relation / linked $\mathfrak{k}$-components} &$\alpha_i$'s, $\beta_i$'s generate & adding $\Delta(\mathfrak{k})$ generates&Non-zero scalar products\\\hline\multicolumn{5}{c}{$\mathfrak{k}$-semisimple type: $A_4$}\\
\hline\hline\begin{tabular}{cc}$\alpha_1$+ & $\alpha_2$\\\tiny{ $A_4$ } & \tiny{  }\end{tabular} &\begin{tabular}{c} = \\~ \end{tabular} & \begin{tabular}{c}$\beta_1$\\\tiny{ $A_4$ }\end{tabular} & $A_2$ & $E_6$ & $\langle\alpha_1, \alpha_2\rangle=-1$,  $\langle\alpha_1, \beta_1\rangle=1$, $\langle\alpha_2, \beta_1\rangle=1$, 
\\
\begin{tabular}{cc}$\alpha_1$+ & $\alpha_2$\\\tiny{ $A_4$ } & \tiny{  }\end{tabular} &\begin{tabular}{c} = \\~ \end{tabular} & \begin{tabular}{cc}$\beta_1$+ & $\beta_2$\\\tiny{ $A_4$ } & \tiny{ $A_4$ }\end{tabular} & $A_3$ & $E_6$ &  $\langle\alpha_1, \beta_1\rangle=1$, $\langle\alpha_1, \beta_2\rangle=1$, $\langle\alpha_2, \beta_1\rangle=1$, $\langle\alpha_2, \beta_2\rangle=1$, 
\\
\hline\multicolumn{5}{c}{$\mathfrak{k}$-semisimple type: $A_3$+$A_1$}\\
\hline\hline\begin{tabular}{cc}$\alpha_1$+ & $\alpha_2$\\\tiny{ $A_3$ } & \tiny{  }\end{tabular} &\begin{tabular}{c} = \\~ \end{tabular} & \begin{tabular}{c}$\beta_1$\\\tiny{ $A_3$ }\end{tabular} & $A_2$ & $A_5$+$A_1$ & $\langle\alpha_1, \alpha_2\rangle=-1$,  $\langle\alpha_1, \beta_1\rangle=1$, $\langle\alpha_2, \beta_1\rangle=1$, 
\\
\begin{tabular}{cc}$\alpha_1$+ & $\alpha_2$\\\tiny{ $A_3$ } & \tiny{ $A_3$ }\end{tabular} &\begin{tabular}{c} = \\~ \end{tabular} & \begin{tabular}{cc}$\beta_1$+ & $\beta_2$\\\tiny{ $A_3$+$A_1$ } & \tiny{ $A_3$+$A_1$ }\end{tabular} & $A_3$ & $E_6$ &  $\langle\alpha_1, \beta_1\rangle=1$, $\langle\alpha_1, \beta_2\rangle=1$, $\langle\alpha_2, \beta_1\rangle=1$, $\langle\alpha_2, \beta_2\rangle=1$, 
\\
\hline\multicolumn{5}{c}{$\mathfrak{k}$-semisimple type: $A_2$+$A_2$}\\
\hline\hline\begin{tabular}{cc}$\alpha_1$+ & $\alpha_2$\\\tiny{  } & \tiny{  }\end{tabular} &\begin{tabular}{c} = \\~ \end{tabular} & \begin{tabular}{ccc}$\beta_1$+ & $\beta_2$+ & $\beta_3$\\\tiny{ $A_2$+$A_2$' } & \tiny{ $A_2$+$A_2$' } & \tiny{ $A_2$+$A_2$' }\end{tabular} & $D_4$ & $E_6$ & $\langle\alpha_1, \alpha_2\rangle=1$,  $\langle\alpha_1, \beta_1\rangle=1$, $\langle\alpha_1, \beta_2\rangle=1$, $\langle\alpha_1, \beta_3\rangle=1$, $\langle\alpha_2, \beta_1\rangle=1$, $\langle\alpha_2, \beta_2\rangle=1$, $\langle\alpha_2, \beta_3\rangle=1$, 
\\
\begin{tabular}{cc}$\alpha_1$+ & $\alpha_2$\\\tiny{ $A_2$+$A_2$' } & \tiny{  }\end{tabular} &\begin{tabular}{c} = \\~ \end{tabular} & \begin{tabular}{c}$\beta_1$\\\tiny{ $A_2$+$A_2$' }\end{tabular} & $A_2$ & $E_6$ & $\langle\alpha_1, \alpha_2\rangle=-1$,  $\langle\alpha_1, \beta_1\rangle=1$, $\langle\alpha_2, \beta_1\rangle=1$, 
\\
\begin{tabular}{cc}$\alpha_1$+ & $\alpha_2$\\\tiny{ $A_2$+$A_2$' } & \tiny{  }\end{tabular} &\begin{tabular}{c} = \\~ \end{tabular} & \begin{tabular}{cc}$\beta_1$+ & $\beta_2$\\\tiny{ $A_2$+$A_2$' } & \tiny{ $A_2$+$A_2$' }\end{tabular} & $A_3$ & $E_6$ &  $\langle\alpha_1, \beta_1\rangle=1$, $\langle\alpha_1, \beta_2\rangle=1$, $\langle\alpha_2, \beta_1\rangle=1$, $\langle\alpha_2, \beta_2\rangle=1$, 
\\
\hline\multicolumn{5}{c}{$\mathfrak{k}$-semisimple type: $A_2$+$A_1$+$A_1$}\\
\hline\hline\begin{tabular}{cc}$\alpha_1$+ & $\alpha_2$\\\tiny{ $A_2$+$A_1$ } & \tiny{ $A_1$' }\end{tabular} &\begin{tabular}{c} = \\~ \end{tabular} & \begin{tabular}{c}$\beta_1$\\\tiny{ $A_2$+$A_1$+$A_1$' }\end{tabular} & $A_2$ & $E_6$ & $\langle\alpha_1, \alpha_2\rangle=-1$,  $\langle\alpha_1, \beta_1\rangle=1$, $\langle\alpha_2, \beta_1\rangle=1$, 
\\
\begin{tabular}{cc}$\alpha_1$+ & $\alpha_2$\\\tiny{ $A_2$+$A_1$ } & \tiny{ $A_1$ }\end{tabular} &\begin{tabular}{c} = \\~ \end{tabular} & \begin{tabular}{cc}$\beta_1$+ & $\beta_2$\\\tiny{ $A_2$+$A_1$+$A_1$' } & \tiny{ $A_2$+$A_1$+$A_1$' }\end{tabular} & $A_3$ & $E_6$ &  $\langle\alpha_1, \beta_1\rangle=1$, $\langle\alpha_1, \beta_2\rangle=1$, $\langle\alpha_2, \beta_1\rangle=1$, $\langle\alpha_2, \beta_2\rangle=1$, 
\\
\begin{tabular}{cc}$\alpha_1$+ & $\alpha_2$\\\tiny{ $A_1$ } & \tiny{ $A_2$ }\end{tabular} &\begin{tabular}{c} = \\~ \end{tabular} & \begin{tabular}{c}$\beta_1$\\\tiny{ $A_2$+$A_1$ }\end{tabular} & $A_2$ & $A_5$+$A_1$ & $\langle\alpha_1, \alpha_2\rangle=-1$,  $\langle\alpha_1, \beta_1\rangle=1$, $\langle\alpha_2, \beta_1\rangle=1$, 
\\
\begin{tabular}{cc}$\alpha_1$+ & $\alpha_2$\\\tiny{ $A_1$ } & \tiny{ $A_1$' }\end{tabular} &\begin{tabular}{c} = \\~ \end{tabular} & \begin{tabular}{cc}$\beta_1$+ & $\beta_2$\\\tiny{ $A_2$+$A_1$ } & \tiny{ $A_2$+$A_1$' }\end{tabular} & $A_3$ & $E_6$ &  $\langle\alpha_1, \beta_1\rangle=1$, $\langle\alpha_1, \beta_2\rangle=1$, $\langle\alpha_2, \beta_1\rangle=1$, $\langle\alpha_2, \beta_2\rangle=1$, 
\\
\begin{tabular}{cc}$\alpha_1$+ & $\alpha_2$\\\tiny{ $A_1$ } & \tiny{ $A_1$' }\end{tabular} &\begin{tabular}{c} = \\~ \end{tabular} & \begin{tabular}{cc}$\beta_1$+ & $\beta_2$\\\tiny{ $A_2$+$A_1$+$A_1$' } & \tiny{ $A_2$ }\end{tabular} & $A_3$ & $E_6$ &  $\langle\alpha_1, \beta_1\rangle=1$, $\langle\alpha_1, \beta_2\rangle=1$, $\langle\alpha_2, \beta_1\rangle=1$, $\langle\alpha_2, \beta_2\rangle=1$, 
\\
\hline\multicolumn{5}{c}{$\mathfrak{k}$-semisimple type: $A_3$}\\
\hline\hline\begin{tabular}{cc}$\alpha_1$+ & $\alpha_2$\\\tiny{  } & \tiny{  }\end{tabular} &\begin{tabular}{c} = \\~ \end{tabular} & \begin{tabular}{cc}$\beta_1$+ & $\beta_2$\\\tiny{ $A_3$ } & \tiny{ $A_3$ }\end{tabular} & $A_3$ & $D_5$ &  $\langle\alpha_1, \beta_1\rangle=1$, $\langle\alpha_1, \beta_2\rangle=1$, $\langle\alpha_2, \beta_1\rangle=1$, $\langle\alpha_2, \beta_2\rangle=1$, 
\\
\begin{tabular}{cc}$\alpha_1$+ & $\alpha_2$\\\tiny{ $A_3$ } & \tiny{  }\end{tabular} &\begin{tabular}{c} = \\~ \end{tabular} & \begin{tabular}{c}$\beta_1$\\\tiny{ $A_3$ }\end{tabular} & $A_2$ & $A_5$ & $\langle\alpha_1, \alpha_2\rangle=-1$,  $\langle\alpha_1, \beta_1\rangle=1$, $\langle\alpha_2, \beta_1\rangle=1$, 
\\
\begin{tabular}{cc}$\alpha_1$+ & $\alpha_2$\\\tiny{ $A_3$ } & \tiny{  }\end{tabular} &\begin{tabular}{c} = \\~ \end{tabular} & \begin{tabular}{cc}$\beta_1$+ & $\beta_2$\\\tiny{ $A_3$ } & \tiny{ $A_3$ }\end{tabular} & $A_3$ & $E_6$ &  $\langle\alpha_1, \beta_1\rangle=1$, $\langle\alpha_1, \beta_2\rangle=1$, $\langle\alpha_2, \beta_1\rangle=1$, $\langle\alpha_2, \beta_2\rangle=1$, 
\\
\begin{tabular}{cc}$\alpha_1$+ & $\alpha_2$\\\tiny{ $A_3$ } & \tiny{  }\end{tabular} &\begin{tabular}{c} = \\~ \end{tabular} & \begin{tabular}{c}$\beta_1$\\\tiny{ $A_3$ }\end{tabular} & $A_2$ & $D_5$ & $\langle\alpha_1, \alpha_2\rangle=-1$,  $\langle\alpha_1, \beta_1\rangle=1$, $\langle\alpha_2, \beta_1\rangle=1$, 
\\
\begin{tabular}{cc}$\alpha_1$+ & $\alpha_2$\\\tiny{ $A_3$ } & \tiny{ $A_3$ }\end{tabular} &\begin{tabular}{c} = \\~ \end{tabular} & \begin{tabular}{cc}$\beta_1$+ & $\beta_2$\\\tiny{ $A_3$ } & \tiny{ $A_3$ }\end{tabular} & $A_3$ & $E_6$ &  $\langle\alpha_1, \beta_1\rangle=1$, $\langle\alpha_1, \beta_2\rangle=1$, $\langle\alpha_2, \beta_1\rangle=1$, $\langle\alpha_2, \beta_2\rangle=1$, 
\\
\hline\multicolumn{5}{c}{$\mathfrak{k}$-semisimple type: $A_2$+$A_1$}\\
\hline\hline\begin{tabular}{cc}$\alpha_1$+ & $\alpha_2$\\\tiny{ $A_2$+$A_1$ } & \tiny{  }\end{tabular} &\begin{tabular}{c} = \\~ \end{tabular} & \begin{tabular}{c}$\beta_1$\\\tiny{ $A_2$+$A_1$ }\end{tabular} & $A_2$ & $D_5$ & $\langle\alpha_1, \alpha_2\rangle=-1$,  $\langle\alpha_1, \beta_1\rangle=1$, $\langle\alpha_2, \beta_1\rangle=1$, 
\\
\begin{tabular}{cc}$\alpha_1$+ & $\alpha_2$\\\tiny{ $A_2$ } & \tiny{  }\end{tabular} &\begin{tabular}{c} = \\~ \end{tabular} & \begin{tabular}{cc}$\beta_1$+ & $\beta_2$\\\tiny{ $A_2$+$A_1$ } & \tiny{ $A_2$+$A_1$ }\end{tabular} & $A_3$ & $D_5$ &  $\langle\alpha_1, \beta_1\rangle=1$, $\langle\alpha_1, \beta_2\rangle=1$, $\langle\alpha_2, \beta_1\rangle=1$, $\langle\alpha_2, \beta_2\rangle=1$, 
\\
\begin{tabular}{cc}$\alpha_1$+ & $\alpha_2$\\\tiny{ $A_2$+$A_1$ } & \tiny{ $A_2$ }\end{tabular} &\begin{tabular}{c} = \\~ \end{tabular} & \begin{tabular}{cc}$\beta_1$+ & $\beta_2$\\\tiny{ $A_2$+$A_1$ } & \tiny{ $A_2$ }\end{tabular} & $A_3$ & $E_6$ &  $\langle\alpha_1, \beta_1\rangle=1$, $\langle\alpha_1, \beta_2\rangle=1$, $\langle\alpha_2, \beta_1\rangle=1$, $\langle\alpha_2, \beta_2\rangle=1$, 
\\
\begin{tabular}{cc}$\alpha_1$+ & $\alpha_2$\\\tiny{ $A_2$ } & \tiny{  }\end{tabular} &\begin{tabular}{c} = \\~ \end{tabular} & \begin{tabular}{c}$\beta_1$\\\tiny{ $A_2$ }\end{tabular} & $A_2$ & $A_4$+$A_1$ & $\langle\alpha_1, \alpha_2\rangle=-1$,  $\langle\alpha_1, \beta_1\rangle=1$, $\langle\alpha_2, \beta_1\rangle=1$, 
\\
\begin{tabular}{cc}$\alpha_1$+ & $\alpha_2$\\\tiny{  } & \tiny{  }\end{tabular} &\begin{tabular}{c} = \\~ \end{tabular} & \begin{tabular}{ccc}$\beta_1$+ & $\beta_2$+ & $\beta_3$\\\tiny{ $A_2$+$A_1$ } & \tiny{ $A_2$+$A_1$ } & \tiny{ $A_2$ }\end{tabular} & $D_4$ & $E_6$ & $\langle\alpha_1, \alpha_2\rangle=1$,  $\langle\alpha_1, \beta_1\rangle=1$, $\langle\alpha_1, \beta_2\rangle=1$, $\langle\alpha_1, \beta_3\rangle=1$, $\langle\alpha_2, \beta_1\rangle=1$, $\langle\alpha_2, \beta_2\rangle=1$, $\langle\alpha_2, \beta_3\rangle=1$, 
\\
\begin{tabular}{cc}$\alpha_1$+ & $\alpha_2$\\\tiny{ $A_2$ } & \tiny{ $A_2$ }\end{tabular} &\begin{tabular}{c} = \\~ \end{tabular} & \begin{tabular}{ccc}$\beta_1$+ & $\beta_2$+ & $\beta_3$\\\tiny{ $A_2$+$A_1$ } & \tiny{ $A_2$+$A_1$ } & \tiny{ $A_2$ }\end{tabular} & $D_4$ & $E_6$ & $\langle\alpha_1, \alpha_2\rangle=1$,  $\langle\alpha_1, \beta_1\rangle=1$, $\langle\alpha_1, \beta_2\rangle=1$, $\langle\alpha_1, \beta_3\rangle=1$, $\langle\alpha_2, \beta_1\rangle=1$, $\langle\alpha_2, \beta_2\rangle=1$, $\langle\alpha_2, \beta_3\rangle=1$, 
\\
\begin{tabular}{cc}$\alpha_1$+ & $\alpha_2$\\\tiny{ $A_1$ } & \tiny{  }\end{tabular} &\begin{tabular}{c} = \\~ \end{tabular} & \begin{tabular}{cc}$\beta_1$+ & $\beta_2$\\\tiny{ $A_2$+$A_1$ } & \tiny{ $A_2$ }\end{tabular} & $A_3$ & $E_6$ &  $\langle\alpha_1, \beta_1\rangle=1$, $\langle\alpha_1, \beta_2\rangle=1$, $\langle\alpha_2, \beta_1\rangle=1$, $\langle\alpha_2, \beta_2\rangle=1$, 
\\
\begin{tabular}{cc}$\alpha_1$+ & $\alpha_2$\\\tiny{ $A_1$ } & \tiny{ $A_2$ }\end{tabular} &\begin{tabular}{c} = \\~ \end{tabular} & \begin{tabular}{c}$\beta_1$\\\tiny{ $A_2$+$A_1$ }\end{tabular} & $A_2$ & $A_5$ & $\langle\alpha_1, \alpha_2\rangle=-1$,  $\langle\alpha_1, \beta_1\rangle=1$, $\langle\alpha_2, \beta_1\rangle=1$, 
\\
\begin{tabular}{cc}$\alpha_1$+ & $\alpha_2$\\\tiny{ $A_2$+$A_1$ } & \tiny{  }\end{tabular} &\begin{tabular}{c} = \\~ \end{tabular} & \begin{tabular}{cc}$\beta_1$+ & $\beta_2$\\\tiny{ $A_2$+$A_1$ } & \tiny{ $A_2$ }\end{tabular} & $A_3$ & $E_6$ &  $\langle\alpha_1, \beta_1\rangle=1$, $\langle\alpha_1, \beta_2\rangle=1$, $\langle\alpha_2, \beta_1\rangle=1$, $\langle\alpha_2, \beta_2\rangle=1$, 
\\
\begin{tabular}{cc}$\alpha_1$+ & $\alpha_2$\\\tiny{ $A_2$+$A_1$ } & \tiny{ $A_1$ }\end{tabular} &\begin{tabular}{c} = \\~ \end{tabular} & \begin{tabular}{cc}$\beta_1$+ & $\beta_2$\\\tiny{ $A_2$+$A_1$ } & \tiny{ $A_2$+$A_1$ }\end{tabular} & $A_3$ & $E_6$ &  $\langle\alpha_1, \beta_1\rangle=1$, $\langle\alpha_1, \beta_2\rangle=1$, $\langle\alpha_2, \beta_1\rangle=1$, $\langle\alpha_2, \beta_2\rangle=1$, 
\\
\begin{tabular}{cc}$\alpha_1$+ & $\alpha_2$\\\tiny{ $A_2$ } & \tiny{ $A_2$ }\end{tabular} &\begin{tabular}{c} = \\~ \end{tabular} & \begin{tabular}{cc}$\beta_1$+ & $\beta_2$\\\tiny{ $A_2$+$A_1$ } & \tiny{ $A_2$+$A_1$ }\end{tabular} & $A_3$ & $E_6$ &  $\langle\alpha_1, \beta_1\rangle=1$, $\langle\alpha_1, \beta_2\rangle=1$, $\langle\alpha_2, \beta_1\rangle=1$, $\langle\alpha_2, \beta_2\rangle=1$, 
\\
\hline\multicolumn{5}{c}{$\mathfrak{k}$-semisimple type: $A_1$+$A_1$+$A_1$}\\
\hline\hline\begin{tabular}{cc}$\alpha_1$+ & $\alpha_2$\\\tiny{ $A_1$ } & \tiny{  }\end{tabular} &\begin{tabular}{c} = \\~ \end{tabular} & \begin{tabular}{c}$\beta_1$\\\tiny{ $A_1$ }\end{tabular} & $A_2$ & $A_3$+2$A_1$ & $\langle\alpha_1, \alpha_2\rangle=-1$,  $\langle\alpha_1, \beta_1\rangle=1$, $\langle\alpha_2, \beta_1\rangle=1$, 
\\
\begin{tabular}{cc}$\alpha_1$+ & $\alpha_2$\\\tiny{ $A_1$ } & \tiny{ $A_1$' }\end{tabular} &\begin{tabular}{c} = \\~ \end{tabular} & \begin{tabular}{c}$\beta_1$\\\tiny{ $A_1$+$A_1$' }\end{tabular} & $A_2$ & $A_4$+$A_1$ & $\langle\alpha_1, \alpha_2\rangle=-1$,  $\langle\alpha_1, \beta_1\rangle=1$, $\langle\alpha_2, \beta_1\rangle=1$, 
\\
\begin{tabular}{cc}$\alpha_1$+ & $\alpha_2$\\\tiny{ $A_1$+$A_1$'+$A_1$'' } & \tiny{ $A_1$' }\end{tabular} &\begin{tabular}{c} = \\~ \end{tabular} & \begin{tabular}{cc}$\beta_1$+ & $\beta_2$\\\tiny{ $A_1$+$A_1$' } & \tiny{ $A_1$''+$A_1$ }\end{tabular} & $A_3$ & $E_6$ &  $\langle\alpha_1, \beta_1\rangle=1$, $\langle\alpha_1, \beta_2\rangle=1$, $\langle\alpha_2, \beta_1\rangle=1$, $\langle\alpha_2, \beta_2\rangle=1$, 
\\
\begin{tabular}{cc}$\alpha_1$+ & $\alpha_2$\\\tiny{ $A_1$ } & \tiny{ $A_1$' }\end{tabular} &\begin{tabular}{c} = \\~ \end{tabular} & \begin{tabular}{cc}$\beta_1$+ & $\beta_2$\\\tiny{ $A_1$+$A_1$' } & \tiny{ $A_1$''+$A_1$ }\end{tabular} & $A_3$ & $E_6$ &  $\langle\alpha_1, \beta_1\rangle=1$, $\langle\alpha_1, \beta_2\rangle=1$, $\langle\alpha_2, \beta_1\rangle=1$, $\langle\alpha_2, \beta_2\rangle=1$, 
\\
\begin{tabular}{cc}$\alpha_1$+ & $\alpha_2$\\\tiny{ $A_1$ } & \tiny{ $A_1$' }\end{tabular} &\begin{tabular}{c} = \\~ \end{tabular} & \begin{tabular}{cc}$\beta_1$+ & $\beta_2$\\\tiny{ $A_1$ } & \tiny{ $A_1$' }\end{tabular} & $A_3$ & $A_5$+$A_1$ &  $\langle\alpha_1, \beta_1\rangle=1$, $\langle\alpha_1, \beta_2\rangle=1$, $\langle\alpha_2, \beta_1\rangle=1$, $\langle\alpha_2, \beta_2\rangle=1$, 
\\
\begin{tabular}{cc}$\alpha_1$+ & $\alpha_2$\\\tiny{ $A_1$+$A_1$' } & \tiny{  }\end{tabular} &\begin{tabular}{c} = \\~ \end{tabular} & \begin{tabular}{cc}$\beta_1$+ & $\beta_2$\\\tiny{ $A_1$ } & \tiny{ $A_1$' }\end{tabular} & $A_3$ & $A_5$+$A_1$ &  $\langle\alpha_1, \beta_1\rangle=1$, $\langle\alpha_1, \beta_2\rangle=1$, $\langle\alpha_2, \beta_1\rangle=1$, $\langle\alpha_2, \beta_2\rangle=1$, 
\\
\begin{tabular}{cc}$\alpha_1$+ & $\alpha_2$\\\tiny{ $A_1$+$A_1$' } & \tiny{ $A_1$'' }\end{tabular} &\begin{tabular}{c} = \\~ \end{tabular} & \begin{tabular}{ccc}$\beta_1$+ & $\beta_2$+ & $\beta_3$\\\tiny{ $A_1$ } & \tiny{ $A_1$' } & \tiny{ $A_1$'' }\end{tabular} & $D_4$ & $E_6$ & $\langle\alpha_1, \alpha_2\rangle=1$,  $\langle\alpha_1, \beta_1\rangle=1$, $\langle\alpha_1, \beta_2\rangle=1$, $\langle\alpha_1, \beta_3\rangle=1$, $\langle\alpha_2, \beta_1\rangle=1$, $\langle\alpha_2, \beta_2\rangle=1$, $\langle\alpha_2, \beta_3\rangle=1$, 
\\
\begin{tabular}{cc}$\alpha_1$+ & $\alpha_2$\\\tiny{ $A_1$+$A_1$'+$A_1$'' } & \tiny{  }\end{tabular} &\begin{tabular}{c} = \\~ \end{tabular} & \begin{tabular}{ccc}$\beta_1$+ & $\beta_2$+ & $\beta_3$\\\tiny{ $A_1$ } & \tiny{ $A_1$' } & \tiny{ $A_1$'' }\end{tabular} & $D_4$ & $E_6$ & $\langle\alpha_1, \alpha_2\rangle=1$,  $\langle\alpha_1, \beta_1\rangle=1$, $\langle\alpha_1, \beta_2\rangle=1$, $\langle\alpha_1, \beta_3\rangle=1$, $\langle\alpha_2, \beta_1\rangle=1$, $\langle\alpha_2, \beta_2\rangle=1$, $\langle\alpha_2, \beta_3\rangle=1$, 
\\
\begin{tabular}{cc}$\alpha_1$+ & $\alpha_2$\\\tiny{ $A_1$ } & \tiny{ $A_1$ }\end{tabular} &\begin{tabular}{c} = \\~ \end{tabular} & \begin{tabular}{ccc}$\beta_1$+ & $\beta_2$+ & $\beta_3$\\\tiny{ $A_1$+$A_1$' } & \tiny{ $A_1$''+$A_1$' } & \tiny{ $A_1$''+$A_1$ }\end{tabular} & $D_4$ & $E_6$ & $\langle\alpha_1, \alpha_2\rangle=1$,  $\langle\alpha_1, \beta_1\rangle=1$, $\langle\alpha_1, \beta_2\rangle=1$, $\langle\alpha_1, \beta_3\rangle=1$, $\langle\alpha_2, \beta_1\rangle=1$, $\langle\alpha_2, \beta_2\rangle=1$, $\langle\alpha_2, \beta_3\rangle=1$, 
\\
\begin{tabular}{cc}$\alpha_1$+ & $\alpha_2$\\\tiny{ $A_1$ } & \tiny{ $A_1$' }\end{tabular} &\begin{tabular}{c} = \\~ \end{tabular} & \begin{tabular}{cc}$\beta_1$+ & $\beta_2$\\\tiny{ $A_1$+$A_1$' } & \tiny{ $A_1$+$A_1$'' }\end{tabular} & $A_3$ & $E_6$ &  $\langle\alpha_1, \beta_1\rangle=1$, $\langle\alpha_1, \beta_2\rangle=1$, $\langle\alpha_2, \beta_1\rangle=1$, $\langle\alpha_2, \beta_2\rangle=1$, 
\\
\begin{tabular}{cc}$\alpha_1$+ & $\alpha_2$\\\tiny{ $A_1$+$A_1$' } & \tiny{ $A_1$'' }\end{tabular} &\begin{tabular}{c} = \\~ \end{tabular} & \begin{tabular}{c}$\beta_1$\\\tiny{ $A_1$+$A_1$'+$A_1$'' }\end{tabular} & $A_2$ & $D_5$ & $\langle\alpha_1, \alpha_2\rangle=-1$,  $\langle\alpha_1, \beta_1\rangle=1$, $\langle\alpha_2, \beta_1\rangle=1$, 
\\
\begin{tabular}{cc}$\alpha_1$+ & $\alpha_2$\\\tiny{ $A_1$ } & \tiny{ $A_1$ }\end{tabular} &\begin{tabular}{c} = \\~ \end{tabular} & \begin{tabular}{cc}$\beta_1$+ & $\beta_2$\\\tiny{ $A_1$+$A_1$'+$A_1$'' } & \tiny{ $A_1$+$A_1$'+$A_1$'' }\end{tabular} & $A_3$ & $D_5$ &  $\langle\alpha_1, \beta_1\rangle=1$, $\langle\alpha_1, \beta_2\rangle=1$, $\langle\alpha_2, \beta_1\rangle=1$, $\langle\alpha_2, \beta_2\rangle=1$, 
\\
\begin{tabular}{cc}$\alpha_1$+ & $\alpha_2$\\\tiny{ $A_1$+$A_1$' } & \tiny{ $A_1$''+$A_1$' }\end{tabular} &\begin{tabular}{c} = \\~ \end{tabular} & \begin{tabular}{cc}$\beta_1$+ & $\beta_2$\\\tiny{ $A_1$+$A_1$'+$A_1$'' } & \tiny{ $A_1$'' }\end{tabular} & $A_3$ & $E_6$ &  $\langle\alpha_1, \beta_1\rangle=1$, $\langle\alpha_1, \beta_2\rangle=1$, $\langle\alpha_2, \beta_1\rangle=1$, $\langle\alpha_2, \beta_2\rangle=1$, 
\\
\begin{tabular}{cc}$\alpha_1$+ & $\alpha_2$\\\tiny{ $A_1$+$A_1$' } & \tiny{ $A_1$ }\end{tabular} &\begin{tabular}{c} = \\~ \end{tabular} & \begin{tabular}{cc}$\beta_1$+ & $\beta_2$\\\tiny{ $A_1$+$A_1$'+$A_1$'' } & \tiny{ $A_1$+$A_1$' }\end{tabular} & $A_3$ & $E_6$ &  $\langle\alpha_1, \beta_1\rangle=1$, $\langle\alpha_1, \beta_2\rangle=1$, $\langle\alpha_2, \beta_1\rangle=1$, $\langle\alpha_2, \beta_2\rangle=1$, 
\\
\begin{tabular}{cc}$\alpha_1$+ & $\alpha_2$\\\tiny{ $A_1$ } & \tiny{ $A_1$' }\end{tabular} &\begin{tabular}{c} = \\~ \end{tabular} & \begin{tabular}{cc}$\beta_1$+ & $\beta_2$\\\tiny{ $A_1$+$A_1$' } & \tiny{  }\end{tabular} & $A_3$ & $A_5$+$A_1$ &  $\langle\alpha_1, \beta_1\rangle=1$, $\langle\alpha_1, \beta_2\rangle=1$, $\langle\alpha_2, \beta_1\rangle=1$, $\langle\alpha_2, \beta_2\rangle=1$, 
\\
\begin{tabular}{cc}$\alpha_1$+ & $\alpha_2$\\\tiny{ $A_1$ } & \tiny{ $A_1$' }\end{tabular} &\begin{tabular}{c} = \\~ \end{tabular} & \begin{tabular}{cc}$\beta_1$+ & $\beta_2$\\\tiny{ $A_1$+$A_1$'+$A_1$'' } & \tiny{ $A_1$'' }\end{tabular} & $A_3$ & $E_6$ &  $\langle\alpha_1, \beta_1\rangle=1$, $\langle\alpha_1, \beta_2\rangle=1$, $\langle\alpha_2, \beta_1\rangle=1$, $\langle\alpha_2, \beta_2\rangle=1$, 
\\
\begin{tabular}{cc}$\alpha_1$+ & $\alpha_2$\\\tiny{ $A_1$ } & \tiny{ $A_1$' }\end{tabular} &\begin{tabular}{c} = \\~ \end{tabular} & \begin{tabular}{cc}$\beta_1$+ & $\beta_2$\\\tiny{ $A_1$+$A_1$'+$A_1$'' } & \tiny{ $A_1$ }\end{tabular} & $A_3$ & $E_6$ &  $\langle\alpha_1, \beta_1\rangle=1$, $\langle\alpha_1, \beta_2\rangle=1$, $\langle\alpha_2, \beta_1\rangle=1$, $\langle\alpha_2, \beta_2\rangle=1$, 
\\
\begin{tabular}{c}2$\alpha_1$\\\tiny{ $A_1$ }\end{tabular} &\begin{tabular}{c} = \\~ \end{tabular} & \begin{tabular}{cccc}$\beta_1$+ & $\beta_2$+ & $\beta_3$+ & $\beta_4$\\\tiny{ $A_1$+$A_1$' } & \tiny{ $A_1$''+$A_1$' } & \tiny{ $A_1$''+$A_1$ } & \tiny{  }\end{tabular} & $D_4$ & $E_6$ &  $\langle\alpha_1, \beta_1\rangle=1$, $\langle\alpha_1, \beta_2\rangle=1$, $\langle\alpha_1, \beta_3\rangle=1$, $\langle\alpha_1, \beta_4\rangle=1$, 
\\
\begin{tabular}{cc}$\alpha_1$+ & $\alpha_2$\\\tiny{ $A_1$ } & \tiny{ $A_1$' }\end{tabular} &\begin{tabular}{c} = \\~ \end{tabular} & \begin{tabular}{ccc}$\beta_1$+ & $\beta_2$+ & $\beta_3$\\\tiny{ $A_1$+$A_1$' } & \tiny{ $A_1$+$A_1$'' } & \tiny{  }\end{tabular} & $D_4$ & $E_6$ & $\langle\alpha_1, \alpha_2\rangle=1$,  $\langle\alpha_1, \beta_1\rangle=1$, $\langle\alpha_1, \beta_2\rangle=1$, $\langle\alpha_1, \beta_3\rangle=1$, $\langle\alpha_2, \beta_1\rangle=1$, $\langle\alpha_2, \beta_2\rangle=1$, $\langle\alpha_2, \beta_3\rangle=1$, 
\\
\begin{tabular}{ccc}$\alpha_1$+ & $\alpha_2$+ & $\alpha_3$\\\tiny{ $A_1$ } & \tiny{ $A_1$' } & \tiny{ $A_1$'' }\end{tabular} &\begin{tabular}{c} = \\~ \end{tabular} & \begin{tabular}{ccc}$\beta_1$+ & $\beta_2$+ & $\beta_3$\\\tiny{ $A_1$ } & \tiny{ $A_1$' } & \tiny{ $A_1$'' }\end{tabular} & $A_5$ & $E_6$ &  $\langle\alpha_1, \beta_2\rangle=1$, $\langle\alpha_1, \beta_3\rangle=1$, $\langle\alpha_2, \beta_1\rangle=1$, $\langle\alpha_2, \beta_3\rangle=1$, $\langle\alpha_3, \beta_1\rangle=1$, $\langle\alpha_3, \beta_2\rangle=1$, 
\\
\hline\multicolumn{5}{c}{$\mathfrak{k}$-semisimple type: $A_2$}\\
\hline\hline\begin{tabular}{cc}$\alpha_1$+ & $\alpha_2$\\\tiny{  } & \tiny{  }\end{tabular} &\begin{tabular}{c} = \\~ \end{tabular} & \begin{tabular}{cc}$\beta_1$+ & $\beta_2$\\\tiny{ $A_2$ } & \tiny{ $A_2$ }\end{tabular} & $A_3$ & $D_5$ &  $\langle\alpha_1, \beta_1\rangle=1$, $\langle\alpha_1, \beta_2\rangle=1$, $\langle\alpha_2, \beta_1\rangle=1$, $\langle\alpha_2, \beta_2\rangle=1$, 
\\
\begin{tabular}{cc}$\alpha_1$+ & $\alpha_2$\\\tiny{ $A_2$ } & \tiny{  }\end{tabular} &\begin{tabular}{c} = \\~ \end{tabular} & \begin{tabular}{c}$\beta_1$\\\tiny{ $A_2$ }\end{tabular} & $A_2$ & $A_4$ & $\langle\alpha_1, \alpha_2\rangle=-1$,  $\langle\alpha_1, \beta_1\rangle=1$, $\langle\alpha_2, \beta_1\rangle=1$, 
\\
\begin{tabular}{cc}$\alpha_1$+ & $\alpha_2$\\\tiny{ $A_2$ } & \tiny{ $A_2$ }\end{tabular} &\begin{tabular}{c} = \\~ \end{tabular} & \begin{tabular}{cc}$\beta_1$+ & $\beta_2$\\\tiny{ $A_2$ } & \tiny{ $A_2$ }\end{tabular} & $A_3$ & $D_5$ &  $\langle\alpha_1, \beta_1\rangle=1$, $\langle\alpha_1, \beta_2\rangle=1$, $\langle\alpha_2, \beta_1\rangle=1$, $\langle\alpha_2, \beta_2\rangle=1$, 
\\
\begin{tabular}{cc}$\alpha_1$+ & $\alpha_2$\\\tiny{  } & \tiny{  }\end{tabular} &\begin{tabular}{c} = \\~ \end{tabular} & \begin{tabular}{ccc}$\beta_1$+ & $\beta_2$+ & $\beta_3$\\\tiny{ $A_2$ } & \tiny{ $A_2$ } & \tiny{ $A_2$ }\end{tabular} & $D_4$ & $E_6$ & $\langle\alpha_1, \alpha_2\rangle=1$,  $\langle\alpha_1, \beta_1\rangle=1$, $\langle\alpha_1, \beta_2\rangle=1$, $\langle\alpha_1, \beta_3\rangle=1$, $\langle\alpha_2, \beta_1\rangle=1$, $\langle\alpha_2, \beta_2\rangle=1$, $\langle\alpha_2, \beta_3\rangle=1$, 
\\
\begin{tabular}{cc}$\alpha_1$+ & $\alpha_2$\\\tiny{ $A_2$ } & \tiny{ $A_2$ }\end{tabular} &\begin{tabular}{c} = \\~ \end{tabular} & \begin{tabular}{ccc}$\beta_1$+ & $\beta_2$+ & $\beta_3$\\\tiny{ $A_2$ } & \tiny{ $A_2$ } & \tiny{ $A_2$ }\end{tabular} & $D_4$ & $E_6$ & $\langle\alpha_1, \alpha_2\rangle=1$,  $\langle\alpha_1, \beta_1\rangle=1$, $\langle\alpha_1, \beta_2\rangle=1$, $\langle\alpha_1, \beta_3\rangle=1$, $\langle\alpha_2, \beta_1\rangle=1$, $\langle\alpha_2, \beta_2\rangle=1$, $\langle\alpha_2, \beta_3\rangle=1$, 
\\
\begin{tabular}{cc}$\alpha_1$+ & $\alpha_2$\\\tiny{ $A_2$ } & \tiny{  }\end{tabular} &\begin{tabular}{c} = \\~ \end{tabular} & \begin{tabular}{cc}$\beta_1$+ & $\beta_2$\\\tiny{ $A_2$ } & \tiny{ $A_2$ }\end{tabular} & $A_3$ & $D_5$ &  $\langle\alpha_1, \beta_1\rangle=1$, $\langle\alpha_1, \beta_2\rangle=1$, $\langle\alpha_2, \beta_1\rangle=1$, $\langle\alpha_2, \beta_2\rangle=1$, 
\\
\hline\multicolumn{5}{c}{$\mathfrak{k}$-semisimple type: $A_1$+$A_1$}\\
\hline\hline\begin{tabular}{cc}$\alpha_1$+ & $\alpha_2$\\\tiny{  } & \tiny{  }\end{tabular} &\begin{tabular}{c} = \\~ \end{tabular} & \begin{tabular}{cc}$\beta_1$+ & $\beta_2$\\\tiny{ $A_1$+$A_1$' } & \tiny{ $A_1$+$A_1$' }\end{tabular} & $A_3$ & $D_4$ &  $\langle\alpha_1, \beta_1\rangle=1$, $\langle\alpha_1, \beta_2\rangle=1$, $\langle\alpha_2, \beta_1\rangle=1$, $\langle\alpha_2, \beta_2\rangle=1$, 
\\
\begin{tabular}{cc}$\alpha_1$+ & $\alpha_2$\\\tiny{ $A_1$ } & \tiny{  }\end{tabular} &\begin{tabular}{c} = \\~ \end{tabular} & \begin{tabular}{c}$\beta_1$\\\tiny{ $A_1$ }\end{tabular} & $A_2$ & $A_3$+$A_1$ & $\langle\alpha_1, \alpha_2\rangle=-1$,  $\langle\alpha_1, \beta_1\rangle=1$, $\langle\alpha_2, \beta_1\rangle=1$, 
\\
\begin{tabular}{cc}$\alpha_1$+ & $\alpha_2$\\\tiny{ $A_1$ } & \tiny{  }\end{tabular} &\begin{tabular}{c} = \\~ \end{tabular} & \begin{tabular}{cc}$\beta_1$+ & $\beta_2$\\\tiny{ $A_1$+$A_1$' } & \tiny{ $A_1$' }\end{tabular} & $A_3$ & $D_5$ &  $\langle\alpha_1, \beta_1\rangle=1$, $\langle\alpha_1, \beta_2\rangle=1$, $\langle\alpha_2, \beta_1\rangle=1$, $\langle\alpha_2, \beta_2\rangle=1$, 
\\
\begin{tabular}{cc}$\alpha_1$+ & $\alpha_2$\\\tiny{ $A_1$ } & \tiny{ $A_1$' }\end{tabular} &\begin{tabular}{c} = \\~ \end{tabular} & \begin{tabular}{cc}$\beta_1$+ & $\beta_2$\\\tiny{ $A_1$ } & \tiny{ $A_1$' }\end{tabular} & $A_3$ & $A_5$ &  $\langle\alpha_1, \beta_1\rangle=1$, $\langle\alpha_1, \beta_2\rangle=1$, $\langle\alpha_2, \beta_1\rangle=1$, $\langle\alpha_2, \beta_2\rangle=1$, 
\\
\begin{tabular}{cc}$\alpha_1$+ & $\alpha_2$\\\tiny{ $A_1$ } & \tiny{ $A_1$ }\end{tabular} &\begin{tabular}{c} = \\~ \end{tabular} & \begin{tabular}{ccc}$\beta_1$+ & $\beta_2$+ & $\beta_3$\\\tiny{ $A_1$+$A_1$' } & \tiny{ $A_1$' } & \tiny{ $A_1$ }\end{tabular} & $D_4$ & $E_6$ & $\langle\alpha_1, \alpha_2\rangle=1$,  $\langle\alpha_1, \beta_1\rangle=1$, $\langle\alpha_1, \beta_2\rangle=1$, $\langle\alpha_1, \beta_3\rangle=1$, $\langle\alpha_2, \beta_1\rangle=1$, $\langle\alpha_2, \beta_2\rangle=1$, $\langle\alpha_2, \beta_3\rangle=1$, 
\\
\begin{tabular}{cc}$\alpha_1$+ & $\alpha_2$\\\tiny{  } & \tiny{  }\end{tabular} &\begin{tabular}{c} = \\~ \end{tabular} & \begin{tabular}{ccc}$\beta_1$+ & $\beta_2$+ & $\beta_3$\\\tiny{ $A_1$+$A_1$' } & \tiny{ $A_1$' } & \tiny{ $A_1$ }\end{tabular} & $D_4$ & $E_6$ & $\langle\alpha_1, \alpha_2\rangle=1$,  $\langle\alpha_1, \beta_1\rangle=1$, $\langle\alpha_1, \beta_2\rangle=1$, $\langle\alpha_1, \beta_3\rangle=1$, $\langle\alpha_2, \beta_1\rangle=1$, $\langle\alpha_2, \beta_2\rangle=1$, $\langle\alpha_2, \beta_3\rangle=1$, 
\\
\begin{tabular}{cc}$\alpha_1$+ & $\alpha_2$\\\tiny{ $A_1$ } & \tiny{ $A_1$' }\end{tabular} &\begin{tabular}{c} = \\~ \end{tabular} & \begin{tabular}{c}$\beta_1$\\\tiny{ $A_1$+$A_1$' }\end{tabular} & $A_2$ & $A_4$ & $\langle\alpha_1, \alpha_2\rangle=-1$,  $\langle\alpha_1, \beta_1\rangle=1$, $\langle\alpha_2, \beta_1\rangle=1$, 
\\
\begin{tabular}{cc}$\alpha_1$+ & $\alpha_2$\\\tiny{ $A_1$+$A_1$' } & \tiny{  }\end{tabular} &\begin{tabular}{c} = \\~ \end{tabular} & \begin{tabular}{cc}$\beta_1$+ & $\beta_2$\\\tiny{ $A_1$ } & \tiny{ $A_1$' }\end{tabular} & $A_3$ & $A_5$ &  $\langle\alpha_1, \beta_1\rangle=1$, $\langle\alpha_1, \beta_2\rangle=1$, $\langle\alpha_2, \beta_1\rangle=1$, $\langle\alpha_2, \beta_2\rangle=1$, 
\\
\begin{tabular}{cc}$\alpha_1$+ & $\alpha_2$\\\tiny{ $A_1$+$A_1$' } & \tiny{  }\end{tabular} &\begin{tabular}{c} = \\~ \end{tabular} & \begin{tabular}{c}$\beta_1$\\\tiny{ $A_1$+$A_1$' }\end{tabular} & $A_2$ & $D_4$ & $\langle\alpha_1, \alpha_2\rangle=-1$,  $\langle\alpha_1, \beta_1\rangle=1$, $\langle\alpha_2, \beta_1\rangle=1$, 
\\
\begin{tabular}{cc}$\alpha_1$+ & $\alpha_2$\\\tiny{ $A_1$ } & \tiny{  }\end{tabular} &\begin{tabular}{c} = \\~ \end{tabular} & \begin{tabular}{cc}$\beta_1$+ & $\beta_2$\\\tiny{ $A_1$+$A_1$' } & \tiny{ $A_1$ }\end{tabular} & $A_3$ & $D_5$ &  $\langle\alpha_1, \beta_1\rangle=1$, $\langle\alpha_1, \beta_2\rangle=1$, $\langle\alpha_2, \beta_1\rangle=1$, $\langle\alpha_2, \beta_2\rangle=1$, 
\\
\begin{tabular}{cc}$\alpha_1$+ & $\alpha_2$\\\tiny{ $A_1$+$A_1$' } & \tiny{ $A_1$' }\end{tabular} &\begin{tabular}{c} = \\~ \end{tabular} & \begin{tabular}{cc}$\beta_1$+ & $\beta_2$\\\tiny{ $A_1$+$A_1$' } & \tiny{ $A_1$' }\end{tabular} & $A_3$ & $D_5$ &  $\langle\alpha_1, \beta_1\rangle=1$, $\langle\alpha_1, \beta_2\rangle=1$, $\langle\alpha_2, \beta_1\rangle=1$, $\langle\alpha_2, \beta_2\rangle=1$, 
\\
\begin{tabular}{cc}$\alpha_1$+ & $\alpha_2$\\\tiny{ $A_1$+$A_1$' } & \tiny{ $A_1$ }\end{tabular} &\begin{tabular}{c} = \\~ \end{tabular} & \begin{tabular}{cc}$\beta_1$+ & $\beta_2$\\\tiny{ $A_1$+$A_1$' } & \tiny{ $A_1$ }\end{tabular} & $A_3$ & $D_5$ &  $\langle\alpha_1, \beta_1\rangle=1$, $\langle\alpha_1, \beta_2\rangle=1$, $\langle\alpha_2, \beta_1\rangle=1$, $\langle\alpha_2, \beta_2\rangle=1$, 
\\
\begin{tabular}{cc}$\alpha_1$+ & $\alpha_2$\\\tiny{ $A_1$ } & \tiny{ $A_1$ }\end{tabular} &\begin{tabular}{c} = \\~ \end{tabular} & \begin{tabular}{cc}$\beta_1$+ & $\beta_2$\\\tiny{ $A_1$+$A_1$' } & \tiny{ $A_1$+$A_1$' }\end{tabular} & $A_3$ & $D_5$ &  $\langle\alpha_1, \beta_1\rangle=1$, $\langle\alpha_1, \beta_2\rangle=1$, $\langle\alpha_2, \beta_1\rangle=1$, $\langle\alpha_2, \beta_2\rangle=1$, 
\\
\begin{tabular}{cc}$\alpha_1$+ & $\alpha_2$\\\tiny{ $A_1$+$A_1$' } & \tiny{ $A_1$+$A_1$' }\end{tabular} &\begin{tabular}{c} = \\~ \end{tabular} & \begin{tabular}{ccc}$\beta_1$+ & $\beta_2$+ & $\beta_3$\\\tiny{ $A_1$+$A_1$' } & \tiny{ $A_1$' } & \tiny{ $A_1$ }\end{tabular} & $D_4$ & $E_6$ & $\langle\alpha_1, \alpha_2\rangle=1$,  $\langle\alpha_1, \beta_1\rangle=1$, $\langle\alpha_1, \beta_2\rangle=1$, $\langle\alpha_1, \beta_3\rangle=1$, $\langle\alpha_2, \beta_1\rangle=1$, $\langle\alpha_2, \beta_2\rangle=1$, $\langle\alpha_2, \beta_3\rangle=1$, 
\\
\begin{tabular}{cc}$\alpha_1$+ & $\alpha_2$\\\tiny{ $A_1$ } & \tiny{ $A_1$' }\end{tabular} &\begin{tabular}{c} = \\~ \end{tabular} & \begin{tabular}{cc}$\beta_1$+ & $\beta_2$\\\tiny{ $A_1$+$A_1$' } & \tiny{  }\end{tabular} & $A_3$ & $A_5$ &  $\langle\alpha_1, \beta_1\rangle=1$, $\langle\alpha_1, \beta_2\rangle=1$, $\langle\alpha_2, \beta_1\rangle=1$, $\langle\alpha_2, \beta_2\rangle=1$, 
\\
\hline\multicolumn{5}{c}{$\mathfrak{k}$-semisimple type: $A_1$}\\
\hline\hline\begin{tabular}{cc}$\alpha_1$+ & $\alpha_2$\\\tiny{  } & \tiny{  }\end{tabular} &\begin{tabular}{c} = \\~ \end{tabular} & \begin{tabular}{cc}$\beta_1$+ & $\beta_2$\\\tiny{ $A_1$ } & \tiny{ $A_1$ }\end{tabular} & $A_3$ & $D_4$ &  $\langle\alpha_1, \beta_1\rangle=1$, $\langle\alpha_1, \beta_2\rangle=1$, $\langle\alpha_2, \beta_1\rangle=1$, $\langle\alpha_2, \beta_2\rangle=1$, 
\\
\begin{tabular}{cc}$\alpha_1$+ & $\alpha_2$\\\tiny{ $A_1$ } & \tiny{  }\end{tabular} &\begin{tabular}{c} = \\~ \end{tabular} & \begin{tabular}{c}$\beta_1$\\\tiny{ $A_1$ }\end{tabular} & $A_2$ & $A_3$ & $\langle\alpha_1, \alpha_2\rangle=-1$,  $\langle\alpha_1, \beta_1\rangle=1$, $\langle\alpha_2, \beta_1\rangle=1$, 
\\
\begin{tabular}{cc}$\alpha_1$+ & $\alpha_2$\\\tiny{ $A_1$ } & \tiny{ $A_1$ }\end{tabular} &\begin{tabular}{c} = \\~ \end{tabular} & \begin{tabular}{cc}$\beta_1$+ & $\beta_2$\\\tiny{ $A_1$ } & \tiny{ $A_1$ }\end{tabular} & $A_3$ & $D_4$ &  $\langle\alpha_1, \beta_1\rangle=1$, $\langle\alpha_1, \beta_2\rangle=1$, $\langle\alpha_2, \beta_1\rangle=1$, $\langle\alpha_2, \beta_2\rangle=1$, 
\\
\end{longtable}

\normalsize

\subsection{$E_7$: $\lL$-strictly infinite weights and corresponding relations}
\noindent Number of different non-solvable subalgebras up to $\gG$-automorphism such that $\nn\cap C(\kk_{ss})$ is a nilradical of a parabolic subalgebra of $\kk$ containing $C(\kk_{ss})\cap \hh$: 73834


\noindent Among them 7427 satisfy the cone condition and 66407 do not.
\tiny

\begin{longtable}{rcl p{1cm}p{1cm}p{3cm} } \multicolumn{3}{c}{ Relation / linked $\mathfrak{k}$-components} &$\alpha_i$'s, $\beta_i$'s generate & adding $\Delta(\mathfrak{k})$ generates&Non-zero scalar products\endhead
\hline\multicolumn{5}{c}{$\mathfrak{k}$-semisimple type: $A_3$+$A_3$}\\
\hline\hline\begin{tabular}{c}2$\alpha_1$\\\tiny{  }\end{tabular} &\begin{tabular}{c} = \\~ \end{tabular} & \begin{tabular}{cccc}$\beta_1$+ & $\beta_2$+ & $\beta_3$+ & $\beta_4$\\\tiny{ $A_3$+$A_3$' } & \tiny{ $A_3$+$A_3$' } & \tiny{ $A_3$+$A_3$' } & \tiny{ $A_3$+$A_3$' }\end{tabular} & $D_4$ & $E_7$ &  $\langle\alpha_1, \beta_1\rangle=1$, $\langle\alpha_1, \beta_2\rangle=1$, $\langle\alpha_1, \beta_3\rangle=1$, $\langle\alpha_1, \beta_4\rangle=1$, 
\\
\hline\multicolumn{5}{c}{$\mathfrak{k}$-semisimple type: $D_5$}\\
\hline\hline\begin{tabular}{cc}$\alpha_1$+ & $\alpha_2$\\\tiny{ $D_5$ } & \tiny{  }\end{tabular} &\begin{tabular}{c} = \\~ \end{tabular} & \begin{tabular}{c}$\beta_1$\\\tiny{ $D_5$ }\end{tabular} & $A_2$ & $E_7$ & $\langle\alpha_1, \alpha_2\rangle=-1$,  $\langle\alpha_1, \beta_1\rangle=1$, $\langle\alpha_2, \beta_1\rangle=1$, 
\\
\begin{tabular}{cc}$\alpha_1$+ & $\alpha_2$\\\tiny{ $D_5$ } & \tiny{  }\end{tabular} &\begin{tabular}{c} = \\~ \end{tabular} & \begin{tabular}{cc}$\beta_1$+ & $\beta_2$\\\tiny{ $D_5$ } & \tiny{ $D_5$ }\end{tabular} & $A_3$ & $E_7$ &  $\langle\alpha_1, \beta_1\rangle=1$, $\langle\alpha_1, \beta_2\rangle=1$, $\langle\alpha_2, \beta_1\rangle=1$, $\langle\alpha_2, \beta_2\rangle=1$, 
\\
\hline\multicolumn{5}{c}{$\mathfrak{k}$-semisimple type: $A_5$}\\
\hline\hline\begin{tabular}{cc}$\alpha_1$+ & $\alpha_2$\\\tiny{ $A_5$ } & \tiny{  }\end{tabular} &\begin{tabular}{c} = \\~ \end{tabular} & \begin{tabular}{c}$\beta_1$\\\tiny{ $A_5$ }\end{tabular} & $A_2$ & $A_7$ & $\langle\alpha_1, \alpha_2\rangle=-1$,  $\langle\alpha_1, \beta_1\rangle=1$, $\langle\alpha_2, \beta_1\rangle=1$, 
\\
\begin{tabular}{cc}$\alpha_1$+ & $\alpha_2$\\\tiny{ $A_5$ } & \tiny{ $A_5$ }\end{tabular} &\begin{tabular}{c} = \\~ \end{tabular} & \begin{tabular}{cc}$\beta_1$+ & $\beta_2$\\\tiny{ $A_5$ } & \tiny{ $A_5$ }\end{tabular} & $A_3$ & $E_7$ &  $\langle\alpha_1, \beta_1\rangle=1$, $\langle\alpha_1, \beta_2\rangle=1$, $\langle\alpha_2, \beta_1\rangle=1$, $\langle\alpha_2, \beta_2\rangle=1$, 
\\
\hline\multicolumn{5}{c}{$\mathfrak{k}$-semisimple type: $A_5$}\\
\hline\hline\begin{tabular}{cc}$\alpha_1$+ & $\alpha_2$\\\tiny{  } & \tiny{  }\end{tabular} &\begin{tabular}{c} = \\~ \end{tabular} & \begin{tabular}{ccc}$\beta_1$+ & $\beta_2$+ & $\beta_3$\\\tiny{ $A_5$ } & \tiny{ $A_5$ } & \tiny{ $A_5$ }\end{tabular} & $D_4$ & $E_7$ & $\langle\alpha_1, \alpha_2\rangle=1$,  $\langle\alpha_1, \beta_1\rangle=1$, $\langle\alpha_1, \beta_2\rangle=1$, $\langle\alpha_1, \beta_3\rangle=1$, $\langle\alpha_2, \beta_1\rangle=1$, $\langle\alpha_2, \beta_2\rangle=1$, $\langle\alpha_2, \beta_3\rangle=1$, 
\\
\begin{tabular}{cc}$\alpha_1$+ & $\alpha_2$\\\tiny{ $A_5$ } & \tiny{  }\end{tabular} &\begin{tabular}{c} = \\~ \end{tabular} & \begin{tabular}{c}$\beta_1$\\\tiny{ $A_5$ }\end{tabular} & $A_2$ & $E_7$ & $\langle\alpha_1, \alpha_2\rangle=-1$,  $\langle\alpha_1, \beta_1\rangle=1$, $\langle\alpha_2, \beta_1\rangle=1$, 
\\
\begin{tabular}{cc}$\alpha_1$+ & $\alpha_2$\\\tiny{ $A_5$ } & \tiny{  }\end{tabular} &\begin{tabular}{c} = \\~ \end{tabular} & \begin{tabular}{cc}$\beta_1$+ & $\beta_2$\\\tiny{ $A_5$ } & \tiny{ $A_5$ }\end{tabular} & $A_3$ & $E_7$ &  $\langle\alpha_1, \beta_1\rangle=1$, $\langle\alpha_1, \beta_2\rangle=1$, $\langle\alpha_2, \beta_1\rangle=1$, $\langle\alpha_2, \beta_2\rangle=1$, 
\\
\hline\multicolumn{5}{c}{$\mathfrak{k}$-semisimple type: $D_4$+$A_1$}\\
\hline\hline\begin{tabular}{cc}$\alpha_1$+ & $\alpha_2$\\\tiny{  } & \tiny{  }\end{tabular} &\begin{tabular}{c} = \\~ \end{tabular} & \begin{tabular}{cc}$\beta_1$+ & $\beta_2$\\\tiny{ $D_4$ } & \tiny{ $D_4$ }\end{tabular} & $A_3$ & $D_6$+$A_1$ &  $\langle\alpha_1, \beta_1\rangle=1$, $\langle\alpha_1, \beta_2\rangle=1$, $\langle\alpha_2, \beta_1\rangle=1$, $\langle\alpha_2, \beta_2\rangle=1$, 
\\
\begin{tabular}{cc}$\alpha_1$+ & $\alpha_2$\\\tiny{ $D_4$ } & \tiny{  }\end{tabular} &\begin{tabular}{c} = \\~ \end{tabular} & \begin{tabular}{c}$\beta_1$\\\tiny{ $D_4$ }\end{tabular} & $A_2$ & $D_6$+$A_1$ & $\langle\alpha_1, \alpha_2\rangle=-1$,  $\langle\alpha_1, \beta_1\rangle=1$, $\langle\alpha_2, \beta_1\rangle=1$, 
\\
\hline\multicolumn{5}{c}{$\mathfrak{k}$-semisimple type: $A_4$+$A_1$}\\
\hline\hline\begin{tabular}{cc}$\alpha_1$+ & $\alpha_2$\\\tiny{ $A_1$ } & \tiny{ $A_4$ }\end{tabular} &\begin{tabular}{c} = \\~ \end{tabular} & \begin{tabular}{c}$\beta_1$\\\tiny{ $A_4$+$A_1$ }\end{tabular} & $A_2$ & $E_7$ & $\langle\alpha_1, \alpha_2\rangle=-1$,  $\langle\alpha_1, \beta_1\rangle=1$, $\langle\alpha_2, \beta_1\rangle=1$, 
\\
\begin{tabular}{cc}$\alpha_1$+ & $\alpha_2$\\\tiny{ $A_4$+$A_1$ } & \tiny{ $A_1$ }\end{tabular} &\begin{tabular}{c} = \\~ \end{tabular} & \begin{tabular}{cc}$\beta_1$+ & $\beta_2$\\\tiny{ $A_4$+$A_1$ } & \tiny{ $A_4$+$A_1$ }\end{tabular} & $A_3$ & $E_7$ &  $\langle\alpha_1, \beta_1\rangle=1$, $\langle\alpha_1, \beta_2\rangle=1$, $\langle\alpha_2, \beta_1\rangle=1$, $\langle\alpha_2, \beta_2\rangle=1$, 
\\
\begin{tabular}{cc}$\alpha_1$+ & $\alpha_2$\\\tiny{ $A_4$ } & \tiny{ $A_4$ }\end{tabular} &\begin{tabular}{c} = \\~ \end{tabular} & \begin{tabular}{cc}$\beta_1$+ & $\beta_2$\\\tiny{ $A_4$+$A_1$ } & \tiny{ $A_4$+$A_1$ }\end{tabular} & $A_3$ & $E_7$ &  $\langle\alpha_1, \beta_1\rangle=1$, $\langle\alpha_1, \beta_2\rangle=1$, $\langle\alpha_2, \beta_1\rangle=1$, $\langle\alpha_2, \beta_2\rangle=1$, 
\\
\begin{tabular}{cc}$\alpha_1$+ & $\alpha_2$\\\tiny{ $A_1$ } & \tiny{ $A_4$ }\end{tabular} &\begin{tabular}{c} = \\~ \end{tabular} & \begin{tabular}{cc}$\beta_1$+ & $\beta_2$\\\tiny{ $A_4$+$A_1$ } & \tiny{ $A_4$ }\end{tabular} & $A_3$ & $E_7$ &  $\langle\alpha_1, \beta_1\rangle=1$, $\langle\alpha_1, \beta_2\rangle=1$, $\langle\alpha_2, \beta_1\rangle=1$, $\langle\alpha_2, \beta_2\rangle=1$, 
\\
\begin{tabular}{cc}$\alpha_1$+ & $\alpha_2$\\\tiny{ $A_1$ } & \tiny{ $A_4$ }\end{tabular} &\begin{tabular}{c} = \\~ \end{tabular} & \begin{tabular}{c}$\beta_1$\\\tiny{ $A_4$+$A_1$ }\end{tabular} & $A_2$ & $A_7$ & $\langle\alpha_1, \alpha_2\rangle=-1$,  $\langle\alpha_1, \beta_1\rangle=1$, $\langle\alpha_2, \beta_1\rangle=1$, 
\\
\hline\multicolumn{5}{c}{$\mathfrak{k}$-semisimple type: $A_3$+$A_2$}\\
\hline\hline\begin{tabular}{cc}$\alpha_1$+ & $\alpha_2$\\\tiny{ $A_3$+$A_2$ } & \tiny{  }\end{tabular} &\begin{tabular}{c} = \\~ \end{tabular} & \begin{tabular}{c}$\beta_1$\\\tiny{ $A_3$+$A_2$ }\end{tabular} & $A_2$ & $E_7$ & $\langle\alpha_1, \alpha_2\rangle=-1$,  $\langle\alpha_1, \beta_1\rangle=1$, $\langle\alpha_2, \beta_1\rangle=1$, 
\\
\begin{tabular}{cc}$\alpha_1$+ & $\alpha_2$\\\tiny{ $A_3$ } & \tiny{  }\end{tabular} &\begin{tabular}{c} = \\~ \end{tabular} & \begin{tabular}{c}$\beta_1$\\\tiny{ $A_3$ }\end{tabular} & $A_2$ & $A_5$+$A_2$ & $\langle\alpha_1, \alpha_2\rangle=-1$,  $\langle\alpha_1, \beta_1\rangle=1$, $\langle\alpha_2, \beta_1\rangle=1$, 
\\
\begin{tabular}{cc}$\alpha_1$+ & $\alpha_2$\\\tiny{ $A_3$+$A_2$ } & \tiny{  }\end{tabular} &\begin{tabular}{c} = \\~ \end{tabular} & \begin{tabular}{cc}$\beta_1$+ & $\beta_2$\\\tiny{ $A_3$+$A_2$ } & \tiny{ $A_3$+$A_2$ }\end{tabular} & $A_3$ & $E_7$ &  $\langle\alpha_1, \beta_1\rangle=1$, $\langle\alpha_1, \beta_2\rangle=1$, $\langle\alpha_2, \beta_1\rangle=1$, $\langle\alpha_2, \beta_2\rangle=1$, 
\\
\begin{tabular}{cc}$\alpha_1$+ & $\alpha_2$\\\tiny{ $A_2$ } & \tiny{  }\end{tabular} &\begin{tabular}{c} = \\~ \end{tabular} & \begin{tabular}{cc}$\beta_1$+ & $\beta_2$\\\tiny{ $A_3$+$A_2$ } & \tiny{ $A_3$ }\end{tabular} & $A_3$ & $E_7$ &  $\langle\alpha_1, \beta_1\rangle=1$, $\langle\alpha_1, \beta_2\rangle=1$, $\langle\alpha_2, \beta_1\rangle=1$, $\langle\alpha_2, \beta_2\rangle=1$, 
\\
\begin{tabular}{cc}$\alpha_1$+ & $\alpha_2$\\\tiny{ $A_3$+$A_2$ } & \tiny{ $A_3$ }\end{tabular} &\begin{tabular}{c} = \\~ \end{tabular} & \begin{tabular}{cc}$\beta_1$+ & $\beta_2$\\\tiny{ $A_3$+$A_2$ } & \tiny{ $A_3$+$A_2$ }\end{tabular} & $A_3$ & $E_7$ &  $\langle\alpha_1, \beta_1\rangle=1$, $\langle\alpha_1, \beta_2\rangle=1$, $\langle\alpha_2, \beta_1\rangle=1$, $\langle\alpha_2, \beta_2\rangle=1$, 
\\
\begin{tabular}{c}2$\alpha_1$\\\tiny{  }\end{tabular} &\begin{tabular}{c} = \\~ \end{tabular} & \begin{tabular}{cccc}$\beta_1$+ & $\beta_2$+ & $\beta_3$+ & $\beta_4$\\\tiny{ $A_3$+$A_2$ } & \tiny{ $A_3$+$A_2$ } & \tiny{ $A_3$+$A_2$ } & \tiny{ $A_3$ }\end{tabular} & $D_4$ & $E_7$ &  $\langle\alpha_1, \beta_1\rangle=1$, $\langle\alpha_1, \beta_2\rangle=1$, $\langle\alpha_1, \beta_3\rangle=1$, $\langle\alpha_1, \beta_4\rangle=1$, 
\\
\begin{tabular}{c}2$\alpha_1$\\\tiny{ $A_3$ }\end{tabular} &\begin{tabular}{c} = \\~ \end{tabular} & \begin{tabular}{cccc}$\beta_1$+ & $\beta_2$+ & $\beta_3$+ & $\beta_4$\\\tiny{ $A_3$+$A_2$ } & \tiny{ $A_3$+$A_2$ } & \tiny{ $A_3$+$A_2$ } & \tiny{ $A_3$ }\end{tabular} & $D_4$ & $E_7$ &  $\langle\alpha_1, \beta_1\rangle=1$, $\langle\alpha_1, \beta_2\rangle=1$, $\langle\alpha_1, \beta_3\rangle=1$, $\langle\alpha_1, \beta_4\rangle=1$, 
\\
\begin{tabular}{cc}$\alpha_1$+ & $\alpha_2$\\\tiny{ $A_3$ } & \tiny{ $A_3$ }\end{tabular} &\begin{tabular}{c} = \\~ \end{tabular} & \begin{tabular}{ccc}$\beta_1$+ & $\beta_2$+ & $\beta_3$\\\tiny{ $A_3$+$A_2$ } & \tiny{ $A_3$+$A_2$ } & \tiny{ $A_3$+$A_2$ }\end{tabular} & $D_4$ & $E_7$ & $\langle\alpha_1, \alpha_2\rangle=1$,  $\langle\alpha_1, \beta_1\rangle=1$, $\langle\alpha_1, \beta_2\rangle=1$, $\langle\alpha_1, \beta_3\rangle=1$, $\langle\alpha_2, \beta_1\rangle=1$, $\langle\alpha_2, \beta_2\rangle=1$, $\langle\alpha_2, \beta_3\rangle=1$, 
\\
\begin{tabular}{cc}$\alpha_1$+ & $\alpha_2$\\\tiny{ $A_3$ } & \tiny{  }\end{tabular} &\begin{tabular}{c} = \\~ \end{tabular} & \begin{tabular}{ccc}$\beta_1$+ & $\beta_2$+ & $\beta_3$\\\tiny{ $A_3$+$A_2$ } & \tiny{ $A_3$+$A_2$ } & \tiny{ $A_3$+$A_2$ }\end{tabular} & $D_4$ & $E_7$ & $\langle\alpha_1, \alpha_2\rangle=1$,  $\langle\alpha_1, \beta_1\rangle=1$, $\langle\alpha_1, \beta_2\rangle=1$, $\langle\alpha_1, \beta_3\rangle=1$, $\langle\alpha_2, \beta_1\rangle=1$, $\langle\alpha_2, \beta_2\rangle=1$, $\langle\alpha_2, \beta_3\rangle=1$, 
\\
\begin{tabular}{cc}$\alpha_1$+ & $\alpha_2$\\\tiny{ $A_2$ } & \tiny{ $A_3$ }\end{tabular} &\begin{tabular}{c} = \\~ \end{tabular} & \begin{tabular}{c}$\beta_1$\\\tiny{ $A_3$+$A_2$ }\end{tabular} & $A_2$ & $A_7$ & $\langle\alpha_1, \alpha_2\rangle=-1$,  $\langle\alpha_1, \beta_1\rangle=1$, $\langle\alpha_2, \beta_1\rangle=1$, 
\\
\begin{tabular}{c}2$\alpha_1$\\\tiny{ $A_3$ }\end{tabular} &\begin{tabular}{c} = \\~ \end{tabular} & \begin{tabular}{cccc}$\beta_1$+ & $\beta_2$+ & $\beta_3$+ & $\beta_4$\\\tiny{ $A_3$+$A_2$ } & \tiny{ $A_3$+$A_2$ } & \tiny{ $A_3$+$A_2$ } & \tiny{  }\end{tabular} & $D_4$ & $E_7$ &  $\langle\alpha_1, \beta_1\rangle=1$, $\langle\alpha_1, \beta_2\rangle=1$, $\langle\alpha_1, \beta_3\rangle=1$, $\langle\alpha_1, \beta_4\rangle=1$, 
\\
\begin{tabular}{cc}$\alpha_1$+ & $\alpha_2$\\\tiny{ $A_3$ } & \tiny{ $A_3$ }\end{tabular} &\begin{tabular}{c} = \\~ \end{tabular} & \begin{tabular}{cc}$\beta_1$+ & $\beta_2$\\\tiny{ $A_3$+$A_2$ } & \tiny{ $A_3$+$A_2$ }\end{tabular} & $A_3$ & $E_7$ &  $\langle\alpha_1, \beta_1\rangle=1$, $\langle\alpha_1, \beta_2\rangle=1$, $\langle\alpha_2, \beta_1\rangle=1$, $\langle\alpha_2, \beta_2\rangle=1$, 
\\
\hline\multicolumn{5}{c}{$\mathfrak{k}$-semisimple type: $A_3$+$A_1$+$A_1$}\\
\hline\hline\begin{tabular}{cc}$\alpha_1$+ & $\alpha_2$\\\tiny{ $A_3$+$A_1$ } & \tiny{  }\end{tabular} &\begin{tabular}{c} = \\~ \end{tabular} & \begin{tabular}{c}$\beta_1$\\\tiny{ $A_3$+$A_1$ }\end{tabular} & $A_2$ & $D_6$+$A_1$ & $\langle\alpha_1, \alpha_2\rangle=-1$,  $\langle\alpha_1, \beta_1\rangle=1$, $\langle\alpha_2, \beta_1\rangle=1$, 
\\
\begin{tabular}{cc}$\alpha_1$+ & $\alpha_2$\\\tiny{ $A_3$+$A_1$ } & \tiny{ $A_1$ }\end{tabular} &\begin{tabular}{c} = \\~ \end{tabular} & \begin{tabular}{cc}$\beta_1$+ & $\beta_2$\\\tiny{ $A_3$+$A_1$+$A_1$' } & \tiny{ $A_3$+$A_1$+$A_1$' }\end{tabular} & $A_3$ & $E_7$ &  $\langle\alpha_1, \beta_1\rangle=1$, $\langle\alpha_1, \beta_2\rangle=1$, $\langle\alpha_2, \beta_1\rangle=1$, $\langle\alpha_2, \beta_2\rangle=1$, 
\\
\begin{tabular}{cc}$\alpha_1$+ & $\alpha_2$\\\tiny{ $A_1$ } & \tiny{  }\end{tabular} &\begin{tabular}{c} = \\~ \end{tabular} & \begin{tabular}{c}$\beta_1$\\\tiny{ $A_1$ }\end{tabular} & $A_2$ & 2$A_3$+$A_1$ & $\langle\alpha_1, \alpha_2\rangle=-1$,  $\langle\alpha_1, \beta_1\rangle=1$, $\langle\alpha_2, \beta_1\rangle=1$, 
\\
\begin{tabular}{cc}$\alpha_1$+ & $\alpha_2$\\\tiny{ $A_3$+$A_1$ } & \tiny{ $A_1$' }\end{tabular} &\begin{tabular}{c} = \\~ \end{tabular} & \begin{tabular}{cc}$\beta_1$+ & $\beta_2$\\\tiny{ $A_3$+$A_1$ } & \tiny{ $A_1$' }\end{tabular} & $A_3$ & $E_7$ &  $\langle\alpha_1, \beta_1\rangle=1$, $\langle\alpha_1, \beta_2\rangle=1$, $\langle\alpha_2, \beta_1\rangle=1$, $\langle\alpha_2, \beta_2\rangle=1$, 
\\
\begin{tabular}{cc}$\alpha_1$+ & $\alpha_2$\\\tiny{ $A_3$+$A_1$ } & \tiny{ $A_1$' }\end{tabular} &\begin{tabular}{c} = \\~ \end{tabular} & \begin{tabular}{cc}$\beta_1$+ & $\beta_2$\\\tiny{ $A_3$+$A_1$+$A_1$' } & \tiny{ $A_3$ }\end{tabular} & $A_3$ & $E_7$ &  $\langle\alpha_1, \beta_1\rangle=1$, $\langle\alpha_1, \beta_2\rangle=1$, $\langle\alpha_2, \beta_1\rangle=1$, $\langle\alpha_2, \beta_2\rangle=1$, 
\\
\begin{tabular}{cc}$\alpha_1$+ & $\alpha_2$\\\tiny{ $A_1$ } & \tiny{ $A_1$ }\end{tabular} &\begin{tabular}{c} = \\~ \end{tabular} & \begin{tabular}{ccc}$\beta_1$+ & $\beta_2$+ & $\beta_3$\\\tiny{ $A_3$+$A_1$+$A_1$' } & \tiny{ $A_3$+$A_1$+$A_1$' } & \tiny{ $A_3$ }\end{tabular} & $D_4$ & $E_7$ & $\langle\alpha_1, \alpha_2\rangle=1$,  $\langle\alpha_1, \beta_1\rangle=1$, $\langle\alpha_1, \beta_2\rangle=1$, $\langle\alpha_1, \beta_3\rangle=1$, $\langle\alpha_2, \beta_1\rangle=1$, $\langle\alpha_2, \beta_2\rangle=1$, $\langle\alpha_2, \beta_3\rangle=1$, 
\\
\begin{tabular}{cc}$\alpha_1$+ & $\alpha_2$\\\tiny{ $A_3$+$A_1$ } & \tiny{  }\end{tabular} &\begin{tabular}{c} = \\~ \end{tabular} & \begin{tabular}{ccc}$\beta_1$+ & $\beta_2$+ & $\beta_3$\\\tiny{ $A_3$+$A_1$ } & \tiny{ $A_3$+$A_1$ } & \tiny{ $A_1$' }\end{tabular} & $D_4$ & $E_7$ & $\langle\alpha_1, \alpha_2\rangle=1$,  $\langle\alpha_1, \beta_1\rangle=1$, $\langle\alpha_1, \beta_2\rangle=1$, $\langle\alpha_1, \beta_3\rangle=1$, $\langle\alpha_2, \beta_1\rangle=1$, $\langle\alpha_2, \beta_2\rangle=1$, $\langle\alpha_2, \beta_3\rangle=1$, 
\\
\begin{tabular}{cc}$\alpha_1$+ & $\alpha_2$\\\tiny{ $A_3$ } & \tiny{  }\end{tabular} &\begin{tabular}{c} = \\~ \end{tabular} & \begin{tabular}{cc}$\beta_1$+ & $\beta_2$\\\tiny{ $A_3$+$A_1$ } & \tiny{ $A_3$+$A_1$ }\end{tabular} & $A_3$ & $D_6$+$A_1$ &  $\langle\alpha_1, \beta_1\rangle=1$, $\langle\alpha_1, \beta_2\rangle=1$, $\langle\alpha_2, \beta_1\rangle=1$, $\langle\alpha_2, \beta_2\rangle=1$, 
\\
\begin{tabular}{cc}$\alpha_1$+ & $\alpha_2$\\\tiny{ $A_1$ } & \tiny{ $A_3$ }\end{tabular} &\begin{tabular}{c} = \\~ \end{tabular} & \begin{tabular}{ccc}$\beta_1$+ & $\beta_2$+ & $\beta_3$\\\tiny{ $A_3$+$A_1$ } & \tiny{ $A_3$+$A_1$ } & \tiny{ $A_1$' }\end{tabular} & $D_4$ & $E_7$ & $\langle\alpha_1, \alpha_2\rangle=1$,  $\langle\alpha_1, \beta_1\rangle=1$, $\langle\alpha_1, \beta_2\rangle=1$, $\langle\alpha_1, \beta_3\rangle=1$, $\langle\alpha_2, \beta_1\rangle=1$, $\langle\alpha_2, \beta_2\rangle=1$, $\langle\alpha_2, \beta_3\rangle=1$, 
\\
\begin{tabular}{cc}$\alpha_1$+ & $\alpha_2$\\\tiny{ $A_3$+$A_1$ } & \tiny{ $A_1$' }\end{tabular} &\begin{tabular}{c} = \\~ \end{tabular} & \begin{tabular}{c}$\beta_1$\\\tiny{ $A_3$+$A_1$+$A_1$' }\end{tabular} & $A_2$ & $E_7$ & $\langle\alpha_1, \alpha_2\rangle=-1$,  $\langle\alpha_1, \beta_1\rangle=1$, $\langle\alpha_2, \beta_1\rangle=1$, 
\\
\begin{tabular}{c}2$\alpha_1$\\\tiny{ $A_1$ }\end{tabular} &\begin{tabular}{c} = \\~ \end{tabular} & \begin{tabular}{cccc}$\beta_1$+ & $\beta_2$+ & $\beta_3$+ & $\beta_4$\\\tiny{ $A_3$+$A_1$+$A_1$' } & \tiny{ $A_3$+$A_1$+$A_1$' } & \tiny{ $A_3$ } & \tiny{  }\end{tabular} & $D_4$ & $E_7$ &  $\langle\alpha_1, \beta_1\rangle=1$, $\langle\alpha_1, \beta_2\rangle=1$, $\langle\alpha_1, \beta_3\rangle=1$, $\langle\alpha_1, \beta_4\rangle=1$, 
\\
\begin{tabular}{cc}$\alpha_1$+ & $\alpha_2$\\\tiny{ $A_1$ } & \tiny{ $A_3$ }\end{tabular} &\begin{tabular}{c} = \\~ \end{tabular} & \begin{tabular}{c}$\beta_1$\\\tiny{ $A_3$+$A_1$ }\end{tabular} & $A_2$ & $D_6$+$A_1$ & $\langle\alpha_1, \alpha_2\rangle=-1$,  $\langle\alpha_1, \beta_1\rangle=1$, $\langle\alpha_2, \beta_1\rangle=1$, 
\\
\begin{tabular}{cc}$\alpha_1$+ & $\alpha_2$\\\tiny{ $A_3$+$A_1$ } & \tiny{ $A_1$' }\end{tabular} &\begin{tabular}{c} = \\~ \end{tabular} & \begin{tabular}{cc}$\beta_1$+ & $\beta_2$\\\tiny{ $A_3$+$A_1$ } & \tiny{ $A_3$+$A_1$' }\end{tabular} & $A_3$ & $E_7$ &  $\langle\alpha_1, \beta_1\rangle=1$, $\langle\alpha_1, \beta_2\rangle=1$, $\langle\alpha_2, \beta_1\rangle=1$, $\langle\alpha_2, \beta_2\rangle=1$, 
\\
\begin{tabular}{cc}$\alpha_1$+ & $\alpha_2$\\\tiny{ $A_1$ } & \tiny{ $A_1$ }\end{tabular} &\begin{tabular}{c} = \\~ \end{tabular} & \begin{tabular}{cc}$\beta_1$+ & $\beta_2$\\\tiny{ $A_3$+$A_1$ } & \tiny{ $A_3$+$A_1$ }\end{tabular} & $A_3$ & $D_6$+$A_1$ &  $\langle\alpha_1, \beta_1\rangle=1$, $\langle\alpha_1, \beta_2\rangle=1$, $\langle\alpha_2, \beta_1\rangle=1$, $\langle\alpha_2, \beta_2\rangle=1$, 
\\
\hline\multicolumn{5}{c}{$\mathfrak{k}$-semisimple type: $A_3$+$A_1$+$A_1$}\\
\hline\hline\begin{tabular}{c}2$\alpha_1$\\\tiny{  }\end{tabular} &\begin{tabular}{c} = \\~ \end{tabular} & \begin{tabular}{cccc}$\beta_1$+ & $\beta_2$+ & $\beta_3$+ & $\beta_4$\\\tiny{ $A_3$+$A_1$ } & \tiny{ $A_3$+$A_1$ } & \tiny{ $A_3$+$A_1$' } & \tiny{ $A_3$+$A_1$' }\end{tabular} & $D_4$ & $E_7$ &  $\langle\alpha_1, \beta_1\rangle=1$, $\langle\alpha_1, \beta_2\rangle=1$, $\langle\alpha_1, \beta_3\rangle=1$, $\langle\alpha_1, \beta_4\rangle=1$, 
\\
\begin{tabular}{cc}$\alpha_1$+ & $\alpha_2$\\\tiny{ $A_1$+$A_1$' } & \tiny{  }\end{tabular} &\begin{tabular}{c} = \\~ \end{tabular} & \begin{tabular}{cc}$\beta_1$+ & $\beta_2$\\\tiny{ $A_3$+$A_1$ } & \tiny{ $A_3$+$A_1$' }\end{tabular} & $A_3$ & $E_7$ &  $\langle\alpha_1, \beta_1\rangle=1$, $\langle\alpha_1, \beta_2\rangle=1$, $\langle\alpha_2, \beta_1\rangle=1$, $\langle\alpha_2, \beta_2\rangle=1$, 
\\
\begin{tabular}{cc}$\alpha_1$+ & $\alpha_2$\\\tiny{ $A_3$+$A_1$ } & \tiny{ $A_3$+$A_1$' }\end{tabular} &\begin{tabular}{c} = \\~ \end{tabular} & \begin{tabular}{cc}$\beta_1$+ & $\beta_2$\\\tiny{ $A_3$+$A_1$ } & \tiny{ $A_3$+$A_1$' }\end{tabular} & $A_3$ & $E_7$ &  $\langle\alpha_1, \beta_1\rangle=1$, $\langle\alpha_1, \beta_2\rangle=1$, $\langle\alpha_2, \beta_1\rangle=1$, $\langle\alpha_2, \beta_2\rangle=1$, 
\\
\begin{tabular}{c}2$\alpha_1$\\\tiny{ $A_3$ }\end{tabular} &\begin{tabular}{c} = \\~ \end{tabular} & \begin{tabular}{cccc}$\beta_1$+ & $\beta_2$+ & $\beta_3$+ & $\beta_4$\\\tiny{ $A_3$+$A_1$ } & \tiny{ $A_3$+$A_1$ } & \tiny{ $A_3$+$A_1$' } & \tiny{ $A_3$+$A_1$' }\end{tabular} & $D_4$ & $E_7$ &  $\langle\alpha_1, \beta_1\rangle=1$, $\langle\alpha_1, \beta_2\rangle=1$, $\langle\alpha_1, \beta_3\rangle=1$, $\langle\alpha_1, \beta_4\rangle=1$, 
\\
\hline\multicolumn{5}{c}{$\mathfrak{k}$-semisimple type: $A_2$+$A_2$+$A_1$}\\
\hline\hline\begin{tabular}{cc}$\alpha_1$+ & $\alpha_2$\\\tiny{ $A_2$+$A_1$ } & \tiny{ $A_2$' }\end{tabular} &\begin{tabular}{c} = \\~ \end{tabular} & \begin{tabular}{c}$\beta_1$\\\tiny{ $A_2$+$A_2$'+$A_1$ }\end{tabular} & $A_2$ & $E_7$ & $\langle\alpha_1, \alpha_2\rangle=-1$,  $\langle\alpha_1, \beta_1\rangle=1$, $\langle\alpha_2, \beta_1\rangle=1$, 
\\
\begin{tabular}{cc}$\alpha_1$+ & $\alpha_2$\\\tiny{ $A_1$ } & \tiny{ $A_2$ }\end{tabular} &\begin{tabular}{c} = \\~ \end{tabular} & \begin{tabular}{c}$\beta_1$\\\tiny{ $A_2$+$A_1$ }\end{tabular} & $A_2$ & $A_5$+$A_2$ & $\langle\alpha_1, \alpha_2\rangle=-1$,  $\langle\alpha_1, \beta_1\rangle=1$, $\langle\alpha_2, \beta_1\rangle=1$, 
\\
\begin{tabular}{cc}$\alpha_1$+ & $\alpha_2$\\\tiny{ $A_2$ } & \tiny{ $A_2$' }\end{tabular} &\begin{tabular}{c} = \\~ \end{tabular} & \begin{tabular}{cc}$\beta_1$+ & $\beta_2$\\\tiny{ $A_2$+$A_2$' } & \tiny{ $A_2$+$A_2$' }\end{tabular} & $A_3$ & $D_6$+$A_1$ &  $\langle\alpha_1, \beta_1\rangle=1$, $\langle\alpha_1, \beta_2\rangle=1$, $\langle\alpha_2, \beta_1\rangle=1$, $\langle\alpha_2, \beta_2\rangle=1$, 
\\
\begin{tabular}{cc}$\alpha_1$+ & $\alpha_2$\\\tiny{ $A_2$ } & \tiny{ $A_2$' }\end{tabular} &\begin{tabular}{c} = \\~ \end{tabular} & \begin{tabular}{cc}$\beta_1$+ & $\beta_2$\\\tiny{ $A_2$+$A_1$ } & \tiny{ $A_2$'+$A_1$ }\end{tabular} & $A_3$ & $E_7$ &  $\langle\alpha_1, \beta_1\rangle=1$, $\langle\alpha_1, \beta_2\rangle=1$, $\langle\alpha_2, \beta_1\rangle=1$, $\langle\alpha_2, \beta_2\rangle=1$, 
\\
\begin{tabular}{cc}$\alpha_1$+ & $\alpha_2$\\\tiny{ $A_2$+$A_2$' } & \tiny{ $A_2$ }\end{tabular} &\begin{tabular}{c} = \\~ \end{tabular} & \begin{tabular}{cc}$\beta_1$+ & $\beta_2$\\\tiny{ $A_2$+$A_2$'+$A_1$ } & \tiny{ $A_2$+$A_2$'+$A_1$ }\end{tabular} & $A_3$ & $E_7$ &  $\langle\alpha_1, \beta_1\rangle=1$, $\langle\alpha_1, \beta_2\rangle=1$, $\langle\alpha_2, \beta_1\rangle=1$, $\langle\alpha_2, \beta_2\rangle=1$, 
\\
\begin{tabular}{c}2$\alpha_1$\\\tiny{ $A_1$ }\end{tabular} &\begin{tabular}{c} = \\~ \end{tabular} & \begin{tabular}{cccc}$\beta_1$+ & $\beta_2$+ & $\beta_3$+ & $\beta_4$\\\tiny{ $A_2$+$A_1$ } & \tiny{ $A_2$'+$A_1$ } & \tiny{ $A_2$'+$A_2$ } & \tiny{ $A_2$'+$A_2$ }\end{tabular} & $D_4$ & $E_7$ &  $\langle\alpha_1, \beta_1\rangle=1$, $\langle\alpha_1, \beta_2\rangle=1$, $\langle\alpha_1, \beta_3\rangle=1$, $\langle\alpha_1, \beta_4\rangle=1$, 
\\
\begin{tabular}{cc}$\alpha_1$+ & $\alpha_2$\\\tiny{ $A_2$+$A_2$' } & \tiny{ $A_1$ }\end{tabular} &\begin{tabular}{c} = \\~ \end{tabular} & \begin{tabular}{cc}$\beta_1$+ & $\beta_2$\\\tiny{ $A_2$+$A_1$ } & \tiny{ $A_2$+$A_2$' }\end{tabular} & $A_3$ & $E_7$ &  $\langle\alpha_1, \beta_1\rangle=1$, $\langle\alpha_1, \beta_2\rangle=1$, $\langle\alpha_2, \beta_1\rangle=1$, $\langle\alpha_2, \beta_2\rangle=1$, 
\\
\begin{tabular}{c}2$\alpha_1$\\\tiny{ $A_2$ }\end{tabular} &\begin{tabular}{c} = \\~ \end{tabular} & \begin{tabular}{cccc}$\beta_1$+ & $\beta_2$+ & $\beta_3$+ & $\beta_4$\\\tiny{ $A_2$+$A_1$ } & \tiny{ $A_2$'+$A_1$ } & \tiny{ $A_2$'+$A_2$ } & \tiny{ $A_2$'+$A_2$ }\end{tabular} & $D_4$ & $E_7$ &  $\langle\alpha_1, \beta_1\rangle=1$, $\langle\alpha_1, \beta_2\rangle=1$, $\langle\alpha_1, \beta_3\rangle=1$, $\langle\alpha_1, \beta_4\rangle=1$, 
\\
\begin{tabular}{cc}$\alpha_1$+ & $\alpha_2$\\\tiny{ $A_1$ } & \tiny{ $A_2$ }\end{tabular} &\begin{tabular}{c} = \\~ \end{tabular} & \begin{tabular}{ccc}$\beta_1$+ & $\beta_2$+ & $\beta_3$\\\tiny{ $A_2$+$A_1$ } & \tiny{ $A_2$+$A_2$' } & \tiny{ $A_2$+$A_2$' }\end{tabular} & $D_4$ & $E_7$ & $\langle\alpha_1, \alpha_2\rangle=1$,  $\langle\alpha_1, \beta_1\rangle=1$, $\langle\alpha_1, \beta_2\rangle=1$, $\langle\alpha_1, \beta_3\rangle=1$, $\langle\alpha_2, \beta_1\rangle=1$, $\langle\alpha_2, \beta_2\rangle=1$, $\langle\alpha_2, \beta_3\rangle=1$, 
\\
\begin{tabular}{cc}$\alpha_1$+ & $\alpha_2$\\\tiny{ $A_2$+$A_1$ } & \tiny{ $A_2$' }\end{tabular} &\begin{tabular}{c} = \\~ \end{tabular} & \begin{tabular}{ccc}$\beta_1$+ & $\beta_2$+ & $\beta_3$\\\tiny{ $A_2$+$A_2$' } & \tiny{ $A_2$+$A_2$' } & \tiny{ $A_1$ }\end{tabular} & $D_4$ & $E_7$ & $\langle\alpha_1, \alpha_2\rangle=1$,  $\langle\alpha_1, \beta_1\rangle=1$, $\langle\alpha_1, \beta_2\rangle=1$, $\langle\alpha_1, \beta_3\rangle=1$, $\langle\alpha_2, \beta_1\rangle=1$, $\langle\alpha_2, \beta_2\rangle=1$, $\langle\alpha_2, \beta_3\rangle=1$, 
\\
\begin{tabular}{cc}$\alpha_1$+ & $\alpha_2$\\\tiny{ $A_2$+$A_1$ } & \tiny{ $A_2$'+$A_1$ }\end{tabular} &\begin{tabular}{c} = \\~ \end{tabular} & \begin{tabular}{cc}$\beta_1$+ & $\beta_2$\\\tiny{ $A_2$+$A_2$'+$A_1$ } & \tiny{ $A_2$+$A_2$'+$A_1$ }\end{tabular} & $A_3$ & $E_7$ &  $\langle\alpha_1, \beta_1\rangle=1$, $\langle\alpha_1, \beta_2\rangle=1$, $\langle\alpha_2, \beta_1\rangle=1$, $\langle\alpha_2, \beta_2\rangle=1$, 
\\
\begin{tabular}{cc}$\alpha_1$+ & $\alpha_2$\\\tiny{ $A_2$+$A_1$ } & \tiny{ $A_2$ }\end{tabular} &\begin{tabular}{c} = \\~ \end{tabular} & \begin{tabular}{cc}$\beta_1$+ & $\beta_2$\\\tiny{ $A_2$+$A_2$'+$A_1$ } & \tiny{ $A_2$+$A_2$' }\end{tabular} & $A_3$ & $E_7$ &  $\langle\alpha_1, \beta_1\rangle=1$, $\langle\alpha_1, \beta_2\rangle=1$, $\langle\alpha_2, \beta_1\rangle=1$, $\langle\alpha_2, \beta_2\rangle=1$, 
\\
\begin{tabular}{cc}$\alpha_1$+ & $\alpha_2$\\\tiny{ $A_2$ } & \tiny{ $A_2$' }\end{tabular} &\begin{tabular}{c} = \\~ \end{tabular} & \begin{tabular}{cc}$\beta_1$+ & $\beta_2$\\\tiny{ $A_2$+$A_2$'+$A_1$ } & \tiny{ $A_1$ }\end{tabular} & $A_3$ & $E_7$ &  $\langle\alpha_1, \beta_1\rangle=1$, $\langle\alpha_1, \beta_2\rangle=1$, $\langle\alpha_2, \beta_1\rangle=1$, $\langle\alpha_2, \beta_2\rangle=1$, 
\\
\hline\multicolumn{5}{c}{$\mathfrak{k}$-semisimple type: $A_2$+$A_1$+$A_1$+$A_1$}\\
\hline\hline\begin{tabular}{cc}$\alpha_1$+ & $\alpha_2$\\\tiny{ $A_1$+$A_1$' } & \tiny{ $A_2$ }\end{tabular} &\begin{tabular}{c} = \\~ \end{tabular} & \begin{tabular}{c}$\beta_1$\\\tiny{ $A_2$+$A_1$+$A_1$' }\end{tabular} & $A_2$ & $D_6$+$A_1$ & $\langle\alpha_1, \alpha_2\rangle=-1$,  $\langle\alpha_1, \beta_1\rangle=1$, $\langle\alpha_2, \beta_1\rangle=1$, 
\\
\begin{tabular}{cc}$\alpha_1$+ & $\alpha_2$\\\tiny{ $A_1$+$A_1$' } & \tiny{ $A_1$+$A_1$'' }\end{tabular} &\begin{tabular}{c} = \\~ \end{tabular} & \begin{tabular}{cc}$\beta_1$+ & $\beta_2$\\\tiny{ $A_2$+$A_1$+$A_1$' } & \tiny{ $A_2$+$A_1$+$A_1$'' }\end{tabular} & $A_3$ & $E_7$ &  $\langle\alpha_1, \beta_1\rangle=1$, $\langle\alpha_1, \beta_2\rangle=1$, $\langle\alpha_2, \beta_1\rangle=1$, $\langle\alpha_2, \beta_2\rangle=1$, 
\\
\begin{tabular}{cc}$\alpha_1$+ & $\alpha_2$\\\tiny{ $A_2$ } & \tiny{ $A_2$ }\end{tabular} &\begin{tabular}{c} = \\~ \end{tabular} & \begin{tabular}{cc}$\beta_1$+ & $\beta_2$\\\tiny{ $A_2$+$A_1$+$A_1$' } & \tiny{ $A_2$+$A_1$+$A_1$' }\end{tabular} & $A_3$ & $D_6$+$A_1$ &  $\langle\alpha_1, \beta_1\rangle=1$, $\langle\alpha_1, \beta_2\rangle=1$, $\langle\alpha_2, \beta_1\rangle=1$, $\langle\alpha_2, \beta_2\rangle=1$, 
\\
\begin{tabular}{cc}$\alpha_1$+ & $\alpha_2$\\\tiny{ $A_1$+$A_1$' } & \tiny{ $A_2$ }\end{tabular} &\begin{tabular}{c} = \\~ \end{tabular} & \begin{tabular}{cc}$\beta_1$+ & $\beta_2$\\\tiny{ $A_2$+$A_1$+$A_1$' } & \tiny{ $A_1$'+$A_1$'' }\end{tabular} & $A_3$ & $E_7$ &  $\langle\alpha_1, \beta_1\rangle=1$, $\langle\alpha_1, \beta_2\rangle=1$, $\langle\alpha_2, \beta_1\rangle=1$, $\langle\alpha_2, \beta_2\rangle=1$, 
\\
\hline\multicolumn{5}{c}{$\mathfrak{k}$-semisimple type: $A_1$+$A_1$+$A_1$+$A_1$+$A_1$}\\
\hline\hline\begin{tabular}{c}2$\alpha_1$\\\tiny{  }\end{tabular} &\begin{tabular}{c} = \\~ \end{tabular} & \begin{tabular}{cccc}$\beta_1$+ & $\beta_2$+ & $\beta_3$+ & $\beta_4$\\\tiny{ $A_1$+$A_1$' } & \tiny{ $A_1$+$A_1$' } & \tiny{ $A_1$''+$A_1$''' } & \tiny{ $A_1$''+$A_1$''' }\end{tabular} & $D_4$ & $D_6$+$A_1$ &  $\langle\alpha_1, \beta_1\rangle=1$, $\langle\alpha_1, \beta_2\rangle=1$, $\langle\alpha_1, \beta_3\rangle=1$, $\langle\alpha_1, \beta_4\rangle=1$, 
\\
\begin{tabular}{cc}$\alpha_1$+ & $\alpha_2$\\\tiny{  } & \tiny{  }\end{tabular} &\begin{tabular}{c} = \\~ \end{tabular} & \begin{tabular}{cc}$\beta_1$+ & $\beta_2$\\\tiny{ $A_1$+$A_1$' } & \tiny{ $A_1$+$A_1$' }\end{tabular} & $A_3$ & $D_4$+3$A_1$ &  $\langle\alpha_1, \beta_1\rangle=1$, $\langle\alpha_1, \beta_2\rangle=1$, $\langle\alpha_2, \beta_1\rangle=1$, $\langle\alpha_2, \beta_2\rangle=1$, 
\\
\begin{tabular}{cc}$\alpha_1$+ & $\alpha_2$\\\tiny{ $A_1$+$A_1$' } & \tiny{  }\end{tabular} &\begin{tabular}{c} = \\~ \end{tabular} & \begin{tabular}{c}$\beta_1$\\\tiny{ $A_1$+$A_1$' }\end{tabular} & $A_2$ & $D_4$+3$A_1$ & $\langle\alpha_1, \alpha_2\rangle=-1$,  $\langle\alpha_1, \beta_1\rangle=1$, $\langle\alpha_2, \beta_1\rangle=1$, 
\\
\begin{tabular}{cc}$\alpha_1$+ & $\alpha_2$\\\tiny{ $A_1$+$A_1$' } & \tiny{ $A_1$''+$A_1$''' }\end{tabular} &\begin{tabular}{c} = \\~ \end{tabular} & \begin{tabular}{cc}$\beta_1$+ & $\beta_2$\\\tiny{ $A_1$+$A_1$'+$A_1$'' } & \tiny{ $A_1$'''+$A_1$''''+$A_1$' }\end{tabular} & $A_3$ & $E_7$ &  $\langle\alpha_1, \beta_1\rangle=1$, $\langle\alpha_1, \beta_2\rangle=1$, $\langle\alpha_2, \beta_1\rangle=1$, $\langle\alpha_2, \beta_2\rangle=1$, 
\\
\begin{tabular}{cc}$\alpha_1$+ & $\alpha_2$\\\tiny{ $A_1$+$A_1$' } & \tiny{ $A_1$''+$A_1$''' }\end{tabular} &\begin{tabular}{c} = \\~ \end{tabular} & \begin{tabular}{cc}$\beta_1$+ & $\beta_2$\\\tiny{ $A_1$+$A_1$' } & \tiny{ $A_1$''+$A_1$''' }\end{tabular} & $A_3$ & $D_6$+$A_1$ &  $\langle\alpha_1, \beta_1\rangle=1$, $\langle\alpha_1, \beta_2\rangle=1$, $\langle\alpha_2, \beta_1\rangle=1$, $\langle\alpha_2, \beta_2\rangle=1$, 
\\
\begin{tabular}{cc}$\alpha_1$+ & $\alpha_2$\\\tiny{ $A_1$+$A_1$' } & \tiny{ $A_1$''+$A_1$''' }\end{tabular} &\begin{tabular}{c} = \\~ \end{tabular} & \begin{tabular}{ccc}$\beta_1$+ & $\beta_2$+ & $\beta_3$\\\tiny{ $A_1$+$A_1$'+$A_1$'' } & \tiny{ $A_1$'''+$A_1$''''+$A_1$' } & \tiny{  }\end{tabular} & $D_4$ & $E_7$ & $\langle\alpha_1, \alpha_2\rangle=1$,  $\langle\alpha_1, \beta_1\rangle=1$, $\langle\alpha_1, \beta_2\rangle=1$, $\langle\alpha_1, \beta_3\rangle=1$, $\langle\alpha_2, \beta_1\rangle=1$, $\langle\alpha_2, \beta_2\rangle=1$, $\langle\alpha_2, \beta_3\rangle=1$, 
\\
\hline\multicolumn{5}{c}{$\mathfrak{k}$-semisimple type: $D_4$}\\
\hline\hline\begin{tabular}{cc}$\alpha_1$+ & $\alpha_2$\\\tiny{  } & \tiny{  }\end{tabular} &\begin{tabular}{c} = \\~ \end{tabular} & \begin{tabular}{cc}$\beta_1$+ & $\beta_2$\\\tiny{ $D_4$ } & \tiny{ $D_4$ }\end{tabular} & $A_3$ & $D_6$ &  $\langle\alpha_1, \beta_1\rangle=1$, $\langle\alpha_1, \beta_2\rangle=1$, $\langle\alpha_2, \beta_1\rangle=1$, $\langle\alpha_2, \beta_2\rangle=1$, 
\\
\begin{tabular}{cc}$\alpha_1$+ & $\alpha_2$\\\tiny{ $D_4$ } & \tiny{  }\end{tabular} &\begin{tabular}{c} = \\~ \end{tabular} & \begin{tabular}{cc}$\beta_1$+ & $\beta_2$\\\tiny{ $D_4$ } & \tiny{ $D_4$ }\end{tabular} & $A_3$ & $E_7$ &  $\langle\alpha_1, \beta_1\rangle=1$, $\langle\alpha_1, \beta_2\rangle=1$, $\langle\alpha_2, \beta_1\rangle=1$, $\langle\alpha_2, \beta_2\rangle=1$, 
\\
\begin{tabular}{cc}$\alpha_1$+ & $\alpha_2$\\\tiny{ $D_4$ } & \tiny{  }\end{tabular} &\begin{tabular}{c} = \\~ \end{tabular} & \begin{tabular}{c}$\beta_1$\\\tiny{ $D_4$ }\end{tabular} & $A_2$ & $D_6$ & $\langle\alpha_1, \alpha_2\rangle=-1$,  $\langle\alpha_1, \beta_1\rangle=1$, $\langle\alpha_2, \beta_1\rangle=1$, 
\\
\begin{tabular}{cc}$\alpha_1$+ & $\alpha_2$\\\tiny{ $D_4$ } & \tiny{ $D_4$ }\end{tabular} &\begin{tabular}{c} = \\~ \end{tabular} & \begin{tabular}{cc}$\beta_1$+ & $\beta_2$\\\tiny{ $D_4$ } & \tiny{ $D_4$ }\end{tabular} & $A_3$ & $E_7$ &  $\langle\alpha_1, \beta_1\rangle=1$, $\langle\alpha_1, \beta_2\rangle=1$, $\langle\alpha_2, \beta_1\rangle=1$, $\langle\alpha_2, \beta_2\rangle=1$, 
\\
\hline\multicolumn{5}{c}{$\mathfrak{k}$-semisimple type: $A_4$}\\
\hline\hline\begin{tabular}{cc}$\alpha_1$+ & $\alpha_2$\\\tiny{ $A_4$ } & \tiny{  }\end{tabular} &\begin{tabular}{c} = \\~ \end{tabular} & \begin{tabular}{c}$\beta_1$\\\tiny{ $A_4$ }\end{tabular} & $A_2$ & $E_6$ & $\langle\alpha_1, \alpha_2\rangle=-1$,  $\langle\alpha_1, \beta_1\rangle=1$, $\langle\alpha_2, \beta_1\rangle=1$, 
\\
\begin{tabular}{cc}$\alpha_1$+ & $\alpha_2$\\\tiny{ $A_4$ } & \tiny{  }\end{tabular} &\begin{tabular}{c} = \\~ \end{tabular} & \begin{tabular}{cc}$\beta_1$+ & $\beta_2$\\\tiny{ $A_4$ } & \tiny{ $A_4$ }\end{tabular} & $A_3$ & $E_6$ &  $\langle\alpha_1, \beta_1\rangle=1$, $\langle\alpha_1, \beta_2\rangle=1$, $\langle\alpha_2, \beta_1\rangle=1$, $\langle\alpha_2, \beta_2\rangle=1$, 
\\
\begin{tabular}{cc}$\alpha_1$+ & $\alpha_2$\\\tiny{ $A_4$ } & \tiny{ $A_4$ }\end{tabular} &\begin{tabular}{c} = \\~ \end{tabular} & \begin{tabular}{cc}$\beta_1$+ & $\beta_2$\\\tiny{ $A_4$ } & \tiny{ $A_4$ }\end{tabular} & $A_3$ & $E_7$ &  $\langle\alpha_1, \beta_1\rangle=1$, $\langle\alpha_1, \beta_2\rangle=1$, $\langle\alpha_2, \beta_1\rangle=1$, $\langle\alpha_2, \beta_2\rangle=1$, 
\\
\begin{tabular}{cc}$\alpha_1$+ & $\alpha_2$\\\tiny{ $A_4$ } & \tiny{  }\end{tabular} &\begin{tabular}{c} = \\~ \end{tabular} & \begin{tabular}{c}$\beta_1$\\\tiny{ $A_4$ }\end{tabular} & $A_2$ & $A_6$ & $\langle\alpha_1, \alpha_2\rangle=-1$,  $\langle\alpha_1, \beta_1\rangle=1$, $\langle\alpha_2, \beta_1\rangle=1$, 
\\
\begin{tabular}{cc}$\alpha_1$+ & $\alpha_2$\\\tiny{ $A_4$ } & \tiny{  }\end{tabular} &\begin{tabular}{c} = \\~ \end{tabular} & \begin{tabular}{cc}$\beta_1$+ & $\beta_2$\\\tiny{ $A_4$ } & \tiny{ $A_4$ }\end{tabular} & $A_3$ & $E_7$ &  $\langle\alpha_1, \beta_1\rangle=1$, $\langle\alpha_1, \beta_2\rangle=1$, $\langle\alpha_2, \beta_1\rangle=1$, $\langle\alpha_2, \beta_2\rangle=1$, 
\\
\begin{tabular}{cc}$\alpha_1$+ & $\alpha_2$\\\tiny{  } & \tiny{  }\end{tabular} &\begin{tabular}{c} = \\~ \end{tabular} & \begin{tabular}{ccc}$\beta_1$+ & $\beta_2$+ & $\beta_3$\\\tiny{ $A_4$ } & \tiny{ $A_4$ } & \tiny{ $A_4$ }\end{tabular} & $D_4$ & $E_7$ & $\langle\alpha_1, \alpha_2\rangle=1$,  $\langle\alpha_1, \beta_1\rangle=1$, $\langle\alpha_1, \beta_2\rangle=1$, $\langle\alpha_1, \beta_3\rangle=1$, $\langle\alpha_2, \beta_1\rangle=1$, $\langle\alpha_2, \beta_2\rangle=1$, $\langle\alpha_2, \beta_3\rangle=1$, 
\\
\begin{tabular}{cc}$\alpha_1$+ & $\alpha_2$\\\tiny{ $A_4$ } & \tiny{ $A_4$ }\end{tabular} &\begin{tabular}{c} = \\~ \end{tabular} & \begin{tabular}{ccc}$\beta_1$+ & $\beta_2$+ & $\beta_3$\\\tiny{ $A_4$ } & \tiny{ $A_4$ } & \tiny{ $A_4$ }\end{tabular} & $D_4$ & $E_7$ & $\langle\alpha_1, \alpha_2\rangle=1$,  $\langle\alpha_1, \beta_1\rangle=1$, $\langle\alpha_1, \beta_2\rangle=1$, $\langle\alpha_1, \beta_3\rangle=1$, $\langle\alpha_2, \beta_1\rangle=1$, $\langle\alpha_2, \beta_2\rangle=1$, $\langle\alpha_2, \beta_3\rangle=1$, 
\\
\hline\multicolumn{5}{c}{$\mathfrak{k}$-semisimple type: $A_3$+$A_1$}\\
\hline\hline\begin{tabular}{cc}$\alpha_1$+ & $\alpha_2$\\\tiny{ $A_3$+$A_1$ } & \tiny{  }\end{tabular} &\begin{tabular}{c} = \\~ \end{tabular} & \begin{tabular}{cc}$\beta_1$+ & $\beta_2$\\\tiny{ $A_3$+$A_1$ } & \tiny{ $A_3$ }\end{tabular} & $A_3$ & $E_7$ &  $\langle\alpha_1, \beta_1\rangle=1$, $\langle\alpha_1, \beta_2\rangle=1$, $\langle\alpha_2, \beta_1\rangle=1$, $\langle\alpha_2, \beta_2\rangle=1$, 
\\
\begin{tabular}{cc}$\alpha_1$+ & $\alpha_2$\\\tiny{ $A_1$ } & \tiny{  }\end{tabular} &\begin{tabular}{c} = \\~ \end{tabular} & \begin{tabular}{cc}$\beta_1$+ & $\beta_2$\\\tiny{ $A_3$+$A_1$ } & \tiny{ $A_3$ }\end{tabular} & $A_3$ & $E_7$ &  $\langle\alpha_1, \beta_1\rangle=1$, $\langle\alpha_1, \beta_2\rangle=1$, $\langle\alpha_2, \beta_1\rangle=1$, $\langle\alpha_2, \beta_2\rangle=1$, 
\\
\begin{tabular}{cc}$\alpha_1$+ & $\alpha_2$\\\tiny{ $A_1$ } & \tiny{  }\end{tabular} &\begin{tabular}{c} = \\~ \end{tabular} & \begin{tabular}{c}$\beta_1$\\\tiny{ $A_1$ }\end{tabular} & $A_2$ & 2$A_3$ & $\langle\alpha_1, \alpha_2\rangle=-1$,  $\langle\alpha_1, \beta_1\rangle=1$, $\langle\alpha_2, \beta_1\rangle=1$, 
\\
\begin{tabular}{cc}$\alpha_1$+ & $\alpha_2$\\\tiny{ $A_3$+$A_1$ } & \tiny{  }\end{tabular} &\begin{tabular}{c} = \\~ \end{tabular} & \begin{tabular}{cc}$\beta_1$+ & $\beta_2$\\\tiny{ $A_1$ } & \tiny{ $A_3$ }\end{tabular} & $A_3$ & $A_7$ &  $\langle\alpha_1, \beta_1\rangle=1$, $\langle\alpha_1, \beta_2\rangle=1$, $\langle\alpha_2, \beta_1\rangle=1$, $\langle\alpha_2, \beta_2\rangle=1$, 
\\
\begin{tabular}{cc}$\alpha_1$+ & $\alpha_2$\\\tiny{  } & \tiny{  }\end{tabular} &\begin{tabular}{c} = \\~ \end{tabular} & \begin{tabular}{cc}$\beta_1$+ & $\beta_2$\\\tiny{ $A_3$ } & \tiny{ $A_3$ }\end{tabular} & $A_3$ & $D_6$+$A_1$ &  $\langle\alpha_1, \beta_1\rangle=1$, $\langle\alpha_1, \beta_2\rangle=1$, $\langle\alpha_2, \beta_1\rangle=1$, $\langle\alpha_2, \beta_2\rangle=1$, 
\\
\begin{tabular}{cc}$\alpha_1$+ & $\alpha_2$\\\tiny{ $A_3$ } & \tiny{  }\end{tabular} &\begin{tabular}{c} = \\~ \end{tabular} & \begin{tabular}{c}$\beta_1$\\\tiny{ $A_3$ }\end{tabular} & $A_2$ & $A_5$+$A_1$ & $\langle\alpha_1, \alpha_2\rangle=-1$,  $\langle\alpha_1, \beta_1\rangle=1$, $\langle\alpha_2, \beta_1\rangle=1$, 
\\
\begin{tabular}{cc}$\alpha_1$+ & $\alpha_2$\\\tiny{ $A_3$+$A_1$ } & \tiny{  }\end{tabular} &\begin{tabular}{c} = \\~ \end{tabular} & \begin{tabular}{c}$\beta_1$\\\tiny{ $A_3$+$A_1$ }\end{tabular} & $A_2$ & $D_6$ & $\langle\alpha_1, \alpha_2\rangle=-1$,  $\langle\alpha_1, \beta_1\rangle=1$, $\langle\alpha_2, \beta_1\rangle=1$, 
\\
\begin{tabular}{cc}$\alpha_1$+ & $\alpha_2$\\\tiny{ $A_3$ } & \tiny{ $A_3$ }\end{tabular} &\begin{tabular}{c} = \\~ \end{tabular} & \begin{tabular}{cc}$\beta_1$+ & $\beta_2$\\\tiny{ $A_3$ } & \tiny{ $A_3$ }\end{tabular} & $A_3$ & $D_6$+$A_1$ &  $\langle\alpha_1, \beta_1\rangle=1$, $\langle\alpha_1, \beta_2\rangle=1$, $\langle\alpha_2, \beta_1\rangle=1$, $\langle\alpha_2, \beta_2\rangle=1$, 
\\
\begin{tabular}{cc}$\alpha_1$+ & $\alpha_2$\\\tiny{ $A_3$+$A_1$ } & \tiny{ $A_3$ }\end{tabular} &\begin{tabular}{c} = \\~ \end{tabular} & \begin{tabular}{cc}$\beta_1$+ & $\beta_2$\\\tiny{ $A_3$+$A_1$ } & \tiny{ $A_3$ }\end{tabular} & $A_3$ & $E_7$ &  $\langle\alpha_1, \beta_1\rangle=1$, $\langle\alpha_1, \beta_2\rangle=1$, $\langle\alpha_2, \beta_1\rangle=1$, $\langle\alpha_2, \beta_2\rangle=1$, 
\\
\begin{tabular}{cc}$\alpha_1$+ & $\alpha_2$\\\tiny{ $A_1$ } & \tiny{ $A_3$ }\end{tabular} &\begin{tabular}{c} = \\~ \end{tabular} & \begin{tabular}{cc}$\beta_1$+ & $\beta_2$\\\tiny{ $A_1$ } & \tiny{ $A_3$ }\end{tabular} & $A_3$ & $A_7$ &  $\langle\alpha_1, \beta_1\rangle=1$, $\langle\alpha_1, \beta_2\rangle=1$, $\langle\alpha_2, \beta_1\rangle=1$, $\langle\alpha_2, \beta_2\rangle=1$, 
\\
\begin{tabular}{cc}$\alpha_1$+ & $\alpha_2$\\\tiny{ $A_3$+$A_1$ } & \tiny{ $A_1$ }\end{tabular} &\begin{tabular}{c} = \\~ \end{tabular} & \begin{tabular}{cc}$\beta_1$+ & $\beta_2$\\\tiny{ $A_3$+$A_1$ } & \tiny{ $A_3$+$A_1$ }\end{tabular} & $A_3$ & $E_7$ &  $\langle\alpha_1, \beta_1\rangle=1$, $\langle\alpha_1, \beta_2\rangle=1$, $\langle\alpha_2, \beta_1\rangle=1$, $\langle\alpha_2, \beta_2\rangle=1$, 
\\
\begin{tabular}{cc}$\alpha_1$+ & $\alpha_2$\\\tiny{ $A_1$ } & \tiny{ $A_3$ }\end{tabular} &\begin{tabular}{c} = \\~ \end{tabular} & \begin{tabular}{cc}$\beta_1$+ & $\beta_2$\\\tiny{ $A_3$+$A_1$ } & \tiny{ $A_3$ }\end{tabular} & $A_3$ & $E_7$ &  $\langle\alpha_1, \beta_1\rangle=1$, $\langle\alpha_1, \beta_2\rangle=1$, $\langle\alpha_2, \beta_1\rangle=1$, $\langle\alpha_2, \beta_2\rangle=1$, 
\\
\begin{tabular}{cc}$\alpha_1$+ & $\alpha_2$\\\tiny{ $A_1$ } & \tiny{ $A_3$ }\end{tabular} &\begin{tabular}{c} = \\~ \end{tabular} & \begin{tabular}{c}$\beta_1$\\\tiny{ $A_3$+$A_1$ }\end{tabular} & $A_2$ & $A_6$ & $\langle\alpha_1, \alpha_2\rangle=-1$,  $\langle\alpha_1, \beta_1\rangle=1$, $\langle\alpha_2, \beta_1\rangle=1$, 
\\
\begin{tabular}{cc}$\alpha_1$+ & $\alpha_2$\\\tiny{ $A_3$ } & \tiny{  }\end{tabular} &\begin{tabular}{c} = \\~ \end{tabular} & \begin{tabular}{cc}$\beta_1$+ & $\beta_2$\\\tiny{ $A_3$+$A_1$ } & \tiny{ $A_3$+$A_1$ }\end{tabular} & $A_3$ & $D_6$ &  $\langle\alpha_1, \beta_1\rangle=1$, $\langle\alpha_1, \beta_2\rangle=1$, $\langle\alpha_2, \beta_1\rangle=1$, $\langle\alpha_2, \beta_2\rangle=1$, 
\\
\begin{tabular}{cc}$\alpha_1$+ & $\alpha_2$\\\tiny{ $A_3$ } & \tiny{  }\end{tabular} &\begin{tabular}{c} = \\~ \end{tabular} & \begin{tabular}{cc}$\beta_1$+ & $\beta_2$\\\tiny{ $A_3$ } & \tiny{ $A_3$ }\end{tabular} & $A_3$ & $D_6$+$A_1$ &  $\langle\alpha_1, \beta_1\rangle=1$, $\langle\alpha_1, \beta_2\rangle=1$, $\langle\alpha_2, \beta_1\rangle=1$, $\langle\alpha_2, \beta_2\rangle=1$, 
\\
\begin{tabular}{c}2$\alpha_1$\\\tiny{  }\end{tabular} &\begin{tabular}{c} = \\~ \end{tabular} & \begin{tabular}{cccc}$\beta_1$+ & $\beta_2$+ & $\beta_3$+ & $\beta_4$\\\tiny{ $A_3$+$A_1$ } & \tiny{ $A_3$+$A_1$ } & \tiny{ $A_3$ } & \tiny{ $A_3$ }\end{tabular} & $D_4$ & $E_7$ &  $\langle\alpha_1, \beta_1\rangle=1$, $\langle\alpha_1, \beta_2\rangle=1$, $\langle\alpha_1, \beta_3\rangle=1$, $\langle\alpha_1, \beta_4\rangle=1$, 
\\
\begin{tabular}{c}2$\alpha_1$\\\tiny{ $A_3$ }\end{tabular} &\begin{tabular}{c} = \\~ \end{tabular} & \begin{tabular}{cccc}$\beta_1$+ & $\beta_2$+ & $\beta_3$+ & $\beta_4$\\\tiny{ $A_3$+$A_1$ } & \tiny{ $A_3$+$A_1$ } & \tiny{ $A_3$ } & \tiny{ $A_3$ }\end{tabular} & $D_4$ & $E_7$ &  $\langle\alpha_1, \beta_1\rangle=1$, $\langle\alpha_1, \beta_2\rangle=1$, $\langle\alpha_1, \beta_3\rangle=1$, $\langle\alpha_1, \beta_4\rangle=1$, 
\\
\begin{tabular}{cc}$\alpha_1$+ & $\alpha_2$\\\tiny{ $A_3$ } & \tiny{ $A_3$ }\end{tabular} &\begin{tabular}{c} = \\~ \end{tabular} & \begin{tabular}{ccc}$\beta_1$+ & $\beta_2$+ & $\beta_3$\\\tiny{ $A_3$+$A_1$ } & \tiny{ $A_3$+$A_1$ } & \tiny{ $A_3$ }\end{tabular} & $D_4$ & $E_7$ & $\langle\alpha_1, \alpha_2\rangle=1$,  $\langle\alpha_1, \beta_1\rangle=1$, $\langle\alpha_1, \beta_2\rangle=1$, $\langle\alpha_1, \beta_3\rangle=1$, $\langle\alpha_2, \beta_1\rangle=1$, $\langle\alpha_2, \beta_2\rangle=1$, $\langle\alpha_2, \beta_3\rangle=1$, 
\\
\begin{tabular}{cc}$\alpha_1$+ & $\alpha_2$\\\tiny{ $A_1$ } & \tiny{ $A_3$ }\end{tabular} &\begin{tabular}{c} = \\~ \end{tabular} & \begin{tabular}{ccc}$\beta_1$+ & $\beta_2$+ & $\beta_3$\\\tiny{ $A_1$ } & \tiny{ $A_3$ } & \tiny{ $A_3$ }\end{tabular} & $D_4$ & $E_7$ & $\langle\alpha_1, \alpha_2\rangle=1$,  $\langle\alpha_1, \beta_1\rangle=1$, $\langle\alpha_1, \beta_2\rangle=1$, $\langle\alpha_1, \beta_3\rangle=1$, $\langle\alpha_2, \beta_1\rangle=1$, $\langle\alpha_2, \beta_2\rangle=1$, $\langle\alpha_2, \beta_3\rangle=1$, 
\\
\begin{tabular}{cc}$\alpha_1$+ & $\alpha_2$\\\tiny{ $A_3$+$A_1$ } & \tiny{  }\end{tabular} &\begin{tabular}{c} = \\~ \end{tabular} & \begin{tabular}{ccc}$\beta_1$+ & $\beta_2$+ & $\beta_3$\\\tiny{ $A_1$ } & \tiny{ $A_3$ } & \tiny{ $A_3$ }\end{tabular} & $D_4$ & $E_7$ & $\langle\alpha_1, \alpha_2\rangle=1$,  $\langle\alpha_1, \beta_1\rangle=1$, $\langle\alpha_1, \beta_2\rangle=1$, $\langle\alpha_1, \beta_3\rangle=1$, $\langle\alpha_2, \beta_1\rangle=1$, $\langle\alpha_2, \beta_2\rangle=1$, $\langle\alpha_2, \beta_3\rangle=1$, 
\\
\begin{tabular}{cc}$\alpha_1$+ & $\alpha_2$\\\tiny{ $A_1$ } & \tiny{ $A_1$ }\end{tabular} &\begin{tabular}{c} = \\~ \end{tabular} & \begin{tabular}{ccc}$\beta_1$+ & $\beta_2$+ & $\beta_3$\\\tiny{ $A_3$+$A_1$ } & \tiny{ $A_3$+$A_1$ } & \tiny{ $A_3$ }\end{tabular} & $D_4$ & $E_7$ & $\langle\alpha_1, \alpha_2\rangle=1$,  $\langle\alpha_1, \beta_1\rangle=1$, $\langle\alpha_1, \beta_2\rangle=1$, $\langle\alpha_1, \beta_3\rangle=1$, $\langle\alpha_2, \beta_1\rangle=1$, $\langle\alpha_2, \beta_2\rangle=1$, $\langle\alpha_2, \beta_3\rangle=1$, 
\\
\begin{tabular}{cc}$\alpha_1$+ & $\alpha_2$\\\tiny{ $A_3$ } & \tiny{ $A_3$ }\end{tabular} &\begin{tabular}{c} = \\~ \end{tabular} & \begin{tabular}{cc}$\beta_1$+ & $\beta_2$\\\tiny{ $A_3$+$A_1$ } & \tiny{ $A_3$+$A_1$ }\end{tabular} & $A_3$ & $E_7$ &  $\langle\alpha_1, \beta_1\rangle=1$, $\langle\alpha_1, \beta_2\rangle=1$, $\langle\alpha_2, \beta_1\rangle=1$, $\langle\alpha_2, \beta_2\rangle=1$, 
\\
\begin{tabular}{cc}$\alpha_1$+ & $\alpha_2$\\\tiny{ $A_3$+$A_1$ } & \tiny{ $A_3$ }\end{tabular} &\begin{tabular}{c} = \\~ \end{tabular} & \begin{tabular}{ccc}$\beta_1$+ & $\beta_2$+ & $\beta_3$\\\tiny{ $A_1$ } & \tiny{ $A_3$ } & \tiny{ $A_3$ }\end{tabular} & $D_4$ & $E_7$ & $\langle\alpha_1, \alpha_2\rangle=1$,  $\langle\alpha_1, \beta_1\rangle=1$, $\langle\alpha_1, \beta_2\rangle=1$, $\langle\alpha_1, \beta_3\rangle=1$, $\langle\alpha_2, \beta_1\rangle=1$, $\langle\alpha_2, \beta_2\rangle=1$, $\langle\alpha_2, \beta_3\rangle=1$, 
\\
\begin{tabular}{cc}$\alpha_1$+ & $\alpha_2$\\\tiny{ $A_3$ } & \tiny{ $A_3$ }\end{tabular} &\begin{tabular}{c} = \\~ \end{tabular} & \begin{tabular}{cc}$\beta_1$+ & $\beta_2$\\\tiny{ $A_3$+$A_1$ } & \tiny{ $A_3$+$A_1$ }\end{tabular} & $A_3$ & $E_6$ &  $\langle\alpha_1, \beta_1\rangle=1$, $\langle\alpha_1, \beta_2\rangle=1$, $\langle\alpha_2, \beta_1\rangle=1$, $\langle\alpha_2, \beta_2\rangle=1$, 
\\
\begin{tabular}{cc}$\alpha_1$+ & $\alpha_2$\\\tiny{ $A_1$ } & \tiny{ $A_3$ }\end{tabular} &\begin{tabular}{c} = \\~ \end{tabular} & \begin{tabular}{c}$\beta_1$\\\tiny{ $A_3$+$A_1$ }\end{tabular} & $A_2$ & $D_6$ & $\langle\alpha_1, \alpha_2\rangle=-1$,  $\langle\alpha_1, \beta_1\rangle=1$, $\langle\alpha_2, \beta_1\rangle=1$, 
\\
\begin{tabular}{cc}$\alpha_1$+ & $\alpha_2$\\\tiny{ $A_1$ } & \tiny{ $A_1$ }\end{tabular} &\begin{tabular}{c} = \\~ \end{tabular} & \begin{tabular}{cc}$\beta_1$+ & $\beta_2$\\\tiny{ $A_3$+$A_1$ } & \tiny{ $A_3$+$A_1$ }\end{tabular} & $A_3$ & $D_6$ &  $\langle\alpha_1, \beta_1\rangle=1$, $\langle\alpha_1, \beta_2\rangle=1$, $\langle\alpha_2, \beta_1\rangle=1$, $\langle\alpha_2, \beta_2\rangle=1$, 
\\
\begin{tabular}{cc}$\alpha_1$+ & $\alpha_2$\\\tiny{ $A_1$ } & \tiny{ $A_3$ }\end{tabular} &\begin{tabular}{c} = \\~ \end{tabular} & \begin{tabular}{cc}$\beta_1$+ & $\beta_2$\\\tiny{ $A_3$+$A_1$ } & \tiny{  }\end{tabular} & $A_3$ & $A_7$ &  $\langle\alpha_1, \beta_1\rangle=1$, $\langle\alpha_1, \beta_2\rangle=1$, $\langle\alpha_2, \beta_1\rangle=1$, $\langle\alpha_2, \beta_2\rangle=1$, 
\\
\begin{tabular}{c}2$\alpha_1$\\\tiny{ $A_1$ }\end{tabular} &\begin{tabular}{c} = \\~ \end{tabular} & \begin{tabular}{cccc}$\beta_1$+ & $\beta_2$+ & $\beta_3$+ & $\beta_4$\\\tiny{ $A_3$+$A_1$ } & \tiny{ $A_3$+$A_1$ } & \tiny{ $A_3$ } & \tiny{  }\end{tabular} & $D_4$ & $E_7$ &  $\langle\alpha_1, \beta_1\rangle=1$, $\langle\alpha_1, \beta_2\rangle=1$, $\langle\alpha_1, \beta_3\rangle=1$, $\langle\alpha_1, \beta_4\rangle=1$, 
\\
\begin{tabular}{cc}$\alpha_1$+ & $\alpha_2$\\\tiny{ $A_1$ } & \tiny{ $A_3$ }\end{tabular} &\begin{tabular}{c} = \\~ \end{tabular} & \begin{tabular}{ccc}$\beta_1$+ & $\beta_2$+ & $\beta_3$\\\tiny{ $A_3$+$A_1$ } & \tiny{ $A_3$ } & \tiny{  }\end{tabular} & $D_4$ & $E_7$ & $\langle\alpha_1, \alpha_2\rangle=1$,  $\langle\alpha_1, \beta_1\rangle=1$, $\langle\alpha_1, \beta_2\rangle=1$, $\langle\alpha_1, \beta_3\rangle=1$, $\langle\alpha_2, \beta_1\rangle=1$, $\langle\alpha_2, \beta_2\rangle=1$, $\langle\alpha_2, \beta_3\rangle=1$, 
\\
\begin{tabular}{cc}$\alpha_1$+ & $\alpha_2$\\\tiny{ $A_3$ } & \tiny{  }\end{tabular} &\begin{tabular}{c} = \\~ \end{tabular} & \begin{tabular}{ccc}$\beta_1$+ & $\beta_2$+ & $\beta_3$\\\tiny{ $A_3$+$A_1$ } & \tiny{ $A_3$+$A_1$ } & \tiny{ $A_3$ }\end{tabular} & $D_4$ & $E_7$ & $\langle\alpha_1, \alpha_2\rangle=1$,  $\langle\alpha_1, \beta_1\rangle=1$, $\langle\alpha_1, \beta_2\rangle=1$, $\langle\alpha_1, \beta_3\rangle=1$, $\langle\alpha_2, \beta_1\rangle=1$, $\langle\alpha_2, \beta_2\rangle=1$, $\langle\alpha_2, \beta_3\rangle=1$, 
\\
\begin{tabular}{c}2$\alpha_1$\\\tiny{ $A_3$ }\end{tabular} &\begin{tabular}{c} = \\~ \end{tabular} & \begin{tabular}{cccc}$\beta_1$+ & $\beta_2$+ & $\beta_3$+ & $\beta_4$\\\tiny{ $A_3$+$A_1$ } & \tiny{ $A_3$+$A_1$ } & \tiny{ $A_3$ } & \tiny{  }\end{tabular} & $D_4$ & $E_7$ &  $\langle\alpha_1, \beta_1\rangle=1$, $\langle\alpha_1, \beta_2\rangle=1$, $\langle\alpha_1, \beta_3\rangle=1$, $\langle\alpha_1, \beta_4\rangle=1$, 
\\
\begin{tabular}{cc}$\alpha_1$+ & $\alpha_2$\\\tiny{ $A_3$ } & \tiny{ $A_3$ }\end{tabular} &\begin{tabular}{c} = \\~ \end{tabular} & \begin{tabular}{ccc}$\beta_1$+ & $\beta_2$+ & $\beta_3$\\\tiny{ $A_3$+$A_1$ } & \tiny{ $A_3$+$A_1$ } & \tiny{  }\end{tabular} & $D_4$ & $E_7$ & $\langle\alpha_1, \alpha_2\rangle=1$,  $\langle\alpha_1, \beta_1\rangle=1$, $\langle\alpha_1, \beta_2\rangle=1$, $\langle\alpha_1, \beta_3\rangle=1$, $\langle\alpha_2, \beta_1\rangle=1$, $\langle\alpha_2, \beta_2\rangle=1$, $\langle\alpha_2, \beta_3\rangle=1$, 
\\
\hline\multicolumn{5}{c}{$\mathfrak{k}$-semisimple type: $A_3$+$A_1$}\\
\hline\hline\begin{tabular}{cc}$\alpha_1$+ & $\alpha_2$\\\tiny{  } & \tiny{  }\end{tabular} &\begin{tabular}{c} = \\~ \end{tabular} & \begin{tabular}{ccc}$\beta_1$+ & $\beta_2$+ & $\beta_3$\\\tiny{ $A_3$+$A_1$ } & \tiny{ $A_3$+$A_1$ } & \tiny{ $A_3$ }\end{tabular} & $D_4$ & $E_7$ & $\langle\alpha_1, \alpha_2\rangle=1$,  $\langle\alpha_1, \beta_1\rangle=1$, $\langle\alpha_1, \beta_2\rangle=1$, $\langle\alpha_1, \beta_3\rangle=1$, $\langle\alpha_2, \beta_1\rangle=1$, $\langle\alpha_2, \beta_2\rangle=1$, $\langle\alpha_2, \beta_3\rangle=1$, 
\\
\begin{tabular}{cc}$\alpha_1$+ & $\alpha_2$\\\tiny{  } & \tiny{  }\end{tabular} &\begin{tabular}{c} = \\~ \end{tabular} & \begin{tabular}{cc}$\beta_1$+ & $\beta_2$\\\tiny{ $A_3$ } & \tiny{ $A_3$ }\end{tabular} & $A_3$ & $D_5$+$A_1$ &  $\langle\alpha_1, \beta_1\rangle=1$, $\langle\alpha_1, \beta_2\rangle=1$, $\langle\alpha_2, \beta_1\rangle=1$, $\langle\alpha_2, \beta_2\rangle=1$, 
\\
\begin{tabular}{cc}$\alpha_1$+ & $\alpha_2$\\\tiny{ $A_3$ } & \tiny{  }\end{tabular} &\begin{tabular}{c} = \\~ \end{tabular} & \begin{tabular}{c}$\beta_1$\\\tiny{ $A_3$ }\end{tabular} & $A_2$ & $D_5$+$A_1$ & $\langle\alpha_1, \alpha_2\rangle=-1$,  $\langle\alpha_1, \beta_1\rangle=1$, $\langle\alpha_2, \beta_1\rangle=1$, 
\\
\hline\multicolumn{5}{c}{$\mathfrak{k}$-semisimple type: $A_2$+$A_2$}\\
\hline\hline\begin{tabular}{cc}$\alpha_1$+ & $\alpha_2$\\\tiny{  } & \tiny{  }\end{tabular} &\begin{tabular}{c} = \\~ \end{tabular} & \begin{tabular}{ccc}$\beta_1$+ & $\beta_2$+ & $\beta_3$\\\tiny{ $A_2$+$A_2$' } & \tiny{ $A_2$+$A_2$' } & \tiny{ $A_2$+$A_2$' }\end{tabular} & $D_4$ & $E_6$ & $\langle\alpha_1, \alpha_2\rangle=1$,  $\langle\alpha_1, \beta_1\rangle=1$, $\langle\alpha_1, \beta_2\rangle=1$, $\langle\alpha_1, \beta_3\rangle=1$, $\langle\alpha_2, \beta_1\rangle=1$, $\langle\alpha_2, \beta_2\rangle=1$, $\langle\alpha_2, \beta_3\rangle=1$, 
\\
\begin{tabular}{cc}$\alpha_1$+ & $\alpha_2$\\\tiny{ $A_2$ } & \tiny{  }\end{tabular} &\begin{tabular}{c} = \\~ \end{tabular} & \begin{tabular}{c}$\beta_1$\\\tiny{ $A_2$ }\end{tabular} & $A_2$ & $A_4$+$A_2$ & $\langle\alpha_1, \alpha_2\rangle=-1$,  $\langle\alpha_1, \beta_1\rangle=1$, $\langle\alpha_2, \beta_1\rangle=1$, 
\\
\begin{tabular}{cc}$\alpha_1$+ & $\alpha_2$\\\tiny{ $A_2$+$A_2$' } & \tiny{  }\end{tabular} &\begin{tabular}{c} = \\~ \end{tabular} & \begin{tabular}{cc}$\beta_1$+ & $\beta_2$\\\tiny{ $A_2$+$A_2$' } & \tiny{ $A_2$ }\end{tabular} & $A_3$ & $E_7$ &  $\langle\alpha_1, \beta_1\rangle=1$, $\langle\alpha_1, \beta_2\rangle=1$, $\langle\alpha_2, \beta_1\rangle=1$, $\langle\alpha_2, \beta_2\rangle=1$, 
\\
\begin{tabular}{c}2$\alpha_1$\\\tiny{ $A_2$ }\end{tabular} &\begin{tabular}{c} = \\~ \end{tabular} & \begin{tabular}{cccc}$\beta_1$+ & $\beta_2$+ & $\beta_3$+ & $\beta_4$\\\tiny{ $A_2$+$A_2$' } & \tiny{ $A_2$+$A_2$' } & \tiny{ $A_2$' } & \tiny{ $A_2$ }\end{tabular} & $D_4$ & $E_7$ &  $\langle\alpha_1, \beta_1\rangle=1$, $\langle\alpha_1, \beta_2\rangle=1$, $\langle\alpha_1, \beta_3\rangle=1$, $\langle\alpha_1, \beta_4\rangle=1$, 
\\
\begin{tabular}{c}2$\alpha_1$\\\tiny{  }\end{tabular} &\begin{tabular}{c} = \\~ \end{tabular} & \begin{tabular}{cccc}$\beta_1$+ & $\beta_2$+ & $\beta_3$+ & $\beta_4$\\\tiny{ $A_2$+$A_2$' } & \tiny{ $A_2$+$A_2$' } & \tiny{ $A_2$' } & \tiny{ $A_2$ }\end{tabular} & $D_4$ & $E_7$ &  $\langle\alpha_1, \beta_1\rangle=1$, $\langle\alpha_1, \beta_2\rangle=1$, $\langle\alpha_1, \beta_3\rangle=1$, $\langle\alpha_1, \beta_4\rangle=1$, 
\\
\begin{tabular}{cc}$\alpha_1$+ & $\alpha_2$\\\tiny{ $A_2$ } & \tiny{ $A_2$' }\end{tabular} &\begin{tabular}{c} = \\~ \end{tabular} & \begin{tabular}{c}$\beta_1$\\\tiny{ $A_2$+$A_2$' }\end{tabular} & $A_2$ & $A_6$ & $\langle\alpha_1, \alpha_2\rangle=-1$,  $\langle\alpha_1, \beta_1\rangle=1$, $\langle\alpha_2, \beta_1\rangle=1$, 
\\
\begin{tabular}{cc}$\alpha_1$+ & $\alpha_2$\\\tiny{ $A_2$ } & \tiny{ $A_2$' }\end{tabular} &\begin{tabular}{c} = \\~ \end{tabular} & \begin{tabular}{cc}$\beta_1$+ & $\beta_2$\\\tiny{ $A_2$ } & \tiny{ $A_2$' }\end{tabular} & $A_3$ & $A_7$ &  $\langle\alpha_1, \beta_1\rangle=1$, $\langle\alpha_1, \beta_2\rangle=1$, $\langle\alpha_2, \beta_1\rangle=1$, $\langle\alpha_2, \beta_2\rangle=1$, 
\\
\begin{tabular}{cc}$\alpha_1$+ & $\alpha_2$\\\tiny{ $A_2$ } & \tiny{ $A_2$ }\end{tabular} &\begin{tabular}{c} = \\~ \end{tabular} & \begin{tabular}{ccc}$\beta_1$+ & $\beta_2$+ & $\beta_3$\\\tiny{ $A_2$+$A_2$' } & \tiny{ $A_2$' } & \tiny{ $A_2$ }\end{tabular} & $D_4$ & $E_7$ & $\langle\alpha_1, \alpha_2\rangle=1$,  $\langle\alpha_1, \beta_1\rangle=1$, $\langle\alpha_1, \beta_2\rangle=1$, $\langle\alpha_1, \beta_3\rangle=1$, $\langle\alpha_2, \beta_1\rangle=1$, $\langle\alpha_2, \beta_2\rangle=1$, $\langle\alpha_2, \beta_3\rangle=1$, 
\\
\begin{tabular}{cc}$\alpha_1$+ & $\alpha_2$\\\tiny{ $A_2$ } & \tiny{ $A_2$' }\end{tabular} &\begin{tabular}{c} = \\~ \end{tabular} & \begin{tabular}{cc}$\beta_1$+ & $\beta_2$\\\tiny{ $A_2$+$A_2$' } & \tiny{ $A_2$+$A_2$' }\end{tabular} & $A_3$ & $D_6$ &  $\langle\alpha_1, \beta_1\rangle=1$, $\langle\alpha_1, \beta_2\rangle=1$, $\langle\alpha_2, \beta_1\rangle=1$, $\langle\alpha_2, \beta_2\rangle=1$, 
\\
\begin{tabular}{cc}$\alpha_1$+ & $\alpha_2$\\\tiny{  } & \tiny{  }\end{tabular} &\begin{tabular}{c} = \\~ \end{tabular} & \begin{tabular}{ccc}$\beta_1$+ & $\beta_2$+ & $\beta_3$\\\tiny{ $A_2$+$A_2$' } & \tiny{ $A_2$' } & \tiny{ $A_2$ }\end{tabular} & $D_4$ & $E_7$ & $\langle\alpha_1, \alpha_2\rangle=1$,  $\langle\alpha_1, \beta_1\rangle=1$, $\langle\alpha_1, \beta_2\rangle=1$, $\langle\alpha_1, \beta_3\rangle=1$, $\langle\alpha_2, \beta_1\rangle=1$, $\langle\alpha_2, \beta_2\rangle=1$, $\langle\alpha_2, \beta_3\rangle=1$, 
\\
\begin{tabular}{cc}$\alpha_1$+ & $\alpha_2$\\\tiny{ $A_2$ } & \tiny{ $A_2$ }\end{tabular} &\begin{tabular}{c} = \\~ \end{tabular} & \begin{tabular}{ccc}$\beta_1$+ & $\beta_2$+ & $\beta_3$\\\tiny{ $A_2$+$A_2$' } & \tiny{ $A_2$+$A_2$' } & \tiny{ $A_2$+$A_2$' }\end{tabular} & $D_4$ & $E_7$ & $\langle\alpha_1, \alpha_2\rangle=1$,  $\langle\alpha_1, \beta_1\rangle=1$, $\langle\alpha_1, \beta_2\rangle=1$, $\langle\alpha_1, \beta_3\rangle=1$, $\langle\alpha_2, \beta_1\rangle=1$, $\langle\alpha_2, \beta_2\rangle=1$, $\langle\alpha_2, \beta_3\rangle=1$, 
\\
\begin{tabular}{cc}$\alpha_1$+ & $\alpha_2$\\\tiny{ $A_2$ } & \tiny{  }\end{tabular} &\begin{tabular}{c} = \\~ \end{tabular} & \begin{tabular}{cc}$\beta_1$+ & $\beta_2$\\\tiny{ $A_2$+$A_2$' } & \tiny{ $A_2$ }\end{tabular} & $A_3$ & $E_7$ &  $\langle\alpha_1, \beta_1\rangle=1$, $\langle\alpha_1, \beta_2\rangle=1$, $\langle\alpha_2, \beta_1\rangle=1$, $\langle\alpha_2, \beta_2\rangle=1$, 
\\
\begin{tabular}{cc}$\alpha_1$+ & $\alpha_2$\\\tiny{ $A_2$ } & \tiny{  }\end{tabular} &\begin{tabular}{c} = \\~ \end{tabular} & \begin{tabular}{ccc}$\beta_1$+ & $\beta_2$+ & $\beta_3$\\\tiny{ $A_2$+$A_2$' } & \tiny{ $A_2$+$A_2$' } & \tiny{ $A_2$' }\end{tabular} & $D_4$ & $E_7$ & $\langle\alpha_1, \alpha_2\rangle=1$,  $\langle\alpha_1, \beta_1\rangle=1$, $\langle\alpha_1, \beta_2\rangle=1$, $\langle\alpha_1, \beta_3\rangle=1$, $\langle\alpha_2, \beta_1\rangle=1$, $\langle\alpha_2, \beta_2\rangle=1$, $\langle\alpha_2, \beta_3\rangle=1$, 
\\
\begin{tabular}{cc}$\alpha_1$+ & $\alpha_2$\\\tiny{ $A_2$ } & \tiny{  }\end{tabular} &\begin{tabular}{c} = \\~ \end{tabular} & \begin{tabular}{cc}$\beta_1$+ & $\beta_2$\\\tiny{ $A_2$+$A_2$' } & \tiny{ $A_2$' }\end{tabular} & $A_3$ & $E_7$ &  $\langle\alpha_1, \beta_1\rangle=1$, $\langle\alpha_1, \beta_2\rangle=1$, $\langle\alpha_2, \beta_1\rangle=1$, $\langle\alpha_2, \beta_2\rangle=1$, 
\\
\begin{tabular}{cc}$\alpha_1$+ & $\alpha_2$\\\tiny{ $A_2$+$A_2$' } & \tiny{  }\end{tabular} &\begin{tabular}{c} = \\~ \end{tabular} & \begin{tabular}{cc}$\beta_1$+ & $\beta_2$\\\tiny{ $A_2$ } & \tiny{ $A_2$' }\end{tabular} & $A_3$ & $A_7$ &  $\langle\alpha_1, \beta_1\rangle=1$, $\langle\alpha_1, \beta_2\rangle=1$, $\langle\alpha_2, \beta_1\rangle=1$, $\langle\alpha_2, \beta_2\rangle=1$, 
\\
\begin{tabular}{cc}$\alpha_1$+ & $\alpha_2$\\\tiny{ $A_2$+$A_2$' } & \tiny{  }\end{tabular} &\begin{tabular}{c} = \\~ \end{tabular} & \begin{tabular}{c}$\beta_1$\\\tiny{ $A_2$+$A_2$' }\end{tabular} & $A_2$ & $E_6$ & $\langle\alpha_1, \alpha_2\rangle=-1$,  $\langle\alpha_1, \beta_1\rangle=1$, $\langle\alpha_2, \beta_1\rangle=1$, 
\\
\begin{tabular}{cc}$\alpha_1$+ & $\alpha_2$\\\tiny{ $A_2$+$A_2$' } & \tiny{  }\end{tabular} &\begin{tabular}{c} = \\~ \end{tabular} & \begin{tabular}{cc}$\beta_1$+ & $\beta_2$\\\tiny{ $A_2$+$A_2$' } & \tiny{ $A_2$' }\end{tabular} & $A_3$ & $E_7$ &  $\langle\alpha_1, \beta_1\rangle=1$, $\langle\alpha_1, \beta_2\rangle=1$, $\langle\alpha_2, \beta_1\rangle=1$, $\langle\alpha_2, \beta_2\rangle=1$, 
\\
\begin{tabular}{cc}$\alpha_1$+ & $\alpha_2$\\\tiny{ $A_2$+$A_2$' } & \tiny{ $A_2$' }\end{tabular} &\begin{tabular}{c} = \\~ \end{tabular} & \begin{tabular}{cc}$\beta_1$+ & $\beta_2$\\\tiny{ $A_2$+$A_2$' } & \tiny{ $A_2$' }\end{tabular} & $A_3$ & $E_7$ &  $\langle\alpha_1, \beta_1\rangle=1$, $\langle\alpha_1, \beta_2\rangle=1$, $\langle\alpha_2, \beta_1\rangle=1$, $\langle\alpha_2, \beta_2\rangle=1$, 
\\
\begin{tabular}{cc}$\alpha_1$+ & $\alpha_2$\\\tiny{ $A_2$+$A_2$' } & \tiny{ $A_2$ }\end{tabular} &\begin{tabular}{c} = \\~ \end{tabular} & \begin{tabular}{cc}$\beta_1$+ & $\beta_2$\\\tiny{ $A_2$+$A_2$' } & \tiny{ $A_2$+$A_2$' }\end{tabular} & $A_3$ & $E_7$ &  $\langle\alpha_1, \beta_1\rangle=1$, $\langle\alpha_1, \beta_2\rangle=1$, $\langle\alpha_2, \beta_1\rangle=1$, $\langle\alpha_2, \beta_2\rangle=1$, 
\\
\begin{tabular}{cc}$\alpha_1$+ & $\alpha_2$\\\tiny{ $A_2$+$A_2$' } & \tiny{  }\end{tabular} &\begin{tabular}{c} = \\~ \end{tabular} & \begin{tabular}{cc}$\beta_1$+ & $\beta_2$\\\tiny{ $A_2$+$A_2$' } & \tiny{ $A_2$+$A_2$' }\end{tabular} & $A_3$ & $E_6$ &  $\langle\alpha_1, \beta_1\rangle=1$, $\langle\alpha_1, \beta_2\rangle=1$, $\langle\alpha_2, \beta_1\rangle=1$, $\langle\alpha_2, \beta_2\rangle=1$, 
\\
\begin{tabular}{cc}$\alpha_1$+ & $\alpha_2$\\\tiny{ $A_2$+$A_2$' } & \tiny{ $A_2$' }\end{tabular} &\begin{tabular}{c} = \\~ \end{tabular} & \begin{tabular}{cc}$\beta_1$+ & $\beta_2$\\\tiny{ $A_2$+$A_2$' } & \tiny{ $A_2$+$A_2$' }\end{tabular} & $A_3$ & $E_7$ &  $\langle\alpha_1, \beta_1\rangle=1$, $\langle\alpha_1, \beta_2\rangle=1$, $\langle\alpha_2, \beta_1\rangle=1$, $\langle\alpha_2, \beta_2\rangle=1$, 
\\
\begin{tabular}{cc}$\alpha_1$+ & $\alpha_2$\\\tiny{ $A_2$ } & \tiny{ $A_2$ }\end{tabular} &\begin{tabular}{c} = \\~ \end{tabular} & \begin{tabular}{cc}$\beta_1$+ & $\beta_2$\\\tiny{ $A_2$+$A_2$' } & \tiny{ $A_2$+$A_2$' }\end{tabular} & $A_3$ & $E_7$ &  $\langle\alpha_1, \beta_1\rangle=1$, $\langle\alpha_1, \beta_2\rangle=1$, $\langle\alpha_2, \beta_1\rangle=1$, $\langle\alpha_2, \beta_2\rangle=1$, 
\\
\begin{tabular}{cc}$\alpha_1$+ & $\alpha_2$\\\tiny{ $A_2$+$A_2$' } & \tiny{ $A_2$ }\end{tabular} &\begin{tabular}{c} = \\~ \end{tabular} & \begin{tabular}{cc}$\beta_1$+ & $\beta_2$\\\tiny{ $A_2$+$A_2$' } & \tiny{ $A_2$ }\end{tabular} & $A_3$ & $E_7$ &  $\langle\alpha_1, \beta_1\rangle=1$, $\langle\alpha_1, \beta_2\rangle=1$, $\langle\alpha_2, \beta_1\rangle=1$, $\langle\alpha_2, \beta_2\rangle=1$, 
\\
\begin{tabular}{cc}$\alpha_1$+ & $\alpha_2$\\\tiny{ $A_2$ } & \tiny{ $A_2$' }\end{tabular} &\begin{tabular}{c} = \\~ \end{tabular} & \begin{tabular}{cc}$\beta_1$+ & $\beta_2$\\\tiny{ $A_2$+$A_2$' } & \tiny{  }\end{tabular} & $A_3$ & $A_7$ &  $\langle\alpha_1, \beta_1\rangle=1$, $\langle\alpha_1, \beta_2\rangle=1$, $\langle\alpha_2, \beta_1\rangle=1$, $\langle\alpha_2, \beta_2\rangle=1$, 
\\
\begin{tabular}{c}2$\alpha_1$\\\tiny{ $A_2$+$A_2$' }\end{tabular} &\begin{tabular}{c} = \\~ \end{tabular} & \begin{tabular}{cccc}$\beta_1$+ & $\beta_2$+ & $\beta_3$+ & $\beta_4$\\\tiny{ $A_2$+$A_2$' } & \tiny{ $A_2$+$A_2$' } & \tiny{ $A_2$' } & \tiny{ $A_2$ }\end{tabular} & $D_4$ & $E_7$ &  $\langle\alpha_1, \beta_1\rangle=1$, $\langle\alpha_1, \beta_2\rangle=1$, $\langle\alpha_1, \beta_3\rangle=1$, $\langle\alpha_1, \beta_4\rangle=1$, 
\\
\begin{tabular}{cc}$\alpha_1$+ & $\alpha_2$\\\tiny{ $A_2$+$A_2$' } & \tiny{ $A_2$ }\end{tabular} &\begin{tabular}{c} = \\~ \end{tabular} & \begin{tabular}{ccc}$\beta_1$+ & $\beta_2$+ & $\beta_3$\\\tiny{ $A_2$+$A_2$' } & \tiny{ $A_2$+$A_2$' } & \tiny{ $A_2$ }\end{tabular} & $D_4$ & $E_7$ & $\langle\alpha_1, \alpha_2\rangle=1$,  $\langle\alpha_1, \beta_1\rangle=1$, $\langle\alpha_1, \beta_2\rangle=1$, $\langle\alpha_1, \beta_3\rangle=1$, $\langle\alpha_2, \beta_1\rangle=1$, $\langle\alpha_2, \beta_2\rangle=1$, $\langle\alpha_2, \beta_3\rangle=1$, 
\\
\begin{tabular}{c}2$\alpha_1$\\\tiny{ $A_2$ }\end{tabular} &\begin{tabular}{c} = \\~ \end{tabular} & \begin{tabular}{cccc}$\beta_1$+ & $\beta_2$+ & $\beta_3$+ & $\beta_4$\\\tiny{ $A_2$+$A_2$' } & \tiny{ $A_2$+$A_2$' } & \tiny{ $A_2$+$A_2$' } & \tiny{  }\end{tabular} & $D_4$ & $E_7$ &  $\langle\alpha_1, \beta_1\rangle=1$, $\langle\alpha_1, \beta_2\rangle=1$, $\langle\alpha_1, \beta_3\rangle=1$, $\langle\alpha_1, \beta_4\rangle=1$, 
\\
\begin{tabular}{cc}$\alpha_1$+ & $\alpha_2$\\\tiny{ $A_2$ } & \tiny{ $A_2$' }\end{tabular} &\begin{tabular}{c} = \\~ \end{tabular} & \begin{tabular}{ccc}$\beta_1$+ & $\beta_2$+ & $\beta_3$\\\tiny{ $A_2$+$A_2$' } & \tiny{ $A_2$+$A_2$' } & \tiny{  }\end{tabular} & $D_4$ & $E_7$ & $\langle\alpha_1, \alpha_2\rangle=1$,  $\langle\alpha_1, \beta_1\rangle=1$, $\langle\alpha_1, \beta_2\rangle=1$, $\langle\alpha_1, \beta_3\rangle=1$, $\langle\alpha_2, \beta_1\rangle=1$, $\langle\alpha_2, \beta_2\rangle=1$, $\langle\alpha_2, \beta_3\rangle=1$, 
\\
\hline\multicolumn{5}{c}{$\mathfrak{k}$-semisimple type: $A_2$+$A_1$+$A_1$}\\
\hline\hline\begin{tabular}{c}2$\alpha_1$\\\tiny{  }\end{tabular} &\begin{tabular}{c} = \\~ \end{tabular} & \begin{tabular}{cccc}$\beta_1$+ & $\beta_2$+ & $\beta_3$+ & $\beta_4$\\\tiny{ $A_2$+$A_1$ } & \tiny{ $A_2$+$A_1$' } & \tiny{ $A_1$ } & \tiny{ $A_1$' }\end{tabular} & $D_4$ & $E_7$ &  $\langle\alpha_1, \beta_1\rangle=1$, $\langle\alpha_1, \beta_2\rangle=1$, $\langle\alpha_1, \beta_3\rangle=1$, $\langle\alpha_1, \beta_4\rangle=1$, 
\\
\begin{tabular}{cc}$\alpha_1$+ & $\alpha_2$\\\tiny{ $A_2$ } & \tiny{  }\end{tabular} &\begin{tabular}{c} = \\~ \end{tabular} & \begin{tabular}{cc}$\beta_1$+ & $\beta_2$\\\tiny{ $A_2$+$A_1$ } & \tiny{ $A_2$+$A_1$ }\end{tabular} & $A_3$ & $D_5$+$A_1$ &  $\langle\alpha_1, \beta_1\rangle=1$, $\langle\alpha_1, \beta_2\rangle=1$, $\langle\alpha_2, \beta_1\rangle=1$, $\langle\alpha_2, \beta_2\rangle=1$, 
\\
\begin{tabular}{cc}$\alpha_1$+ & $\alpha_2$\\\tiny{ $A_1$ } & \tiny{  }\end{tabular} &\begin{tabular}{c} = \\~ \end{tabular} & \begin{tabular}{c}$\beta_1$\\\tiny{ $A_1$ }\end{tabular} & $A_2$ & $A_3$+$A_2$+$A_1$ & $\langle\alpha_1, \alpha_2\rangle=-1$,  $\langle\alpha_1, \beta_1\rangle=1$, $\langle\alpha_2, \beta_1\rangle=1$, 
\\
\begin{tabular}{cc}$\alpha_1$+ & $\alpha_2$\\\tiny{ $A_2$+$A_1$ } & \tiny{  }\end{tabular} &\begin{tabular}{c} = \\~ \end{tabular} & \begin{tabular}{c}$\beta_1$\\\tiny{ $A_2$+$A_1$ }\end{tabular} & $A_2$ & $D_5$+$A_1$ & $\langle\alpha_1, \alpha_2\rangle=-1$,  $\langle\alpha_1, \beta_1\rangle=1$, $\langle\alpha_2, \beta_1\rangle=1$, 
\\
\begin{tabular}{cc}$\alpha_1$+ & $\alpha_2$\\\tiny{ $A_2$+$A_1$+$A_1$' } & \tiny{  }\end{tabular} &\begin{tabular}{c} = \\~ \end{tabular} & \begin{tabular}{cc}$\beta_1$+ & $\beta_2$\\\tiny{ $A_2$+$A_1$ } & \tiny{ $A_1$' }\end{tabular} & $A_3$ & $E_7$ &  $\langle\alpha_1, \beta_1\rangle=1$, $\langle\alpha_1, \beta_2\rangle=1$, $\langle\alpha_2, \beta_1\rangle=1$, $\langle\alpha_2, \beta_2\rangle=1$, 
\\
\begin{tabular}{cc}$\alpha_1$+ & $\alpha_2$\\\tiny{ $A_2$ } & \tiny{  }\end{tabular} &\begin{tabular}{c} = \\~ \end{tabular} & \begin{tabular}{cc}$\beta_1$+ & $\beta_2$\\\tiny{ $A_2$+$A_1$ } & \tiny{ $A_1$ }\end{tabular} & $A_3$ & $D_6$+$A_1$ &  $\langle\alpha_1, \beta_1\rangle=1$, $\langle\alpha_1, \beta_2\rangle=1$, $\langle\alpha_2, \beta_1\rangle=1$, $\langle\alpha_2, \beta_2\rangle=1$, 
\\
\begin{tabular}{cc}$\alpha_1$+ & $\alpha_2$\\\tiny{ $A_1$+$A_1$' } & \tiny{  }\end{tabular} &\begin{tabular}{c} = \\~ \end{tabular} & \begin{tabular}{cc}$\beta_1$+ & $\beta_2$\\\tiny{ $A_2$+$A_1$ } & \tiny{ $A_2$+$A_1$' }\end{tabular} & $A_3$ & $E_7$ &  $\langle\alpha_1, \beta_1\rangle=1$, $\langle\alpha_1, \beta_2\rangle=1$, $\langle\alpha_2, \beta_1\rangle=1$, $\langle\alpha_2, \beta_2\rangle=1$, 
\\
\begin{tabular}{cc}$\alpha_1$+ & $\alpha_2$\\\tiny{ $A_1$+$A_1$' } & \tiny{ $A_2$ }\end{tabular} &\begin{tabular}{c} = \\~ \end{tabular} & \begin{tabular}{cc}$\beta_1$+ & $\beta_2$\\\tiny{ $A_2$+$A_1$ } & \tiny{ $A_1$' }\end{tabular} & $A_3$ & $A_7$ &  $\langle\alpha_1, \beta_1\rangle=1$, $\langle\alpha_1, \beta_2\rangle=1$, $\langle\alpha_2, \beta_1\rangle=1$, $\langle\alpha_2, \beta_2\rangle=1$, 
\\
\begin{tabular}{c}2$\alpha_1$\\\tiny{ $A_1$ }\end{tabular} &\begin{tabular}{c} = \\~ \end{tabular} & \begin{tabular}{cccc}$\beta_1$+ & $\beta_2$+ & $\beta_3$+ & $\beta_4$\\\tiny{ $A_2$+$A_1$ } & \tiny{ $A_2$+$A_1$' } & \tiny{ $A_1$ } & \tiny{ $A_1$' }\end{tabular} & $D_4$ & $E_7$ &  $\langle\alpha_1, \beta_1\rangle=1$, $\langle\alpha_1, \beta_2\rangle=1$, $\langle\alpha_1, \beta_3\rangle=1$, $\langle\alpha_1, \beta_4\rangle=1$, 
\\
\begin{tabular}{cc}$\alpha_1$+ & $\alpha_2$\\\tiny{ $A_1$ } & \tiny{  }\end{tabular} &\begin{tabular}{c} = \\~ \end{tabular} & \begin{tabular}{ccc}$\beta_1$+ & $\beta_2$+ & $\beta_3$\\\tiny{ $A_2$+$A_1$ } & \tiny{ $A_2$+$A_1$' } & \tiny{ $A_1$ }\end{tabular} & $D_4$ & $E_7$ & $\langle\alpha_1, \alpha_2\rangle=1$,  $\langle\alpha_1, \beta_1\rangle=1$, $\langle\alpha_1, \beta_2\rangle=1$, $\langle\alpha_1, \beta_3\rangle=1$, $\langle\alpha_2, \beta_1\rangle=1$, $\langle\alpha_2, \beta_2\rangle=1$, $\langle\alpha_2, \beta_3\rangle=1$, 
\\
\begin{tabular}{cc}$\alpha_1$+ & $\alpha_2$\\\tiny{ $A_1$ } & \tiny{ $A_2$ }\end{tabular} &\begin{tabular}{c} = \\~ \end{tabular} & \begin{tabular}{c}$\beta_1$\\\tiny{ $A_2$+$A_1$ }\end{tabular} & $A_2$ & $A_5$+$A_1$ & $\langle\alpha_1, \alpha_2\rangle=-1$,  $\langle\alpha_1, \beta_1\rangle=1$, $\langle\alpha_2, \beta_1\rangle=1$, 
\\
\begin{tabular}{cc}$\alpha_1$+ & $\alpha_2$\\\tiny{ $A_2$+$A_1$ } & \tiny{  }\end{tabular} &\begin{tabular}{c} = \\~ \end{tabular} & \begin{tabular}{ccc}$\beta_1$+ & $\beta_2$+ & $\beta_3$\\\tiny{ $A_2$+$A_1$ } & \tiny{ $A_2$+$A_1$ } & \tiny{ $A_1$' }\end{tabular} & $D_4$ & $E_7$ & $\langle\alpha_1, \alpha_2\rangle=1$,  $\langle\alpha_1, \beta_1\rangle=1$, $\langle\alpha_1, \beta_2\rangle=1$, $\langle\alpha_1, \beta_3\rangle=1$, $\langle\alpha_2, \beta_1\rangle=1$, $\langle\alpha_2, \beta_2\rangle=1$, $\langle\alpha_2, \beta_3\rangle=1$, 
\\
\begin{tabular}{cc}$\alpha_1$+ & $\alpha_2$\\\tiny{ $A_2$+$A_1$ } & \tiny{ $A_1$' }\end{tabular} &\begin{tabular}{c} = \\~ \end{tabular} & \begin{tabular}{c}$\beta_1$\\\tiny{ $A_2$+$A_1$+$A_1$' }\end{tabular} & $A_2$ & $E_6$ & $\langle\alpha_1, \alpha_2\rangle=-1$,  $\langle\alpha_1, \beta_1\rangle=1$, $\langle\alpha_2, \beta_1\rangle=1$, 
\\
\begin{tabular}{cc}$\alpha_1$+ & $\alpha_2$\\\tiny{ $A_2$+$A_1$ } & \tiny{ $A_2$+$A_1$' }\end{tabular} &\begin{tabular}{c} = \\~ \end{tabular} & \begin{tabular}{cc}$\beta_1$+ & $\beta_2$\\\tiny{ $A_2$+$A_1$+$A_1$' } & \tiny{ $A_2$ }\end{tabular} & $A_3$ & $E_7$ &  $\langle\alpha_1, \beta_1\rangle=1$, $\langle\alpha_1, \beta_2\rangle=1$, $\langle\alpha_2, \beta_1\rangle=1$, $\langle\alpha_2, \beta_2\rangle=1$, 
\\
\begin{tabular}{cc}$\alpha_1$+ & $\alpha_2$\\\tiny{ $A_2$+$A_1$ } & \tiny{ $A_1$ }\end{tabular} &\begin{tabular}{c} = \\~ \end{tabular} & \begin{tabular}{cc}$\beta_1$+ & $\beta_2$\\\tiny{ $A_2$+$A_1$+$A_1$' } & \tiny{ $A_1$+$A_1$' }\end{tabular} & $A_3$ & $E_7$ &  $\langle\alpha_1, \beta_1\rangle=1$, $\langle\alpha_1, \beta_2\rangle=1$, $\langle\alpha_2, \beta_1\rangle=1$, $\langle\alpha_2, \beta_2\rangle=1$, 
\\
\begin{tabular}{cc}$\alpha_1$+ & $\alpha_2$\\\tiny{ $A_1$+$A_1$' } & \tiny{ $A_2$ }\end{tabular} &\begin{tabular}{c} = \\~ \end{tabular} & \begin{tabular}{ccc}$\beta_1$+ & $\beta_2$+ & $\beta_3$\\\tiny{ $A_2$+$A_1$ } & \tiny{ $A_1$' } & \tiny{ $A_2$ }\end{tabular} & $D_4$ & $E_7$ & $\langle\alpha_1, \alpha_2\rangle=1$,  $\langle\alpha_1, \beta_1\rangle=1$, $\langle\alpha_1, \beta_2\rangle=1$, $\langle\alpha_1, \beta_3\rangle=1$, $\langle\alpha_2, \beta_1\rangle=1$, $\langle\alpha_2, \beta_2\rangle=1$, $\langle\alpha_2, \beta_3\rangle=1$, 
\\
\begin{tabular}{cc}$\alpha_1$+ & $\alpha_2$\\\tiny{ $A_1$ } & \tiny{ $A_1$' }\end{tabular} &\begin{tabular}{c} = \\~ \end{tabular} & \begin{tabular}{cc}$\beta_1$+ & $\beta_2$\\\tiny{ $A_1$ } & \tiny{ $A_1$' }\end{tabular} & $A_3$ & $A_5$+$A_2$ &  $\langle\alpha_1, \beta_1\rangle=1$, $\langle\alpha_1, \beta_2\rangle=1$, $\langle\alpha_2, \beta_1\rangle=1$, $\langle\alpha_2, \beta_2\rangle=1$, 
\\
\begin{tabular}{cc}$\alpha_1$+ & $\alpha_2$\\\tiny{ $A_2$+$A_1$ } & \tiny{ $A_1$' }\end{tabular} &\begin{tabular}{c} = \\~ \end{tabular} & \begin{tabular}{cc}$\beta_1$+ & $\beta_2$\\\tiny{ $A_2$+$A_1$ } & \tiny{ $A_1$' }\end{tabular} & $A_3$ & $E_7$ &  $\langle\alpha_1, \beta_1\rangle=1$, $\langle\alpha_1, \beta_2\rangle=1$, $\langle\alpha_2, \beta_1\rangle=1$, $\langle\alpha_2, \beta_2\rangle=1$, 
\\
\begin{tabular}{cc}$\alpha_1$+ & $\alpha_2$\\\tiny{ $A_1$ } & \tiny{ $A_2$ }\end{tabular} &\begin{tabular}{c} = \\~ \end{tabular} & \begin{tabular}{ccc}$\beta_1$+ & $\beta_2$+ & $\beta_3$\\\tiny{ $A_2$+$A_1$ } & \tiny{ $A_1$ } & \tiny{ $A_1$' }\end{tabular} & $D_4$ & $E_7$ & $\langle\alpha_1, \alpha_2\rangle=1$,  $\langle\alpha_1, \beta_1\rangle=1$, $\langle\alpha_1, \beta_2\rangle=1$, $\langle\alpha_1, \beta_3\rangle=1$, $\langle\alpha_2, \beta_1\rangle=1$, $\langle\alpha_2, \beta_2\rangle=1$, $\langle\alpha_2, \beta_3\rangle=1$, 
\\
\begin{tabular}{cc}$\alpha_1$+ & $\alpha_2$\\\tiny{ $A_1$ } & \tiny{ $A_2$ }\end{tabular} &\begin{tabular}{c} = \\~ \end{tabular} & \begin{tabular}{ccc}$\beta_1$+ & $\beta_2$+ & $\beta_3$\\\tiny{ $A_2$+$A_1$ } & \tiny{ $A_2$+$A_1$ } & \tiny{ $A_1$' }\end{tabular} & $D_4$ & $E_7$ & $\langle\alpha_1, \alpha_2\rangle=1$,  $\langle\alpha_1, \beta_1\rangle=1$, $\langle\alpha_1, \beta_2\rangle=1$, $\langle\alpha_1, \beta_3\rangle=1$, $\langle\alpha_2, \beta_1\rangle=1$, $\langle\alpha_2, \beta_2\rangle=1$, $\langle\alpha_2, \beta_3\rangle=1$, 
\\
\begin{tabular}{cc}$\alpha_1$+ & $\alpha_2$\\\tiny{ $A_1$+$A_1$' } & \tiny{  }\end{tabular} &\begin{tabular}{c} = \\~ \end{tabular} & \begin{tabular}{cc}$\beta_1$+ & $\beta_2$\\\tiny{ $A_1$ } & \tiny{ $A_1$' }\end{tabular} & $A_3$ & $A_5$+$A_2$ &  $\langle\alpha_1, \beta_1\rangle=1$, $\langle\alpha_1, \beta_2\rangle=1$, $\langle\alpha_2, \beta_1\rangle=1$, $\langle\alpha_2, \beta_2\rangle=1$, 
\\
\begin{tabular}{c}2$\alpha_1$\\\tiny{  }\end{tabular} &\begin{tabular}{c} = \\~ \end{tabular} & \begin{tabular}{cccc}$\beta_1$+ & $\beta_2$+ & $\beta_3$+ & $\beta_4$\\\tiny{ $A_2$+$A_1$ } & \tiny{ $A_2$+$A_1$ } & \tiny{ $A_2$+$A_1$' } & \tiny{ $A_1$' }\end{tabular} & $D_4$ & $E_7$ &  $\langle\alpha_1, \beta_1\rangle=1$, $\langle\alpha_1, \beta_2\rangle=1$, $\langle\alpha_1, \beta_3\rangle=1$, $\langle\alpha_1, \beta_4\rangle=1$, 
\\
\begin{tabular}{cc}$\alpha_1$+ & $\alpha_2$\\\tiny{ $A_2$+$A_1$+$A_1$' } & \tiny{  }\end{tabular} &\begin{tabular}{c} = \\~ \end{tabular} & \begin{tabular}{cc}$\beta_1$+ & $\beta_2$\\\tiny{ $A_2$+$A_1$ } & \tiny{ $A_2$+$A_1$' }\end{tabular} & $A_3$ & $E_7$ &  $\langle\alpha_1, \beta_1\rangle=1$, $\langle\alpha_1, \beta_2\rangle=1$, $\langle\alpha_2, \beta_1\rangle=1$, $\langle\alpha_2, \beta_2\rangle=1$, 
\\
\begin{tabular}{c}2$\alpha_1$\\\tiny{ $A_1$ }\end{tabular} &\begin{tabular}{c} = \\~ \end{tabular} & \begin{tabular}{cccc}$\beta_1$+ & $\beta_2$+ & $\beta_3$+ & $\beta_4$\\\tiny{ $A_2$+$A_1$ } & \tiny{ $A_2$+$A_1$ } & \tiny{ $A_2$+$A_1$' } & \tiny{ $A_1$' }\end{tabular} & $D_4$ & $E_7$ &  $\langle\alpha_1, \beta_1\rangle=1$, $\langle\alpha_1, \beta_2\rangle=1$, $\langle\alpha_1, \beta_3\rangle=1$, $\langle\alpha_1, \beta_4\rangle=1$, 
\\
\begin{tabular}{cc}$\alpha_1$+ & $\alpha_2$\\\tiny{ $A_1$ } & \tiny{  }\end{tabular} &\begin{tabular}{c} = \\~ \end{tabular} & \begin{tabular}{ccc}$\beta_1$+ & $\beta_2$+ & $\beta_3$\\\tiny{ $A_2$+$A_1$ } & \tiny{ $A_2$+$A_1$ } & \tiny{ $A_2$+$A_1$' }\end{tabular} & $D_4$ & $E_7$ & $\langle\alpha_1, \alpha_2\rangle=1$,  $\langle\alpha_1, \beta_1\rangle=1$, $\langle\alpha_1, \beta_2\rangle=1$, $\langle\alpha_1, \beta_3\rangle=1$, $\langle\alpha_2, \beta_1\rangle=1$, $\langle\alpha_2, \beta_2\rangle=1$, $\langle\alpha_2, \beta_3\rangle=1$, 
\\
\begin{tabular}{cc}$\alpha_1$+ & $\alpha_2$\\\tiny{ $A_1$ } & \tiny{ $A_1$' }\end{tabular} &\begin{tabular}{c} = \\~ \end{tabular} & \begin{tabular}{c}$\beta_1$\\\tiny{ $A_1$+$A_1$' }\end{tabular} & $A_2$ & $A_4$+$A_2$ & $\langle\alpha_1, \alpha_2\rangle=-1$,  $\langle\alpha_1, \beta_1\rangle=1$, $\langle\alpha_2, \beta_1\rangle=1$, 
\\
\begin{tabular}{cc}$\alpha_1$+ & $\alpha_2$\\\tiny{ $A_1$+$A_1$' } & \tiny{ $A_2$ }\end{tabular} &\begin{tabular}{c} = \\~ \end{tabular} & \begin{tabular}{c}$\beta_1$\\\tiny{ $A_2$+$A_1$+$A_1$' }\end{tabular} & $A_2$ & $D_6$ & $\langle\alpha_1, \alpha_2\rangle=-1$,  $\langle\alpha_1, \beta_1\rangle=1$, $\langle\alpha_2, \beta_1\rangle=1$, 
\\
\begin{tabular}{cc}$\alpha_1$+ & $\alpha_2$\\\tiny{ $A_1$+$A_1$' } & \tiny{ $A_1$ }\end{tabular} &\begin{tabular}{c} = \\~ \end{tabular} & \begin{tabular}{cc}$\beta_1$+ & $\beta_2$\\\tiny{ $A_2$+$A_1$+$A_1$' } & \tiny{ $A_2$+$A_1$ }\end{tabular} & $A_3$ & $E_7$ &  $\langle\alpha_1, \beta_1\rangle=1$, $\langle\alpha_1, \beta_2\rangle=1$, $\langle\alpha_2, \beta_1\rangle=1$, $\langle\alpha_2, \beta_2\rangle=1$, 
\\
\begin{tabular}{cc}$\alpha_1$+ & $\alpha_2$\\\tiny{ $A_1$ } & \tiny{ $A_2$ }\end{tabular} &\begin{tabular}{c} = \\~ \end{tabular} & \begin{tabular}{cc}$\beta_1$+ & $\beta_2$\\\tiny{ $A_2$+$A_1$+$A_1$' } & \tiny{ $A_1$ }\end{tabular} & $A_3$ & $E_7$ &  $\langle\alpha_1, \beta_1\rangle=1$, $\langle\alpha_1, \beta_2\rangle=1$, $\langle\alpha_2, \beta_1\rangle=1$, $\langle\alpha_2, \beta_2\rangle=1$, 
\\
\begin{tabular}{cc}$\alpha_1$+ & $\alpha_2$\\\tiny{ $A_2$+$A_1$ } & \tiny{  }\end{tabular} &\begin{tabular}{c} = \\~ \end{tabular} & \begin{tabular}{ccc}$\beta_1$+ & $\beta_2$+ & $\beta_3$\\\tiny{ $A_2$+$A_1$ } & \tiny{ $A_1$ } & \tiny{ $A_1$' }\end{tabular} & $D_4$ & $E_7$ & $\langle\alpha_1, \alpha_2\rangle=1$,  $\langle\alpha_1, \beta_1\rangle=1$, $\langle\alpha_1, \beta_2\rangle=1$, $\langle\alpha_1, \beta_3\rangle=1$, $\langle\alpha_2, \beta_1\rangle=1$, $\langle\alpha_2, \beta_2\rangle=1$, $\langle\alpha_2, \beta_3\rangle=1$, 
\\
\begin{tabular}{cc}$\alpha_1$+ & $\alpha_2$\\\tiny{ $A_2$+$A_1$ } & \tiny{ $A_2$+$A_1$' }\end{tabular} &\begin{tabular}{c} = \\~ \end{tabular} & \begin{tabular}{cc}$\beta_1$+ & $\beta_2$\\\tiny{ $A_2$+$A_1$ } & \tiny{ $A_2$+$A_1$' }\end{tabular} & $A_3$ & $E_7$ &  $\langle\alpha_1, \beta_1\rangle=1$, $\langle\alpha_1, \beta_2\rangle=1$, $\langle\alpha_2, \beta_1\rangle=1$, $\langle\alpha_2, \beta_2\rangle=1$, 
\\
\begin{tabular}{cc}$\alpha_1$+ & $\alpha_2$\\\tiny{ $A_2$+$A_1$ } & \tiny{ $A_2$ }\end{tabular} &\begin{tabular}{c} = \\~ \end{tabular} & \begin{tabular}{cc}$\beta_1$+ & $\beta_2$\\\tiny{ $A_2$+$A_1$+$A_1$' } & \tiny{ $A_2$+$A_1$' }\end{tabular} & $A_3$ & $E_7$ &  $\langle\alpha_1, \beta_1\rangle=1$, $\langle\alpha_1, \beta_2\rangle=1$, $\langle\alpha_2, \beta_1\rangle=1$, $\langle\alpha_2, \beta_2\rangle=1$, 
\\
\begin{tabular}{cc}$\alpha_1$+ & $\alpha_2$\\\tiny{ $A_2$+$A_1$ } & \tiny{ $A_1$ }\end{tabular} &\begin{tabular}{c} = \\~ \end{tabular} & \begin{tabular}{cc}$\beta_1$+ & $\beta_2$\\\tiny{ $A_2$+$A_1$ } & \tiny{ $A_1$ }\end{tabular} & $A_3$ & $D_6$+$A_1$ &  $\langle\alpha_1, \beta_1\rangle=1$, $\langle\alpha_1, \beta_2\rangle=1$, $\langle\alpha_2, \beta_1\rangle=1$, $\langle\alpha_2, \beta_2\rangle=1$, 
\\
\begin{tabular}{cc}$\alpha_1$+ & $\alpha_2$\\\tiny{ $A_2$+$A_1$ } & \tiny{ $A_1$' }\end{tabular} &\begin{tabular}{c} = \\~ \end{tabular} & \begin{tabular}{cc}$\beta_1$+ & $\beta_2$\\\tiny{ $A_2$+$A_1$+$A_1$' } & \tiny{ $A_2$ }\end{tabular} & $A_3$ & $E_7$ &  $\langle\alpha_1, \beta_1\rangle=1$, $\langle\alpha_1, \beta_2\rangle=1$, $\langle\alpha_2, \beta_1\rangle=1$, $\langle\alpha_2, \beta_2\rangle=1$, 
\\
\begin{tabular}{cc}$\alpha_1$+ & $\alpha_2$\\\tiny{ $A_2$+$A_1$ } & \tiny{ $A_1$ }\end{tabular} &\begin{tabular}{c} = \\~ \end{tabular} & \begin{tabular}{cc}$\beta_1$+ & $\beta_2$\\\tiny{ $A_2$+$A_1$+$A_1$' } & \tiny{ $A_2$+$A_1$+$A_1$' }\end{tabular} & $A_3$ & $E_6$ &  $\langle\alpha_1, \beta_1\rangle=1$, $\langle\alpha_1, \beta_2\rangle=1$, $\langle\alpha_2, \beta_1\rangle=1$, $\langle\alpha_2, \beta_2\rangle=1$, 
\\
\begin{tabular}{cc}$\alpha_1$+ & $\alpha_2$\\\tiny{ $A_2$+$A_1$+$A_1$' } & \tiny{ $A_2$ }\end{tabular} &\begin{tabular}{c} = \\~ \end{tabular} & \begin{tabular}{ccc}$\beta_1$+ & $\beta_2$+ & $\beta_3$\\\tiny{ $A_2$+$A_1$ } & \tiny{ $A_1$' } & \tiny{ $A_2$ }\end{tabular} & $D_4$ & $E_7$ & $\langle\alpha_1, \alpha_2\rangle=1$,  $\langle\alpha_1, \beta_1\rangle=1$, $\langle\alpha_1, \beta_2\rangle=1$, $\langle\alpha_1, \beta_3\rangle=1$, $\langle\alpha_2, \beta_1\rangle=1$, $\langle\alpha_2, \beta_2\rangle=1$, $\langle\alpha_2, \beta_3\rangle=1$, 
\\
\begin{tabular}{cc}$\alpha_1$+ & $\alpha_2$\\\tiny{ $A_1$ } & \tiny{ $A_2$ }\end{tabular} &\begin{tabular}{c} = \\~ \end{tabular} & \begin{tabular}{cc}$\beta_1$+ & $\beta_2$\\\tiny{ $A_2$+$A_1$ } & \tiny{ $A_2$ }\end{tabular} & $A_3$ & $D_6$+$A_1$ &  $\langle\alpha_1, \beta_1\rangle=1$, $\langle\alpha_1, \beta_2\rangle=1$, $\langle\alpha_2, \beta_1\rangle=1$, $\langle\alpha_2, \beta_2\rangle=1$, 
\\
\begin{tabular}{cc}$\alpha_1$+ & $\alpha_2$\\\tiny{ $A_2$ } & \tiny{ $A_2$ }\end{tabular} &\begin{tabular}{c} = \\~ \end{tabular} & \begin{tabular}{cc}$\beta_1$+ & $\beta_2$\\\tiny{ $A_2$+$A_1$+$A_1$' } & \tiny{ $A_2$+$A_1$+$A_1$' }\end{tabular} & $A_3$ & $D_6$ &  $\langle\alpha_1, \beta_1\rangle=1$, $\langle\alpha_1, \beta_2\rangle=1$, $\langle\alpha_2, \beta_1\rangle=1$, $\langle\alpha_2, \beta_2\rangle=1$, 
\\
\begin{tabular}{cc}$\alpha_1$+ & $\alpha_2$\\\tiny{ $A_1$ } & \tiny{ $A_2$ }\end{tabular} &\begin{tabular}{c} = \\~ \end{tabular} & \begin{tabular}{cc}$\beta_1$+ & $\beta_2$\\\tiny{ $A_2$+$A_1$+$A_1$' } & \tiny{ $A_1$' }\end{tabular} & $A_3$ & $E_7$ &  $\langle\alpha_1, \beta_1\rangle=1$, $\langle\alpha_1, \beta_2\rangle=1$, $\langle\alpha_2, \beta_1\rangle=1$, $\langle\alpha_2, \beta_2\rangle=1$, 
\\
\begin{tabular}{cc}$\alpha_1$+ & $\alpha_2$\\\tiny{ $A_2$+$A_1$+$A_1$' } & \tiny{ $A_2$ }\end{tabular} &\begin{tabular}{c} = \\~ \end{tabular} & \begin{tabular}{cc}$\beta_1$+ & $\beta_2$\\\tiny{ $A_2$+$A_1$ } & \tiny{ $A_2$+$A_1$' }\end{tabular} & $A_3$ & $E_7$ &  $\langle\alpha_1, \beta_1\rangle=1$, $\langle\alpha_1, \beta_2\rangle=1$, $\langle\alpha_2, \beta_1\rangle=1$, $\langle\alpha_2, \beta_2\rangle=1$, 
\\
\begin{tabular}{cc}$\alpha_1$+ & $\alpha_2$\\\tiny{ $A_1$ } & \tiny{ $A_2$ }\end{tabular} &\begin{tabular}{c} = \\~ \end{tabular} & \begin{tabular}{cc}$\beta_1$+ & $\beta_2$\\\tiny{ $A_1$+$A_1$' } & \tiny{ $A_2$+$A_1$' }\end{tabular} & $A_3$ & $E_7$ &  $\langle\alpha_1, \beta_1\rangle=1$, $\langle\alpha_1, \beta_2\rangle=1$, $\langle\alpha_2, \beta_1\rangle=1$, $\langle\alpha_2, \beta_2\rangle=1$, 
\\
\begin{tabular}{cc}$\alpha_1$+ & $\alpha_2$\\\tiny{ $A_1$+$A_1$' } & \tiny{ $A_1$' }\end{tabular} &\begin{tabular}{c} = \\~ \end{tabular} & \begin{tabular}{cc}$\beta_1$+ & $\beta_2$\\\tiny{ $A_2$+$A_1$+$A_1$' } & \tiny{ $A_2$+$A_1$' }\end{tabular} & $A_3$ & $E_7$ &  $\langle\alpha_1, \beta_1\rangle=1$, $\langle\alpha_1, \beta_2\rangle=1$, $\langle\alpha_2, \beta_1\rangle=1$, $\langle\alpha_2, \beta_2\rangle=1$, 
\\
\begin{tabular}{cc}$\alpha_1$+ & $\alpha_2$\\\tiny{ $A_1$ } & \tiny{ $A_1$' }\end{tabular} &\begin{tabular}{c} = \\~ \end{tabular} & \begin{tabular}{cc}$\beta_1$+ & $\beta_2$\\\tiny{ $A_2$+$A_1$ } & \tiny{ $A_2$+$A_1$' }\end{tabular} & $A_3$ & $E_6$ &  $\langle\alpha_1, \beta_1\rangle=1$, $\langle\alpha_1, \beta_2\rangle=1$, $\langle\alpha_2, \beta_1\rangle=1$, $\langle\alpha_2, \beta_2\rangle=1$, 
\\
\begin{tabular}{cc}$\alpha_1$+ & $\alpha_2$\\\tiny{ $A_1$ } & \tiny{ $A_1$' }\end{tabular} &\begin{tabular}{c} = \\~ \end{tabular} & \begin{tabular}{cc}$\beta_1$+ & $\beta_2$\\\tiny{ $A_2$+$A_1$+$A_1$' } & \tiny{ $A_2$ }\end{tabular} & $A_3$ & $E_6$ &  $\langle\alpha_1, \beta_1\rangle=1$, $\langle\alpha_1, \beta_2\rangle=1$, $\langle\alpha_2, \beta_1\rangle=1$, $\langle\alpha_2, \beta_2\rangle=1$, 
\\
\begin{tabular}{cc}$\alpha_1$+ & $\alpha_2$\\\tiny{ $A_1$ } & \tiny{ $A_2$ }\end{tabular} &\begin{tabular}{c} = \\~ \end{tabular} & \begin{tabular}{cc}$\beta_1$+ & $\beta_2$\\\tiny{ $A_1$+$A_1$' } & \tiny{ $A_2$+$A_1$ }\end{tabular} & $A_3$ & $E_7$ &  $\langle\alpha_1, \beta_1\rangle=1$, $\langle\alpha_1, \beta_2\rangle=1$, $\langle\alpha_2, \beta_1\rangle=1$, $\langle\alpha_2, \beta_2\rangle=1$, 
\\
\begin{tabular}{cc}$\alpha_1$+ & $\alpha_2$\\\tiny{ $A_2$ } & \tiny{ $A_2$ }\end{tabular} &\begin{tabular}{c} = \\~ \end{tabular} & \begin{tabular}{cc}$\beta_1$+ & $\beta_2$\\\tiny{ $A_2$+$A_1$ } & \tiny{ $A_2$+$A_1$ }\end{tabular} & $A_3$ & $D_6$+$A_1$ &  $\langle\alpha_1, \beta_1\rangle=1$, $\langle\alpha_1, \beta_2\rangle=1$, $\langle\alpha_2, \beta_1\rangle=1$, $\langle\alpha_2, \beta_2\rangle=1$, 
\\
\begin{tabular}{c}2$\alpha_1$\\\tiny{ $A_2$ }\end{tabular} &\begin{tabular}{c} = \\~ \end{tabular} & \begin{tabular}{cccc}$\beta_1$+ & $\beta_2$+ & $\beta_3$+ & $\beta_4$\\\tiny{ $A_2$+$A_1$ } & \tiny{ $A_2$+$A_1$ } & \tiny{ $A_2$+$A_1$' } & \tiny{ $A_1$' }\end{tabular} & $D_4$ & $E_7$ &  $\langle\alpha_1, \beta_1\rangle=1$, $\langle\alpha_1, \beta_2\rangle=1$, $\langle\alpha_1, \beta_3\rangle=1$, $\langle\alpha_1, \beta_4\rangle=1$, 
\\
\begin{tabular}{c}2$\alpha_1$\\\tiny{ $A_2$+$A_1$ }\end{tabular} &\begin{tabular}{c} = \\~ \end{tabular} & \begin{tabular}{cccc}$\beta_1$+ & $\beta_2$+ & $\beta_3$+ & $\beta_4$\\\tiny{ $A_2$+$A_1$ } & \tiny{ $A_2$+$A_1$ } & \tiny{ $A_2$+$A_1$' } & \tiny{ $A_1$' }\end{tabular} & $D_4$ & $E_7$ &  $\langle\alpha_1, \beta_1\rangle=1$, $\langle\alpha_1, \beta_2\rangle=1$, $\langle\alpha_1, \beta_3\rangle=1$, $\langle\alpha_1, \beta_4\rangle=1$, 
\\
\begin{tabular}{cc}$\alpha_1$+ & $\alpha_2$\\\tiny{ $A_2$+$A_1$ } & \tiny{ $A_2$ }\end{tabular} &\begin{tabular}{c} = \\~ \end{tabular} & \begin{tabular}{ccc}$\beta_1$+ & $\beta_2$+ & $\beta_3$\\\tiny{ $A_2$+$A_1$ } & \tiny{ $A_2$+$A_1$ } & \tiny{ $A_2$+$A_1$' }\end{tabular} & $D_4$ & $E_7$ & $\langle\alpha_1, \alpha_2\rangle=1$,  $\langle\alpha_1, \beta_1\rangle=1$, $\langle\alpha_1, \beta_2\rangle=1$, $\langle\alpha_1, \beta_3\rangle=1$, $\langle\alpha_2, \beta_1\rangle=1$, $\langle\alpha_2, \beta_2\rangle=1$, $\langle\alpha_2, \beta_3\rangle=1$, 
\\
\begin{tabular}{cc}$\alpha_1$+ & $\alpha_2$\\\tiny{ $A_2$+$A_1$ } & \tiny{ $A_1$' }\end{tabular} &\begin{tabular}{c} = \\~ \end{tabular} & \begin{tabular}{cc}$\beta_1$+ & $\beta_2$\\\tiny{ $A_1$+$A_1$' } & \tiny{ $A_2$ }\end{tabular} & $A_3$ & $A_7$ &  $\langle\alpha_1, \beta_1\rangle=1$, $\langle\alpha_1, \beta_2\rangle=1$, $\langle\alpha_2, \beta_1\rangle=1$, $\langle\alpha_2, \beta_2\rangle=1$, 
\\
\begin{tabular}{cc}$\alpha_1$+ & $\alpha_2$\\\tiny{ $A_2$+$A_1$ } & \tiny{ $A_2$+$A_1$ }\end{tabular} &\begin{tabular}{c} = \\~ \end{tabular} & \begin{tabular}{ccc}$\beta_1$+ & $\beta_2$+ & $\beta_3$\\\tiny{ $A_2$+$A_1$+$A_1$' } & \tiny{ $A_2$+$A_1$+$A_1$' } & \tiny{ $A_2$ }\end{tabular} & $D_4$ & $E_7$ & $\langle\alpha_1, \alpha_2\rangle=1$,  $\langle\alpha_1, \beta_1\rangle=1$, $\langle\alpha_1, \beta_2\rangle=1$, $\langle\alpha_1, \beta_3\rangle=1$, $\langle\alpha_2, \beta_1\rangle=1$, $\langle\alpha_2, \beta_2\rangle=1$, $\langle\alpha_2, \beta_3\rangle=1$, 
\\
\begin{tabular}{cc}$\alpha_1$+ & $\alpha_2$\\\tiny{ $A_1$ } & \tiny{ $A_1$ }\end{tabular} &\begin{tabular}{c} = \\~ \end{tabular} & \begin{tabular}{ccc}$\beta_1$+ & $\beta_2$+ & $\beta_3$\\\tiny{ $A_2$+$A_1$+$A_1$' } & \tiny{ $A_2$+$A_1$+$A_1$' } & \tiny{ $A_2$ }\end{tabular} & $D_4$ & $E_7$ & $\langle\alpha_1, \alpha_2\rangle=1$,  $\langle\alpha_1, \beta_1\rangle=1$, $\langle\alpha_1, \beta_2\rangle=1$, $\langle\alpha_1, \beta_3\rangle=1$, $\langle\alpha_2, \beta_1\rangle=1$, $\langle\alpha_2, \beta_2\rangle=1$, $\langle\alpha_2, \beta_3\rangle=1$, 
\\
\begin{tabular}{c}2$\alpha_1$\\\tiny{ $A_1$+$A_1$' }\end{tabular} &\begin{tabular}{c} = \\~ \end{tabular} & \begin{tabular}{cccc}$\beta_1$+ & $\beta_2$+ & $\beta_3$+ & $\beta_4$\\\tiny{ $A_2$+$A_1$ } & \tiny{ $A_2$+$A_1$' } & \tiny{ $A_1$ } & \tiny{ $A_1$' }\end{tabular} & $D_4$ & $E_7$ &  $\langle\alpha_1, \beta_1\rangle=1$, $\langle\alpha_1, \beta_2\rangle=1$, $\langle\alpha_1, \beta_3\rangle=1$, $\langle\alpha_1, \beta_4\rangle=1$, 
\\
\begin{tabular}{cc}$\alpha_1$+ & $\alpha_2$\\\tiny{ $A_1$+$A_1$' } & \tiny{ $A_2$+$A_1$ }\end{tabular} &\begin{tabular}{c} = \\~ \end{tabular} & \begin{tabular}{ccc}$\beta_1$+ & $\beta_2$+ & $\beta_3$\\\tiny{ $A_2$+$A_1$ } & \tiny{ $A_1$' } & \tiny{ $A_1$ }\end{tabular} & $D_4$ & $E_7$ & $\langle\alpha_1, \alpha_2\rangle=1$,  $\langle\alpha_1, \beta_1\rangle=1$, $\langle\alpha_1, \beta_2\rangle=1$, $\langle\alpha_1, \beta_3\rangle=1$, $\langle\alpha_2, \beta_1\rangle=1$, $\langle\alpha_2, \beta_2\rangle=1$, $\langle\alpha_2, \beta_3\rangle=1$, 
\\
\begin{tabular}{cc}$\alpha_1$+ & $\alpha_2$\\\tiny{ $A_1$+$A_1$' } & \tiny{ $A_1$' }\end{tabular} &\begin{tabular}{c} = \\~ \end{tabular} & \begin{tabular}{ccc}$\beta_1$+ & $\beta_2$+ & $\beta_3$\\\tiny{ $A_2$+$A_1$ } & \tiny{ $A_2$+$A_1$' } & \tiny{ $A_1$ }\end{tabular} & $D_4$ & $E_7$ & $\langle\alpha_1, \alpha_2\rangle=1$,  $\langle\alpha_1, \beta_1\rangle=1$, $\langle\alpha_1, \beta_2\rangle=1$, $\langle\alpha_1, \beta_3\rangle=1$, $\langle\alpha_2, \beta_1\rangle=1$, $\langle\alpha_2, \beta_2\rangle=1$, $\langle\alpha_2, \beta_3\rangle=1$, 
\\
\begin{tabular}{cc}$\alpha_1$+ & $\alpha_2$\\\tiny{ $A_2$+$A_1$ } & \tiny{ $A_1$' }\end{tabular} &\begin{tabular}{c} = \\~ \end{tabular} & \begin{tabular}{ccc}$\beta_1$+ & $\beta_2$+ & $\beta_3$\\\tiny{ $A_2$+$A_1$ } & \tiny{ $A_1$' } & \tiny{ $A_2$ }\end{tabular} & $D_4$ & $E_7$ & $\langle\alpha_1, \alpha_2\rangle=1$,  $\langle\alpha_1, \beta_1\rangle=1$, $\langle\alpha_1, \beta_2\rangle=1$, $\langle\alpha_1, \beta_3\rangle=1$, $\langle\alpha_2, \beta_1\rangle=1$, $\langle\alpha_2, \beta_2\rangle=1$, $\langle\alpha_2, \beta_3\rangle=1$, 
\\
\begin{tabular}{c}2$\alpha_1$\\\tiny{  }\end{tabular} &\begin{tabular}{c} = \\~ \end{tabular} & \begin{tabular}{cccc}$\beta_1$+ & $\beta_2$+ & $\beta_3$+ & $\beta_4$\\\tiny{ $A_2$+$A_1$ } & \tiny{ $A_2$+$A_1$' } & \tiny{ $A_2$+$A_1$' } & \tiny{ $A_1$ }\end{tabular} & $D_4$ & $E_7$ &  $\langle\alpha_1, \beta_1\rangle=1$, $\langle\alpha_1, \beta_2\rangle=1$, $\langle\alpha_1, \beta_3\rangle=1$, $\langle\alpha_1, \beta_4\rangle=1$, 
\\
\begin{tabular}{c}2$\alpha_1$\\\tiny{ $A_1$ }\end{tabular} &\begin{tabular}{c} = \\~ \end{tabular} & \begin{tabular}{cccc}$\beta_1$+ & $\beta_2$+ & $\beta_3$+ & $\beta_4$\\\tiny{ $A_1$+$A_1$' } & \tiny{ $A_2$+$A_1$' } & \tiny{ $A_2$+$A_1$ } & \tiny{ $A_2$ }\end{tabular} & $D_4$ & $E_7$ &  $\langle\alpha_1, \beta_1\rangle=1$, $\langle\alpha_1, \beta_2\rangle=1$, $\langle\alpha_1, \beta_3\rangle=1$, $\langle\alpha_1, \beta_4\rangle=1$, 
\\
\begin{tabular}{cc}$\alpha_1$+ & $\alpha_2$\\\tiny{ $A_1$ } & \tiny{ $A_2$ }\end{tabular} &\begin{tabular}{c} = \\~ \end{tabular} & \begin{tabular}{ccc}$\beta_1$+ & $\beta_2$+ & $\beta_3$\\\tiny{ $A_1$+$A_1$' } & \tiny{ $A_2$+$A_1$' } & \tiny{ $A_2$ }\end{tabular} & $D_4$ & $E_7$ & $\langle\alpha_1, \alpha_2\rangle=1$,  $\langle\alpha_1, \beta_1\rangle=1$, $\langle\alpha_1, \beta_2\rangle=1$, $\langle\alpha_1, \beta_3\rangle=1$, $\langle\alpha_2, \beta_1\rangle=1$, $\langle\alpha_2, \beta_2\rangle=1$, $\langle\alpha_2, \beta_3\rangle=1$, 
\\
\begin{tabular}{c}2$\alpha_1$\\\tiny{ $A_2$ }\end{tabular} &\begin{tabular}{c} = \\~ \end{tabular} & \begin{tabular}{cccc}$\beta_1$+ & $\beta_2$+ & $\beta_3$+ & $\beta_4$\\\tiny{ $A_2$+$A_1$ } & \tiny{ $A_2$+$A_1$' } & \tiny{ $A_2$+$A_1$' } & \tiny{ $A_1$ }\end{tabular} & $D_4$ & $E_7$ &  $\langle\alpha_1, \beta_1\rangle=1$, $\langle\alpha_1, \beta_2\rangle=1$, $\langle\alpha_1, \beta_3\rangle=1$, $\langle\alpha_1, \beta_4\rangle=1$, 
\\
\begin{tabular}{cc}$\alpha_1$+ & $\alpha_2$\\\tiny{ $A_2$+$A_1$ } & \tiny{ $A_2$ }\end{tabular} &\begin{tabular}{c} = \\~ \end{tabular} & \begin{tabular}{ccc}$\beta_1$+ & $\beta_2$+ & $\beta_3$\\\tiny{ $A_1$+$A_1$' } & \tiny{ $A_2$+$A_1$' } & \tiny{ $A_2$ }\end{tabular} & $D_4$ & $E_7$ & $\langle\alpha_1, \alpha_2\rangle=1$,  $\langle\alpha_1, \beta_1\rangle=1$, $\langle\alpha_1, \beta_2\rangle=1$, $\langle\alpha_1, \beta_3\rangle=1$, $\langle\alpha_2, \beta_1\rangle=1$, $\langle\alpha_2, \beta_2\rangle=1$, $\langle\alpha_2, \beta_3\rangle=1$, 
\\
\begin{tabular}{c}2$\alpha_1$\\\tiny{ $A_1$ }\end{tabular} &\begin{tabular}{c} = \\~ \end{tabular} & \begin{tabular}{cccc}$\beta_1$+ & $\beta_2$+ & $\beta_3$+ & $\beta_4$\\\tiny{ $A_2$+$A_1$ } & \tiny{ $A_2$+$A_1$' } & \tiny{ $A_2$+$A_1$' } & \tiny{ $A_1$ }\end{tabular} & $D_4$ & $E_7$ &  $\langle\alpha_1, \beta_1\rangle=1$, $\langle\alpha_1, \beta_2\rangle=1$, $\langle\alpha_1, \beta_3\rangle=1$, $\langle\alpha_1, \beta_4\rangle=1$, 
\\
\begin{tabular}{cc}$\alpha_1$+ & $\alpha_2$\\\tiny{ $A_1$ } & \tiny{ $A_1$' }\end{tabular} &\begin{tabular}{c} = \\~ \end{tabular} & \begin{tabular}{ccc}$\beta_1$+ & $\beta_2$+ & $\beta_3$\\\tiny{ $A_2$+$A_1$ } & \tiny{ $A_2$+$A_1$' } & \tiny{ $A_2$ }\end{tabular} & $D_4$ & $E_7$ & $\langle\alpha_1, \alpha_2\rangle=1$,  $\langle\alpha_1, \beta_1\rangle=1$, $\langle\alpha_1, \beta_2\rangle=1$, $\langle\alpha_1, \beta_3\rangle=1$, $\langle\alpha_2, \beta_1\rangle=1$, $\langle\alpha_2, \beta_2\rangle=1$, $\langle\alpha_2, \beta_3\rangle=1$, 
\\
\begin{tabular}{cc}$\alpha_1$+ & $\alpha_2$\\\tiny{ $A_1$ } & \tiny{ $A_1$ }\end{tabular} &\begin{tabular}{c} = \\~ \end{tabular} & \begin{tabular}{cc}$\beta_1$+ & $\beta_2$\\\tiny{ $A_2$+$A_1$ } & \tiny{ $A_2$+$A_1$ }\end{tabular} & $A_3$ & $D_6$+$A_1$ &  $\langle\alpha_1, \beta_1\rangle=1$, $\langle\alpha_1, \beta_2\rangle=1$, $\langle\alpha_2, \beta_1\rangle=1$, $\langle\alpha_2, \beta_2\rangle=1$, 
\\
\begin{tabular}{cc}$\alpha_1$+ & $\alpha_2$\\\tiny{ $A_2$+$A_1$ } & \tiny{ $A_1$' }\end{tabular} &\begin{tabular}{c} = \\~ \end{tabular} & \begin{tabular}{cc}$\beta_1$+ & $\beta_2$\\\tiny{ $A_2$+$A_1$ } & \tiny{ $A_1$' }\end{tabular} & $A_3$ & $A_7$ &  $\langle\alpha_1, \beta_1\rangle=1$, $\langle\alpha_1, \beta_2\rangle=1$, $\langle\alpha_2, \beta_1\rangle=1$, $\langle\alpha_2, \beta_2\rangle=1$, 
\\
\begin{tabular}{cc}$\alpha_1$+ & $\alpha_2$\\\tiny{ $A_1$+$A_1$' } & \tiny{ $A_2$+$A_1$' }\end{tabular} &\begin{tabular}{c} = \\~ \end{tabular} & \begin{tabular}{cc}$\beta_1$+ & $\beta_2$\\\tiny{ $A_2$+$A_1$+$A_1$' } & \tiny{ $A_1$' }\end{tabular} & $A_3$ & $E_7$ &  $\langle\alpha_1, \beta_1\rangle=1$, $\langle\alpha_1, \beta_2\rangle=1$, $\langle\alpha_2, \beta_1\rangle=1$, $\langle\alpha_2, \beta_2\rangle=1$, 
\\
\begin{tabular}{cc}$\alpha_1$+ & $\alpha_2$\\\tiny{ $A_1$+$A_1$' } & \tiny{ $A_2$+$A_1$ }\end{tabular} &\begin{tabular}{c} = \\~ \end{tabular} & \begin{tabular}{cc}$\beta_1$+ & $\beta_2$\\\tiny{ $A_2$+$A_1$+$A_1$' } & \tiny{ $A_2$+$A_1$ }\end{tabular} & $A_3$ & $E_7$ &  $\langle\alpha_1, \beta_1\rangle=1$, $\langle\alpha_1, \beta_2\rangle=1$, $\langle\alpha_2, \beta_1\rangle=1$, $\langle\alpha_2, \beta_2\rangle=1$, 
\\
\begin{tabular}{cc}$\alpha_1$+ & $\alpha_2$\\\tiny{ $A_2$+$A_1$ } & \tiny{ $A_2$ }\end{tabular} &\begin{tabular}{c} = \\~ \end{tabular} & \begin{tabular}{cc}$\beta_1$+ & $\beta_2$\\\tiny{ $A_2$+$A_1$+$A_1$' } & \tiny{ $A_2$+$A_1$ }\end{tabular} & $A_3$ & $E_7$ &  $\langle\alpha_1, \beta_1\rangle=1$, $\langle\alpha_1, \beta_2\rangle=1$, $\langle\alpha_2, \beta_1\rangle=1$, $\langle\alpha_2, \beta_2\rangle=1$, 
\\
\begin{tabular}{cc}$\alpha_1$+ & $\alpha_2$\\\tiny{ $A_2$+$A_1$+$A_1$' } & \tiny{ $A_1$ }\end{tabular} &\begin{tabular}{c} = \\~ \end{tabular} & \begin{tabular}{cc}$\beta_1$+ & $\beta_2$\\\tiny{ $A_1$+$A_1$' } & \tiny{ $A_2$+$A_1$ }\end{tabular} & $A_3$ & $E_7$ &  $\langle\alpha_1, \beta_1\rangle=1$, $\langle\alpha_1, \beta_2\rangle=1$, $\langle\alpha_2, \beta_1\rangle=1$, $\langle\alpha_2, \beta_2\rangle=1$, 
\\
\begin{tabular}{cc}$\alpha_1$+ & $\alpha_2$\\\tiny{ $A_2$+$A_1$ } & \tiny{ $A_2$ }\end{tabular} &\begin{tabular}{c} = \\~ \end{tabular} & \begin{tabular}{cc}$\beta_1$+ & $\beta_2$\\\tiny{ $A_2$+$A_1$ } & \tiny{ $A_2$ }\end{tabular} & $A_3$ & $D_6$+$A_1$ &  $\langle\alpha_1, \beta_1\rangle=1$, $\langle\alpha_1, \beta_2\rangle=1$, $\langle\alpha_2, \beta_1\rangle=1$, $\langle\alpha_2, \beta_2\rangle=1$, 
\\
\begin{tabular}{cc}$\alpha_1$+ & $\alpha_2$\\\tiny{ $A_2$+$A_1$ } & \tiny{ $A_1$' }\end{tabular} &\begin{tabular}{c} = \\~ \end{tabular} & \begin{tabular}{cc}$\beta_1$+ & $\beta_2$\\\tiny{ $A_2$+$A_1$ } & \tiny{ $A_2$+$A_1$' }\end{tabular} & $A_3$ & $E_7$ &  $\langle\alpha_1, \beta_1\rangle=1$, $\langle\alpha_1, \beta_2\rangle=1$, $\langle\alpha_2, \beta_1\rangle=1$, $\langle\alpha_2, \beta_2\rangle=1$, 
\\
\begin{tabular}{cc}$\alpha_1$+ & $\alpha_2$\\\tiny{ $A_2$+$A_1$ } & \tiny{ $A_2$+$A_1$' }\end{tabular} &\begin{tabular}{c} = \\~ \end{tabular} & \begin{tabular}{ccc}$\beta_1$+ & $\beta_2$+ & $\beta_3$\\\tiny{ $A_2$+$A_1$ } & \tiny{ $A_1$' } & \tiny{ $A_2$ }\end{tabular} & $D_4$ & $E_7$ & $\langle\alpha_1, \alpha_2\rangle=1$,  $\langle\alpha_1, \beta_1\rangle=1$, $\langle\alpha_1, \beta_2\rangle=1$, $\langle\alpha_1, \beta_3\rangle=1$, $\langle\alpha_2, \beta_1\rangle=1$, $\langle\alpha_2, \beta_2\rangle=1$, $\langle\alpha_2, \beta_3\rangle=1$, 
\\
\begin{tabular}{ccc}$\alpha_1$+ & $\alpha_2$+ & $\alpha_3$\\\tiny{ $A_2$+$A_1$ } & \tiny{ $A_1$' } & \tiny{ $A_2$ }\end{tabular} &\begin{tabular}{c} = \\~ \end{tabular} & \begin{tabular}{ccc}$\beta_1$+ & $\beta_2$+ & $\beta_3$\\\tiny{ $A_2$+$A_1$ } & \tiny{ $A_1$' } & \tiny{ $A_2$ }\end{tabular} & $A_5$ & $E_7$ &  $\langle\alpha_1, \beta_1\rangle=1$, $\langle\alpha_1, \beta_3\rangle=1$, $\langle\alpha_2, \beta_2\rangle=1$, $\langle\alpha_2, \beta_3\rangle=1$, $\langle\alpha_3, \beta_1\rangle=1$, $\langle\alpha_3, \beta_2\rangle=1$, 
\\
\begin{tabular}{cc}$\alpha_1$+ & $\alpha_2$\\\tiny{ $A_2$+$A_1$ } & \tiny{ $A_1$' }\end{tabular} &\begin{tabular}{c} = \\~ \end{tabular} & \begin{tabular}{cc}$\beta_1$+ & $\beta_2$\\\tiny{ $A_2$+$A_1$+$A_1$' } & \tiny{  }\end{tabular} & $A_3$ & $E_7$ &  $\langle\alpha_1, \beta_1\rangle=1$, $\langle\alpha_1, \beta_2\rangle=1$, $\langle\alpha_2, \beta_1\rangle=1$, $\langle\alpha_2, \beta_2\rangle=1$, 
\\
\begin{tabular}{c}2$\alpha_1$\\\tiny{ $A_2$+$A_1$ }\end{tabular} &\begin{tabular}{c} = \\~ \end{tabular} & \begin{tabular}{cccc}$\beta_1$+ & $\beta_2$+ & $\beta_3$+ & $\beta_4$\\\tiny{ $A_2$+$A_1$+$A_1$' } & \tiny{ $A_2$+$A_1$+$A_1$' } & \tiny{ $A_2$ } & \tiny{  }\end{tabular} & $D_4$ & $E_7$ &  $\langle\alpha_1, \beta_1\rangle=1$, $\langle\alpha_1, \beta_2\rangle=1$, $\langle\alpha_1, \beta_3\rangle=1$, $\langle\alpha_1, \beta_4\rangle=1$, 
\\
\begin{tabular}{c}2$\alpha_1$\\\tiny{ $A_1$ }\end{tabular} &\begin{tabular}{c} = \\~ \end{tabular} & \begin{tabular}{cccc}$\beta_1$+ & $\beta_2$+ & $\beta_3$+ & $\beta_4$\\\tiny{ $A_2$+$A_1$+$A_1$' } & \tiny{ $A_2$+$A_1$+$A_1$' } & \tiny{ $A_2$ } & \tiny{  }\end{tabular} & $D_4$ & $E_7$ &  $\langle\alpha_1, \beta_1\rangle=1$, $\langle\alpha_1, \beta_2\rangle=1$, $\langle\alpha_1, \beta_3\rangle=1$, $\langle\alpha_1, \beta_4\rangle=1$, 
\\
\begin{tabular}{cc}$\alpha_1$+ & $\alpha_2$\\\tiny{ $A_2$+$A_1$ } & \tiny{ $A_1$ }\end{tabular} &\begin{tabular}{c} = \\~ \end{tabular} & \begin{tabular}{ccc}$\beta_1$+ & $\beta_2$+ & $\beta_3$\\\tiny{ $A_2$+$A_1$+$A_1$' } & \tiny{ $A_2$+$A_1$+$A_1$' } & \tiny{  }\end{tabular} & $D_4$ & $E_7$ & $\langle\alpha_1, \alpha_2\rangle=1$,  $\langle\alpha_1, \beta_1\rangle=1$, $\langle\alpha_1, \beta_2\rangle=1$, $\langle\alpha_1, \beta_3\rangle=1$, $\langle\alpha_2, \beta_1\rangle=1$, $\langle\alpha_2, \beta_2\rangle=1$, $\langle\alpha_2, \beta_3\rangle=1$, 
\\
\begin{tabular}{c}2$\alpha_1$\\\tiny{ $A_2$+$A_1$ }\end{tabular} &\begin{tabular}{c} = \\~ \end{tabular} & \begin{tabular}{cccc}$\beta_1$+ & $\beta_2$+ & $\beta_3$+ & $\beta_4$\\\tiny{ $A_1$+$A_1$' } & \tiny{ $A_2$+$A_1$' } & \tiny{ $A_2$+$A_1$ } & \tiny{ $A_2$ }\end{tabular} & $D_4$ & $E_7$ &  $\langle\alpha_1, \beta_1\rangle=1$, $\langle\alpha_1, \beta_2\rangle=1$, $\langle\alpha_1, \beta_3\rangle=1$, $\langle\alpha_1, \beta_4\rangle=1$, 
\\
\begin{tabular}{cc}$\alpha_1$+ & $\alpha_2$\\\tiny{ $A_2$+$A_1$ } & \tiny{ $A_1$ }\end{tabular} &\begin{tabular}{c} = \\~ \end{tabular} & \begin{tabular}{ccc}$\beta_1$+ & $\beta_2$+ & $\beta_3$\\\tiny{ $A_1$+$A_1$' } & \tiny{ $A_2$+$A_1$' } & \tiny{ $A_2$+$A_1$ }\end{tabular} & $D_4$ & $E_7$ & $\langle\alpha_1, \alpha_2\rangle=1$,  $\langle\alpha_1, \beta_1\rangle=1$, $\langle\alpha_1, \beta_2\rangle=1$, $\langle\alpha_1, \beta_3\rangle=1$, $\langle\alpha_2, \beta_1\rangle=1$, $\langle\alpha_2, \beta_2\rangle=1$, $\langle\alpha_2, \beta_3\rangle=1$, 
\\
\begin{tabular}{cc}$\alpha_1$+ & $\alpha_2$\\\tiny{ $A_2$+$A_1$ } & \tiny{ $A_2$+$A_1$' }\end{tabular} &\begin{tabular}{c} = \\~ \end{tabular} & \begin{tabular}{ccc}$\beta_1$+ & $\beta_2$+ & $\beta_3$\\\tiny{ $A_2$+$A_1$ } & \tiny{ $A_2$+$A_1$' } & \tiny{ $A_2$ }\end{tabular} & $D_4$ & $E_7$ & $\langle\alpha_1, \alpha_2\rangle=1$,  $\langle\alpha_1, \beta_1\rangle=1$, $\langle\alpha_1, \beta_2\rangle=1$, $\langle\alpha_1, \beta_3\rangle=1$, $\langle\alpha_2, \beta_1\rangle=1$, $\langle\alpha_2, \beta_2\rangle=1$, $\langle\alpha_2, \beta_3\rangle=1$, 
\\
\begin{tabular}{cc}$\alpha_1$+ & $\alpha_2$\\\tiny{ $A_1$ } & \tiny{ $A_1$' }\end{tabular} &\begin{tabular}{c} = \\~ \end{tabular} & \begin{tabular}{cc}$\beta_1$+ & $\beta_2$\\\tiny{ $A_1$+$A_1$' } & \tiny{  }\end{tabular} & $A_3$ & $A_5$+$A_2$ &  $\langle\alpha_1, \beta_1\rangle=1$, $\langle\alpha_1, \beta_2\rangle=1$, $\langle\alpha_2, \beta_1\rangle=1$, $\langle\alpha_2, \beta_2\rangle=1$, 
\\
\begin{tabular}{cc}$\alpha_1$+ & $\alpha_2$\\\tiny{ $A_2$+$A_1$ } & \tiny{ $A_1$' }\end{tabular} &\begin{tabular}{c} = \\~ \end{tabular} & \begin{tabular}{ccc}$\beta_1$+ & $\beta_2$+ & $\beta_3$\\\tiny{ $A_1$+$A_1$' } & \tiny{ $A_2$ } & \tiny{  }\end{tabular} & $D_4$ & $E_7$ & $\langle\alpha_1, \alpha_2\rangle=1$,  $\langle\alpha_1, \beta_1\rangle=1$, $\langle\alpha_1, \beta_2\rangle=1$, $\langle\alpha_1, \beta_3\rangle=1$, $\langle\alpha_2, \beta_1\rangle=1$, $\langle\alpha_2, \beta_2\rangle=1$, $\langle\alpha_2, \beta_3\rangle=1$, 
\\
\begin{tabular}{cc}$\alpha_1$+ & $\alpha_2$\\\tiny{ $A_1$ } & \tiny{ $A_1$' }\end{tabular} &\begin{tabular}{c} = \\~ \end{tabular} & \begin{tabular}{ccc}$\beta_1$+ & $\beta_2$+ & $\beta_3$\\\tiny{ $A_2$+$A_1$+$A_1$' } & \tiny{ $A_2$ } & \tiny{  }\end{tabular} & $D_4$ & $E_7$ & $\langle\alpha_1, \alpha_2\rangle=1$,  $\langle\alpha_1, \beta_1\rangle=1$, $\langle\alpha_1, \beta_2\rangle=1$, $\langle\alpha_1, \beta_3\rangle=1$, $\langle\alpha_2, \beta_1\rangle=1$, $\langle\alpha_2, \beta_2\rangle=1$, $\langle\alpha_2, \beta_3\rangle=1$, 
\\
\hline\multicolumn{5}{c}{$\mathfrak{k}$-semisimple type: $A_1$+$A_1$+$A_1$+$A_1$}\\
\hline\hline\begin{tabular}{cc}$\alpha_1$+ & $\alpha_2$\\\tiny{ $A_1$+$A_1$'+$A_1$'' } & \tiny{  }\end{tabular} &\begin{tabular}{c} = \\~ \end{tabular} & \begin{tabular}{cc}$\beta_1$+ & $\beta_2$\\\tiny{ $A_1$+$A_1$' } & \tiny{ $A_1$'' }\end{tabular} & $A_3$ & $D_6$+$A_1$ &  $\langle\alpha_1, \beta_1\rangle=1$, $\langle\alpha_1, \beta_2\rangle=1$, $\langle\alpha_2, \beta_1\rangle=1$, $\langle\alpha_2, \beta_2\rangle=1$, 
\\
\begin{tabular}{cc}$\alpha_1$+ & $\alpha_2$\\\tiny{ $A_1$+$A_1$'+$A_1$'' } & \tiny{  }\end{tabular} &\begin{tabular}{c} = \\~ \end{tabular} & \begin{tabular}{ccc}$\beta_1$+ & $\beta_2$+ & $\beta_3$\\\tiny{ $A_1$+$A_1$' } & \tiny{ $A_1$''+$A_1$' } & \tiny{ $A_1$''' }\end{tabular} & $D_4$ & $E_7$ & $\langle\alpha_1, \alpha_2\rangle=1$,  $\langle\alpha_1, \beta_1\rangle=1$, $\langle\alpha_1, \beta_2\rangle=1$, $\langle\alpha_1, \beta_3\rangle=1$, $\langle\alpha_2, \beta_1\rangle=1$, $\langle\alpha_2, \beta_2\rangle=1$, $\langle\alpha_2, \beta_3\rangle=1$, 
\\
\begin{tabular}{cc}$\alpha_1$+ & $\alpha_2$\\\tiny{ $A_1$ } & \tiny{  }\end{tabular} &\begin{tabular}{c} = \\~ \end{tabular} & \begin{tabular}{c}$\beta_1$\\\tiny{ $A_1$ }\end{tabular} & $A_2$ & $A_3$+3$A_1$ & $\langle\alpha_1, \alpha_2\rangle=-1$,  $\langle\alpha_1, \beta_1\rangle=1$, $\langle\alpha_2, \beta_1\rangle=1$, 
\\
\begin{tabular}{cc}$\alpha_1$+ & $\alpha_2$\\\tiny{ $A_1$ } & \tiny{  }\end{tabular} &\begin{tabular}{c} = \\~ \end{tabular} & \begin{tabular}{ccc}$\beta_1$+ & $\beta_2$+ & $\beta_3$\\\tiny{ $A_1$+$A_1$' } & \tiny{ $A_1$+$A_1$' } & \tiny{ $A_1$'' }\end{tabular} & $D_4$ & $D_6$+$A_1$ & $\langle\alpha_1, \alpha_2\rangle=1$,  $\langle\alpha_1, \beta_1\rangle=1$, $\langle\alpha_1, \beta_2\rangle=1$, $\langle\alpha_1, \beta_3\rangle=1$, $\langle\alpha_2, \beta_1\rangle=1$, $\langle\alpha_2, \beta_2\rangle=1$, $\langle\alpha_2, \beta_3\rangle=1$, 
\\
\begin{tabular}{cc}$\alpha_1$+ & $\alpha_2$\\\tiny{ $A_1$+$A_1$' } & \tiny{  }\end{tabular} &\begin{tabular}{c} = \\~ \end{tabular} & \begin{tabular}{cc}$\beta_1$+ & $\beta_2$\\\tiny{ $A_1$+$A_1$' } & \tiny{ $A_1$''+$A_1$' }\end{tabular} & $A_3$ & $D_6$+$A_1$ &  $\langle\alpha_1, \beta_1\rangle=1$, $\langle\alpha_1, \beta_2\rangle=1$, $\langle\alpha_2, \beta_1\rangle=1$, $\langle\alpha_2, \beta_2\rangle=1$, 
\\
\begin{tabular}{cc}$\alpha_1$+ & $\alpha_2$\\\tiny{ $A_1$+$A_1$' } & \tiny{ $A_1$'' }\end{tabular} &\begin{tabular}{c} = \\~ \end{tabular} & \begin{tabular}{ccc}$\beta_1$+ & $\beta_2$+ & $\beta_3$\\\tiny{ $A_1$+$A_1$' } & \tiny{ $A_1$''+$A_1$ } & \tiny{ $A_1$''' }\end{tabular} & $D_4$ & $E_7$ & $\langle\alpha_1, \alpha_2\rangle=1$,  $\langle\alpha_1, \beta_1\rangle=1$, $\langle\alpha_1, \beta_2\rangle=1$, $\langle\alpha_1, \beta_3\rangle=1$, $\langle\alpha_2, \beta_1\rangle=1$, $\langle\alpha_2, \beta_2\rangle=1$, $\langle\alpha_2, \beta_3\rangle=1$, 
\\
\begin{tabular}{c}2$\alpha_1$\\\tiny{  }\end{tabular} &\begin{tabular}{c} = \\~ \end{tabular} & \begin{tabular}{cccc}$\beta_1$+ & $\beta_2$+ & $\beta_3$+ & $\beta_4$\\\tiny{ $A_1$+$A_1$' } & \tiny{ $A_1$+$A_1$' } & \tiny{ $A_1$'' } & \tiny{ $A_1$'' }\end{tabular} & $D_4$ & $D_6$+$A_1$ &  $\langle\alpha_1, \beta_1\rangle=1$, $\langle\alpha_1, \beta_2\rangle=1$, $\langle\alpha_1, \beta_3\rangle=1$, $\langle\alpha_1, \beta_4\rangle=1$, 
\\
\begin{tabular}{cc}$\alpha_1$+ & $\alpha_2$\\\tiny{  } & \tiny{  }\end{tabular} &\begin{tabular}{c} = \\~ \end{tabular} & \begin{tabular}{cc}$\beta_1$+ & $\beta_2$\\\tiny{ $A_1$ } & \tiny{ $A_1$ }\end{tabular} & $A_3$ & $D_4$+3$A_1$ &  $\langle\alpha_1, \beta_1\rangle=1$, $\langle\alpha_1, \beta_2\rangle=1$, $\langle\alpha_2, \beta_1\rangle=1$, $\langle\alpha_2, \beta_2\rangle=1$, 
\\
\begin{tabular}{cc}$\alpha_1$+ & $\alpha_2$\\\tiny{  } & \tiny{  }\end{tabular} &\begin{tabular}{c} = \\~ \end{tabular} & \begin{tabular}{cc}$\beta_1$+ & $\beta_2$\\\tiny{ $A_1$+$A_1$' } & \tiny{ $A_1$+$A_1$' }\end{tabular} & $A_3$ & $D_4$+2$A_1$ &  $\langle\alpha_1, \beta_1\rangle=1$, $\langle\alpha_1, \beta_2\rangle=1$, $\langle\alpha_2, \beta_1\rangle=1$, $\langle\alpha_2, \beta_2\rangle=1$, 
\\
\begin{tabular}{cc}$\alpha_1$+ & $\alpha_2$\\\tiny{ $A_1$+$A_1$' } & \tiny{  }\end{tabular} &\begin{tabular}{c} = \\~ \end{tabular} & \begin{tabular}{c}$\beta_1$\\\tiny{ $A_1$+$A_1$' }\end{tabular} & $A_2$ & $D_4$+2$A_1$ & $\langle\alpha_1, \alpha_2\rangle=-1$,  $\langle\alpha_1, \beta_1\rangle=1$, $\langle\alpha_2, \beta_1\rangle=1$, 
\\
\begin{tabular}{cc}$\alpha_1$+ & $\alpha_2$\\\tiny{ $A_1$+$A_1$' } & \tiny{  }\end{tabular} &\begin{tabular}{c} = \\~ \end{tabular} & \begin{tabular}{cc}$\beta_1$+ & $\beta_2$\\\tiny{ $A_1$+$A_1$' } & \tiny{ $A_1$+$A_1$'' }\end{tabular} & $A_3$ & $D_6$+$A_1$ &  $\langle\alpha_1, \beta_1\rangle=1$, $\langle\alpha_1, \beta_2\rangle=1$, $\langle\alpha_2, \beta_1\rangle=1$, $\langle\alpha_2, \beta_2\rangle=1$, 
\\
\begin{tabular}{cc}$\alpha_1$+ & $\alpha_2$\\\tiny{ $A_1$ } & \tiny{ $A_1$ }\end{tabular} &\begin{tabular}{c} = \\~ \end{tabular} & \begin{tabular}{cc}$\beta_1$+ & $\beta_2$\\\tiny{ $A_1$ } & \tiny{ $A_1$ }\end{tabular} & $A_3$ & $D_4$+3$A_1$ &  $\langle\alpha_1, \beta_1\rangle=1$, $\langle\alpha_1, \beta_2\rangle=1$, $\langle\alpha_2, \beta_1\rangle=1$, $\langle\alpha_2, \beta_2\rangle=1$, 
\\
\begin{tabular}{cc}$\alpha_1$+ & $\alpha_2$\\\tiny{ $A_1$+$A_1$' } & \tiny{ $A_1$'' }\end{tabular} &\begin{tabular}{c} = \\~ \end{tabular} & \begin{tabular}{cc}$\beta_1$+ & $\beta_2$\\\tiny{ $A_1$+$A_1$' } & \tiny{ $A_1$'' }\end{tabular} & $A_3$ & $D_6$+$A_1$ &  $\langle\alpha_1, \beta_1\rangle=1$, $\langle\alpha_1, \beta_2\rangle=1$, $\langle\alpha_2, \beta_1\rangle=1$, $\langle\alpha_2, \beta_2\rangle=1$, 
\\
\begin{tabular}{c}2$\alpha_1$\\\tiny{ $A_1$ }\end{tabular} &\begin{tabular}{c} = \\~ \end{tabular} & \begin{tabular}{cccc}$\beta_1$+ & $\beta_2$+ & $\beta_3$+ & $\beta_4$\\\tiny{ $A_1$+$A_1$' } & \tiny{ $A_1$+$A_1$' } & \tiny{ $A_1$'' } & \tiny{ $A_1$'' }\end{tabular} & $D_4$ & $D_6$+$A_1$ &  $\langle\alpha_1, \beta_1\rangle=1$, $\langle\alpha_1, \beta_2\rangle=1$, $\langle\alpha_1, \beta_3\rangle=1$, $\langle\alpha_1, \beta_4\rangle=1$, 
\\
\begin{tabular}{cc}$\alpha_1$+ & $\alpha_2$\\\tiny{ $A_1$+$A_1$' } & \tiny{ $A_1$'' }\end{tabular} &\begin{tabular}{c} = \\~ \end{tabular} & \begin{tabular}{c}$\beta_1$\\\tiny{ $A_1$+$A_1$'+$A_1$'' }\end{tabular} & $A_2$ & $D_5$+$A_1$ & $\langle\alpha_1, \alpha_2\rangle=-1$,  $\langle\alpha_1, \beta_1\rangle=1$, $\langle\alpha_2, \beta_1\rangle=1$, 
\\
\begin{tabular}{cc}$\alpha_1$+ & $\alpha_2$\\\tiny{ $A_1$+$A_1$'+$A_1$'' } & \tiny{ $A_1$'+$A_1$''' }\end{tabular} &\begin{tabular}{c} = \\~ \end{tabular} & \begin{tabular}{ccc}$\beta_1$+ & $\beta_2$+ & $\beta_3$\\\tiny{ $A_1$+$A_1$' } & \tiny{ $A_1$+$A_1$'' } & \tiny{ $A_1$''' }\end{tabular} & $D_4$ & $E_7$ & $\langle\alpha_1, \alpha_2\rangle=1$,  $\langle\alpha_1, \beta_1\rangle=1$, $\langle\alpha_1, \beta_2\rangle=1$, $\langle\alpha_1, \beta_3\rangle=1$, $\langle\alpha_2, \beta_1\rangle=1$, $\langle\alpha_2, \beta_2\rangle=1$, $\langle\alpha_2, \beta_3\rangle=1$, 
\\
\begin{tabular}{cc}$\alpha_1$+ & $\alpha_2$\\\tiny{ $A_1$+$A_1$' } & \tiny{  }\end{tabular} &\begin{tabular}{c} = \\~ \end{tabular} & \begin{tabular}{cc}$\beta_1$+ & $\beta_2$\\\tiny{ $A_1$+$A_1$' } & \tiny{ $A_1$''+$A_1$ }\end{tabular} & $A_3$ & $D_6$+$A_1$ &  $\langle\alpha_1, \beta_1\rangle=1$, $\langle\alpha_1, \beta_2\rangle=1$, $\langle\alpha_2, \beta_1\rangle=1$, $\langle\alpha_2, \beta_2\rangle=1$, 
\\
\begin{tabular}{cc}$\alpha_1$+ & $\alpha_2$\\\tiny{ $A_1$+$A_1$' } & \tiny{ $A_1$'' }\end{tabular} &\begin{tabular}{c} = \\~ \end{tabular} & \begin{tabular}{cc}$\beta_1$+ & $\beta_2$\\\tiny{ $A_1$+$A_1$'+$A_1$'' } & \tiny{ $A_1$''+$A_1$''' }\end{tabular} & $A_3$ & $E_7$ &  $\langle\alpha_1, \beta_1\rangle=1$, $\langle\alpha_1, \beta_2\rangle=1$, $\langle\alpha_2, \beta_1\rangle=1$, $\langle\alpha_2, \beta_2\rangle=1$, 
\\
\begin{tabular}{cc}$\alpha_1$+ & $\alpha_2$\\\tiny{ $A_1$ } & \tiny{ $A_1$ }\end{tabular} &\begin{tabular}{c} = \\~ \end{tabular} & \begin{tabular}{cc}$\beta_1$+ & $\beta_2$\\\tiny{ $A_1$+$A_1$'+$A_1$'' } & \tiny{ $A_1$+$A_1$'+$A_1$'' }\end{tabular} & $A_3$ & $D_5$+$A_1$ &  $\langle\alpha_1, \beta_1\rangle=1$, $\langle\alpha_1, \beta_2\rangle=1$, $\langle\alpha_2, \beta_1\rangle=1$, $\langle\alpha_2, \beta_2\rangle=1$, 
\\
\begin{tabular}{cc}$\alpha_1$+ & $\alpha_2$\\\tiny{ $A_1$+$A_1$'+$A_1$'' } & \tiny{ $A_1$''+$A_1$''' }\end{tabular} &\begin{tabular}{c} = \\~ \end{tabular} & \begin{tabular}{ccc}$\beta_1$+ & $\beta_2$+ & $\beta_3$\\\tiny{ $A_1$+$A_1$' } & \tiny{ $A_1$''+$A_1$ } & \tiny{ $A_1$''' }\end{tabular} & $D_4$ & $E_7$ & $\langle\alpha_1, \alpha_2\rangle=1$,  $\langle\alpha_1, \beta_1\rangle=1$, $\langle\alpha_1, \beta_2\rangle=1$, $\langle\alpha_1, \beta_3\rangle=1$, $\langle\alpha_2, \beta_1\rangle=1$, $\langle\alpha_2, \beta_2\rangle=1$, $\langle\alpha_2, \beta_3\rangle=1$, 
\\
\begin{tabular}{cc}$\alpha_1$+ & $\alpha_2$\\\tiny{ $A_1$+$A_1$'+$A_1$'' } & \tiny{  }\end{tabular} &\begin{tabular}{c} = \\~ \end{tabular} & \begin{tabular}{ccc}$\beta_1$+ & $\beta_2$+ & $\beta_3$\\\tiny{ $A_1$+$A_1$' } & \tiny{ $A_1$+$A_1$'' } & \tiny{ $A_1$''' }\end{tabular} & $D_4$ & $E_7$ & $\langle\alpha_1, \alpha_2\rangle=1$,  $\langle\alpha_1, \beta_1\rangle=1$, $\langle\alpha_1, \beta_2\rangle=1$, $\langle\alpha_1, \beta_3\rangle=1$, $\langle\alpha_2, \beta_1\rangle=1$, $\langle\alpha_2, \beta_2\rangle=1$, $\langle\alpha_2, \beta_3\rangle=1$, 
\\
\begin{tabular}{cc}$\alpha_1$+ & $\alpha_2$\\\tiny{ $A_1$+$A_1$'+$A_1$'' } & \tiny{ $A_1$'''+$A_1$'' }\end{tabular} &\begin{tabular}{c} = \\~ \end{tabular} & \begin{tabular}{cc}$\beta_1$+ & $\beta_2$\\\tiny{ $A_1$+$A_1$'+$A_1$'' } & \tiny{ $A_1$'''+$A_1$'' }\end{tabular} & $A_3$ & $E_7$ &  $\langle\alpha_1, \beta_1\rangle=1$, $\langle\alpha_1, \beta_2\rangle=1$, $\langle\alpha_2, \beta_1\rangle=1$, $\langle\alpha_2, \beta_2\rangle=1$, 
\\
\begin{tabular}{cc}$\alpha_1$+ & $\alpha_2$\\\tiny{ $A_1$+$A_1$' } & \tiny{ $A_1$+$A_1$'' }\end{tabular} &\begin{tabular}{c} = \\~ \end{tabular} & \begin{tabular}{cc}$\beta_1$+ & $\beta_2$\\\tiny{ $A_1$+$A_1$' } & \tiny{ $A_1$+$A_1$'' }\end{tabular} & $A_3$ & $D_6$+$A_1$ &  $\langle\alpha_1, \beta_1\rangle=1$, $\langle\alpha_1, \beta_2\rangle=1$, $\langle\alpha_2, \beta_1\rangle=1$, $\langle\alpha_2, \beta_2\rangle=1$, 
\\
\begin{tabular}{cc}$\alpha_1$+ & $\alpha_2$\\\tiny{ $A_1$+$A_1$' } & \tiny{ $A_1$+$A_1$'' }\end{tabular} &\begin{tabular}{c} = \\~ \end{tabular} & \begin{tabular}{cc}$\beta_1$+ & $\beta_2$\\\tiny{ $A_1$+$A_1$'+$A_1$'' } & \tiny{ $A_1$+$A_1$'+$A_1$''' }\end{tabular} & $A_3$ & $E_7$ &  $\langle\alpha_1, \beta_1\rangle=1$, $\langle\alpha_1, \beta_2\rangle=1$, $\langle\alpha_2, \beta_1\rangle=1$, $\langle\alpha_2, \beta_2\rangle=1$, 
\\
\begin{tabular}{cc}$\alpha_1$+ & $\alpha_2$\\\tiny{ $A_1$+$A_1$' } & \tiny{ $A_1$'' }\end{tabular} &\begin{tabular}{c} = \\~ \end{tabular} & \begin{tabular}{cc}$\beta_1$+ & $\beta_2$\\\tiny{ $A_1$+$A_1$'+$A_1$'' } & \tiny{ $A_1$'+$A_1$''' }\end{tabular} & $A_3$ & $E_7$ &  $\langle\alpha_1, \beta_1\rangle=1$, $\langle\alpha_1, \beta_2\rangle=1$, $\langle\alpha_2, \beta_1\rangle=1$, $\langle\alpha_2, \beta_2\rangle=1$, 
\\
\begin{tabular}{cc}$\alpha_1$+ & $\alpha_2$\\\tiny{ $A_1$+$A_1$'+$A_1$'' } & \tiny{ $A_1$ }\end{tabular} &\begin{tabular}{c} = \\~ \end{tabular} & \begin{tabular}{cc}$\beta_1$+ & $\beta_2$\\\tiny{ $A_1$+$A_1$'+$A_1$'' } & \tiny{ $A_1$+$A_1$'+$A_1$''' }\end{tabular} & $A_3$ & $E_7$ &  $\langle\alpha_1, \beta_1\rangle=1$, $\langle\alpha_1, \beta_2\rangle=1$, $\langle\alpha_2, \beta_1\rangle=1$, $\langle\alpha_2, \beta_2\rangle=1$, 
\\
\begin{tabular}{cc}$\alpha_1$+ & $\alpha_2$\\\tiny{ $A_1$+$A_1$' } & \tiny{ $A_1$'' }\end{tabular} &\begin{tabular}{c} = \\~ \end{tabular} & \begin{tabular}{cc}$\beta_1$+ & $\beta_2$\\\tiny{ $A_1$+$A_1$'+$A_1$'' } & \tiny{  }\end{tabular} & $A_3$ & $D_6$+$A_1$ &  $\langle\alpha_1, \beta_1\rangle=1$, $\langle\alpha_1, \beta_2\rangle=1$, $\langle\alpha_2, \beta_1\rangle=1$, $\langle\alpha_2, \beta_2\rangle=1$, 
\\
\begin{tabular}{cc}$\alpha_1$+ & $\alpha_2$\\\tiny{ $A_1$+$A_1$' } & \tiny{ $A_1$'' }\end{tabular} &\begin{tabular}{c} = \\~ \end{tabular} & \begin{tabular}{cc}$\beta_1$+ & $\beta_2$\\\tiny{ $A_1$+$A_1$'+$A_1$'' } & \tiny{ $A_1$'''+$A_1$'' }\end{tabular} & $A_3$ & $E_7$ &  $\langle\alpha_1, \beta_1\rangle=1$, $\langle\alpha_1, \beta_2\rangle=1$, $\langle\alpha_2, \beta_1\rangle=1$, $\langle\alpha_2, \beta_2\rangle=1$, 
\\
\begin{tabular}{cc}$\alpha_1$+ & $\alpha_2$\\\tiny{ $A_1$+$A_1$' } & \tiny{ $A_1$'' }\end{tabular} &\begin{tabular}{c} = \\~ \end{tabular} & \begin{tabular}{ccc}$\beta_1$+ & $\beta_2$+ & $\beta_3$\\\tiny{ $A_1$+$A_1$'+$A_1$'' } & \tiny{ $A_1$'+$A_1$''' } & \tiny{  }\end{tabular} & $D_4$ & $E_7$ & $\langle\alpha_1, \alpha_2\rangle=1$,  $\langle\alpha_1, \beta_1\rangle=1$, $\langle\alpha_1, \beta_2\rangle=1$, $\langle\alpha_1, \beta_3\rangle=1$, $\langle\alpha_2, \beta_1\rangle=1$, $\langle\alpha_2, \beta_2\rangle=1$, $\langle\alpha_2, \beta_3\rangle=1$, 
\\
\begin{tabular}{cc}$\alpha_1$+ & $\alpha_2$\\\tiny{ $A_1$+$A_1$' } & \tiny{ $A_1$+$A_1$' }\end{tabular} &\begin{tabular}{c} = \\~ \end{tabular} & \begin{tabular}{ccc}$\beta_1$+ & $\beta_2$+ & $\beta_3$\\\tiny{ $A_1$+$A_1$'+$A_1$'' } & \tiny{ $A_1$+$A_1$'+$A_1$''' } & \tiny{ $A_1$'''+$A_1$'' }\end{tabular} & $D_4$ & $E_7$ & $\langle\alpha_1, \alpha_2\rangle=1$,  $\langle\alpha_1, \beta_1\rangle=1$, $\langle\alpha_1, \beta_2\rangle=1$, $\langle\alpha_1, \beta_3\rangle=1$, $\langle\alpha_2, \beta_1\rangle=1$, $\langle\alpha_2, \beta_2\rangle=1$, $\langle\alpha_2, \beta_3\rangle=1$, 
\\
\begin{tabular}{cc}$\alpha_1$+ & $\alpha_2$\\\tiny{ $A_1$+$A_1$' } & \tiny{ $A_1$'' }\end{tabular} &\begin{tabular}{c} = \\~ \end{tabular} & \begin{tabular}{ccc}$\beta_1$+ & $\beta_2$+ & $\beta_3$\\\tiny{ $A_1$+$A_1$' } & \tiny{ $A_1$+$A_1$'' } & \tiny{ $A_1$''' }\end{tabular} & $D_4$ & $E_7$ & $\langle\alpha_1, \alpha_2\rangle=1$,  $\langle\alpha_1, \beta_1\rangle=1$, $\langle\alpha_1, \beta_2\rangle=1$, $\langle\alpha_1, \beta_3\rangle=1$, $\langle\alpha_2, \beta_1\rangle=1$, $\langle\alpha_2, \beta_2\rangle=1$, $\langle\alpha_2, \beta_3\rangle=1$, 
\\
\begin{tabular}{cc}$\alpha_1$+ & $\alpha_2$\\\tiny{ $A_1$ } & \tiny{ $A_1$ }\end{tabular} &\begin{tabular}{c} = \\~ \end{tabular} & \begin{tabular}{ccc}$\beta_1$+ & $\beta_2$+ & $\beta_3$\\\tiny{ $A_1$+$A_1$'+$A_1$'' } & \tiny{ $A_1$+$A_1$'+$A_1$'' } & \tiny{  }\end{tabular} & $D_4$ & $D_6$+$A_1$ & $\langle\alpha_1, \alpha_2\rangle=1$,  $\langle\alpha_1, \beta_1\rangle=1$, $\langle\alpha_1, \beta_2\rangle=1$, $\langle\alpha_1, \beta_3\rangle=1$, $\langle\alpha_2, \beta_1\rangle=1$, $\langle\alpha_2, \beta_2\rangle=1$, $\langle\alpha_2, \beta_3\rangle=1$, 
\\
\begin{tabular}{cc}$\alpha_1$+ & $\alpha_2$\\\tiny{ $A_1$+$A_1$'+$A_1$'' } & \tiny{ $A_1$ }\end{tabular} &\begin{tabular}{c} = \\~ \end{tabular} & \begin{tabular}{cc}$\beta_1$+ & $\beta_2$\\\tiny{ $A_1$+$A_1$'+$A_1$'' } & \tiny{ $A_1$+$A_1$'''+$A_1$'' }\end{tabular} & $A_3$ & $E_7$ &  $\langle\alpha_1, \beta_1\rangle=1$, $\langle\alpha_1, \beta_2\rangle=1$, $\langle\alpha_2, \beta_1\rangle=1$, $\langle\alpha_2, \beta_2\rangle=1$, 
\\
\begin{tabular}{cc}$\alpha_1$+ & $\alpha_2$\\\tiny{ $A_1$+$A_1$' } & \tiny{ $A_1$''+$A_1$ }\end{tabular} &\begin{tabular}{c} = \\~ \end{tabular} & \begin{tabular}{cc}$\beta_1$+ & $\beta_2$\\\tiny{ $A_1$+$A_1$'+$A_1$'' } & \tiny{ $A_1$+$A_1$'''+$A_1$' }\end{tabular} & $A_3$ & $E_7$ &  $\langle\alpha_1, \beta_1\rangle=1$, $\langle\alpha_1, \beta_2\rangle=1$, $\langle\alpha_2, \beta_1\rangle=1$, $\langle\alpha_2, \beta_2\rangle=1$, 
\\
\begin{tabular}{cc}$\alpha_1$+ & $\alpha_2$\\\tiny{ $A_1$+$A_1$' } & \tiny{ $A_1$'' }\end{tabular} &\begin{tabular}{c} = \\~ \end{tabular} & \begin{tabular}{cc}$\beta_1$+ & $\beta_2$\\\tiny{ $A_1$+$A_1$'+$A_1$'' } & \tiny{ $A_1$'''+$A_1$' }\end{tabular} & $A_3$ & $E_7$ &  $\langle\alpha_1, \beta_1\rangle=1$, $\langle\alpha_1, \beta_2\rangle=1$, $\langle\alpha_2, \beta_1\rangle=1$, $\langle\alpha_2, \beta_2\rangle=1$, 
\\
\begin{tabular}{c}2$\alpha_1$\\\tiny{ $A_1$+$A_1$' }\end{tabular} &\begin{tabular}{c} = \\~ \end{tabular} & \begin{tabular}{cccc}$\beta_1$+ & $\beta_2$+ & $\beta_3$+ & $\beta_4$\\\tiny{ $A_1$+$A_1$'+$A_1$'' } & \tiny{ $A_1$+$A_1$'''+$A_1$' } & \tiny{ $A_1$'''+$A_1$'' } & \tiny{  }\end{tabular} & $D_4$ & $E_7$ &  $\langle\alpha_1, \beta_1\rangle=1$, $\langle\alpha_1, \beta_2\rangle=1$, $\langle\alpha_1, \beta_3\rangle=1$, $\langle\alpha_1, \beta_4\rangle=1$, 
\\
\begin{tabular}{cc}$\alpha_1$+ & $\alpha_2$\\\tiny{ $A_1$+$A_1$' } & \tiny{ $A_1$''+$A_1$ }\end{tabular} &\begin{tabular}{c} = \\~ \end{tabular} & \begin{tabular}{ccc}$\beta_1$+ & $\beta_2$+ & $\beta_3$\\\tiny{ $A_1$+$A_1$'+$A_1$'' } & \tiny{ $A_1$+$A_1$'''+$A_1$' } & \tiny{  }\end{tabular} & $D_4$ & $E_7$ & $\langle\alpha_1, \alpha_2\rangle=1$,  $\langle\alpha_1, \beta_1\rangle=1$, $\langle\alpha_1, \beta_2\rangle=1$, $\langle\alpha_1, \beta_3\rangle=1$, $\langle\alpha_2, \beta_1\rangle=1$, $\langle\alpha_2, \beta_2\rangle=1$, $\langle\alpha_2, \beta_3\rangle=1$, 
\\
\begin{tabular}{cc}$\alpha_1$+ & $\alpha_2$\\\tiny{ $A_1$+$A_1$' } & \tiny{ $A_1$'' }\end{tabular} &\begin{tabular}{c} = \\~ \end{tabular} & \begin{tabular}{ccc}$\beta_1$+ & $\beta_2$+ & $\beta_3$\\\tiny{ $A_1$+$A_1$'+$A_1$'' } & \tiny{ $A_1$'''+$A_1$'' } & \tiny{  }\end{tabular} & $D_4$ & $E_7$ & $\langle\alpha_1, \alpha_2\rangle=1$,  $\langle\alpha_1, \beta_1\rangle=1$, $\langle\alpha_1, \beta_2\rangle=1$, $\langle\alpha_1, \beta_3\rangle=1$, $\langle\alpha_2, \beta_1\rangle=1$, $\langle\alpha_2, \beta_2\rangle=1$, $\langle\alpha_2, \beta_3\rangle=1$, 
\\
\begin{tabular}{cc}$\alpha_1$+ & $\alpha_2$\\\tiny{ $A_1$+$A_1$'+$A_1$'' } & \tiny{ $A_1$ }\end{tabular} &\begin{tabular}{c} = \\~ \end{tabular} & \begin{tabular}{cc}$\beta_1$+ & $\beta_2$\\\tiny{ $A_1$+$A_1$'+$A_1$'' } & \tiny{ $A_1$+$A_1$'''+$A_1$' }\end{tabular} & $A_3$ & $E_7$ &  $\langle\alpha_1, \beta_1\rangle=1$, $\langle\alpha_1, \beta_2\rangle=1$, $\langle\alpha_2, \beta_1\rangle=1$, $\langle\alpha_2, \beta_2\rangle=1$, 
\\
\begin{tabular}{c}2$\alpha_1$\\\tiny{ $A_1$ }\end{tabular} &\begin{tabular}{c} = \\~ \end{tabular} & \begin{tabular}{cccc}$\beta_1$+ & $\beta_2$+ & $\beta_3$+ & $\beta_4$\\\tiny{ $A_1$+$A_1$'+$A_1$'' } & \tiny{ $A_1$+$A_1$'+$A_1$'' } & \tiny{  } & \tiny{  }\end{tabular} & $D_4$ & $D_6$+$A_1$ &  $\langle\alpha_1, \beta_1\rangle=1$, $\langle\alpha_1, \beta_2\rangle=1$, $\langle\alpha_1, \beta_3\rangle=1$, $\langle\alpha_1, \beta_4\rangle=1$, 
\\
\hline\multicolumn{5}{c}{$\mathfrak{k}$-semisimple type: $A_1$+$A_1$+$A_1$+$A_1$}\\
\hline\hline\begin{tabular}{c}2$\alpha_1$\\\tiny{  }\end{tabular} &\begin{tabular}{c} = \\~ \end{tabular} & \begin{tabular}{cccc}$\beta_1$+ & $\beta_2$+ & $\beta_3$+ & $\beta_4$\\\tiny{ $A_1$+$A_1$' } & \tiny{ $A_1$+$A_1$' } & \tiny{ $A_1$''+$A_1$''' } & \tiny{ $A_1$''+$A_1$''' }\end{tabular} & $D_4$ & $D_6$ &  $\langle\alpha_1, \beta_1\rangle=1$, $\langle\alpha_1, \beta_2\rangle=1$, $\langle\alpha_1, \beta_3\rangle=1$, $\langle\alpha_1, \beta_4\rangle=1$, 
\\
\begin{tabular}{c}2$\alpha_1$\\\tiny{ $A_1$+$A_1$' }\end{tabular} &\begin{tabular}{c} = \\~ \end{tabular} & \begin{tabular}{cccc}$\beta_1$+ & $\beta_2$+ & $\beta_3$+ & $\beta_4$\\\tiny{ $A_1$+$A_1$' } & \tiny{ $A_1$''+$A_1$ } & \tiny{ $A_1$'''+$A_1$' } & \tiny{ $A_1$'''+$A_1$'' }\end{tabular} & $D_4$ & $E_7$ &  $\langle\alpha_1, \beta_1\rangle=1$, $\langle\alpha_1, \beta_2\rangle=1$, $\langle\alpha_1, \beta_3\rangle=1$, $\langle\alpha_1, \beta_4\rangle=1$, 
\\
\begin{tabular}{cc}$\alpha_1$+ & $\alpha_2$\\\tiny{ $A_1$+$A_1$' } & \tiny{  }\end{tabular} &\begin{tabular}{c} = \\~ \end{tabular} & \begin{tabular}{ccc}$\beta_1$+ & $\beta_2$+ & $\beta_3$\\\tiny{ $A_1$+$A_1$' } & \tiny{ $A_1$+$A_1$'' } & \tiny{ $A_1$'''+$A_1$' }\end{tabular} & $D_4$ & $E_7$ & $\langle\alpha_1, \alpha_2\rangle=1$,  $\langle\alpha_1, \beta_1\rangle=1$, $\langle\alpha_1, \beta_2\rangle=1$, $\langle\alpha_1, \beta_3\rangle=1$, $\langle\alpha_2, \beta_1\rangle=1$, $\langle\alpha_2, \beta_2\rangle=1$, $\langle\alpha_2, \beta_3\rangle=1$, 
\\
\begin{tabular}{cc}$\alpha_1$+ & $\alpha_2$\\\tiny{ $A_1$+$A_1$' } & \tiny{  }\end{tabular} &\begin{tabular}{c} = \\~ \end{tabular} & \begin{tabular}{ccc}$\beta_1$+ & $\beta_2$+ & $\beta_3$\\\tiny{ $A_1$+$A_1$' } & \tiny{ $A_1$''+$A_1$' } & \tiny{ $A_1$'''+$A_1$'' }\end{tabular} & $D_4$ & $E_7$ & $\langle\alpha_1, \alpha_2\rangle=1$,  $\langle\alpha_1, \beta_1\rangle=1$, $\langle\alpha_1, \beta_2\rangle=1$, $\langle\alpha_1, \beta_3\rangle=1$, $\langle\alpha_2, \beta_1\rangle=1$, $\langle\alpha_2, \beta_2\rangle=1$, $\langle\alpha_2, \beta_3\rangle=1$, 
\\
\begin{tabular}{cc}$\alpha_1$+ & $\alpha_2$\\\tiny{ $A_1$+$A_1$' } & \tiny{ $A_1$''+$A_1$''' }\end{tabular} &\begin{tabular}{c} = \\~ \end{tabular} & \begin{tabular}{cc}$\beta_1$+ & $\beta_2$\\\tiny{ $A_1$+$A_1$' } & \tiny{ $A_1$''+$A_1$''' }\end{tabular} & $A_3$ & $A_7$ &  $\langle\alpha_1, \beta_1\rangle=1$, $\langle\alpha_1, \beta_2\rangle=1$, $\langle\alpha_2, \beta_1\rangle=1$, $\langle\alpha_2, \beta_2\rangle=1$, 
\\
\begin{tabular}{cc}$\alpha_1$+ & $\alpha_2$\\\tiny{ $A_1$+$A_1$' } & \tiny{ $A_1$''+$A_1$''' }\end{tabular} &\begin{tabular}{c} = \\~ \end{tabular} & \begin{tabular}{cc}$\beta_1$+ & $\beta_2$\\\tiny{ $A_1$+$A_1$' } & \tiny{ $A_1$''+$A_1$''' }\end{tabular} & $A_3$ & $D_6$ &  $\langle\alpha_1, \beta_1\rangle=1$, $\langle\alpha_1, \beta_2\rangle=1$, $\langle\alpha_2, \beta_1\rangle=1$, $\langle\alpha_2, \beta_2\rangle=1$, 
\\
\begin{tabular}{c}2$\alpha_1$\\\tiny{ $A_1$+$A_1$' }\end{tabular} &\begin{tabular}{c} = \\~ \end{tabular} & \begin{tabular}{cccc}$\beta_1$+ & $\beta_2$+ & $\beta_3$+ & $\beta_4$\\\tiny{ $A_1$+$A_1$' } & \tiny{ $A_1$+$A_1$'' } & \tiny{ $A_1$'''+$A_1$' } & \tiny{ $A_1$'''+$A_1$'' }\end{tabular} & $D_4$ & $E_7$ &  $\langle\alpha_1, \beta_1\rangle=1$, $\langle\alpha_1, \beta_2\rangle=1$, $\langle\alpha_1, \beta_3\rangle=1$, $\langle\alpha_1, \beta_4\rangle=1$, 
\\
\begin{tabular}{cc}$\alpha_1$+ & $\alpha_2$\\\tiny{ $A_1$+$A_1$' } & \tiny{ $A_1$+$A_1$'' }\end{tabular} &\begin{tabular}{c} = \\~ \end{tabular} & \begin{tabular}{ccc}$\beta_1$+ & $\beta_2$+ & $\beta_3$\\\tiny{ $A_1$+$A_1$' } & \tiny{ $A_1$''+$A_1$' } & \tiny{ $A_1$''+$A_1$''' }\end{tabular} & $D_4$ & $E_7$ & $\langle\alpha_1, \alpha_2\rangle=1$,  $\langle\alpha_1, \beta_1\rangle=1$, $\langle\alpha_1, \beta_2\rangle=1$, $\langle\alpha_1, \beta_3\rangle=1$, $\langle\alpha_2, \beta_1\rangle=1$, $\langle\alpha_2, \beta_2\rangle=1$, $\langle\alpha_2, \beta_3\rangle=1$, 
\\
\begin{tabular}{cc}$\alpha_1$+ & $\alpha_2$\\\tiny{ $A_1$+$A_1$' } & \tiny{ $A_1$''+$A_1$''' }\end{tabular} &\begin{tabular}{c} = \\~ \end{tabular} & \begin{tabular}{ccc}$\beta_1$+ & $\beta_2$+ & $\beta_3$\\\tiny{ $A_1$+$A_1$' } & \tiny{ $A_1$''+$A_1$''' } & \tiny{  }\end{tabular} & $D_4$ & $E_7$ & $\langle\alpha_1, \alpha_2\rangle=1$,  $\langle\alpha_1, \beta_1\rangle=1$, $\langle\alpha_1, \beta_2\rangle=1$, $\langle\alpha_1, \beta_3\rangle=1$, $\langle\alpha_2, \beta_1\rangle=1$, $\langle\alpha_2, \beta_2\rangle=1$, $\langle\alpha_2, \beta_3\rangle=1$, 
\\
\begin{tabular}{cc}$\alpha_1$+ & $\alpha_2$\\\tiny{ $A_1$+$A_1$' } & \tiny{ $A_1$''+$A_1$' }\end{tabular} &\begin{tabular}{c} = \\~ \end{tabular} & \begin{tabular}{cc}$\beta_1$+ & $\beta_2$\\\tiny{ $A_1$+$A_1$' } & \tiny{ $A_1$''+$A_1$' }\end{tabular} & $A_3$ & $D_6$+$A_1$ &  $\langle\alpha_1, \beta_1\rangle=1$, $\langle\alpha_1, \beta_2\rangle=1$, $\langle\alpha_2, \beta_1\rangle=1$, $\langle\alpha_2, \beta_2\rangle=1$, 
\\
\begin{tabular}{cc}$\alpha_1$+ & $\alpha_2$\\\tiny{ $A_1$+$A_1$' } & \tiny{ $A_1$''+$A_1$' }\end{tabular} &\begin{tabular}{c} = \\~ \end{tabular} & \begin{tabular}{ccc}$\beta_1$+ & $\beta_2$+ & $\beta_3$\\\tiny{ $A_1$+$A_1$' } & \tiny{ $A_1$+$A_1$'' } & \tiny{ $A_1$'''+$A_1$' }\end{tabular} & $D_4$ & $E_7$ & $\langle\alpha_1, \alpha_2\rangle=1$,  $\langle\alpha_1, \beta_1\rangle=1$, $\langle\alpha_1, \beta_2\rangle=1$, $\langle\alpha_1, \beta_3\rangle=1$, $\langle\alpha_2, \beta_1\rangle=1$, $\langle\alpha_2, \beta_2\rangle=1$, $\langle\alpha_2, \beta_3\rangle=1$, 
\\
\hline\multicolumn{5}{c}{$\mathfrak{k}$-semisimple type: $A_3$}\\
\hline\hline\begin{tabular}{cc}$\alpha_1$+ & $\alpha_2$\\\tiny{  } & \tiny{  }\end{tabular} &\begin{tabular}{c} = \\~ \end{tabular} & \begin{tabular}{cc}$\beta_1$+ & $\beta_2$\\\tiny{ $A_3$ } & \tiny{ $A_3$ }\end{tabular} & $A_3$ & $D_6$ &  $\langle\alpha_1, \beta_1\rangle=1$, $\langle\alpha_1, \beta_2\rangle=1$, $\langle\alpha_2, \beta_1\rangle=1$, $\langle\alpha_2, \beta_2\rangle=1$, 
\\
\begin{tabular}{cc}$\alpha_1$+ & $\alpha_2$\\\tiny{ $A_3$ } & \tiny{  }\end{tabular} &\begin{tabular}{c} = \\~ \end{tabular} & \begin{tabular}{c}$\beta_1$\\\tiny{ $A_3$ }\end{tabular} & $A_2$ & $A_5$ & $\langle\alpha_1, \alpha_2\rangle=-1$,  $\langle\alpha_1, \beta_1\rangle=1$, $\langle\alpha_2, \beta_1\rangle=1$, 
\\
\begin{tabular}{cc}$\alpha_1$+ & $\alpha_2$\\\tiny{ $A_3$ } & \tiny{  }\end{tabular} &\begin{tabular}{c} = \\~ \end{tabular} & \begin{tabular}{cc}$\beta_1$+ & $\beta_2$\\\tiny{ $A_3$ } & \tiny{ $A_3$ }\end{tabular} & $A_3$ & $E_6$ &  $\langle\alpha_1, \beta_1\rangle=1$, $\langle\alpha_1, \beta_2\rangle=1$, $\langle\alpha_2, \beta_1\rangle=1$, $\langle\alpha_2, \beta_2\rangle=1$, 
\\
\begin{tabular}{cc}$\alpha_1$+ & $\alpha_2$\\\tiny{  } & \tiny{  }\end{tabular} &\begin{tabular}{c} = \\~ \end{tabular} & \begin{tabular}{ccc}$\beta_1$+ & $\beta_2$+ & $\beta_3$\\\tiny{ $A_3$ } & \tiny{ $A_3$ } & \tiny{ $A_3$ }\end{tabular} & $D_4$ & $E_7$ & $\langle\alpha_1, \alpha_2\rangle=1$,  $\langle\alpha_1, \beta_1\rangle=1$, $\langle\alpha_1, \beta_2\rangle=1$, $\langle\alpha_1, \beta_3\rangle=1$, $\langle\alpha_2, \beta_1\rangle=1$, $\langle\alpha_2, \beta_2\rangle=1$, $\langle\alpha_2, \beta_3\rangle=1$, 
\\
\begin{tabular}{cc}$\alpha_1$+ & $\alpha_2$\\\tiny{ $A_3$ } & \tiny{  }\end{tabular} &\begin{tabular}{c} = \\~ \end{tabular} & \begin{tabular}{cc}$\beta_1$+ & $\beta_2$\\\tiny{ $A_3$ } & \tiny{ $A_3$ }\end{tabular} & $A_3$ & $D_6$ &  $\langle\alpha_1, \beta_1\rangle=1$, $\langle\alpha_1, \beta_2\rangle=1$, $\langle\alpha_2, \beta_1\rangle=1$, $\langle\alpha_2, \beta_2\rangle=1$, 
\\
\begin{tabular}{c}2$\alpha_1$\\\tiny{  }\end{tabular} &\begin{tabular}{c} = \\~ \end{tabular} & \begin{tabular}{cccc}$\beta_1$+ & $\beta_2$+ & $\beta_3$+ & $\beta_4$\\\tiny{ $A_3$ } & \tiny{ $A_3$ } & \tiny{ $A_3$ } & \tiny{ $A_3$ }\end{tabular} & $D_4$ & $E_7$ &  $\langle\alpha_1, \beta_1\rangle=1$, $\langle\alpha_1, \beta_2\rangle=1$, $\langle\alpha_1, \beta_3\rangle=1$, $\langle\alpha_1, \beta_4\rangle=1$, 
\\
\begin{tabular}{cc}$\alpha_1$+ & $\alpha_2$\\\tiny{  } & \tiny{  }\end{tabular} &\begin{tabular}{c} = \\~ \end{tabular} & \begin{tabular}{cc}$\beta_1$+ & $\beta_2$\\\tiny{ $A_3$ } & \tiny{ $A_3$ }\end{tabular} & $A_3$ & $D_5$ &  $\langle\alpha_1, \beta_1\rangle=1$, $\langle\alpha_1, \beta_2\rangle=1$, $\langle\alpha_2, \beta_1\rangle=1$, $\langle\alpha_2, \beta_2\rangle=1$, 
\\
\begin{tabular}{cc}$\alpha_1$+ & $\alpha_2$\\\tiny{ $A_3$ } & \tiny{ $A_3$ }\end{tabular} &\begin{tabular}{c} = \\~ \end{tabular} & \begin{tabular}{cc}$\beta_1$+ & $\beta_2$\\\tiny{ $A_3$ } & \tiny{ $A_3$ }\end{tabular} & $A_3$ & $D_6$ &  $\langle\alpha_1, \beta_1\rangle=1$, $\langle\alpha_1, \beta_2\rangle=1$, $\langle\alpha_2, \beta_1\rangle=1$, $\langle\alpha_2, \beta_2\rangle=1$, 
\\
\begin{tabular}{cc}$\alpha_1$+ & $\alpha_2$\\\tiny{ $A_3$ } & \tiny{  }\end{tabular} &\begin{tabular}{c} = \\~ \end{tabular} & \begin{tabular}{c}$\beta_1$\\\tiny{ $A_3$ }\end{tabular} & $A_2$ & $D_5$ & $\langle\alpha_1, \alpha_2\rangle=-1$,  $\langle\alpha_1, \beta_1\rangle=1$, $\langle\alpha_2, \beta_1\rangle=1$, 
\\
\begin{tabular}{cc}$\alpha_1$+ & $\alpha_2$\\\tiny{ $A_3$ } & \tiny{ $A_3$ }\end{tabular} &\begin{tabular}{c} = \\~ \end{tabular} & \begin{tabular}{cc}$\beta_1$+ & $\beta_2$\\\tiny{ $A_3$ } & \tiny{ $A_3$ }\end{tabular} & $A_3$ & $E_6$ &  $\langle\alpha_1, \beta_1\rangle=1$, $\langle\alpha_1, \beta_2\rangle=1$, $\langle\alpha_2, \beta_1\rangle=1$, $\langle\alpha_2, \beta_2\rangle=1$, 
\\
\begin{tabular}{cc}$\alpha_1$+ & $\alpha_2$\\\tiny{ $A_3$ } & \tiny{ $A_3$ }\end{tabular} &\begin{tabular}{c} = \\~ \end{tabular} & \begin{tabular}{ccc}$\beta_1$+ & $\beta_2$+ & $\beta_3$\\\tiny{ $A_3$ } & \tiny{ $A_3$ } & \tiny{ $A_3$ }\end{tabular} & $D_4$ & $E_7$ & $\langle\alpha_1, \alpha_2\rangle=1$,  $\langle\alpha_1, \beta_1\rangle=1$, $\langle\alpha_1, \beta_2\rangle=1$, $\langle\alpha_1, \beta_3\rangle=1$, $\langle\alpha_2, \beta_1\rangle=1$, $\langle\alpha_2, \beta_2\rangle=1$, $\langle\alpha_2, \beta_3\rangle=1$, 
\\
\begin{tabular}{c}2$\alpha_1$\\\tiny{ $A_3$ }\end{tabular} &\begin{tabular}{c} = \\~ \end{tabular} & \begin{tabular}{cccc}$\beta_1$+ & $\beta_2$+ & $\beta_3$+ & $\beta_4$\\\tiny{ $A_3$ } & \tiny{ $A_3$ } & \tiny{ $A_3$ } & \tiny{ $A_3$ }\end{tabular} & $D_4$ & $E_7$ &  $\langle\alpha_1, \beta_1\rangle=1$, $\langle\alpha_1, \beta_2\rangle=1$, $\langle\alpha_1, \beta_3\rangle=1$, $\langle\alpha_1, \beta_4\rangle=1$, 
\\
\begin{tabular}{cc}$\alpha_1$+ & $\alpha_2$\\\tiny{ $A_3$ } & \tiny{  }\end{tabular} &\begin{tabular}{c} = \\~ \end{tabular} & \begin{tabular}{ccc}$\beta_1$+ & $\beta_2$+ & $\beta_3$\\\tiny{ $A_3$ } & \tiny{ $A_3$ } & \tiny{ $A_3$ }\end{tabular} & $D_4$ & $E_7$ & $\langle\alpha_1, \alpha_2\rangle=1$,  $\langle\alpha_1, \beta_1\rangle=1$, $\langle\alpha_1, \beta_2\rangle=1$, $\langle\alpha_1, \beta_3\rangle=1$, $\langle\alpha_2, \beta_1\rangle=1$, $\langle\alpha_2, \beta_2\rangle=1$, $\langle\alpha_2, \beta_3\rangle=1$, 
\\
\begin{tabular}{c}2$\alpha_1$\\\tiny{ $A_3$ }\end{tabular} &\begin{tabular}{c} = \\~ \end{tabular} & \begin{tabular}{cccc}$\beta_1$+ & $\beta_2$+ & $\beta_3$+ & $\beta_4$\\\tiny{ $A_3$ } & \tiny{ $A_3$ } & \tiny{ $A_3$ } & \tiny{  }\end{tabular} & $D_4$ & $E_7$ &  $\langle\alpha_1, \beta_1\rangle=1$, $\langle\alpha_1, \beta_2\rangle=1$, $\langle\alpha_1, \beta_3\rangle=1$, $\langle\alpha_1, \beta_4\rangle=1$, 
\\
\begin{tabular}{cc}$\alpha_1$+ & $\alpha_2$\\\tiny{ $A_3$ } & \tiny{ $A_3$ }\end{tabular} &\begin{tabular}{c} = \\~ \end{tabular} & \begin{tabular}{ccc}$\beta_1$+ & $\beta_2$+ & $\beta_3$\\\tiny{ $A_3$ } & \tiny{ $A_3$ } & \tiny{  }\end{tabular} & $D_4$ & $E_7$ & $\langle\alpha_1, \alpha_2\rangle=1$,  $\langle\alpha_1, \beta_1\rangle=1$, $\langle\alpha_1, \beta_2\rangle=1$, $\langle\alpha_1, \beta_3\rangle=1$, $\langle\alpha_2, \beta_1\rangle=1$, $\langle\alpha_2, \beta_2\rangle=1$, $\langle\alpha_2, \beta_3\rangle=1$, 
\\
\hline\multicolumn{5}{c}{$\mathfrak{k}$-semisimple type: $A_2$+$A_1$}\\
\hline\hline\begin{tabular}{cc}$\alpha_1$+ & $\alpha_2$\\\tiny{ $A_2$+$A_1$ } & \tiny{  }\end{tabular} &\begin{tabular}{c} = \\~ \end{tabular} & \begin{tabular}{c}$\beta_1$\\\tiny{ $A_2$+$A_1$ }\end{tabular} & $A_2$ & $D_5$ & $\langle\alpha_1, \alpha_2\rangle=-1$,  $\langle\alpha_1, \beta_1\rangle=1$, $\langle\alpha_2, \beta_1\rangle=1$, 
\\
\begin{tabular}{cc}$\alpha_1$+ & $\alpha_2$\\\tiny{ $A_2$ } & \tiny{  }\end{tabular} &\begin{tabular}{c} = \\~ \end{tabular} & \begin{tabular}{cc}$\beta_1$+ & $\beta_2$\\\tiny{ $A_2$+$A_1$ } & \tiny{ $A_2$+$A_1$ }\end{tabular} & $A_3$ & $D_5$ &  $\langle\alpha_1, \beta_1\rangle=1$, $\langle\alpha_1, \beta_2\rangle=1$, $\langle\alpha_2, \beta_1\rangle=1$, $\langle\alpha_2, \beta_2\rangle=1$, 
\\
\begin{tabular}{cc}$\alpha_1$+ & $\alpha_2$\\\tiny{ $A_2$+$A_1$ } & \tiny{ $A_2$ }\end{tabular} &\begin{tabular}{c} = \\~ \end{tabular} & \begin{tabular}{cc}$\beta_1$+ & $\beta_2$\\\tiny{ $A_2$+$A_1$ } & \tiny{ $A_2$ }\end{tabular} & $A_3$ & $E_6$ &  $\langle\alpha_1, \beta_1\rangle=1$, $\langle\alpha_1, \beta_2\rangle=1$, $\langle\alpha_2, \beta_1\rangle=1$, $\langle\alpha_2, \beta_2\rangle=1$, 
\\
\begin{tabular}{cc}$\alpha_1$+ & $\alpha_2$\\\tiny{ $A_1$ } & \tiny{  }\end{tabular} &\begin{tabular}{c} = \\~ \end{tabular} & \begin{tabular}{cc}$\beta_1$+ & $\beta_2$\\\tiny{ $A_2$+$A_1$ } & \tiny{ $A_2$ }\end{tabular} & $A_3$ & $E_6$ &  $\langle\alpha_1, \beta_1\rangle=1$, $\langle\alpha_1, \beta_2\rangle=1$, $\langle\alpha_2, \beta_1\rangle=1$, $\langle\alpha_2, \beta_2\rangle=1$, 
\\
\begin{tabular}{cc}$\alpha_1$+ & $\alpha_2$\\\tiny{ $A_2$ } & \tiny{  }\end{tabular} &\begin{tabular}{c} = \\~ \end{tabular} & \begin{tabular}{c}$\beta_1$\\\tiny{ $A_2$ }\end{tabular} & $A_2$ & $A_4$+$A_1$ & $\langle\alpha_1, \alpha_2\rangle=-1$,  $\langle\alpha_1, \beta_1\rangle=1$, $\langle\alpha_2, \beta_1\rangle=1$, 
\\
\begin{tabular}{cc}$\alpha_1$+ & $\alpha_2$\\\tiny{ $A_1$ } & \tiny{  }\end{tabular} &\begin{tabular}{c} = \\~ \end{tabular} & \begin{tabular}{c}$\beta_1$\\\tiny{ $A_1$ }\end{tabular} & $A_2$ & $A_3$+$A_2$ & $\langle\alpha_1, \alpha_2\rangle=-1$,  $\langle\alpha_1, \beta_1\rangle=1$, $\langle\alpha_2, \beta_1\rangle=1$, 
\\
\begin{tabular}{cc}$\alpha_1$+ & $\alpha_2$\\\tiny{ $A_2$ } & \tiny{ $A_2$ }\end{tabular} &\begin{tabular}{c} = \\~ \end{tabular} & \begin{tabular}{ccc}$\beta_1$+ & $\beta_2$+ & $\beta_3$\\\tiny{ $A_2$+$A_1$ } & \tiny{ $A_2$+$A_1$ } & \tiny{ $A_2$ }\end{tabular} & $D_4$ & $E_6$ & $\langle\alpha_1, \alpha_2\rangle=1$,  $\langle\alpha_1, \beta_1\rangle=1$, $\langle\alpha_1, \beta_2\rangle=1$, $\langle\alpha_1, \beta_3\rangle=1$, $\langle\alpha_2, \beta_1\rangle=1$, $\langle\alpha_2, \beta_2\rangle=1$, $\langle\alpha_2, \beta_3\rangle=1$, 
\\
\begin{tabular}{cc}$\alpha_1$+ & $\alpha_2$\\\tiny{ $A_2$ } & \tiny{ $A_2$ }\end{tabular} &\begin{tabular}{c} = \\~ \end{tabular} & \begin{tabular}{cc}$\beta_1$+ & $\beta_2$\\\tiny{ $A_2$ } & \tiny{ $A_2$ }\end{tabular} & $A_3$ & $D_5$+$A_1$ &  $\langle\alpha_1, \beta_1\rangle=1$, $\langle\alpha_1, \beta_2\rangle=1$, $\langle\alpha_2, \beta_1\rangle=1$, $\langle\alpha_2, \beta_2\rangle=1$, 
\\
\begin{tabular}{cc}$\alpha_1$+ & $\alpha_2$\\\tiny{ $A_1$ } & \tiny{ $A_2$ }\end{tabular} &\begin{tabular}{c} = \\~ \end{tabular} & \begin{tabular}{cc}$\beta_1$+ & $\beta_2$\\\tiny{ $A_1$ } & \tiny{ $A_2$ }\end{tabular} & $A_3$ & $A_6$ &  $\langle\alpha_1, \beta_1\rangle=1$, $\langle\alpha_1, \beta_2\rangle=1$, $\langle\alpha_2, \beta_1\rangle=1$, $\langle\alpha_2, \beta_2\rangle=1$, 
\\
\begin{tabular}{cc}$\alpha_1$+ & $\alpha_2$\\\tiny{ $A_2$ } & \tiny{  }\end{tabular} &\begin{tabular}{c} = \\~ \end{tabular} & \begin{tabular}{cc}$\beta_1$+ & $\beta_2$\\\tiny{ $A_2$+$A_1$ } & \tiny{ $A_1$ }\end{tabular} & $A_3$ & $D_6$ &  $\langle\alpha_1, \beta_1\rangle=1$, $\langle\alpha_1, \beta_2\rangle=1$, $\langle\alpha_2, \beta_1\rangle=1$, $\langle\alpha_2, \beta_2\rangle=1$, 
\\
\begin{tabular}{c}2$\alpha_1$\\\tiny{ $A_2$ }\end{tabular} &\begin{tabular}{c} = \\~ \end{tabular} & \begin{tabular}{cccc}$\beta_1$+ & $\beta_2$+ & $\beta_3$+ & $\beta_4$\\\tiny{ $A_2$+$A_1$ } & \tiny{ $A_1$ } & \tiny{ $A_2$ } & \tiny{ $A_2$ }\end{tabular} & $D_4$ & $E_7$ &  $\langle\alpha_1, \beta_1\rangle=1$, $\langle\alpha_1, \beta_2\rangle=1$, $\langle\alpha_1, \beta_3\rangle=1$, $\langle\alpha_1, \beta_4\rangle=1$, 
\\
\begin{tabular}{cc}$\alpha_1$+ & $\alpha_2$\\\tiny{ $A_2$ } & \tiny{ $A_2$ }\end{tabular} &\begin{tabular}{c} = \\~ \end{tabular} & \begin{tabular}{ccc}$\beta_1$+ & $\beta_2$+ & $\beta_3$\\\tiny{ $A_2$+$A_1$ } & \tiny{ $A_1$ } & \tiny{ $A_2$ }\end{tabular} & $D_4$ & $E_7$ & $\langle\alpha_1, \alpha_2\rangle=1$,  $\langle\alpha_1, \beta_1\rangle=1$, $\langle\alpha_1, \beta_2\rangle=1$, $\langle\alpha_1, \beta_3\rangle=1$, $\langle\alpha_2, \beta_1\rangle=1$, $\langle\alpha_2, \beta_2\rangle=1$, $\langle\alpha_2, \beta_3\rangle=1$, 
\\
\begin{tabular}{cc}$\alpha_1$+ & $\alpha_2$\\\tiny{ $A_2$+$A_1$ } & \tiny{ $A_1$ }\end{tabular} &\begin{tabular}{c} = \\~ \end{tabular} & \begin{tabular}{cc}$\beta_1$+ & $\beta_2$\\\tiny{ $A_2$+$A_1$ } & \tiny{ $A_1$ }\end{tabular} & $A_3$ & $D_6$ &  $\langle\alpha_1, \beta_1\rangle=1$, $\langle\alpha_1, \beta_2\rangle=1$, $\langle\alpha_2, \beta_1\rangle=1$, $\langle\alpha_2, \beta_2\rangle=1$, 
\\
\begin{tabular}{cc}$\alpha_1$+ & $\alpha_2$\\\tiny{ $A_1$ } & \tiny{ $A_2$ }\end{tabular} &\begin{tabular}{c} = \\~ \end{tabular} & \begin{tabular}{cc}$\beta_1$+ & $\beta_2$\\\tiny{ $A_2$+$A_1$ } & \tiny{ $A_2$ }\end{tabular} & $A_3$ & $D_6$ &  $\langle\alpha_1, \beta_1\rangle=1$, $\langle\alpha_1, \beta_2\rangle=1$, $\langle\alpha_2, \beta_1\rangle=1$, $\langle\alpha_2, \beta_2\rangle=1$, 
\\
\begin{tabular}{cc}$\alpha_1$+ & $\alpha_2$\\\tiny{ $A_2$ } & \tiny{  }\end{tabular} &\begin{tabular}{c} = \\~ \end{tabular} & \begin{tabular}{cc}$\beta_1$+ & $\beta_2$\\\tiny{ $A_2$ } & \tiny{ $A_2$ }\end{tabular} & $A_3$ & $D_5$+$A_1$ &  $\langle\alpha_1, \beta_1\rangle=1$, $\langle\alpha_1, \beta_2\rangle=1$, $\langle\alpha_2, \beta_1\rangle=1$, $\langle\alpha_2, \beta_2\rangle=1$, 
\\
\begin{tabular}{cc}$\alpha_1$+ & $\alpha_2$\\\tiny{ $A_2$+$A_1$ } & \tiny{ $A_2$ }\end{tabular} &\begin{tabular}{c} = \\~ \end{tabular} & \begin{tabular}{ccc}$\beta_1$+ & $\beta_2$+ & $\beta_3$\\\tiny{ $A_2$+$A_1$ } & \tiny{ $A_2$ } & \tiny{ $A_2$ }\end{tabular} & $D_4$ & $E_7$ & $\langle\alpha_1, \alpha_2\rangle=1$,  $\langle\alpha_1, \beta_1\rangle=1$, $\langle\alpha_1, \beta_2\rangle=1$, $\langle\alpha_1, \beta_3\rangle=1$, $\langle\alpha_2, \beta_1\rangle=1$, $\langle\alpha_2, \beta_2\rangle=1$, $\langle\alpha_2, \beta_3\rangle=1$, 
\\
\begin{tabular}{cc}$\alpha_1$+ & $\alpha_2$\\\tiny{ $A_1$ } & \tiny{ $A_1$ }\end{tabular} &\begin{tabular}{c} = \\~ \end{tabular} & \begin{tabular}{ccc}$\beta_1$+ & $\beta_2$+ & $\beta_3$\\\tiny{ $A_2$+$A_1$ } & \tiny{ $A_1$ } & \tiny{ $A_2$ }\end{tabular} & $D_4$ & $E_7$ & $\langle\alpha_1, \alpha_2\rangle=1$,  $\langle\alpha_1, \beta_1\rangle=1$, $\langle\alpha_1, \beta_2\rangle=1$, $\langle\alpha_1, \beta_3\rangle=1$, $\langle\alpha_2, \beta_1\rangle=1$, $\langle\alpha_2, \beta_2\rangle=1$, $\langle\alpha_2, \beta_3\rangle=1$, 
\\
\begin{tabular}{cc}$\alpha_1$+ & $\alpha_2$\\\tiny{  } & \tiny{  }\end{tabular} &\begin{tabular}{c} = \\~ \end{tabular} & \begin{tabular}{cc}$\beta_1$+ & $\beta_2$\\\tiny{ $A_2$ } & \tiny{ $A_2$ }\end{tabular} & $A_3$ & $D_5$+$A_1$ &  $\langle\alpha_1, \beta_1\rangle=1$, $\langle\alpha_1, \beta_2\rangle=1$, $\langle\alpha_2, \beta_1\rangle=1$, $\langle\alpha_2, \beta_2\rangle=1$, 
\\
\begin{tabular}{cc}$\alpha_1$+ & $\alpha_2$\\\tiny{ $A_2$+$A_1$ } & \tiny{  }\end{tabular} &\begin{tabular}{c} = \\~ \end{tabular} & \begin{tabular}{cc}$\beta_1$+ & $\beta_2$\\\tiny{ $A_1$ } & \tiny{ $A_2$ }\end{tabular} & $A_3$ & $A_6$ &  $\langle\alpha_1, \beta_1\rangle=1$, $\langle\alpha_1, \beta_2\rangle=1$, $\langle\alpha_2, \beta_1\rangle=1$, $\langle\alpha_2, \beta_2\rangle=1$, 
\\
\begin{tabular}{c}2$\alpha_1$\\\tiny{  }\end{tabular} &\begin{tabular}{c} = \\~ \end{tabular} & \begin{tabular}{cccc}$\beta_1$+ & $\beta_2$+ & $\beta_3$+ & $\beta_4$\\\tiny{ $A_2$+$A_1$ } & \tiny{ $A_1$ } & \tiny{ $A_2$ } & \tiny{ $A_2$ }\end{tabular} & $D_4$ & $E_7$ &  $\langle\alpha_1, \beta_1\rangle=1$, $\langle\alpha_1, \beta_2\rangle=1$, $\langle\alpha_1, \beta_3\rangle=1$, $\langle\alpha_1, \beta_4\rangle=1$, 
\\
\begin{tabular}{cc}$\alpha_1$+ & $\alpha_2$\\\tiny{ $A_1$ } & \tiny{ $A_2$ }\end{tabular} &\begin{tabular}{c} = \\~ \end{tabular} & \begin{tabular}{c}$\beta_1$\\\tiny{ $A_2$+$A_1$ }\end{tabular} & $A_2$ & $A_5$ & $\langle\alpha_1, \alpha_2\rangle=-1$,  $\langle\alpha_1, \beta_1\rangle=1$, $\langle\alpha_2, \beta_1\rangle=1$, 
\\
\begin{tabular}{cc}$\alpha_1$+ & $\alpha_2$\\\tiny{ $A_2$ } & \tiny{ $A_2$ }\end{tabular} &\begin{tabular}{c} = \\~ \end{tabular} & \begin{tabular}{cc}$\beta_1$+ & $\beta_2$\\\tiny{ $A_2$+$A_1$ } & \tiny{ $A_2$+$A_1$ }\end{tabular} & $A_3$ & $D_6$ &  $\langle\alpha_1, \beta_1\rangle=1$, $\langle\alpha_1, \beta_2\rangle=1$, $\langle\alpha_2, \beta_1\rangle=1$, $\langle\alpha_2, \beta_2\rangle=1$, 
\\
\begin{tabular}{cc}$\alpha_1$+ & $\alpha_2$\\\tiny{ $A_1$ } & \tiny{ $A_1$ }\end{tabular} &\begin{tabular}{c} = \\~ \end{tabular} & \begin{tabular}{ccc}$\beta_1$+ & $\beta_2$+ & $\beta_3$\\\tiny{ $A_2$+$A_1$ } & \tiny{ $A_2$+$A_1$ } & \tiny{ $A_2$ }\end{tabular} & $D_4$ & $E_7$ & $\langle\alpha_1, \alpha_2\rangle=1$,  $\langle\alpha_1, \beta_1\rangle=1$, $\langle\alpha_1, \beta_2\rangle=1$, $\langle\alpha_1, \beta_3\rangle=1$, $\langle\alpha_2, \beta_1\rangle=1$, $\langle\alpha_2, \beta_2\rangle=1$, $\langle\alpha_2, \beta_3\rangle=1$, 
\\
\begin{tabular}{cc}$\alpha_1$+ & $\alpha_2$\\\tiny{  } & \tiny{  }\end{tabular} &\begin{tabular}{c} = \\~ \end{tabular} & \begin{tabular}{ccc}$\beta_1$+ & $\beta_2$+ & $\beta_3$\\\tiny{ $A_2$+$A_1$ } & \tiny{ $A_2$+$A_1$ } & \tiny{ $A_2$ }\end{tabular} & $D_4$ & $E_6$ & $\langle\alpha_1, \alpha_2\rangle=1$,  $\langle\alpha_1, \beta_1\rangle=1$, $\langle\alpha_1, \beta_2\rangle=1$, $\langle\alpha_1, \beta_3\rangle=1$, $\langle\alpha_2, \beta_1\rangle=1$, $\langle\alpha_2, \beta_2\rangle=1$, $\langle\alpha_2, \beta_3\rangle=1$, 
\\
\begin{tabular}{c}2$\alpha_1$\\\tiny{ $A_1$ }\end{tabular} &\begin{tabular}{c} = \\~ \end{tabular} & \begin{tabular}{cccc}$\beta_1$+ & $\beta_2$+ & $\beta_3$+ & $\beta_4$\\\tiny{ $A_2$+$A_1$ } & \tiny{ $A_1$ } & \tiny{ $A_2$ } & \tiny{ $A_2$ }\end{tabular} & $D_4$ & $E_7$ &  $\langle\alpha_1, \beta_1\rangle=1$, $\langle\alpha_1, \beta_2\rangle=1$, $\langle\alpha_1, \beta_3\rangle=1$, $\langle\alpha_1, \beta_4\rangle=1$, 
\\
\begin{tabular}{cc}$\alpha_1$+ & $\alpha_2$\\\tiny{ $A_1$ } & \tiny{  }\end{tabular} &\begin{tabular}{c} = \\~ \end{tabular} & \begin{tabular}{ccc}$\beta_1$+ & $\beta_2$+ & $\beta_3$\\\tiny{ $A_2$+$A_1$ } & \tiny{ $A_2$ } & \tiny{ $A_2$ }\end{tabular} & $D_4$ & $E_7$ & $\langle\alpha_1, \alpha_2\rangle=1$,  $\langle\alpha_1, \beta_1\rangle=1$, $\langle\alpha_1, \beta_2\rangle=1$, $\langle\alpha_1, \beta_3\rangle=1$, $\langle\alpha_2, \beta_1\rangle=1$, $\langle\alpha_2, \beta_2\rangle=1$, $\langle\alpha_2, \beta_3\rangle=1$, 
\\
\begin{tabular}{cc}$\alpha_1$+ & $\alpha_2$\\\tiny{  } & \tiny{  }\end{tabular} &\begin{tabular}{c} = \\~ \end{tabular} & \begin{tabular}{ccc}$\beta_1$+ & $\beta_2$+ & $\beta_3$\\\tiny{ $A_2$+$A_1$ } & \tiny{ $A_1$ } & \tiny{ $A_2$ }\end{tabular} & $D_4$ & $E_7$ & $\langle\alpha_1, \alpha_2\rangle=1$,  $\langle\alpha_1, \beta_1\rangle=1$, $\langle\alpha_1, \beta_2\rangle=1$, $\langle\alpha_1, \beta_3\rangle=1$, $\langle\alpha_2, \beta_1\rangle=1$, $\langle\alpha_2, \beta_2\rangle=1$, $\langle\alpha_2, \beta_3\rangle=1$, 
\\
\begin{tabular}{cc}$\alpha_1$+ & $\alpha_2$\\\tiny{ $A_1$ } & \tiny{ $A_2$ }\end{tabular} &\begin{tabular}{c} = \\~ \end{tabular} & \begin{tabular}{ccc}$\beta_1$+ & $\beta_2$+ & $\beta_3$\\\tiny{ $A_1$ } & \tiny{ $A_2$ } & \tiny{ $A_2$ }\end{tabular} & $D_4$ & $E_7$ & $\langle\alpha_1, \alpha_2\rangle=1$,  $\langle\alpha_1, \beta_1\rangle=1$, $\langle\alpha_1, \beta_2\rangle=1$, $\langle\alpha_1, \beta_3\rangle=1$, $\langle\alpha_2, \beta_1\rangle=1$, $\langle\alpha_2, \beta_2\rangle=1$, $\langle\alpha_2, \beta_3\rangle=1$, 
\\
\begin{tabular}{cc}$\alpha_1$+ & $\alpha_2$\\\tiny{ $A_2$+$A_1$ } & \tiny{  }\end{tabular} &\begin{tabular}{c} = \\~ \end{tabular} & \begin{tabular}{ccc}$\beta_1$+ & $\beta_2$+ & $\beta_3$\\\tiny{ $A_1$ } & \tiny{ $A_2$ } & \tiny{ $A_2$ }\end{tabular} & $D_4$ & $E_7$ & $\langle\alpha_1, \alpha_2\rangle=1$,  $\langle\alpha_1, \beta_1\rangle=1$, $\langle\alpha_1, \beta_2\rangle=1$, $\langle\alpha_1, \beta_3\rangle=1$, $\langle\alpha_2, \beta_1\rangle=1$, $\langle\alpha_2, \beta_2\rangle=1$, $\langle\alpha_2, \beta_3\rangle=1$, 
\\
\begin{tabular}{cc}$\alpha_1$+ & $\alpha_2$\\\tiny{ $A_2$+$A_1$ } & \tiny{ $A_1$ }\end{tabular} &\begin{tabular}{c} = \\~ \end{tabular} & \begin{tabular}{cc}$\beta_1$+ & $\beta_2$\\\tiny{ $A_2$+$A_1$ } & \tiny{ $A_2$+$A_1$ }\end{tabular} & $A_3$ & $E_6$ &  $\langle\alpha_1, \beta_1\rangle=1$, $\langle\alpha_1, \beta_2\rangle=1$, $\langle\alpha_2, \beta_1\rangle=1$, $\langle\alpha_2, \beta_2\rangle=1$, 
\\
\begin{tabular}{cc}$\alpha_1$+ & $\alpha_2$\\\tiny{ $A_2$+$A_1$ } & \tiny{ $A_2$ }\end{tabular} &\begin{tabular}{c} = \\~ \end{tabular} & \begin{tabular}{cc}$\beta_1$+ & $\beta_2$\\\tiny{ $A_2$+$A_1$ } & \tiny{ $A_2$ }\end{tabular} & $A_3$ & $D_6$ &  $\langle\alpha_1, \beta_1\rangle=1$, $\langle\alpha_1, \beta_2\rangle=1$, $\langle\alpha_2, \beta_1\rangle=1$, $\langle\alpha_2, \beta_2\rangle=1$, 
\\
\begin{tabular}{cc}$\alpha_1$+ & $\alpha_2$\\\tiny{ $A_2$+$A_1$ } & \tiny{  }\end{tabular} &\begin{tabular}{c} = \\~ \end{tabular} & \begin{tabular}{cc}$\beta_1$+ & $\beta_2$\\\tiny{ $A_2$+$A_1$ } & \tiny{ $A_2$ }\end{tabular} & $A_3$ & $E_6$ &  $\langle\alpha_1, \beta_1\rangle=1$, $\langle\alpha_1, \beta_2\rangle=1$, $\langle\alpha_2, \beta_1\rangle=1$, $\langle\alpha_2, \beta_2\rangle=1$, 
\\
\begin{tabular}{cc}$\alpha_1$+ & $\alpha_2$\\\tiny{ $A_2$ } & \tiny{ $A_2$ }\end{tabular} &\begin{tabular}{c} = \\~ \end{tabular} & \begin{tabular}{cc}$\beta_1$+ & $\beta_2$\\\tiny{ $A_2$+$A_1$ } & \tiny{ $A_2$+$A_1$ }\end{tabular} & $A_3$ & $E_6$ &  $\langle\alpha_1, \beta_1\rangle=1$, $\langle\alpha_1, \beta_2\rangle=1$, $\langle\alpha_2, \beta_1\rangle=1$, $\langle\alpha_2, \beta_2\rangle=1$, 
\\
\begin{tabular}{cc}$\alpha_1$+ & $\alpha_2$\\\tiny{ $A_2$ } & \tiny{  }\end{tabular} &\begin{tabular}{c} = \\~ \end{tabular} & \begin{tabular}{ccc}$\beta_1$+ & $\beta_2$+ & $\beta_3$\\\tiny{ $A_2$+$A_1$ } & \tiny{ $A_1$ } & \tiny{ $A_2$ }\end{tabular} & $D_4$ & $E_7$ & $\langle\alpha_1, \alpha_2\rangle=1$,  $\langle\alpha_1, \beta_1\rangle=1$, $\langle\alpha_1, \beta_2\rangle=1$, $\langle\alpha_1, \beta_3\rangle=1$, $\langle\alpha_2, \beta_1\rangle=1$, $\langle\alpha_2, \beta_2\rangle=1$, $\langle\alpha_2, \beta_3\rangle=1$, 
\\
\begin{tabular}{cc}$\alpha_1$+ & $\alpha_2$\\\tiny{ $A_2$+$A_1$ } & \tiny{ $A_2$+$A_1$ }\end{tabular} &\begin{tabular}{c} = \\~ \end{tabular} & \begin{tabular}{ccc}$\beta_1$+ & $\beta_2$+ & $\beta_3$\\\tiny{ $A_2$+$A_1$ } & \tiny{ $A_1$ } & \tiny{ $A_2$ }\end{tabular} & $D_4$ & $E_7$ & $\langle\alpha_1, \alpha_2\rangle=1$,  $\langle\alpha_1, \beta_1\rangle=1$, $\langle\alpha_1, \beta_2\rangle=1$, $\langle\alpha_1, \beta_3\rangle=1$, $\langle\alpha_2, \beta_1\rangle=1$, $\langle\alpha_2, \beta_2\rangle=1$, $\langle\alpha_2, \beta_3\rangle=1$, 
\\
\begin{tabular}{c}2$\alpha_1$\\\tiny{ $A_2$+$A_1$ }\end{tabular} &\begin{tabular}{c} = \\~ \end{tabular} & \begin{tabular}{cccc}$\beta_1$+ & $\beta_2$+ & $\beta_3$+ & $\beta_4$\\\tiny{ $A_2$+$A_1$ } & \tiny{ $A_1$ } & \tiny{ $A_2$ } & \tiny{ $A_2$ }\end{tabular} & $D_4$ & $E_7$ &  $\langle\alpha_1, \beta_1\rangle=1$, $\langle\alpha_1, \beta_2\rangle=1$, $\langle\alpha_1, \beta_3\rangle=1$, $\langle\alpha_1, \beta_4\rangle=1$, 
\\
\begin{tabular}{cc}$\alpha_1$+ & $\alpha_2$\\\tiny{ $A_2$+$A_1$ } & \tiny{ $A_1$ }\end{tabular} &\begin{tabular}{c} = \\~ \end{tabular} & \begin{tabular}{ccc}$\beta_1$+ & $\beta_2$+ & $\beta_3$\\\tiny{ $A_2$+$A_1$ } & \tiny{ $A_1$ } & \tiny{ $A_2$ }\end{tabular} & $D_4$ & $E_7$ & $\langle\alpha_1, \alpha_2\rangle=1$,  $\langle\alpha_1, \beta_1\rangle=1$, $\langle\alpha_1, \beta_2\rangle=1$, $\langle\alpha_1, \beta_3\rangle=1$, $\langle\alpha_2, \beta_1\rangle=1$, $\langle\alpha_2, \beta_2\rangle=1$, $\langle\alpha_2, \beta_3\rangle=1$, 
\\
\begin{tabular}{cc}$\alpha_1$+ & $\alpha_2$\\\tiny{ $A_2$ } & \tiny{ $A_2$ }\end{tabular} &\begin{tabular}{c} = \\~ \end{tabular} & \begin{tabular}{ccc}$\beta_1$+ & $\beta_2$+ & $\beta_3$\\\tiny{ $A_2$+$A_1$ } & \tiny{ $A_2$+$A_1$ } & \tiny{ $A_2$ }\end{tabular} & $D_4$ & $E_7$ & $\langle\alpha_1, \alpha_2\rangle=1$,  $\langle\alpha_1, \beta_1\rangle=1$, $\langle\alpha_1, \beta_2\rangle=1$, $\langle\alpha_1, \beta_3\rangle=1$, $\langle\alpha_2, \beta_1\rangle=1$, $\langle\alpha_2, \beta_2\rangle=1$, $\langle\alpha_2, \beta_3\rangle=1$, 
\\
\begin{tabular}{cc}$\alpha_1$+ & $\alpha_2$\\\tiny{ $A_1$ } & \tiny{ $A_1$ }\end{tabular} &\begin{tabular}{c} = \\~ \end{tabular} & \begin{tabular}{cc}$\beta_1$+ & $\beta_2$\\\tiny{ $A_2$+$A_1$ } & \tiny{ $A_2$+$A_1$ }\end{tabular} & $A_3$ & $D_6$ &  $\langle\alpha_1, \beta_1\rangle=1$, $\langle\alpha_1, \beta_2\rangle=1$, $\langle\alpha_2, \beta_1\rangle=1$, $\langle\alpha_2, \beta_2\rangle=1$, 
\\
\begin{tabular}{cc}$\alpha_1$+ & $\alpha_2$\\\tiny{ $A_1$ } & \tiny{ $A_2$ }\end{tabular} &\begin{tabular}{c} = \\~ \end{tabular} & \begin{tabular}{cc}$\beta_1$+ & $\beta_2$\\\tiny{ $A_2$+$A_1$ } & \tiny{  }\end{tabular} & $A_3$ & $A_6$ &  $\langle\alpha_1, \beta_1\rangle=1$, $\langle\alpha_1, \beta_2\rangle=1$, $\langle\alpha_2, \beta_1\rangle=1$, $\langle\alpha_2, \beta_2\rangle=1$, 
\\
\begin{tabular}{c}2$\alpha_1$\\\tiny{ $A_1$ }\end{tabular} &\begin{tabular}{c} = \\~ \end{tabular} & \begin{tabular}{cccc}$\beta_1$+ & $\beta_2$+ & $\beta_3$+ & $\beta_4$\\\tiny{ $A_2$+$A_1$ } & \tiny{ $A_2$+$A_1$ } & \tiny{ $A_2$ } & \tiny{  }\end{tabular} & $D_4$ & $E_7$ &  $\langle\alpha_1, \beta_1\rangle=1$, $\langle\alpha_1, \beta_2\rangle=1$, $\langle\alpha_1, \beta_3\rangle=1$, $\langle\alpha_1, \beta_4\rangle=1$, 
\\
\begin{tabular}{cc}$\alpha_1$+ & $\alpha_2$\\\tiny{ $A_1$ } & \tiny{ $A_2$ }\end{tabular} &\begin{tabular}{c} = \\~ \end{tabular} & \begin{tabular}{ccc}$\beta_1$+ & $\beta_2$+ & $\beta_3$\\\tiny{ $A_2$+$A_1$ } & \tiny{ $A_2$ } & \tiny{  }\end{tabular} & $D_4$ & $E_7$ & $\langle\alpha_1, \alpha_2\rangle=1$,  $\langle\alpha_1, \beta_1\rangle=1$, $\langle\alpha_1, \beta_2\rangle=1$, $\langle\alpha_1, \beta_3\rangle=1$, $\langle\alpha_2, \beta_1\rangle=1$, $\langle\alpha_2, \beta_2\rangle=1$, $\langle\alpha_2, \beta_3\rangle=1$, 
\\
\begin{tabular}{c}2$\alpha_1$\\\tiny{ $A_2$+$A_1$ }\end{tabular} &\begin{tabular}{c} = \\~ \end{tabular} & \begin{tabular}{cccc}$\beta_1$+ & $\beta_2$+ & $\beta_3$+ & $\beta_4$\\\tiny{ $A_2$+$A_1$ } & \tiny{ $A_2$+$A_1$ } & \tiny{ $A_2$ } & \tiny{  }\end{tabular} & $D_4$ & $E_7$ &  $\langle\alpha_1, \beta_1\rangle=1$, $\langle\alpha_1, \beta_2\rangle=1$, $\langle\alpha_1, \beta_3\rangle=1$, $\langle\alpha_1, \beta_4\rangle=1$, 
\\
\begin{tabular}{cc}$\alpha_1$+ & $\alpha_2$\\\tiny{ $A_2$+$A_1$ } & \tiny{ $A_1$ }\end{tabular} &\begin{tabular}{c} = \\~ \end{tabular} & \begin{tabular}{ccc}$\beta_1$+ & $\beta_2$+ & $\beta_3$\\\tiny{ $A_2$+$A_1$ } & \tiny{ $A_2$+$A_1$ } & \tiny{  }\end{tabular} & $D_4$ & $E_7$ & $\langle\alpha_1, \alpha_2\rangle=1$,  $\langle\alpha_1, \beta_1\rangle=1$, $\langle\alpha_1, \beta_2\rangle=1$, $\langle\alpha_1, \beta_3\rangle=1$, $\langle\alpha_2, \beta_1\rangle=1$, $\langle\alpha_2, \beta_2\rangle=1$, $\langle\alpha_2, \beta_3\rangle=1$, 
\\
\begin{tabular}{cc}$\alpha_1$+ & $\alpha_2$\\\tiny{ $A_2$+$A_1$ } & \tiny{ $A_2$ }\end{tabular} &\begin{tabular}{c} = \\~ \end{tabular} & \begin{tabular}{ccc}$\beta_1$+ & $\beta_2$+ & $\beta_3$\\\tiny{ $A_2$+$A_1$ } & \tiny{ $A_2$ } & \tiny{  }\end{tabular} & $D_4$ & $E_7$ & $\langle\alpha_1, \alpha_2\rangle=1$,  $\langle\alpha_1, \beta_1\rangle=1$, $\langle\alpha_1, \beta_2\rangle=1$, $\langle\alpha_1, \beta_3\rangle=1$, $\langle\alpha_2, \beta_1\rangle=1$, $\langle\alpha_2, \beta_2\rangle=1$, $\langle\alpha_2, \beta_3\rangle=1$, 
\\
\begin{tabular}{cc}$\alpha_1$+ & $\alpha_2$\\\tiny{ $A_2$+$A_1$ } & \tiny{ $A_2$ }\end{tabular} &\begin{tabular}{c} = \\~ \end{tabular} & \begin{tabular}{ccc}$\beta_1$+ & $\beta_2$+ & $\beta_3$\\\tiny{ $A_1$ } & \tiny{ $A_2$ } & \tiny{ $A_2$ }\end{tabular} & $D_4$ & $E_7$ & $\langle\alpha_1, \alpha_2\rangle=1$,  $\langle\alpha_1, \beta_1\rangle=1$, $\langle\alpha_1, \beta_2\rangle=1$, $\langle\alpha_1, \beta_3\rangle=1$, $\langle\alpha_2, \beta_1\rangle=1$, $\langle\alpha_2, \beta_2\rangle=1$, $\langle\alpha_2, \beta_3\rangle=1$, 
\\
\begin{tabular}{c}2$\alpha_1$\\\tiny{ $A_2$ }\end{tabular} &\begin{tabular}{c} = \\~ \end{tabular} & \begin{tabular}{cccc}$\beta_1$+ & $\beta_2$+ & $\beta_3$+ & $\beta_4$\\\tiny{ $A_2$+$A_1$ } & \tiny{ $A_2$+$A_1$ } & \tiny{ $A_2$ } & \tiny{  }\end{tabular} & $D_4$ & $E_7$ &  $\langle\alpha_1, \beta_1\rangle=1$, $\langle\alpha_1, \beta_2\rangle=1$, $\langle\alpha_1, \beta_3\rangle=1$, $\langle\alpha_1, \beta_4\rangle=1$, 
\\
\begin{tabular}{cc}$\alpha_1$+ & $\alpha_2$\\\tiny{ $A_2$ } & \tiny{ $A_2$ }\end{tabular} &\begin{tabular}{c} = \\~ \end{tabular} & \begin{tabular}{ccc}$\beta_1$+ & $\beta_2$+ & $\beta_3$\\\tiny{ $A_2$+$A_1$ } & \tiny{ $A_2$+$A_1$ } & \tiny{  }\end{tabular} & $D_4$ & $E_7$ & $\langle\alpha_1, \alpha_2\rangle=1$,  $\langle\alpha_1, \beta_1\rangle=1$, $\langle\alpha_1, \beta_2\rangle=1$, $\langle\alpha_1, \beta_3\rangle=1$, $\langle\alpha_2, \beta_1\rangle=1$, $\langle\alpha_2, \beta_2\rangle=1$, $\langle\alpha_2, \beta_3\rangle=1$, 
\\
\begin{tabular}{cc}$\alpha_1$+ & $\alpha_2$\\\tiny{ $A_2$+$A_1$ } & \tiny{ $A_2$+$A_1$ }\end{tabular} &\begin{tabular}{c} = \\~ \end{tabular} & \begin{tabular}{ccc}$\beta_1$+ & $\beta_2$+ & $\beta_3$\\\tiny{ $A_2$+$A_1$ } & \tiny{ $A_2$+$A_1$ } & \tiny{ $A_2$ }\end{tabular} & $D_4$ & $E_7$ & $\langle\alpha_1, \alpha_2\rangle=1$,  $\langle\alpha_1, \beta_1\rangle=1$, $\langle\alpha_1, \beta_2\rangle=1$, $\langle\alpha_1, \beta_3\rangle=1$, $\langle\alpha_2, \beta_1\rangle=1$, $\langle\alpha_2, \beta_2\rangle=1$, $\langle\alpha_2, \beta_3\rangle=1$, 
\\
\hline\multicolumn{5}{c}{$\mathfrak{k}$-semisimple type: $A_1$+$A_1$+$A_1$}\\
\hline\hline\begin{tabular}{c}2$\alpha_1$\\\tiny{  }\end{tabular} &\begin{tabular}{c} = \\~ \end{tabular} & \begin{tabular}{cccc}$\beta_1$+ & $\beta_2$+ & $\beta_3$+ & $\beta_4$\\\tiny{ $A_1$+$A_1$' } & \tiny{ $A_1$''+$A_1$' } & \tiny{ $A_1$ } & \tiny{ $A_1$'' }\end{tabular} & $D_4$ & $E_7$ &  $\langle\alpha_1, \beta_1\rangle=1$, $\langle\alpha_1, \beta_2\rangle=1$, $\langle\alpha_1, \beta_3\rangle=1$, $\langle\alpha_1, \beta_4\rangle=1$, 
\\
\begin{tabular}{c}2$\alpha_1$\\\tiny{ $A_1$ }\end{tabular} &\begin{tabular}{c} = \\~ \end{tabular} & \begin{tabular}{cccc}$\beta_1$+ & $\beta_2$+ & $\beta_3$+ & $\beta_4$\\\tiny{ $A_1$+$A_1$' } & \tiny{ $A_1$''+$A_1$ } & \tiny{ $A_1$' } & \tiny{ $A_1$'' }\end{tabular} & $D_4$ & $E_7$ &  $\langle\alpha_1, \beta_1\rangle=1$, $\langle\alpha_1, \beta_2\rangle=1$, $\langle\alpha_1, \beta_3\rangle=1$, $\langle\alpha_1, \beta_4\rangle=1$, 
\\
\begin{tabular}{cc}$\alpha_1$+ & $\alpha_2$\\\tiny{ $A_1$ } & \tiny{  }\end{tabular} &\begin{tabular}{c} = \\~ \end{tabular} & \begin{tabular}{cc}$\beta_1$+ & $\beta_2$\\\tiny{ $A_1$+$A_1$' } & \tiny{ $A_1$ }\end{tabular} & $A_3$ & $D_5$+$A_1$ &  $\langle\alpha_1, \beta_1\rangle=1$, $\langle\alpha_1, \beta_2\rangle=1$, $\langle\alpha_2, \beta_1\rangle=1$, $\langle\alpha_2, \beta_2\rangle=1$, 
\\
\begin{tabular}{cc}$\alpha_1$+ & $\alpha_2$\\\tiny{ $A_1$ } & \tiny{ $A_1$' }\end{tabular} &\begin{tabular}{c} = \\~ \end{tabular} & \begin{tabular}{c}$\beta_1$\\\tiny{ $A_1$+$A_1$' }\end{tabular} & $A_2$ & $A_4$+$A_1$ & $\langle\alpha_1, \alpha_2\rangle=-1$,  $\langle\alpha_1, \beta_1\rangle=1$, $\langle\alpha_2, \beta_1\rangle=1$, 
\\
\begin{tabular}{cc}$\alpha_1$+ & $\alpha_2$\\\tiny{ $A_1$ } & \tiny{  }\end{tabular} &\begin{tabular}{c} = \\~ \end{tabular} & \begin{tabular}{ccc}$\beta_1$+ & $\beta_2$+ & $\beta_3$\\\tiny{ $A_1$+$A_1$' } & \tiny{ $A_1$+$A_1$' } & \tiny{ $A_1$'' }\end{tabular} & $D_4$ & $D_6$ & $\langle\alpha_1, \alpha_2\rangle=1$,  $\langle\alpha_1, \beta_1\rangle=1$, $\langle\alpha_1, \beta_2\rangle=1$, $\langle\alpha_1, \beta_3\rangle=1$, $\langle\alpha_2, \beta_1\rangle=1$, $\langle\alpha_2, \beta_2\rangle=1$, $\langle\alpha_2, \beta_3\rangle=1$, 
\\
\begin{tabular}{cc}$\alpha_1$+ & $\alpha_2$\\\tiny{ $A_1$+$A_1$'+$A_1$'' } & \tiny{  }\end{tabular} &\begin{tabular}{c} = \\~ \end{tabular} & \begin{tabular}{cc}$\beta_1$+ & $\beta_2$\\\tiny{ $A_1$+$A_1$' } & \tiny{ $A_1$'' }\end{tabular} & $A_3$ & $D_6$ &  $\langle\alpha_1, \beta_1\rangle=1$, $\langle\alpha_1, \beta_2\rangle=1$, $\langle\alpha_2, \beta_1\rangle=1$, $\langle\alpha_2, \beta_2\rangle=1$, 
\\
\begin{tabular}{cc}$\alpha_1$+ & $\alpha_2$\\\tiny{  } & \tiny{  }\end{tabular} &\begin{tabular}{c} = \\~ \end{tabular} & \begin{tabular}{cc}$\beta_1$+ & $\beta_2$\\\tiny{ $A_1$ } & \tiny{ $A_1$ }\end{tabular} & $A_3$ & $D_4$+2$A_1$ &  $\langle\alpha_1, \beta_1\rangle=1$, $\langle\alpha_1, \beta_2\rangle=1$, $\langle\alpha_2, \beta_1\rangle=1$, $\langle\alpha_2, \beta_2\rangle=1$, 
\\
\begin{tabular}{cc}$\alpha_1$+ & $\alpha_2$\\\tiny{ $A_1$ } & \tiny{ $A_1$' }\end{tabular} &\begin{tabular}{c} = \\~ \end{tabular} & \begin{tabular}{cc}$\beta_1$+ & $\beta_2$\\\tiny{ $A_1$ } & \tiny{ $A_1$' }\end{tabular} & $A_3$ & $A_5$+$A_1$ &  $\langle\alpha_1, \beta_1\rangle=1$, $\langle\alpha_1, \beta_2\rangle=1$, $\langle\alpha_2, \beta_1\rangle=1$, $\langle\alpha_2, \beta_2\rangle=1$, 
\\
\begin{tabular}{cc}$\alpha_1$+ & $\alpha_2$\\\tiny{ $A_1$ } & \tiny{  }\end{tabular} &\begin{tabular}{c} = \\~ \end{tabular} & \begin{tabular}{c}$\beta_1$\\\tiny{ $A_1$ }\end{tabular} & $A_2$ & $A_3$+2$A_1$ & $\langle\alpha_1, \alpha_2\rangle=-1$,  $\langle\alpha_1, \beta_1\rangle=1$, $\langle\alpha_2, \beta_1\rangle=1$, 
\\
\begin{tabular}{c}2$\alpha_1$\\\tiny{  }\end{tabular} &\begin{tabular}{c} = \\~ \end{tabular} & \begin{tabular}{cccc}$\beta_1$+ & $\beta_2$+ & $\beta_3$+ & $\beta_4$\\\tiny{ $A_1$+$A_1$' } & \tiny{ $A_1$+$A_1$' } & \tiny{ $A_1$'' } & \tiny{ $A_1$'' }\end{tabular} & $D_4$ & $D_6$ &  $\langle\alpha_1, \beta_1\rangle=1$, $\langle\alpha_1, \beta_2\rangle=1$, $\langle\alpha_1, \beta_3\rangle=1$, $\langle\alpha_1, \beta_4\rangle=1$, 
\\
\begin{tabular}{cc}$\alpha_1$+ & $\alpha_2$\\\tiny{ $A_1$ } & \tiny{ $A_1$' }\end{tabular} &\begin{tabular}{c} = \\~ \end{tabular} & \begin{tabular}{cc}$\beta_1$+ & $\beta_2$\\\tiny{ $A_1$+$A_1$' } & \tiny{ $A_1$''+$A_1$ }\end{tabular} & $A_3$ & $E_6$ &  $\langle\alpha_1, \beta_1\rangle=1$, $\langle\alpha_1, \beta_2\rangle=1$, $\langle\alpha_2, \beta_1\rangle=1$, $\langle\alpha_2, \beta_2\rangle=1$, 
\\
\begin{tabular}{cc}$\alpha_1$+ & $\alpha_2$\\\tiny{ $A_1$+$A_1$' } & \tiny{  }\end{tabular} &\begin{tabular}{c} = \\~ \end{tabular} & \begin{tabular}{cc}$\beta_1$+ & $\beta_2$\\\tiny{ $A_1$ } & \tiny{ $A_1$' }\end{tabular} & $A_3$ & $A_5$+$A_1$ &  $\langle\alpha_1, \beta_1\rangle=1$, $\langle\alpha_1, \beta_2\rangle=1$, $\langle\alpha_2, \beta_1\rangle=1$, $\langle\alpha_2, \beta_2\rangle=1$, 
\\
\begin{tabular}{cc}$\alpha_1$+ & $\alpha_2$\\\tiny{  } & \tiny{  }\end{tabular} &\begin{tabular}{c} = \\~ \end{tabular} & \begin{tabular}{cc}$\beta_1$+ & $\beta_2$\\\tiny{ $A_1$+$A_1$' } & \tiny{ $A_1$+$A_1$' }\end{tabular} & $A_3$ & $D_4$+$A_1$ &  $\langle\alpha_1, \beta_1\rangle=1$, $\langle\alpha_1, \beta_2\rangle=1$, $\langle\alpha_2, \beta_1\rangle=1$, $\langle\alpha_2, \beta_2\rangle=1$, 
\\
\begin{tabular}{cc}$\alpha_1$+ & $\alpha_2$\\\tiny{ $A_1$+$A_1$' } & \tiny{  }\end{tabular} &\begin{tabular}{c} = \\~ \end{tabular} & \begin{tabular}{cc}$\beta_1$+ & $\beta_2$\\\tiny{ $A_1$+$A_1$' } & \tiny{ $A_1$+$A_1$'' }\end{tabular} & $A_3$ & $D_6$ &  $\langle\alpha_1, \beta_1\rangle=1$, $\langle\alpha_1, \beta_2\rangle=1$, $\langle\alpha_2, \beta_1\rangle=1$, $\langle\alpha_2, \beta_2\rangle=1$, 
\\
\begin{tabular}{cc}$\alpha_1$+ & $\alpha_2$\\\tiny{ $A_1$ } & \tiny{  }\end{tabular} &\begin{tabular}{c} = \\~ \end{tabular} & \begin{tabular}{cc}$\beta_1$+ & $\beta_2$\\\tiny{ $A_1$+$A_1$' } & \tiny{ $A_1$' }\end{tabular} & $A_3$ & $D_5$+$A_1$ &  $\langle\alpha_1, \beta_1\rangle=1$, $\langle\alpha_1, \beta_2\rangle=1$, $\langle\alpha_2, \beta_1\rangle=1$, $\langle\alpha_2, \beta_2\rangle=1$, 
\\
\begin{tabular}{cc}$\alpha_1$+ & $\alpha_2$\\\tiny{ $A_1$+$A_1$' } & \tiny{  }\end{tabular} &\begin{tabular}{c} = \\~ \end{tabular} & \begin{tabular}{ccc}$\beta_1$+ & $\beta_2$+ & $\beta_3$\\\tiny{ $A_1$+$A_1$' } & \tiny{ $A_1$ } & \tiny{ $A_1$'' }\end{tabular} & $D_4$ & $E_7$ & $\langle\alpha_1, \alpha_2\rangle=1$,  $\langle\alpha_1, \beta_1\rangle=1$, $\langle\alpha_1, \beta_2\rangle=1$, $\langle\alpha_1, \beta_3\rangle=1$, $\langle\alpha_2, \beta_1\rangle=1$, $\langle\alpha_2, \beta_2\rangle=1$, $\langle\alpha_2, \beta_3\rangle=1$, 
\\
\begin{tabular}{cc}$\alpha_1$+ & $\alpha_2$\\\tiny{ $A_1$+$A_1$' } & \tiny{ $A_1$ }\end{tabular} &\begin{tabular}{c} = \\~ \end{tabular} & \begin{tabular}{ccc}$\beta_1$+ & $\beta_2$+ & $\beta_3$\\\tiny{ $A_1$+$A_1$' } & \tiny{ $A_1$''+$A_1$ } & \tiny{ $A_1$'' }\end{tabular} & $D_4$ & $E_7$ & $\langle\alpha_1, \alpha_2\rangle=1$,  $\langle\alpha_1, \beta_1\rangle=1$, $\langle\alpha_1, \beta_2\rangle=1$, $\langle\alpha_1, \beta_3\rangle=1$, $\langle\alpha_2, \beta_1\rangle=1$, $\langle\alpha_2, \beta_2\rangle=1$, $\langle\alpha_2, \beta_3\rangle=1$, 
\\
\begin{tabular}{c}2$\alpha_1$\\\tiny{ $A_1$+$A_1$' }\end{tabular} &\begin{tabular}{c} = \\~ \end{tabular} & \begin{tabular}{cccc}$\beta_1$+ & $\beta_2$+ & $\beta_3$+ & $\beta_4$\\\tiny{ $A_1$+$A_1$' } & \tiny{ $A_1$''+$A_1$ } & \tiny{ $A_1$' } & \tiny{ $A_1$'' }\end{tabular} & $D_4$ & $E_7$ &  $\langle\alpha_1, \beta_1\rangle=1$, $\langle\alpha_1, \beta_2\rangle=1$, $\langle\alpha_1, \beta_3\rangle=1$, $\langle\alpha_1, \beta_4\rangle=1$, 
\\
\begin{tabular}{c}2$\alpha_1$\\\tiny{  }\end{tabular} &\begin{tabular}{c} = \\~ \end{tabular} & \begin{tabular}{cccc}$\beta_1$+ & $\beta_2$+ & $\beta_3$+ & $\beta_4$\\\tiny{ $A_1$+$A_1$' } & \tiny{ $A_1$+$A_1$'' } & \tiny{ $A_1$' } & \tiny{ $A_1$'' }\end{tabular} & $D_4$ & $E_7$ &  $\langle\alpha_1, \beta_1\rangle=1$, $\langle\alpha_1, \beta_2\rangle=1$, $\langle\alpha_1, \beta_3\rangle=1$, $\langle\alpha_1, \beta_4\rangle=1$, 
\\
\begin{tabular}{cc}$\alpha_1$+ & $\alpha_2$\\\tiny{ $A_1$+$A_1$' } & \tiny{  }\end{tabular} &\begin{tabular}{c} = \\~ \end{tabular} & \begin{tabular}{ccc}$\beta_1$+ & $\beta_2$+ & $\beta_3$\\\tiny{ $A_1$+$A_1$' } & \tiny{ $A_1$'' } & \tiny{ $A_1$' }\end{tabular} & $D_4$ & $E_7$ & $\langle\alpha_1, \alpha_2\rangle=1$,  $\langle\alpha_1, \beta_1\rangle=1$, $\langle\alpha_1, \beta_2\rangle=1$, $\langle\alpha_1, \beta_3\rangle=1$, $\langle\alpha_2, \beta_1\rangle=1$, $\langle\alpha_2, \beta_2\rangle=1$, $\langle\alpha_2, \beta_3\rangle=1$, 
\\
\begin{tabular}{cc}$\alpha_1$+ & $\alpha_2$\\\tiny{ $A_1$+$A_1$' } & \tiny{  }\end{tabular} &\begin{tabular}{c} = \\~ \end{tabular} & \begin{tabular}{cc}$\beta_1$+ & $\beta_2$\\\tiny{ $A_1$+$A_1$' } & \tiny{ $A_1$''+$A_1$' }\end{tabular} & $A_3$ & $D_6$ &  $\langle\alpha_1, \beta_1\rangle=1$, $\langle\alpha_1, \beta_2\rangle=1$, $\langle\alpha_2, \beta_1\rangle=1$, $\langle\alpha_2, \beta_2\rangle=1$, 
\\
\begin{tabular}{c}2$\alpha_1$\\\tiny{ $A_1$ }\end{tabular} &\begin{tabular}{c} = \\~ \end{tabular} & \begin{tabular}{cccc}$\beta_1$+ & $\beta_2$+ & $\beta_3$+ & $\beta_4$\\\tiny{ $A_1$+$A_1$' } & \tiny{ $A_1$+$A_1$'' } & \tiny{ $A_1$' } & \tiny{ $A_1$'' }\end{tabular} & $D_4$ & $E_7$ &  $\langle\alpha_1, \beta_1\rangle=1$, $\langle\alpha_1, \beta_2\rangle=1$, $\langle\alpha_1, \beta_3\rangle=1$, $\langle\alpha_1, \beta_4\rangle=1$, 
\\
\begin{tabular}{cc}$\alpha_1$+ & $\alpha_2$\\\tiny{ $A_1$ } & \tiny{  }\end{tabular} &\begin{tabular}{c} = \\~ \end{tabular} & \begin{tabular}{ccc}$\beta_1$+ & $\beta_2$+ & $\beta_3$\\\tiny{ $A_1$+$A_1$' } & \tiny{ $A_1$''+$A_1$ } & \tiny{ $A_1$'' }\end{tabular} & $D_4$ & $E_7$ & $\langle\alpha_1, \alpha_2\rangle=1$,  $\langle\alpha_1, \beta_1\rangle=1$, $\langle\alpha_1, \beta_2\rangle=1$, $\langle\alpha_1, \beta_3\rangle=1$, $\langle\alpha_2, \beta_1\rangle=1$, $\langle\alpha_2, \beta_2\rangle=1$, $\langle\alpha_2, \beta_3\rangle=1$, 
\\
\begin{tabular}{cc}$\alpha_1$+ & $\alpha_2$\\\tiny{ $A_1$ } & \tiny{ $A_1$' }\end{tabular} &\begin{tabular}{c} = \\~ \end{tabular} & \begin{tabular}{ccc}$\beta_1$+ & $\beta_2$+ & $\beta_3$\\\tiny{ $A_1$+$A_1$' } & \tiny{ $A_1$' } & \tiny{ $A_1$'' }\end{tabular} & $D_4$ & $E_7$ & $\langle\alpha_1, \alpha_2\rangle=1$,  $\langle\alpha_1, \beta_1\rangle=1$, $\langle\alpha_1, \beta_2\rangle=1$, $\langle\alpha_1, \beta_3\rangle=1$, $\langle\alpha_2, \beta_1\rangle=1$, $\langle\alpha_2, \beta_2\rangle=1$, $\langle\alpha_2, \beta_3\rangle=1$, 
\\
\begin{tabular}{cc}$\alpha_1$+ & $\alpha_2$\\\tiny{ $A_1$ } & \tiny{ $A_1$' }\end{tabular} &\begin{tabular}{c} = \\~ \end{tabular} & \begin{tabular}{cc}$\beta_1$+ & $\beta_2$\\\tiny{ $A_1$+$A_1$' } & \tiny{ $A_1$''+$A_1$' }\end{tabular} & $A_3$ & $E_6$ &  $\langle\alpha_1, \beta_1\rangle=1$, $\langle\alpha_1, \beta_2\rangle=1$, $\langle\alpha_2, \beta_1\rangle=1$, $\langle\alpha_2, \beta_2\rangle=1$, 
\\
\begin{tabular}{cc}$\alpha_1$+ & $\alpha_2$\\\tiny{ $A_1$+$A_1$' } & \tiny{  }\end{tabular} &\begin{tabular}{c} = \\~ \end{tabular} & \begin{tabular}{ccc}$\beta_1$+ & $\beta_2$+ & $\beta_3$\\\tiny{ $A_1$+$A_1$' } & \tiny{ $A_1$'' } & \tiny{ $A_1$ }\end{tabular} & $D_4$ & $E_7$ & $\langle\alpha_1, \alpha_2\rangle=1$,  $\langle\alpha_1, \beta_1\rangle=1$, $\langle\alpha_1, \beta_2\rangle=1$, $\langle\alpha_1, \beta_3\rangle=1$, $\langle\alpha_2, \beta_1\rangle=1$, $\langle\alpha_2, \beta_2\rangle=1$, $\langle\alpha_2, \beta_3\rangle=1$, 
\\
\begin{tabular}{cc}$\alpha_1$+ & $\alpha_2$\\\tiny{ $A_1$ } & \tiny{  }\end{tabular} &\begin{tabular}{c} = \\~ \end{tabular} & \begin{tabular}{ccc}$\beta_1$+ & $\beta_2$+ & $\beta_3$\\\tiny{ $A_1$ } & \tiny{ $A_1$' } & \tiny{ $A_1$' }\end{tabular} & $D_4$ & $D_6$+$A_1$ & $\langle\alpha_1, \alpha_2\rangle=1$,  $\langle\alpha_1, \beta_1\rangle=1$, $\langle\alpha_1, \beta_2\rangle=1$, $\langle\alpha_1, \beta_3\rangle=1$, $\langle\alpha_2, \beta_1\rangle=1$, $\langle\alpha_2, \beta_2\rangle=1$, $\langle\alpha_2, \beta_3\rangle=1$, 
\\
\begin{tabular}{cc}$\alpha_1$+ & $\alpha_2$\\\tiny{ $A_1$+$A_1$' } & \tiny{ $A_1$'' }\end{tabular} &\begin{tabular}{c} = \\~ \end{tabular} & \begin{tabular}{cc}$\beta_1$+ & $\beta_2$\\\tiny{ $A_1$+$A_1$' } & \tiny{ $A_1$'' }\end{tabular} & $A_3$ & $A_6$ &  $\langle\alpha_1, \beta_1\rangle=1$, $\langle\alpha_1, \beta_2\rangle=1$, $\langle\alpha_2, \beta_1\rangle=1$, $\langle\alpha_2, \beta_2\rangle=1$, 
\\
\begin{tabular}{cc}$\alpha_1$+ & $\alpha_2$\\\tiny{ $A_1$ } & \tiny{  }\end{tabular} &\begin{tabular}{c} = \\~ \end{tabular} & \begin{tabular}{ccc}$\beta_1$+ & $\beta_2$+ & $\beta_3$\\\tiny{ $A_1$ } & \tiny{ $A_1$ } & \tiny{ $A_1$' }\end{tabular} & $D_4$ & $D_6$+$A_1$ & $\langle\alpha_1, \alpha_2\rangle=1$,  $\langle\alpha_1, \beta_1\rangle=1$, $\langle\alpha_1, \beta_2\rangle=1$, $\langle\alpha_1, \beta_3\rangle=1$, $\langle\alpha_2, \beta_1\rangle=1$, $\langle\alpha_2, \beta_2\rangle=1$, $\langle\alpha_2, \beta_3\rangle=1$, 
\\
\begin{tabular}{cc}$\alpha_1$+ & $\alpha_2$\\\tiny{ $A_1$+$A_1$' } & \tiny{  }\end{tabular} &\begin{tabular}{c} = \\~ \end{tabular} & \begin{tabular}{cc}$\beta_1$+ & $\beta_2$\\\tiny{ $A_1$+$A_1$' } & \tiny{ $A_1$''+$A_1$ }\end{tabular} & $A_3$ & $D_6$ &  $\langle\alpha_1, \beta_1\rangle=1$, $\langle\alpha_1, \beta_2\rangle=1$, $\langle\alpha_2, \beta_1\rangle=1$, $\langle\alpha_2, \beta_2\rangle=1$, 
\\
\begin{tabular}{cc}$\alpha_1$+ & $\alpha_2$\\\tiny{ $A_1$+$A_1$' } & \tiny{ $A_1$'' }\end{tabular} &\begin{tabular}{c} = \\~ \end{tabular} & \begin{tabular}{ccc}$\beta_1$+ & $\beta_2$+ & $\beta_3$\\\tiny{ $A_1$ } & \tiny{ $A_1$' } & \tiny{ $A_1$'' }\end{tabular} & $D_4$ & $E_6$ & $\langle\alpha_1, \alpha_2\rangle=1$,  $\langle\alpha_1, \beta_1\rangle=1$, $\langle\alpha_1, \beta_2\rangle=1$, $\langle\alpha_1, \beta_3\rangle=1$, $\langle\alpha_2, \beta_1\rangle=1$, $\langle\alpha_2, \beta_2\rangle=1$, $\langle\alpha_2, \beta_3\rangle=1$, 
\\
\begin{tabular}{c}2$\alpha_1$\\\tiny{  }\end{tabular} &\begin{tabular}{c} = \\~ \end{tabular} & \begin{tabular}{cccc}$\beta_1$+ & $\beta_2$+ & $\beta_3$+ & $\beta_4$\\\tiny{ $A_1$ } & \tiny{ $A_1$ } & \tiny{ $A_1$' } & \tiny{ $A_1$' }\end{tabular} & $D_4$ & $D_6$+$A_1$ &  $\langle\alpha_1, \beta_1\rangle=1$, $\langle\alpha_1, \beta_2\rangle=1$, $\langle\alpha_1, \beta_3\rangle=1$, $\langle\alpha_1, \beta_4\rangle=1$, 
\\
\begin{tabular}{cc}$\alpha_1$+ & $\alpha_2$\\\tiny{ $A_1$+$A_1$'+$A_1$'' } & \tiny{ $A_1$ }\end{tabular} &\begin{tabular}{c} = \\~ \end{tabular} & \begin{tabular}{cc}$\beta_1$+ & $\beta_2$\\\tiny{ $A_1$+$A_1$' } & \tiny{ $A_1$+$A_1$'' }\end{tabular} & $A_3$ & $E_6$ &  $\langle\alpha_1, \beta_1\rangle=1$, $\langle\alpha_1, \beta_2\rangle=1$, $\langle\alpha_2, \beta_1\rangle=1$, $\langle\alpha_2, \beta_2\rangle=1$, 
\\
\begin{tabular}{cc}$\alpha_1$+ & $\alpha_2$\\\tiny{ $A_1$ } & \tiny{ $A_1$ }\end{tabular} &\begin{tabular}{c} = \\~ \end{tabular} & \begin{tabular}{ccc}$\beta_1$+ & $\beta_2$+ & $\beta_3$\\\tiny{ $A_1$+$A_1$' } & \tiny{ $A_1$''+$A_1$' } & \tiny{ $A_1$''+$A_1$ }\end{tabular} & $D_4$ & $E_6$ & $\langle\alpha_1, \alpha_2\rangle=1$,  $\langle\alpha_1, \beta_1\rangle=1$, $\langle\alpha_1, \beta_2\rangle=1$, $\langle\alpha_1, \beta_3\rangle=1$, $\langle\alpha_2, \beta_1\rangle=1$, $\langle\alpha_2, \beta_2\rangle=1$, $\langle\alpha_2, \beta_3\rangle=1$, 
\\
\begin{tabular}{cc}$\alpha_1$+ & $\alpha_2$\\\tiny{ $A_1$+$A_1$'+$A_1$'' } & \tiny{ $A_1$' }\end{tabular} &\begin{tabular}{c} = \\~ \end{tabular} & \begin{tabular}{cc}$\beta_1$+ & $\beta_2$\\\tiny{ $A_1$+$A_1$' } & \tiny{ $A_1$''+$A_1$ }\end{tabular} & $A_3$ & $E_6$ &  $\langle\alpha_1, \beta_1\rangle=1$, $\langle\alpha_1, \beta_2\rangle=1$, $\langle\alpha_2, \beta_1\rangle=1$, $\langle\alpha_2, \beta_2\rangle=1$, 
\\
\begin{tabular}{cc}$\alpha_1$+ & $\alpha_2$\\\tiny{ $A_1$ } & \tiny{ $A_1$' }\end{tabular} &\begin{tabular}{c} = \\~ \end{tabular} & \begin{tabular}{cc}$\beta_1$+ & $\beta_2$\\\tiny{ $A_1$+$A_1$' } & \tiny{ $A_1$+$A_1$'' }\end{tabular} & $A_3$ & $E_6$ &  $\langle\alpha_1, \beta_1\rangle=1$, $\langle\alpha_1, \beta_2\rangle=1$, $\langle\alpha_2, \beta_1\rangle=1$, $\langle\alpha_2, \beta_2\rangle=1$, 
\\
\begin{tabular}{cc}$\alpha_1$+ & $\alpha_2$\\\tiny{ $A_1$+$A_1$'+$A_1$'' } & \tiny{  }\end{tabular} &\begin{tabular}{c} = \\~ \end{tabular} & \begin{tabular}{ccc}$\beta_1$+ & $\beta_2$+ & $\beta_3$\\\tiny{ $A_1$ } & \tiny{ $A_1$' } & \tiny{ $A_1$'' }\end{tabular} & $D_4$ & $E_6$ & $\langle\alpha_1, \alpha_2\rangle=1$,  $\langle\alpha_1, \beta_1\rangle=1$, $\langle\alpha_1, \beta_2\rangle=1$, $\langle\alpha_1, \beta_3\rangle=1$, $\langle\alpha_2, \beta_1\rangle=1$, $\langle\alpha_2, \beta_2\rangle=1$, $\langle\alpha_2, \beta_3\rangle=1$, 
\\
\begin{tabular}{cc}$\alpha_1$+ & $\alpha_2$\\\tiny{ $A_1$+$A_1$' } & \tiny{ $A_1$' }\end{tabular} &\begin{tabular}{c} = \\~ \end{tabular} & \begin{tabular}{ccc}$\beta_1$+ & $\beta_2$+ & $\beta_3$\\\tiny{ $A_1$+$A_1$' } & \tiny{ $A_1$+$A_1$'' } & \tiny{ $A_1$' }\end{tabular} & $D_4$ & $E_7$ & $\langle\alpha_1, \alpha_2\rangle=1$,  $\langle\alpha_1, \beta_1\rangle=1$, $\langle\alpha_1, \beta_2\rangle=1$, $\langle\alpha_1, \beta_3\rangle=1$, $\langle\alpha_2, \beta_1\rangle=1$, $\langle\alpha_2, \beta_2\rangle=1$, $\langle\alpha_2, \beta_3\rangle=1$, 
\\
\begin{tabular}{cc}$\alpha_1$+ & $\alpha_2$\\\tiny{ $A_1$ } & \tiny{  }\end{tabular} &\begin{tabular}{c} = \\~ \end{tabular} & \begin{tabular}{ccc}$\beta_1$+ & $\beta_2$+ & $\beta_3$\\\tiny{ $A_1$+$A_1$' } & \tiny{ $A_1$''+$A_1$ } & \tiny{ $A_1$' }\end{tabular} & $D_4$ & $E_7$ & $\langle\alpha_1, \alpha_2\rangle=1$,  $\langle\alpha_1, \beta_1\rangle=1$, $\langle\alpha_1, \beta_2\rangle=1$, $\langle\alpha_1, \beta_3\rangle=1$, $\langle\alpha_2, \beta_1\rangle=1$, $\langle\alpha_2, \beta_2\rangle=1$, $\langle\alpha_2, \beta_3\rangle=1$, 
\\
\begin{tabular}{cc}$\alpha_1$+ & $\alpha_2$\\\tiny{ $A_1$ } & \tiny{ $A_1$' }\end{tabular} &\begin{tabular}{c} = \\~ \end{tabular} & \begin{tabular}{ccc}$\beta_1$+ & $\beta_2$+ & $\beta_3$\\\tiny{ $A_1$+$A_1$' } & \tiny{ $A_1$'' } & \tiny{ $A_1$' }\end{tabular} & $D_4$ & $E_7$ & $\langle\alpha_1, \alpha_2\rangle=1$,  $\langle\alpha_1, \beta_1\rangle=1$, $\langle\alpha_1, \beta_2\rangle=1$, $\langle\alpha_1, \beta_3\rangle=1$, $\langle\alpha_2, \beta_1\rangle=1$, $\langle\alpha_2, \beta_2\rangle=1$, $\langle\alpha_2, \beta_3\rangle=1$, 
\\
\begin{tabular}{cc}$\alpha_1$+ & $\alpha_2$\\\tiny{ $A_1$+$A_1$'+$A_1$'' } & \tiny{ $A_1$ }\end{tabular} &\begin{tabular}{c} = \\~ \end{tabular} & \begin{tabular}{ccc}$\beta_1$+ & $\beta_2$+ & $\beta_3$\\\tiny{ $A_1$+$A_1$' } & \tiny{ $A_1$'' } & \tiny{ $A_1$ }\end{tabular} & $D_4$ & $E_7$ & $\langle\alpha_1, \alpha_2\rangle=1$,  $\langle\alpha_1, \beta_1\rangle=1$, $\langle\alpha_1, \beta_2\rangle=1$, $\langle\alpha_1, \beta_3\rangle=1$, $\langle\alpha_2, \beta_1\rangle=1$, $\langle\alpha_2, \beta_2\rangle=1$, $\langle\alpha_2, \beta_3\rangle=1$, 
\\
\begin{tabular}{cc}$\alpha_1$+ & $\alpha_2$\\\tiny{ $A_1$+$A_1$' } & \tiny{ $A_1$' }\end{tabular} &\begin{tabular}{c} = \\~ \end{tabular} & \begin{tabular}{ccc}$\beta_1$+ & $\beta_2$+ & $\beta_3$\\\tiny{ $A_1$+$A_1$' } & \tiny{ $A_1$''+$A_1$ } & \tiny{ $A_1$' }\end{tabular} & $D_4$ & $E_7$ & $\langle\alpha_1, \alpha_2\rangle=1$,  $\langle\alpha_1, \beta_1\rangle=1$, $\langle\alpha_1, \beta_2\rangle=1$, $\langle\alpha_1, \beta_3\rangle=1$, $\langle\alpha_2, \beta_1\rangle=1$, $\langle\alpha_2, \beta_2\rangle=1$, $\langle\alpha_2, \beta_3\rangle=1$, 
\\
\begin{tabular}{cc}$\alpha_1$+ & $\alpha_2$\\\tiny{ $A_1$+$A_1$'+$A_1$'' } & \tiny{ $A_1$'' }\end{tabular} &\begin{tabular}{c} = \\~ \end{tabular} & \begin{tabular}{cc}$\beta_1$+ & $\beta_2$\\\tiny{ $A_1$+$A_1$' } & \tiny{ $A_1$''+$A_1$' }\end{tabular} & $A_3$ & $E_6$ &  $\langle\alpha_1, \beta_1\rangle=1$, $\langle\alpha_1, \beta_2\rangle=1$, $\langle\alpha_2, \beta_1\rangle=1$, $\langle\alpha_2, \beta_2\rangle=1$, 
\\
\begin{tabular}{cc}$\alpha_1$+ & $\alpha_2$\\\tiny{ $A_1$+$A_1$' } & \tiny{  }\end{tabular} &\begin{tabular}{c} = \\~ \end{tabular} & \begin{tabular}{c}$\beta_1$\\\tiny{ $A_1$+$A_1$' }\end{tabular} & $A_2$ & $D_4$+$A_1$ & $\langle\alpha_1, \alpha_2\rangle=-1$,  $\langle\alpha_1, \beta_1\rangle=1$, $\langle\alpha_2, \beta_1\rangle=1$, 
\\
\begin{tabular}{cc}$\alpha_1$+ & $\alpha_2$\\\tiny{ $A_1$+$A_1$' } & \tiny{ $A_1$'' }\end{tabular} &\begin{tabular}{c} = \\~ \end{tabular} & \begin{tabular}{cc}$\beta_1$+ & $\beta_2$\\\tiny{ $A_1$+$A_1$' } & \tiny{ $A_1$'' }\end{tabular} & $A_3$ & $D_6$ &  $\langle\alpha_1, \beta_1\rangle=1$, $\langle\alpha_1, \beta_2\rangle=1$, $\langle\alpha_2, \beta_1\rangle=1$, $\langle\alpha_2, \beta_2\rangle=1$, 
\\
\begin{tabular}{cc}$\alpha_1$+ & $\alpha_2$\\\tiny{ $A_1$ } & \tiny{ $A_1$ }\end{tabular} &\begin{tabular}{c} = \\~ \end{tabular} & \begin{tabular}{cc}$\beta_1$+ & $\beta_2$\\\tiny{ $A_1$ } & \tiny{ $A_1$ }\end{tabular} & $A_3$ & $D_4$+2$A_1$ &  $\langle\alpha_1, \beta_1\rangle=1$, $\langle\alpha_1, \beta_2\rangle=1$, $\langle\alpha_2, \beta_1\rangle=1$, $\langle\alpha_2, \beta_2\rangle=1$, 
\\
\begin{tabular}{cc}$\alpha_1$+ & $\alpha_2$\\\tiny{ $A_1$+$A_1$' } & \tiny{ $A_1$ }\end{tabular} &\begin{tabular}{c} = \\~ \end{tabular} & \begin{tabular}{cc}$\beta_1$+ & $\beta_2$\\\tiny{ $A_1$+$A_1$' } & \tiny{ $A_1$ }\end{tabular} & $A_3$ & $D_5$+$A_1$ &  $\langle\alpha_1, \beta_1\rangle=1$, $\langle\alpha_1, \beta_2\rangle=1$, $\langle\alpha_2, \beta_1\rangle=1$, $\langle\alpha_2, \beta_2\rangle=1$, 
\\
\begin{tabular}{cc}$\alpha_1$+ & $\alpha_2$\\\tiny{ $A_1$ } & \tiny{ $A_1$ }\end{tabular} &\begin{tabular}{c} = \\~ \end{tabular} & \begin{tabular}{cc}$\beta_1$+ & $\beta_2$\\\tiny{ $A_1$+$A_1$' } & \tiny{ $A_1$+$A_1$' }\end{tabular} & $A_3$ & $D_5$+$A_1$ &  $\langle\alpha_1, \beta_1\rangle=1$, $\langle\alpha_1, \beta_2\rangle=1$, $\langle\alpha_2, \beta_1\rangle=1$, $\langle\alpha_2, \beta_2\rangle=1$, 
\\
\begin{tabular}{c}2$\alpha_1$\\\tiny{ $A_1$+$A_1$' }\end{tabular} &\begin{tabular}{c} = \\~ \end{tabular} & \begin{tabular}{cccc}$\beta_1$+ & $\beta_2$+ & $\beta_3$+ & $\beta_4$\\\tiny{ $A_1$+$A_1$' } & \tiny{ $A_1$''+$A_1$' } & \tiny{ $A_1$ } & \tiny{ $A_1$'' }\end{tabular} & $D_4$ & $E_7$ &  $\langle\alpha_1, \beta_1\rangle=1$, $\langle\alpha_1, \beta_2\rangle=1$, $\langle\alpha_1, \beta_3\rangle=1$, $\langle\alpha_1, \beta_4\rangle=1$, 
\\
\begin{tabular}{cc}$\alpha_1$+ & $\alpha_2$\\\tiny{ $A_1$+$A_1$' } & \tiny{ $A_1$' }\end{tabular} &\begin{tabular}{c} = \\~ \end{tabular} & \begin{tabular}{ccc}$\beta_1$+ & $\beta_2$+ & $\beta_3$\\\tiny{ $A_1$ } & \tiny{ $A_1$ } & \tiny{ $A_1$' }\end{tabular} & $D_4$ & $D_6$+$A_1$ & $\langle\alpha_1, \alpha_2\rangle=1$,  $\langle\alpha_1, \beta_1\rangle=1$, $\langle\alpha_1, \beta_2\rangle=1$, $\langle\alpha_1, \beta_3\rangle=1$, $\langle\alpha_2, \beta_1\rangle=1$, $\langle\alpha_2, \beta_2\rangle=1$, $\langle\alpha_2, \beta_3\rangle=1$, 
\\
\begin{tabular}{cc}$\alpha_1$+ & $\alpha_2$\\\tiny{ $A_1$+$A_1$' } & \tiny{ $A_1$+$A_1$'' }\end{tabular} &\begin{tabular}{c} = \\~ \end{tabular} & \begin{tabular}{ccc}$\beta_1$+ & $\beta_2$+ & $\beta_3$\\\tiny{ $A_1$+$A_1$' } & \tiny{ $A_1$'' } & \tiny{ $A_1$ }\end{tabular} & $D_4$ & $E_7$ & $\langle\alpha_1, \alpha_2\rangle=1$,  $\langle\alpha_1, \beta_1\rangle=1$, $\langle\alpha_1, \beta_2\rangle=1$, $\langle\alpha_1, \beta_3\rangle=1$, $\langle\alpha_2, \beta_1\rangle=1$, $\langle\alpha_2, \beta_2\rangle=1$, $\langle\alpha_2, \beta_3\rangle=1$, 
\\
\begin{tabular}{c}2$\alpha_1$\\\tiny{ $A_1$ }\end{tabular} &\begin{tabular}{c} = \\~ \end{tabular} & \begin{tabular}{cccc}$\beta_1$+ & $\beta_2$+ & $\beta_3$+ & $\beta_4$\\\tiny{ $A_1$ } & \tiny{ $A_1$ } & \tiny{ $A_1$' } & \tiny{ $A_1$' }\end{tabular} & $D_4$ & $D_6$+$A_1$ &  $\langle\alpha_1, \beta_1\rangle=1$, $\langle\alpha_1, \beta_2\rangle=1$, $\langle\alpha_1, \beta_3\rangle=1$, $\langle\alpha_1, \beta_4\rangle=1$, 
\\
\begin{tabular}{cc}$\alpha_1$+ & $\alpha_2$\\\tiny{ $A_1$ } & \tiny{ $A_1$' }\end{tabular} &\begin{tabular}{c} = \\~ \end{tabular} & \begin{tabular}{cc}$\beta_1$+ & $\beta_2$\\\tiny{ $A_1$+$A_1$' } & \tiny{  }\end{tabular} & $A_3$ & $A_5$+$A_1$ &  $\langle\alpha_1, \beta_1\rangle=1$, $\langle\alpha_1, \beta_2\rangle=1$, $\langle\alpha_2, \beta_1\rangle=1$, $\langle\alpha_2, \beta_2\rangle=1$, 
\\
\begin{tabular}{cc}$\alpha_1$+ & $\alpha_2$\\\tiny{ $A_1$+$A_1$' } & \tiny{ $A_1$' }\end{tabular} &\begin{tabular}{c} = \\~ \end{tabular} & \begin{tabular}{ccc}$\beta_1$+ & $\beta_2$+ & $\beta_3$\\\tiny{ $A_1$+$A_1$' } & \tiny{ $A_1$''+$A_1$' } & \tiny{ $A_1$ }\end{tabular} & $D_4$ & $E_7$ & $\langle\alpha_1, \alpha_2\rangle=1$,  $\langle\alpha_1, \beta_1\rangle=1$, $\langle\alpha_1, \beta_2\rangle=1$, $\langle\alpha_1, \beta_3\rangle=1$, $\langle\alpha_2, \beta_1\rangle=1$, $\langle\alpha_2, \beta_2\rangle=1$, $\langle\alpha_2, \beta_3\rangle=1$, 
\\
\begin{tabular}{cc}$\alpha_1$+ & $\alpha_2$\\\tiny{ $A_1$+$A_1$' } & \tiny{ $A_1$'' }\end{tabular} &\begin{tabular}{c} = \\~ \end{tabular} & \begin{tabular}{ccc}$\beta_1$+ & $\beta_2$+ & $\beta_3$\\\tiny{ $A_1$+$A_1$' } & \tiny{ $A_1$'' } & \tiny{  }\end{tabular} & $D_4$ & $E_7$ & $\langle\alpha_1, \alpha_2\rangle=1$,  $\langle\alpha_1, \beta_1\rangle=1$, $\langle\alpha_1, \beta_2\rangle=1$, $\langle\alpha_1, \beta_3\rangle=1$, $\langle\alpha_2, \beta_1\rangle=1$, $\langle\alpha_2, \beta_2\rangle=1$, $\langle\alpha_2, \beta_3\rangle=1$, 
\\
\begin{tabular}{c}2$\alpha_1$\\\tiny{ $A_1$ }\end{tabular} &\begin{tabular}{c} = \\~ \end{tabular} & \begin{tabular}{cccc}$\beta_1$+ & $\beta_2$+ & $\beta_3$+ & $\beta_4$\\\tiny{ $A_1$+$A_1$' } & \tiny{ $A_1$''+$A_1$' } & \tiny{ $A_1$ } & \tiny{ $A_1$'' }\end{tabular} & $D_4$ & $E_7$ &  $\langle\alpha_1, \beta_1\rangle=1$, $\langle\alpha_1, \beta_2\rangle=1$, $\langle\alpha_1, \beta_3\rangle=1$, $\langle\alpha_1, \beta_4\rangle=1$, 
\\
\begin{tabular}{cc}$\alpha_1$+ & $\alpha_2$\\\tiny{ $A_1$ } & \tiny{ $A_1$' }\end{tabular} &\begin{tabular}{c} = \\~ \end{tabular} & \begin{tabular}{ccc}$\beta_1$+ & $\beta_2$+ & $\beta_3$\\\tiny{ $A_1$+$A_1$' } & \tiny{ $A_1$'' } & \tiny{ $A_1$ }\end{tabular} & $D_4$ & $E_7$ & $\langle\alpha_1, \alpha_2\rangle=1$,  $\langle\alpha_1, \beta_1\rangle=1$, $\langle\alpha_1, \beta_2\rangle=1$, $\langle\alpha_1, \beta_3\rangle=1$, $\langle\alpha_2, \beta_1\rangle=1$, $\langle\alpha_2, \beta_2\rangle=1$, $\langle\alpha_2, \beta_3\rangle=1$, 
\\
\begin{tabular}{cc}$\alpha_1$+ & $\alpha_2$\\\tiny{ $A_1$+$A_1$'+$A_1$'' } & \tiny{ $A_1$+$A_1$' }\end{tabular} &\begin{tabular}{c} = \\~ \end{tabular} & \begin{tabular}{ccc}$\beta_1$+ & $\beta_2$+ & $\beta_3$\\\tiny{ $A_1$+$A_1$' } & \tiny{ $A_1$''+$A_1$ } & \tiny{ $A_1$'' }\end{tabular} & $D_4$ & $E_7$ & $\langle\alpha_1, \alpha_2\rangle=1$,  $\langle\alpha_1, \beta_1\rangle=1$, $\langle\alpha_1, \beta_2\rangle=1$, $\langle\alpha_1, \beta_3\rangle=1$, $\langle\alpha_2, \beta_1\rangle=1$, $\langle\alpha_2, \beta_2\rangle=1$, $\langle\alpha_2, \beta_3\rangle=1$, 
\\
\begin{tabular}{cc}$\alpha_1$+ & $\alpha_2$\\\tiny{ $A_1$+$A_1$' } & \tiny{ $A_1$ }\end{tabular} &\begin{tabular}{c} = \\~ \end{tabular} & \begin{tabular}{ccc}$\beta_1$+ & $\beta_2$+ & $\beta_3$\\\tiny{ $A_1$+$A_1$' } & \tiny{ $A_1$''+$A_1$' } & \tiny{ $A_1$'' }\end{tabular} & $D_4$ & $E_7$ & $\langle\alpha_1, \alpha_2\rangle=1$,  $\langle\alpha_1, \beta_1\rangle=1$, $\langle\alpha_1, \beta_2\rangle=1$, $\langle\alpha_1, \beta_3\rangle=1$, $\langle\alpha_2, \beta_1\rangle=1$, $\langle\alpha_2, \beta_2\rangle=1$, $\langle\alpha_2, \beta_3\rangle=1$, 
\\
\begin{tabular}{cc}$\alpha_1$+ & $\alpha_2$\\\tiny{ $A_1$ } & \tiny{ $A_1$' }\end{tabular} &\begin{tabular}{c} = \\~ \end{tabular} & \begin{tabular}{ccc}$\beta_1$+ & $\beta_2$+ & $\beta_3$\\\tiny{ $A_1$+$A_1$' } & \tiny{ $A_1$ } & \tiny{ $A_1$'' }\end{tabular} & $D_4$ & $E_7$ & $\langle\alpha_1, \alpha_2\rangle=1$,  $\langle\alpha_1, \beta_1\rangle=1$, $\langle\alpha_1, \beta_2\rangle=1$, $\langle\alpha_1, \beta_3\rangle=1$, $\langle\alpha_2, \beta_1\rangle=1$, $\langle\alpha_2, \beta_2\rangle=1$, $\langle\alpha_2, \beta_3\rangle=1$, 
\\
\begin{tabular}{cc}$\alpha_1$+ & $\alpha_2$\\\tiny{ $A_1$+$A_1$' } & \tiny{ $A_1$''+$A_1$' }\end{tabular} &\begin{tabular}{c} = \\~ \end{tabular} & \begin{tabular}{cc}$\beta_1$+ & $\beta_2$\\\tiny{ $A_1$+$A_1$' } & \tiny{ $A_1$''+$A_1$' }\end{tabular} & $A_3$ & $D_6$ &  $\langle\alpha_1, \beta_1\rangle=1$, $\langle\alpha_1, \beta_2\rangle=1$, $\langle\alpha_2, \beta_1\rangle=1$, $\langle\alpha_2, \beta_2\rangle=1$, 
\\
\begin{tabular}{cc}$\alpha_1$+ & $\alpha_2$\\\tiny{ $A_1$+$A_1$' } & \tiny{ $A_1$' }\end{tabular} &\begin{tabular}{c} = \\~ \end{tabular} & \begin{tabular}{ccc}$\beta_1$+ & $\beta_2$+ & $\beta_3$\\\tiny{ $A_1$+$A_1$' } & \tiny{ $A_1$''+$A_1$' } & \tiny{ $A_1$'' }\end{tabular} & $D_4$ & $E_7$ & $\langle\alpha_1, \alpha_2\rangle=1$,  $\langle\alpha_1, \beta_1\rangle=1$, $\langle\alpha_1, \beta_2\rangle=1$, $\langle\alpha_1, \beta_3\rangle=1$, $\langle\alpha_2, \beta_1\rangle=1$, $\langle\alpha_2, \beta_2\rangle=1$, $\langle\alpha_2, \beta_3\rangle=1$, 
\\
\begin{tabular}{cc}$\alpha_1$+ & $\alpha_2$\\\tiny{ $A_1$+$A_1$'+$A_1$'' } & \tiny{ $A_1$'+$A_1$'' }\end{tabular} &\begin{tabular}{c} = \\~ \end{tabular} & \begin{tabular}{ccc}$\beta_1$+ & $\beta_2$+ & $\beta_3$\\\tiny{ $A_1$+$A_1$' } & \tiny{ $A_1$''+$A_1$' } & \tiny{ $A_1$ }\end{tabular} & $D_4$ & $E_7$ & $\langle\alpha_1, \alpha_2\rangle=1$,  $\langle\alpha_1, \beta_1\rangle=1$, $\langle\alpha_1, \beta_2\rangle=1$, $\langle\alpha_1, \beta_3\rangle=1$, $\langle\alpha_2, \beta_1\rangle=1$, $\langle\alpha_2, \beta_2\rangle=1$, $\langle\alpha_2, \beta_3\rangle=1$, 
\\
\begin{tabular}{cc}$\alpha_1$+ & $\alpha_2$\\\tiny{ $A_1$ } & \tiny{ $A_1$ }\end{tabular} &\begin{tabular}{c} = \\~ \end{tabular} & \begin{tabular}{ccc}$\beta_1$+ & $\beta_2$+ & $\beta_3$\\\tiny{ $A_1$+$A_1$' } & \tiny{ $A_1$+$A_1$' } & \tiny{  }\end{tabular} & $D_4$ & $D_6$+$A_1$ & $\langle\alpha_1, \alpha_2\rangle=1$,  $\langle\alpha_1, \beta_1\rangle=1$, $\langle\alpha_1, \beta_2\rangle=1$, $\langle\alpha_1, \beta_3\rangle=1$, $\langle\alpha_2, \beta_1\rangle=1$, $\langle\alpha_2, \beta_2\rangle=1$, $\langle\alpha_2, \beta_3\rangle=1$, 
\\
\begin{tabular}{c}2$\alpha_1$\\\tiny{ $A_1$+$A_1$' }\end{tabular} &\begin{tabular}{c} = \\~ \end{tabular} & \begin{tabular}{cccc}$\beta_1$+ & $\beta_2$+ & $\beta_3$+ & $\beta_4$\\\tiny{ $A_1$ } & \tiny{ $A_1$ } & \tiny{ $A_1$' } & \tiny{ $A_1$' }\end{tabular} & $D_4$ & $D_6$+$A_1$ &  $\langle\alpha_1, \beta_1\rangle=1$, $\langle\alpha_1, \beta_2\rangle=1$, $\langle\alpha_1, \beta_3\rangle=1$, $\langle\alpha_1, \beta_4\rangle=1$, 
\\
\begin{tabular}{cc}$\alpha_1$+ & $\alpha_2$\\\tiny{ $A_1$+$A_1$' } & \tiny{ $A_1$ }\end{tabular} &\begin{tabular}{c} = \\~ \end{tabular} & \begin{tabular}{ccc}$\beta_1$+ & $\beta_2$+ & $\beta_3$\\\tiny{ $A_1$ } & \tiny{ $A_1$' } & \tiny{ $A_1$' }\end{tabular} & $D_4$ & $D_6$+$A_1$ & $\langle\alpha_1, \alpha_2\rangle=1$,  $\langle\alpha_1, \beta_1\rangle=1$, $\langle\alpha_1, \beta_2\rangle=1$, $\langle\alpha_1, \beta_3\rangle=1$, $\langle\alpha_2, \beta_1\rangle=1$, $\langle\alpha_2, \beta_2\rangle=1$, $\langle\alpha_2, \beta_3\rangle=1$, 
\\
\begin{tabular}{cc}$\alpha_1$+ & $\alpha_2$\\\tiny{ $A_1$ } & \tiny{  }\end{tabular} &\begin{tabular}{c} = \\~ \end{tabular} & \begin{tabular}{ccc}$\beta_1$+ & $\beta_2$+ & $\beta_3$\\\tiny{ $A_1$+$A_1$' } & \tiny{ $A_1$''+$A_1$' } & \tiny{ $A_1$ }\end{tabular} & $D_4$ & $E_7$ & $\langle\alpha_1, \alpha_2\rangle=1$,  $\langle\alpha_1, \beta_1\rangle=1$, $\langle\alpha_1, \beta_2\rangle=1$, $\langle\alpha_1, \beta_3\rangle=1$, $\langle\alpha_2, \beta_1\rangle=1$, $\langle\alpha_2, \beta_2\rangle=1$, $\langle\alpha_2, \beta_3\rangle=1$, 
\\
\begin{tabular}{cc}$\alpha_1$+ & $\alpha_2$\\\tiny{ $A_1$+$A_1$' } & \tiny{ $A_1$''+$A_1$ }\end{tabular} &\begin{tabular}{c} = \\~ \end{tabular} & \begin{tabular}{ccc}$\beta_1$+ & $\beta_2$+ & $\beta_3$\\\tiny{ $A_1$+$A_1$' } & \tiny{ $A_1$ } & \tiny{ $A_1$'' }\end{tabular} & $D_4$ & $E_7$ & $\langle\alpha_1, \alpha_2\rangle=1$,  $\langle\alpha_1, \beta_1\rangle=1$, $\langle\alpha_1, \beta_2\rangle=1$, $\langle\alpha_1, \beta_3\rangle=1$, $\langle\alpha_2, \beta_1\rangle=1$, $\langle\alpha_2, \beta_2\rangle=1$, $\langle\alpha_2, \beta_3\rangle=1$, 
\\
\begin{tabular}{c}2$\alpha_1$\\\tiny{ $A_1$ }\end{tabular} &\begin{tabular}{c} = \\~ \end{tabular} & \begin{tabular}{cccc}$\beta_1$+ & $\beta_2$+ & $\beta_3$+ & $\beta_4$\\\tiny{ $A_1$+$A_1$' } & \tiny{ $A_1$''+$A_1$' } & \tiny{ $A_1$''+$A_1$ } & \tiny{  }\end{tabular} & $D_4$ & $E_6$ &  $\langle\alpha_1, \beta_1\rangle=1$, $\langle\alpha_1, \beta_2\rangle=1$, $\langle\alpha_1, \beta_3\rangle=1$, $\langle\alpha_1, \beta_4\rangle=1$, 
\\
\begin{tabular}{c}2$\alpha_1$\\\tiny{ $A_1$ }\end{tabular} &\begin{tabular}{c} = \\~ \end{tabular} & \begin{tabular}{cccc}$\beta_1$+ & $\beta_2$+ & $\beta_3$+ & $\beta_4$\\\tiny{ $A_1$+$A_1$' } & \tiny{ $A_1$+$A_1$' } & \tiny{ $A_1$'' } & \tiny{ $A_1$'' }\end{tabular} & $D_4$ & $D_6$ &  $\langle\alpha_1, \beta_1\rangle=1$, $\langle\alpha_1, \beta_2\rangle=1$, $\langle\alpha_1, \beta_3\rangle=1$, $\langle\alpha_1, \beta_4\rangle=1$, 
\\
\begin{tabular}{cc}$\alpha_1$+ & $\alpha_2$\\\tiny{ $A_1$ } & \tiny{ $A_1$ }\end{tabular} &\begin{tabular}{c} = \\~ \end{tabular} & \begin{tabular}{cc}$\beta_1$+ & $\beta_2$\\\tiny{ $A_1$+$A_1$'+$A_1$'' } & \tiny{ $A_1$+$A_1$'+$A_1$'' }\end{tabular} & $A_3$ & $D_5$ &  $\langle\alpha_1, \beta_1\rangle=1$, $\langle\alpha_1, \beta_2\rangle=1$, $\langle\alpha_2, \beta_1\rangle=1$, $\langle\alpha_2, \beta_2\rangle=1$, 
\\
\begin{tabular}{cc}$\alpha_1$+ & $\alpha_2$\\\tiny{ $A_1$+$A_1$' } & \tiny{ $A_1$' }\end{tabular} &\begin{tabular}{c} = \\~ \end{tabular} & \begin{tabular}{cc}$\beta_1$+ & $\beta_2$\\\tiny{ $A_1$+$A_1$'+$A_1$'' } & \tiny{ $A_1$+$A_1$'' }\end{tabular} & $A_3$ & $E_6$ &  $\langle\alpha_1, \beta_1\rangle=1$, $\langle\alpha_1, \beta_2\rangle=1$, $\langle\alpha_2, \beta_1\rangle=1$, $\langle\alpha_2, \beta_2\rangle=1$, 
\\
\begin{tabular}{cc}$\alpha_1$+ & $\alpha_2$\\\tiny{ $A_1$+$A_1$' } & \tiny{ $A_1$'' }\end{tabular} &\begin{tabular}{c} = \\~ \end{tabular} & \begin{tabular}{c}$\beta_1$\\\tiny{ $A_1$+$A_1$'+$A_1$'' }\end{tabular} & $A_2$ & $D_5$ & $\langle\alpha_1, \alpha_2\rangle=-1$,  $\langle\alpha_1, \beta_1\rangle=1$, $\langle\alpha_2, \beta_1\rangle=1$, 
\\
\begin{tabular}{cc}$\alpha_1$+ & $\alpha_2$\\\tiny{ $A_1$ } & \tiny{ $A_1$' }\end{tabular} &\begin{tabular}{c} = \\~ \end{tabular} & \begin{tabular}{ccc}$\beta_1$+ & $\beta_2$+ & $\beta_3$\\\tiny{ $A_1$+$A_1$' } & \tiny{ $A_1$''+$A_1$ } & \tiny{  }\end{tabular} & $D_4$ & $E_6$ & $\langle\alpha_1, \alpha_2\rangle=1$,  $\langle\alpha_1, \beta_1\rangle=1$, $\langle\alpha_1, \beta_2\rangle=1$, $\langle\alpha_1, \beta_3\rangle=1$, $\langle\alpha_2, \beta_1\rangle=1$, $\langle\alpha_2, \beta_2\rangle=1$, $\langle\alpha_2, \beta_3\rangle=1$, 
\\
\begin{tabular}{cc}$\alpha_1$+ & $\alpha_2$\\\tiny{ $A_1$ } & \tiny{  }\end{tabular} &\begin{tabular}{c} = \\~ \end{tabular} & \begin{tabular}{ccc}$\beta_1$+ & $\beta_2$+ & $\beta_3$\\\tiny{ $A_1$+$A_1$' } & \tiny{ $A_1$+$A_1$'' } & \tiny{ $A_1$'' }\end{tabular} & $D_4$ & $E_7$ & $\langle\alpha_1, \alpha_2\rangle=1$,  $\langle\alpha_1, \beta_1\rangle=1$, $\langle\alpha_1, \beta_2\rangle=1$, $\langle\alpha_1, \beta_3\rangle=1$, $\langle\alpha_2, \beta_1\rangle=1$, $\langle\alpha_2, \beta_2\rangle=1$, $\langle\alpha_2, \beta_3\rangle=1$, 
\\
\begin{tabular}{cc}$\alpha_1$+ & $\alpha_2$\\\tiny{ $A_1$+$A_1$' } & \tiny{ $A_1$ }\end{tabular} &\begin{tabular}{c} = \\~ \end{tabular} & \begin{tabular}{cc}$\beta_1$+ & $\beta_2$\\\tiny{ $A_1$+$A_1$'+$A_1$'' } & \tiny{ $A_1$+$A_1$'' }\end{tabular} & $A_3$ & $E_6$ &  $\langle\alpha_1, \beta_1\rangle=1$, $\langle\alpha_1, \beta_2\rangle=1$, $\langle\alpha_2, \beta_1\rangle=1$, $\langle\alpha_2, \beta_2\rangle=1$, 
\\
\begin{tabular}{cc}$\alpha_1$+ & $\alpha_2$\\\tiny{ $A_1$+$A_1$' } & \tiny{ $A_1$''+$A_1$ }\end{tabular} &\begin{tabular}{c} = \\~ \end{tabular} & \begin{tabular}{cc}$\beta_1$+ & $\beta_2$\\\tiny{ $A_1$+$A_1$'+$A_1$'' } & \tiny{ $A_1$' }\end{tabular} & $A_3$ & $E_6$ &  $\langle\alpha_1, \beta_1\rangle=1$, $\langle\alpha_1, \beta_2\rangle=1$, $\langle\alpha_2, \beta_1\rangle=1$, $\langle\alpha_2, \beta_2\rangle=1$, 
\\
\begin{tabular}{c}2$\alpha_1$\\\tiny{ $A_1$+$A_1$'+$A_1$'' }\end{tabular} &\begin{tabular}{c} = \\~ \end{tabular} & \begin{tabular}{cccc}$\beta_1$+ & $\beta_2$+ & $\beta_3$+ & $\beta_4$\\\tiny{ $A_1$+$A_1$' } & \tiny{ $A_1$+$A_1$'' } & \tiny{ $A_1$' } & \tiny{ $A_1$'' }\end{tabular} & $D_4$ & $E_7$ &  $\langle\alpha_1, \beta_1\rangle=1$, $\langle\alpha_1, \beta_2\rangle=1$, $\langle\alpha_1, \beta_3\rangle=1$, $\langle\alpha_1, \beta_4\rangle=1$, 
\\
\begin{tabular}{cc}$\alpha_1$+ & $\alpha_2$\\\tiny{ $A_1$ } & \tiny{ $A_1$' }\end{tabular} &\begin{tabular}{c} = \\~ \end{tabular} & \begin{tabular}{cc}$\beta_1$+ & $\beta_2$\\\tiny{ $A_1$+$A_1$'+$A_1$'' } & \tiny{ $A_1$ }\end{tabular} & $A_3$ & $E_6$ &  $\langle\alpha_1, \beta_1\rangle=1$, $\langle\alpha_1, \beta_2\rangle=1$, $\langle\alpha_2, \beta_1\rangle=1$, $\langle\alpha_2, \beta_2\rangle=1$, 
\\
\begin{tabular}{cc}$\alpha_1$+ & $\alpha_2$\\\tiny{ $A_1$ } & \tiny{ $A_1$' }\end{tabular} &\begin{tabular}{c} = \\~ \end{tabular} & \begin{tabular}{cc}$\beta_1$+ & $\beta_2$\\\tiny{ $A_1$+$A_1$'+$A_1$'' } & \tiny{ $A_1$' }\end{tabular} & $A_3$ & $E_6$ &  $\langle\alpha_1, \beta_1\rangle=1$, $\langle\alpha_1, \beta_2\rangle=1$, $\langle\alpha_2, \beta_1\rangle=1$, $\langle\alpha_2, \beta_2\rangle=1$, 
\\
\begin{tabular}{cc}$\alpha_1$+ & $\alpha_2$\\\tiny{ $A_1$ } & \tiny{  }\end{tabular} &\begin{tabular}{c} = \\~ \end{tabular} & \begin{tabular}{ccc}$\beta_1$+ & $\beta_2$+ & $\beta_3$\\\tiny{ $A_1$+$A_1$' } & \tiny{ $A_1$''+$A_1$' } & \tiny{ $A_1$'' }\end{tabular} & $D_4$ & $E_7$ & $\langle\alpha_1, \alpha_2\rangle=1$,  $\langle\alpha_1, \beta_1\rangle=1$, $\langle\alpha_1, \beta_2\rangle=1$, $\langle\alpha_1, \beta_3\rangle=1$, $\langle\alpha_2, \beta_1\rangle=1$, $\langle\alpha_2, \beta_2\rangle=1$, $\langle\alpha_2, \beta_3\rangle=1$, 
\\
\begin{tabular}{cc}$\alpha_1$+ & $\alpha_2$\\\tiny{ $A_1$+$A_1$' } & \tiny{ $A_1$ }\end{tabular} &\begin{tabular}{c} = \\~ \end{tabular} & \begin{tabular}{cc}$\beta_1$+ & $\beta_2$\\\tiny{ $A_1$+$A_1$'+$A_1$'' } & \tiny{ $A_1$+$A_1$' }\end{tabular} & $A_3$ & $E_6$ &  $\langle\alpha_1, \beta_1\rangle=1$, $\langle\alpha_1, \beta_2\rangle=1$, $\langle\alpha_2, \beta_1\rangle=1$, $\langle\alpha_2, \beta_2\rangle=1$, 
\\
\begin{tabular}{cc}$\alpha_1$+ & $\alpha_2$\\\tiny{ $A_1$ } & \tiny{ $A_1$' }\end{tabular} &\begin{tabular}{c} = \\~ \end{tabular} & \begin{tabular}{cc}$\beta_1$+ & $\beta_2$\\\tiny{ $A_1$+$A_1$'+$A_1$'' } & \tiny{ $A_1$'' }\end{tabular} & $A_3$ & $E_6$ &  $\langle\alpha_1, \beta_1\rangle=1$, $\langle\alpha_1, \beta_2\rangle=1$, $\langle\alpha_2, \beta_1\rangle=1$, $\langle\alpha_2, \beta_2\rangle=1$, 
\\
\begin{tabular}{cc}$\alpha_1$+ & $\alpha_2$\\\tiny{ $A_1$+$A_1$'+$A_1$'' } & \tiny{ $A_1$' }\end{tabular} &\begin{tabular}{c} = \\~ \end{tabular} & \begin{tabular}{ccc}$\beta_1$+ & $\beta_2$+ & $\beta_3$\\\tiny{ $A_1$+$A_1$' } & \tiny{ $A_1$'' } & \tiny{ $A_1$' }\end{tabular} & $D_4$ & $E_7$ & $\langle\alpha_1, \alpha_2\rangle=1$,  $\langle\alpha_1, \beta_1\rangle=1$, $\langle\alpha_1, \beta_2\rangle=1$, $\langle\alpha_1, \beta_3\rangle=1$, $\langle\alpha_2, \beta_1\rangle=1$, $\langle\alpha_2, \beta_2\rangle=1$, $\langle\alpha_2, \beta_3\rangle=1$, 
\\
\begin{tabular}{cc}$\alpha_1$+ & $\alpha_2$\\\tiny{ $A_1$ } & \tiny{ $A_1$ }\end{tabular} &\begin{tabular}{c} = \\~ \end{tabular} & \begin{tabular}{ccc}$\beta_1$+ & $\beta_2$+ & $\beta_3$\\\tiny{ $A_1$+$A_1$'+$A_1$'' } & \tiny{ $A_1$+$A_1$' } & \tiny{ $A_1$'' }\end{tabular} & $D_4$ & $E_7$ & $\langle\alpha_1, \alpha_2\rangle=1$,  $\langle\alpha_1, \beta_1\rangle=1$, $\langle\alpha_1, \beta_2\rangle=1$, $\langle\alpha_1, \beta_3\rangle=1$, $\langle\alpha_2, \beta_1\rangle=1$, $\langle\alpha_2, \beta_2\rangle=1$, $\langle\alpha_2, \beta_3\rangle=1$, 
\\
\begin{tabular}{cc}$\alpha_1$+ & $\alpha_2$\\\tiny{ $A_1$ } & \tiny{ $A_1$' }\end{tabular} &\begin{tabular}{c} = \\~ \end{tabular} & \begin{tabular}{ccc}$\beta_1$+ & $\beta_2$+ & $\beta_3$\\\tiny{ $A_1$+$A_1$' } & \tiny{ $A_1$+$A_1$'' } & \tiny{  }\end{tabular} & $D_4$ & $E_6$ & $\langle\alpha_1, \alpha_2\rangle=1$,  $\langle\alpha_1, \beta_1\rangle=1$, $\langle\alpha_1, \beta_2\rangle=1$, $\langle\alpha_1, \beta_3\rangle=1$, $\langle\alpha_2, \beta_1\rangle=1$, $\langle\alpha_2, \beta_2\rangle=1$, $\langle\alpha_2, \beta_3\rangle=1$, 
\\
\begin{tabular}{cc}$\alpha_1$+ & $\alpha_2$\\\tiny{ $A_1$ } & \tiny{ $A_1$' }\end{tabular} &\begin{tabular}{c} = \\~ \end{tabular} & \begin{tabular}{ccc}$\beta_1$+ & $\beta_2$+ & $\beta_3$\\\tiny{ $A_1$+$A_1$' } & \tiny{ $A_1$''+$A_1$' } & \tiny{  }\end{tabular} & $D_4$ & $E_6$ & $\langle\alpha_1, \alpha_2\rangle=1$,  $\langle\alpha_1, \beta_1\rangle=1$, $\langle\alpha_1, \beta_2\rangle=1$, $\langle\alpha_1, \beta_3\rangle=1$, $\langle\alpha_2, \beta_1\rangle=1$, $\langle\alpha_2, \beta_2\rangle=1$, $\langle\alpha_2, \beta_3\rangle=1$, 
\\
\begin{tabular}{cc}$\alpha_1$+ & $\alpha_2$\\\tiny{ $A_1$ } & \tiny{ $A_1$ }\end{tabular} &\begin{tabular}{c} = \\~ \end{tabular} & \begin{tabular}{ccc}$\beta_1$+ & $\beta_2$+ & $\beta_3$\\\tiny{ $A_1$+$A_1$'+$A_1$'' } & \tiny{ $A_1$'+$A_1$'' } & \tiny{ $A_1$ }\end{tabular} & $D_4$ & $E_7$ & $\langle\alpha_1, \alpha_2\rangle=1$,  $\langle\alpha_1, \beta_1\rangle=1$, $\langle\alpha_1, \beta_2\rangle=1$, $\langle\alpha_1, \beta_3\rangle=1$, $\langle\alpha_2, \beta_1\rangle=1$, $\langle\alpha_2, \beta_2\rangle=1$, $\langle\alpha_2, \beta_3\rangle=1$, 
\\
\begin{tabular}{cc}$\alpha_1$+ & $\alpha_2$\\\tiny{ $A_1$+$A_1$' } & \tiny{ $A_1$''+$A_1$' }\end{tabular} &\begin{tabular}{c} = \\~ \end{tabular} & \begin{tabular}{cc}$\beta_1$+ & $\beta_2$\\\tiny{ $A_1$+$A_1$'+$A_1$'' } & \tiny{ $A_1$'' }\end{tabular} & $A_3$ & $E_6$ &  $\langle\alpha_1, \beta_1\rangle=1$, $\langle\alpha_1, \beta_2\rangle=1$, $\langle\alpha_2, \beta_1\rangle=1$, $\langle\alpha_2, \beta_2\rangle=1$, 
\\
\begin{tabular}{cc}$\alpha_1$+ & $\alpha_2$\\\tiny{ $A_1$+$A_1$' } & \tiny{ $A_1$'' }\end{tabular} &\begin{tabular}{c} = \\~ \end{tabular} & \begin{tabular}{cc}$\beta_1$+ & $\beta_2$\\\tiny{ $A_1$+$A_1$'+$A_1$'' } & \tiny{  }\end{tabular} & $A_3$ & $D_6$ &  $\langle\alpha_1, \beta_1\rangle=1$, $\langle\alpha_1, \beta_2\rangle=1$, $\langle\alpha_2, \beta_1\rangle=1$, $\langle\alpha_2, \beta_2\rangle=1$, 
\\
\begin{tabular}{cc}$\alpha_1$+ & $\alpha_2$\\\tiny{ $A_1$ } & \tiny{ $A_1$ }\end{tabular} &\begin{tabular}{c} = \\~ \end{tabular} & \begin{tabular}{ccc}$\beta_1$+ & $\beta_2$+ & $\beta_3$\\\tiny{ $A_1$+$A_1$'+$A_1$'' } & \tiny{ $A_1$+$A_1$'+$A_1$'' } & \tiny{  }\end{tabular} & $D_4$ & $D_6$ & $\langle\alpha_1, \alpha_2\rangle=1$,  $\langle\alpha_1, \beta_1\rangle=1$, $\langle\alpha_1, \beta_2\rangle=1$, $\langle\alpha_1, \beta_3\rangle=1$, $\langle\alpha_2, \beta_1\rangle=1$, $\langle\alpha_2, \beta_2\rangle=1$, $\langle\alpha_2, \beta_3\rangle=1$, 
\\
\begin{tabular}{cc}$\alpha_1$+ & $\alpha_2$\\\tiny{ $A_1$+$A_1$' } & \tiny{ $A_1$+$A_1$' }\end{tabular} &\begin{tabular}{c} = \\~ \end{tabular} & \begin{tabular}{ccc}$\beta_1$+ & $\beta_2$+ & $\beta_3$\\\tiny{ $A_1$+$A_1$'+$A_1$'' } & \tiny{ $A_1$+$A_1$'' } & \tiny{ $A_1$' }\end{tabular} & $D_4$ & $E_7$ & $\langle\alpha_1, \alpha_2\rangle=1$,  $\langle\alpha_1, \beta_1\rangle=1$, $\langle\alpha_1, \beta_2\rangle=1$, $\langle\alpha_1, \beta_3\rangle=1$, $\langle\alpha_2, \beta_1\rangle=1$, $\langle\alpha_2, \beta_2\rangle=1$, $\langle\alpha_2, \beta_3\rangle=1$, 
\\
\begin{tabular}{c}2$\alpha_1$\\\tiny{  }\end{tabular} &\begin{tabular}{c} = \\~ \end{tabular} & \begin{tabular}{cccc}$\beta_1$+ & $\beta_2$+ & $\beta_3$+ & $\beta_4$\\\tiny{ $A_1$+$A_1$' } & \tiny{ $A_1$''+$A_1$ } & \tiny{ $A_1$' } & \tiny{ $A_1$'' }\end{tabular} & $D_4$ & $E_7$ &  $\langle\alpha_1, \beta_1\rangle=1$, $\langle\alpha_1, \beta_2\rangle=1$, $\langle\alpha_1, \beta_3\rangle=1$, $\langle\alpha_1, \beta_4\rangle=1$, 
\\
\begin{tabular}{c}2$\alpha_1$\\\tiny{ $A_1$+$A_1$'+$A_1$'' }\end{tabular} &\begin{tabular}{c} = \\~ \end{tabular} & \begin{tabular}{cccc}$\beta_1$+ & $\beta_2$+ & $\beta_3$+ & $\beta_4$\\\tiny{ $A_1$+$A_1$' } & \tiny{ $A_1$''+$A_1$ } & \tiny{ $A_1$' } & \tiny{ $A_1$'' }\end{tabular} & $D_4$ & $E_7$ &  $\langle\alpha_1, \beta_1\rangle=1$, $\langle\alpha_1, \beta_2\rangle=1$, $\langle\alpha_1, \beta_3\rangle=1$, $\langle\alpha_1, \beta_4\rangle=1$, 
\\
\begin{tabular}{cc}$\alpha_1$+ & $\alpha_2$\\\tiny{ $A_1$+$A_1$' } & \tiny{ $A_1$' }\end{tabular} &\begin{tabular}{c} = \\~ \end{tabular} & \begin{tabular}{cc}$\beta_1$+ & $\beta_2$\\\tiny{ $A_1$+$A_1$'+$A_1$'' } & \tiny{ $A_1$'+$A_1$'' }\end{tabular} & $A_3$ & $E_6$ &  $\langle\alpha_1, \beta_1\rangle=1$, $\langle\alpha_1, \beta_2\rangle=1$, $\langle\alpha_2, \beta_1\rangle=1$, $\langle\alpha_2, \beta_2\rangle=1$, 
\\
\begin{tabular}{cc}$\alpha_1$+ & $\alpha_2$\\\tiny{ $A_1$+$A_1$'+$A_1$'' } & \tiny{ $A_1$' }\end{tabular} &\begin{tabular}{c} = \\~ \end{tabular} & \begin{tabular}{ccc}$\beta_1$+ & $\beta_2$+ & $\beta_3$\\\tiny{ $A_1$+$A_1$' } & \tiny{ $A_1$ } & \tiny{ $A_1$'' }\end{tabular} & $D_4$ & $E_7$ & $\langle\alpha_1, \alpha_2\rangle=1$,  $\langle\alpha_1, \beta_1\rangle=1$, $\langle\alpha_1, \beta_2\rangle=1$, $\langle\alpha_1, \beta_3\rangle=1$, $\langle\alpha_2, \beta_1\rangle=1$, $\langle\alpha_2, \beta_2\rangle=1$, $\langle\alpha_2, \beta_3\rangle=1$, 
\\
\begin{tabular}{cc}$\alpha_1$+ & $\alpha_2$\\\tiny{ $A_1$ } & \tiny{ $A_1$' }\end{tabular} &\begin{tabular}{c} = \\~ \end{tabular} & \begin{tabular}{ccc}$\beta_1$+ & $\beta_2$+ & $\beta_3$\\\tiny{ $A_1$+$A_1$' } & \tiny{  } & \tiny{  }\end{tabular} & $D_4$ & $D_6$+$A_1$ & $\langle\alpha_1, \alpha_2\rangle=1$,  $\langle\alpha_1, \beta_1\rangle=1$, $\langle\alpha_1, \beta_2\rangle=1$, $\langle\alpha_1, \beta_3\rangle=1$, $\langle\alpha_2, \beta_1\rangle=1$, $\langle\alpha_2, \beta_2\rangle=1$, $\langle\alpha_2, \beta_3\rangle=1$, 
\\
\begin{tabular}{c}2$\alpha_1$\\\tiny{ $A_1$ }\end{tabular} &\begin{tabular}{c} = \\~ \end{tabular} & \begin{tabular}{cccc}$\beta_1$+ & $\beta_2$+ & $\beta_3$+ & $\beta_4$\\\tiny{ $A_1$+$A_1$'+$A_1$'' } & \tiny{ $A_1$'+$A_1$'' } & \tiny{ $A_1$ } & \tiny{  }\end{tabular} & $D_4$ & $E_7$ &  $\langle\alpha_1, \beta_1\rangle=1$, $\langle\alpha_1, \beta_2\rangle=1$, $\langle\alpha_1, \beta_3\rangle=1$, $\langle\alpha_1, \beta_4\rangle=1$, 
\\
\begin{tabular}{cc}$\alpha_1$+ & $\alpha_2$\\\tiny{ $A_1$+$A_1$' } & \tiny{ $A_1$' }\end{tabular} &\begin{tabular}{c} = \\~ \end{tabular} & \begin{tabular}{ccc}$\beta_1$+ & $\beta_2$+ & $\beta_3$\\\tiny{ $A_1$+$A_1$'+$A_1$'' } & \tiny{ $A_1$'+$A_1$'' } & \tiny{  }\end{tabular} & $D_4$ & $E_7$ & $\langle\alpha_1, \alpha_2\rangle=1$,  $\langle\alpha_1, \beta_1\rangle=1$, $\langle\alpha_1, \beta_2\rangle=1$, $\langle\alpha_1, \beta_3\rangle=1$, $\langle\alpha_2, \beta_1\rangle=1$, $\langle\alpha_2, \beta_2\rangle=1$, $\langle\alpha_2, \beta_3\rangle=1$, 
\\
\begin{tabular}{cc}$\alpha_1$+ & $\alpha_2$\\\tiny{ $A_1$+$A_1$'+$A_1$'' } & \tiny{ $A_1$'' }\end{tabular} &\begin{tabular}{c} = \\~ \end{tabular} & \begin{tabular}{ccc}$\beta_1$+ & $\beta_2$+ & $\beta_3$\\\tiny{ $A_1$+$A_1$' } & \tiny{ $A_1$' } & \tiny{ $A_1$'' }\end{tabular} & $D_4$ & $E_7$ & $\langle\alpha_1, \alpha_2\rangle=1$,  $\langle\alpha_1, \beta_1\rangle=1$, $\langle\alpha_1, \beta_2\rangle=1$, $\langle\alpha_1, \beta_3\rangle=1$, $\langle\alpha_2, \beta_1\rangle=1$, $\langle\alpha_2, \beta_2\rangle=1$, $\langle\alpha_2, \beta_3\rangle=1$, 
\\
\begin{tabular}{cc}$\alpha_1$+ & $\alpha_2$\\\tiny{ $A_1$ } & \tiny{ $A_1$' }\end{tabular} &\begin{tabular}{c} = \\~ \end{tabular} & \begin{tabular}{ccc}$\beta_1$+ & $\beta_2$+ & $\beta_3$\\\tiny{ $A_1$+$A_1$'+$A_1$'' } & \tiny{ $A_1$ } & \tiny{  }\end{tabular} & $D_4$ & $E_7$ & $\langle\alpha_1, \alpha_2\rangle=1$,  $\langle\alpha_1, \beta_1\rangle=1$, $\langle\alpha_1, \beta_2\rangle=1$, $\langle\alpha_1, \beta_3\rangle=1$, $\langle\alpha_2, \beta_1\rangle=1$, $\langle\alpha_2, \beta_2\rangle=1$, $\langle\alpha_2, \beta_3\rangle=1$, 
\\
\begin{tabular}{cc}$\alpha_1$+ & $\alpha_2$\\\tiny{ $A_1$+$A_1$' } & \tiny{ $A_1$+$A_1$'' }\end{tabular} &\begin{tabular}{c} = \\~ \end{tabular} & \begin{tabular}{ccc}$\beta_1$+ & $\beta_2$+ & $\beta_3$\\\tiny{ $A_1$+$A_1$'+$A_1$'' } & \tiny{ $A_1$ } & \tiny{  }\end{tabular} & $D_4$ & $E_7$ & $\langle\alpha_1, \alpha_2\rangle=1$,  $\langle\alpha_1, \beta_1\rangle=1$, $\langle\alpha_1, \beta_2\rangle=1$, $\langle\alpha_1, \beta_3\rangle=1$, $\langle\alpha_2, \beta_1\rangle=1$, $\langle\alpha_2, \beta_2\rangle=1$, $\langle\alpha_2, \beta_3\rangle=1$, 
\\
\begin{tabular}{c}2$\alpha_1$\\\tiny{ $A_1$+$A_1$' }\end{tabular} &\begin{tabular}{c} = \\~ \end{tabular} & \begin{tabular}{cccc}$\beta_1$+ & $\beta_2$+ & $\beta_3$+ & $\beta_4$\\\tiny{ $A_1$+$A_1$'+$A_1$'' } & \tiny{ $A_1$+$A_1$'' } & \tiny{ $A_1$' } & \tiny{  }\end{tabular} & $D_4$ & $E_7$ &  $\langle\alpha_1, \beta_1\rangle=1$, $\langle\alpha_1, \beta_2\rangle=1$, $\langle\alpha_1, \beta_3\rangle=1$, $\langle\alpha_1, \beta_4\rangle=1$, 
\\
\begin{tabular}{cc}$\alpha_1$+ & $\alpha_2$\\\tiny{ $A_1$+$A_1$' } & \tiny{ $A_1$''+$A_1$ }\end{tabular} &\begin{tabular}{c} = \\~ \end{tabular} & \begin{tabular}{ccc}$\beta_1$+ & $\beta_2$+ & $\beta_3$\\\tiny{ $A_1$+$A_1$'+$A_1$'' } & \tiny{ $A_1$' } & \tiny{  }\end{tabular} & $D_4$ & $E_7$ & $\langle\alpha_1, \alpha_2\rangle=1$,  $\langle\alpha_1, \beta_1\rangle=1$, $\langle\alpha_1, \beta_2\rangle=1$, $\langle\alpha_1, \beta_3\rangle=1$, $\langle\alpha_2, \beta_1\rangle=1$, $\langle\alpha_2, \beta_2\rangle=1$, $\langle\alpha_2, \beta_3\rangle=1$, 
\\
\begin{tabular}{c}2$\alpha_1$\\\tiny{ $A_1$ }\end{tabular} &\begin{tabular}{c} = \\~ \end{tabular} & \begin{tabular}{cccc}$\beta_1$+ & $\beta_2$+ & $\beta_3$+ & $\beta_4$\\\tiny{ $A_1$+$A_1$'+$A_1$'' } & \tiny{ $A_1$+$A_1$'+$A_1$'' } & \tiny{  } & \tiny{  }\end{tabular} & $D_4$ & $D_6$ &  $\langle\alpha_1, \beta_1\rangle=1$, $\langle\alpha_1, \beta_2\rangle=1$, $\langle\alpha_1, \beta_3\rangle=1$, $\langle\alpha_1, \beta_4\rangle=1$, 
\\
\begin{tabular}{cc}$\alpha_1$+ & $\alpha_2$\\\tiny{ $A_1$+$A_1$' } & \tiny{ $A_1$' }\end{tabular} &\begin{tabular}{c} = \\~ \end{tabular} & \begin{tabular}{ccc}$\beta_1$+ & $\beta_2$+ & $\beta_3$\\\tiny{ $A_1$+$A_1$'+$A_1$'' } & \tiny{ $A_1$+$A_1$'' } & \tiny{  }\end{tabular} & $D_4$ & $E_7$ & $\langle\alpha_1, \alpha_2\rangle=1$,  $\langle\alpha_1, \beta_1\rangle=1$, $\langle\alpha_1, \beta_2\rangle=1$, $\langle\alpha_1, \beta_3\rangle=1$, $\langle\alpha_2, \beta_1\rangle=1$, $\langle\alpha_2, \beta_2\rangle=1$, $\langle\alpha_2, \beta_3\rangle=1$, 
\\
\begin{tabular}{cc}$\alpha_1$+ & $\alpha_2$\\\tiny{ $A_1$+$A_1$' } & \tiny{ $A_1$+$A_1$'' }\end{tabular} &\begin{tabular}{c} = \\~ \end{tabular} & \begin{tabular}{cc}$\beta_1$+ & $\beta_2$\\\tiny{ $A_1$+$A_1$'+$A_1$'' } & \tiny{ $A_1$ }\end{tabular} & $A_3$ & $E_6$ &  $\langle\alpha_1, \beta_1\rangle=1$, $\langle\alpha_1, \beta_2\rangle=1$, $\langle\alpha_2, \beta_1\rangle=1$, $\langle\alpha_2, \beta_2\rangle=1$, 
\\
\begin{tabular}{c}2$\alpha_1$\\\tiny{ $A_1$ }\end{tabular} &\begin{tabular}{c} = \\~ \end{tabular} & \begin{tabular}{cccc}$\beta_1$+ & $\beta_2$+ & $\beta_3$+ & $\beta_4$\\\tiny{ $A_1$+$A_1$' } & \tiny{ $A_1$+$A_1$' } & \tiny{  } & \tiny{  }\end{tabular} & $D_4$ & $D_6$+$A_1$ &  $\langle\alpha_1, \beta_1\rangle=1$, $\langle\alpha_1, \beta_2\rangle=1$, $\langle\alpha_1, \beta_3\rangle=1$, $\langle\alpha_1, \beta_4\rangle=1$, 
\\
\begin{tabular}{cc}$\alpha_1$+ & $\alpha_2$\\\tiny{ $A_1$+$A_1$' } & \tiny{ $A_1$ }\end{tabular} &\begin{tabular}{c} = \\~ \end{tabular} & \begin{tabular}{ccc}$\beta_1$+ & $\beta_2$+ & $\beta_3$\\\tiny{ $A_1$+$A_1$'+$A_1$'' } & \tiny{ $A_1$+$A_1$'' } & \tiny{  }\end{tabular} & $D_4$ & $E_7$ & $\langle\alpha_1, \alpha_2\rangle=1$,  $\langle\alpha_1, \beta_1\rangle=1$, $\langle\alpha_1, \beta_2\rangle=1$, $\langle\alpha_1, \beta_3\rangle=1$, $\langle\alpha_2, \beta_1\rangle=1$, $\langle\alpha_2, \beta_2\rangle=1$, $\langle\alpha_2, \beta_3\rangle=1$, 
\\
\begin{tabular}{c}2$\alpha_1$\\\tiny{ $A_1$+$A_1$'+$A_1$'' }\end{tabular} &\begin{tabular}{c} = \\~ \end{tabular} & \begin{tabular}{cccc}$\beta_1$+ & $\beta_2$+ & $\beta_3$+ & $\beta_4$\\\tiny{ $A_1$+$A_1$' } & \tiny{ $A_1$''+$A_1$' } & \tiny{ $A_1$ } & \tiny{ $A_1$'' }\end{tabular} & $D_4$ & $E_7$ &  $\langle\alpha_1, \beta_1\rangle=1$, $\langle\alpha_1, \beta_2\rangle=1$, $\langle\alpha_1, \beta_3\rangle=1$, $\langle\alpha_1, \beta_4\rangle=1$, 
\\
\begin{tabular}{cc}$\alpha_1$+ & $\alpha_2$\\\tiny{ $A_1$+$A_1$'+$A_1$'' } & \tiny{ $A_1$+$A_1$'' }\end{tabular} &\begin{tabular}{c} = \\~ \end{tabular} & \begin{tabular}{ccc}$\beta_1$+ & $\beta_2$+ & $\beta_3$\\\tiny{ $A_1$+$A_1$' } & \tiny{ $A_1$''+$A_1$' } & \tiny{ $A_1$'' }\end{tabular} & $D_4$ & $E_7$ & $\langle\alpha_1, \alpha_2\rangle=1$,  $\langle\alpha_1, \beta_1\rangle=1$, $\langle\alpha_1, \beta_2\rangle=1$, $\langle\alpha_1, \beta_3\rangle=1$, $\langle\alpha_2, \beta_1\rangle=1$, $\langle\alpha_2, \beta_2\rangle=1$, $\langle\alpha_2, \beta_3\rangle=1$, 
\\
\begin{tabular}{c}2$\alpha_1$\\\tiny{ $A_1$ }\end{tabular} &\begin{tabular}{c} = \\~ \end{tabular} & \begin{tabular}{cccc}$\beta_1$+ & $\beta_2$+ & $\beta_3$+ & $\beta_4$\\\tiny{ $A_1$+$A_1$'+$A_1$'' } & \tiny{ $A_1$+$A_1$'' } & \tiny{ $A_1$' } & \tiny{  }\end{tabular} & $D_4$ & $E_7$ &  $\langle\alpha_1, \beta_1\rangle=1$, $\langle\alpha_1, \beta_2\rangle=1$, $\langle\alpha_1, \beta_3\rangle=1$, $\langle\alpha_1, \beta_4\rangle=1$, 
\\
\begin{tabular}{cc}$\alpha_1$+ & $\alpha_2$\\\tiny{ $A_1$ } & \tiny{ $A_1$' }\end{tabular} &\begin{tabular}{c} = \\~ \end{tabular} & \begin{tabular}{ccc}$\beta_1$+ & $\beta_2$+ & $\beta_3$\\\tiny{ $A_1$+$A_1$'+$A_1$'' } & \tiny{ $A_1$' } & \tiny{  }\end{tabular} & $D_4$ & $E_7$ & $\langle\alpha_1, \alpha_2\rangle=1$,  $\langle\alpha_1, \beta_1\rangle=1$, $\langle\alpha_1, \beta_2\rangle=1$, $\langle\alpha_1, \beta_3\rangle=1$, $\langle\alpha_2, \beta_1\rangle=1$, $\langle\alpha_2, \beta_2\rangle=1$, $\langle\alpha_2, \beta_3\rangle=1$, 
\\
\begin{tabular}{cc}$\alpha_1$+ & $\alpha_2$\\\tiny{ $A_1$+$A_1$' } & \tiny{ $A_1$+$A_1$' }\end{tabular} &\begin{tabular}{c} = \\~ \end{tabular} & \begin{tabular}{ccc}$\beta_1$+ & $\beta_2$+ & $\beta_3$\\\tiny{ $A_1$+$A_1$'+$A_1$'' } & \tiny{ $A_1$+$A_1$' } & \tiny{ $A_1$'' }\end{tabular} & $D_4$ & $E_7$ & $\langle\alpha_1, \alpha_2\rangle=1$,  $\langle\alpha_1, \beta_1\rangle=1$, $\langle\alpha_1, \beta_2\rangle=1$, $\langle\alpha_1, \beta_3\rangle=1$, $\langle\alpha_2, \beta_1\rangle=1$, $\langle\alpha_2, \beta_2\rangle=1$, $\langle\alpha_2, \beta_3\rangle=1$, 
\\
\begin{tabular}{ccc}$\alpha_1$+ & $\alpha_2$+ & $\alpha_3$\\\tiny{ $A_1$ } & \tiny{ $A_1$' } & \tiny{ $A_1$'' }\end{tabular} &\begin{tabular}{c} = \\~ \end{tabular} & \begin{tabular}{ccc}$\beta_1$+ & $\beta_2$+ & $\beta_3$\\\tiny{ $A_1$ } & \tiny{ $A_1$' } & \tiny{ $A_1$'' }\end{tabular} & $A_5$ & $E_6$ &  $\langle\alpha_1, \beta_2\rangle=1$, $\langle\alpha_1, \beta_3\rangle=1$, $\langle\alpha_2, \beta_1\rangle=1$, $\langle\alpha_2, \beta_3\rangle=1$, $\langle\alpha_3, \beta_1\rangle=1$, $\langle\alpha_3, \beta_2\rangle=1$, 
\\
\hline\multicolumn{5}{c}{$\mathfrak{k}$-semisimple type: $A_1$+$A_1$+$A_1$}\\
\hline\hline\begin{tabular}{cc}$\alpha_1$+ & $\alpha_2$\\\tiny{ $A_1$+$A_1$' } & \tiny{ $A_1$''+$A_1$ }\end{tabular} &\begin{tabular}{c} = \\~ \end{tabular} & \begin{tabular}{cc}$\beta_1$+ & $\beta_2$\\\tiny{ $A_1$+$A_1$' } & \tiny{ $A_1$''+$A_1$ }\end{tabular} & $A_3$ & $D_6$ &  $\langle\alpha_1, \beta_1\rangle=1$, $\langle\alpha_1, \beta_2\rangle=1$, $\langle\alpha_2, \beta_1\rangle=1$, $\langle\alpha_2, \beta_2\rangle=1$, 
\\
\begin{tabular}{cc}$\alpha_1$+ & $\alpha_2$\\\tiny{ $A_1$+$A_1$' } & \tiny{ $A_1$+$A_1$' }\end{tabular} &\begin{tabular}{c} = \\~ \end{tabular} & \begin{tabular}{ccc}$\beta_1$+ & $\beta_2$+ & $\beta_3$\\\tiny{ $A_1$+$A_1$' } & \tiny{ $A_1$''+$A_1$' } & \tiny{ $A_1$''+$A_1$ }\end{tabular} & $D_4$ & $E_7$ & $\langle\alpha_1, \alpha_2\rangle=1$,  $\langle\alpha_1, \beta_1\rangle=1$, $\langle\alpha_1, \beta_2\rangle=1$, $\langle\alpha_1, \beta_3\rangle=1$, $\langle\alpha_2, \beta_1\rangle=1$, $\langle\alpha_2, \beta_2\rangle=1$, $\langle\alpha_2, \beta_3\rangle=1$, 
\\
\begin{tabular}{cc}$\alpha_1$+ & $\alpha_2$\\\tiny{ $A_1$+$A_1$' } & \tiny{ $A_1$+$A_1$'' }\end{tabular} &\begin{tabular}{c} = \\~ \end{tabular} & \begin{tabular}{cc}$\beta_1$+ & $\beta_2$\\\tiny{ $A_1$+$A_1$' } & \tiny{ $A_1$+$A_1$'' }\end{tabular} & $A_3$ & $D_6$ &  $\langle\alpha_1, \beta_1\rangle=1$, $\langle\alpha_1, \beta_2\rangle=1$, $\langle\alpha_2, \beta_1\rangle=1$, $\langle\alpha_2, \beta_2\rangle=1$, 
\\
\end{longtable}

% % % % % % % % % % % % % % % % % % % % % % % % % % % %
% % % % % % % % % % % % % % % % % % % % % % % % % % % % % % % % %

\begin{longtable}{rcl p{1cm}p{1cm}p{3cm} } \multicolumn{3}{c}{ Relation / linked $\mathfrak{k}$-components} &$\alpha_i$'s, $\beta_i$'s generate & adding $\Delta(\mathfrak{k})$ generates&Non-zero scalar products\endhead
\\\hline\multicolumn{5}{c}{$\mathfrak{k}$-semisimple type: $A_2$}\\
\hline\hline\begin{tabular}{cc}$\alpha_1$+ & $\alpha_2$\\\tiny{  } & \tiny{  }\end{tabular} &\begin{tabular}{c} = \\~ \end{tabular} & \begin{tabular}{cc}$\beta_1$+ & $\beta_2$\\\tiny{ $A_2$ } & \tiny{ $A_2$ }\end{tabular} & $A_3$ & $D_5$ &  $\langle\alpha_1, \beta_1\rangle=1$, $\langle\alpha_1, \beta_2\rangle=1$, $\langle\alpha_2, \beta_1\rangle=1$, $\langle\alpha_2, \beta_2\rangle=1$, 
\\
\begin{tabular}{cc}$\alpha_1$+ & $\alpha_2$\\\tiny{ $A_2$ } & \tiny{  }\end{tabular} &\begin{tabular}{c} = \\~ \end{tabular} & \begin{tabular}{c}$\beta_1$\\\tiny{ $A_2$ }\end{tabular} & $A_2$ & $A_4$ & $\langle\alpha_1, \alpha_2\rangle=-1$,  $\langle\alpha_1, \beta_1\rangle=1$, $\langle\alpha_2, \beta_1\rangle=1$, 
\\
\begin{tabular}{cc}$\alpha_1$+ & $\alpha_2$\\\tiny{ $A_2$ } & \tiny{ $A_2$ }\end{tabular} &\begin{tabular}{c} = \\~ \end{tabular} & \begin{tabular}{cc}$\beta_1$+ & $\beta_2$\\\tiny{ $A_2$ } & \tiny{ $A_2$ }\end{tabular} & $A_3$ & $D_5$ &  $\langle\alpha_1, \beta_1\rangle=1$, $\langle\alpha_1, \beta_2\rangle=1$, $\langle\alpha_2, \beta_1\rangle=1$, $\langle\alpha_2, \beta_2\rangle=1$, 
\\
\begin{tabular}{cc}$\alpha_1$+ & $\alpha_2$\\\tiny{  } & \tiny{  }\end{tabular} &\begin{tabular}{c} = \\~ \end{tabular} & \begin{tabular}{ccc}$\beta_1$+ & $\beta_2$+ & $\beta_3$\\\tiny{ $A_2$ } & \tiny{ $A_2$ } & \tiny{ $A_2$ }\end{tabular} & $D_4$ & $E_6$ & $\langle\alpha_1, \alpha_2\rangle=1$,  $\langle\alpha_1, \beta_1\rangle=1$, $\langle\alpha_1, \beta_2\rangle=1$, $\langle\alpha_1, \beta_3\rangle=1$, $\langle\alpha_2, \beta_1\rangle=1$, $\langle\alpha_2, \beta_2\rangle=1$, $\langle\alpha_2, \beta_3\rangle=1$, 
\\
\begin{tabular}{cc}$\alpha_1$+ & $\alpha_2$\\\tiny{ $A_2$ } & \tiny{  }\end{tabular} &\begin{tabular}{c} = \\~ \end{tabular} & \begin{tabular}{cc}$\beta_1$+ & $\beta_2$\\\tiny{ $A_2$ } & \tiny{ $A_2$ }\end{tabular} & $A_3$ & $D_5$ &  $\langle\alpha_1, \beta_1\rangle=1$, $\langle\alpha_1, \beta_2\rangle=1$, $\langle\alpha_2, \beta_1\rangle=1$, $\langle\alpha_2, \beta_2\rangle=1$, 
\\
\begin{tabular}{cc}$\alpha_1$+ & $\alpha_2$\\\tiny{ $A_2$ } & \tiny{ $A_2$ }\end{tabular} &\begin{tabular}{c} = \\~ \end{tabular} & \begin{tabular}{ccc}$\beta_1$+ & $\beta_2$+ & $\beta_3$\\\tiny{ $A_2$ } & \tiny{ $A_2$ } & \tiny{ $A_2$ }\end{tabular} & $D_4$ & $E_6$ & $\langle\alpha_1, \alpha_2\rangle=1$,  $\langle\alpha_1, \beta_1\rangle=1$, $\langle\alpha_1, \beta_2\rangle=1$, $\langle\alpha_1, \beta_3\rangle=1$, $\langle\alpha_2, \beta_1\rangle=1$, $\langle\alpha_2, \beta_2\rangle=1$, $\langle\alpha_2, \beta_3\rangle=1$, 
\\
\hline\multicolumn{5}{c}{$\mathfrak{k}$-semisimple type: $A_1$+$A_1$}\\
\hline\hline\begin{tabular}{cc}$\alpha_1$+ & $\alpha_2$\\\tiny{  } & \tiny{  }\end{tabular} &\begin{tabular}{c} = \\~ \end{tabular} & \begin{tabular}{cc}$\beta_1$+ & $\beta_2$\\\tiny{ $A_1$ } & \tiny{ $A_1$ }\end{tabular} & $A_3$ & $D_4$+$A_1$ &  $\langle\alpha_1, \beta_1\rangle=1$, $\langle\alpha_1, \beta_2\rangle=1$, $\langle\alpha_2, \beta_1\rangle=1$, $\langle\alpha_2, \beta_2\rangle=1$, 
\\
\begin{tabular}{cc}$\alpha_1$+ & $\alpha_2$\\\tiny{ $A_1$ } & \tiny{  }\end{tabular} &\begin{tabular}{c} = \\~ \end{tabular} & \begin{tabular}{ccc}$\beta_1$+ & $\beta_2$+ & $\beta_3$\\\tiny{ $A_1$ } & \tiny{ $A_1$ } & \tiny{ $A_1$' }\end{tabular} & $D_4$ & $D_6$ & $\langle\alpha_1, \alpha_2\rangle=1$,  $\langle\alpha_1, \beta_1\rangle=1$, $\langle\alpha_1, \beta_2\rangle=1$, $\langle\alpha_1, \beta_3\rangle=1$, $\langle\alpha_2, \beta_1\rangle=1$, $\langle\alpha_2, \beta_2\rangle=1$, $\langle\alpha_2, \beta_3\rangle=1$, 
\\
\begin{tabular}{cc}$\alpha_1$+ & $\alpha_2$\\\tiny{ $A_1$ } & \tiny{  }\end{tabular} &\begin{tabular}{c} = \\~ \end{tabular} & \begin{tabular}{c}$\beta_1$\\\tiny{ $A_1$ }\end{tabular} & $A_2$ & $A_3$+$A_1$ & $\langle\alpha_1, \alpha_2\rangle=-1$,  $\langle\alpha_1, \beta_1\rangle=1$, $\langle\alpha_2, \beta_1\rangle=1$, 
\\
\begin{tabular}{cc}$\alpha_1$+ & $\alpha_2$\\\tiny{ $A_1$ } & \tiny{ $A_1$' }\end{tabular} &\begin{tabular}{c} = \\~ \end{tabular} & \begin{tabular}{cc}$\beta_1$+ & $\beta_2$\\\tiny{ $A_1$ } & \tiny{ $A_1$' }\end{tabular} & $A_3$ & $A_5$ &  $\langle\alpha_1, \beta_1\rangle=1$, $\langle\alpha_1, \beta_2\rangle=1$, $\langle\alpha_2, \beta_1\rangle=1$, $\langle\alpha_2, \beta_2\rangle=1$, 
\\
\begin{tabular}{cc}$\alpha_1$+ & $\alpha_2$\\\tiny{ $A_1$ } & \tiny{ $A_1$ }\end{tabular} &\begin{tabular}{c} = \\~ \end{tabular} & \begin{tabular}{ccc}$\beta_1$+ & $\beta_2$+ & $\beta_3$\\\tiny{ $A_1$+$A_1$' } & \tiny{ $A_1$' } & \tiny{ $A_1$ }\end{tabular} & $D_4$ & $E_6$ & $\langle\alpha_1, \alpha_2\rangle=1$,  $\langle\alpha_1, \beta_1\rangle=1$, $\langle\alpha_1, \beta_2\rangle=1$, $\langle\alpha_1, \beta_3\rangle=1$, $\langle\alpha_2, \beta_1\rangle=1$, $\langle\alpha_2, \beta_2\rangle=1$, $\langle\alpha_2, \beta_3\rangle=1$, 
\\
\begin{tabular}{cc}$\alpha_1$+ & $\alpha_2$\\\tiny{ $A_1$ } & \tiny{  }\end{tabular} &\begin{tabular}{c} = \\~ \end{tabular} & \begin{tabular}{cc}$\beta_1$+ & $\beta_2$\\\tiny{ $A_1$+$A_1$' } & \tiny{ $A_1$' }\end{tabular} & $A_3$ & $D_5$ &  $\langle\alpha_1, \beta_1\rangle=1$, $\langle\alpha_1, \beta_2\rangle=1$, $\langle\alpha_2, \beta_1\rangle=1$, $\langle\alpha_2, \beta_2\rangle=1$, 
\\
\begin{tabular}{cc}$\alpha_1$+ & $\alpha_2$\\\tiny{  } & \tiny{  }\end{tabular} &\begin{tabular}{c} = \\~ \end{tabular} & \begin{tabular}{ccc}$\beta_1$+ & $\beta_2$+ & $\beta_3$\\\tiny{ $A_1$+$A_1$' } & \tiny{ $A_1$' } & \tiny{ $A_1$ }\end{tabular} & $D_4$ & $E_6$ & $\langle\alpha_1, \alpha_2\rangle=1$,  $\langle\alpha_1, \beta_1\rangle=1$, $\langle\alpha_1, \beta_2\rangle=1$, $\langle\alpha_1, \beta_3\rangle=1$, $\langle\alpha_2, \beta_1\rangle=1$, $\langle\alpha_2, \beta_2\rangle=1$, $\langle\alpha_2, \beta_3\rangle=1$, 
\\
\begin{tabular}{cc}$\alpha_1$+ & $\alpha_2$\\\tiny{ $A_1$ } & \tiny{  }\end{tabular} &\begin{tabular}{c} = \\~ \end{tabular} & \begin{tabular}{cc}$\beta_1$+ & $\beta_2$\\\tiny{ $A_1$+$A_1$' } & \tiny{ $A_1$ }\end{tabular} & $A_3$ & $D_5$ &  $\langle\alpha_1, \beta_1\rangle=1$, $\langle\alpha_1, \beta_2\rangle=1$, $\langle\alpha_2, \beta_1\rangle=1$, $\langle\alpha_2, \beta_2\rangle=1$, 
\\
\begin{tabular}{c}2$\alpha_1$\\\tiny{  }\end{tabular} &\begin{tabular}{c} = \\~ \end{tabular} & \begin{tabular}{cccc}$\beta_1$+ & $\beta_2$+ & $\beta_3$+ & $\beta_4$\\\tiny{ $A_1$ } & \tiny{ $A_1$ } & \tiny{ $A_1$' } & \tiny{ $A_1$' }\end{tabular} & $D_4$ & $D_6$ &  $\langle\alpha_1, \beta_1\rangle=1$, $\langle\alpha_1, \beta_2\rangle=1$, $\langle\alpha_1, \beta_3\rangle=1$, $\langle\alpha_1, \beta_4\rangle=1$, 
\\
\begin{tabular}{cc}$\alpha_1$+ & $\alpha_2$\\\tiny{ $A_1$ } & \tiny{  }\end{tabular} &\begin{tabular}{c} = \\~ \end{tabular} & \begin{tabular}{ccc}$\beta_1$+ & $\beta_2$+ & $\beta_3$\\\tiny{ $A_1$ } & \tiny{ $A_1$' } & \tiny{ $A_1$' }\end{tabular} & $D_4$ & $D_6$ & $\langle\alpha_1, \alpha_2\rangle=1$,  $\langle\alpha_1, \beta_1\rangle=1$, $\langle\alpha_1, \beta_2\rangle=1$, $\langle\alpha_1, \beta_3\rangle=1$, $\langle\alpha_2, \beta_1\rangle=1$, $\langle\alpha_2, \beta_2\rangle=1$, $\langle\alpha_2, \beta_3\rangle=1$, 
\\
\begin{tabular}{cc}$\alpha_1$+ & $\alpha_2$\\\tiny{ $A_1$ } & \tiny{ $A_1$' }\end{tabular} &\begin{tabular}{c} = \\~ \end{tabular} & \begin{tabular}{c}$\beta_1$\\\tiny{ $A_1$+$A_1$' }\end{tabular} & $A_2$ & $A_4$ & $\langle\alpha_1, \alpha_2\rangle=-1$,  $\langle\alpha_1, \beta_1\rangle=1$, $\langle\alpha_2, \beta_1\rangle=1$, 
\\
\begin{tabular}{cc}$\alpha_1$+ & $\alpha_2$\\\tiny{ $A_1$+$A_1$' } & \tiny{  }\end{tabular} &\begin{tabular}{c} = \\~ \end{tabular} & \begin{tabular}{cc}$\beta_1$+ & $\beta_2$\\\tiny{ $A_1$ } & \tiny{ $A_1$' }\end{tabular} & $A_3$ & $A_5$ &  $\langle\alpha_1, \beta_1\rangle=1$, $\langle\alpha_1, \beta_2\rangle=1$, $\langle\alpha_2, \beta_1\rangle=1$, $\langle\alpha_2, \beta_2\rangle=1$, 
\\
\begin{tabular}{cc}$\alpha_1$+ & $\alpha_2$\\\tiny{  } & \tiny{  }\end{tabular} &\begin{tabular}{c} = \\~ \end{tabular} & \begin{tabular}{cc}$\beta_1$+ & $\beta_2$\\\tiny{ $A_1$+$A_1$' } & \tiny{ $A_1$+$A_1$' }\end{tabular} & $A_3$ & $D_4$ &  $\langle\alpha_1, \beta_1\rangle=1$, $\langle\alpha_1, \beta_2\rangle=1$, $\langle\alpha_2, \beta_1\rangle=1$, $\langle\alpha_2, \beta_2\rangle=1$, 
\\
\begin{tabular}{cc}$\alpha_1$+ & $\alpha_2$\\\tiny{ $A_1$+$A_1$' } & \tiny{  }\end{tabular} &\begin{tabular}{c} = \\~ \end{tabular} & \begin{tabular}{c}$\beta_1$\\\tiny{ $A_1$+$A_1$' }\end{tabular} & $A_2$ & $D_4$ & $\langle\alpha_1, \alpha_2\rangle=-1$,  $\langle\alpha_1, \beta_1\rangle=1$, $\langle\alpha_2, \beta_1\rangle=1$, 
\\
\begin{tabular}{cc}$\alpha_1$+ & $\alpha_2$\\\tiny{ $A_1$+$A_1$' } & \tiny{ $A_1$ }\end{tabular} &\begin{tabular}{c} = \\~ \end{tabular} & \begin{tabular}{cc}$\beta_1$+ & $\beta_2$\\\tiny{ $A_1$+$A_1$' } & \tiny{ $A_1$ }\end{tabular} & $A_3$ & $D_5$ &  $\langle\alpha_1, \beta_1\rangle=1$, $\langle\alpha_1, \beta_2\rangle=1$, $\langle\alpha_2, \beta_1\rangle=1$, $\langle\alpha_2, \beta_2\rangle=1$, 
\\
\begin{tabular}{cc}$\alpha_1$+ & $\alpha_2$\\\tiny{ $A_1$ } & \tiny{ $A_1$ }\end{tabular} &\begin{tabular}{c} = \\~ \end{tabular} & \begin{tabular}{cc}$\beta_1$+ & $\beta_2$\\\tiny{ $A_1$+$A_1$' } & \tiny{ $A_1$+$A_1$' }\end{tabular} & $A_3$ & $D_5$ &  $\langle\alpha_1, \beta_1\rangle=1$, $\langle\alpha_1, \beta_2\rangle=1$, $\langle\alpha_2, \beta_1\rangle=1$, $\langle\alpha_2, \beta_2\rangle=1$, 
\\
\begin{tabular}{cc}$\alpha_1$+ & $\alpha_2$\\\tiny{ $A_1$+$A_1$' } & \tiny{ $A_1$' }\end{tabular} &\begin{tabular}{c} = \\~ \end{tabular} & \begin{tabular}{cc}$\beta_1$+ & $\beta_2$\\\tiny{ $A_1$+$A_1$' } & \tiny{ $A_1$' }\end{tabular} & $A_3$ & $D_5$ &  $\langle\alpha_1, \beta_1\rangle=1$, $\langle\alpha_1, \beta_2\rangle=1$, $\langle\alpha_2, \beta_1\rangle=1$, $\langle\alpha_2, \beta_2\rangle=1$, 
\\
\begin{tabular}{c}2$\alpha_1$\\\tiny{ $A_1$ }\end{tabular} &\begin{tabular}{c} = \\~ \end{tabular} & \begin{tabular}{cccc}$\beta_1$+ & $\beta_2$+ & $\beta_3$+ & $\beta_4$\\\tiny{ $A_1$ } & \tiny{ $A_1$ } & \tiny{ $A_1$' } & \tiny{ $A_1$' }\end{tabular} & $D_4$ & $D_6$ &  $\langle\alpha_1, \beta_1\rangle=1$, $\langle\alpha_1, \beta_2\rangle=1$, $\langle\alpha_1, \beta_3\rangle=1$, $\langle\alpha_1, \beta_4\rangle=1$, 
\\
\begin{tabular}{cc}$\alpha_1$+ & $\alpha_2$\\\tiny{ $A_1$ } & \tiny{ $A_1$' }\end{tabular} &\begin{tabular}{c} = \\~ \end{tabular} & \begin{tabular}{cc}$\beta_1$+ & $\beta_2$\\\tiny{ $A_1$+$A_1$' } & \tiny{  }\end{tabular} & $A_3$ & $A_5$ &  $\langle\alpha_1, \beta_1\rangle=1$, $\langle\alpha_1, \beta_2\rangle=1$, $\langle\alpha_2, \beta_1\rangle=1$, $\langle\alpha_2, \beta_2\rangle=1$, 
\\
\begin{tabular}{cc}$\alpha_1$+ & $\alpha_2$\\\tiny{ $A_1$ } & \tiny{ $A_1$ }\end{tabular} &\begin{tabular}{c} = \\~ \end{tabular} & \begin{tabular}{cc}$\beta_1$+ & $\beta_2$\\\tiny{ $A_1$ } & \tiny{ $A_1$ }\end{tabular} & $A_3$ & $D_4$+$A_1$ &  $\langle\alpha_1, \beta_1\rangle=1$, $\langle\alpha_1, \beta_2\rangle=1$, $\langle\alpha_2, \beta_1\rangle=1$, $\langle\alpha_2, \beta_2\rangle=1$, 
\\
\begin{tabular}{cc}$\alpha_1$+ & $\alpha_2$\\\tiny{ $A_1$+$A_1$' } & \tiny{ $A_1$' }\end{tabular} &\begin{tabular}{c} = \\~ \end{tabular} & \begin{tabular}{ccc}$\beta_1$+ & $\beta_2$+ & $\beta_3$\\\tiny{ $A_1$ } & \tiny{ $A_1$ } & \tiny{ $A_1$' }\end{tabular} & $D_4$ & $D_6$ & $\langle\alpha_1, \alpha_2\rangle=1$,  $\langle\alpha_1, \beta_1\rangle=1$, $\langle\alpha_1, \beta_2\rangle=1$, $\langle\alpha_1, \beta_3\rangle=1$, $\langle\alpha_2, \beta_1\rangle=1$, $\langle\alpha_2, \beta_2\rangle=1$, $\langle\alpha_2, \beta_3\rangle=1$, 
\\
\begin{tabular}{cc}$\alpha_1$+ & $\alpha_2$\\\tiny{ $A_1$+$A_1$' } & \tiny{ $A_1$+$A_1$' }\end{tabular} &\begin{tabular}{c} = \\~ \end{tabular} & \begin{tabular}{ccc}$\beta_1$+ & $\beta_2$+ & $\beta_3$\\\tiny{ $A_1$+$A_1$' } & \tiny{ $A_1$' } & \tiny{ $A_1$ }\end{tabular} & $D_4$ & $E_6$ & $\langle\alpha_1, \alpha_2\rangle=1$,  $\langle\alpha_1, \beta_1\rangle=1$, $\langle\alpha_1, \beta_2\rangle=1$, $\langle\alpha_1, \beta_3\rangle=1$, $\langle\alpha_2, \beta_1\rangle=1$, $\langle\alpha_2, \beta_2\rangle=1$, $\langle\alpha_2, \beta_3\rangle=1$, 
\\
\begin{tabular}{cc}$\alpha_1$+ & $\alpha_2$\\\tiny{ $A_1$ } & \tiny{ $A_1$ }\end{tabular} &\begin{tabular}{c} = \\~ \end{tabular} & \begin{tabular}{ccc}$\beta_1$+ & $\beta_2$+ & $\beta_3$\\\tiny{ $A_1$+$A_1$' } & \tiny{ $A_1$+$A_1$' } & \tiny{  }\end{tabular} & $D_4$ & $D_6$ & $\langle\alpha_1, \alpha_2\rangle=1$,  $\langle\alpha_1, \beta_1\rangle=1$, $\langle\alpha_1, \beta_2\rangle=1$, $\langle\alpha_1, \beta_3\rangle=1$, $\langle\alpha_2, \beta_1\rangle=1$, $\langle\alpha_2, \beta_2\rangle=1$, $\langle\alpha_2, \beta_3\rangle=1$, 
\\
\begin{tabular}{c}2$\alpha_1$\\\tiny{ $A_1$ }\end{tabular} &\begin{tabular}{c} = \\~ \end{tabular} & \begin{tabular}{cccc}$\beta_1$+ & $\beta_2$+ & $\beta_3$+ & $\beta_4$\\\tiny{ $A_1$+$A_1$' } & \tiny{ $A_1$+$A_1$' } & \tiny{  } & \tiny{  }\end{tabular} & $D_4$ & $D_6$ &  $\langle\alpha_1, \beta_1\rangle=1$, $\langle\alpha_1, \beta_2\rangle=1$, $\langle\alpha_1, \beta_3\rangle=1$, $\langle\alpha_1, \beta_4\rangle=1$, 
\\
\begin{tabular}{cc}$\alpha_1$+ & $\alpha_2$\\\tiny{ $A_1$ } & \tiny{ $A_1$' }\end{tabular} &\begin{tabular}{c} = \\~ \end{tabular} & \begin{tabular}{ccc}$\beta_1$+ & $\beta_2$+ & $\beta_3$\\\tiny{ $A_1$+$A_1$' } & \tiny{  } & \tiny{  }\end{tabular} & $D_4$ & $D_6$ & $\langle\alpha_1, \alpha_2\rangle=1$,  $\langle\alpha_1, \beta_1\rangle=1$, $\langle\alpha_1, \beta_2\rangle=1$, $\langle\alpha_1, \beta_3\rangle=1$, $\langle\alpha_2, \beta_1\rangle=1$, $\langle\alpha_2, \beta_2\rangle=1$, $\langle\alpha_2, \beta_3\rangle=1$, 
\\
\begin{tabular}{c}2$\alpha_1$\\\tiny{ $A_1$+$A_1$' }\end{tabular} &\begin{tabular}{c} = \\~ \end{tabular} & \begin{tabular}{cccc}$\beta_1$+ & $\beta_2$+ & $\beta_3$+ & $\beta_4$\\\tiny{ $A_1$ } & \tiny{ $A_1$ } & \tiny{ $A_1$' } & \tiny{ $A_1$' }\end{tabular} & $D_4$ & $D_6$ &  $\langle\alpha_1, \beta_1\rangle=1$, $\langle\alpha_1, \beta_2\rangle=1$, $\langle\alpha_1, \beta_3\rangle=1$, $\langle\alpha_1, \beta_4\rangle=1$, 
\\
\begin{tabular}{cc}$\alpha_1$+ & $\alpha_2$\\\tiny{ $A_1$+$A_1$' } & \tiny{ $A_1$ }\end{tabular} &\begin{tabular}{c} = \\~ \end{tabular} & \begin{tabular}{ccc}$\beta_1$+ & $\beta_2$+ & $\beta_3$\\\tiny{ $A_1$ } & \tiny{ $A_1$' } & \tiny{ $A_1$' }\end{tabular} & $D_4$ & $D_6$ & $\langle\alpha_1, \alpha_2\rangle=1$,  $\langle\alpha_1, \beta_1\rangle=1$, $\langle\alpha_1, \beta_2\rangle=1$, $\langle\alpha_1, \beta_3\rangle=1$, $\langle\alpha_2, \beta_1\rangle=1$, $\langle\alpha_2, \beta_2\rangle=1$, $\langle\alpha_2, \beta_3\rangle=1$, 
\\
\end{longtable}

% % % % % % % % % % % % % % % % % % % % % % % % % % % %
% % % % % % % % % % % % % % % % % % % % % % % % % % % % % % % % %
\newpage
\begin{longtable}{rcl p{1cm}p{1cm}p{3cm} } \multicolumn{3}{c}{ Relation / linked $\mathfrak{k}$-components} &$\alpha_i$'s, $\beta_i$'s generate & adding $\Delta(\mathfrak{k})$ generates&Non-zero scalar products\endhead
\hline\multicolumn{5}{c}{$\mathfrak{k}$-semisimple type: $A_1$}\\
\hline\hline\begin{tabular}{cc}$\alpha_1$+ & $\alpha_2$\\\tiny{  } & \tiny{  }\end{tabular} &\begin{tabular}{c} = \\~ \end{tabular} & \begin{tabular}{cc}$\beta_1$+ & $\beta_2$\\\tiny{ $A_1$ } & \tiny{ $A_1$ }\end{tabular} & $A_3$ & $D_4$ &  $\langle\alpha_1, \beta_1\rangle=1$, $\langle\alpha_1, \beta_2\rangle=1$, $\langle\alpha_2, \beta_1\rangle=1$, $\langle\alpha_2, \beta_2\rangle=1$, 
\\
\begin{tabular}{cc}$\alpha_1$+ & $\alpha_2$\\\tiny{ $A_1$ } & \tiny{  }\end{tabular} &\begin{tabular}{c} = \\~ \end{tabular} & \begin{tabular}{c}$\beta_1$\\\tiny{ $A_1$ }\end{tabular} & $A_2$ & $A_3$ & $\langle\alpha_1, \alpha_2\rangle=-1$,  $\langle\alpha_1, \beta_1\rangle=1$, $\langle\alpha_2, \beta_1\rangle=1$, 
\\
\begin{tabular}{cc}$\alpha_1$+ & $\alpha_2$\\\tiny{ $A_1$ } & \tiny{ $A_1$ }\end{tabular} &\begin{tabular}{c} = \\~ \end{tabular} & \begin{tabular}{cc}$\beta_1$+ & $\beta_2$\\\tiny{ $A_1$ } & \tiny{ $A_1$ }\end{tabular} & $A_3$ & $D_4$ &  $\langle\alpha_1, \beta_1\rangle=1$, $\langle\alpha_1, \beta_2\rangle=1$, $\langle\alpha_2, \beta_1\rangle=1$, $\langle\alpha_2, \beta_2\rangle=1$, 
\\
\end{longtable}


} %end of \arxivVersion command

\end{appendix}

\end{document}
